% Generated mechanically from the frozen more-morphisms.tex target.
% Do not edit this generated file; edit the bound locale source instead.

% OK, start here.
%

\setcounter{chapter}{36}
\chapter{态射续篇}



\phantomsection
\label{more-morphisms-section-phantom}


\section{引言}
\label{more-morphisms-section-introduction}

\noindent
本章继续研究概形态射的性质。基本参考文献是 \cite{EGA}。






\section{加厚}
\label{more-morphisms-section-thickenings}

\noindent
下列术语未必完全标准，但使用起来很方便。

\begin{definition}
\label{more-morphisms-definition-thickening}
加厚。
\begin{enumerate}
\item 若 $X$ 是 $X'$ 的闭子概形，且二者的底层拓扑空间相同，
则称概形 $X'$ 是概形 $X$ 的一个{\it 加厚}。
\item 若 $X$ 是 $X'$ 的闭子概形，且定义 $X$ 的拟凝聚理想层
$\mathcal{I} \subset \mathcal{O}_{X'}$ 平方为零，
则称概形 $X'$ 是概形 $X$ 的一个{\it 一阶加厚}。
\item 若 $X$ 是 $X'$ 的闭子概形，且定义 $X$ 的拟凝聚理想层
$\mathcal{I} \subset \mathcal{O}_{X'}$ 幂零，即存在整数
$n \geq 0$ 使得 $\mathcal{I}^{n + 1} = 0$，
则称概形 $X'$ 是概形 $X$ 的一个{\it 有限阶加厚}。
\item 若 $X$ 是 $X$ 的闭子概形，且
$\mathcal{I}^{n + 1} = 0$，其中
$\mathcal{I} \subset \mathcal{O}_{X'}$ 是定义 $X$
的拟凝聚理想层，则称概形 $X'$ 是概形 $X$ 的一个
{\it $n$ 阶加厚}。
% P05-ERR-CH37-0001: frozen source says X is a closed subscheme of X.
\item 给定两个加厚 $X \subset X'$ 与 $Y \subset Y'$，
所谓{\it 加厚的态射}，是满足 $f'(X) \subset Y$
的态射 $f' : X' \to Y'$；换言之，$f'|_X$
经闭子概形 $Y$ 分解。在此情形中，置
$f = f'|_X : X \to Y$，并称
$(f, f') : (X \subset X') \to (Y \subset Y')$
为加厚的态射。
\item 设 $S$ 是概形。类似地定义{\it $S$ 上的加厚}及
{\it $S$ 上加厚的态射}。这意味着上述概形
$X, X', Y, Y'$ 都是 $S$-概形，而且态射
$X \to X'$、$Y \to Y'$ 与 $f' : X' \to Y'$ 都是 $S$-态射。
\end{enumerate}
\end{definition}

\noindent
设 $i_X : X \to X'$ 是加厚。核
$\mathcal{I} = \Ker(i_X^\sharp)$ 的任意局部截面都是局部幂零的。
这使我们能在一定程度上控制 $\mathcal{O}_{X'}$ 的结构，
但从技术上说，有限阶加厚 $X \subset X'$ 这一类要容易处理得多。
具体地，若 $X'$ 是 $n$ 阶加厚，则有滤过
% P05-ERR-CH37-0002: frozen source says “nth order thinking”.
$$
0 = \mathcal{I}^{n + 1} \subset
\mathcal{I}^n \subset
\mathcal{I}^{n - 1} \subset \ldots \subset
\mathcal{I} \subset \mathcal{O}_{X'}
$$
并且可见 $X'$ 具有闭子空间滤过
$$
X = X_1 \subset X_2 \subset \ldots \subset X_n \subset X_{n + 1} = X'
$$
，使每一对 $X_i \subset X_{i + 1}$ 都是 $S$ 上的一阶加厚。
利用简单的归纳论证，许多针对一阶加厚证明的结果，
都可改述为关于有限阶加厚的结果。

\medskip\noindent
一阶加厚可描述如下（见模，引理
\ref{modules-lemma-double-structure-gives-derivation}）。

\begin{lemma}
\label{more-morphisms-lemma-first-order-thickening}
设 $X$ 是基概形 $S$ 上的概形。考虑 $X$ 上层的短正合列
$$
0 \to \mathcal{I} \to \mathcal{A} \to \mathcal{O}_X \to 0
$$
，其中 $\mathcal{A}$ 是 $f^{-1}\mathcal{O}_S$-代数层，
$\mathcal{A} \to \mathcal{O}_X$ 是
$f^{-1}\mathcal{O}_S$-代数层的满射，而 $\mathcal{I}$ 是其核。
若
\begin{enumerate}
\item $\mathcal{I}$ 是 $\mathcal{A}$ 中平方为零的理想；并且
\item $\mathcal{I}$ 作为 $\mathcal{O}_X$-模是拟凝聚的，
\end{enumerate}
则 $X' = (X, \mathcal{A})$ 是概形，且 $X \to X'$ 是
$S$ 上的一阶加厚。此外，$S$ 上任意一阶加厚都具有这种形式。
\end{lemma}

\begin{proof}
显然 $X'$ 是局部赋环空间。设 $U = \Spec(B)$ 是 $X$
的仿射开集，并置 $A = \Gamma(U, \mathcal{A})$。
注意，由于 $H^1(U, \mathcal{I}) = 0$
（见概形上同调，引理
\ref{coherent-lemma-quasi-coherent-affine-cohomology-zero}），
映射 $A \to B$ 是满射。按假设，其核
$I = \mathcal{I}(U)$ 是环 $A$ 中平方为零的理想。
由概形，引理 \ref{schemes-lemma-morphism-into-affine}，
存在局部赋环空间的典范态射
$$
(U, \mathcal{A}|_U) \longrightarrow \Spec(A)
$$
，它来自映射 $B \to \Gamma(U, \mathcal{A})$。
% P05-ERR-CH37-0003: direction conflicts with the preceding A -> B surjection.
由于此态射嵌入交换图
$$
\xymatrix{
(U, \mathcal{O}_X|_U) \ar[d] \ar[r] & \Spec(B) \ar[d] \\
(U, \mathcal{A}|_U) \ar[r] & \Spec(A)
}
$$
，可见它在底层拓扑空间上是同胚。因此为证明它是同构，
只需检查它在局部环上诱导同构。
对与素理想 $\mathfrak p \subset A$ 对应的 $u \in U$，
得到短正合列的交换图
$$
\xymatrix{
0 \ar[r] &
I_{\mathfrak p} \ar[r] \ar[d] &
A_{\mathfrak p} \ar[r] \ar[d] &
B_{\mathfrak p} \ar[r] \ar[d] & 0 \\
0 \ar[r] &
\mathcal{I}_u \ar[r] &
\mathcal{A}_u \ar[r] &
\mathcal{O}_{X, u} \ar[r] & 0.
}
$$
由于 $\mathcal{I}$ 与 $\mathcal{O}_X$ 都是拟凝聚层，
左右两条竖直箭头均为同构，故中间的映射也是同构。
所以 $X' = (X, \mathcal{A})$ 的每一点都有仿射邻域，
而 $X'$ 如所需是概形。
\end{proof}

\begin{lemma}
\label{more-morphisms-lemma-thickening-affine-scheme}
\begin{slogan}
仿射性不受加厚影响
\end{slogan}
\begin{reference}
有限阶加厚的情形见 \cite[Proposition 5.1.9]{EGA1}。
\end{reference}
仿射概形的任意加厚都是仿射的。
\end{lemma}

\begin{proof}
这是极限，命题 \ref{limits-proposition-affine} 的特殊情形。
\end{proof}

\begin{proof}[有限阶加厚情形的证明]
设 $X \subset X'$ 是有限阶加厚，且 $X$ 仿射。
可以利用 Serre 判别法证明 $X'$ 仿射。更确切地，将使用概形上同调，引理
\ref{coherent-lemma-quasi-compact-h1-zero-covering}。
设 $\mathcal{F}$ 是拟凝聚 $\mathcal{O}_{X'}$-模。
只需证明 $H^1(X', \mathcal{F}) = 0$。
以 $i : X \to X'$ 表示给定的闭浸入，并记
$\mathcal{I} = \Ker(i^\sharp : \mathcal{O}_{X'} \to i_*\mathcal{O}_X)$.
由定义 \ref{more-morphisms-definition-thickening} 之后关于有限阶加厚的讨论，
存在 $n \geq 0$ 以及滤过
$$
0 = \mathcal{F}_{n + 1} \subset \mathcal{F}_n \subset
\mathcal{F}_{n - 1} \subset \ldots \subset
\mathcal{F}_0 = \mathcal{F}
$$
，其各项为拟凝聚子模，且
$\mathcal{F}_a/\mathcal{F}_{a + 1}$ 被 $\mathcal{I}$ 零化。
具体可取 $\mathcal{F}_a = \mathcal{I}^a\mathcal{F}$。
于是对某个拟凝聚 $\mathcal{O}_X$-模 $\mathcal{G}_a$，有
$\mathcal{F}_a/\mathcal{F}_{a + 1} = i_*\mathcal{G}_a$，
见态射，引理 \ref{morphisms-lemma-i-star-equivalence}。由此得到
$$
H^1(X', \mathcal{F}_a/\mathcal{F}_{a + 1}) =
H^1(X', i_*\mathcal{G}_a) = H^1(X, \mathcal{G}_a) = 0
$$
第二个等号来自概形上同调，引理
\ref{coherent-lemma-relative-affine-cohomology}，
最后一个等号来自概形上同调，引理
\ref{coherent-lemma-quasi-coherent-affine-cohomology-zero}。
因此 $\mathcal{F}$ 有一个有限滤过，其逐次商的第一同调均消失；
简单归纳即可推出 $H^1(X', \mathcal{F}) = 0$。
\end{proof}

\begin{lemma}
\label{more-morphisms-lemma-base-change-thickening}
设 $S \subset S'$ 是概形的加厚，给定态射 $X' \to S'$，
并置 $X = S \times_{S'} X'$。则
$(X \subset X') \to (S \subset S')$ 是加厚的态射。
若 $S \subset S'$ 是一阶（分别地，有限阶）加厚，
则 $X \subset X'$ 也是一阶（分别地，有限阶）加厚。
\end{lemma}

\begin{proof}
略。
\end{proof}

\begin{lemma}
\label{more-morphisms-lemma-composition-thickening}
\begin{slogan}
加厚的复合仍是加厚
\end{slogan}
若 $S \subset S'$ 与 $S' \subset S''$ 都是加厚，
则 $S \subset S''$ 也是加厚。
\end{lemma}

\begin{proof}
略。
\end{proof}

\begin{lemma}
\label{more-morphisms-lemma-descending-property-thickening}
成为加厚这一性质关于 fpqc 是局部的。
一阶加厚也同样如此。
\end{lemma}

\begin{proof}
该陈述的含义如下。设 $X \to X'$ 是概形态射，
$\{g_i : X'_i \to X'\}$ 是 fpqc 覆盖，
使得对所有 $i$，基变换 $X_i \to X'_i$ 都是加厚。
则 $X \to X'$ 是加厚。由于诸态射 $g_i$ 联合满射，
可知 $X \to X'$ 满射。由下降，引理
\ref{descent-lemma-descending-property-closed-immersion}，
$X \to X'$ 是闭浸入。因此 $X \to X'$ 是加厚。
略去一阶加厚情形的证明。
\end{proof}









\section{加厚的态射}
\label{more-morphisms-section-morphisms-thickenings}

\noindent
若 $(f, f') : (X \subset X') \to (Y \subset Y')$
是概形加厚的态射，则态射 $f$ 的性质往往由 $f'$ 继承。
这有若干种变体。

\begin{lemma}
\label{more-morphisms-lemma-thicken-property-morphisms}
设 $(f, f') : (X \subset X') \to (S \subset S')$
是加厚的态射。则
\begin{enumerate}
\item $f$ 是仿射态射，当且仅当 $f'$ 是仿射态射；
\item $f$ 是满射，当且仅当 $f'$ 是满射；
\item $f$ 拟紧，当且仅当 $f'$ 拟紧；
\item $f$ 泛闭，当且仅当 $f'$ 泛闭；
\item $f$ 整，当且仅当 $f'$ 整；
\item $f$（拟）分离，当且仅当 $f'$（拟）分离；
\item $f$ 泛单射，当且仅当 $f'$ 泛单射；
\item $f$ 泛开，当且仅当 $f'$ 泛开；
\item $f$ 拟仿射，当且仅当 $f'$ 拟仿射；以及
\item 此处添加更多内容。
\end{enumerate}
\end{lemma}

\begin{proof}
注意，$S \to S'$ 与 $X \to X'$ 都是泛同胚
（例如见态射，引理
\ref{morphisms-lemma-reduction-universal-homeomorphism}）。
这立即蕴含 (2)、(3)、(4)、(7)、(8)。
(1) 由引理 \ref{more-morphisms-lemma-thickening-affine-scheme} 得出；
该引理说明 $S$ 与 $S'$ 的仿射开集之间，以及
$X$ 与 $X'$ 的仿射开集之间，都存在一一对应。
(9) 由极限，引理 \ref{limits-lemma-thickening-quasi-affine}
以及刚才关于 $S$ 与 $S'$ 的仿射开集所作评注得出。
由态射，引理
\ref{morphisms-lemma-integral-universally-closed}，
(5) 由 (1) 与 (4) 得出。最后注意
$$
S \times_X S = S \times_{X'} S \to S \times_{X'} S' \to S' \times_{X'} S'
$$
是加厚（由引理 \ref{more-morphisms-lemma-base-change-thickening}，
两条箭头都是加厚）。
% P05-ERR-CH37-0004: frozen source fibre products reverse X/S roles.
所以把 (3) 与 (4) 应用于态射
$(S \subset S') \to (S \times_X S \to S' \times_{X'} S')$
，就得到 (6)。
\end{proof}

\begin{lemma}
\label{more-morphisms-lemma-thicken-property-relatively-ample}
设 $(f, f') : (X \subset X') \to (S \subset S')$
是加厚的态射。设 $\mathcal{L}'$ 是 $X'$ 上的可逆层，
并以 $\mathcal{L}$ 表示它到 $X$ 的限制。
则 $\mathcal{L}'$ 是 $f'$-丰沛的，当且仅当
$\mathcal{L}$ 是 $f$-丰沛的。
\end{lemma}

\begin{proof}
回忆，相对丰沛性是关于基底中每个仿射开集的条件，
见态射，定义 \ref{morphisms-definition-relatively-ample}。
由引理 \ref{more-morphisms-lemma-thickening-affine-scheme}，
$S$ 与 $S'$ 的仿射开集之间存在一一对应。
所以可假设 $S$ 与 $S'$ 都仿射，并归结为证明：
$\mathcal{L}'$ 丰沛，当且仅当 $\mathcal{L}$ 丰沛。
这就是极限，引理 \ref{limits-lemma-ample-on-reduction}。
\end{proof}

\begin{lemma}
\label{more-morphisms-lemma-thicken-property-morphisms-cartesian}
设 $(f, f') : (X \subset X') \to (S \subset S')$
是加厚的态射，并且 $X = S \times_{S'} X'$。
若 $S \subset S'$ 是有限阶加厚，则
\begin{enumerate}
\item $f$ 是闭浸入，当且仅当 $f'$ 是闭浸入；
\item $f$ 局部有限型，当且仅当 $f'$ 局部有限型；
\item $f$ 局部拟有限，当且仅当 $f'$ 局部拟有限；
\item $f$ 是相对维数为 $d$ 的局部有限型态射，当且仅当
$f'$ 是相对维数为 $d$ 的局部有限型态射；
\item $\Omega_{X/S} = 0$，当且仅当 $\Omega_{X'/S'} = 0$；
\item $f$ 非分歧，当且仅当 $f'$ 非分歧；
\item $f$ 固有，当且仅当 $f'$ 固有；
\item $f$ 有限，当且仅当 $f'$ 有限；
\item $f$ 是单态射，当且仅当 $f'$ 是单态射；
\item $f$ 是浸入，当且仅当 $f'$ 是浸入；以及
\item 此处添加更多内容。
\end{enumerate}
\end{lemma}

\begin{proof}
引理所列性质 $\mathcal{P}$ 在基变换下都稳定，
所以若 $f'$ 具有性质 $\mathcal{P}$，则 $f$ 也具有。
见概形，诸引理 \ref{schemes-lemma-base-change-immersion} 与
\ref{schemes-lemma-base-change-monomorphism}，以及态射，诸引理
\ref{morphisms-lemma-base-change-finite-type},
\ref{morphisms-lemma-base-change-quasi-finite},
\ref{morphisms-lemma-base-change-relative-dimension-d},
\ref{morphisms-lemma-base-change-differentials},
\ref{morphisms-lemma-base-change-unramified},
\ref{morphisms-lemma-base-change-proper}, 及
\ref{morphisms-lemma-base-change-finite}.

\medskip\noindent
因此每种情形中有内容的方向，都是从 $f$ 具有该性质出发，
推出 $f'$ 也具有。对加厚的阶数作归纳，可假设
$S \subset S'$ 是一阶加厚，见定义
\ref{more-morphisms-definition-thickening} 之后紧接着的讨论。

\medskip\noindent
大多数证明都要归结到仿射情形。设 $U' \subset S'$ 是仿射开集，
$V' \subset X'$ 是位于 $U'$ 上方的仿射开集。
令 $U' = \Spec(A')$，并以 $I \subset A'$ 表示定义闭子概形
$U' \cap S$ 的理想。设 $V' = \Spec(B')$。
于是 $V' \cap X = \Spec(B'/IB')$。置
$A = A'/I$、$B = B'/IB'$，得到交换图
$$
\xymatrix{
0 \ar[r] &
IB' \ar[r] &
B' \ar[r] &
B \ar[r] & 0 \\
0 \ar[r] &
IA' \ar[r] \ar[u] &
A' \ar[r] \ar[u] &
A \ar[r] \ar[u] & 0
}
$$
，其两行正合，且 $I^2 = 0$。

\medskip\noindent
把 (1) 译成代数：若 $A \to B$ 满射，则 $A' \to B'$ 满射。
这由 Nakayama 引理（代数，引理 \ref{algebra-lemma-NAK}）得出。

\medskip\noindent
把 (2) 译成代数：若 $A \to B$ 是有限型环映射，则
$A' \to B'$ 也是有限型环映射。这可将 Nakayama 引理
（代数，引理 \ref{algebra-lemma-NAK}）
应用于映射 $A'[x_1, \ldots, x_n] \to B'$ 得出，
其中 $A[x_1, \ldots, x_n] \to B$ 满射。

\medskip\noindent
(3) 的证明。它由 (2) 以及如下事实得出：局部有限型态射的拟有限性
可在纤维上检验，见态射，引理
\ref{morphisms-lemma-quasi-finite-at-point-characterize}。

\medskip\noindent
(4) 的证明。它由 (2) 以及如下事实得出：附加性质“相对维数为
$d$”可在纤维上检验（按定义如此，见态射，定义
\ref{morphisms-definition-relative-dimension-d}。
% P05-ERR-CH37-0005: frozen source is missing the closing parenthesis.

\medskip\noindent
把 (5) 译成代数：若 $\Omega_{B/A} = 0$，则
$\Omega_{B'/A'} = 0$。由代数，引理
\ref{algebra-lemma-differentials-base-change}，有
$0 = \Omega_{B/A} = \Omega_{B'/A'}/I\Omega_{B'/A'}$。
因此由 Nakayama 引理（代数，引理 \ref{algebra-lemma-NAK}），
$\Omega_{B'/A'} = 0$。

\medskip\noindent
把 (6) 译成代数：若 $A \to B$ 是非分歧映射，则
$A' \to B'$ 非分歧。由于 $A \to B$ 有限型，
由上面的 (2)，$A' \to B'$ 有限型。
又因 $A \to B$ 非分歧，有 $\Omega_{B/A} = 0$；
由 (5)，有 $\Omega_{B'/A'} = 0$。所以 $A' \to B'$ 非分歧。

\medskip\noindent
(7) 的证明。把 (2)、引理
\ref{more-morphisms-lemma-thicken-property-morphisms} 的结果，
以及“固有等价于拟紧 $+$ 分离 $+$ 局部有限型 $+$ 泛闭”这一事实结合即可。

\medskip\noindent
(8) 的证明。把 (2) 与引理
\ref{more-morphisms-lemma-thicken-property-morphisms} 的结果结合，
并使用“有限等价于整 $+$ 局部有限型”这一事实即可
（态射，引理 \ref{morphisms-lemma-finite-integral}）。

\medskip\noindent
(9) 的证明。由于 $f$ 是单态射，有 $X = X \times_S X$。
可把迄今证明的结果应用于加厚的态射
$(X \subset X') \to (X \times_S X \subset X' \times_{S'} X')$.
。由 (1)，得到 $X' \to X' \times_{S'} X'$ 是闭浸入。
事实上它是一阶加厚，因为定义闭浸入
$X' \to X' \times_{S'} X'$ 的理想，包含在理想
$\mathcal{I} \subset \mathcal{O}_{S'}$ 的拉回中；
该理想定义 $S'$ 中的 $S$。确实，
$X = X \times_S X = (X' \times_{S'} X') \times_{S'} S$
包含于 $X'$。所以由态射，引理
\ref{morphisms-lemma-differentials-diagonal}，
只需证明 $\Omega_{X'/S'} = 0$；
这由 (5) 以及 $X/S$ 的相应陈述得出。

\medskip\noindent
(10) 的证明。若 $f : X \to S$ 是浸入，则它分解为
$X \to U \to S$，其中 $U \to S$ 是开浸入，
$X \to U$ 是闭浸入。令 $U' \subset S'$ 为底层拓扑空间
与 $U$ 相同的开子概形。于是 $X' \to S'$ 经 $U'$ 分解，
由 (1)，$X' \to U'$ 是闭浸入。证明完毕。
\end{proof}

\noindent
下一个引理是前一引理的变体。这里不再假设所涉加厚为有限阶
（该假设使我们能把局部有限型这一性质从 $f$ 传递给 $f'$），
而是直接假定 $f$ 与 $f'$ 都局部有限型。

\begin{lemma}
\label{more-morphisms-lemma-properties-that-extend-over-thickenings}
设 $(f, f') : (X \subset X') \to (Y \to Y')$ 是加厚的态射。
% P05-ERR-CH37-0006: frozen source uses Y -> Y' instead of Y subset Y'.
假设 $f$ 与 $f'$ 局部有限型，且 $X = Y \times_{Y'} X'$。则
\begin{enumerate}
\item $f$ 局部拟有限，当且仅当 $f'$ 局部拟有限；
\item $f$ 有限，当且仅当 $f'$ 有限；
\item $f$ 是闭浸入，当且仅当 $f'$ 是闭浸入；
\item $\Omega_{X/Y} = 0$，当且仅当 $\Omega_{X'/Y'} = 0$；
\item $f$ 非分歧，当且仅当 $f'$ 非分歧；
\item $f$ 是单态射，当且仅当 $f'$ 是单态射；
\item $f$ 是浸入，当且仅当 $f'$ 是浸入；
\item $f$ 固有，当且仅当 $f'$ 固有；以及
\item 此处添加更多内容。
\end{enumerate}
\end{lemma}

\begin{proof}
引理所列性质 $\mathcal{P}$ 在基变换下都稳定，
所以若 $f'$ 具有性质 $\mathcal{P}$，则 $f$ 也具有。
见概形，诸引理 \ref{schemes-lemma-base-change-immersion} 与
\ref{schemes-lemma-base-change-monomorphism}，以及态射，诸引理
\ref{morphisms-lemma-base-change-quasi-finite},
\ref{morphisms-lemma-base-change-relative-dimension-d},
\ref{morphisms-lemma-base-change-differentials},
\ref{morphisms-lemma-base-change-unramified},
\ref{morphisms-lemma-base-change-proper}, 及
\ref{morphisms-lemma-base-change-finite}.
因此每种情形中只需证明：若 $f$ 具有所需性质，
则 $f'$ 也具有。

\medskip\noindent
一个态射局部拟有限，当且仅当它局部有限型且概形论纤维为离散空间，
见态射，引理 \ref{morphisms-lemma-locally-quasi-finite-fibres}。
由于传到加厚不会改变底层拓扑空间，可知若（且仅若）
$f$ 局部拟有限，则 $f'$ 局部拟有限。这证明了 (1)。

\medskip\noindent
(2) 由引理 \ref{more-morphisms-lemma-thicken-property-morphisms} 的 (5)
以及下列事实得出：有限态射恰好是局部有限型的整态射
（态射，引理 \ref{morphisms-lemma-finite-integral}）。

\medskip\noindent
(3)。这立即由态射，引理
\ref{morphisms-lemma-check-closed-infinitesimally} 得出。

\medskip\noindent
(4) 由下列代数陈述得出：设 $A$ 是环，
$I \subset A$ 是局部幂零理想，$B$ 是有限型 $A$-代数。
若 $\Omega_{(B/IB)/(A/I)} = 0$，则 $\Omega_{B/A} = 0$。
具体地，该假设意味着 $I\Omega_{B/A} = 0$，见代数，引理
\ref{algebra-lemma-differentials-base-change}。
另一方面，$\Omega_{B/A}$ 是有限 $B$-模，见代数，引理
\ref{algebra-lemma-differentials-finitely-generated}。
所以 $\Omega_{B/A}$ 的消失性由 Nakayama 引理
（代数，引理 \ref{algebra-lemma-NAK}）以及
$IB$ 包含于 $B$ 的 Jacobson 根这一事实得出。

\medskip\noindent
(5) 立即由 (4) 与态射，引理
\ref{morphisms-lemma-unramified-omega-zero} 得出。

\medskip\noindent
(6) 的证明。由于 $f$ 是单态射，有 $X = X \times_Y X$。
可把迄今证明的结果应用于加厚的态射
$(X \subset X') \to (X \times_Y X \subset X' \times_{Y'} X')$.
。由 (3)，得到
$\Delta_{X'/Y'} : X' \to X' \times_{Y'} X'$ 是闭浸入。
事实上，$\Delta_{X'/Y'}$ 在底层集合上是双射，
故 $\Delta_{X'/Y'}$ 是加厚。另一方面，由态射，引理
\ref{morphisms-lemma-diagonal-morphism-finite-type}，
$\Delta_{X'/Y'}$ 局部有限表示。换言之，
$\Delta_{X'/Y'}(X')$ 由有限型拟凝聚理想层
$\mathcal{J} \subset \mathcal{O}_{X' \times_{Y'} X'}$ 截出。
由 (5)，$\Omega_{X'/Y'} = 0$；再由态射，引理
\ref{morphisms-lemma-differentials-diagonal}，
$X' \to X' \times_{Y'} X'$ 的余法层为零。
换言之，$\mathcal{J}/\mathcal{J}^2 = 0$。
例如由代数，引理
\ref{algebra-lemma-ideal-is-squared-union-connected}，
这蕴含 $\Delta_{X'/Y'}$ 是同构。

\medskip\noindent
(7) 的证明。若 $f : X \to Y$ 是浸入，则它分解为
$X \to V \to Y$，其中 $V \to Y$ 是开浸入，
$X \to V$ 是闭浸入。令 $V' \subset Y'$ 为底层拓扑空间
与 $V$ 相同的开子概形。于是 $X' \to V'$ 经 $V'$ 分解，
% P05-ERR-CH37-0007: frozen source's factorization sentence is tautological.
并由 (3) 得到 $X' \to V'$ 是闭浸入。

\medskip\noindent
(8) 由引理 \ref{more-morphisms-lemma-thicken-property-morphisms}
以及固有态射的定义得出：它是拟紧、泛闭、分离且局部有限型的态射。
\end{proof}








\section{加厚的 Picard 群}
\label{more-morphisms-section-picard-group-thickening}

\noindent
这里给出一些关于加厚的 Picard 群的材料。

\begin{lemma}
\label{more-morphisms-lemma-picard-group-first-order-thickening}
设 $X \subset X'$ 是理想层为 $\mathcal{I}$ 的一阶加厚。则有阿贝尔群的典范正合列
$$
\xymatrix{
0 \ar[r] &
H^0(X, \mathcal{I}) \ar[r] &
H^0(X', \mathcal{O}_{X'}^*) \ar[r] &
H^0(X, \mathcal{O}^*_X) \ar `r[d] `d[l] `l[llld] `d[dll] [dll] \\
& H^1(X, \mathcal{I}) \ar[r] &
\Pic(X') \ar[r] &
\Pic(X) \ar `r[d] `d[l] `l[llld] `d[dll] [dll] \\
& H^2(X, \mathcal{I}) \ar[r] & \ldots \ar[r] & \ldots
}
$$
。
\end{lemma}

\begin{proof}
这是与阿贝尔群层的短正合列
$$
0 \to \mathcal{I} \to \mathcal{O}_{X'}^* \to \mathcal{O}_X^* \to 0
$$
相伴的长上同调正合列，其中第一个映射把 $\mathcal{I}$ 的局部截面 $f$
送到 $\mathcal{O}_{X'}$ 的可逆截面 $1 + f$。
这里还使用了环化空间的 Picard 群与其可逆函数层的第一上同调群之间的等同；
见 上同调章引理 \ref{cohomology-lemma-h1-invertible}。
\end{proof}

\begin{lemma}
\label{more-morphisms-lemma-torsion-pic-thickening}
设 $X \subset X'$ 是一个加厚，且整数 $n$ 在 $\mathcal{O}_X$ 中可逆。
则映射 $\Pic(X')[n] \to \Pic(X)[n]$ 是双射。
\end{lemma}

\begin{proof}[有限阶加厚情形的证明]
由定义 \ref{more-morphisms-definition-thickening} 后所述的一般原理，问题可归结为一阶加厚的情形。
此时由引理 \ref{more-morphisms-lemma-picard-group-first-order-thickening}，只需证明
$H^1(X, \mathcal{I})[n]$、$H^1(X, \mathcal{I})/n$ 和
$H^2(X, \mathcal{I})[n]$ 均为零。
这是因为 $\mathcal{I}$ 作为 $\mathcal{O}_X$-模，其上的乘 $n$ 映射是同构。
\end{proof}

\begin{proof}[一般情形的证明]
设 $\mathcal{I} \subset \mathcal{O}_{X'}$ 是截出 $X$ 的拟凝聚理想层。
则有阿贝尔群层的短正合列
$$
0 \to (1 + \mathcal{I})^* \to \mathcal{O}_{X'}^* \to \mathcal{O}_X^* \to 0
$$
从而得到如引理 \ref{more-morphisms-lemma-picard-group-first-order-thickening} 所述的长上同调正合列，
只是把 $H^i(X, \mathcal{I})$ 换成 $H^i(X, (1 + \mathcal{I})^*)$。
因此只需证明取 $n$ 次幂给出同构
$(1 + \mathcal{I})^* \to (1 + \mathcal{I})^*$。
在仿射开集上取截面后，这由
代数章引理 \ref{algebra-lemma-lift-nth-roots} 得出。
\end{proof}





\section{无穷小邻域}
\label{more-morphisms-section-first-order-infinitesimal-neighbourhood}

\noindent
有限阶加厚有如下自然构造。设 $i : Z \to X$ 是概形的浸入。
选择开子概形 $U \subset X$，使得 $i$ 将 $Z$ 等同于闭子概形
$Z \subset U$。设 $\mathcal{I} \subset \mathcal{O}_U$
是在 $U$ 中定义 $Z$ 的拟凝聚理想层。对 $n \geq 1$，
可以考虑由拟凝聚理想层 $\mathcal{I}^{n + 1}$ 定义的闭子概形
$Z_n \subset U$。

\begin{definition}
\label{more-morphisms-definition-first-order-infinitesimal-neighbourhood}
设 $i : Z \to X$ 是概形的浸入。
\begin{enumerate}
\item $Z$ 在 $X$ 中的{\it 一阶无穷小邻域}是上述 $X$ 上的一阶加厚
$Z \subset Z_1$。
\item $Z$ 在 $X$ 中的{\it $n$ 阶无穷小邻域}是上述 $X$ 上的
$n$ 阶加厚 $Z \subset Z_n$。
\end{enumerate}
\end{definition}

\noindent
这些加厚具有如下泛性质（它也消除了对上述构造依赖于开集 $U$ 之选取的疑虑）。

\begin{lemma}
\label{more-morphisms-lemma-first-order-infinitesimal-neighbourhood}
设 $i : Z \to X$ 是概形的浸入。
\begin{enumerate}
\item $Z$ 在 $X$ 中的一阶无穷小邻域 $Z'$ 具有如下泛性质：给定任意交换图
$$
\xymatrix{
Z \ar[d]_i & T \ar[l]^a \ar[d] \\
X & T' \ar[l]_b
}
$$
其中 $T \subset T'$ 是 $X$ 上的一阶加厚，则存在唯一的 $X$ 上加厚的态射
$(a', a) : (T \subset T') \to (Z \subset Z')$。
\item 对 $n \geq 1$，$Z$ 在 $X$ 中的 $n$ 阶无穷小邻域 $Z_n$
具有如下泛性质：给定任意交换图
$$
\xymatrix{
Z \ar[d]_i & T \ar[l]^a \ar[d] \\
X & T' \ar[l]_b
}
$$
其中 $T \subset T'$ 是 $X$ 上的 $n$ 阶加厚，则存在唯一的 $X$ 上加厚的态射
$(a', a) : (T \subset T') \to (Z \subset Z_n)$。
\end{enumerate}
\end{lemma}

\begin{proof}
只证明 (1)。设 $U \subset X$ 是构造 $Z'$ 时所用的开集；也就是说，
$Z$ 被等同于 $U$ 中由拟凝聚理想层 $\mathcal{I}$ 截出的闭子概形。
由于 $|T| = |T'|$，可知 $b(T') \subset U$，因而可把 $b$ 看作到 $U$ 的态射。
设 $\mathcal{J} \subset \mathcal{O}_{T'}$ 是截出 $T$ 的理想。
由上图中的 $b(T) \subset Z$，可知
$b^\sharp(b^{-1}\mathcal{I}) \subset \mathcal{J}$。
由于 $T'$ 是 $T$ 的一阶加厚，有 $\mathcal{J}^2 = 0$，故
$b^\sharp(b^{-1}(\mathcal{I}^2)) = 0$。由
概形章引理 \ref{schemes-lemma-characterize-closed-subspace}，
这表明 $b$ 经由 $Z'$ 分解。以 $a' : T' \to Z'$ 表示此分解，结论即明。
\end{proof}

\begin{lemma}
\label{more-morphisms-lemma-infinitesimal-neighbourhood-conormal}
设 $i : Z \to X$ 是概形的浸入，并设 $Z \subset Z'$ 是 $Z$ 在 $X$
中的一阶无穷小邻域。则图
$$
\xymatrix{
Z \ar[r] \ar[d] & Z' \ar[d] \\
Z \ar[r] & X
}
$$
由 态射章引理 \ref{morphisms-lemma-conormal-functorial}
诱导余法层的映射 $\mathcal{C}_{Z/X} \to \mathcal{C}_{Z/Z'}$。
此映射是同构。
\end{lemma}

\begin{proof}
这由上述 $Z'$ 的构造立即可见。
\end{proof}
















\section{形式非分歧态射}
\label{more-morphisms-section-formally-unramified}

\noindent
回顾：若对每个实线部分交换的图
$$
\xymatrix{
A \ar[r] \ar@{-->}[rd] & B/I \\
R \ar[r] \ar[u] & B \ar[u]
}
$$
（其中 $I \subset B$ 是平方为零的理想），至多存在一条使该图交换的虚线箭头，
则称环映射 $R \to A$ 是{\it 形式非分歧的}；见
代数章定义 \ref{algebra-definition-formally-unramified}。
这引出概形态射的如下类似概念。

\begin{definition}
\label{more-morphisms-definition-formally-unramified}
设 $f : X \to S$ 是概形态射。若给定任意实线部分交换的图
$$
\xymatrix{
X \ar[d]_f & T \ar[d]^i \ar[l] \\
S & T' \ar[l] \ar@{-->}[lu]
}
$$
其中 $T \subset T'$ 是 $S$ 上仿射概形的一阶加厚时，至多存在一条使该图交换的
虚线箭头，则称 $f$ 是{\it 形式非分歧的}。
\end{definition}

\noindent
先证明一些形式性引理，也就是只需画出相应交换图即可证明的引理。

\begin{lemma}
\label{more-morphisms-lemma-formally-unramified-not-affine}
若 $f : X \to S$ 是形式非分歧态射，则给定任意实线部分交换的图
$$
\xymatrix{
X \ar[d]_f & T \ar[d]^i \ar[l] \\
S & T' \ar[l] \ar@{-->}[lu]
}
$$
其中 $T \subset T'$ 是 $S$ 上概形的一阶加厚，至多存在一条使该图交换的虚线箭头。
换言之，在定义 \ref{more-morphisms-definition-formally-unramified} 中可以去掉 $T$ 为仿射的条件。
\end{lemma}

\begin{proof}
这是因为一个态射由它在各仿射开集上的限制所决定。
\end{proof}

\begin{lemma}
\label{more-morphisms-lemma-composition-formally-unramified}
形式非分歧态射的复合仍是形式非分歧的。
\end{lemma}

\begin{proof}
这是形式性的。
\end{proof}

\begin{lemma}
\label{more-morphisms-lemma-base-change-formally-unramified}
形式非分歧态射的基变换仍是形式非分歧的。
\end{lemma}

\begin{proof}
这是形式性的。
\end{proof}

\begin{lemma}
\label{more-morphisms-lemma-formally-unramified-on-opens}
设 $f : X \to S$ 是概形态射，并设 $U \subset X$ 与 $V \subset S$
是满足 $f(U) \subset V$ 的开集。若 $f$ 形式非分歧，则
$f|_U : U \to V$ 也形式非分歧。
\end{lemma}

\begin{proof}
考虑定义 \ref{more-morphisms-definition-formally-unramified} 中那样的实线部分交换图
$$
\xymatrix{
U \ar[d]_{f|_U} & T \ar[d]^i \ar[l]^a \\
V & T' \ar[l] \ar@{-->}[lu]
}
$$
。
% P05-ERR-CH37-0008: Frozen source says ``formally ramified'' here; retained literally pending canon review.
若 $f$ 形式分歧，则至多存在一个满足 $a'|_T = a$ 的 $S$-态射
$a' : T' \to X$。因此显然至多存在一个这样的到 $U$ 的态射。
\end{proof}

\begin{lemma}
\label{more-morphisms-lemma-affine-formally-unramified}
设 $f : X \to S$ 是概形态射，并假设 $X$ 与 $S$ 均为仿射。
则 $f$ 形式非分歧，当且仅当
$\mathcal{O}_S(S) \to \mathcal{O}_X(X)$ 是形式非分歧环映射。
\end{lemma}

\begin{proof}
由环范畴与仿射概形范畴之间的等价，这立即从定义
（定义 \ref{more-morphisms-definition-formally-unramified} 与
代数章定义 \ref{algebra-definition-formally-unramified}）得出；见
概形章引理 \ref{schemes-lemma-category-affine-schemes}。
\end{proof}

\noindent
下面用微分层给出一个刻画。

\begin{lemma}
\label{more-morphisms-lemma-formally-unramified-differentials}
设 $f : X \to S$ 是概形态射。则 $f$ 形式非分歧，当且仅当
$\Omega_{X/S} = 0$。
\end{lemma}

\begin{proof}
回顾 态射章引理 \ref{morphisms-lemma-differentials-affine}
证明中的部分论证。取开集 $W \subset X \times_S X$，使对角态射
$\Delta : X \to X \times_S X$ 诱导到 $W$ 的闭浸入。
设 $\mathcal{J} \subset \mathcal{O}_W$ 是此闭浸入的理想层，并设
$X' \subset W$ 是由拟凝聚理想层 $\mathcal{J}^2$ 定义的闭子概形。
考虑由两个投影 $X \times_S X \to X$ 诱导的两个态射
$p_1, p_2 : X' \to X$。注意，$p_1$ 与 $p_2$ 同
$\Delta : X \to X'$ 复合后相等，而且 $X \to X'$ 是由平方为零的理想
定义的闭浸入。此外有短正合列
$$
0 \to \mathcal{J}/\mathcal{J}^2 \to \mathcal{O}_{X'} \to \mathcal{O}_X \to 0
$$
且 $\Omega_{X/S} = \mathcal{J}/\mathcal{J}^2$。再者，
$\mathcal{J}/\mathcal{J}^2$ 由局部截面
$p_1^\sharp(f) - p_2^\sharp(f)$ 生成，其中 $f$ 遍历
$\mathcal{O}_X$ 的局部截面。

\medskip\noindent
假设 $f : X \to S$ 形式非分歧。按假设，这意味着 $p_1 = p_2$
限制到任意仿射开集 $T' \subset X'$ 后相等，故 $p_1 = p_2$。
由上可得 $\Omega_{X/S} = \mathcal{J}/\mathcal{J}^2 = 0$。

\medskip\noindent
反之，假设 $\Omega_{X/S} = 0$，于是 $X' = X$。任取定义
\ref{more-morphisms-definition-formally-unramified} 中图的两条虚线箭头
$f'_1, f'_2 : T' \to X$。它们给出态射
$(f'_1, f'_2) : T' \to X \times_S X$。
由于 $f'_1|_T = f'_2|_T$ 且 $|T| =|T'|$，$T'$ 在
$(f'_1, f'_2)$ 下的像包含于上面选取的开集 $W$ 中。
又因 $(f'_1, f'_2)(T) \subset \Delta(X)$，且 $T$ 在 $T'$ 中由平方为零的
理想定义，可知 $(f'_1, f'_2)$ 经由 $X'$ 分解。因 $X' = X$，
遂得所需的 $f_1' = f'_2$。
\end{proof}

\begin{lemma}
\label{more-morphisms-lemma-formally-unramified}
设 $f : X \to Y$ 是概形态射。以下条件等价：
\begin{enumerate}
\item $f$ 形式非分歧；
\item 对每个 $x \in X$，存在满足 $f(U) \subset V$ 的开集
$x \in U \subset X$ 与 $f(x) \in V \subset Y$，使得
$f|_U : U \to V$ 形式非分歧；
\item 对每一对满足 $f(U) \subset V$ 的仿射开集 $U \subset X$ 与
$V \subset Y$，环映射 $\mathcal{O}_Y(V) \to \mathcal{O}_X(U)$
形式非分歧；
\item 存在仿射开覆盖 $Y = \bigcup V_j$，且对每个 $j$ 存在仿射开覆盖
$f^{-1}(V_j) = \bigcup U_{ji}$，使得对所有 $j$ 与 $i$，环映射
$\mathcal{O}_Y(V) \to \mathcal{O}_X(U)$ 形式非分歧。
\end{enumerate}
\end{lemma}

\begin{proof}
结合引理 \ref{more-morphisms-lemma-formally-unramified-on-opens} 与
\ref{more-morphisms-lemma-affine-formally-unramified} 可得 (1) $\Rightarrow$ (3)。
(3) $\Rightarrow$ (4) 是立即的。由引理
\ref{more-morphisms-lemma-affine-formally-unramified} 得 (4) $\Rightarrow$ (2)。
若 (2) 成立，则由引理 \ref{more-morphisms-lemma-formally-unramified-differentials}
（这里还使用 态射章引理
\ref{morphisms-lemma-differentials-restrict-open}），
$\Omega_{X/S}$ 在每一点的某个邻域上消失；于是再由引理
\ref{more-morphisms-lemma-formally-unramified-differentials}，$f$ 形式非分歧。
\end{proof}

\begin{lemma}
\label{more-morphisms-lemma-unramified-formally-unramified}
\begin{slogan}
非分歧态射恰为局部有限型的形式非分歧态射。
\end{slogan}
设 $f : X \to S$ 是概形态射。以下条件等价：
\begin{enumerate}
\item 态射 $f$ 非分歧（分别地，G-非分歧）；
\item 态射 $f$ 局部有限型（分别地，局部有限表示）且形式非分歧。
\end{enumerate}
\end{lemma}

\begin{proof}
使用引理 \ref{more-morphisms-lemma-formally-unramified-differentials} 与
态射章引理 \ref{morphisms-lemma-unramified-omega-zero}。
\end{proof}









\section{泛一阶加厚}
\label{more-morphisms-section-universal-thickening}

\noindent
设 $h : Z \to X$ 是概形态射。$Z$ 在 $X$ 上的一个{\it 泛一阶加厚}
是 $X$ 上的一阶加厚 $Z \subset Z'$，使得给定 $X$ 上的任意一阶加厚
$T \subset T'$ 和实线部分交换图
$$
\xymatrix{
& Z \ar[ld] & & T \ar[rd] \ar[ll]^a \\
Z' \ar[rrd] & & & & T' \ar@{..>}[llll]_{a'} \ar[lld]^b \\
 & & X
}
$$
时，存在唯一的虚线箭头使该图交换。注意，在这种情形下
$(a, a') : (T \subset T') \to (Z \subset Z')$ 是 $X$ 上加厚的态射。
因此若泛一阶加厚存在，则它在唯一同构的意义下唯一。
一般而言泛一阶加厚未必存在，但当 $h$ 形式非分歧时它确实存在。

\begin{lemma}
\label{more-morphisms-lemma-universal-thickening}
设 $h : Z \to X$ 是形式非分歧的概形态射。则存在 $Z$ 在 $X$ 上的
泛一阶加厚 $Z \subset Z'$。
\end{lemma}

\begin{proof}
在本证明中，若 $Z \subset Z'$ 满足引理中的条件，就称它为 $Z$ 在 $X$
上的泛一阶加厚。我们将通过粘合构造泛一阶加厚
$Z \subset Z'$，先从仿射情形开始；该情形即
代数章引理 \ref{algebra-lemma-universal-thickening}。
首先作一些一般性说明。

\medskip\noindent
若 $Z$ 在 $X$ 上的泛一阶加厚存在，则它在唯一同构的意义下唯一。
再设 $V \subset Z$ 与 $U \subset X$ 是满足 $h(V) \subset U$ 的开子概形，
并设 $Z \subset Z'$ 是 $Z$ 在 $X$ 上的泛一阶加厚。令
$V' \subset Z'$ 为满足 $V = Z \cap V'$ 的开子概形。我们断言
$V \subset V'$ 是 $V$ 在 $U$ 上的泛一阶加厚。事实上，设给定任意图
$$
\xymatrix{
V \ar[d]_h & T \ar[l]^a \ar[d] \\
U & T' \ar[l]_b
}
$$
其中 $T \subset T'$ 是 $U$ 上的一阶加厚。由 $Z'$ 的泛性质，得到
$(a, a') : (T \subset T') \to (Z \subset Z')$。
但由于底拓扑空间相等 $|T| = |T'|$，可知 $a'(T') \subset V'$。
因此可把 $(a, a')$ 看作 $U$ 上加厚的态射
$(a, a') : (T \subset T') \to (V \subset V')$，唯一性也很清楚。
完全类似地可证：若 $h(Z) \subset U$，且 $Z \subset Z'$ 是 $U$ 上的
泛一阶加厚，则 $Z \subset Z'$ 也是 $X$ 上的泛一阶加厚。

\medskip\noindent
在粘合各仿射片之前，先证明 $Z$ 与 $X$ 均为仿射时引理成立。
写 $X = \Spec(R)$、$Z = \Spec(S)$。由
代数章引理 \ref{algebra-lemma-universal-thickening}，
存在 $X$ 上的一阶加厚 $Z \subset Z'$，它对如下图具有引理中的泛性质：
$$
\xymatrix{
Z \ar[d]_h & T \ar[l]^a \ar[d] \\
X & T' \ar[l]_b
}
$$
其中 $T, T'$ 为仿射。对一般的图，可选择仿射开覆盖
$T' = \bigcup T'_i$，并得到 $X$ 上的态射 $a'_i : T'_i \to Z'$，
满足 $a'_i|_{T_i} = a|_{T_i}$。由唯一性，$a'_i$ 与 $a'_j$
在 $T'_i \cap T'_j$ 的任意仿射开集上相等。因此诸态射 $a'_i$
粘合成所需的 $X$ 上整体态射 $a' : T' \to Z'$。
故当 $X$ 与 $Z$ 均为仿射时引理成立。

\medskip\noindent
选择仿射开覆盖 $Z = \bigcup Z_i$，使每个 $Z_i$ 都映入 $X$ 的某个仿射开集
$U_i$。由引理 \ref{more-morphisms-lemma-formally-unramified-on-opens}，态射
$Z_i \to U_i$ 形式非分歧。因此由仿射情形，得到 $U_i$ 上的泛一阶加厚
$Z_i \subset Z_i'$。由上述一般性说明，$Z_i \subset Z_i'$ 也是 $Z_i$
在 $X$ 上的泛一阶加厚。令 $Z'_{i, j} \subset Z'_i$ 为满足
$Z_i \cap Z_j = Z'_{i, j} \cap Z_i$ 的开子概形。由一般性说明，
$Z'_{i, j}$ 与 $Z'_{j, i}$ 都是 $Z_i \cap Z_j$ 在 $X$ 上的泛一阶加厚。
于是由第一条一般性说明，存在典范同构
$\varphi_{ij} : Z'_{i, j} \to Z'_{j, i}$，它在 $Z_i \cap Z_j$ 上诱导恒等映射。
我们断言这些态射满足
概形章第 \ref{schemes-section-glueing-schemes} 节 的上闭链条件。
（验证从略。提示：使用 $Z'_{i, j} \cap Z'_{i, k}$ 是
$Z_i \cap Z_j \cap Z_k$ 在 $X$ 上的泛一阶加厚，并由上文知它在唯一同构的
意义下被确定。）因此可以使用
概形章第 \ref{schemes-section-glueing-schemes} 节 的结果，
得到 $X$ 上的一阶加厚 $Z \subset Z'$，使满足
$Z_i = Z'_i \cap Z$ 的开子概形 $Z'_i \subset Z'$ 是 $Z_i$
在 $X$ 上的泛一阶加厚。

\medskip\noindent
由此可形式地推出 $Z'$ 是 $Z$ 在 $X$ 上的泛一阶加厚。事实上，
对于任意如下图，我们已经有泛性质：
$$
\xymatrix{
Z \ar[d]_h & T \ar[l]^a \ar[d] \\
X & T' \ar[l]_b
}
$$
其中 $a(T)$ 包含于某个 $Z_i$。给定一般的图，可选择开覆盖
$T' = \bigcup T'_i$，使 $a(T_i) \subset Z_i$。于是得到 $X$ 上的态射
$a'_i : T'_i \to Z'$，满足 $a'_i|_{T_i} = a|_{T_i}$。
$a'_i$ 与 $a'_j$ 在 $T'_i \cap T'_j$ 上必然相等，因为
$a'_i|_{T'_i \cap T'_j}$ 与 $a'_j|_{T'_i \cap T'_j}$ 都解出了映入
$Z_i \cap Z_j$ 在 $X$ 上的泛一阶加厚 $Z'_i \cap Z'_j$ 的问题。
因此诸态射 $a'_i$ 粘合成所需的 $X$ 上整体态射
$a' : T' \to Z'$。证明完毕。
\end{proof}

\begin{definition}
\label{more-morphisms-definition-universal-thickening}
设 $h : Z \to X$ 是形式非分歧的概形态射。
\begin{enumerate}
\item $Z$ 在 $X$ 上的{\it 泛一阶加厚}是引理
\ref{more-morphisms-lemma-universal-thickening} 中构造的加厚 $Z \subset Z'$。
\item $Z$ 在 $X$ 上的{\it 余法层}是 $Z$ 在其 $X$ 上泛一阶加厚
$Z'$ 中的余法层。
\end{enumerate}
在这种情形下，我们常将该余法层记为 $\mathcal{C}_{Z/X}$。
\end{definition}

\noindent
因此在 $Z$ 上有层的短正合列
$$
0 \to \mathcal{C}_{Z/X} \to \mathcal{O}_{Z'} \to \mathcal{O}_Z \to 0
$$
。下面的引理证明，此定义与 $Z \to X$ 为浸入时的定义并不冲突。

\begin{lemma}
\label{more-morphisms-lemma-immersion-universal-thickening}
设 $i : Z \to X$ 是概形的浸入。则
\begin{enumerate}
\item $i$ 形式非分歧；
\item $Z$ 在 $X$ 上的泛一阶加厚就是定义
\ref{more-morphisms-definition-first-order-infinitesimal-neighbourhood} 中
$Z$ 在 $X$ 中的一阶无穷小邻域；
\item 态射章定义 \ref{morphisms-definition-conormal-sheaf}
意义下 $i$ 的余法层与定义 \ref{more-morphisms-definition-universal-thickening}
意义下 $i$ 的余法层一致。
\end{enumerate}
\end{lemma}

\begin{proof}
由 态射章引理 \ref{morphisms-lemma-open-immersion-unramified} 与
\ref{morphisms-lemma-closed-immersion-unramified}，浸入是非分歧的，
故由引理 \ref{more-morphisms-lemma-unramified-formally-unramified} 它形式非分歧。
其余断言由引理 \ref{more-morphisms-lemma-first-order-infinitesimal-neighbourhood}、
\ref{more-morphisms-lemma-infinitesimal-neighbourhood-conormal} 与诸定义合并得出。
\end{proof}

\begin{lemma}
\label{more-morphisms-lemma-universal-thickening-unramified}
设 $Z \to X$ 是形式非分歧的概形态射。则泛一阶加厚 $Z'$
在 $X$ 上形式非分歧。
\end{lemma}

\begin{proof}
有两个证明。第一个是在仿射局部运算并应用
代数章引理 \ref{algebra-lemma-differentials-universal-thickening}
来证明 $\Omega_{Z'/X} = 0$；此后引理
\ref{more-morphisms-lemma-formally-unramified-differentials} 即给出所需结论。
第二个是如下直接论证。

\medskip\noindent
设 $T \subset T'$ 是一阶加厚，并设
$$
\xymatrix{
Z' \ar[d] & T \ar[l]^c \ar[d] \\
X & T' \ar[l] \ar[lu]^{a, b}
}
$$
为交换图。考虑图中的两个态射 $a, b : T' \to Z'$。令
$T_0 = c^{-1}(Z) \subset T$，并令
$T'_a = a^{-1}(Z)$（均取概形论意义）。
由于 $Z'$ 是 $Z$ 的一阶加厚，$T'$ 是 $T'_a$ 的一阶加厚。
又因 $c = a|_T$，有 $T_0 = T \cap T'_a$（概形论意义）。
由于 $T'$ 是 $T$ 的一阶加厚，遂知 $T'_a$ 是 $T_0$ 的一阶加厚。
现在 $a|_{T'_a}$ 与 $b|_{T'_a}$ 是 $X$ 上从 $T'_a$ 到 $Z'$ 的态射，
且它们限制到 $T_0$ 后作为到 $Z$ 的态射相等。因此由 $Z'$ 的泛性质，
$a|_{T'_a} = b|_{T'_a}$。于是 $a$ 与 $b$ 是从 $T'_a$ 的一阶加厚
$T'$ 出发的态射，它们在 $T'_a$ 上的限制作为到 $Z$ 的态射相等。
再次使用 $Z'$ 的泛性质，得到 $a = b$。换言之，形式非分歧态射的定义性质
对 $Z' \to X$ 成立，正是所求。
\end{proof}

\begin{lemma}
\label{more-morphisms-lemma-universal-thickening-functorial}
考虑概形的交换图
$$
\xymatrix{
Z \ar[r]_h \ar[d]_f & X \ar[d]^g \\
W \ar[r]^{h'} & Y
}
$$
其中 $h$ 与 $h'$ 形式非分歧。设 $Z \subset Z'$ 是 $Z$ 在 $X$ 上的泛一阶加厚，
$W \subset W'$ 是 $W$ 在 $Y$ 上的泛一阶加厚。存在 $Y$ 上加厚的典范态射
$(f, f') : (Z, Z') \to (W, W')$，它嵌入如下交换图：
$$
\xymatrix{
& & & Z' \ar[ld] \ar[d]^{f'} \\
Z \ar[rr] \ar[d]_f \ar[rrru] & & X \ar[d] & W' \ar[ld] \\
W \ar[rrru]|!{[rr];[rruu]}\hole \ar[rr] & & Y
}
$$
特别地，加厚的态射 $(f, f')$ 诱导余法层的态射
$f^*\mathcal{C}_{W/Y} \to \mathcal{C}_{Z/X}$。
\end{lemma}

\begin{proof}
第一个断言由 $W'$ 的泛性质立即可见。余法层上的诱导映射，是将
态射章引理 \ref{morphisms-lemma-conormal-functorial}
应用于 $(Z \subset Z') \to (W \subset W')$ 所得的映射。
\end{proof}

\begin{lemma}
\label{more-morphisms-lemma-universal-thickening-fibre-product}
设
$$
\xymatrix{
Z \ar[r]_h \ar[d]_f & X \ar[d]^g \\
W \ar[r]^{h'} & Y
}
$$
是概形范畴中的纤维积图，并设 $h'$ 形式非分歧。则 $h$ 形式非分歧；
若 $W \subset W'$ 是 $W$ 在 $Y$ 上的泛一阶加厚，则
$Z = X \times_Y W \subset X \times_Y W'$ 是 $Z$ 在 $X$ 上的泛一阶加厚。
特别地，引理 \ref{more-morphisms-lemma-universal-thickening-functorial} 的典范映射
$f^*\mathcal{C}_{W/Y} \to \mathcal{C}_{Z/X}$ 是满射。
\end{lemma}

\begin{proof}
由引理 \ref{more-morphisms-lemma-base-change-formally-unramified}，态射 $h$ 形式非分歧。
显然 $X \times_Y W'$ 是一阶加厚。因为 $W'$ 具有泛性质，
由纤维积的映射性质容易检验它也具有泛性质。
关于这为何推出余法层之间的映射为满射，见
态射章引理 \ref{morphisms-lemma-conormal-functorial-flat}。
\end{proof}

\begin{lemma}
\label{more-morphisms-lemma-universal-thickening-fibre-product-flat}
设
$$
\xymatrix{
Z \ar[r]_h \ar[d]_f & X \ar[d]^g \\
W \ar[r]^{h'} & Y
}
$$
是概形范畴中的纤维积图，其中 $h'$ 形式非分歧且 $g$ 平坦。
在这种情形下，相应的泛一阶加厚之间的映射 $Z' \to W'$ 平坦，并且
$f^*\mathcal{C}_{W/Y} \to \mathcal{C}_{Z/X}$ 是同构。
\end{lemma}

\begin{proof}
平坦性在基变换下保持；见
态射章引理 \ref{morphisms-lemma-base-change-flat}。
因此第一个断言由引理 \ref{more-morphisms-lemma-universal-thickening-fibre-product}
中对 $W'$ 的描述得出。显然 $X \times_Y W'$ 是一阶加厚。
因为 $W'$ 具有泛性质，由纤维积的映射性质容易检验
它也具有泛性质。关于这为何推出余法层之间的映射是同构，见
态射章引理 \ref{morphisms-lemma-conormal-functorial-flat}。
\end{proof}

\begin{lemma}
\label{more-morphisms-lemma-universal-thickening-localize}
取泛一阶加厚与取开集可交换。更精确地，设 $h : Z \to X$
是形式非分歧的概形态射，$V \subset Z$、$U \subset X$ 是满足
$h(V) \subset U$ 的开集，并设 $Z'$ 是 $Z$ 在 $X$ 上的泛一阶加厚。
则 $h|_V : V \to U$ 形式非分歧，而且 $V$ 在 $U$ 上的泛一阶加厚
是满足 $V = Z \cap V'$ 的开子概形 $V' \subset Z'$。特别地，
$\mathcal{C}_{Z/X}|_V = \mathcal{C}_{V/U}$。
\end{lemma}

\begin{proof}
第一个断言即引理 \ref{more-morphisms-lemma-formally-unramified-on-opens}。
泛加厚的相容性可以从引理 \ref{more-morphisms-lemma-universal-thickening} 的证明、
代数章引理 \ref{algebra-lemma-universal-thickening-localize}，
或引理 \ref{more-morphisms-lemma-universal-thickening-fibre-product-flat} 推出。
\end{proof}

\begin{lemma}
\label{more-morphisms-lemma-differentials-universally-unramified}
设 $h : Z \to X$ 是 $S$ 上形式非分歧的概形态射。设 $Z \subset Z'$
是 $Z$ 在 $X$ 上的泛一阶加厚，其结构态射为 $h' : Z' \to X$。典范映射
$$
c_{h'} : (h')^*\Omega_{X/S} \longrightarrow \Omega_{Z'/S}
$$
诱导同构 $h^*\Omega_{X/S} \to \Omega_{Z'/S} \otimes \mathcal{O}_Z$。
\end{lemma}

\begin{proof}
映射 $c_{h'}$ 是 态射章引理
\ref{morphisms-lemma-functoriality-differentials} 中定义的映射。
若 $i : Z \to Z'$ 是给定的闭浸入，则 $i^*c_{h'}$ 是映射
$h^*\Omega_{X/S} \to \Omega_{Z'/S} \otimes \mathcal{O}_Z$。
通过局部化，检验它是同构可归结为仿射情形；见引理
\ref{more-morphisms-lemma-universal-thickening-localize} 与
态射章引理 \ref{morphisms-lemma-differentials-restrict-open}。
在仿射情形下，结论即
代数章引理 \ref{algebra-lemma-differentials-universal-thickening}。
\end{proof}

\begin{lemma}
\label{more-morphisms-lemma-universally-unramified-differentials-sequence}
设 $h : Z \to X$ 是 $S$ 上形式非分歧的概形态射。则有典范正合列
$$
\mathcal{C}_{Z/X} \to h^*\Omega_{X/S} \to \Omega_{Z/S} \to 0.
$$
第一条箭头由 $\text{d}_{Z'/S}$ 诱导，其中
% P05-ERR-CH37-0009: Frozen source calls Z' a universal first order ``neighbourhood'' here; retained pending canon review.
$Z'$ 是 $Z$ 在 $X$ 上的泛一阶邻域。
\end{lemma}

\begin{proof}
我们知道有典范正合列
$$
\mathcal{C}_{Z/Z'} \to
\Omega_{Z'/S} \otimes \mathcal{O}_Z \to
\Omega_{Z/S} \to 0.
$$
；见 态射章引理
\ref{morphisms-lemma-differentials-relative-immersion}。
因此应用引理 \ref{more-morphisms-lemma-differentials-universally-unramified} 即得结论。
\end{proof}

\begin{lemma}
\label{more-morphisms-lemma-two-unramified-morphisms}
设
$$
\xymatrix{
Z \ar[r]_i \ar[rd]_j & X \ar[d] \\
& Y
}
$$
是概形的交换图，其中 $i$ 与 $j$ 形式非分歧。则有典范正合列
$$
\mathcal{C}_{Z/Y} \to
\mathcal{C}_{Z/X} \to
i^*\Omega_{X/Y} \to 0
$$
其中第一条箭头来自引理 \ref{more-morphisms-lemma-universal-thickening-functorial}，
第二条来自引理 \ref{more-morphisms-lemma-universally-unramified-differentials-sequence}。
\end{lemma}

\begin{proof}
以 $Z \to Z'$ 表示 $Z$ 在 $X$ 上的泛一阶加厚，以 $Z \to Z''$
表示 $Z$ 在 $Y$ 上的泛一阶加厚。
% P05-ERR-CH37-0010: Frozen source cites the differential-sequence lemma for this functorial morphism; retained pending canon review.
由引理 \ref{more-morphisms-lemma-universally-unramified-differentials-sequence}，
有典范态射 $Z' \to Z''$，从而有交换图
$$
\xymatrix{
Z \ar[r]_{i'} \ar[rd]_{j'} & Z' \ar[r] \ar[d] & X \ar[d] \\
& Z'' \ar[r] & Y
}
$$
将 态射章引理 \ref{morphisms-lemma-two-immersions}
应用于左边的三角形，得到正合列
$$
\mathcal{C}_{Z/Z''} \to
\mathcal{C}_{Z/Z'} \to
(i')^*\Omega_{Z'/Z''} \to 0
$$
由于 $Z''$ 在 $Y$ 上形式非分歧（见引理
\ref{more-morphisms-lemma-universal-thickening-unramified}），有
% P05-ERR-CH37-0011: Frozen source writes Omega_{Z/Y}; the next sentence uses Omega_{Z'/Y}. Formula retained pending canon review.
$\Omega_{Z'/Z''} = \Omega_{Z/Y}$（结合引理
\ref{more-morphisms-lemma-formally-unramified-differentials} 与
态射章引理 \ref{morphisms-lemma-triangle-differentials}）。
再由引理 \ref{more-morphisms-lemma-differentials-universally-unramified}，有
$(i')^*\Omega_{Z'/Y} = i^*\Omega_{X/Y}$。
\end{proof}

\begin{lemma}
\label{more-morphisms-lemma-transitivity-conormal}
设 $Z \to Y \to X$ 是形式非分歧的概形态射。
\begin{enumerate}
\item 若 $Z \subset Z'$ 是 $Z$ 在 $X$ 上的泛一阶加厚，且
$Y \subset Y'$ 是 $Y$ 在 $X$ 上的泛一阶加厚，则存在态射
$Z' \to Y'$，并且 $Y \times_{Y'} Z'$ 是 $Z$ 在 $Y$ 上的泛一阶加厚。
\item 有典范正合列
$$
i^*\mathcal{C}_{Y/X} \to
\mathcal{C}_{Z/X} \to
\mathcal{C}_{Z/Y} \to 0
$$
其中各映射来自引理 \ref{more-morphisms-lemma-universal-thickening-functorial}，
且 $i : Z \to Y$ 是第一个态射。
\end{enumerate}
\end{lemma}

\begin{proof}
(1) 中的映射 $h : Z' \to Y'$ 来自引理
\ref{more-morphisms-lemma-universal-thickening-functorial}。由 $Z'$ 与 $Y'$ 的泛性质，
$Y \times_{Y'} Z'$ 是 $Z$ 在 $Y$ 上的泛一阶加厚这一断言立即可见。
由 态射章引理 \ref{morphisms-lemma-transitivity-conormal}，有正合列
$$
(i')^*\mathcal{C}_{Y \times_{Y'} Z'/Z'} \to
\mathcal{C}_{Z/Z'} \to
\mathcal{C}_{Z/Y \times_{Y'} Z'} \to 0
$$
其中 $i' : Z \to Y \times_{Y'} Z'$ 是给定态射。由
态射章引理 \ref{morphisms-lemma-conormal-functorial-flat}，
存在满射
$h^*\mathcal{C}_{Y/Y'} \to \mathcal{C}_{Y \times_{Y'} Z'/Z'}$。
再结合等式 $\mathcal{C}_{Y/Y'} = \mathcal{C}_{Y/X}$、
$\mathcal{C}_{Z/Z'} = \mathcal{C}_{Z/X}$ 与
$\mathcal{C}_{Z/Y \times_{Y'} Z'} = \mathcal{C}_{Z/Y}$，即证引理。
\end{proof}











\section{形式平展态射}
\label{more-morphisms-section-formally-etale}

\noindent
回顾：若对每个实线部分交换的图
$$
\xymatrix{
A \ar[r] \ar@{-->}[rd] & B/I \\
R \ar[r] \ar[u] & B \ar[u]
}
$$
（其中 $I \subset B$ 是平方为零的理想），恰好存在一条使该图交换的虚线箭头，
则称环映射 $R \to A$ 是{\it 形式平展的}；见
代数章定义 \ref{algebra-definition-formally-etale}。
这引出概形态射的如下类似概念。

\begin{definition}
\label{more-morphisms-definition-formally-etale}
设 $f : X \to S$ 是概形态射。若给定任意实线部分交换的图
$$
\xymatrix{
X \ar[d]_f & T \ar[d]^i \ar[l] \\
S & T' \ar[l] \ar@{-->}[lu]
}
$$
其中 $T \subset T'$ 是 $S$ 上仿射概形的一阶加厚时，恰好存在一条使该图交换的
虚线箭头，则称 $f$ 是{\it 形式平展的}。
\end{definition}

\noindent
显然，形式平展态射是形式非分歧的。因此若 $f : X \to S$ 形式平展，
则 $\Omega_{X/S}$ 为零；见引理
\ref{more-morphisms-lemma-formally-unramified-differentials}。

\begin{lemma}
\label{more-morphisms-lemma-formally-etale-not-affine}
若 $f : X \to S$ 是形式平展态射，则给定任意实线部分交换的图
$$
\xymatrix{
X \ar[d]_f & T \ar[d]^i \ar[l] \\
S & T' \ar[l] \ar@{-->}[lu]
}
$$
其中 $T \subset T'$ 是 $S$ 上概形的一阶加厚，恰好存在一条使该图交换的虚线箭头。
换言之，在定义 \ref{more-morphisms-definition-formally-etale} 中可以去掉 $T$ 为仿射的条件。
\end{lemma}

\begin{proof}
设 $T' = \bigcup T'_i$ 是仿射开覆盖，并令 $T_i = T \cap T'_i$。
于是得到嵌入图中的态射 $a'_i : T'_i \to X$。由唯一性，
$a'_i$ 与 $a'_j$ 在 $T'_i \cap T'_j$ 的任意仿射开子概形上相等，
故 $a'_i$ 与 $a'_j$ 在 $T'_i \cap T'_j$ 上相等。因此诸态射 $a'_i$ 粘合成整体态射
$a' : T' \to X$。$a'$ 的唯一性已见于引理
\ref{more-morphisms-lemma-formally-unramified-not-affine}。
\end{proof}

\begin{lemma}
\label{more-morphisms-lemma-composition-formally-etale}
形式平展态射的复合仍是形式平展的。
\end{lemma}

\begin{proof}
这是形式性的。
\end{proof}

\begin{lemma}
\label{more-morphisms-lemma-base-change-formally-etale}
形式平展态射的基变换仍是形式平展的。
\end{lemma}

\begin{proof}
这是形式性的。
\end{proof}

\begin{lemma}
\label{more-morphisms-lemma-formally-etale-on-opens}
设 $f : X \to S$ 是概形态射，并设 $U \subset X$ 与 $V \subset S$
是满足 $f(U) \subset V$ 的开子概形。若 $f$ 形式平展，
则 $f|_U : U \to V$ 也形式平展。
\end{lemma}

\begin{proof}
考虑定义 \ref{more-morphisms-definition-formally-etale} 中那样的实线部分交换图
$$
\xymatrix{
U \ar[d]_{f|_U} & T \ar[d]^i \ar[l]^a \\
V & T' \ar[l] \ar@{-->}[lu]
}
$$
。
% P05-ERR-CH37-0012: Frozen source says ``formally ramified'' here; retained literally pending canon review.
若 $f$ 形式分歧，则恰好存在一个满足 $a'|_T = a$ 的 $S$-态射
$a' : T' \to X$。由于 $|T'| = |T|$，可知 $a'(T') \subset U$，
这就给出从 $T'$ 到 $U$ 的唯一态射。
\end{proof}

\begin{lemma}
\label{more-morphisms-lemma-characterize-formally-etale}
设 $f : X \to S$ 是概形态射。以下条件等价：
\begin{enumerate}
\item $f$ 形式平展；
\item $f$ 形式非分歧，且 $X$ 在 $S$ 上的泛一阶加厚等于 $X$；
\item $f$ 形式非分歧，且 $\mathcal{C}_{X/S} = 0$；
\item $\Omega_{X/S} = 0$ 且 $\mathcal{C}_{X/S} = 0$。
\end{enumerate}
\end{lemma}

\begin{proof}
事实上，最后一个断言之所以有意义，是因为 $\Omega_{X/S} = 0$
意味着可由引理 \ref{more-morphisms-lemma-formally-unramified-differentials} 与
定义 \ref{more-morphisms-definition-universal-thickening} 定义 $\mathcal{C}_{X/S}$。
这也说明 (3) 与 (4) 等价。

\medskip\noindent
假设 (1)、(2)、(3) 中任一个，都推出 $f$ 形式非分歧。因此可以假设
$f$ 形式非分歧。(1)、(2)、(3) 的等价性来自 $X$ 在 $S$ 上的泛一阶加厚
$X'$ 的泛性质以及如下事实：
$X = X' \Leftrightarrow \mathcal{C}_{X/S} = 0$；毕竟按定义，
$\mathcal{C}_{X/S} = \mathcal{C}_{X/X'}$ 是 $X$ 在 $X'$ 中的理想层。
\end{proof}

\begin{lemma}
\label{more-morphisms-lemma-unramified-flat-formally-etale}
非分歧平坦态射是形式平展的。
\end{lemma}

\begin{proof}
设 $X \to S$ 非分歧且平坦。则 $\Delta : X \to X \times_S X$
是开浸入；见 态射章引理
\ref{morphisms-lemma-diagonal-unramified-morphism}。
我们须证明 $\mathcal{C}_{X/S}$ 为零。考虑两个投影
$p, q : X \times_S X \to X$。由于 $f$ 形式非分歧（见引理
\ref{more-morphisms-lemma-unramified-formally-unramified}），$q$ 形式非分歧（见引理
\ref{more-morphisms-lemma-base-change-formally-unramified}）。由于 $f$ 平坦，$p$ 平坦；
见 态射章引理 \ref{morphisms-lemma-base-change-flat}。
因此由引理 \ref{more-morphisms-lemma-universal-thickening-fibre-product-flat}，
$p^*\mathcal{C}_{X/S} = \mathcal{C}_q$，其中 $\mathcal{C}_q$
表示形式非分歧态射 $q : X \times_S X \to X$ 的余法层。
但 $\Delta(X) \subset X \times_S X$ 是经 $q$ 同构地映到 $X$ 的开子概形。
故由引理 \ref{more-morphisms-lemma-universal-thickening-localize}，
$\mathcal{C}_q|_{\Delta(X)} = \mathcal{C}_{X/X} = 0$。
换言之，$\mathcal{C}_{X/S}$ 经恒等态射拉回到 $X$ 上为零，即
$\mathcal{C}_{X/S} = 0$。
\end{proof}

\begin{lemma}
\label{more-morphisms-lemma-affine-formally-etale}
设 $f : X \to S$ 是概形态射，并假设 $X$ 与 $S$ 均为仿射。
则 $f$ 形式平展，当且仅当
$\mathcal{O}_S(S) \to \mathcal{O}_X(X)$ 是形式平展环映射。
\end{lemma}

\begin{proof}
由环范畴与仿射概形范畴之间的等价，这立即从定义
（定义 \ref{more-morphisms-definition-formally-etale} 与
代数章定义 \ref{algebra-definition-formally-etale}）得出；见
概形章引理 \ref{schemes-lemma-category-affine-schemes}。
\end{proof}

\begin{lemma}
\label{more-morphisms-lemma-formally-etale}
设 $f : X \to Y$ 是概形态射。以下条件等价：
\begin{enumerate}
\item $f$ 形式平展；
\item 对每个 $x \in X$，存在满足 $f(U) \subset V$ 的开集
$x \in U \subset X$ 与 $f(x) \in V \subset Y$，使得
$f|_U : U \to V$ 形式平展；
\item 对每一对满足 $f(U) \subset V$ 的仿射开集 $U \subset X$ 与
$V \subset Y$，环映射 $\mathcal{O}_Y(V) \to \mathcal{O}_X(U)$ 形式平展；
\item 存在仿射开覆盖 $Y = \bigcup V_j$，且对每个 $j$ 存在仿射开覆盖
$f^{-1}(V_j) = \bigcup U_{ji}$，使得对所有 $j$ 与 $i$，环映射
$\mathcal{O}_Y(V) \to \mathcal{O}_X(U)$ 形式平展。
\end{enumerate}
\end{lemma}

\begin{proof}
结合引理 \ref{more-morphisms-lemma-formally-etale-on-opens} 与
\ref{more-morphisms-lemma-affine-formally-etale} 得 (1) $\Rightarrow$ (3)。
(3) $\Rightarrow$ (4) 是立即的。由引理 \ref{more-morphisms-lemma-affine-formally-etale}
得 (4) $\Rightarrow$ (2)。若 (2) 成立，则由引理
\ref{more-morphisms-lemma-characterize-formally-etale}（这里还使用 态射章引理
\ref{morphisms-lemma-differentials-restrict-open} 与
引理 \ref{more-morphisms-lemma-universal-thickening-localize}），
$\Omega_{X/S}$ 与 $\mathcal{C}_{X/S}$ 在每一点的某个邻域上消失；于是由引理
\ref{more-morphisms-lemma-characterize-formally-etale}，
% P05-ERR-CH37-0013: Frozen source concludes ``formally unramified'' rather than ``formally etale''; retained pending canon review.
$f$ 形式非分歧。
\end{proof}

\begin{lemma}
\label{more-morphisms-lemma-etale-formally-etale}
设 $f : X \to S$ 是概形态射。以下条件等价：
\begin{enumerate}
\item 态射 $f$ 平展；
\item 态射 $f$ 局部有限表示且形式平展。
\end{enumerate}
\end{lemma}

\begin{proof}
假设 $f$ 平展。平展态射局部有限表示、平坦且非分歧；见
态射章第 \ref{morphisms-section-etale} 节。
故 $f$ 局部有限表示且形式平展；见引理
\ref{more-morphisms-lemma-unramified-flat-formally-etale}。

\medskip\noindent
反之，假设 $f$ 局部有限表示且形式平展。平展性在 $X$ 与 $S$ 的 Zariski
拓扑上是局部的；见 态射章引理
\ref{morphisms-lemma-etale-characterize}。由引理
\ref{more-morphisms-lemma-formally-etale-on-opens}，可用映入仿射开集 $V$ 的仿射开集
$U$ 覆盖 $X$，使 $U \to V$ 形式平展（且有限表示；见 Morphisms,
引理 \ref{morphisms-lemma-locally-finite-presentation-characterize}）。
由引理 \ref{more-morphisms-lemma-affine-formally-etale}，环映射
$\mathcal{O}(V) \to \mathcal{O}(U)$ 形式平展（且有限表示）。
于是由 代数章引理 \ref{algebra-lemma-formally-etale-etale} 得证。
（讨论形式光滑态射时，我们还将给出这一蕴含的另一个证明。）
\end{proof}













\section{映射的无穷小形变}
\label{more-morphisms-section-action-by-derivations}

\noindent
本节说明如何用导子使映射作无穷小移动。贯穿本节，我们使用如下事实：
$X$ 的加厚 $X'$ 上的层可以视为 $X$ 上的层。

\begin{lemma}
\label{more-morphisms-lemma-difference-derivation}
设 $S$ 是概形，$X \subset X'$ 与 $Y \subset Y'$ 是 $S$ 上的两个一阶加厚，
且 $(a, a'), (b, b') : (X \subset X') \to (Y \subset Y')$
是 $S$ 上加厚的两个态射。假设
\begin{enumerate}
\item $a = b$；
\item 两个映射 $a^*\mathcal{C}_{Y/Y'} \to \mathcal{C}_{X/X'}$
（态射章引理 \ref{morphisms-lemma-conormal-functorial}）相等。
\end{enumerate}
则映射 $(a')^\sharp - (b')^\sharp$ 分解为
$$
\mathcal{O}_{Y'} \to \mathcal{O}_Y \xrightarrow{D}
a_*\mathcal{C}_{X/X'} \to a_*\mathcal{O}_{X'}
$$
其中 $D$ 是 $\mathcal{O}_S$-导子。
\end{lemma}

\begin{proof}
不在 $Y$ 上而在 $X$ 上运算；其优点是拉回函子 $a^{-1}$ 正合。
由 (1) 与 (2)，得到行正合的交换图
$$
\xymatrix{
0 \ar[r] &
\mathcal{C}_{X/X'} \ar[r] &
\mathcal{O}_{X'} \ar[r] &
\mathcal{O}_X \ar[r] & 0 \\
0 \ar[r] &
a^{-1}\mathcal{C}_{Y/Y'} \ar[r] \ar[u] &
a^{-1}\mathcal{O}_{Y'}
\ar[r] \ar@<1ex>[u]^{(a')^\sharp} \ar@<-1ex>[u]_{(b')^\sharp} &
a^{-1}\mathcal{O}_Y \ar[r] \ar[u] & 0
}
$$
一般而言，在这种情形下，$\mathcal{O}_S$-代数映射
$(a')^\sharp$ 与 $(b')^\sharp$ 之差是从 $a^{-1}\mathcal{O}_Y$ 到
$\mathcal{C}_{X/X'}$ 的 $\mathcal{O}_S$-导子。由函子 $a^{-1}$ 与 $a_*$
的伴随性，这等价于从 $\mathcal{O}_Y$ 到
$a_*\mathcal{C}_{X/X'}$ 的 $\mathcal{O}_S$-导子。略去一些细节。
\end{proof}

\noindent
注意，在上述引理的情形下，可以把 $D$ 写成
\begin{equation}
\label{more-morphisms-equation-D}
D = \text{d}_{Y/S} \circ \theta
\end{equation}
其中 $\theta$ 是 $\mathcal{O}_Y$-线性映射
$\theta : \Omega_{Y/S} \to a_*\mathcal{C}_{X/X'}$。
当然，再次由伴随性，也可把 $\theta$ 看作 $\mathcal{O}_X$-线性映射
$\theta : a^*\Omega_{Y/S} \to \mathcal{C}_{X/X'}$。

\begin{lemma}
\label{more-morphisms-lemma-action-by-derivations}
设 $S$ 是概形，且
$(a, a') : (X \subset X') \to (Y \subset Y')$
是 $S$ 上一阶加厚的态射。设
$$
\theta : a^*\Omega_{Y/S} \to \mathcal{C}_{X/X'}
$$
是 $\mathcal{O}_X$-线性映射。则存在唯一的对的态射
$(b, b') : (X \subset X') \to (Y \subset Y')$，使引理
\ref{more-morphisms-lemma-difference-derivation} 的 (1)、(2) 成立，并且导子 $D$
与 $\theta$ 由等式 (\ref{more-morphisms-equation-D}) 联系。
\end{lemma}

\begin{proof}
只需令 $b = a$，并将 $(b')^\sharp$ 定义为映射
$$
(a')^\sharp + D : a^{-1}\mathcal{O}_{Y'} \to \mathcal{O}_{X'}
$$
其中 $D$ 如等式 (\ref{more-morphisms-equation-D}) 所示。略去如下验证：
$(b')^\sharp$ 是 $\mathcal{O}_S$-代数层的映射，且引理
\ref{more-morphisms-lemma-difference-derivation} 的 (1)、(2) 成立。
等式 (\ref{more-morphisms-equation-D}) 由构造成立。
\end{proof}

\begin{remark}
\label{more-morphisms-remark-action-by-derivations}
假设与记号同引理 \ref{more-morphisms-lemma-action-by-derivations}。
局部截面 $\theta$ 在 $a'$ 上的作用有时记为 $\theta \cdot a'$。
注意，这无非表示 $(a')^\sharp$ 与 $(\theta \cdot a')^\sharp$ 之差为导子
$D$，它与 $\theta$ 由等式 (\ref{more-morphisms-equation-D}) 联系。
\end{remark}

\begin{lemma}
\label{more-morphisms-lemma-sheaf}
设 $S$ 是概形，$X \subset X'$ 与 $Y \subset Y'$ 是 $S$ 上的一阶加厚。
假设给定态射 $a : X \to Y$ 以及 $\mathcal{O}_X$-模的映射
$A : a^*\mathcal{C}_{Y/Y'} \to \mathcal{C}_{X/X'}$。
对开子概形 $U' \subset X'$，考虑满足下列条件的态射 $a' : U' \to Y'$：
\begin{enumerate}
\item $a'$ 是 $S$ 上的态射；
\item $a'|_U = a|_U$；
\item 诱导映射 $a^*\mathcal{C}_{Y/Y'}|_U \to \mathcal{C}_{X/X'}|_U$
是 $A$ 在 $U$ 上的限制。
\end{enumerate}
这里 $U = X \cap U'$。则规则
\begin{equation}
\label{more-morphisms-equation-sheaf}
U' \mapsto
\{a' : U' \to Y'\text{ such that (1), (2), (3) hold.}\}
\end{equation}
定义了 $X'$ 上的集合层。
\end{lemma}

\begin{proof}
用 $\mathcal{F}$ 表示引理中的规则。对
$V' \subset U' \subset X'$，$\mathcal{F}$ 的限制映射
$\mathcal{F}(U') \to \mathcal{F}(V')$ 就是限制映射
$a' \mapsto a'|_{V'}$。在此定义下，由于态射是局部定义的，
$\mathcal{F}$ 显然是层。
\end{proof}

\noindent
在下一个引理中，我们将 $X$ 上的层同 $X$ 的任意加厚上的层等同起来。

\begin{lemma}
\label{more-morphisms-lemma-action-sheaf}
记号与假设同引理 \ref{more-morphisms-lemma-sheaf}。层
$$
\SheafHom_{\mathcal{O}_X}(a^*\Omega_{Y/S}, \mathcal{C}_{X/X'})
$$
作用在层 (\ref{more-morphisms-equation-sheaf}) 上。此外，对层
(\ref{more-morphisms-equation-sheaf}) 在其上有截面的任意开集 $U' \subset X'$，
该作用是单传递的。
\end{lemma}

\begin{proof}
这是引理 \ref{more-morphisms-lemma-difference-derivation}、
\ref{more-morphisms-lemma-action-by-derivations} 与 \ref{more-morphisms-lemma-sheaf} 的合并。
\end{proof}

\begin{remark}
\label{more-morphisms-remark-special-case}
引理 \ref{more-morphisms-lemma-difference-derivation}、
\ref{more-morphisms-lemma-action-by-derivations}、\ref{more-morphisms-lemma-sheaf} 与
\ref{more-morphisms-lemma-action-sheaf} 的一个特殊情形是 $Y = Y'$。
此时映射 $A$ 总为零，引理 \ref{more-morphisms-lemma-sheaf} 的层就是由规则
$$
U' \mapsto
\{a' : U' \to Y\text{ 相对于 }S\text{ 其中 } a'|_U = a|_U\}
$$
给出的层，而层
$\SheafHom_{\mathcal{O}_X}(a^*\Omega_{Y/S}, \mathcal{C}_{X/X'})$
作用在其上。
\end{remark}

\begin{remark}
\label{more-morphisms-remark-another-special-case}
引理 \ref{more-morphisms-lemma-difference-derivation}、
\ref{more-morphisms-lemma-action-by-derivations}、\ref{more-morphisms-lemma-sheaf} 与
\ref{more-morphisms-lemma-action-sheaf} 的另一个特殊情形是：$S$ 本身为加厚
$Z \subset Z' = S$，且 $Y = Z \times_{Z'} Y'$。图示为
$$
\xymatrix{
(X \subset X') \ar@{..>}[rr]_{(a, ?)} \ar[rd]_{(g, g')} & &
(Y \subset Y') \ar[ld]^{(h, h')} \\
& (Z \subset Z')
}
$$
此时映射 $A : a^*\mathcal{C}_{Y/Y'} \to \mathcal{C}_{X/X'}$
由 $a$ 决定：映射 $h^*\mathcal{C}_{Z/Z'} \to \mathcal{C}_{Y/Y'}$
是满射（因为假设了 $Y = Z \times_{Z'} Y'$），所以拉回
$g^*\mathcal{C}_{Z/Z'} = a^*h^*\mathcal{C}_{Z/Z'} \to
a^*\mathcal{C}_{Y/Y'}$ 是满射；而复合
$g^*\mathcal{C}_{Z/Z'} \to a^*\mathcal{C}_{Y/Y'} \to \mathcal{C}_{X/X'}$
必须是由 $g'$ 诱导的典范映射。因此引理 \ref{more-morphisms-lemma-sheaf} 的层
就是由规则
$$
U' \mapsto
\{a' : U' \to Y'\text{ 相对于 }Z'\text{ 其中 } a'|_U = a|_U\}
$$
给出的层，而层
$\SheafHom_{\mathcal{O}_X}(a^*\Omega_{Y/Z}, \mathcal{C}_{X/X'})$
作用在其上。
\end{remark}

\begin{lemma}
\label{more-morphisms-lemma-omega-deformation}
设 $S$ 是概形，$X \subset X'$ 是 $S$ 上的一阶加厚，$Y$ 是 $S$ 上的概形。
设 $a', b' : X' \to Y$ 是 $S$ 上满足
$a = a'|_X = b'|_X$ 的两个态射。由此得到交换图
$$
\xymatrix{
X \ar[r] \ar[d]_a & X' \ar[d]^{(b', a')} \\
Y \ar[r]^-{\Delta_{Y/S}} & Y \times_S Y
}
$$
由于水平箭头都是浸入，其余法层分别为
$\mathcal{C}_{X/X'}$ 与 $\Omega_{Y/S}$，由
态射章引理 \ref{morphisms-lemma-conormal-functorial}，
得到映射 $\theta : a^*\Omega_{Y/S} \to \mathcal{C}_{X/X'}$。
于是此 $\theta$ 与引理 \ref{more-morphisms-lemma-difference-derivation} 的导子 $D$
由等式 (\ref{more-morphisms-equation-D}) 联系。
\end{lemma}

\begin{proof}
从略。提示：可在仿射开集上检验该等式，此时它来自如下计算。
若 $f$ 是 $\mathcal{O}_Y$ 的局部截面，则
$1 \otimes f - f \otimes 1$ 是 $\mathcal{C}_{Y/(Y \times_S Y)}$
中对应于 $\text{d}_{Y/S}(f)$ 的局部截面。它被映到
$\mathcal{C}_{X/X'}$ 的局部截面
$(a')^\sharp(f) - (b')^\sharp(f) = D(f)$。换言之，
$\theta(\text{d}_{Y/S}(f)) = D(f)$。
\end{proof}

\noindent
为后续用途，需要一个粗略地说如下的结果：引理
\ref{more-morphisms-lemma-action-by-derivations} 的构造与平展局部化相容。

\begin{lemma}
\label{more-morphisms-lemma-sheaf-differentials-etale-localization}
设
$$
\xymatrix{
X_1 \ar[d] & X_2 \ar[l]^f \ar[d] \\
S_1 & S_2 \ar[l]
}
$$
是概形的交换图，其中 $X_2 \to X_1$ 与 $S_2 \to S_1$ 平展。则
态射章引理 \ref{morphisms-lemma-functoriality-differentials}
的映射 $c_f : f^*\Omega_{X_1/S_1} \to \Omega_{X_2/S_2}$ 是同构。
\end{lemma}

\begin{proof}
回顾，平展态射 $U \to V$ 是满足 $\Omega_{U/V} = 0$ 的光滑态射。
由此可见，态射章引理
\ref{morphisms-lemma-triangle-differentials} 推出
$\Omega_{X_2/S_2} = \Omega_{X_2/S_1}$，而 态射章引理
\ref{morphisms-lemma-triangle-differentials-smooth} 推出映射
$f^*\Omega_{X_1/S_1} \to \Omega_{X_2/S_1}$（将态射 $f$ 看作
$S_1$ 上的态射）是同构。故引理成立。
\end{proof}

\begin{lemma}
\label{more-morphisms-lemma-action-by-derivations-etale-localization}
考虑一阶加厚的交换图
$$
\vcenter{
\xymatrix{
(T_2 \subset T_2') \ar[d]_{(h, h')} \ar[rr]_{(a_2, a_2')} & &
(X_2 \subset X_2') \ar[d]^{(f, f')} \\
(T_1 \subset T_1') \ar[rr]^{(a_1, a_1')} & &
(X_1 \subset X_1')
}
}
\quad
\begin{matrix}
\text{and a commutative} \\
\text{diagram of schemes}
\end{matrix}
\quad
\vcenter{
\xymatrix{
X_2' \ar[r] \ar[d] & S_2 \ar[d] \\
X_1' \ar[r] & S_1
}
}
$$
其中 $X_2 \to X_1$ 与 $S_2 \to S_1$ 平展。对任意
$\mathcal{O}_{T_1}$-线性映射
$\theta_1 : a_1^*\Omega_{X_1/S_1} \to \mathcal{C}_{T_1/T'_1}$，
令 $\theta_2$ 为复合
$$
\xymatrix{
a_2^*\Omega_{X_2/S_2} \ar@{=}[r] &
h^*a_1^*\Omega_{X_1/S_1} \ar[r]^-{h^*\theta_1} &
h^*\mathcal{C}_{T_1/T'_1} \ar[r] &
\mathcal{C}_{T_2/T'_2}
}
$$
（等号将在证明中说明）。则图
$$
\xymatrix{
T_2' \ar[rr]_{\theta_2 \cdot a_2'} \ar[d] & & X'_2 \ar[d] \\
T_1' \ar[rr]^{\theta_1 \cdot a_1'} & & X'_1
}
$$
交换，其中作用 $\theta_2 \cdot a_2'$ 与 $\theta_1 \cdot a_1'$
如注记 \ref{more-morphisms-remark-action-by-derivations} 所述。
\end{lemma}

\begin{proof}
该等号来自引理 \ref{more-morphisms-lemma-sheaf-differentials-etale-localization}
的等同 $f^*\Omega_{X_1/S_1} = \Omega_{X_2/S_2}$。事实上，由此以及
$f \circ a_2 = a_1 \circ h$，有
$a_2^*\Omega_{X_2/S_2} = a_2^*f^*\Omega_{X_1/S_1} =
h^*a_1^*\Omega_{X_1/S_1}$。说明这一点后，可以在仿射开集上检验该图的
交换性。因此可假设最初大图中的概形都是仿射的。于是得到交换图
$$
\vcenter{
\xymatrix{
(B'_2, I_2) & & (A'_2, J_2) \ar[ll]^{a_2'} \\
(B'_1, I_1) \ar[u]^{h'} & & (A'_1, J_1) \ar[ll]_{a_1'} \ar[u]_{f'}
}
}
\quad\text{且}\quad
\vcenter{
\xymatrix{
A'_2 & & R_2 \ar[ll] \\
A'_1 \ar[u] & & R_1 \ar[ll] \ar[u]
}
}
$$
该记号表示 $I_1, I_2, J_1, J_2$ 是平方为零的理想，而对之间的映射是把理想
送入理想的环映射。令 $A_1 = A'_1/J_1$、$A_2 = A'_2/J_2$、
$B_1 = B'_1/I_1$ 且 $B_2 = B'_2/I_2$。已知
$$
A_2 \otimes_{A_1} \Omega_{A_1/R_1} \longrightarrow \Omega_{A_2/R_2}
$$
是同构。映射
$\theta_1 : B_1 \otimes_{A_1} \Omega_{A_1/R_1} \to I_1$
是 $B_1$-线性的，因而给出 $R_1$-导子
$D_1 = \theta_1 \circ \text{d}_{A_1/R_1} : A_1 \to I_1$。
类似地，映射
$\theta_2 : B_2 \otimes_{A_2} \Omega_{A_2/R_2} \to I_2$
给出 $R_2$-导子
$D_2 = \theta_2 \circ \text{d}_{A_2/R_2} : A_2 \to I_2$。
$\theta_2$ 的构造蕴含 $\theta_1$ 与 $\theta_2$ 的如下相容性：
对每个 $x \in A_1$，有
$$
h'(D_1(x)) = D_2(f'(x))
$$
，作为 $I_2$ 的元素。通过
$A'_1 \to A_1 \xrightarrow{D_1} I_1 \to B_1$，可把 $D_1$
看作映射 $A'_1 \to B'_1$；类似地，可把 $D_2$ 看作映射
$A'_2 \to B'_2$。于是上述等式对 $x \in A'_1$ 成立。
由引理 \ref{more-morphisms-lemma-action-by-derivations} 与注记
\ref{more-morphisms-remark-action-by-derivations} 中作用的构造，
$\theta_1 \cdot a_1'$ 对应于环映射 $a_1' + D_1 : A'_1 \to B'_1$，
而 $\theta_2 \cdot a_2'$ 对应于环映射
$a_2' + D_2 : A'_2 \to B'_2$。由上述等式得到
$h' \circ (a_1' + D_1) = (a_2' + D_2) \circ f'$
，正是所需结论。
\end{proof}

\begin{remark}
\label{more-morphisms-remark-tiny-improvement}
引理 \ref{more-morphisms-lemma-action-by-derivations-etale-localization}
可作如下改进。假设有引理
\ref{more-morphisms-lemma-action-by-derivations-etale-localization} 中的交换图，但不假设
$X_2 \to X_1$ 与 $S_2 \to S_1$ 平展。再假设有
$\theta_1 : a_1^*\Omega_{X_1/S_1} \to \mathcal{C}_{T_1/T'_1}$
及
$\theta_2 : a_2^*\Omega_{X_2/S_2} \to \mathcal{C}_{T_2/T'_2}$
使得
$$
\xymatrix{
f_*\mathcal{O}_{X_2} \ar[rr]_{f_*D_2} & &
f_*a_{2, *}\mathcal{C}_{T_2/T_2'} \\
\mathcal{O}_{X_1} \ar[rr]^{D_1} \ar[u]^{f^\sharp} & &
a_{1, *}\mathcal{C}_{T_1/T_1'} \ar[u]_{\text{由下列对象诱导： }(h')^\sharp}
}
$$
交换，其中 $D_i$ 如等式 (\ref{more-morphisms-equation-D}) 那样对应于 $\theta_i$。
则引理 \ref{more-morphisms-lemma-action-by-derivations-etale-localization} 的结论成立。
$X_2 \to X_1$ 与 $S_2 \to S_1$ 均平展这一条件的重要性在于，
它使我们能够从 $\theta_1$ 构造 $\theta_2$。
\end{remark}








\section{概形的无穷小形变}
\label{more-morphisms-section-deform}

\noindent
下面的简单引理常是检验映射的无穷小形变是否平坦的便利工具。

\begin{lemma}
\label{more-morphisms-lemma-deform}
设 $(f, f') : (X \subset X') \to (S \subset S')$ 是一阶加厚的态射，
并假设 $f$ 平坦。则以下条件等价：
\begin{enumerate}
\item $f'$ 平坦，且 $X = S \times_{S'} X'$；
\item 典范映射 $f^*\mathcal{C}_{S/S'} \to \mathcal{C}_{X/X'}$ 是同构。
\end{enumerate}
\end{lemma}

\begin{proof}
由于问题在 $X'$ 上是局部的，可假设 $X, X', S, S'$ 都是仿射概形。
写 $S' = \Spec(A')$、$X' = \Spec(B')$、$S = \Spec(A)$、
$X = \Spec(B)$，其中对某些平方为零的理想有
$A = A'/I$ 且 $B = B'/J$。于是得到交换图
$$
\xymatrix{
0 \ar[r] &
J \ar[r] &
B' \ar[r] &
B \ar[r] & 0 \\
0 \ar[r] &
I \ar[r] \ar[u] &
A' \ar[r] \ar[u] &
A \ar[r] \ar[u] & 0
}
$$
，其行正合。引理中的典范映射是
$$
I \otimes_A B = I \otimes_{A'} B' \longrightarrow J.
$$
$f$ 平坦这一假设表示 $A \to B$ 平坦。

\medskip\noindent
假设 (1)。则 $A' \to B'$ 平坦且 $J = IB'$。平坦性推出
$\text{Tor}_1^{A'}(B', A) = 0$（见
代数章引理 \ref{algebra-lemma-characterize-flat}）。
这表示 $I \otimes_{A'} B' \to B'$ 为单射（见
代数章说明 \ref{algebra-remark-Tor-ring-mod-ideal}）。
故 $I \otimes_A B \to J$ 是同构。

\medskip\noindent
假设 (2)。由此 $J = IB'$，故 $X = S \times_{S'} X'$。
此外，反向使用上一段的诸蕴含，得到
$\text{Tor}_1^{A'}(B', A'/I) = 0$。因此由
代数章引理 \ref{algebra-lemma-what-does-it-mean}，
$B'$ 在 $A'$ 上平坦。
\end{proof}

\noindent
下面的引理是“逐纤维平坦性判据”（``crit\`ere de platitude par fibres''）
的“幂零”版本；见节 \ref{more-morphisms-section-criterion-flat-fibres}。

\begin{lemma}
\label{more-morphisms-lemma-flatness-morphism-thickenings}
考虑交换图
$$
\xymatrix{
(X \subset X') \ar[rr]_{(f, f')} \ar[rd] & & (Y \subset Y') \ar[ld] \\
& (S \subset S')
}
$$
，其中各对象为加厚。假设
\begin{enumerate}
\item $X'$ 在 $S'$ 上平坦；
\item $f$ 平坦；
\item $S \subset S'$ 是有限阶加厚；
\item $X = S \times_{S'} X'$ 且 $Y = S \times_{S'} Y'$。
\end{enumerate}
则 $f'$ 平坦，并且 $Y'$ 在 $f'$ 的像中的每一点处都在 $S'$ 上平坦。
\end{lemma}

\begin{proof}
这是 代数章引理
\ref{algebra-lemma-criterion-flatness-fibre-nilpotent} 的直接推论。
\end{proof}

\noindent
概形态射的许多性质在平坦形变下保持。

\begin{lemma}
\label{more-morphisms-lemma-deform-property}
考虑交换图
$$
\xymatrix{
(X \subset X') \ar[rr]_{(f, f')} \ar[rd] & & (Y \subset Y') \ar[ld] \\
& (S \subset S')
}
$$
，其中各对象为加厚。假设 $S \subset S'$ 是有限阶加厚，$X'$ 在 $S'$ 上平坦，
$X = S \times_{S'} X'$，且 $Y = S \times_{S'} Y'$。则
\begin{enumerate}
\item $f$ 平坦，当且仅当 $f'$ 平坦；
\label{more-morphisms-item-flat}
\item $f$ 是同构，当且仅当 $f'$ 是同构；
\label{more-morphisms-item-isomorphism}
\item $f$ 是开浸入，当且仅当 $f'$ 是开浸入；
\label{more-morphisms-item-open-immersion}
\item $f$ 拟紧，当且仅当 $f'$ 拟紧；
\label{more-morphisms-item-quasi-compact}
\item $f$ 泛闭，当且仅当 $f'$ 泛闭；
\label{more-morphisms-item-universally-closed}
\item $f$（拟）分离，当且仅当 $f'$（拟）分离；
\label{more-morphisms-item-separated}
\item $f$ 是单态射，当且仅当 $f'$ 是单态射；
\label{more-morphisms-item-monomorphism}
\item $f$ 是满射，当且仅当 $f'$ 是满射；
\label{more-morphisms-item-surjective}
\item $f$ 泛单射，当且仅当 $f'$ 泛单射；
\label{more-morphisms-item-universally-injective}
\item $f$ 是仿射态射，当且仅当 $f'$ 是仿射态射；
\label{more-morphisms-item-affine}
\item
\label{more-morphisms-item-finite-type}
$f$ 局部有限型，当且仅当 $f'$ 局部有限型；
\item $f$ 局部拟有限，当且仅当 $f'$ 局部拟有限；
\label{more-morphisms-item-quasi-finite}
\item
\label{more-morphisms-item-finite-presentation}
$f$ 局部有限表示，当且仅当 $f'$ 局部有限表示；
\item
\label{more-morphisms-item-relative-dimension-d}
$f$ 局部有限型且相对维数为 $d$，当且仅当 $f'$ 局部有限型且相对维数为 $d$；
\item $f$ 泛开，当且仅当 $f'$ 泛开；
\label{more-morphisms-item-universally-open}
\item $f$ 是 syntomic 态射，当且仅当 $f'$ 是 syntomic 态射；
\label{more-morphisms-item-syntomic}
\item $f$ 光滑，当且仅当 $f'$ 光滑；
\label{more-morphisms-item-smooth}
\item $f$ 非分歧，当且仅当 $f'$ 非分歧；
\label{more-morphisms-item-unramified}
\item $f$ 平展，当且仅当 $f'$ 平展；
\label{more-morphisms-item-etale}
\item $f$ 固有，当且仅当 $f'$ 固有；
\label{more-morphisms-item-proper}
\item $f$ 是整态射，当且仅当 $f'$ 是整态射；
\label{more-morphisms-item-integral}
\item $f$ 是有限态射，当且仅当 $f'$ 是有限态射；
\label{more-morphisms-item-finite}
\item
\label{more-morphisms-item-finite-locally-free}
$f$ 是有限局部自由态射（秩为 $d$），当且仅当 $f'$ 是有限局部自由态射
（秩为 $d$）；
\item 此处还应添加更多性质。
\end{enumerate}
\end{lemma}

\begin{proof}
关于 $X$ 与 $Y$ 的假设意味着，$f$ 是 $f'$ 经 $X \to X'$ 的基变换。
上述 (1)--(23) 中列出的性质 $\mathcal{P}$ 在基变换下都保持，
因此若 $f'$ 具有性质 $\mathcal{P}$，则 $f$ 也具有。见
概形章引理 \ref{schemes-lemma-base-change-immersion},
\ref{schemes-lemma-quasi-compact-preserved-base-change},
\ref{schemes-lemma-separated-permanence}, 及
\ref{schemes-lemma-base-change-monomorphism}
及
态射章引理
\ref{morphisms-lemma-base-change-surjective},
\ref{morphisms-lemma-base-change-universally-injective},
\ref{morphisms-lemma-base-change-affine},
\ref{morphisms-lemma-base-change-finite-type},
\ref{morphisms-lemma-base-change-quasi-finite},
\ref{morphisms-lemma-base-change-finite-presentation},
\ref{morphisms-lemma-base-change-relative-dimension-d},
\ref{morphisms-lemma-base-change-syntomic},
\ref{morphisms-lemma-base-change-smooth},
\ref{morphisms-lemma-base-change-unramified},
\ref{morphisms-lemma-base-change-etale},
\ref{morphisms-lemma-base-change-proper},
\ref{morphisms-lemma-base-change-finite}, 及
\ref{morphisms-lemma-base-change-finite-locally-free}.

\medskip\noindent
因此每种情形中有趣的方向都是：假设 $f$ 具有该性质，并推出 $f'$ 也具有。
对加厚的阶数作归纳，可假设 $S \subset S'$ 是一阶加厚；见紧接定义
\ref{more-morphisms-definition-thickening} 之后的讨论。先作几条一般性说明，
以下论证中将不再逐次提及。
(I) 设 $W' \subset S'$ 是仿射开集，并设 $U' \subset X'$ 与
$V' \subset Y'$ 是位于 $W'$ 上方、满足 $f'(U') \subset V'$ 的仿射开集。
令 $W' = \Spec(R')$，并以 $I \subset R'$ 表示定义闭子概形
$W' \cap S$ 的理想。写 $U' = \Spec(B')$ 且 $V' = \Spec(A')$。
于是得到交换图
$$
\xymatrix{
0 \ar[r] &
IB' \ar[r] &
B' \ar[r] &
B \ar[r] & 0 \\
0 \ar[r] &
IA' \ar[r] \ar[u] &
A' \ar[r] \ar[u] &
A \ar[r] \ar[u] & 0
}
$$
，其行正合。此外 $IB' \cong I \otimes_R B$；见引理
\ref{more-morphisms-lemma-deform} 的证明。
(II) 态射 $X \to X'$ 与 $Y \to Y'$ 是泛同胚。因此映射 $f$ 与 $f'$
（在任意基变换后）的拓扑相同。
(III) 若 $f$ 平坦，则 $f'$ 平坦，且 $Y' \to S'$ 在 $f'$ 的像中的每一点处
都平坦；见引理 \ref{more-morphisms-lemma-flatness-morphism-thickenings}。

\medskip\noindent
关于 (\ref{more-morphisms-item-flat})。这就是一般性说明 (III)。

\medskip\noindent
关于 (\ref{more-morphisms-item-isomorphism})。假设 $f$ 是同构。由 (III)，
$Y' \to S'$ 平坦。选择仿射开集 $V' \subset Y'$，并令
$U' = (f')^{-1}(V')$。于是 $V = Y \cap V'$ 是仿射的，从而
$V \cong f^{-1}(V) = U = Y \times_{Y'} U'$ 是仿射的。由引理
\ref{more-morphisms-lemma-thickening-affine-scheme}，$U'$ 是仿射的。
因此有一般性说明 (I) 中那样的图；又因 $R' \to A'$ 平坦，
$IA \cong I \otimes_R A$。于是
$IB' \cong I \otimes_R B \cong I \otimes_R A \cong IA'$，
且 $A \cong B$。由上图各行的正合性，$A' \cong B'$，即
$U' \cong V'$。故 $f'$ 是同构。

\medskip\noindent
关于 (\ref{more-morphisms-item-open-immersion})。假设 $f$ 是开浸入。
则 $f$ 给出 $X$ 与开子概形 $V \subset Y$ 之间的同构。
令 $V' \subset Y'$ 为底拓扑空间等于 $V$ 的开子概形。于是 $f'$
是从 $X'$ 到 $V'$ 的映射，由 (\ref{more-morphisms-item-isomorphism}) 它是同构。
故 $f'$ 是开浸入。

\medskip\noindent
关于 (\ref{more-morphisms-item-quasi-compact})。由说明 (II) 立即可得。更一般的陈述亦见
引理 \ref{more-morphisms-lemma-thicken-property-morphisms}。

\medskip\noindent
关于 (\ref{more-morphisms-item-universally-closed})。由说明 (II) 立即可得。更一般的陈述亦见
引理 \ref{more-morphisms-lemma-thicken-property-morphisms}。

\medskip\noindent
关于 (\ref{more-morphisms-item-separated})。注意
$X \times_Y X = Y \times_{Y'} (X' \times_{Y'} X')$，故
$X' \times_{Y'} X'$ 是 $X \times_Y X$ 的一个加厚。
映射 $\Delta_{X/Y}$ 与 $\Delta_{X'/Y'}$
的拓扑相符，故得证。更一般的陈述亦见引理
\ref{more-morphisms-lemma-thicken-property-morphisms}。

\medskip\noindent
关于 (\ref{more-morphisms-item-monomorphism})。假设 $f$ 是单态射。考虑对角态射
$\Delta_{X'/Y'} : X' \to X' \times_{Y'} X'$。
$\Delta_{X'/Y'}$ 经 $S \to S'$ 的基变换为 $\Delta_{X/Y}$，
后者按假设是同构。由 (\ref{more-morphisms-item-isomorphism})，
$\Delta_{X'/Y'}$ 是同构。

\medskip\noindent
关于 (\ref{more-morphisms-item-surjective})。这是显然的。更一般的陈述亦见引理
\ref{more-morphisms-lemma-thicken-property-morphisms}。

\medskip\noindent
关于 (\ref{more-morphisms-item-universally-injective})。由说明 (II) 立即可得。
更一般的陈述亦见引理 \ref{more-morphisms-lemma-thicken-property-morphisms}。

\medskip\noindent
关于 (\ref{more-morphisms-item-affine})。假设 $f$ 是仿射态射。选择仿射开集
$V' \subset Y'$ 并令 $U' = (f')^{-1}(V')$。于是
$V = Y \cap V'$ 是仿射的，从而 $U = Y \times_{Y'} U'$ 是仿射的。
由引理 \ref{more-morphisms-lemma-thickening-affine-scheme}，$U'$ 是仿射的。
故 $f'$ 是仿射态射。更一般的陈述亦见引理
\ref{more-morphisms-lemma-thicken-property-morphisms}。

\medskip\noindent
关于 (\ref{more-morphisms-item-finite-type})。由说明 (I)，问题化为证明：若
$A \to B$ 有限型，则 $A' \to B'$ 有限型。设
$x_1, \ldots, x_n \in B'$ 在 $B$ 中的像作为 $A$-代数生成 $B$。
% P05-ERR-CH37-0014: Frozen source repeats the surjectivity to B and appears to omit the intended conclusion for B'; retained pending canon review.
则 $A'[x_1, \ldots, x_n] \to B$ 是满射，因为
$A'[x_1, \ldots, x_n] \to B$ 与
$I \otimes_R A[x_1, \ldots, x_n] \to I \otimes_R B$ 都是满射。
更一般的陈述亦见引理 \ref{more-morphisms-lemma-thicken-property-morphisms-cartesian}。

\medskip\noindent
关于 (\ref{more-morphisms-item-quasi-finite})。由 (\ref{more-morphisms-item-finite-type}) 以及有限型态射的
拟有限性可在纤维上检验这一事实得出；见 态射章引理
\ref{morphisms-lemma-quasi-finite-at-point-characterize}。
更一般的陈述亦见引理 \ref{more-morphisms-lemma-thicken-property-morphisms-cartesian}。

\medskip\noindent
关于 (\ref{more-morphisms-item-finite-presentation})。由说明 (I)，问题化为证明：若
$A \to B$ 有限表示，则 $A' \to B'$ 有限表示。由
(\ref{more-morphisms-item-finite-type})，可假设对某个理想 $K'$ 有
$B' = A'[x_1, \ldots, x_n]/K'$。得到短正合列
$$
0 \to K' \to A'[x_1, \ldots, x_n] \to B' \to 0
$$
由于 $B'$ 在 $R'$ 上平坦，$K' \otimes_{R'} R$ 是满射
$A[x_1, \ldots, x_n] \to B$ 的核。由关于 $A \to B$ 的假设，
存在有限多个 $f'_1, \ldots, f'_m \in K'$，其在
$A[x_1, \ldots, x_n]$ 中的像生成此核。由于 $I$ 幂零，
由 Nakayama 引理可知 $f'_1, \ldots, f'_m$ 生成 $K'$；见
代数章引理 \ref{algebra-lemma-NAK}。

\medskip\noindent
关于 (\ref{more-morphisms-item-relative-dimension-d})。由 (\ref{more-morphisms-item-finite-type})
与一般性说明 (II) 得出。更一般的陈述亦见引理
\ref{more-morphisms-lemma-thicken-property-morphisms-cartesian}。

\medskip\noindent
关于 (\ref{more-morphisms-item-universally-open})。由一般性说明 (II) 立即可得。
更一般的陈述亦见引理 \ref{more-morphisms-lemma-thicken-property-morphisms}。

\medskip\noindent
关于 (\ref{more-morphisms-item-syntomic})。假设 $f$ 是 syntomic 态射。由
(\ref{more-morphisms-item-finite-presentation})，$f'$ 局部有限表示；由一般性说明 (III)，
$f'$ 平坦，且 $f'$ 的纤维就是 $f$ 的纤维。因此由
态射章引理 \ref{morphisms-lemma-syntomic-flat-fibres}，
$f'$ 是 syntomic 态射。

\medskip\noindent
关于 (\ref{more-morphisms-item-smooth})。假设 $f$ 光滑。由
(\ref{more-morphisms-item-finite-presentation})，$f'$ 局部有限表示；由一般性说明 (III)，
$f'$ 平坦，且 $f'$ 的纤维就是 $f$ 的纤维。因此由
态射章引理 \ref{morphisms-lemma-smooth-flat-smooth-fibres}，
$f'$ 光滑。

\medskip\noindent
关于 (\ref{more-morphisms-item-unramified})。假设 $f$ 非分歧。由
(\ref{more-morphisms-item-finite-type})，$f'$ 局部有限型，且 $f'$ 的纤维就是 $f$ 的纤维。
因此由 态射章引理
\ref{morphisms-lemma-unramified-etale-fibres}，$f'$ 非分歧。
更一般的陈述亦见引理 \ref{more-morphisms-lemma-thicken-property-morphisms-cartesian}。

\medskip\noindent
关于 (\ref{more-morphisms-item-etale})。假设 $f$ 平展。由
(\ref{more-morphisms-item-finite-presentation})，$f'$ 局部有限表示；由一般性说明 (III)，
$f'$ 平坦，且 $f'$ 的纤维就是 $f$ 的纤维。因此由
态射章引理 \ref{morphisms-lemma-etale-flat-etale-fibres}，
$f'$ 平展。

\medskip\noindent
关于 (\ref{more-morphisms-item-proper})。合并 (\ref{more-morphisms-item-separated})、
(\ref{more-morphisms-item-finite-type})、(\ref{more-morphisms-item-quasi-compact}) 与
(\ref{more-morphisms-item-universally-closed}) 即得。更一般的陈述亦见引理
\ref{more-morphisms-lemma-thicken-property-morphisms-cartesian}。

\medskip\noindent
关于 (\ref{more-morphisms-item-integral})。将 (\ref{more-morphisms-item-universally-closed})、
(\ref{more-morphisms-item-affine}) 与 态射章引理
\ref{morphisms-lemma-integral-universally-closed} 合并。更一般的陈述亦见引理
\ref{more-morphisms-lemma-thicken-property-morphisms}。

\medskip\noindent
关于 (\ref{more-morphisms-item-finite})。将 (\ref{more-morphisms-item-integral})、
(\ref{more-morphisms-item-finite-type}) 与 态射章引理
\ref{morphisms-lemma-finite-integral} 合并。更一般的陈述亦见引理
\ref{more-morphisms-lemma-thicken-property-morphisms-cartesian}。

\medskip\noindent
关于 (\ref{more-morphisms-item-finite-locally-free})。假设 $f$ 有限局部自由。
由 (\ref{more-morphisms-item-finite})，$f'$ 有限；由一般性说明 (III)，$f'$ 平坦；
由 (\ref{more-morphisms-item-finite-presentation})，$f'$ 局部有限表示。因此由
态射章引理 \ref{morphisms-lemma-finite-flat}，
$f'$ 有限局部自由。
\end{proof}

\noindent
下面的引理是“逐纤维平坦性判据”（``crit\`ere de platitude par fibres''）
的“局部幂零”版本；见节 \ref{more-morphisms-section-criterion-flat-fibres}。

\begin{lemma}
\label{more-morphisms-lemma-flatness-morphism-thickenings-fp-over-ft}
考虑交换图
$$
\xymatrix{
(X \subset X') \ar[rr]_{(f, f')} \ar[rd] & & (Y \subset Y') \ar[ld] \\
& (S \subset S')
}
$$
，其中各对象为加厚。假设
\begin{enumerate}
\item $Y' \to S'$ 局部有限型；
\item $X' \to S'$ 平坦且局部有限表示；
\item $f$ 平坦；
\item $X = S \times_{S'} X'$ 且 $Y = S \times_{S'} Y'$。
\end{enumerate}
则 $f'$ 平坦，并且对 $f'$ 的像中的每个 $y' \in Y'$，局部环
$\mathcal{O}_{Y', y'}$ 在 $\mathcal{O}_{S', s'}$ 上平坦且本质有限表示。
\end{lemma}

\begin{proof}
这是 代数章引理
\ref{algebra-lemma-criterion-flatness-fibre-locally-nilpotent} 的直接推论。
\end{proof}

\noindent
概形态射的许多性质在上述引理那样的平坦形变下保持。

\begin{lemma}
\label{more-morphisms-lemma-deform-property-fp-over-ft}
考虑交换图
$$
\xymatrix{
(X \subset X') \ar[rr]_{(f, f')} \ar[rd] & & (Y \subset Y') \ar[ld] \\
& (S \subset S')
}
$$
，其中各对象为加厚。假设 $Y' \to S'$ 局部有限型，
$X' \to S'$ 平坦且局部有限表示，$X = S \times_{S'} X'$，
且 $Y = S \times_{S'} Y'$。则
\begin{enumerate}
\item $f$ 平坦，当且仅当 $f'$ 平坦；
\label{more-morphisms-item-flat-fp-over-ft}
\item $f$ 是同构，当且仅当 $f'$ 是同构；
\label{more-morphisms-item-isomorphism-fp-over-ft}
\item $f$ 是开浸入，当且仅当 $f'$ 是开浸入；
\label{more-morphisms-item-open-immersion-fp-over-ft}
\item $f$ 拟紧，当且仅当 $f'$ 拟紧；
\label{more-morphisms-item-quasi-compact-fp-over-ft}
\item $f$ 泛闭，当且仅当 $f'$ 泛闭；
\label{more-morphisms-item-universally-closed-fp-over-ft}
\item $f$（拟）分离，当且仅当 $f'$（拟）分离；
\label{more-morphisms-item-separated-fp-over-ft}
\item $f$ 是单态射，当且仅当 $f'$ 是单态射；
\label{more-morphisms-item-monomorphism-fp-over-ft}
\item $f$ 是满射，当且仅当 $f'$ 是满射；
\label{more-morphisms-item-surjective-fp-over-ft}
\item $f$ 泛单射，当且仅当 $f'$ 泛单射；
\label{more-morphisms-item-universally-injective-fp-over-ft}
\item $f$ 是仿射态射，当且仅当 $f'$ 是仿射态射；
\label{more-morphisms-item-affine-fp-over-ft}
\item $f$ 局部拟有限，当且仅当 $f'$ 局部拟有限；
\label{more-morphisms-item-quasi-finite-fp-over-ft}
\item
\label{more-morphisms-item-relative-dimension-d-fp-over-ft}
$f$ 局部有限型且相对维数为 $d$，当且仅当 $f'$ 局部有限型且相对维数为 $d$；
\item $f$ 泛开，当且仅当 $f'$ 泛开；
\label{more-morphisms-item-universally-open-fp-over-ft}
\item $f$ 是 syntomic 态射，当且仅当 $f'$ 是 syntomic 态射；
\label{more-morphisms-item-syntomic-fp-over-ft}
\item $f$ 光滑，当且仅当 $f'$ 光滑；
\label{more-morphisms-item-smooth-fp-over-ft}
\item $f$ 非分歧，当且仅当 $f'$ 非分歧；
\label{more-morphisms-item-unramified-fp-over-ft}
\item $f$ 平展，当且仅当 $f'$ 平展；
\label{more-morphisms-item-etale-fp-over-ft}
\item $f$ 固有，当且仅当 $f'$ 固有；
\label{more-morphisms-item-proper-fp-over-ft}
\item $f$ 是有限态射，当且仅当 $f'$ 是有限态射；
\label{more-morphisms-item-finite-fp-over-ft}
\item
\label{more-morphisms-item-finite-locally-free-fp-over-ft}
$f$ 是有限局部自由态射（秩为 $d$），当且仅当 $f'$ 是有限局部自由态射
（秩为 $d$）；
\item 此处还应添加更多性质。
\end{enumerate}
\end{lemma}

\begin{proof}
关于 $X$ 与 $Y$ 的假设意味着，$f$ 是 $f'$ 经 $X \to X'$ 的基变换。
上述 (1)--(20) 中列出的性质 $\mathcal{P}$ 在基变换下都保持，
因此若 $f'$ 具有性质 $\mathcal{P}$，则 $f$ 也具有。见
概形章引理 \ref{schemes-lemma-base-change-immersion},
\ref{schemes-lemma-quasi-compact-preserved-base-change},
\ref{schemes-lemma-separated-permanence}, 及
\ref{schemes-lemma-base-change-monomorphism}
及
态射章引理
\ref{morphisms-lemma-base-change-surjective},
\ref{morphisms-lemma-base-change-universally-injective},
\ref{morphisms-lemma-base-change-affine},
\ref{morphisms-lemma-base-change-quasi-finite},
\ref{morphisms-lemma-base-change-relative-dimension-d},
\ref{morphisms-lemma-base-change-syntomic},
\ref{morphisms-lemma-base-change-smooth},
\ref{morphisms-lemma-base-change-unramified},
\ref{morphisms-lemma-base-change-etale},
\ref{morphisms-lemma-base-change-proper},
\ref{morphisms-lemma-base-change-finite}, 及
\ref{morphisms-lemma-base-change-finite-locally-free}.

\medskip\noindent
因此每种情形中有趣的方向都是：假设 $f$ 具有该性质，并推出 $f'$ 也具有。
先作几条一般性说明，以下论证中将不再逐次提及。
(I) 设 $W' \subset S'$ 是仿射开集，并设 $U' \subset X'$ 与
$V' \subset Y'$ 是位于 $W'$ 上方、满足 $f'(U') \subset V'$ 的仿射开集。
令 $W' = \Spec(R')$，并以 $I \subset R'$ 表示定义闭子概形
$W' \cap S$ 的理想。写 $U' = \Spec(B')$ 且 $V' = \Spec(A')$。
于是得到交换图
$$
\xymatrix{
0 \ar[r] &
IB' \ar[r] &
B' \ar[r] &
B \ar[r] & 0 \\
0 \ar[r] &
IA' \ar[r] \ar[u] &
A' \ar[r] \ar[u] &
A \ar[r] \ar[u] & 0
}
$$
，其行正合。
(II) 态射 $X \to X'$ 与 $Y \to Y'$ 是泛同胚。因此映射 $f$ 与 $f'$
（在任意基变换后）的拓扑相同。
(III) 若 $f$ 平坦，则 $f'$ 平坦，且 $Y' \to S'$ 在 $f'$ 的像中的每一点处
都平坦；见引理 \ref{more-morphisms-lemma-flatness-morphism-thickenings}。

\medskip\noindent
关于 (\ref{more-morphisms-item-flat-fp-over-ft})。这就是一般性说明 (III)。

\medskip\noindent
关于 (\ref{more-morphisms-item-isomorphism-fp-over-ft})。假设 $f$ 是同构。
选择仿射开集 $V' \subset Y'$，并令 $U' = (f')^{-1}(V')$。
于是 $V = Y \cap V'$ 是仿射的，从而
$V \cong f^{-1}(V) = U = Y \times_{Y'} U'$ 是仿射的。由引理
\ref{more-morphisms-lemma-thickening-affine-scheme}，$U'$ 是仿射的。
因此有一般性说明 (I) 中那样的图。由 代数章引理
\ref{algebra-lemma-isomorphism-modulo-locally-nilpotent}，
$A' \to B'$ 是同构，即 $U' \cong V'$。故 $f'$ 是同构。

\medskip\noindent
关于 (\ref{more-morphisms-item-open-immersion-fp-over-ft})。假设 $f$ 是开浸入。
则 $f$ 给出 $X$ 与开子概形 $V \subset Y$ 之间的同构。
令 $V' \subset Y'$ 为底拓扑空间等于 $V$ 的开子概形。于是 $f'$
是从 $X'$ 到 $V'$ 的映射，由 (\ref{more-morphisms-item-isomorphism-fp-over-ft})
它是同构。故 $f'$ 是开浸入。

\medskip\noindent
关于 (\ref{more-morphisms-item-quasi-compact-fp-over-ft})。由说明 (II) 立即可得。
更一般的陈述亦见引理 \ref{more-morphisms-lemma-thicken-property-morphisms}。

\medskip\noindent
关于 (\ref{more-morphisms-item-universally-closed-fp-over-ft})。由说明 (II) 立即可得。
更一般的陈述亦见引理 \ref{more-morphisms-lemma-thicken-property-morphisms}。

\medskip\noindent
关于 (\ref{more-morphisms-item-separated-fp-over-ft})。注意
$X \times_Y X = Y \times_{Y'} (X' \times_{Y'} X')$，故
$X' \times_{Y'} X'$ 是 $X \times_Y X$ 的一个加厚。
映射 $\Delta_{X/Y}$ 与 $\Delta_{X'/Y'}$
的拓扑相符，故得证。更一般的陈述亦见引理
\ref{more-morphisms-lemma-thicken-property-morphisms}。

\medskip\noindent
关于 (\ref{more-morphisms-item-monomorphism-fp-over-ft})。假设 $f$ 是单态射。
考虑对角态射 $\Delta_{X'/Y'} : X' \to X' \times_{Y'} X'$。
注意 $X' \times_{Y'} X' \to S'$ 局部有限型。
$\Delta_{X'/Y'}$ 经 $S \to S'$ 的基变换为 $\Delta_{X/Y}$，
后者按假设是同构。由 (\ref{more-morphisms-item-isomorphism-fp-over-ft})，
$\Delta_{X'/Y'}$ 是同构。

\medskip\noindent
关于 (\ref{more-morphisms-item-surjective-fp-over-ft})。这是显然的。更一般的陈述亦见引理
\ref{more-morphisms-lemma-thicken-property-morphisms}。

\medskip\noindent
关于 (\ref{more-morphisms-item-universally-injective-fp-over-ft})。由说明 (II) 立即可得。
更一般的陈述亦见引理 \ref{more-morphisms-lemma-thicken-property-morphisms}。

\medskip\noindent
关于 (\ref{more-morphisms-item-affine-fp-over-ft})。假设 $f$ 是仿射态射。
选择仿射开集 $V' \subset Y'$，并令 $U' = (f')^{-1}(V')$。
于是 $V = Y \cap V'$ 是仿射的，从而 $U = Y \times_{Y'} U'$ 是仿射的。
由引理 \ref{more-morphisms-lemma-thickening-affine-scheme}，$U'$ 是仿射的。
故 $f'$ 是仿射态射。更一般的陈述亦见引理
\ref{more-morphisms-lemma-thicken-property-morphisms}。

\medskip\noindent
关于 (\ref{more-morphisms-item-quasi-finite-fp-over-ft})。这是因为 $f'$ 局部有限型
（由 态射章引理 \ref{morphisms-lemma-permanence-finite-type}），
而有限型态射的拟有限性可在纤维上检验；见 态射章引理
\ref{morphisms-lemma-quasi-finite-at-point-characterize}。

\medskip\noindent
关于 (\ref{more-morphisms-item-relative-dimension-d-fp-over-ft})。由一般性说明 (II)
以及 $f'$ 局部有限型这一事实得出（态射章引理
\ref{morphisms-lemma-permanence-finite-type}）。

\medskip\noindent
关于 (\ref{more-morphisms-item-universally-open-fp-over-ft})。由一般性说明 (II) 立即可得。
更一般的陈述亦见引理 \ref{more-morphisms-lemma-thicken-property-morphisms}。

\medskip\noindent
关于 (\ref{more-morphisms-item-syntomic-fp-over-ft})。假设 $f$ 是 syntomic 态射。
由 态射章引理 \ref{morphisms-lemma-finite-presentation-permanence}，
$f'$ 局部有限表示。由一般性说明 (III)，$f'$ 平坦。
$f'$ 的纤维就是 $f$ 的纤维。因此由 态射章引理
\ref{morphisms-lemma-syntomic-flat-fibres}，$f'$ 是 syntomic 态射。

\medskip\noindent
关于 (\ref{more-morphisms-item-smooth-fp-over-ft})。假设 $f$ 光滑。由
态射章引理 \ref{morphisms-lemma-finite-presentation-permanence}，
$f'$ 局部有限表示。由一般性说明 (III)，$f'$ 平坦。
$f'$ 的纤维就是 $f$ 的纤维。因此由 态射章引理
\ref{morphisms-lemma-smooth-flat-smooth-fibres}，$f'$ 光滑。

\medskip\noindent
关于 (\ref{more-morphisms-item-unramified-fp-over-ft})。假设 $f$ 非分歧。由
态射章引理 \ref{morphisms-lemma-permanence-finite-type}，
$f'$ 局部有限型。$f'$ 的纤维就是 $f$ 的纤维。因此由
态射章引理 \ref{morphisms-lemma-unramified-etale-fibres}，
$f'$ 非分歧。

\medskip\noindent
关于 (\ref{more-morphisms-item-etale-fp-over-ft})。假设 $f$ 平展。由
态射章引理 \ref{morphisms-lemma-finite-presentation-permanence}，
$f'$ 局部有限表示。由一般性说明 (III)，$f'$ 平坦。
$f'$ 的纤维就是 $f$ 的纤维。因此由 态射章引理
\ref{morphisms-lemma-etale-flat-etale-fibres}，$f'$ 平展。

\medskip\noindent
关于 (\ref{more-morphisms-item-proper-fp-over-ft})。合并
(\ref{more-morphisms-item-separated-fp-over-ft})、$f$ 局部有限型这一事实
（态射章引理 \ref{morphisms-lemma-permanence-finite-type}）、
(\ref{more-morphisms-item-quasi-compact-fp-over-ft}) 与
(\ref{more-morphisms-item-universally-closed-fp-over-ft}) 即得。

\medskip\noindent
关于 (\ref{more-morphisms-item-finite-fp-over-ft})。合并
(\ref{more-morphisms-item-universally-closed-fp-over-ft})、
(\ref{more-morphisms-item-affine-fp-over-ft})、态射章引理
\ref{morphisms-lemma-integral-universally-closed}、$f$ 局部有限型这一事实
（态射章引理 \ref{morphisms-lemma-permanence-finite-type}）以及
态射章引理 \ref{morphisms-lemma-finite-integral}。

\medskip\noindent
关于 (\ref{more-morphisms-item-finite-locally-free-fp-over-ft})。
假设 $f$ 有限局部自由。由 (\ref{more-morphisms-item-finite-fp-over-ft})，$f'$ 有限。
由一般性说明 (III)，$f'$ 平坦。由 态射章引理
\ref{morphisms-lemma-finite-presentation-permanence}，$f'$ 局部有限表示。
因此由 态射章引理 \ref{morphisms-lemma-finite-flat}，
$f'$ 有限局部自由。
\end{proof}

\begin{lemma}[射影概形的形变]
\label{more-morphisms-lemma-deform-projective}
设 $f : X \to S$ 是固有、平坦且有限表示的概形态射，
$\mathcal{L}$ 是 $f$-丰沛的，并假设 $S$ 拟紧。则存在 $d_0 \geq 0$，
使得对每个笛卡尔图
$$
\vcenter{
\xymatrix{
X \ar[r]_{i'} \ar[d]_f & X' \ar[d]^{f'} \\
S \ar[r]^i & S'
}
}
\quad\text{且}\quad
\begin{matrix}
\text{invertible }\mathcal{O}_{X'}\text{-模}\\
\mathcal{L}'\text{ 其中 }\mathcal{L} \cong (i')^*\mathcal{L}'
\end{matrix}
$$
，其中 $S \subset S'$ 是加厚且 $f'$ 固有、平坦、有限表示，都有
\begin{enumerate}
\item $R^p(f')_*(\mathcal{L}')^{\otimes d} = 0$
对所有 $p > 0$ 与 $d \geq d_0$ 成立；
\item $\mathcal{A}'_d = (f')_*(\mathcal{L}')^{\otimes d}$
在 $d \geq d_0$ 时有限局部自由；
\item $\mathcal{A}' =
\mathcal{O}_{S'} \oplus \bigoplus_{d \geq d_0} \mathcal{A}'_d$
是有限表示的拟凝聚 $\mathcal{O}_{S'}$-代数；
\item 存在典范同构
$r' : X' \to \underline{\text{Proj}}_{S'}(\mathcal{A}')$；
\item 存在典范同构
$\theta' : (r')^*\mathcal{O}_{\underline{\text{Proj}}_{S'}(\mathcal{A}')}(1)
\to \mathcal{L}'$.
\end{enumerate}
$\mathcal{A}'$、$r'$、$\theta'$ 的构造关于数据
$(X', S', i, i', f', \mathcal{L}')$ 是函子性的。
\end{lemma}

\begin{proof}
先描述映射 $r'$ 与 $\theta'$。注意 $\mathcal{L}'$ 是 $f'$-丰沛的；
见引理 \ref{more-morphisms-lemma-thicken-property-relatively-ample}。
有拟凝聚分次 $\mathcal{O}_{S'}$-代数的典范映射
$\mathcal{A}' \to \bigoplus_{d \geq 0} (f')_*(\mathcal{L}')^{\otimes d}$
它在次数 $\geq d_0$ 上为同构。因此，它在相对 Proj 上诱导同构，
并与结构层的 Serre 扭曲相容；见 构造章引理
\ref{constructions-lemma-eventual-iso-graded-rings-map-relative-proj}。
故由 态射章引理
\ref{morphisms-lemma-characterize-relatively-ample} 得到态射 $r'$
（该引理又援用了 构造章引理
\ref{constructions-lemma-invertible-map-into-relative-proj} 的构造），
而由 态射章引理 \ref{morphisms-lemma-proper-ample-is-proj}，
它是同构。映射 $\theta'$ 来自 构造章引理
\ref{constructions-lemma-invertible-map-into-relative-proj}。
由 性质章引理 \ref{properties-lemma-ample-gcd-is-one}，
$\theta'$ 是同构（这里还使用了如下事实：在 $S'$ 的各仿射开集上，
态射 $r'$ 与 性质章引理 \ref{properties-lemma-map-into-proj}
中的态射相同；这已在 态射章引理
\ref{morphisms-lemma-proper-ample-is-proj} 的证明中说明）。

\medskip\noindent
假设 (1)、(2) 中所述的消失与局部自由性成立，则该构造的函子性可如下看出。
设 $h : T \to S'$ 是概形态射，以 $f_T : X'_T \to T$ 表示 $f'$ 的基变换，
并以 $\mathcal{L}_T$ 表示 $\mathcal{L}$ 到 $X'_T$ 的拉回。
% P05-ERR-CH37-0015: Frozen source says pullback of L although the deformation datum uses L'; retained pending canon review.
由上同调与基变换（例如 概形的导出范畴章引理
\ref{perfect-lemma-compare-base-change} 的表述），在 $T$ 上有相应的消失，并且
$h^*\mathcal{A}'_d = f_{T, *}\mathcal{L}_T^{\otimes d}$
（因而也有推出层的局部自由性以及相应分次
$\mathcal{O}_T$-代数 $\mathcal{A}_T$ 的有限生成性）。因此态射
$r_T : X_T \to
\underline{\text{Proj}}_T(\bigoplus f_{T, *}\mathcal{L}_T^{\otimes d})$
% P05-ERR-CH37-0016: Frozen source writes X_T after introducing X'_T; formula retained pending canon review.
就是 $r'$ 到 $T$ 的基变换，而 $\theta'$ 的拉回就是映射 $\theta_T$。

\medskip\noindent
综上，只需证明 (1)、(2)、(3)。选择 $d_0$，使对所有
$d \geq d_0$ 与 $p > 0$ 都有
$R^pf_*\mathcal{L}^{\otimes d} = 0$；见 概形上同调章引理
\ref{coherent-lemma-coherent-proper-ample}。我们断言此 $d_0$ 可用。

\medskip\noindent
由上同调与基变换（概形的导出范畴章引理
\ref{perfect-lemma-flat-proper-perfect-direct-image-general}），
$E'_d = Rf'_*(\mathcal{L}')^{\otimes d}$ 是
$D(\mathcal{O}_{S'})$ 的完美对象，且其构造与任意基变换可交换。
特别地，$E_d = Li^*E'_d = Rf_*\mathcal{L}^{\otimes d}$。
由 概形的导出范畴章引理
\ref{perfect-lemma-vanishing-implies-locally-free}，当 $d \geq d_0$ 时，
复形 $E_d$ 同构于置于上同调次数 $0$ 的有限局部自由
$\mathcal{O}_S$-模 $f_*\mathcal{L}^{\otimes d}$。再由
概形的导出范畴章引理
\ref{perfect-lemma-open-where-cohomology-in-degree-i-rank-r}，
$E'_d$ 同构于置于上同调次数 $0$ 的有限局部自由模。
这当然表示 $E'_d = \mathcal{A}'_d[0]$，对 $p > 0$ 有
$R^pf'_*(\mathcal{L}')^{\otimes d} = 0$，且 $\mathcal{A}'_d$
有限局部自由。这证明了 (1) 与 (2)。

\medskip\noindent
最后还须证明，$\mathcal{A}'$ 作为 $\mathcal{O}_{S'}$-代数层有限表示
（此概念在 Properties, Section
\ref{properties-section-extending-quasi-coherent-sheaves} 中引入）。
设 $U' = \Spec(R') \subset S'$ 是仿射开集。则
$A' = \mathcal{A}'(U')$ 是分次 $R'$-代数，其各分次部分都是有限投射
$R'$-模。须证明 $A'$ 是有限表示的 $R'$-代数。
我们将把它归结到 Noether 情形。具体而言，可以找到有限型
$\mathbf{Z}$-子代数 $R'_0 \subset R'$ 以及 $R'_0$ 上的一对
\footnote{它具有与 $X' \to S'$ 及 $\mathcal{L}'$ 相同的性质；即
$X'_0 \to \Spec(R'_0)$ 平坦且固有，而 $\mathcal{L}'_0$ 丰沛。}
$(X'_0, \mathcal{L}'_0)$，其基变换为
$(X'_{U'}, \mathcal{L}'|_{X'_{U'}})$；见
极限章引理
\ref{limits-lemma-descend-modules-finite-presentation},
\ref{limits-lemma-descend-invertible-modules},
\ref{limits-lemma-eventually-proper},
\ref{limits-lemma-descend-flat-finite-presentation}, 及
\ref{limits-lemma-limit-ample}.
概形上同调章引理 \ref{coherent-lemma-coherent-proper-ample}
推出
$A'_0 = \bigoplus_{d \geq 0} H^0(X'_0, (\mathcal{L}'_0)^{\otimes d})$
是有限生成的分次 $R'_0$-代数，并推出存在 $d'_0$，使得
$H^p(X'_0, (\mathcal{L}'_0)^{\otimes d}) = 0$, $p > 0$ for $d \geq d'_0$.
。将上述论证应用于 $X'_0 \to \Spec(R'_0)$ 与 $\mathcal{L}'_0$，
可知 $(A'_0)_d$ 是有限投射 $R'_0$-模，并且
$$
A'_d = \mathcal{A}'_d(U') =
H^0(X'_{U'}, (\mathcal{L}')^{\otimes d}|_{X'_{U'}}) =
H^0(X'_0, (\mathcal{L}'_0)^{\otimes d}) \otimes_{R'_0} R' =
(A'_0)_d \otimes_{R'_0} R'
$$
对 $d \geq d'_0$ 成立。论证中的一个小转折是：我们不知道能否选择
$d'_0$ 等于 $d_0$\footnote{事实上，通过更多极限论证可以归结到这一情形。}。
为绕开这一点，用以下一系列论证完成证明：
\begin{enumerate}
\item[(a)] 代数
$B = R'_0 \oplus \bigoplus_{d \geq \max(d_0, d'_0)} (A'_0)_d$
是有限型 $R'_0$-代数：对 $B \subset A'_0$ 应用 Artin--Tate 引理；见
代数章引理 \ref{algebra-lemma-Artin-Tate}。
\item[(b)] 由于 $R'_0$ 是 Noether 环，$B$ 是有限表示的 $R'_0$-代数。
\item[(c)] 由张量积的右正合性，$B \otimes_{R'_0} R'$
是有限表示的 $R'$-代数。
\item[(d)] 由上述等式，这恰好表示
$C = R' \oplus \bigoplus_{d \geq \max(d_0, d'_0)} A'_d$
是有限表示的 $R'$-代数。
\item[(e)] 商 $A'/C$ 是有限投射 $R'$-模 $A'_d$（其中
% P05-ERR-CH37-0017: Frozen source includes the upper endpoint although C already contains that degree; retained pending canon review.
$d_0 \leq d \leq \max(d_0, d'_0)$）的直和，故作为 $R'$-模有限表示。
\item[(f)] 由 代数章引理
\ref{algebra-lemma-finitely-presented-over-subring}，商 $A'/C$
作为 $C$-模有限表示。
\item[(g)] 因此由 代数章引理 \ref{algebra-lemma-extension}，
$A'$ 作为 $C$-模有限表示。
\item[(h)] 由 代数章引理 \ref{algebra-lemma-finite-finite-type}，
这推出 $A'$ 作为 $C$-代数有限表示。
\item[(i)] 最后，将 代数章引理
\ref{algebra-lemma-compose-finite-type} 应用于 $R' \to C \to A'$，
这推出 $A'$ 作为 $R'$-代数有限表示。
\end{enumerate}
证明完毕。
\end{proof}









\section{形式光滑态射}
\label{more-morphisms-section-formally-smooth}

\noindent
Michael Artin 对光滑性的微分判据（例如 态射章引理
\ref{morphisms-lemma-smooth-at-point}）的看法是，它们（在实践中）基本无用。
本节引入形式光滑态射 $X \to S$ 的概念。这种态射由如下性质刻画：
只要 $T$ 仿射，$X$ 的 $T$-值点就能提升到 $T$ 的无穷小加厚上。
主要结果是：形式光滑且局部有限表示的态射是光滑的；见引理
\ref{more-morphisms-lemma-smooth-formally-smooth}。事实证明，这一判据往往比上述微分判据
更容易使用。

\medskip\noindent
回顾：若对每个实线部分交换的图
$$
\xymatrix{
A \ar[r] \ar@{-->}[rd] & B/I \\
R \ar[r] \ar[u] & B \ar[u]
}
$$
（其中 $I \subset B$ 是平方为零的理想），存在一条使该图交换的虚线箭头，
则称环映射 $R \to A$ 是{\it 形式光滑的}；见
代数章定义 \ref{algebra-definition-formally-smooth}。
这引出概形态射的如下类似概念。

\begin{definition}
\label{more-morphisms-definition-formally-smooth}
设 $f : X \to S$ 是概形态射。若给定任意实线部分交换的图
$$
\xymatrix{
X \ar[d]_f & T \ar[d]^i \ar[l] \\
S & T' \ar[l] \ar@{-->}[lu]
}
$$
其中 $T \subset T'$ 是 $S$ 上仿射概形的一阶加厚时，存在一条使该图交换的
虚线箭头，则称 $f$ 是{\it 形式光滑的}。
\end{definition}

\noindent
对形式非分歧与形式平展态射，可以去掉 $T'$ 为仿射的条件；见引理
\ref{more-morphisms-lemma-formally-unramified-not-affine} 与
\ref{more-morphisms-lemma-formally-etale-not-affine}。对形式光滑态射则不再如此。
事实上，一个稍微更自然的条件是：应能在 $T'$ 上 Zariski 局部地补出虚线箭头。
分析引理 \ref{more-morphisms-lemma-formally-smooth} 的证明可知，这与现有定义等价。
特别地，形式光滑性在源上是 Zariski 局部的（事实上它在源上是光滑局部的，
% P05-ERR-CH37-0018: Frozen source contains the literal placeholder ``insert future reference here''.
此处插入将来的参考文献）。

\begin{lemma}
\label{more-morphisms-lemma-composition-formally-smooth}
形式光滑态射的复合仍是形式光滑的。
\end{lemma}

\begin{proof}
从略。
\end{proof}

\begin{lemma}
\label{more-morphisms-lemma-base-change-formally-smooth}
形式光滑态射的基变换仍是形式光滑的。
\end{lemma}

\begin{proof}
从略；其代数版本见 代数章引理 \ref{algebra-lemma-base-change-fs}。
\end{proof}

\begin{lemma}
\label{more-morphisms-lemma-formally-etale-unramified-smooth}
设 $f : X \to S$ 是概形态射。则 $f$ 形式平展，当且仅当
$f$ 形式光滑且形式非分歧。
\end{lemma}

\begin{proof}
从略。
\end{proof}

\begin{lemma}
\label{more-morphisms-lemma-formally-smooth-on-opens}
设 $f : X \to S$ 是概形态射，并设 $U \subset X$ 与 $V \subset S$
是满足 $f(U) \subset V$ 的开子概形。若 $f$ 形式光滑，
则 $f|_U : U \to V$ 也形式光滑。
\end{lemma}

\begin{proof}
考虑定义 \ref{more-morphisms-definition-formally-smooth} 中那样的实线部分交换图
$$
\xymatrix{
U \ar[d]_{f|_U} & T \ar[d]^i \ar[l]^a \\
V & T' \ar[l] \ar@{-->}[lu]
}
$$
。若 $f$ 形式光滑，则存在满足 $a'|_T = a$ 的 $S$-态射
$a' : T' \to X$。由于 $T$ 与 $T'$ 的底集合相同，$a'$ 是到 $U$ 的态射
（见 概形章第 \ref{schemes-section-open-immersion} 节），
而且它显然也是 $V$-态射。因此得到所需的上述虚线箭头。
\end{proof}

\begin{lemma}
\label{more-morphisms-lemma-affine-formally-smooth}
设 $f : X \to S$ 是概形态射，并假设 $X$ 与 $S$ 均为仿射。
则 $f$ 形式光滑，当且仅当
$\mathcal{O}_S(S) \to \mathcal{O}_X(X)$ 是形式光滑环映射。
\end{lemma}

\begin{proof}
由环范畴与仿射概形范畴之间的等价，这立即从定义
（定义 \ref{more-morphisms-definition-formally-smooth} 与
代数章定义 \ref{algebra-definition-formally-smooth}）得出；见
概形章引理 \ref{schemes-lemma-category-affine-schemes}。
\end{proof}

\noindent
下面的引理是本节的主要结果。它是函子观点的一次胜利：与 Limits,
命题 \ref{limits-proposition-characterize-locally-finite-presentation}
合并后，它表明可用函子 $h_X : \Sch/S \to \textit{Sets}$ 的“简单”性质
来识别态射 $f : X \to S$ 是否光滑。

\begin{lemma}[无穷小提升判据]
\label{more-morphisms-lemma-smooth-formally-smooth}
设 $f : X \to S$ 是概形态射。以下条件等价：
\begin{enumerate}
\item 态射 $f$ 光滑；
\item 态射 $f$ 局部有限表示且形式光滑。
\end{enumerate}
\end{lemma}

\begin{proof}
假设 $f : X \to S$ 局部有限表示且形式光滑。考虑满足
$f(U) \subset V$ 的一对仿射开集 $\Spec(A) = U \subset X$ 与
$\Spec(R) = V \subset S$。由引理 \ref{more-morphisms-lemma-formally-smooth-on-opens}，
$U \to V$ 形式光滑。由引理 \ref{more-morphisms-lemma-affine-formally-smooth}，
$R \to A$ 形式光滑。由 态射章引理
\ref{morphisms-lemma-locally-finite-presentation-characterize}，
$R \to A$ 有限表示。由 代数章命题
\ref{algebra-proposition-smooth-formally-smooth}，$R \to A$ 光滑。
故由光滑态射的定义，$X \to S$ 光滑。

\medskip\noindent
反之，假设 $f : X \to S$ 光滑。考虑实线部分交换图
$$
\xymatrix{
X \ar[d]_f & T \ar[d]^i \ar[l]^a \\
S & T' \ar[l] \ar@{-->}[lu]
}
$$
，如定义 \ref{more-morphisms-definition-formally-smooth} 所述。我们将证明虚线箭头存在，
从而证明 $f$ 形式光滑。

\medskip\noindent
令 $\mathcal{F}$ 为引理 \ref{more-morphisms-lemma-sheaf} 在注记
\ref{more-morphisms-remark-special-case} 所讨论的特殊情形下给出的 $T'$ 上集合层。令
$$
\mathcal{H} =
\SheafHom_{\mathcal{O}_T}(a^*\Omega_{X/S}, \mathcal{C}_{T/T'})
$$
为 $\mathcal{O}_T$-模层，其上有引理 \ref{more-morphisms-lemma-action-sheaf}
中的作用 $\mathcal{H} \times \mathcal{F} \to \mathcal{F}$。
我们的目标就是证明 $\mathcal{F}(T) \not = \emptyset$。换言之，
要证明 $\mathcal{F}$ 是 $T$ 上平凡的 $\mathcal{H}$-挠子（见 Cohomology,
第 \ref{cohomology-section-h1-torsors} 节）。分两步：
(I) 为证明 $\mathcal{F}$ 是挠子，须对所有 $t \in T$ 证明
$\mathcal{F}_t \not = \emptyset$（见 上同调章定义
\ref{cohomology-definition-torsor}）。
(II) 为证明 $\mathcal{F}$ 是平凡挠子，只需证明
$H^1(T, \mathcal{H}) = 0$（见 上同调章引理
\ref{cohomology-lemma-torsors-h1}——可把 $\mathcal{H}$ 看作阿贝尔层或
$\mathcal{O}_T$-模来计算其上同调；见 上同调章引理
\ref{cohomology-lemma-modules-abelian}）。

\medskip\noindent
先证明 (I)。对每个 $t \in T$，可选择 $t$ 的仿射开邻域
$U \subset T$，使 $a(U)$ 包含于仿射开集 $\Spec(A) = W \subset X$，
而后者映到仿射开集 $\Spec(R) = V \subset S$。
由 态射章引理 \ref{morphisms-lemma-smooth-characterize}，
环映射 $R \to A$ 光滑。因此由 代数章命题
\ref{algebra-proposition-smooth-formally-smooth}，环映射 $R \to A$
形式光滑。进而由引理 \ref{more-morphisms-lemma-affine-formally-smooth}，
$W \to V$ 形式光滑。故可将 $a|_U : U \to W$ 提升为 $V$-态射
$a' : U' \to W \subset X$，这表明 $\mathcal{F}(U) \not = \emptyset$。

\medskip\noindent
最后证明 (II)。由 态射章引理
\ref{morphisms-lemma-finite-presentation-differentials}，
$\Omega_{X/S}$ 有限表示（由 态射章引理
\ref{morphisms-lemma-smooth-omega-finite-locally-free}，
它甚至有限局部自由）。因此 $a^*\Omega_{X/S}$ 有限表示（见
模章引理 \ref{modules-lemma-pullback-finite-presentation}）。故层
$\mathcal{H} =
\SheafHom_{\mathcal{O}_T}(a^*\Omega_{X/S}, \mathcal{C}_{T/T'})$
由 概形章第 \ref{schemes-section-quasi-coherent} 节 中的讨论是拟凝聚的。
因此由 概形上同调章引理
\ref{coherent-lemma-quasi-coherent-affine-cohomology-zero}，
有所需的 $H^1(T, \mathcal{H}) = 0$。
\end{proof}

\noindent
局部投射拟凝聚模定义于 Properties, Section
\ref{properties-section-locally-projective}。

\begin{lemma}
\label{more-morphisms-lemma-formally-smooth-sheaf-differentials}
设 $f : X \to Y$ 是形式光滑的概形态射。则 $\Omega_{X/Y}$
在 $X$ 上局部投射。
\end{lemma}

\begin{proof}
选择满足 $f(U) \subset V$ 的仿射开集 $U \subset X$ 与 $V \subset Y$。
由引理 \ref{more-morphisms-lemma-formally-smooth-on-opens}，
$f|_U : U \to V$ 形式光滑。因此由引理
\ref{more-morphisms-lemma-affine-formally-smooth}，环映射
$\Gamma(V, \mathcal{O}_V) \to \Gamma(U, \mathcal{O}_U)$ 形式光滑。
故由 代数章引理 \ref{algebra-lemma-characterize-formally-smooth-again}，
$\Gamma(U, \mathcal{O}_U)$-模
$\Omega_{\Gamma(U, \mathcal{O}_U)/\Gamma(V, \mathcal{O}_V)}$
是投射的。因此 $\Omega_{U/V}$ 局部投射；见 Properties, Section
\ref{properties-section-locally-projective}。
\end{proof}

\begin{lemma}
\label{more-morphisms-lemma-h1-is-zero}
设 $T$ 是仿射概形，$\mathcal{F}$、$\mathcal{G}$ 是拟凝聚
$\mathcal{O}_T$-模。考虑
$\mathcal{H} = \SheafHom_{\mathcal{O}_T}(\mathcal{F}, \mathcal{G})$。
若 $\mathcal{F}$ 局部投射，则 $H^1(T, \mathcal{H}) = 0$。
\end{lemma}

\begin{proof}
由概形上局部投射层的定义（见 性质章定义
\ref{properties-definition-locally-projective}），$\mathcal{F}$ 是自由
$\mathcal{O}_T$-模的直和项。因此可假设
$\mathcal{F} = \bigoplus_{i \in I} \mathcal{O}_T$ 是自由模。
此时 $\mathcal{H} = \prod_{i \in I} \mathcal{G}$ 是拟凝聚模的积。
由 上同调章引理 \ref{cohomology-lemma-cohomology-products}，
$H^1 = 0$，因为仿射概形上拟凝聚层的上同调为零；见
概形上同调章引理
\ref{coherent-lemma-quasi-coherent-affine-cohomology-zero}。
\end{proof}

\begin{lemma}
\label{more-morphisms-lemma-formally-smooth}
设 $f : X \to Y$ 是概形态射。以下条件等价：
\begin{enumerate}
\item $f$ 形式光滑；
\item 对每个 $x \in X$，存在满足 $f(U) \subset V$ 的开集
$x \in U \subset X$ 与 $f(x) \in V \subset Y$，使得
$f|_U : U \to V$ 形式光滑；
\item 对每一对满足 $f(U) \subset V$ 的仿射开集 $U \subset X$ 与
$V \subset Y$，环映射 $\mathcal{O}_Y(V) \to \mathcal{O}_X(U)$ 形式光滑；
\item 存在仿射开覆盖 $Y = \bigcup V_j$，且对每个 $j$ 存在仿射开覆盖
$f^{-1}(V_j) = \bigcup U_{ji}$，使得对所有 $j$ 与 $i$，
% P05-ERR-CH37-0019: In this and the analogous formally unramified/etale item (4), frozen source uses undefined V,U instead of the indexed opens.
环映射 $\mathcal{O}_Y(V) \to \mathcal{O}_X(U)$ 形式光滑。
\end{enumerate}
\end{lemma}

\begin{proof}
蕴含 (1) $\Rightarrow$ (2)、(1) $\Rightarrow$ (3) 与
(2) $\Rightarrow$ (4) 来自引理 \ref{more-morphisms-lemma-formally-smooth-on-opens}。
蕴含 (3) $\Rightarrow$ (4) 是立即的。

\medskip\noindent
假设 (4)。证明 $f$ 形式光滑的论证同引理
\ref{more-morphisms-lemma-smooth-formally-smooth} 证明的第二部分。考虑实线部分交换图
$$
\xymatrix{
X \ar[d]_f & T \ar[d]^i \ar[l]^a \\
Y & T' \ar[l] \ar@{-->}[lu]
}
$$
，如定义 \ref{more-morphisms-definition-formally-smooth} 所述。我们将证明虚线箭头存在，
从而证明 $f$ 形式光滑。令 $\mathcal{F}$ 为引理 \ref{more-morphisms-lemma-sheaf}
在注记 \ref{more-morphisms-remark-special-case} 所讨论的特殊情形下给出的
$T'$ 上集合层。令
$$
\mathcal{H} =
\SheafHom_{\mathcal{O}_T}(a^*\Omega_{X/Y}, \mathcal{C}_{T/T'})
$$
为 $T$ 上的 $\mathcal{O}_T$-模层，其上有引理
\ref{more-morphisms-lemma-action-sheaf} 中的作用
$\mathcal{H} \times \mathcal{F} \to \mathcal{F}$。
此作用 $\mathcal{H} \times \mathcal{F} \to \mathcal{F}$
使 $\mathcal{F}$ 成为伪 $\mathcal{H}$-挠子；见 Cohomology,
定义 \ref{cohomology-definition-torsor}。目标是证明
$\mathcal{F}$ 是平凡的 $\mathcal{H}$-挠子。分两步：
(I) 为证明 $\mathcal{F}$ 是挠子，须证明 $\mathcal{F}$ 局部有截面。
(II) 为证明 $\mathcal{F}$ 是平凡挠子，只需证明
$H^1(T, \mathcal{H}) = 0$；见 上同调章引理
\ref{cohomology-lemma-torsors-h1}。

\medskip\noindent
先证明 (I)。对每个 $t \in T$，可选择 $t$ 的仿射开邻域
$W \subset T$，使对某些 $i, j$ 有 $a(W)$ 包含于 $U_{ji}$。
令 $W' \subset T'$ 为相应的开子概形。由假设 (4)，可将
$a|_W : W \to U_{ji}$ 提升为 $V_j$-态射
$a' : W' \to U_{ji}$，这表明 $\mathcal{F}(W')$ 非空。

\medskip\noindent
最后证明 (II)。由引理 \ref{more-morphisms-lemma-formally-smooth-sheaf-differentials}，
$\Omega_{U_{ji}/V_j}$ 局部投射。因此 $\Omega_{X/Y}$ 局部投射；见
性质章引理 \ref{properties-lemma-locally-projective}。
故 $a^*\Omega_{X/Y}$ 局部投射；见 性质章引理
\ref{properties-lemma-locally-projective-pullback}。因此
$$
H^1(T, \mathcal{H}) =
H^1(T, \SheafHom_{\mathcal{O}_T}(a^*\Omega_{X/Y}, \mathcal{C}_{T/T'})) = 0
$$
，由引理 \ref{more-morphisms-lemma-h1-is-zero} 得到所需结论。
\end{proof}

\begin{lemma}
\label{more-morphisms-lemma-triangle-differentials-formally-smooth}
设 $f : X \to Y$、$g : Y \to S$ 是概形态射，并假设 $f$ 形式光滑。则
$$
0 \to f^*\Omega_{Y/S} \to \Omega_{X/S} \to \Omega_{X/Y} \to 0
$$
（见 态射章引理
\ref{morphisms-lemma-triangle-differentials}）是短正合列。
\end{lemma}

\begin{proof}
此引理的代数版本如下：给定环映射 $A \to B \to C$，其中
$B \to C$ 形式光滑，则
$$
0 \to C \otimes_B \Omega_{B/A} \to \Omega_{C/A} \to \Omega_{C/B} \to 0
$$
这一 代数章引理 \ref{algebra-lemma-exact-sequence-differentials}
中的序列正合。这就是 代数章引理
\ref{algebra-lemma-ses-formally-smooth}。
\end{proof}

\begin{lemma}
\label{more-morphisms-lemma-differentials-formally-unramified-formally-smooth}
设 $h : Z \to X$ 是 $S$ 上形式非分歧的概形态射，并假设 $Z$
在 $S$ 上形式光滑。则典范正合列
$$
0 \to \mathcal{C}_{Z/X} \to h^*\Omega_{X/S} \to \Omega_{Z/S} \to 0
$$
（见引理 \ref{more-morphisms-lemma-universally-unramified-differentials-sequence}）
是短正合列。
\end{lemma}

\begin{proof}
设 $Z \to Z'$ 是 $Z$ 在 $X$ 上的泛一阶加厚。由引理
\ref{more-morphisms-lemma-universally-unramified-differentials-sequence} 的证明，
我们的序列与序列
$$
\mathcal{C}_{Z/Z'} \to \Omega_{Z'/S} \otimes \mathcal{O}_Z \to
\Omega_{Z/S} \to 0.
$$
等同。由于 $Z \to S$ 形式光滑，可在 $Z'$ 上局部地找到包含映射
$Z \to Z'$ 的一个 $S$ 上左逆 $Z' \to Z$。因此该序列局部分裂；见
态射章引理
\ref{morphisms-lemma-differentials-relative-immersion-section}。
\end{proof}

\begin{lemma}
\label{more-morphisms-lemma-two-unramified-morphisms-formally-smooth}
设
$$
\xymatrix{
Z \ar[r]_i \ar[rd]_j & X \ar[d]^f \\
& Y
}
$$
是概形的交换图，其中 $i$ 与 $j$ 形式非分歧，且 $f$ 形式光滑。
则典范正合列
$$
0 \to
\mathcal{C}_{Z/Y} \to
\mathcal{C}_{Z/X} \to
i^*\Omega_{X/Y} \to 0
$$
（见引理 \ref{more-morphisms-lemma-two-unramified-morphisms}）正合且局部分裂。
\end{lemma}

\begin{proof}
以 $Z \to Z'$ 表示 $Z$ 在 $X$ 上的泛一阶加厚，以 $Z \to Z''$
表示 $Z$ 在 $Y$ 上的泛一阶加厚。
% P05-ERR-CH37-0020: Frozen source again cites the differential-sequence lemma for the functorial map Z' -> Z''.
由引理 \ref{more-morphisms-lemma-universally-unramified-differentials-sequence}，
有典范态射 $Z' \to Z''$，从而有交换图
$$
\xymatrix{
Z \ar[r]_{i'} \ar[rd]_{j'} & Z' \ar[r]_a \ar[d]^k & X \ar[d]^f \\
& Z'' \ar[r]^b & Y
}
$$
在引理 \ref{more-morphisms-lemma-two-unramified-morphisms} 的证明中，
我们已将上述序列与序列
$$
\mathcal{C}_{Z/Z''} \to
\mathcal{C}_{Z/Z'} \to
(i')^*\Omega_{Z'/Z''} \to 0
$$
等同。设 $U'' \subset Z''$ 是仿射开集，以 $U \subset Z$ 与
$U' \subset Z'$ 表示相应的仿射开子概形。由于 $f$ 形式光滑，存在态射
$h : U'' \to X$，它在 $U$ 上与 $i$ 相同，且 $f \circ h$ 等于
$b|_{U''}$。由于 $Z'$ 是泛一阶加厚，得到唯一的态射
% P05-ERR-CH37-0021: Frozen source writes the ill-typed equality g = a \circ h; formula retained pending canon review.
$g : U'' \to Z'$，满足 $g = a \circ h$。$Z''$ 的泛性质推出
$k \circ g$ 是包含映射 $U'' \to Z''$。故 $g$ 是 $k$ 的左逆。图示为
$$
\xymatrix{
U \ar[d] \ar[r] & Z' \ar[d]^k \\
U'' \ar[r] \ar[ru]^g & Z''
}
$$
因此 $g$ 诱导映射
$\mathcal{C}_{Z/Z'}|_U \to \mathcal{C}_{Z/Z''}|_U$，
它是 $U$ 上映射 $\mathcal{C}_{Z/Z''} \to \mathcal{C}_{Z/Z'}$ 的左逆。
\end{proof}



































\section{Noether 基底上的光滑性}
\label{more-morphisms-section-smooth-Noetherian}

\noindent
事实证明，若基底是 Noether 的，则在形式光滑性的表述中可以少要求一些。
在某种意义上，下面这些引理是形变理论的开端。

\begin{lemma}
\label{more-morphisms-lemma-lifting-along-artinian-at-point}
设 $f : X \to S$ 是概形态射，$x \in X$。假设 $S$ 局部 Noether，
且 $f$ 局部有限型。以下条件等价：
\begin{enumerate}
\item $f$ 在 $x$ 处光滑；
\item 对每个实线部分交换的图
$$
\xymatrix{
X \ar[d]_f & \Spec(B) \ar[d]^i \ar[l]^-\alpha \\
S & \Spec(B') \ar[l]_-{\beta} \ar@{-->}[lu]
}
$$
，其中 $B' \to B$ 是局部环的满射，$\Ker(B' \to B)$ 平方为零，
且 $\alpha$ 把 $\Spec(B)$ 的闭点映到 $x$，存在一条使该图交换的虚线箭头；
\item 与 (2) 相同，但 $B' \to B$ 遍历小扩张（见 代数章定义
\ref{algebra-definition-small-extension}）；
\item 与 (2) 相同，但 $B' \to B$ 遍历如下小扩张：$\alpha$ 诱导同构
$\kappa(x) \to \kappa(\mathfrak m)$，其中
$\mathfrak m \subset B$ 是极大理想。
\end{enumerate}
\end{lemma}

\begin{proof}
选择 $f(x)$ 的仿射邻域 $V \subset S$，并选择 $x$ 的仿射邻域
$U \subset X$，使得 $f(U) \subset V$。对 (2) 中任意“测试”图，
态射 $\alpha$ 都把 $\Spec(B)$ 映入 $U$，而态射 $\beta$ 把
$\Spec(B')$ 映入 $V$（见 Schemes, Section
\ref{schemes-section-points}）。故引理归结为仿射概形之间的态射
$f|_U : U \to V$。（事实上，$V$ 是 Noether 的，且 $f|_U$ 有限型；见
性质章引理 \ref{properties-lemma-locally-Noetherian} 与
态射章引理 \ref{morphisms-lemma-locally-finite-type-characterize}。）
在此仿射情形下，引理与 代数章引理
\ref{algebra-lemma-smooth-test-artinian} 完全相同。
\end{proof}

\noindent
有时知道如下事实很有用：只需对以有限型点为“中心”的小扩张检验提升判据
（见 态射章第 \ref{morphisms-section-points-finite-type} 节）。

\begin{lemma}
\label{more-morphisms-lemma-lifting-along-artinian}
设 $f : X \to S$ 是概形态射。假设 $S$ 局部 Noether，且 $f$ 局部有限型。
以下条件等价：
\begin{enumerate}
\item $f$ 光滑；
\item 对每个实线部分交换的图
$$
\xymatrix{
X \ar[d]_f & \Spec(B) \ar[d]^i \ar[l]^-\alpha \\
S & \Spec(B') \ar[l]_-{\beta} \ar@{-->}[lu]
}
$$
，其中 $B' \to B$ 是 Artin 局部环的小扩张，且 $\beta$ 有限型（！），
存在一条使该图交换的虚线箭头。
\end{enumerate}
\end{lemma}

\begin{proof}
若 $f$ 光滑，则无穷小提升判据（引理
\ref{more-morphisms-lemma-smooth-formally-smooth}）表明 $f$ 形式光滑，故 (2) 成立。

\medskip\noindent
假设 (2)。$f$ 不光滑的点 $x \in X$ 构成 $X$ 的闭子集 $T$。
由 态射章第 \ref{morphisms-section-points-finite-type} 节 中的讨论，
若 $T \not = \emptyset$，则存在点 $x \in T \subset X$，使态射
$$
\Spec(\kappa(x)) \to X \to S
$$
有限型（具体地，在 $T$ 中任取一个在 $X$ 的某个仿射开集内为闭点的点 $x$）。
由 态射章引理 \ref{morphisms-lemma-artinian-finite-type}，
给定剩余域为 $\kappa(x)$ 的任意 Artin 局部环 $B'$，任意态射
$\beta : \Spec(B') \to S$ 都有限型。因此，引理
\ref{more-morphisms-lemma-lifting-along-artinian-at-point} (4) 中使用的所有图都对应于
本引理 (2) 中的图。故 $X \to S$ 在 $x$ 处光滑，矛盾。
% P05-ERR-CH37-0022: Frozen source has the typo ``smooth a x''; translated as intended while preserving a non-rendering erratum marker.
\end{proof}

\noindent
下面是一个有用的应用。

\begin{lemma}
\label{more-morphisms-lemma-check-smoothness-on-infinitesimal-nbhds}
设 $f : X \to S$ 是局部 Noether 概形之间的有限型态射。
设 $Z \subset S$ 是闭子概形，其 $n$ 阶无穷小邻域为
$Z_n \subset S$。令 $X_n = Z_n \times_S X$。
\begin{enumerate}
\item 若对所有 $n$，$X_n \to Z_n$ 都光滑，则 $f$ 在
$f^{-1}(Z)$ 的每一点处光滑。
\item 若对所有 $n$，$X_n \to Z_n$ 都平展，则 $f$ 在
$f^{-1}(Z)$ 的每一点处平展。
\end{enumerate}
\end{lemma}

\begin{proof}
假设对所有 $n$，$X_n \to Z_n$ 都光滑。设 $x \in X$ 位于 $Z$ 的某一点上方。
给定引理 \ref{more-morphisms-lemma-lifting-along-artinian-at-point} (3) 中那样的小扩张
$B' \to B$ 与态射 $\alpha$、$\beta$，$B'$ 的极大理想幂零（因为 $B'$
是 Artin 的），故对适当的 $n$，态射 $\beta$ 经由 $Z_n$ 分解，而
$\alpha$ 经由 $X_n$ 分解。于是应用 $X_n \to Z_n$ 的提升性质，得到图中
所需的虚线箭头。这证明了 (1)。(2) 由 (1) 以及如下事实得出：一个态射平展，
当且仅当它光滑且相对维数为 $0$。
\end{proof}

\begin{lemma}
\label{more-morphisms-lemma-check-flatness-on-infinitesimal-nbhds}
设 $f : X \to S$ 是局部 Noether 概形之间的态射。设 $Z \subset S$
是闭子概形，其 $n$ 阶无穷小邻域为 $Z_n \subset S$。令
$X_n = Z_n \times_S X$。若对所有 $n$，$X_n \to Z_n$ 都平坦，
则 $f$ 在 $f^{-1}(Z)$ 的每一点处平坦。
\end{lemma}

\begin{proof}
这就是将 代数章引理 \ref{algebra-lemma-flat-module-powers}
翻译成概形的语言。
\end{proof}








\section{朴素余切复形}
\label{more-morphisms-section-netherlander}

\noindent
本节延续 模章第 \ref{modules-section-netherlander} 节，
而后者又延续 代数章第 \ref{algebra-section-netherlander} 节 中的讨论。

\begin{definition}
\label{more-morphisms-definition-netherlander}
设 $f : X \to Y$ 是概形态射。$f$ 的{\it 朴素余切复形}是 Modules,
定义 \ref{modules-definition-cotangent-complex-morphism-ringed-topoi}
中定义的复形。记作 $\NL_f$ 或 $\NL_{X/Y}$。
\end{definition}

\begin{lemma}
\label{more-morphisms-lemma-NL-affine}
设 $f : X \to Y$ 是概形态射。设 $\Spec(A) = U \subset X$ 与
$\Spec(R) = V \subset S$ 是满足 $f(U) \subset V$ 的仿射开集。
存在复形的典范映射
$$
\widetilde{\NL_{A/R}} \longrightarrow \NL_{X/Y}|_U
$$
，它在 $D(\mathcal{O}_U)$ 中是同构。
% P05-ERR-CH37-0023: Frozen source names V as an open of S although the morphism target is Y.
\end{lemma}

\begin{proof}
由 模章第 \ref{modules-section-netherlander} 节 中
$\NL_{X/Y}$ 的构造，有 $A = \mathcal{O}_X(U)$-模复形的典范映射\linebreak[4]
$\NL_{\mathcal{O}_X(U)/f^{-1}\mathcal{O}_Y(U)} \to \NL_{X/Y}(U)$，
并且它与进一步的限制相容。使用典范映射
$R \to f^{-1}\mathcal{O}_Y(U)$，得到 $A$-模复形的典范映射
$\NL_{A/R} \to \NL_{\mathcal{O}_X(U)/f^{-1}\mathcal{O}_Y(U)}$。
使用 $\widetilde{\ }$ 函子的泛性质（见 概形章引理
\ref{schemes-lemma-compare-constructions}），得到引理所述的映射。
要检验此映射在上同调层上为同构，只需检验它在茎上诱导同构。
这由 代数章引理 \ref{algebra-lemma-NL-localize-bottom}、
\ref{algebra-lemma-localize-NL} 与 模章引理
\ref{modules-lemma-stalk-NL} 得出（这里还使用了如下描述：在点
$\mathfrak p \in \Spec(A)$ 处，$\mathcal{O}_X$ 与
$f^{-1}\mathcal{O}_Y$ 的茎分别是 $A_\mathfrak p$ 与 $R_\mathfrak q$，
其中 $\mathfrak q = R \cap \mathfrak p$；所用参考为 概形章引理
\ref{schemes-lemma-spec-sheaves} 与 层章引理
\ref{sheaves-lemma-stalk-pullback}）。
\end{proof}

\begin{lemma}
\label{more-morphisms-lemma-netherlander-quasi-coherent}
设 $f : X \to Y$ 是概形态射。复形 $\NL_{X/Y}$ 的上同调层拟凝聚，
在次数 $-1$、$0$ 之外为零，并且在次数 $0$ 等于 $\Omega_{X/Y}$。
\end{lemma}

\begin{proof}
由 模章第 \ref{modules-section-netherlander} 节 中朴素余切复形的构造，
$\NL_{X/Y}$ 位于次数 $-1$、$0$，且其次数 $0$ 上同调为 $\Omega_{X/Y}$。
微分层拟凝聚（由 态射章引理
\ref{morphisms-lemma-differentials-diagonal}）。因此只需证明
$H^{-1}(\NL_{X/Y})$ 拟凝聚。使用引理 \ref{more-morphisms-lemma-NL-affine}
在各仿射开集上检验即可。
\end{proof}

\begin{lemma}
\label{more-morphisms-lemma-netherlander-fp}
设 $f : X \to Y$ 是概形态射。若 $f$ 局部有限表示，则在 $X$ 上局部地，
$\NL_{X/Y}$ 拟同构于复形
$$
\ldots \to 0 \to \mathcal{F}^{-1} \to \mathcal{F}^0 \to 0 \to \ldots
$$
，其中各项为拟凝聚 $\mathcal{O}_X$-模，$\mathcal{F}^0$ 有限表示，
而 $\mathcal{F}^{-1}$ 有限型。
\end{lemma}

\begin{proof}
由引理 \ref{more-morphisms-lemma-NL-affine}，只需证明当 $R \to A$ 是有限表示环映射时，
$\NL_{A/R}$ 具有这种形状。写
$A = R[x_1, \ldots, x_n]/I$，其中 $I$ 有限生成。则 $I/I^2$
是有限 $A$-模，且 $\NL_{A/R}$ 拟同构于
$$
\ldots \to 0 \to I/I^2 \to
\bigoplus\nolimits_{i = 1, \ldots, n} A\text{d}x_i \to 0 \to \ldots
$$
；这由 代数章第 \ref{algebra-section-netherlander} 节，
特别是 代数章引理 \ref{algebra-lemma-NL-homotopy} 得出。
\end{proof}

\begin{lemma}
\label{more-morphisms-lemma-NL-formally-smooth}
设 $f : X \to Y$ 是概形态射。以下条件等价：
\begin{enumerate}
\item $f$ 形式光滑；
\item $H^{-1}(\NL_{X/Y}) = 0$ 及 $H^0(\NL_{X/Y}) = \Omega_{X/Y}$
局部投射。
\end{enumerate}
\end{lemma}

\begin{proof}
这由 代数章命题
\ref{algebra-proposition-characterize-formally-smooth} 与引理
\ref{more-morphisms-lemma-formally-smooth} 得出。
\end{proof}

\begin{lemma}
\label{more-morphisms-lemma-NL-formally-etale}
设 $f : X \to Y$ 是概形态射。以下条件等价：
\begin{enumerate}
\item $f$ 形式平展；
\item $H^{-1}(\NL_{X/Y}) = H^0(\NL_{X/Y}) = 0$.
\end{enumerate}
\end{lemma}

\begin{proof}
形式平展态射形式光滑，故由引理 \ref{more-morphisms-lemma-NL-formally-smooth}，
$H^{-1}(\NL_{X/Y}) = 0$。另一方面，由引理
\ref{more-morphisms-lemma-characterize-formally-etale}，$\Omega_{X/Y} = 0$。
反之，若 (2) 成立，则由引理 \ref{more-morphisms-lemma-NL-formally-smooth}，
$f$ 形式光滑；由引理 \ref{more-morphisms-lemma-formally-unramified-differentials}
又知其形式非分歧；从而由引理
\ref{more-morphisms-lemma-formally-etale-unramified-smooth}，它形式平展。
\end{proof}

\begin{lemma}
\label{more-morphisms-lemma-NL-smooth}
设 $f : X \to Y$ 是概形态射。以下条件等价：
\begin{enumerate}
\item $f$ 光滑；
\item $f$ 局部有限表示，
$H^{-1}(\NL_{X/Y}) = 0$, 及 $H^0(\NL_{X/Y}) = \Omega_{X/Y}$
有限局部自由。
\end{enumerate}
\end{lemma}

\begin{proof}
这由光滑环同态的定义（代数章定义
\ref{algebra-definition-smooth}）、引理 \ref{more-morphisms-lemma-NL-affine}
以及光滑概形态射的定义（态射章定义
\ref{morphisms-definition-smooth}）得出。这里还使用了如下事实：
对环上的模而言，有限局部自由等价于有限投射（代数章引理
\ref{algebra-lemma-finite-projective}）。
\end{proof}

\begin{lemma}
\label{more-morphisms-lemma-NL-etale}
设 $f : X \to Y$ 是概形态射。以下条件等价：
\begin{enumerate}
\item $f$ 平展；
\item $f$ 局部有限表示，且
$H^{-1}(\NL_{X/Y}) = H^0(\NL_{X/Y}) = 0$.
\end{enumerate}
\end{lemma}

\begin{proof}
这由平展环同态的定义（代数章定义
\ref{algebra-definition-etale}）、引理 \ref{more-morphisms-lemma-NL-affine}
以及平展概形态射的定义（态射章定义
\ref{morphisms-definition-etale}）得出。
\end{proof}


\begin{lemma}
\label{more-morphisms-lemma-NL-immersion}
设 $i : Z \to X$ 是概形的浸入。则在 $D(\mathcal{O}_Z)$ 中，
$\NL_{Z/X}$ 同构于 $\mathcal{C}_{Z/X}[1]$，其中
$\mathcal{C}_{Z/X}$ 是 $Z$ 在 $X$ 中的余法层。
\end{lemma}

\begin{proof}
这由 代数章引理 \ref{algebra-lemma-NL-surjection}、
态射章引理 \ref{morphisms-lemma-affine-conormal} 与引理
\ref{more-morphisms-lemma-NL-affine} 得出。
\end{proof}

\begin{lemma}
\label{more-morphisms-lemma-exact-sequence-NL}
设 $f : X \to Y$ 与 $g : Y \to Z$ 是概形态射。则有上同调层的典范六项正合列
$$
H^{-1}(f^*\NL_{Y/Z}) \to
H^{-1}(\NL_{X/Z}) \to
H^{-1}(\NL_{X/Y}) \to
f^*\Omega_{Y/Z} \to \Omega_{X/Z} \to \Omega_{X/Y} \to 0
$$
。
\end{lemma}

\begin{proof}
这是 模章引理
\ref{modules-lemma-exact-sequence-NL-ringed-topoi} 的特殊情形。
\end{proof}

\begin{lemma}
\label{more-morphisms-lemma-get-triangle-NL}
设 $f : X \to Y$ 与 $Y \to Z$ 是概形态射。假设 $X \to Y$
是完全交态射。则在 $D(\mathcal{O}_X)$ 中有典范可区别三角
$$
f^*\NL_{Y/Z} \to \NL_{X/Z} \to \NL_{X/Y} \to f^*\NL_{Y/Z}[1]
$$
，它恢复引理 \ref{more-morphisms-lemma-exact-sequence-NL} 的 $6$ 项正合列。
\end{lemma}

\begin{proof}
只需证明 模章引理
\ref{modules-lemma-exact-sequence-NL-ringed-topoi} 的典范映射
$$
f^*\NL_{Y/Z} \to \text{Cone}(\NL_{X/Z} \to \NL_{X/Y})[-1]
$$
在 $D(\mathcal{O}_X)$ 中是同构。为此，只需证明该 $6$ 项序列左端为零，
即 $H^{-1}(f^*\NL_{Y/Z}) \to H^{-1}(\NL_{X/Z})$ 是单射。
在仿射局部，这由 代数进阶章引理
\ref{more-algebra-lemma-transitive-lci-at-end} 中相应的代数结果得出。
使用引理 \ref{more-morphisms-lemma-NL-affine} 将其翻译成代数即可。
\end{proof}

\begin{lemma}
\label{more-morphisms-lemma-smooth-etale-permanence}
设 $X \to Y \to Z$ 是概形态射。假设 $X \to Z$ 光滑，
且 $Y \to Z$ 平展。则 $X \to Y$ 光滑。
\end{lemma}

\begin{proof}
由 态射章引理
\ref{morphisms-lemma-finite-presentation-permanence}，态射 $X \to Y$
局部有限表示。由引理 \ref{more-morphisms-lemma-NL-smooth}，
$H^{-1}(\NL_{X/Z}) = 0$，且模 $\Omega_{X/Z}$ 有限局部自由。
由引理 \ref{more-morphisms-lemma-NL-etale}，
$H^{-1}(\NL_{Y/Z}) = H^0(\NL_{Y/Z}) = 0$。
由引理 \ref{more-morphisms-lemma-exact-sequence-NL}，得到
$H^{-1}(\NL_{X/Y}) = 0$，且
$\Omega_{X/Y} \cong \Omega_{X/Z}$ 有限局部自由。
由引理 \ref{more-morphisms-lemma-NL-smooth}，态射 $X \to Y$ 光滑。
\end{proof}

\begin{lemma}
\label{more-morphisms-lemma-get-NL}
设概形态射 $f : X \to Y$ 分解为 $f = g \circ i$，其中 $i$ 是浸入，
而 $g : P \to Y$ 形式光滑（例如光滑）。则在 $D(\mathcal{O}_X)$ 中有典范同构
$$
\NL_{X/Y} \cong \left(\mathcal{C}_{X/P} \to i^*\Omega_{P/Y}\right)
$$
，其中余法层 $\mathcal{C}_{X/P}$ 置于次数 $-1$。
\end{lemma}

\begin{proof}
（括号中的断言见引理 \ref{more-morphisms-lemma-smooth-formally-smooth}。）
由引理 \ref{more-morphisms-lemma-NL-immersion} 与 \ref{more-morphisms-lemma-NL-formally-smooth}，有
$\NL_{X/P} = \mathcal{C}_{X/P}[1]$ 与 $\NL_{P/Y} = \Omega_{P/Y}$，
其中 $\Omega_{P/Y}$ 局部投射。这还推出
$i^*\NL_{P/Y} \to i^*\Omega_{P/Y}$ 是拟同构（略去一个小细节；
理由是 $i^*\NL_{P/Y}$ 与 $\tau_{\geq -1}Li^*\NL_{P/Y}$ 相同；见
代数进阶章引理 \ref{more-algebra-lemma-tensor-NL}）。
因此 模章引理
\ref{modules-lemma-exact-sequence-NL-ringed-topoi} 的典范映射
$$
i^*\NL_{P/Y} \to \text{Cone}(\NL_{X/Y} \to \NL_{X/P})[-1]
$$
在 $D(\mathcal{O}_X)$ 中是同构，因为由上可知上同调群
$H^{-1}(i^*\NL_{P/Y})$ 为零。换言之，有可区别三角
$$
i^*\NL_{P/Y} \to \NL_{X/Y} \to \NL_{X/P} \to i^*\NL_{P/Y}[1]
$$
显然，这表示 $\NL_{X/Y}$ 是映射
$\NL_{X/P}[-1] \to i^*\NL_{P/Y}$ 的锥。根据上面对这些导出范畴对象的
上同调层的计算，这等价于引理的断言。
\end{proof}

\begin{lemma}
\label{more-morphisms-lemma-base-change-NL}
考虑概形的笛卡尔图
$$
\xymatrix{
X' \ar[r]_{g'} \ar[d] & X \ar[d] \\
Y' \ar[r] & Y
}
$$
典范映射 $(g')^*\NL_{X/Y} \to \NL_{X'/Y'}$ 在 $H^0$ 上诱导同构，
在 $H^{-1}$ 上诱导满射。
\end{lemma}

\begin{proof}
翻译成代数后，这就是 代数进阶章引理
\ref{more-algebra-lemma-base-change-NL}。用引理 \ref{more-morphisms-lemma-NL-affine}
完成翻译。
\end{proof}

\begin{lemma}
\label{more-morphisms-lemma-flat-base-change-NL}
考虑概形的笛卡尔图
$$
\xymatrix{
X' \ar[d] \ar[r]_{g'} & X \ar[d] \\
Y' \ar[r] & Y
}
$$
若 $Y' \to Y$ 平坦，则典范映射
$(g')^*\NL_{X/Y} \to \NL_{X'/Y'}$ 是拟同构。
\end{lemma}

\begin{proof}
由引理 \ref{more-morphisms-lemma-NL-affine}，这从 代数章引理
\ref{algebra-lemma-change-base-NL} 得出。
\end{proof}

\begin{lemma}
\label{more-morphisms-lemma-base-change-NL-flat}
考虑概形的笛卡尔图
$$
\xymatrix{
X' \ar[r]_{g'} \ar[d] & X \ar[d] \\
Y' \ar[r] & Y
}
$$
若 $X \to Y$ 平坦，则典范映射
$(g')^*\NL_{X/Y} \to \NL_{X'/Y'}$ 是拟同构。
若此外 $\NL_{X/Y}$ 的 Tor 振幅包含于 $[-1, 0]$，则
$L(g')^*\NL_{X/Y} \to \NL_{X'/Y'}$ 也是拟同构。
\end{lemma}

\begin{proof}
翻译成代数后，这就是 代数进阶章引理
\ref{more-algebra-lemma-base-change-NL-flat}。使用引理
\ref{more-morphisms-lemma-NL-affine} 与 概形的导出范畴章引理
\ref{perfect-lemma-affine-compare-bounded}、
\ref{perfect-lemma-tor-dimension-affine} 完成翻译。
\end{proof}











\section{概形范畴中的推出，I}
\label{more-morphisms-section-pushouts}

\noindent
本节构造 $Y \leftarrow X \rightarrow X'$ 的推出，其中 $X \to Y$ 仿射，
而 $X \to X'$ 是加厚。事实上，这对我们是一个重要情形，值得详细讨论。
在节 \ref{more-morphisms-section-pushouts-II} 中将讨论一个更有趣也更困难的情形。
关于任意范畴中推出的一般讨论，见 Categories, Section
\ref{categories-section-pushouts}。

\begin{lemma}
\label{more-morphisms-lemma-basic-example-pushout}
设 $A' \to A$ 是环的满射，$B \to A$ 是环映射，并设
$B' = B \times_A A'$ 为环的纤维积。令
$S = \Spec(A)$、$S' = \Spec(A')$、$T = \Spec(B)$ 且
$T' = \Spec(B')$。则
$$
\vcenter{
\xymatrix{
S \ar[r]_i \ar[d]_f & S' \ar[d]^{f'} \\
T \ar[r]^{i'} & T'
}
}
\quad\text{corresponding to}\quad
\vcenter{
\xymatrix{
A & A' \ar[l] \\
B \ar[u] & B' \ar[l] \ar[u]
}
}
$$
是概形的推出。
\end{lemma}

\begin{proof}
由 代数进阶章引理
\ref{more-algebra-lemma-points-of-fibre-product}，作为拓扑空间有
$T' = T \amalg_S S'$；也就是说，该图在拓扑空间范畴中是推出。
接着考虑映射
$$
((i')^\sharp, (f')^\sharp) :
\mathcal{O}_{T'}
\longrightarrow
i'_*\mathcal{O}_T \times_{g_*\mathcal{O}_S} f'_*\mathcal{O}_{S'}
$$
，其中 $g = i' \circ f = f' \circ i$。我们断言此映射是环层的同构。
事实上，可把两边都看成拟凝聚 $\mathcal{O}_{T'}$-模（右边使用
概形章引理 \ref{schemes-lemma-push-forward-quasi-coherent}），
且该映射为 $\mathcal{O}_{T'}$-线性的。因此只需证明它在整体截面层面上
是同构（概形章引理
\ref{schemes-lemma-equivalence-quasi-coherent}）。在整体截面上，
恢复引理陈述中的等同 $B' \to B \times_A A'$（这正是选取 $B'$ 的方式）。

\medskip\noindent
设 $X$ 是概形，并假设给定满足 $m' \circ i = n \circ f$ 的概形态射
$m' : S' \to X$ 与 $n : T \to X$（将此复合记为 $m$）。由于
$T' = T \amalg_S S'$（见上文），得到与 $m'$、$n$ 相容的唯一拓扑空间映射
$n' : T' \to X$。由上一段对 $\mathcal{O}_{T'}$ 的描述，
得到环层的唯一同态
$$
(n')^\sharp :
\mathcal{O}_X
\longrightarrow
(n')_*\mathcal{O}_{T'} =
m'_*\mathcal{O}_T \times_{m_*\mathcal{O}_T} n_*\mathcal{O}_S
$$
，它由 $(m')^\sharp$ 与 $n^\sharp$ 给出。因此
$(n', (n')^\sharp)$ 是与 $m'$、$n$ 相容的唯一环化空间态射
$T' \to X$。为完成证明，只需证明 $n'$ 是概形态射，
即局部环化空间的态射。

\medskip\noindent
设 $t' \in T'$ 的像为 $x \in X$。须证明
$\mathcal{O}_{X, x} \to \mathcal{O}_{T', t'}$ 是局部同态。
若 $t' \not \in T$，则 $t'$ 是唯一一点 $s' \in S'$ 的像，且
$\mathcal{O}_{T', t'} = \mathcal{O}_{S', s'}$。事实上，
$S' \setminus S \to T' \setminus T$ 是概形的同构，因为
$B' \to A'$ 诱导同构
$\Ker(B' \to B) = \Ker(A' \to A)$。若 $t'$ 是 $t \in T$ 的像，
则已知复合
$\mathcal{O}_{X, x} \to \mathcal{O}_{T', t'} \to \mathcal{O}_{T, t}$
是局部同态，故同样得到结论。
\end{proof}

\begin{lemma}
\label{more-morphisms-lemma-pushout-fpqc-local}
设 $\mathcal{I} \to (\Sch/S)_{fppf}$、$i \mapsto X_i$ 是概形的图，
而 $(W, X_i \to W)$ 是该图在概形范畴中的余锥（范畴章说明
\ref{categories-remark-cones-and-cocones}）。若存在概形的 fpqc 覆盖
$\{W_a \to W\}_{a \in A}$，使得
\begin{enumerate}
\item 对所有 $a \in A$，在概形范畴中有
$W_a = \colim X_i \times_W W_a$
；
\item 对所有 $a, b \in A$，在概形范畴中有
$W_a \times_W W_b = \colim X_i \times_W W_a \times_W W_b$
；
\end{enumerate}
则在概形范畴中 $W = \colim X_i$。
\end{lemma}

\begin{proof}
事实上，对概形 $T$，一个态射 $W \to T$ 等价于一族态射
$W_a \to T$、$a \in A$，它们在交叠 $W_a \times_W W_b$ 上相等；见
下降章引理 \ref{descent-lemma-fpqc-universal-effective-epimorphisms}。
\end{proof}

\begin{lemma}
\label{more-morphisms-lemma-pushout-along-thickening}
设 $X \to X'$ 是概形的加厚，而 $X \to Y$ 是仿射概形态射。
则在概形范畴中存在推出
$$
\xymatrix{
X \ar[r] \ar[d]_f
&
X' \ar[d]^{f'}
\\
Y \ar[r]
&
Y'
}
$$
。此外，$Y \subset Y'$ 是加厚，$X = Y \times_{Y'} X'$，并且在
$|Y| = |Y'|$ 上有层的等式
$$
\mathcal{O}_{Y'} = \mathcal{O}_Y \times_{f_*\mathcal{O}_X} f'_*\mathcal{O}_{X'}
$$
。
\end{lemma}

\begin{proof}
先将 $Y'$ 构造为环化空间。具体地，作为拓扑空间取 $Y' = Y$。
以 $f' : X' \to Y'$ 表示等于 $f$ 的拓扑空间映射。取引理等式的右边作为
结构层 $\mathcal{O}_{Y'}$。为证明 $Y'$ 是概形，须证明每一点都有仿射邻域。
由于层的纤维积构造与限制到开集可交换，可假设 $Y$ 仿射。
于是 $X$ 仿射（因为 $f$ 仿射），且 $X'$ 也仿射（见引理
\ref{more-morphisms-lemma-thickening-affine-scheme}）。设
$Y \leftarrow X \rightarrow X'$ 对应于 $B \rightarrow A \leftarrow A'$。
令 $B' = B \times_A A'$；这就是 $\mathcal{O}_{Y'}$ 的整体截面。
由于 $A' \to A$ 是核局部幂零的满射，$B' \to B$ 也是核局部幂零的满射。
故作为拓扑空间，$\Spec(B') = \Spec(B)$。我们断言
$Y' = \Spec(B')$。为此，将证明：对 $g' \in B'$ 及其像 $g \in B$，有
$\mathcal{O}_{Y'}(D(g)) = B'_{g'}$。事实上，由 代数进阶章引理
\ref{more-algebra-lemma-diagram-localize} 有
$$
(B')_{g'} = B_g \times_{A_h} A'_{h'}
$$
，其中 $h \in A$、$h' \in A'$ 是 $g'$ 的像。由于 $B_g$、分别地
$A_h$、分别地 $A'_{h'}$ 等于 $\mathcal{O}_Y(D(g))$、分别地
$f_*\mathcal{O}_X(D(g))$、分别地 $f'_*\mathcal{O}_{X'}(D(g))$，
断言得证。

\medskip\noindent
还须证明 $Y'$ 是推出。上述讨论表明，概形 $Y'$ 有仿射开覆盖
$Y' = \bigcup W'_i$，使相应的开集 $U'_i \subset X'$、$W_i \subset Y$
与 $U_i \subset X$ 都仿射。此外，若 $A'_i$、$B_i$、$A_i$
分别是对应于 $U'_i$、$W_i$、$U_i$ 的环，则 $W'_i$ 对应于
$B_i \times_{A_i} A'_i$。因此可应用引理
\ref{more-morphisms-lemma-basic-example-pushout} 与 \ref{more-morphisms-lemma-pushout-fpqc-local}，
得出我们的构造在概形范畴中是推出。
\end{proof}

\noindent
在下面的引理中，使用 范畴章例
\ref{categories-example-2-fibre-product-categories} 中定义的范畴纤维积。

\begin{lemma}
\label{more-morphisms-lemma-equivalence-categories-schemes-over-pushout}
设 $X \to X'$ 是概形的加厚，$X \to Y$ 是仿射概形态射，并设
$Y' = Y \amalg_X X'$ 为推出（见引理
\ref{more-morphisms-lemma-pushout-along-thickening}）。基变换给出函子
$$
F : (\Sch/Y') \longrightarrow (\Sch/Y) \times_{(\Sch/X)} (\Sch/X')
$$
，它由 $V' \longmapsto (V' \times_{Y'} Y, V' \times_{Y'} X', 1)$ 给出，
并有左伴随函子
$$
G : (\Sch/Y) \times_{(\Sch/X)} (\Sch/X') \longrightarrow (\Sch/Y')
$$
，它把三元组 $(V, U', \varphi)$ 送到推出
$V \amalg_{(V \times_Y X)} U'$。最后，$F \circ G$ 同构于恒等函子。
\end{lemma}

\begin{proof}
设 $(V, U', \varphi)$ 是纤维积范畴的对象。令
$U = U' \times_{X'} X$。注意 $U \to U'$ 是加厚。由于
$\varphi : V \times_Y X \to U' \times_{X'} X = U$ 是同构，
有位于 $X \to Y$ 上方的态射 $U \to V$，它将 $U$ 等同于纤维积
$X \times_Y V$。特别地，$U \to V$ 仿射；见 态射章引理
\ref{morphisms-lemma-base-change-affine}。因此可应用引理
\ref{more-morphisms-lemma-pushout-along-thickening} 得到推出 $V' = V \amalg_U U'$。
由于 $V'$ 是推出，而且给定了在 $U$ 上作为到 $Y'$ 的态射相等的态射
$V \to Y$ 与 $U' \to X'$，遂得到态射 $V' \to Y'$。
令 $G(V, U', \varphi) = V'$，便得到函子 $G$。

\medskip\noindent
下面证明 $G$ 是 $F$ 的左伴随。设 $Z$ 是 $Y'$ 上的概形。须证明
$$
\Mor(V', Z) = \Mor((V, U', \varphi), F(Z))
$$
，其中态射集取自各自的范畴。设 $g' : V' \to Z$ 是态射。
以 $\tilde g$、分别地 $\tilde f'$ 表示 $g'$ 与态射 $V \to V'$、
分别地 $U' \to V'$ 的复合。分别经 $Y \to Y'$、$X' \to Y'$ 对
$\tilde g$、$\tilde f'$ 作基变换，得到态射
$g : V \to Z \times_{Y'} Y$、分别地
$f' : U' \to Z \times_{Y'} X'$。于是 $(g, f')$ 是上述等式右边的元素
（略去细节）。反之，设
$(g, f') : (V, U', \varphi) \to F(Z)$ 是右边的元素。
可考虑 $g$、分别地 $f$ 与 $Z \times_{Y'} X' \to Z$、分别地
$Z \times_{Y'} Y \to Z$ 的复合 $\tilde g : V \to Z$、分别地
$\tilde f' : U' \to Z$。则 $\tilde g$ 与 $\tilde f'$ 作为从 $U$
到 $Z$ 的态射相等。由推出的泛性质，得到态射 $g' : V' \to Z$，
即等式左边的元素。略去这些构造互为逆的验证。
% P05-ERR-CH37-0024: Frozen source swaps f/f' and the two projection targets in the converse composition sentence.

\medskip\noindent
为证明 $F \circ G$ 同构于恒等函子，须证明伴随映射
$(V, U', \varphi) \to F(G(V, U', \varphi))$ 是同构。
为此可在仿射局部运算。写 $X = \Spec(A)$、$X' = \Spec(A')$、
$Y = \Spec(B)$。则 $A' \to A$ 与 $B \to A$ 是 More on Algebra,
引理 \ref{more-algebra-lemma-module-over-fibre-product} 中那样的环映射，
且 $Y' = \Spec(B')$，其中 $B' = B \times_A A'$。再设
$V = \Spec(D)$、$U' = \Spec(C')$，而 $\varphi$ 由 $A$-代数同构
$D \otimes_B A \to C' \otimes_{A'} A = C'/IC'$ 给出。
令 $D' = D \times_{C'/IC'} C'$。在这种情形下，须证明
$D' \otimes_{B'} B \cong D$ 且 $D' \otimes_{B'} A' \cong C'$。
这是 代数进阶章引理
\ref{more-algebra-lemma-module-over-fibre-product} 的特殊情形。
\end{proof}

\begin{lemma}
\label{more-morphisms-lemma-scheme-over-pushout-flat-modules}
设 $X \to X'$ 是概形的加厚，$X \to Y$ 是仿射概形态射，并设
$Y' = Y \amalg_X X'$ 为推出（见引理
\ref{more-morphisms-lemma-pushout-along-thickening}）。设 $V' \to Y'$ 是概形态射。令
$V = Y \times_{Y'} V'$、$U' = X' \times_{Y'} V'$ 且
$U = X \times_{Y'} V'$。以下两者的范畴等价：
\begin{enumerate}
\item 在 $Y'$ 上平坦的拟凝聚 $\mathcal{O}_{V'}$-模；
\item 三元组 $(\mathcal{G}, \mathcal{F}', \varphi)$ 的范畴，其中
\begin{enumerate}
\item $\mathcal{G}$ 是在 $Y$ 上平坦的拟凝聚 $\mathcal{O}_V$-模；
\item $\mathcal{F}'$ 是在 $X'$ 上平坦的拟凝聚 $\mathcal{O}_{U'}$-模；
\item $\varphi : (U \to V)^*\mathcal{G} \to (U \to U')^*\mathcal{F}'$
是 $\mathcal{O}_U$-模的同构。
\end{enumerate}
\end{enumerate}
该等价把 $\mathcal{G}'$ 映到
$((V \to V')^*\mathcal{G}', (U' \to V')^*\mathcal{G}', can)$。
假设 $\mathcal{G}'$ 对应于三元组
$(\mathcal{G}, \mathcal{F}', \varphi)$。则
\begin{enumerate}
\item[(a)] $\mathcal{G}'$ 是有限型 $\mathcal{O}_{V'}$-模，当且仅当
$\mathcal{G}$ 与 $\mathcal{F}'$ 分别是有限型 $\mathcal{O}_Y$-模与
$\mathcal{O}_{U'}$-模。
\item[(b)] 若 $V' \to Y'$ 局部有限表示，则 $\mathcal{G}'$ 是有限表示的
$\mathcal{O}_{V'}$-模，当且仅当 $\mathcal{G}$ 与 $\mathcal{F}'$
分别是有限表示的 $\mathcal{O}_Y$-模与 $\mathcal{O}_{U'}$-模。
\end{enumerate}
\end{lemma}

\begin{proof}
一个拟逆函子把三元组 $(\mathcal{G}, \mathcal{F}', \varphi)$ 映到纤维积
$$
(V \to V')_*\mathcal{G}
\times_{(U \to V')_*\mathcal{F}}
(U' \to V')_*\mathcal{F}'
$$
，其中 $\mathcal{F} = (U \to U')^*\mathcal{F}'$。这是可行的，因为在仿射情形下，
恢复 代数进阶章引理
\ref{more-algebra-lemma-relative-flat-module-over-fibre-product} 的等价。
略去一些细节。

\medskip\noindent
(a)、(b) 由 代数进阶章引理
\ref{more-algebra-lemma-relative-finite-module-over-fibre-product} 与
\ref{more-algebra-lemma-relative-finitely-presented-module-over-fibre-product}
得出。
\end{proof}

\begin{lemma}
\label{more-morphisms-lemma-equivalence-categories-schemes-over-pushout-flat}
在引理 \ref{more-morphisms-lemma-equivalence-categories-schemes-over-pushout} 的情形下，
若对某个三元组 $(V, U', \varphi)$ 有 $V' = G(V, U', \varphi)$，则
\begin{enumerate}
\item $V' \to Y'$ 局部有限型，当且仅当 $V \to Y$ 与 $U' \to X'$
局部有限型；
\item $V' \to Y'$ 平坦，当且仅当 $V \to Y$ 与 $U' \to X'$ 平坦；
\item $V' \to Y'$ 平坦且局部有限表示，当且仅当 $V \to Y$ 与
$U' \to X'$ 平坦且局部有限表示；
\item $V' \to Y'$ 光滑，当且仅当 $V \to Y$ 与 $U' \to X'$ 光滑；
\item $V' \to Y'$ 平展，当且仅当 $V \to Y$ 与 $U' \to X'$ 平展；
\item 视需要在此添加更多性质。
\end{enumerate}
若 $W'$ 在 $Y'$ 上平坦，则伴随映射 $G(F(W')) \to W'$ 是同构。
因此 $F$ 与 $G$ 在以下两个范畴之间定义互为拟逆的函子：在 $Y'$ 上平坦的
概形范畴，以及满足 $V \to Y$ 与 $U' \to X'$ 平坦的三元组
$(V, U', \varphi)$ 的范畴。
\end{lemma}

\begin{proof}
逐个考察仿射片后，本引理的断言等价于 代数进阶章引理
\ref{more-algebra-lemma-properties-algebras-over-fibre-product} 中的相应断言。
\end{proof}







\section{平坦轨迹的开性}
\label{more-morphisms-section-open-flat}

\noindent
这个结果的证明需要一些工作，并且（或许）值得单列一节。结果如下。

\begin{theorem}
\label{more-morphisms-theorem-openness-flatness}
\begin{reference}
\cite[IV Theorem 11.3.1]{EGA}
\end{reference}
设 $S$ 是概形。设 $f : X \to S$ 是局部有限表示的态射。
设 $\mathcal{F}$ 是局部有限表示的拟凝聚 $\mathcal{O}_X$-模。则
$$
U = \{x \in X \mid \mathcal{F}\text{ is flat over }S\text{ 在 }x\}
$$
是 $X$ 的开子集。
\end{theorem}

\begin{proof}
开性可以在 $X$ 上局部检验，故可以假设 $f$ 是仿射概形之间的态射。
此时本定理恰为 代数章定理
\ref{algebra-theorem-openness-flatness}。
\end{proof}

\begin{lemma}
\label{more-morphisms-lemma-flat-locus-base-change}
设 $S$ 是概形。设
$$
\xymatrix{
X' \ar[r]_{g'} \ar[d]_{f'} & X \ar[d]^f \\
S' \ar[r]^g & S
}
$$
是概形的笛卡尔图。设 $\mathcal{F}$ 是拟凝聚 $\mathcal{O}_X$-模。
设 $x' \in X'$ 的像为 $x = g'(x')$ 与 $s' = f'(x')$。
\begin{enumerate}
\item 若 $\mathcal{F}$ 在 $x$ 处于 $S$ 上平坦，则
$(g')^*\mathcal{F}$ 在 $x'$ 处于 $S'$ 上平坦。
\item 若 $g$ 在 $s'$ 处平坦，且 $(g')^*\mathcal{F}$ 在 $x'$ 处于
$S'$ 上平坦，则 $\mathcal{F}$ 在 $x$ 处于 $S$ 上平坦。
\end{enumerate}
特别地，若 $g$ 平坦、$f$ 局部有限表示，且 $\mathcal{F}$ 局部有限表示，
则定理 \ref{more-morphisms-theorem-openness-flatness} 中开子集的形成与基变换交换。
\end{lemma}

\begin{proof}
考虑局部环的交换图
$$
\xymatrix{
\mathcal{O}_{X', x'} & \mathcal{O}_{X, x} \ar[l] \\
\mathcal{O}_{S', s'} \ar[u] & \mathcal{O}_{S, s} \ar[l] \ar[u]
}
$$
注意到 $\mathcal{O}_{X', x'}$ 是
$\mathcal{O}_{X, x} \otimes_{\mathcal{O}_{S, s}} \mathcal{O}_{S', s'}$
的一个局部化，并且 $((g')^*\mathcal{F})_{x'}$ 等于
$\mathcal{F}_x \otimes_{\mathcal{O}_{X, x}} \mathcal{O}_{X', x'}$。
故本引理由 代数章引理
\ref{algebra-lemma-base-change-flat-up-down} 得出。
\end{proof}






\section{按纤维判别平坦性}
\label{more-morphisms-section-criterion-flat-fibres}

\noindent
考虑概形的交换图（左图）
$$
\xymatrix{
X \ar[rr]_f \ar[dr] & & Y \ar[dl] \\
& S
}
\quad
\xymatrix{
X_s \ar[rr]_{f_s} \ar[rd] & & Y_s \ar[dl] \\
& \Spec(\kappa(s))
}
$$
以及拟凝聚 $\mathcal{O}_X$-模 $\mathcal{F}$。
给定落在 $s \in S$ 上方的点 $x \in X$，其像为 $y = f(x)$，
我们考虑如下问题：$\mathcal{F}$ 在 $x$ 处是否于 $Y$ 上平坦？
若 $\mathcal{F}$ 在 $x$ 处于 $S$ 上平坦，则定理表明，这一问题与
$\mathcal{F}$ 在纤维上的限制
$$
\mathcal{F}_s = (X_s \to X)^*\mathcal{F}
$$
在 $x$ 处是否于 $Y_s$ 上平坦密切相关。下面将给出一个“诺特”版本和
一个“有限表示”版本；此前我们还处理过一个“幂零”版本，见引理
\ref{more-morphisms-lemma-flatness-morphism-thickenings}。

\begin{theorem}
\label{more-morphisms-theorem-criterion-flatness-fibre-Noetherian}
设 $S$ 是概形。设 $f : X \to Y$ 是 $S$ 上的概形态射。
设 $\mathcal{F}$ 是拟凝聚 $\mathcal{O}_X$-模。设 $x \in X$，令
$y = f(x)$，并令 $s \in S$ 为 $x$ 在 $S$ 中的像。
假设 $S$、$X$、$Y$ 局部诺特，$\mathcal{F}$ 凝聚，且
$\mathcal{F}_x \not = 0$。则以下条件等价：
\begin{enumerate}
\item $\mathcal{F}$ 在 $x$ 处于 $S$ 上平坦，且
$\mathcal{F}_s$ 在 $x$ 处于 $Y_s$ 上平坦；
\item $Y$ 在 $y$ 处于 $S$ 上平坦，且 $\mathcal{F}$ 在 $x$ 处于
$Y$ 上平坦。
\end{enumerate}
\end{theorem}

\begin{proof}
考虑环映射
$$
\mathcal{O}_{S, s} \longrightarrow
\mathcal{O}_{Y, y} \longrightarrow \mathcal{O}_{X, x}
$$
以及模 $\mathcal{F}_x$。$\mathcal{F}_s$ 在 $x$ 处的茎为模
$\mathcal{F}_x/\mathfrak m_s \mathcal{F}_x$，而 $Y_s$ 在 $y$ 处的局部环为
$\mathcal{O}_{Y, y}/\mathfrak m_s \mathcal{O}_{Y, y}$。
因此蕴含 (1) $\Rightarrow$ (2) 即 代数章引理
\ref{algebra-lemma-criterion-flatness-fibre-Noetherian}。
若 (2) 成立，则第一个环映射忠实平坦，且 $\mathcal{F}_x$ 在
$\mathcal{O}_{Y, y}$ 上平坦，故由 代数章引理
\ref{algebra-lemma-composition-flat}，$\mathcal{F}_x$ 在
$\mathcal{O}_{S, s}$ 上平坦。此外，
$\mathcal{F}_x/\mathfrak m_s \mathcal{F}_x$ 是平坦模 $\mathcal{F}_x$
沿 $\mathcal{O}_{Y, y} \to \mathcal{O}_{Y, y}/\mathfrak m_s \mathcal{O}_{Y, y}$
的基变换，故由 代数章引理
\ref{algebra-lemma-flat-base-change}，它也是平坦的。
\end{proof}

\noindent
下面是非诺特版本。

\begin{theorem}
\label{more-morphisms-theorem-criterion-flatness-fibre}
设 $S$ 是概形。设 $f : X \to Y$ 是 $S$ 上的概形态射。
设 $\mathcal{F}$ 是拟凝聚 $\mathcal{O}_X$-模。假设
\begin{enumerate}
\item $X$ 在 $S$ 上局部有限表示；
\item $\mathcal{F}$ 是有限表示的 $\mathcal{O}_X$-模；
\item $Y$ 在 $S$ 上局部有限型。
\end{enumerate}
设 $x \in X$，令 $y = f(x)$，并令 $s \in S$ 为 $x$ 在 $S$ 中的像。
若 $\mathcal{F}_x \not = 0$，则以下条件等价：
\begin{enumerate}
\item $\mathcal{F}$ 在 $x$ 处于 $S$ 上平坦，且
$\mathcal{F}_s$ 在 $x$ 处于 $Y_s$ 上平坦；
\item $Y$ 在 $y$ 处于 $S$ 上平坦，且 $\mathcal{F}$ 在 $x$ 处于
$Y$ 上平坦。
\end{enumerate}
此外，使 (1) 与 (2) 成立的点 $x$ 所成的集合在
$\text{Supp}(\mathcal{F})$ 中是开的。
\end{theorem}

\begin{proof}
考虑环映射
$$
\mathcal{O}_{S, s} \longrightarrow
\mathcal{O}_{Y, y} \longrightarrow \mathcal{O}_{X, x}
$$
以及模 $\mathcal{F}_x$。$\mathcal{F}_s$ 在 $x$ 处的茎为模
$\mathcal{F}_x/\mathfrak m_s \mathcal{F}_x$，而 $Y_s$ 在 $y$ 处的局部环为
$\mathcal{O}_{Y, y}/\mathfrak m_s \mathcal{O}_{Y, y}$。
因此蕴含 (1) $\Rightarrow$ (2) 即 代数章引理
\ref{algebra-lemma-criterion-flatness-fibre-fp-over-ft}。
若 (2) 成立，则第一个环映射忠实平坦，且 $\mathcal{F}_x$ 在
$\mathcal{O}_{Y, y}$ 上平坦，故由 代数章引理
\ref{algebra-lemma-composition-flat}，$\mathcal{F}_x$ 在
$\mathcal{O}_{S, s}$ 上平坦。此外，
$\mathcal{F}_x/\mathfrak m_s \mathcal{F}_x$ 是平坦模 $\mathcal{F}_x$
沿 $\mathcal{O}_{Y, y} \to \mathcal{O}_{Y, y}/\mathfrak m_s \mathcal{O}_{Y, y}$
的基变换，故由 代数章引理
\ref{algebra-lemma-flat-base-change}，它也是平坦的。

\medskip\noindent
由 态射章引理
\ref{morphisms-lemma-finite-presentation-permanence}，态射 $f$
局部有限表示。考虑集合
\begin{equation}
\label{more-morphisms-equation-open}
U = \{x \in X \mid \mathcal{F} \text{ flat at }x
\text{ over both }Y\text{ 且 }S\}.
\end{equation}
由定理 \ref{more-morphisms-theorem-openness-flatness}，此集合在 $X$ 中是开的。
注意，若 $x \in U$，则 $\mathcal{F}_s$ 是平坦模沿态射
$Y_s \to Y$ 的基变换，因而在 $x$ 处于 $Y_s$ 上平坦；见
态射章引理 \ref{morphisms-lemma-base-change-module-flat}。
故 $U \cap \text{Supp}(\mathcal{F})$ 的每一点都满足条件 (1)。
另一方面，显然若 $x \in \text{Supp}(\mathcal{F})$ 满足 (1) 与 (2)，
则 $x \in U$。因此所求开集为 $U \cap \text{Supp}(\mathcal{F})$。
\end{proof}

\noindent
这些定理常以如下简化形式使用。这里只给出整体陈述——当然也有逐点版本。

\begin{lemma}
\label{more-morphisms-lemma-morphism-between-flat-Noetherian}
设 $S$ 是概形。设 $f : X \to Y$ 是 $S$ 上的概形态射。假设
\begin{enumerate}
\item $S$、$X$、$Y$ 局部诺特；
\item $X$ 在 $S$ 上平坦；
\item 对每个 $s \in S$，态射 $f_s : X_s \to Y_s$ 平坦。
\end{enumerate}
则 $f$ 平坦。若 $f$ 还满射，则 $Y$ 在 $S$ 上平坦。
\end{lemma}

\begin{proof}
这是定理 \ref{more-morphisms-theorem-criterion-flatness-fibre-Noetherian} 的特殊情形。
\end{proof}

\begin{lemma}
\label{more-morphisms-lemma-morphism-between-flat}
设 $S$ 是概形。设 $f : X \to Y$ 是 $S$ 上的概形态射。假设
\begin{enumerate}
\item $X$ 在 $S$ 上局部有限表示；
\item $X$ 在 $S$ 上平坦；
\item 对每个 $s \in S$，态射 $f_s : X_s \to Y_s$ 平坦；
\item $Y$ 在 $S$ 上局部有限型。
\end{enumerate}
则 $f$ 平坦。若 $f$ 还满射，则 $Y$ 在 $S$ 上平坦。
\end{lemma}

\begin{proof}
这是定理 \ref{more-morphisms-theorem-criterion-flatness-fibre} 的特殊情形。
\end{proof}

\begin{lemma}
\label{more-morphisms-lemma-base-change-criterion-flatness-fibre}
设 $S$ 是概形。设 $f : X \to Y$ 是 $S$ 上的概形态射。
设 $\mathcal{F}$ 是拟凝聚 $\mathcal{O}_X$-模。假设
\begin{enumerate}
\item $X$ 在 $S$ 上局部有限表示；
\item $\mathcal{F}$ 是有限表示的 $\mathcal{O}_X$-模；
\item $\mathcal{F}$ 在 $S$ 上平坦；
\item $Y$ 在 $S$ 上局部有限型。
\end{enumerate}
则集合
$$
U = \{x \in X \mid \mathcal{F} \text{ flat at }x \text{ 相对于 }Y\}.
$$
在 $X$ 中是开的，且其形成与任意基变换交换：若 $S' \to S$ 是概形态射，
而 $U'$ 是 $X' = X \times_S S'$ 中使
$\mathcal{F}' = \mathcal{F} \times_S S'$ 在 $Y' = Y \times_S S'$ 上平坦
的点集，则 $U' = U \times_S S'$。
\end{lemma}

\begin{proof}
由 态射章引理
\ref{morphisms-lemma-finite-presentation-permanence}，态射 $f$
局部有限表示。故由定理 \ref{more-morphisms-theorem-openness-flatness}，$U$ 是开的。
由于已假设 $\mathcal{F}$ 在 $S$ 上平坦，定理
\ref{more-morphisms-theorem-criterion-flatness-fibre} 给出
$$
U = \{x \in X \mid \mathcal{F}_s \text{ flat at }x \text{ 相对于 }Y_s\}.
$$
其中 $s$ 始终表示 $x$ 在 $S$ 中的像。（当 $\mathcal{F}_x = 0$ 时，
这一描述也平凡成立。）此外，本引理的假设对态射 $f' : X' \to Y'$
与层 $\mathcal{F}'$ 仍然成立，故 $U'$ 有类似的描述。换言之，只需证明：
给定映到 $s \in S$ 的 $s' \in S'$，集合
$$
\{x' \in X'_{s'} \mid
\mathcal{F}'_{s'} \text{ flat at }x' \text{ 相对于 }Y'_{s'}\}
$$
是 $X_s$ 中相应轨迹的逆像。由引理
\ref{more-morphisms-lemma-flat-locus-base-change}，这确实成立，因为在笛卡尔图
$$
\xymatrix{
X'_{s'} \ar[d] \ar[r] & X_s \ar[d] \\
Y'_{s'} \ar[r] & Y_s
}
$$
中，水平态射是平坦态射
$\Spec(\kappa(s')) \to \Spec(\kappa(s))$ 的基变换，因而平坦。
\end{proof}

\begin{lemma}
\label{more-morphisms-lemma-base-change-flatness-fibres}
设 $S$ 是概形。设 $f : X \to Y$ 是 $S$ 上的概形态射。假设
\begin{enumerate}
\item $X$ 在 $S$ 上局部有限表示；
\item $X$ 在 $S$ 上平坦；
\item $Y$ 在 $S$ 上局部有限型。
\end{enumerate}
则集合
$$
U = \{x \in X \mid X\text{ flat at }x \text{ 相对于 }Y\}.
$$
在 $X$ 中是开的，且其形成与任意基变换交换。
\end{lemma}

\begin{proof}
这是引理 \ref{more-morphisms-lemma-base-change-criterion-flatness-fibre} 的特殊情形。
\end{proof}

\noindent
下一个引理是 代数章引理
\ref{algebra-lemma-free-fibre-flat-free} 的一个变体。注意，假设
$(\mathcal{F}_s)_x$ 是平坦 $\mathcal{O}_{X_s, x}$-模，意味着
$(\mathcal{F}_s)_x$ 是自由 $\mathcal{O}_{X_s, x}$-模。若 $x \in X_s$
是 $X_s$ 某个不可约分支的泛点且 $X_s$ 约化，则这一点总是成立
（此时 $\mathcal{O}_{X_s, x}$ 是域；见 代数章引理
\ref{algebra-lemma-minimal-prime-reduced-ring}）。

\begin{lemma}
\label{more-morphisms-lemma-flat-and-free-at-point-fibre}
设 $f : X \to S$ 是有限表示的概形态射，$\mathcal{F}$ 是有限表示的
$\mathcal{O}_X$-模。设 $x \in X$ 的像为 $s \in S$。若 $\mathcal{F}$
在 $x$ 处于 $S$ 上平坦，且 $(\mathcal{F}_s)_x$ 是平坦
$\mathcal{O}_{X_s, x}$-模，则 $\mathcal{F}$ 在 $x$ 的某个邻域上有限自由。
\end{lemma}

\begin{proof}
若 $\mathcal{F}_x \otimes \kappa(x)$ 为零，则由 Nakayama 引理
（代数章引理 \ref{algebra-lemma-NAK}），$\mathcal{F}_x = 0$，
因而 $\mathcal{F}$ 在 $x$ 的某个邻域上为零（模章引理
\ref{modules-lemma-finite-type-stalk-zero}），此时结论成立。因此可假设
$\mathcal{F}_x \otimes \kappa(x)$ 非零，并将定理
\ref{more-morphisms-theorem-criterion-flatness-fibre} 应用于 $f = \text{id} : X \to X$。
由此得出 $\mathcal{F}_x$ 在 $\mathcal{O}_{X, x}$ 上平坦，故
$\mathcal{F}_x$ 自由；例如见 代数章引理
\ref{algebra-lemma-finite-flat-local}。选择开邻域 $x \in U \subset X$
以及截面 $s_1, \ldots, s_r \in \mathcal{F}(U)$，使它们在
$\mathcal{F}_x$ 中的像构成一组基。缩小 $U$ 后，相应映射
$\psi : \mathcal{O}_U^{\oplus r} \to \mathcal{F}|_U$ 是满射
（模章引理 \ref{modules-lemma-finite-type-stalk-zero}）。于是
$\Ker(\psi)$ 有限型（见 模章引理
\ref{modules-lemma-kernel-surjection-finite-free-onto-finite-presentation}），
且 $\Ker(\psi)_x = 0$。再次缩小 $U$ 后，$\psi$ 即为同构。
\end{proof}

\begin{lemma}
\label{more-morphisms-lemma-finite-free-open}
设 $f : X \to S$ 是局部有限表示的概形态射。设 $\mathcal{F}$ 是
在 $S$ 上平坦的有限表示 $\mathcal{O}_X$-模。则集合
$$
\{x \in X : \mathcal{F}\text{ free in a neighbourhood of }x\}
$$
在 $X$ 中是开的，且其形成与任意基变换 $S' \to S$ 交换。
\end{lemma}

\begin{proof}
开性是平凡的。设 $x \in X$ 映到 $s \in S$。由引理
\ref{more-morphisms-lemma-flat-and-free-at-point-fibre}，$x$ 属于上述集合，当且仅当
$\mathcal{F}|_{X_s}$ 在 $x$ 处于 $X_s$ 上平坦。显然这又等价于
$\mathcal{F}$ 在 $x$ 处于 $X$ 上平坦（因为 $\mathcal{F}_x$ 的自由性
蕴含这一陈述，而这一陈述又蕴含 $\mathcal{F}|_{X_s}$ 在 $x$ 处于
$X_s$ 上平坦）。因此，将引理
\ref{more-morphisms-lemma-base-change-criterion-flatness-fibre} 应用于 $S$ 上的
$\text{id} : X \to X$，即得基变换陈述。
\end{proof}






\section{光滑概形之间的闭浸入}
\label{more-morphisms-section-closed-immersions-smooth}

\noindent
这里收录一些不太适合放在其他地方的结果。

\begin{lemma}
\label{more-morphisms-lemma-etale-local-structure}
设 $S$ 是概形。设 $Y \to X$ 是在 $S$ 上光滑的概形之间的闭浸入。
对每个 $y \in Y$，存在整数 $0 \leq m, n$ 以及交换图
$$
\xymatrix{
Y \ar[d] &
V \ar[l] \ar[d] \ar[r] &
\mathbf{A}^m_S
\ar[d]^{(a_1, \ldots, a_m) \mapsto (a_1, \ldots, a_m, 0 \ldots, 0)} \\
X &
U \ar[l] \ar[r]^-\pi &
\mathbf{A}^{m + n}_S
}
$$
其中 $U \subset X$ 是开集，$V = Y \cap U$，$\pi$ 平展，
$V = \pi^{-1}(\mathbf{A}^m_S)$，且 $y \in V$。
\end{lemma}

\begin{proof}
这个问题在 $X$ 上是局部的，故可以用 $y$ 的一个开邻域替换 $X$。
由 除子章引理
\ref{divisors-lemma-immersion-smooth-into-smooth-regular-immersion}，
$Y \to X$ 是正则浸入，因而可以假设 $X = \Spec(A)$ 仿射，并存在正则序列
$f_1, \ldots, f_n \in A$，使得 $Y = V(f_1, \ldots, f_n)$。
进一步缩小 $X$（从而也缩小 $Y$）后，可以假设存在平展态射
$Y \to \mathbf{A}^m_S$；见 态射章引理
\ref{morphisms-lemma-smooth-etale-over-affine-space}。设
$\mathcal{O}_Y(Y)$ 中的 $\overline{g}_1, \ldots, \overline{g}_m$
是此平展态射的坐标函数。选取这些函数的提升
$g_1, \ldots, g_m \in A$，并考虑态射
$$
(g_1, \ldots, g_m, f_1, \ldots, f_n) :
X
\longrightarrow
\mathbf{A}^{m + n}_S
$$
于 $S$ 上。两边概形均在 $S$ 上局部有限表示，故此态射局部有限表示
（态射章引理
\ref{morphisms-lemma-finite-presentation-permanence}）。按构造，此态射在
$\mathbf{A}^m_S \subset \mathbf{A}^{m + n}_S$ 上的限制平展。因此，为证明
$X \to \mathbf{A}^{m + n}_S$ 在 $y$ 处平展，只需证明
$X \to \mathbf{A}^{m + n}_S$ 在 $y$ 处平坦；见 态射章引理
\ref{morphisms-lemma-etale-at-point}。设 $s \in S$ 为 $y$ 的像。由定理
\ref{more-morphisms-theorem-criterion-flatness-fibre}，只需检验
$X_s \to \mathbf{A}^{m + n}_s$ 在 $y$ 处平坦。设
$z \in \mathbf{A}^{m + n}_s$ 为 $y$ 的像。局部环映射
$$
\mathcal{O}_{\mathbf{A}^{m + n}_s, z}
\longrightarrow
\mathcal{O}_{X_s, y}
$$
由 代数章引理 \ref{algebra-lemma-CM-over-regular-flat} 是平坦的。
事实上，域上的光滑概形是正则的，而正则环是 Cohen--Macaulay 环；见
簇章引理 \ref{varieties-lemma-smooth-regular} 与 代数章引理
\ref{algebra-lemma-regular-ring-CM}。因此源与靶都是正则局部环
（从而都是 CM 环）。源与靶的维数相同：由 代数进阶章引理
\ref{more-algebra-lemma-dimension-etale-extension}，有
$\dim(\mathcal{O}_{Y_s, y}) = \dim(\mathcal{O}_{\mathbf{A}^m_s, z})$
；又有 $\dim(\mathcal{O}_{\mathbf{A}^{m + n}_s, z}) = n +
\dim(\mathcal{O}_{\mathbf{A}^m_s, z})$，并且
$\dim(\mathcal{O}_{X_s, y}) = n + \dim(\mathcal{O}_{Y_s, y})$
，因为 $\mathcal{O}_{Y_s, y}$ 是 $\mathcal{O}_{X_s, y}$ 对长度为 $n$
的正则序列 $f_1, \ldots, f_n$ 的商（见 除子章说明
\ref{divisors-remark-relative-regular-immersion-elements}）。最后，由于
$Y_s \to \mathbf{A}^m_s$ 在 $y$ 处平展，上述箭头的纤维环在
$\kappa(z)$ 上有限。证明完毕。
\end{proof}

\begin{remark}
\label{more-morphisms-remark-relative-blowup}
固定一个环 $R$，并令 $S = \Spec(R)$。固定整数 $0 \leq m$ 与
$1 \leq n$。考虑闭浸入
$$
Z = \mathbf{A}^m_S \longrightarrow \mathbf{A}^{m + n}_S = X,\quad
(a_1, \ldots, a_m) \mapsto (a_1, \ldots, a_m, 0, \ldots 0).
$$
我们考察 $X$ 沿闭子概形 $Z$ 的吹起 $X'$。写成
$$
X = \Spec(A)
\quad\text{其中}\quad
A = R[x_1, \ldots, x_m, y_1, \ldots, y_n]
$$
则 $X'$ 是 $A$ 关于理想 $(y_1, \ldots, y_n)$ 的 Rees 代数的 Proj。
此 Rees 代数等于 $B = A[T_1, \ldots, T_n]/(y_iT_j - y_jT_i)$；略去细节。
因此 $X' = \text{Proj}(B)$ 在 $S$ 上光滑，因为它由仿射开集
\begin{align*}
D_+(T_i)
& =
\Spec(B_{(T_i)}) \\
& =
\Spec(A[t_1, \ldots, \hat t_i, \ldots t_n]/(y_j - y_i t_j)) \\
& =
\Spec(R[x_1, \ldots, x_m, y_i, t_1, \ldots, \hat t_i, \ldots, t_n])
\end{align*}
覆盖，而这些开集同构于 $\mathbf{A}^{n + m}_S$。在此坐标图中，
例外除子由令 $y_i = 0$ 截出，故例外除子也在 $S$ 上光滑。
\end{remark}

\begin{lemma}
\label{more-morphisms-lemma-blowup}
设 $S$ 是概形。设 $Z \to X$ 是在 $S$ 上光滑的概形之间的闭浸入。
设 $b : X' \to X$ 是沿 $Z$ 的吹起，例外除子为 $E \subset X'$。
则 $X'$ 与 $E$ 在 $S$ 上光滑。态射 $p : E \to Z$ 典范同构于射影空间丛
$$
\mathbf{P}(\mathcal{I}/\mathcal{I}^2) \longrightarrow Z
$$
，其中 $\mathcal{I} \subset \mathcal{O}_X$ 是 $Z$ 的理想层。
来自射影空间丛结构的相对 $\mathcal{O}_E(1)$ 同构于
$\mathcal{O}_{X'}(-E)$ 在 $E$ 上的限制。
\end{lemma}

\begin{proof}
% P05-ERR-0025: frozen source has typographical "immmersion".
由 除子章引理
\ref{divisors-lemma-immersion-smooth-into-smooth-regular-immersion}，
浸入 $Z \to X$ 是正则浸入，故理想层 $\mathcal{I}$ 有限型，从而 $b$
是射影态射，并带有相对丰沛可逆层
$\mathcal{O}_{X'}(1) = \mathcal{O}_{X'}(-E)$；见 除子章引理
\ref{divisors-lemma-blowing-up-gives-effective-Cartier-divisor} 及
\ref{divisors-lemma-blowing-up-projective}.
典范映射 $\mathcal{I} \to b_*\mathcal{O}_{X'}(1)$ 按照吹起的构造给出闭浸入
$$
X' \longrightarrow
\mathbf{P}\left(\bigoplus\nolimits_{n \geq 0}
\text{Sym}^n_{\mathcal{O}_X}(\mathcal{I})\right)
$$
。此态射在 $E$ 上的限制给出 $Z$ 上的典范映射
$$
E \longrightarrow
\mathbf{P}\left(\bigoplus\nolimits_{n \geq 0}
\text{Sym}^n_{\mathcal{O}_Z}(\mathcal{I}/\mathcal{I}^2)\right)
$$
。由于 $\mathcal{I}/\mathcal{I}^2$ 有限局部自由，若此典范映射是同构，
则引理的最后一部分成立。说完这些之后，问题现在在 $X$ 上是平展局部的。
事实上，由 除子章引理
\ref{divisors-lemma-flat-base-change-blowing-up}，吹起与平坦基变换交换；
而且在预合成一个满平展态射之后可以检验光滑性。因此，利用引理
\ref{more-morphisms-lemma-etale-local-structure} 所给出的 $S$ 上概形闭浸入的平展局部结构，
问题归结为评注 \ref{more-morphisms-remark-relative-blowup} 中讨论的情形。
\end{proof}








\section{平坦模与相对伴随点}
\label{more-morphisms-section-flat-relative-assassin}

\noindent
本节将证明：在几何情形下，平坦模的支集在基底上（在某种意义下）
是等维的。诺特情形可参见 \cite[IV Proposition 12.1.1.5]{EGA}。
首先证明两个辅助引理。

\begin{lemma}
\label{more-morphisms-lemma-mod-injective-valuation-ring}
% P05-ERR-0026: frozen source has ungrammatical "Let A -> B is".
设 $A$ 是赋值环。设 $A \to B$ 是局部环之间的局部同态，且本质有限型。
设 $u : N \to M$ 是有限 $B$-模之间的映射。假设 $M$ 在 $A$ 上平坦，且
$\overline{u} : N/\mathfrak m_A N \to M/\mathfrak m_A M$ 是单射。
则 $u$ 是单射，且 $M/u(N)$ 在 $A$ 上平坦。
\end{lemma}

\begin{proof}
我们将从 代数章引理 \ref{algebra-lemma-mod-injective-general}
推出本引理（请注意，我们交换了 $M$ 与 $N$ 的角色）。为完成化约，
将使用 代数进阶章引理
\ref{more-algebra-lemma-valuation-ring-flat-essentially-finite-type}
把问题化到有限表示情形。

\medskip\noindent
由假设，可以把 $B$ 写成多项式代数 $P = A[x_1, \ldots, x_n]$
在素理想 $\mathfrak q$ 处的局部化的一个商。于是可以把
$u : N \to M$ 看成有限 $P_\mathfrak q$-模之间的映射。因此不妨并且确实
假设 $B$ 在 $A$ 上本质有限表示。

\medskip\noindent
其次，$B$-模 $N$ 有限，但未必有限表示。将 $N = \colim N_\lambda$
写成有限表示 $B$-模的滤过余极限，其转移映射均为满射。例如，选择表示
$0 \to K \to B^{\oplus r} \to N \to 0$，把 $K$ 写成其有限子模
$K_\lambda$ 的并，并令 $N_\lambda = \Coker(K_\lambda \to B^{\oplus r})$。
模 $N/\mathfrak m_A N$ 在诺特环 $B/\mathfrak m_A B$ 上有限表示。
故对充分大的 $\lambda$，有
$N_\lambda/\mathfrak m_A N_\lambda = N/\mathfrak m_A N$。
现在，若能对复合映射 $u_\lambda : N_\lambda \to M$ 证明本引理，
则可推出 $N_\lambda = N$，并且结论对 $u$ 成立。因此不妨并且确实假设
$N$ 在 $B$ 上有限表示。

\medskip\noindent
由 代数进阶章引理
\ref{more-algebra-lemma-valuation-ring-flat-essentially-finite-type}，模 $M$
在 $B$ 上有限表示。因此 代数章引理
\ref{algebra-lemma-mod-injective-general} 的所有假设都成立，结论随之得出。
\end{proof}

\begin{lemma}
\label{more-morphisms-lemma-make-base-change}
\begin{reference}
这可见于 \cite[IV Proposition 12.1.1.5]{EGA} 的证明
\end{reference}
设 $f : X \to S$ 是概形态射。设点 $y \in X$ 的像为 $t \in S$。
以 $Y \subset X$ 表示 $\{y\}$ 的闭包，并赋予它 $X$ 的整闭子概形结构。
设 $s \in S$，并设 $x \in Y_s$ 是 $Y_s$ 某个不可约分支的泛点。
存在笛卡尔图
$$
\xymatrix{
X' \ar[r]_{g'} \ar[d]_{f'} & X \ar[d]^f \\
S' \ar[r]^g & S
}
$$
满足以下性质：
\begin{enumerate}
\item $S'$ 是某个赋值环的谱，其泛点为 $t'$、闭点为 $s'$；
\item $g(t') = t$ 且 $g(s') = s$；
\item 存在点 $y' \in X'_{t'}$，它是
$(S' \times_S Y)_{t'} = Y_t \times_t t'$ 某个不可约分支的泛点，
并满足 $g'(y') = y$；
\item 以 $Y' \subset X'$ 表示 $\{y'\}$ 的闭包，并赋予它 $X'$
的整闭子概形结构；则存在点 $x' \in Y'_{s'}$，它是
$Y'_{s'}$ 某个不可约分支的泛点，并满足 $g'(x') = x$。
\end{enumerate}
\end{lemma}

\begin{proof}
选择赋值环 $R$，令 $S' = \Spec(R)$，其泛点为 $t'$、闭点为 $s'$；
再选择满足 $h(t') = y$ 与 $h(s') = x$ 的态射 $h : S' \to X$。
见 概形章引理 \ref{schemes-lemma-points-specialize}。令
$g = f \circ h$，于是 $g(t') = t$ 且 $g(s') = s$。考虑基变换
$$
\xymatrix{
X' \ar[r]_{g'} \ar[d] & X \ar[d] \\
S' \ar@/^1em/[u]^\sigma \ar[r]^-g & S
}
$$
得到此基变换的一个截面 $\sigma$，满足 $h = g' \circ \sigma$。

\medskip\noindent
由于 $h$ 经由 $Y$ 分解，$\sigma$ 当然经由 $Y$ 的基变换
$S' \times_S Y$ 分解。设 $y' \in X'_{t'} \subset X'$ 是纤维
$$
(S' \times_S Y)_{t'} = Y_t \times_t t'
$$
的某个不可约分支的泛点，该分支包含点 $\sigma(t')$；也就是说，
$y' \leadsto \sigma(t')$。
% P05-ERR-0027: frozen source applies g:S'->S to points of X'; likely g'.
由于 $g'(y') \in Y_t$ 且 $g(y') \leadsto g(\sigma(t')) = y$，而 $y$
是纤维 $Y_t$ 的泛点，故 $g'(y') = y$。以 $Y' \subset X'$ 表示
$\{y'\}$ 在 $X'$ 中的闭包，并赋予它整闭子概形结构。由于
$\sigma(t') \in Y'$，$\sigma$ 经由 $Y'$ 分解。选择 $Y'_{s'}$
某个包含 $\sigma(s')$ 的不可约分支的泛点 $x' \in Y'_{s'}$；也就是说，
$x' \leadsto \sigma(s')$，从而 $g'(x') \leadsto g'(\sigma(s')) = x$。
再次利用假设中 $x$ 是 $Y_s$ 某个不可约分支的泛点，以及
$g'(Y') \subset Y$，得出 $g'(x') = x$。
\end{proof}

\begin{lemma}
\label{more-morphisms-lemma-associated-point-specializes}
设 $f : X \to S$ 是局部有限型的概形态射，$\mathcal{F}$ 是有限型拟凝聚
$\mathcal{O}_X$-模。设 $y \in \text{Ass}_{X/S}(\mathcal{F})$ 的像为
$t \in S$。以 $Y \subset X$ 表示 $\{y\}$ 在 $X$ 中的闭包，并赋予它
整闭子概形结构。设 $s \in S$，并设 $x \in Y_s$ 是 $Y_s$
某个不可约分支的泛点。若 $\mathcal{F}$ 在 $x$ 处于 $S$ 上平坦，则
$x \in \text{Ass}_{X/S}(\mathcal{F})$，且
$\dim_x(Y_s) = \dim(Y_t)$。
\end{lemma}

\begin{proof}
选择引理 \ref{more-morphisms-lemma-make-base-change} 中的图，并令
$\mathcal{F}' = (g')^*\mathcal{F}$。除子章引理
\ref{divisors-lemma-base-change-relative-assassin} 推出
$y' \in \text{Ass}_{X'/S'}(\mathcal{F}')$。由 $y'$ 的选取还可得
$\dim(Y'_{t'}) = \dim(Y_t)$；例如见 代数章引理
\ref{algebra-lemma-inequalities-under-field-extension}。由 代数章引理
\ref{algebra-lemma-finite-type-domain-over-valuation-ring-dim-fibres}，
$Y'_{s'}$ 等维，且其维数等于 $\dim(Y_t)$。由于 $\mathcal{F}$ 在 $x$
处于 $S$ 上平坦，$\mathcal{F}'$ 在 $x'$ 处于 $S'$ 上平坦；见
态射章引理 \ref{morphisms-lemma-base-change-module-flat}。

\medskip\noindent
% P05-ERR-0028: frozen source uses X'/S here, while the surrounding base is S'.
假设能够证明 $x' \in \text{Ass}_{X'/S}(\mathcal{F}')$。则 除子章引理
\ref{divisors-lemma-base-change-relative-assassin} 推出
$x \in \text{Ass}_{X/S}(\mathcal{F})$，并且 $Y'_{s'}$ 中包含 $x'$
的不可约分支 $C'$ 是 $C \times_s s'$ 的一个不可约分支，其中
$C \subset Y_s$ 是包含 $x$ 的不可约分支。因此
$\dim(C) = \dim(C') = \dim(Y_t)$（见上文），证明完成。
这把问题化到下一段讨论的情形。

\medskip\noindent
假设 $S = \Spec(A)$，其中 $A$ 是赋值环，而 $t$ 与 $s$ 分别是 $S$
的泛点与闭点。为导出矛盾，假设
$x \not \in \text{Ass}_{X/S}(\mathcal{F})$。换言之，假设
$x \not \in \text{Ass}_{X_s}(\mathcal{F}_s)$，其中 $\mathcal{F}_s$
是 $\mathcal{F}$ 在 $X_s$ 上的拉回。考虑环映射
$$
A \longrightarrow \mathcal{O}_{X, x} = B
$$
以及 $B = \mathcal{O}_{X, x}$ 上的模 $N = \mathcal{F}_x$。于是
$B/\mathfrak m_A B = \mathcal{O}_{X_s, x}$，而 $N/\mathfrak m_A N$
是 $\mathcal{F}_s$ 在点 $x$ 处的茎。以 $\mathfrak q \subset B$
表示点 $y$ 对应的素理想；见 概形章引理
\ref{schemes-lemma-specialize-points}。由于 $x$ 是 $Y_s$ 的泛点，
$\mathfrak q + \mathfrak m_A B$ 的根为 $\mathfrak m_B$。集合
$\text{Ass}_{B/\mathfrak m_A B}(N/\mathfrak m_A N)$ 是有限素理想集
（代数章引理 \ref{algebra-lemma-finite-ass}），并且由于
$x \not \in \text{Ass}_{X/S}(\mathcal{F})$，它不含
$B/\mathfrak m_A B$ 的极大理想。
% P05-ERR-0029: frozen source has duplicated word "of of".
因此 $\mathfrak q$ 在 $B/\mathfrak m_A B$ 中的像不包含于这些素理想中的
任何一个。由素理想回避（代数章引理 \ref{algebra-lemma-silly}），
可以找到元素 $g \in \mathfrak q$，使其在 $B/\mathfrak m_A B$ 中的像是
$N/\mathfrak m_A N$ 上的非零因子（这里使用 代数章引理
\ref{algebra-lemma-ass-zero-divisors} 对零因子的描述）。由于
$N = \mathcal{F}_x$ 在 $A$ 上平坦，引理
\ref{more-morphisms-lemma-mod-injective-valuation-ring} 表明
$$
g : N \longrightarrow N
$$
是单射。特别地，若 $K = \text{Frac}(A)$ 是 $A$ 的分式域，则
$$
g : N \otimes_A K \longrightarrow N \otimes_A K
$$
是单射。注意，$\mathfrak q$ 对应于 $B \otimes_A K$ 的一个素理想。
以 $\mathcal{F}_t$ 表示 $\mathcal{F}$ 在泛纤维 $X_t$ 上的限制。我们有
$(B \otimes_A K)_{\mathfrak q} = \mathcal{O}_{X_t, y}$，且
$(N \otimes_A K)_\mathfrak q$ 是 $\mathcal{F}_t$ 在 $y$ 处的茎。
因此 $g \in \mathfrak m_y \subset \mathcal{O}_{X_t, y}$ 是茎
$(\mathcal{F}_t)_y$ 上的非零因子，这与假设
$y \in \text{Ass}_{X/S}(\mathcal{F})$ 矛盾。
\end{proof}

\begin{lemma}
\label{more-morphisms-lemma-relative-dimension-support-flat}
设 $f : X \to S$ 是局部有限型的概形态射。设 $\mathcal{F}$ 是在 $S$
上平坦的有限型拟凝聚 $\mathcal{O}_X$-模。假设 $S$ 不可约，泛点为
$\eta$。若 $\dim(\text{Supp}(\mathcal{F}_\eta)) \leq r$，则对所有
$s \in S$ 都有 $\dim(\text{Supp}(\mathcal{F}_s)) \leq r$。
\end{lemma}

\begin{proof}
设 $x \in \text{Supp}(\mathcal{F}_s)$ 是
$\text{Supp}(\mathcal{F}_s)$ 某个不可约分支的泛点。由 代数章引理
\ref{algebra-lemma-going-down-flat-module}，可以在
$\text{Supp}(\mathcal{F})$ 中找到特化 $y \leadsto x$，满足
$f(y) = \eta$。当然可以假设 $y$ 是
$\text{Supp}(\mathcal{F}_\eta)$ 某个不可约分支的泛点。由引理
\ref{more-morphisms-lemma-associated-point-specializes}，$\overline{\{x\}}$ 的维数至多为
$r$。
\end{proof}

\begin{lemma}
\label{more-morphisms-lemma-flat-associated-equidimensional}
设 $f : X \to S$ 是局部有限型的概形态射，$\mathcal{F}$ 是有限型拟凝聚
$\mathcal{O}_X$-模。设 $y \in \text{Ass}_{X/S}(\mathcal{F})$。
以 $Y \subset X$ 表示 $\{y\}$ 在 $X$ 中的闭包，并赋予它整闭子概形结构；
以 $T \subset S$ 表示 $\{f(y)\}$ 的闭包，并赋予它整闭子概形结构。
得到交换图
$$
\xymatrix{
Y \ar[r] \ar[d] & X \ar[d] \\
T \ar[r] & S
}
$$
，其中 $Y \to T$ 是支配态射。假设在 $Y \to T$ 的各纤维的所有不可约
分支的泛点处，$\mathcal{F}$ 都在 $S$ 上平坦（例如，若
$\mathcal{F}$ 在 $S$ 上平坦）。则
\begin{enumerate}
\item 若 $s \in S$，且 $x \in Y_s$ 是 $Y_s$ 某个不可约分支的泛点，
则 $x \in \text{Ass}_{X/S}(\mathcal{F})$；
\item 存在整数 $d \geq 0$，使得 $Y \to T$ 的相对维数为 $d$；见
态射章定义 \ref{morphisms-definition-relative-dimension-d}。
\end{enumerate}
\end{lemma}

\begin{proof}
这立即由逐点版本引理 \ref{more-morphisms-lemma-associated-point-specializes} 得出。
注意，为计算局部代数概形 $Y_s$ 的维数，只需考察泛点附近；见
簇章第 \ref{varieties-section-algebraic-schemes} 节。
\end{proof}

\begin{remark}
\label{more-morphisms-remark-flat-equidimensional}
下面列出一些情形，在这些情形中，上述材料，尤其是引理
\ref{more-morphisms-lemma-flat-associated-equidimensional}，可以用来断定概形态射
$f : X \to S$ 具有 态射章定义
\ref{morphisms-definition-relative-dimension-d} 所定义的相对维数 $d$。
例如，以下条件足够：
\begin{enumerate}
\item $X$ 是整概形，泛点为 $\xi$；
\item $\kappa(\xi)$ 在 $\kappa(f(\xi))$ 上的超越次数为 $d$；
\item $f$ 局部有限型；
\item 存在有限型拟凝聚 $\mathcal{O}_X$-模 $\mathcal{F}$，它在 $S$
上平坦，且 $\text{Supp}(\mathcal{F}) = X$。
\end{enumerate}
另一组足够的假设如下：
\begin{enumerate}
\item $S$ 不可约，泛点为 $\eta$；
\item $X_\eta$ 在 $X$ 中稠密；
\item $X_\eta$ 的每个不可约分支都具有维数 $d$；
\item $f$ 局部有限型；
\item 存在有限型拟凝聚 $\mathcal{O}_X$-模 $\mathcal{F}$，它在 $S$
上平坦，且 $\text{Supp}(\mathcal{F}) = X$。
\end{enumerate}
当然，可以放宽 $\mathcal{F}$ 的平坦性条件，只要求 $\mathcal{F}$
在 $S$ 上余维 $0$ 处平坦；也就是说，$\mathcal{F}$ 在每个纤维的每个泛点处都于
$S$ 上平坦。若以后需要这些结果，我们会在此仔细陈述并证明它们。
\end{remark}











\section{再论正规化}
\label{more-morphisms-section-normalization}

\noindent
正规化与光滑基变换交换。

\begin{lemma}
\label{more-morphisms-lemma-integral-closure-smooth-pullback}
设 $f : Y \to X$ 是概形的光滑态射，$\mathcal{A}$ 是拟凝聚
$\mathcal{O}_X$-代数层。$\mathcal{O}_Y$ 在 $f^*\mathcal{A}$ 中的整闭包
等于 $f^*\mathcal{A}'$，其中 $\mathcal{A}' \subset \mathcal{A}$ 是
$\mathcal{O}_X$ 在 $\mathcal{A}$ 中的整闭包。
\end{lemma}

\begin{proof}
这是用概形语言对 代数章引理
\ref{algebra-lemma-integral-closure-commutes-smooth} 的重述。略去细节。
\end{proof}

\begin{lemma}[正规化与光滑基变换交换]
\label{more-morphisms-lemma-normalization-smooth-localization}
设
$$
\xymatrix{
Y_2 \ar[r] \ar[d]_{f_2} & Y_1 \ar[d]^{f_1} \\
X_2 \ar[r]^\varphi & X_1
}
$$
是概形范畴中的纤维方块。假设 $f_1$ 拟紧且拟分离，并且 $\varphi$ 光滑。
设 $Y_i \to X_i' \to X_i$ 是 $X_i$ 在 $Y_i$ 中的正规化。则
$X_2' \cong X_2 \times_{X_1} X_1'$。
\end{lemma}

\begin{proof}
将分解 $Y_1 \to X_1' \to X_1$ 基变换到 $X_2$，得到分解
$Y_2 \to X_2 \times_{X_1} X_1' \to X_2$，并且
$X_2 \times_{X_1} X_1' \to X_2$ 是整态射（态射章引理
\ref{morphisms-lemma-base-change-finite}）。因此，由 态射章引理
\ref{morphisms-lemma-characterize-normalization} 的泛性质，得到态射
$h : X_2' \to X_2 \times_{X_1} X_1'$。注意，$X_2'$ 是
$\mathcal{O}_{X_2}$ 在 $f_{2, *}\mathcal{O}_{Y_2}$ 中的整闭包的相对谱。
若 $\mathcal{A}' \subset f_{1, *}\mathcal{O}_{Y_1}$ 表示
$\mathcal{O}_{X_1}$ 的整闭包，则 $X_2 \times_{X_1} X_1'$ 是
$\varphi^*\mathcal{A}'$ 的相对谱；见 构造章引理
\ref{constructions-lemma-spec-properties}。由 概形上同调章引理
\ref{coherent-lemma-flat-base-change-cohomology}，有
$f_{2, *}\mathcal{O}_{Y_2} = \varphi^*f_{1, *}\mathcal{O}_{Y_1}$。
故结论由引理 \ref{more-morphisms-lemma-integral-closure-smooth-pullback} 得出。
\end{proof}

\begin{lemma}[正规化与光滑态射]
\label{more-morphisms-lemma-normalization-and-smooth}
设 $X \to Y$ 是概形的光滑态射。假设 $Y$ 的每个拟紧开集都只有有限多个
不可约分支。则 $X$ 也具有此性质，并且在 $X$ 上存在唯一同构
$X^\nu = X \times_Y Y^\nu$，其中 $X^\nu$、$Y^\nu$ 分别是 $X$、$Y$
的正规化。
\end{lemma}

\begin{proof}
由 下降章引理
\ref{descent-lemma-locally-finite-nr-irred-local-fppf}，$X$ 的每个拟紧开集
都只有有限多个不可约分支。注意 $X_{red} = X \times_Y Y_{red}$，因为
约化概形上的光滑概形是约化的；见 下降章引理
\ref{descent-lemma-reduced-local-smooth}。故可假设 $X$ 与 $Y$ 都约化
（按定义，一个概形的正规化等于其约化的正规化）。其次，由 Descent,
引理 \ref{descent-lemma-normal-local-smooth}，
$X' = X \times_Y Y^\nu$ 是正规概形。态射 $X' \to Y^\nu$ 光滑
（从而平坦），因此 $X'$ 各不可约分支的泛点位于 $Y^\nu$ 各不可约
分支的泛点上方。由于 $Y^\nu \to Y$ 是双有理态射，可知
$X' \to X$ 也是双有理态射（因为 $X' \to Y^\nu$ 在 $Y$ 的泛点上方的
纤维上诱导同构）。于是存在分解 $X^\nu \to X' \to X$；见 Morphisms,
引理 \ref{morphisms-lemma-normalization-normal}。由于 $X'$ 正规且在
$X$ 上整，此分解是同构。
\end{proof}

\begin{lemma}[正规化与 Hensel 化]
\label{more-morphisms-lemma-normalization-and-henselization}
设 $X$ 是局部诺特概形，$\nu : X^\nu \to X$ 是正规化态射。则对任意点
$x \in X$，基变换
$$
X^\nu \times_X \Spec(\mathcal{O}_{X, x}^h) \to \Spec(\mathcal{O}_{X, x}^h),
\quad\text{分别}\quad
X^\nu \times_X \Spec(\mathcal{O}_{X, x}^{sh}) \to \Spec(\mathcal{O}_{X, x}^{sh})
$$
分别是 $\Spec(\mathcal{O}_{X, x}^h)$ 与
$\Spec(\mathcal{O}_{X, x}^{sh})$ 的正规化。
\end{lemma}

\begin{proof}
设 $\eta_1, \ldots, \eta_r$ 是 $X$ 中经过 $x$ 的不可约分支的泛点。
将正规化基变换到 $\Spec(\mathcal{O}_{X, x})$，得到
$\mathcal{O}_{X, x}$ 在 $\prod \kappa(\eta_i)$ 中的整闭包的谱。
这由 态射章定义 \ref{morphisms-definition-normalization}
中 $X$ 的正规化构造以及 态射章引理
\ref{morphisms-lemma-integral-closure} 得出；也可使用 态射章引理
\ref{morphisms-lemma-description-normalization} 对正规化的描述。
因此问题化为如下代数问题。设 $A$ 是诺特局部环；回忆这蕴含其 Hensel 化
$A^h$ 与严格 Hensel 化 $A^{sh}$ 也都是诺特的（代数进阶章引理
\ref{more-algebra-lemma-henselization-noetherian}）。设
$\mathfrak p_1, \ldots, \mathfrak p_r$ 是其极小素理想，并设 $A'$ 是
$A$ 在 $\prod \kappa(\mathfrak p_i)$ 中的整闭包。问题是证明：
$A' \otimes_A A^h$ 以及相应的 $A' \otimes_A A^{sh}$，分别可由诺特局部环
$A^h$ 与 $A^{sh}$ 以同样方式构造。

\medskip\noindent
% P05-ERR-0030: frozen source says minimal primes of A (likely A^h) and A^sh.
由于 $A^h$ 以及相应的 $A^{sh}$ 都是平展 $A$-代数的余极限，
$A$ 与 $A^{sh}$ 的极小素理想恰好分别是 $A^h$ 与 $A^{sh}$ 中位于
$A$ 的极小素理想上方的素理想（由下降性质；见 代数章引理
\ref{algebra-lemma-flat-going-down} 与
\ref{algebra-lemma-minimal-prime-image-minimal-prime}）。因此 More on Algebra,
引理 \ref{more-algebra-lemma-fibres-henselization} 告诉我们，
$A^h \otimes_A \prod \kappa(\mathfrak p_i)$,
分别
$A^{sh} \otimes_A \prod \kappa(\mathfrak p_i)$
分别是 $A^h$ 与 $A^{sh}$ 的极小素理想处的剩余域之积。已知在一个上环中
取整闭包与平展基变换交换；见 代数章引理
\ref{algebra-lemma-integral-closure-commutes-etale}。把 $A^h$ 与 $A^{sh}$
写成平展 $A$-代数的极限，可知对到 $A^h$ 与 $A^{sh}$ 的基变换也有同样
结论（也可使用更一般的 代数章引理
\ref{algebra-lemma-integral-closure-commutes-colim-smooth}）。
\end{proof}








\section{正规态射}
\label{more-morphisms-section-normal}

\noindent
Deligne 与 Mumford 的文章 \cite{DM} 提到了正规态射的概念。这只是可以用
类似方式定义的一系列态射类型之一\footnote{其他类型包括 coprof
$\leq k$、Cohen--Macaulay、$(S_k)$、正则、$(R_k)$ 与约化。见
\cite[IV Definition 6.8.1.]{EGA}。Gorenstein 态射将在 Duality for Schemes,
第 \ref{duality-section-gorenstein} 节 中定义。}随着需要，我们会逐步用
各自的章节加入这些类型。

\begin{definition}
\label{more-morphisms-definition-normal}
设 $f : X \to Y$ 是概形态射。假设所有纤维 $X_y$ 都是局部诺特概形。
\begin{enumerate}
\item 设 $x \in X$ 且 $y = f(x)$。若 $f$ 在 $x$ 处平坦，并且概形
$X_y$ 在 $x$ 处于 $\kappa(y)$ 上几何正规（见 簇章定义
\ref{varieties-definition-geometrically-normal}），则称 $f$ 在
$x$ 处{\it 正规}。
\item 若 $f$ 在 $X$ 的每一点处都正规，则称 $f$ 是{\it 正规态射}。
\end{enumerate}
\end{definition}

\noindent
因此，态射 $X \to Y$ 正规这一条件，比仅仅要求所有纤维都是正规局部诺特
概形更强。

\begin{lemma}
\label{more-morphisms-lemma-normal}
设 $f : X \to Y$ 是概形态射。假设 $f$ 的所有纤维都局部诺特。
以下条件等价：
\begin{enumerate}
\item $f$ 正规；
\item $f$ 平坦，且其纤维都是几何正规概形。
\end{enumerate}
\end{lemma}

\begin{proof}
这直接由定义得出。
\end{proof}

\begin{lemma}
\label{more-morphisms-lemma-smooth-normal}
光滑态射是正规的。
\end{lemma}

\begin{proof}
设 $f : X \to Y$ 是光滑态射。由 态射章引理
\ref{morphisms-lemma-smooth-locally-finite-presentation}，$f$ 局部有限表示，
故纤维 $X_y$ 在某个域上局部有限型，从而局部诺特。此外，$f$ 平坦；见
态射章引理 \ref{morphisms-lemma-smooth-flat}。最后，纤维 $X_y$
在域上光滑（由 态射章引理
\ref{morphisms-lemma-base-change-smooth}），故由 簇章引理
\ref{varieties-lemma-smooth-geometrically-normal} 几何正规。因此由引理
\ref{more-morphisms-lemma-normal}，$f$ 正规。
\end{proof}

\noindent
我们要证明，对光滑拓扑而言，这一概念在源与靶上都是局部的。先处理纤维
局部诺特这一性质。

\begin{lemma}
\label{more-morphisms-lemma-locally-Noetherian-fibres-fppf-local-source-and-target}
性质 $\mathcal{P}(f)=$“$f$ 的纤维局部诺特”在源与靶的 fppf 拓扑下都是
局部的。
\end{lemma}

\begin{proof}
设 $f : X \to Y$ 是概形态射。设
$\{\varphi_i : Y_i \to Y\}_{i \in I}$ 是 $Y$ 的 fppf 覆盖，并以
$f_i : X_i \to Y_i$ 表示 $f$ 沿 $\varphi_i$ 的基变换。设 $i \in I$，
且 $y_i \in Y_i$ 是一点。令 $y = \varphi_i(y_i)$。注意
$$
X_{i, y_i} = \Spec(\kappa(y_i)) \times_{\Spec(\kappa(y))} X_y.
$$
此外，由于 $\varphi_i$ 有限表示，域扩张 $\kappa(y_i)/\kappa(y)$ 有限生成。
因此在此情形中，$X_y$ 局部诺特，当且仅当 $X_{i, y_i}$ 局部诺特；见
簇章引理 \ref{varieties-lemma-locally-Noetherian-base-change}。
这就推出在靶上的局部性。

\medskip\noindent
设 $\{X_i \to X\}$ 是 $X$ 的 fppf 覆盖，并设 $y \in Y$。此时
$\{X_{i, y} \to X_y\}$ 是该纤维的 fppf 覆盖。因此源上的局部性由
下降章引理 \ref{descent-lemma-Noetherian-local-fppf} 得出。
\end{proof}

\begin{lemma}
\label{more-morphisms-lemma-normal-fppf-local-source-and-target}
性质 $\mathcal{P}(f)=$“$f$ 的纤维局部诺特且 $f$ 正规”在靶的 fppf
拓扑下是局部的，并在源的光滑拓扑下是局部的。
\end{lemma}

\begin{proof}
我们有
$\mathcal{P}(f) =
\mathcal{P}_1(f) \wedge \mathcal{P}_2(f) \wedge \mathcal{P}_3(f)$
，其中 $\mathcal{P}_1(f)=$“$f$ 的纤维局部诺特”，
$\mathcal{P}_2(f)=$“$f$ 平坦”，而
$\mathcal{P}_3(f)=$“$f$ 的纤维几何正规”。我们已经知道
$\mathcal{P}_1$ 与 $\mathcal{P}_2$ 在源与靶的 fppf 拓扑下都是局部的；见
引理 \ref{more-morphisms-lemma-locally-Noetherian-fibres-fppf-local-source-and-target}
以及 下降章引理 \ref{descent-lemma-descending-property-flat} 与
\ref{descent-lemma-flat-fpqc-local-source}。故只需处理 $\mathcal{P}_3$。

\medskip\noindent
设 $f : X \to Y$ 是概形态射。设
$\{\varphi_i : Y_i \to Y\}_{i \in I}$ 是 $Y$ 的 fpqc 覆盖，并以
$f_i : X_i \to Y_i$ 表示 $f$ 沿 $\varphi_i$ 的基变换。设 $i \in I$，
且 $y_i \in Y_i$ 是一点。令 $y = \varphi_i(y_i)$。注意
$$
X_{i, y_i} = \Spec(\kappa(y_i)) \times_{\Spec(\kappa(y))} X_y.
$$
因此在此情形中，$X_y$ 几何正规，当且仅当 $X_{i, y_i}$ 几何正规；见
簇章引理 \ref{varieties-lemma-geometrically-normal-upstairs}。
这表明 $\mathcal{P}_3$ 在靶的 fpqc 拓扑下是局部的。

\medskip\noindent
设 $\{X_i \to X\}$ 是 $X$ 的光滑覆盖，并设 $y \in Y$。此时
$\{X_{i, y} \to X_y\}$ 是该纤维的光滑覆盖。因此 $\mathcal{P}_3$
在源的光滑拓扑下的局部性由 下降章引理
\ref{descent-lemma-normal-local-smooth} 得出。综合以上结果即得本引理。
\end{proof}



\section{正则态射}
\label{more-morphisms-section-regular}

\noindent
与第 \ref{more-morphisms-section-normal} 节比较。这一概念的代数版本见 More on Algebra,
第 \ref{more-algebra-section-regular} 节。

\begin{definition}
\label{more-morphisms-definition-regular}
设 $f : X \to Y$ 是概形态射。假设所有纤维 $X_y$ 都是局部诺特概形。
\begin{enumerate}
\item 设 $x \in X$ 且 $y = f(x)$。若 $f$ 在 $x$ 处平坦，并且概形
$X_y$ 在 $x$ 处于 $\kappa(y)$ 上几何正则（见 簇章定义
\ref{varieties-definition-geometrically-regular}），则称 $f$ 在
$x$ 处{\it 正则}。
\item 若 $f$ 在 $X$ 的每一点处都正则，则称 $f$ 是{\it 正则态射}。
\end{enumerate}
\end{definition}

\noindent
态射 $X \to Y$ 正则这一条件，比仅仅要求所有纤维都是正则局部诺特概形
更强。

\begin{lemma}
\label{more-morphisms-lemma-regular}
设 $f : X \to Y$ 是概形态射。假设 $f$ 的所有纤维都局部诺特。
以下条件等价：
\begin{enumerate}
\item $f$ 正则；
\item $f$ 平坦，且其纤维都是几何正则概形；
\item 对每一对满足 $f(U) \subset V$ 的仿射开集 $U \subset X$、
$V \subset Y$，环映射 $\mathcal{O}_Y(V) \to \mathcal{O}_X(U)$ 正则；
\item 存在开覆盖 $Y = \bigcup_{j \in J} V_j$ 以及开覆盖
$f^{-1}(V_j) = \bigcup_{i \in I_j} U_i$，使每个态射
$U_i \to V_j$ 都正则；
\item 存在仿射开覆盖 $Y = \bigcup_{j \in J} V_j$ 以及仿射开覆盖
$f^{-1}(V_j) = \bigcup_{i \in I_j} U_i$，使环映射
$\mathcal{O}_Y(V_j) \to \mathcal{O}_X(U_i)$ 都正则。
\end{enumerate}
\end{lemma}

\begin{proof}
由定义立即得到 (1) 与 (2) 等价。设 $x \in X$ 且 $y = f(x)$。
按定义，$f$ 在 $x$ 处平坦，当且仅当
$\mathcal{O}_{Y, y} \to \mathcal{O}_{X, x}$ 是平坦环映射；而 $X_y$
在 $x$ 处于 $\kappa(y)$ 上几何正则，当且仅当
$\mathcal{O}_{X_y, x} = \mathcal{O}_{X, x}/\mathfrak m_y\mathcal{O}_{X, x}$
是 $\kappa(y)$ 上的几何正则代数。因此，
% P05-ERR-0032: frozen source unexpectedly capitalizes "Whether" mid-sentence.
$f$ 在 $x$ 处是否正则，仅取决于局部环之间的局部同态
$\mathcal{O}_{Y, y} \to \mathcal{O}_{X, x}$。故 (1) 与 (4) 显然等价。

\medskip\noindent
回忆 代数进阶章定义 \ref{more-algebra-definition-regular}：
环映射 $A \to B$ 正则，当且仅当它平坦，并且对所有素理想
$\mathfrak p \subset A$，纤维环 $B \otimes_A \kappa(\mathfrak p)$
都是诺特且几何正则的。由 簇章引理
\ref{varieties-lemma-geometrically-regular}，这等价于
$\Spec(B \otimes_A \kappa(\mathfrak p))$ 是 $\kappa(\mathfrak p)$ 上的
几何正则概形。因此 (2) 推出 (3)，而 (3) 显然推出 (5)。最后假设 (5)。
这推出 $f$ 平坦（见 态射章引理
\ref{morphisms-lemma-flat-characterize}）。此外，若 $y \in Y$，则对某个
$j$ 有 $y \in V_j$，并且 $X_y = \bigcup_{i \in I_j} U_{i, y}$；由
簇章引理 \ref{varieties-lemma-geometrically-regular}，每个
$U_{i, y}$ 都在 $\kappa(y)$ 上几何正则。再次应用 簇章引理
\ref{varieties-lemma-geometrically-regular}，可知 $X_y$ 几何正则。
故 (2) 成立，本引理得证。
\end{proof}

\begin{lemma}
\label{more-morphisms-lemma-smooth-regular}
光滑态射是正则的。
\end{lemma}

\begin{proof}
设 $f : X \to Y$ 是光滑态射。由 态射章引理
\ref{morphisms-lemma-smooth-locally-finite-presentation}，$f$ 局部有限表示，
故纤维 $X_y$ 在某个域上局部有限型，从而局部诺特。此外，$f$ 平坦；见
态射章引理 \ref{morphisms-lemma-smooth-flat}。最后，纤维 $X_y$
在域上光滑（由 态射章引理
\ref{morphisms-lemma-base-change-smooth}），并因此由
% P05-ERR-0031: frozen source cites smooth-geometrically-normal for regularity.
簇章引理 \ref{varieties-lemma-smooth-geometrically-normal}
几何正则。于是由引理
\ref{more-morphisms-lemma-regular}，$f$ 正则。
\end{proof}

\begin{lemma}
\label{more-morphisms-lemma-regular-fppf-local-source-and-target}
性质 $\mathcal{P}(f)=$“$f$ 的纤维局部诺特且 $f$ 正则”在靶的 fppf
拓扑下是局部的，并在源的光滑拓扑下是局部的。
\end{lemma}

\begin{proof}
我们有
$\mathcal{P}(f) =
\mathcal{P}_1(f) \wedge \mathcal{P}_2(f) \wedge \mathcal{P}_3(f)$
，其中 $\mathcal{P}_1(f)=$“$f$ 的纤维局部诺特”，
$\mathcal{P}_2(f)=$“$f$ 平坦”，而
$\mathcal{P}_3(f)=$“$f$ 的纤维几何正则”。我们已经知道
$\mathcal{P}_1$ 与 $\mathcal{P}_2$ 在源与靶的 fppf 拓扑下都是局部的；见
引理 \ref{more-morphisms-lemma-locally-Noetherian-fibres-fppf-local-source-and-target}
以及 下降章引理 \ref{descent-lemma-descending-property-flat} 与
\ref{descent-lemma-flat-fpqc-local-source}。故只需处理 $\mathcal{P}_3$。

\medskip\noindent
设 $f : X \to Y$ 是概形态射。设
$\{\varphi_i : Y_i \to Y\}_{i \in I}$ 是 $Y$ 的 fpqc 覆盖，并以
$f_i : X_i \to Y_i$ 表示 $f$ 沿 $\varphi_i$ 的基变换。设 $i \in I$，
且 $y_i \in Y_i$ 是一点。令 $y = \varphi_i(y_i)$。注意
$$
X_{i, y_i} = \Spec(\kappa(y_i)) \times_{\Spec(\kappa(y))} X_y.
$$
因此在此情形中，$X_y$ 几何正则，当且仅当 $X_{i, y_i}$ 几何正则；见
簇章引理 \ref{varieties-lemma-geometrically-regular-upstairs}。
这表明 $\mathcal{P}_3$ 在靶的 fpqc 拓扑下是局部的。

\medskip\noindent
设 $\{X_i \to X\}$ 是 $X$ 的光滑覆盖，并设 $y \in Y$。此时
$\{X_{i, y} \to X_y\}$ 是该纤维的光滑覆盖。因此 $\mathcal{P}_3$
在源的光滑拓扑下的局部性由 下降章引理
\ref{descent-lemma-regular-local-smooth} 得出。综合以上结果即得本引理。
\end{proof}




\section{Cohen--Macaulay 态射}
\label{more-morphisms-section-CM}

\noindent
与第 \ref{more-morphisms-section-normal} 节比较。注意，正如 Algebra, Section
\ref{algebra-section-geometrically-CM} 与 Varieties, Section
\ref{varieties-section-CM} 所指出的，“几何 Cohen--Macaulay”等同于通常的
Cohen--Macaulay。

\begin{definition}
\label{more-morphisms-definition-CM}
设 $f : X \to Y$ 是概形态射。假设所有纤维 $X_y$ 都是局部诺特概形。
\begin{enumerate}
\item 设 $x \in X$ 且 $y = f(x)$。若 $f$ 在 $x$ 处平坦，并且概形
$X_y$ 在 $x$ 处的局部环是 Cohen--Macaulay 的，则称 $f$ 在
$x$ 处{\it Cohen--Macaulay}。
\item 若 $f$ 在 $X$ 的每一点处都 Cohen--Macaulay，则称 $f$ 是
{\it Cohen--Macaulay 态射}。
\end{enumerate}
\end{definition}

\noindent
下面给出一种等价表述。

\begin{lemma}
\label{more-morphisms-lemma-CM}
设 $f : X \to Y$ 是概形态射。假设 $f$ 的所有纤维都局部诺特。
以下条件等价：
\begin{enumerate}
\item $f$ 是 Cohen--Macaulay 态射；
\item $f$ 平坦，且其纤维都是 Cohen--Macaulay 概形。
\end{enumerate}
\end{lemma}

\begin{proof}
这直接由定义得出。
\end{proof}

\begin{lemma}
\label{more-morphisms-lemma-CM-dimension}
设 $f : X \to Y$ 是局部诺特概形之间的局部有限型 Cohen--Macaulay
态射。对 $X$ 中每个在 $Y$ 中的像为 $y$ 的点 $x$，有
$$
\dim_x(X) = \dim_y(Y) + \dim_x(X_y),
$$
，其中 $X_y$ 表示 $y$ 上方的纤维。
\end{lemma}

\begin{proof}
用 $x$ 的一个开邻域替换 $X$ 后，存在自然数 $d$，使 $X \to Y$
的所有纤维在每一点处的维数均为 $d$；见 态射章引理
\ref{morphisms-lemma-flat-finite-presentation-CM-fibres-relative-dimension}。
于是 $f$ 平坦、局部有限型，且相对维数为 $d$。故结论由 Morphisms,
引理 \ref{morphisms-lemma-rel-dimension-dimension} 得出。
\end{proof}

\begin{lemma}
\label{more-morphisms-lemma-composition-CM}
设 $f : X \to Y$ 与 $g : Y \to Z$ 是概形态射。假设 $f$、$g$ 与
$g \circ f$ 的纤维都局部诺特。设 $x \in X$ 的像分别为
$y \in Y$ 与 $z \in Z$。
\begin{enumerate}
\item 若 $f$ 在 $x$ 处 Cohen--Macaulay，且 $g$ 在 $f(x)$ 处
Cohen--Macaulay，则 $g \circ f$ 在 $x$ 处 Cohen--Macaulay。
\item 若 $f$ 与 $g$ 都是 Cohen--Macaulay 态射，则 $g \circ f$
也是 Cohen--Macaulay 态射。
\item 若 $g \circ f$ 在 $x$ 处 Cohen--Macaulay，且 $f$ 在 $x$
处平坦，则 $f$ 在 $x$ 处 Cohen--Macaulay，且 $g$ 在 $f(x)$ 处
Cohen--Macaulay。
\item 若 $g \circ f$ 是 Cohen--Macaulay 态射且 $f$ 平坦，则 $f$
是 Cohen--Macaulay 态射，并且 $g$ 在 $f$ 的像中每一点处都
Cohen--Macaulay。
\end{enumerate}
\end{lemma}

\begin{proof}
考虑诺特局部环的映射
$$
\mathcal{O}_{Y_z, y} \to \mathcal{O}_{X_z, x}
$$
并注意到其纤维为
$$
\mathcal{O}_{X_z, x}/\mathfrak m_{Y_z, y}\mathcal{O}_{X_z, x} =
\mathcal{O}_{X_y, x}
$$
% P05-ERR-0033: frozen source has ungrammatical "Thus the lemma this follows".
因此本引理由 代数章引理 \ref{algebra-lemma-CM-goes-up} 得出。
\end{proof}

\begin{lemma}
\label{more-morphisms-lemma-flat-morphism-from-CM-scheme}
设 $f : X \to Y$ 是局部诺特概形之间的平坦态射。若 $X$ 是
Cohen--Macaulay 概形，则 $f$ 是 Cohen--Macaulay 态射，并且对所有
$x \in X$，$\mathcal{O}_{Y, f(x)}$ 都是 Cohen--Macaulay 环。
\end{lemma}

\begin{proof}
化为代数陈述后，结论由 代数章引理
\ref{algebra-lemma-CM-goes-up} 得出。
\end{proof}

\begin{lemma}
\label{more-morphisms-lemma-base-change-CM}
设 $f : X \to Y$ 是概形态射，并假设所有纤维 $X_y$ 都是局部诺特概形。
设 $Y' \to Y$ 局部有限型，并设 $f' : X' \to Y'$ 是 $f$ 的基变换。
设点 $x' \in X'$ 的像为 $x \in X$。
\begin{enumerate}
\item 若 $f$ 在 $x$ 处 Cohen--Macaulay，则 $f' : X' \to Y'$
在 $x'$ 处 Cohen--Macaulay。
\item 若 $f$ 在 $x$ 处平坦，且 $f'$ 在 $x'$ 处 Cohen--Macaulay，
则 $f$ 在 $x$ 处 Cohen--Macaulay。
\item 若 $Y' \to Y$ 在 $f'(x')$ 处平坦，且 $f'$ 在 $x'$ 处
Cohen--Macaulay，则 $f$ 在 $x$ 处 Cohen--Macaulay。
\end{enumerate}
\end{lemma}

\begin{proof}
注意，关于 $Y' \to Y$ 的假设蕴含：对映到 $y \in Y$ 的
$y' \in Y'$，域扩张 $\kappa(y')/\kappa(y)$ 有限生成。因此所有纤维
$X'_{y'} = (X_y)_{\kappa(y')}$ 也局部诺特；见 簇章引理
\ref{varieties-lemma-locally-Noetherian-base-change}。故本引理的陈述有意义。
令 $y' = f'(x')$ 且 $y = f(x)$，得到局部环的交换图
$$
\xymatrix{
\mathcal{O}_{X', x'} & \mathcal{O}_{X, x} \ar[l] \\
\mathcal{O}_{Y', y'} \ar[u] & \mathcal{O}_{Y, y} \ar[l] \ar[u]
}
$$
，其中左上角是右上角与左下角在右下角上的张量积的一个局部化。

\medskip\noindent
假设 $f$ 在 $x$ 处 Cohen--Macaulay。映射
$\mathcal{O}_{Y, y} \to \mathcal{O}_{X, x}$ 的平坦性蕴含
$\mathcal{O}_{Y', y'} \to \mathcal{O}_{X', x'}$ 的平坦性；见 Algebra,
引理 \ref{algebra-lemma-base-change-flat-up-down}。环
$\mathcal{O}_{X, x}/\mathfrak m_y\mathcal{O}_{X, x}$ 是
Cohen--Macaulay 的，蕴含
$\mathcal{O}_{X', x'}/\mathfrak m_{y'}\mathcal{O}_{X', x'}$ 是
Cohen--Macaulay 的；见 簇章引理
\ref{varieties-lemma-CM-base-change}。故 $f'$ 在 $x'$ 处 Cohen--Macaulay。

\medskip\noindent
假设 $f$ 在 $x$ 处平坦，且 $f'$ 在 $x'$ 处 Cohen--Macaulay。
$\mathcal{O}_{X', x'}/\mathfrak m_{y'}\mathcal{O}_{X', x'}$ 是
Cohen--Macaulay 的，蕴含
$\mathcal{O}_{X, x}/\mathfrak m_y\mathcal{O}_{X, x}$ 是
Cohen--Macaulay 的；见 簇章引理
\ref{varieties-lemma-CM-base-change}。故 $f$ 在 $x$ 处 Cohen--Macaulay。

\medskip\noindent
假设 $Y' \to Y$ 在 $y'$ 处平坦，且 $f'$ 在 $x'$ 处
Cohen--Macaulay。映射
$\mathcal{O}_{Y', y'} \to \mathcal{O}_{X', x'}$ 与
$\mathcal{O}_{Y, y} \to \mathcal{O}_{Y', y'}$ 的平坦性蕴含
$\mathcal{O}_{Y, y} \to \mathcal{O}_{X, x}$ 的平坦性；见 Algebra,
引理 \ref{algebra-lemma-base-change-flat-up-down}。环
$\mathcal{O}_{X', x'}/\mathfrak m_{y'}\mathcal{O}_{X', x'}$ 是
Cohen--Macaulay 的，蕴含
$\mathcal{O}_{X, x}/\mathfrak m_y\mathcal{O}_{X, x}$ 是
Cohen--Macaulay 的；见 簇章引理
\ref{varieties-lemma-CM-base-change}。故 $f$ 在 $x$ 处 Cohen--Macaulay。
\end{proof}

\begin{lemma}
\label{more-morphisms-lemma-flat-finite-presentation-CM-open}
\begin{reference}
\cite[IV Corollary 12.1.7(iii)]{EGA}
\end{reference}
设 $f : X \to S$ 是平坦且局部有限表示的概形态射。令
$$
W = \{x \in X \mid f\text{ is Cohen-Macaulay at }x\}
$$
则
\begin{enumerate}
\item $W = \{x \in X \mid \mathcal{O}_{X_{f(x)}, x}\text{ 为 Cohen--Macaulay}\}$；
\item $W$ 在 $X$ 中是开的；
\item $W$ 在 $X \to S$ 的每个纤维中稠密；
\item $W$ 的形成与 $f$ 的任意基变换交换：对任意态射
$g : S' \to S$，考虑 $f$ 的基变换 $f' : X' \to S'$ 以及投影
$g' : X' \to X$。则态射 $f'$ 对应的集合 $W'$ 等于
$W' = (g')^{-1}(W)$。
\end{enumerate}
\end{lemma}

\begin{proof}
由于 $f$ 平坦且具有局部诺特纤维，(1) 中的等式由定义成立。(2) 与 (3)
由 代数章引理
\ref{algebra-lemma-generic-CM-flat-finite-presentation} 得出。(4) 可由
代数章引理 \ref{algebra-lemma-CM-locus-commutes-base-change} 或
簇章引理 \ref{varieties-lemma-CM-base-change} 得出。
\end{proof}

\begin{lemma}
\label{more-morphisms-lemma-flat-finite-presentation-characterize-CM}
设 $f : X \to S$ 是平坦且局部有限表示的概形态射。设 $x \in X$
的像为 $s \in S$，并令 $d = \dim_x(X_s)$。以下条件等价：
\begin{enumerate}
\item $f$ 在 $x$ 处 Cohen--Macaulay；
\item 存在 $x$ 的开邻域 $U \subset X$ 以及 $S$ 上的局部拟有限态射
$U \to \mathbf{A}^d_S$，它在 $x$ 处平坦；
\item 存在 $x$ 的开邻域 $U \subset X$ 以及 $S$ 上的局部拟有限平坦态射
$U \to \mathbf{A}^d_S$；
\item 对 $x$ 的任意开邻域 $U \subset X$ 及任意 $S$-态射
$g : U \to \mathbf{A}^d_S$，都有：$g$ 在 $x$ 处拟有限
$\Rightarrow$ $g$ 在 $x$ 处平坦。
\end{enumerate}
\end{lemma}

\begin{proof}
平坦性的开性说明 (2) 与 (3) 等价；见定理
\ref{more-morphisms-theorem-openness-flatness}。

\medskip\noindent
选择满足 $x \in U$ 的仿射开集 $U = \Spec(A) \subset X$，以及满足
$f(U) \subset V$ 的仿射开集 $V = \Spec(R) \subset S$。则 $R \to A$
是有限表示平坦环映射。设 $\mathfrak p \subset A$ 是 $x$ 对应的素理想。
用一个主局部化替换 $A$ 后，可以假设存在拟有限映射
$R[x_1, \ldots, x_d]  \to A$；见 代数章引理
\ref{algebra-lemma-quasi-finite-over-polynomial-algebra}。因此至少存在一对
$(U, g)$，由 $x$ 的开邻域 $U \subset X$ 与局部\footnote{若 $S$
拟分离，则 $g$ 将是拟有限的。}拟有限态射
$g : U \to \mathbf{A}^d_S$ 组成。

\medskip\noindent
断言：给定有限表示平坦映射 $R \to A$、素理想 $\mathfrak p \subset A$
以及在 $\mathfrak p$ 处拟有限的
$\varphi : R[x_1, \ldots, x_d] \to A$，则 $\Spec(\varphi)$ 在
$\mathfrak p$ 处平坦，当且仅当 $\Spec(A) \to \Spec(R)$ 在
$\mathfrak p$ 处 Cohen--Macaulay。事实上，由定理
\ref{more-morphisms-theorem-criterion-flatness-fibre}，平坦性可以在纤维上检验。
Cohen--Macaulay 性也如此（因为已经假设 $A$ 在 $R$ 上平坦）。
故该断言由 代数章引理 \ref{algebra-lemma-where-CM} 得出。

\medskip\noindent
该断言说明 (1) 与 (4) 等价；再结合第二段中已经构造出合适的
$(U, g)$ 这一事实，该断言还说明 (1) 与 (2) 等价。
\end{proof}

\begin{lemma}
\label{more-morphisms-lemma-flat-finite-presentation-CM-pieces}
设 $f : X \to S$ 是平坦且局部有限表示的概形态射。对 $d \geq 0$，
存在开集 $U_d \subset X$，满足以下性质：
\begin{enumerate}
\item $W = \bigcup_{d \geq 0} U_d$ 在 $f$ 的每个纤维中稠密；
\item $U_d \to S$ 的相对维数为 $d$（见 态射章定义
\ref{morphisms-definition-relative-dimension-d}）。
\end{enumerate}
\end{lemma}

\begin{proof}
这由引理 \ref{more-morphisms-lemma-flat-finite-presentation-CM-open} 与 态射章引理
\ref{morphisms-lemma-flat-finite-presentation-CM-fibres-relative-dimension}
结合得出。
\end{proof}

\begin{lemma}
\label{more-morphisms-lemma-flat-finite-presentation-specialization-dimension}
设 $f : X \to S$ 是平坦且局部有限表示的概形态射。假设
$x' \leadsto x$ 是 $X$ 中点的特化，其在 $S$ 中的像为
$s' \leadsto s$。若 $x$ 是 $X_s$ 某个不可约分支的泛点，则
$\dim_{x'}(X_{s'}) = \dim_x(X_s)$。
\end{lemma}

\begin{proof}
对某个 $d$，点 $x$ 包含在 $U_d$ 中，其中 $U_d$ 如引理
\ref{more-morphisms-lemma-flat-finite-presentation-CM-pieces} 所述。
\end{proof}

\begin{lemma}
\label{more-morphisms-lemma-CM-local-source-and-target}
性质 $\mathcal{P}(f)=$“$f$ 的纤维局部诺特且 $f$ 是
Cohen--Macaulay 态射”在靶的 fppf 拓扑下是局部的，并在源的 syntomic
拓扑下是局部的。
\end{lemma}

\begin{proof}
我们有
$\mathcal{P}(f) =
\mathcal{P}_1(f) \wedge \mathcal{P}_2(f)$
，其中 $\mathcal{P}_1(f)=$“$f$ 平坦”，而
$\mathcal{P}_2(f)=$“$f$ 的纤维局部诺特且 Cohen--Macaulay”。已知
$\mathcal{P}_1$ 在源与靶的 fppf 拓扑下都是局部的；见 下降章引理
\ref{descent-lemma-descending-property-flat} 与
\ref{descent-lemma-flat-fpqc-local-source}。故只需处理 $\mathcal{P}_2$。

\medskip\noindent
设 $f : X \to Y$ 是概形态射。设
$\{\varphi_i : Y_i \to Y\}_{i \in I}$ 是 $Y$ 的 fppf 覆盖，并以
$f_i : X_i \to Y_i$ 表示 $f$ 沿 $\varphi_i$ 的基变换。设 $i \in I$，
且 $y_i \in Y_i$ 是一点。令 $y = \varphi_i(y_i)$。注意
$$
X_{i, y_i} = \Spec(\kappa(y_i)) \times_{\Spec(\kappa(y))} X_y.
$$
，并且 $\kappa(y_i)/\kappa(y)$ 是有限生成域扩张。因此若 $X_y$
局部诺特，则 $X_{i, y_i}$ 局部诺特；见 簇章引理
\ref{varieties-lemma-locally-Noetherian-base-change}。若此外 $X_y$
是 Cohen--Macaulay 概形，则 $X_{i, y_i}$ 也是 Cohen--Macaulay 概形；
见 簇章引理 \ref{varieties-lemma-CM-base-change}。因此
$\mathcal{P}_2$ 在靶的 fppf 拓扑下是局部的。

\medskip\noindent
设 $\{X_i \to X\}$ 是 $X$ 的 syntomic 覆盖，并设 $y \in Y$。此时
$\{X_{i, y} \to X_y\}$ 是该纤维的 syntomic 覆盖。因此
$\mathcal{P}_2$ 在源的 syntomic 拓扑下的局部性由 下降章引理
\ref{descent-lemma-CM-local-syntomic} 得出。综合以上结果即得本引理。
\end{proof}






\section{切割 Cohen--Macaulay 态射}
\label{more-morphisms-section-slicing-CM}

\noindent
本节的结果最终导向如下断言：fppf 拓扑与“有限表示、平坦、拟有限”拓扑
相同。下一个引理与 除子章引理
\ref{divisors-lemma-fibre-Cartier} 密切相关。

\begin{lemma}
\label{more-morphisms-lemma-slice-given-element}
设 $f : X \to S$ 是概形态射。设点 $x \in X$ 的像为 $s \in S$，并设
$h \in \mathfrak m_x \subset \mathcal{O}_{X, x}$。假设
\begin{enumerate}
\item $f$ 局部有限表示；
\item $f$ 在 $x$ 处平坦；
\item $h$ 在
$\mathcal{O}_{X_s, x} = \mathcal{O}_{X, x}/\mathfrak m_s\mathcal{O}_{X, x}$
中的像 $\overline{h}$ 是非零因子。
\end{enumerate}
则存在 $x$ 的仿射开邻域 $U \subset X$，使 $h$ 来自
$h \in \Gamma(U, \mathcal{O}_U)$，并且 $D = V(h)$ 是 $U$ 中包含
$x \in D$ 的有效 Cartier 除子，且 $D \to S$ 平坦并局部有限表示。
\end{lemma}

\begin{proof}
我们把问题化到诺特情形来证明。由平坦性的开性（见定理
\ref{more-morphisms-theorem-openness-flatness}），用 $x$ 的一个开邻域替换 $X$ 后，
可以假设 $X \to S$ 平坦。还可假设 $X$ 与 $S$ 仿射。稍微缩小 $X$ 后，
可以假设存在 $h \in \Gamma(X, \mathcal{O}_X)$，它映到给定的 $h$。

\medskip\noindent
可写 $S = \Spec(A)$，并把 $A = \colim_i A_i$ 写成有限型
$\mathbf{Z}$-代数的定向余极限。于是，由 代数章引理
\ref{algebra-lemma-flat-finite-presentation-limit-flat}，或 极限章引理
\ref{limits-lemma-descend-finite-presentation}、
\ref{limits-lemma-descend-affine-finite-presentation} 以及
% P05-ERR-0034: frozen source repeats limits-lemma-descend-finite-presentation.
\ref{limits-lemma-descend-finite-presentation}，可以找到笛卡尔图
$$
\xymatrix{
X \ar[r] \ar[d]_f & X_0 \ar[d]^{f_0} \\
S \ar[r] & S_0
}
$$
其中 $f_0$ 平坦且有限表示，$X_0$ 仿射，而 $S_0$ 仿射且诺特。令
$x_0 \in X_0$ 与相应的 $s_0 \in S_0$ 分别为 $x$ 与相应的 $s$ 的像。
还可假设存在元素 $h_0 \in \Gamma(X_0, \mathcal{O}_{X_0})$，其在 $X$
上的限制为 $h$。（若使用上面的代数参考，这是显然的；若使用极限一章
的参考，则把 $h$ 看作态射 $X \to \mathbf{A}^1_S$，结论由 Limits,
引理 \ref{limits-lemma-descend-finite-presentation} 得出。）注意，
$\mathcal{O}_{X_s, x}$ 是
$\mathcal{O}_{(X_0)_{s_0}, x_0} \otimes_{\kappa(s_0)} \kappa(s)$
的一个局部化，故
$\mathcal{O}_{(X_0)_{s_0}, x_0} \to \mathcal{O}_{X_s, x}$ 是平坦局部
环映射，特别是忠实平坦的。因此像
$\overline{h}_0 \in \mathcal{O}_{(X_0)_{s_0}, x_0}$ 属于
$\mathfrak m_{(X_0)_{s_0}, x_0}$，且是非零因子。我们断言：用 $x_0$
的一个主开邻域替换 $X_0$ 后，元素 $h_0$ 是
$B_0 = \Gamma(X_0, \mathcal{O}_{X_0})$ 中的非零因子，并且
$B_0/h_0B_0$ 在 $A_0 = \Gamma(S_0, \mathcal{O}_{S_0})$ 上平坦。
若断言成立，则
$$
0 \to B_0 \xrightarrow{h_0} B_0 \to B_0/h_0B_0 \to 0
$$
是平坦 $A_0$-模的短正合列。故与 $A$ 作张量后仍正合（由 Algebra,
引理 \ref{algebra-lemma-flat-tor-zero}），从而本引理得证。

\medskip\noindent
只剩证明上述断言。相应的代数陈述如下（这里略去下标 ${}_0$）：设
$A \to B$ 是诺特环之间的有限型平坦环映射。设素理想
$\mathfrak q \subset B$ 位于 $\mathfrak p \subset A$ 上方。假设
$h \in \mathfrak q$ 在 $B_{\mathfrak q}/\mathfrak p B_{\mathfrak q}$
中的像是非零因子。目标是证明：
% P05-ERR-0035: frozen source has ungrammatical "after possible replacing".
对某个 $g \in B$、$g \not \in \mathfrak q$，以 $B_g$ 替换 $B$ 后，
元素 $h$ 成为非零因子，而 $B/hB$ 在 $A$ 上平坦。由 代数章引理
\ref{algebra-lemma-grothendieck}，$h$ 是 $B_{\mathfrak q}$ 中的非零因子，
且 $B_{\mathfrak q}/hB_{\mathfrak q}$ 在 $A$ 上平坦。由平坦性的开性，
见 代数章定理 \ref{algebra-theorem-openness-flatness} 或定理
\ref{more-morphisms-theorem-openness-flatness}，对某个 $g \in B$、
$g \not \in \mathfrak q$，以 $B_g$ 替换 $B$ 后，$B/hB$ 在 $A$ 上平坦。
最后，令 $I = \{b \in B \mid hb = 0\}$ 为 $h$ 的零化子。由于 $h$
是 $B_{\mathfrak q}$ 中的非零因子，$IB_{\mathfrak q} = 0$。又由于
$B$ 诺特，$I$ 有限生成。因此存在 $g \in B$、$g \not \in \mathfrak q$，
使 $IB_g = 0$。以 $B_g$ 替换 $B$ 后，$h$ 是非零因子。
\end{proof}

\begin{lemma}
\label{more-morphisms-lemma-slice-given-elements}
设 $f : X \to S$ 是概形态射。设点 $x \in X$ 的像为 $s \in S$，并设
$h_1, \ldots, h_r \in \mathcal{O}_{X, x}$。假设
\begin{enumerate}
\item $f$ 局部有限表示；
\item $f$ 在 $x$ 处平坦；
\item $h_1, \ldots, h_r$ 在
$\mathcal{O}_{X_s, x} = \mathcal{O}_{X, x}/\mathfrak m_s\mathcal{O}_{X, x}$
中的像构成正则序列。
\end{enumerate}
则存在 $x$ 的仿射开邻域 $U \subset X$，使 $h_1, \ldots, h_r$ 来自
$h_1, \ldots, h_r \in \Gamma(U, \mathcal{O}_U)$，并且
$Z = V(h_1, \ldots, h_r) \to U$ 是满足 $x \in Z$ 的正则浸入，而
$Z \to S$ 平坦且局部有限表示。此外，对 $S$ 上的任意概形 $S'$，
基变换 $Z_{S'} \to U_{S'}$ 都是正则浸入。
\end{lemma}

\begin{proof}
（我们关于正则序列的约定蕴含：对每个 $i$，$h_i \in \mathfrak m_x$。）
当 $r = 1$ 时，结论由引理 \ref{more-morphisms-lemma-slice-given-element} 与 Divisors,
引理 \ref{divisors-lemma-relative-Cartier} 结合得出；后一个引理说明
$V(h_1)$ 在基变换后仍是有效 Cartier 除子。当 $r > 1$ 时，对 $r$
作直接归纳即可（恰好应用 $r$ 次 $r = 1$ 的结果；略去细节）。

\medskip\noindent
另一种证明方法使用 Divisors, Section
\ref{divisors-section-relative-regular-immersion} 的材料。首先，由平坦性的
开性（见定理 \ref{more-morphisms-theorem-openness-flatness}），用 $x$ 的一个开邻域替换
$X$ 后，可以假设 $X \to S$ 平坦。还可假设 $X$ 与 $S$ 仿射。稍微缩小
$X$ 后，可以假设 $h_1, \ldots, h_r \in \Gamma(X, \mathcal{O}_X)$。
令 $Z = V(h_1, \ldots, h_r)$。注意，$X_s$ 是诺特概形（因为它是代数
$\kappa(s)$-概形；见 Varieties, Section
\ref{varieties-section-algebraic-schemes}），且 $X_s$ 上的拓扑由 $X$
上的拓扑诱导（见 概形章引理
\ref{schemes-lemma-fibre-topological}）。故进一步缩小 $X$ 后，可以假设
$Z_s \subset X_s$ 是由 $r$ 个元素 $h_i|_{X_s}$ 截出的正则浸入；见
除子章引理 \ref{divisors-lemma-Noetherian-scheme-regular-ideal}
及其证明。还显然有 $r = \dim_x(X_s) - \dim_x(Z_s)$，因为
\begin{align*}
\dim_x(X_s) & = \dim(\mathcal{O}_{X_s, x}) +
\text{trdeg}_{\kappa(s)}(\kappa(x)), \\
\dim_x(Z_s) & = \dim(\mathcal{O}_{Z_s, x}) +
\text{trdeg}_{\kappa(s)}(\kappa(x)), \\
\dim(\mathcal{O}_{X_s, x}) & = \dim(\mathcal{O}_{Z_s, x}) + r
\end{align*}
前两个等式由 代数章引理
\ref{algebra-lemma-dimension-at-a-point-finite-type-field} 得出，
% P05-ERR-0036: frozen source says "the second" after already citing first two.
而第二个等式由把 代数章引理
\ref{algebra-lemma-one-equation} 应用 $r$ 次得出。因此 除子章引理
\ref{divisors-lemma-fibre-quasi-regular} 的 (3) 表明（在 Zariski 意义下
缩小 $X$ 后）态射 $Z \to X$ 是正则浸入，可以对它应用 除子章引理
\ref{divisors-lemma-relative-regular-immersion-flat-in-neighbourhood}
（该引理给出平坦性与基变换陈述）。
\end{proof}

\begin{lemma}
\label{more-morphisms-lemma-slice-once}
设 $f : X \to S$ 是概形态射。设点 $x \in X$ 的像为 $s \in S$。假设
\begin{enumerate}
\item $f$ 局部有限表示；
\item $f$ 在 $x$ 处平坦；
\item $\mathcal{O}_{X_s, x}$ 具有 $\text{depth} \geq 1$。
\end{enumerate}
则存在 $x$ 的仿射开邻域 $U \subset X$，以及包含 $x$ 的有效 Cartier
除子 $D \subset U$，使 $D \to S$ 平坦且有限表示。
\end{lemma}

\begin{proof}
选取任意 $h \in \mathfrak m_x \subset \mathcal{O}_{X, x}$，使它在
$\mathcal{O}_{X_s, x}$ 中的像是非零因子，然后应用引理
\ref{more-morphisms-lemma-slice-given-element}。
\end{proof}

\begin{lemma}
\label{more-morphisms-lemma-slice-CM}
\begin{reference}
\cite[IV Proposition 17.16.1]{EGA}
\end{reference}
设 $f : X \to S$ 是概形态射。设点 $x \in X$ 的像为 $s \in S$。假设
\begin{enumerate}
\item $f$ 局部有限表示；
\item $f$ 在 $x$ 处 Cohen--Macaulay；
\item $x$ 是 $X_s$ 的闭点。
\end{enumerate}
则存在包含 $x$ 的正则浸入 $Z \to X$，使得
\begin{enumerate}
\item[(a)] $Z \to S$ 平坦且局部有限表示；
\item[(b)] $Z \to S$ 局部拟有限；
\item[(c)] 集合论意义下 $Z_s = \{x\}$。
\end{enumerate}
\end{lemma}

\begin{proof}
不妨并且确实用 $s$ 的一个仿射开邻域替换 $S$。对仿射 $S$，我们将对
$d = \dim_x(X_s)$ 作归纳来证明本引理。

\medskip\noindent
情形 $d = 0$。此时证明可以取 $Z$ 为 $x$ 的一个开邻域。（注意，开浸入
是正则浸入。）事实上，若 $d = 0$，则 $X \to S$ 在 $x$ 处拟有限；见
态射章引理
\ref{morphisms-lemma-locally-quasi-finite-rel-dimension-0}。因此存在仿射开邻域
$U \subset X$，使 $U \to S$ 拟有限；见 态射章引理
\ref{morphisms-lemma-quasi-finite-points-open}。用 $U$ 替换 $X$ 后，纤维
$X_s$ 是有限离散集。
% P05-ERR-0037: frozen source says an affine open neighbourhood "of X", likely x.
于是进一步用 $X$ 的一个仿射开邻域替换 $X$ 后，有
$f^{-1}(\{s\}) = \{x\}$（因为 $X_s$ 上的拓扑由 $X$ 上的拓扑诱导；
见 概形章引理 \ref{schemes-lemma-fibre-topological}）。
这证明了该情形。

\medskip\noindent
下面假设 $d > 0$。注意，由于 $x$ 是其纤维的闭点，扩张
$\kappa(x)/\kappa(s)$ 有限（由 Hilbert 零点定理；见 态射章引理
\ref{morphisms-lemma-closed-point-fibre-locally-finite-type}）。因此
$$
\text{depth}(\mathcal{O}_{X_s, x}) = \dim(\mathcal{O}_{X_s, x}) = d > 0
$$
，第一个等式因为 $\mathcal{O}_{X_s, x}$ 是 Cohen--Macaulay 环而成立，
第二个等式由 态射章引理
\ref{morphisms-lemma-dimension-fibre-at-a-point} 得出。因此可应用引理
\ref{more-morphisms-lemma-slice-once} 找到图
$$
\xymatrix{
D \ar[r] \ar[rrd] & U \ar[r] \ar[rd] & X \ar[d] \\
& & S
}
$$
，其中 $x \in D$。注意，对某个非零因子 $\overline{h}$，有
$\mathcal{O}_{D_s, x} = \mathcal{O}_{X_s, x}/(\overline{h})$；见
除子章引理 \ref{divisors-lemma-relative-Cartier}。因此
$\mathcal{O}_{D_s, x}$ 是 Cohen--Macaulay 环，其维数比
$\mathcal{O}_{X_s, x}$ 的维数小一；例如见 代数章引理
\ref{algebra-lemma-reformulate-CM}。于是态射 $D \to S$ 平坦、局部有限表示，
并在 $x$ 处 Cohen--Macaulay，且
$\dim_x(D_s) = \dim_x(X_s) - 1 = d - 1$。由归纳假设，可以找到具有
(a)、(b)、(c) 性质的正则浸入 $Z \to D$。由于 $Z \to D \to U$
都是正则浸入，由 除子章引理
\ref{divisors-lemma-composition-regular-immersion}，$Z \to U$ 也是
正则浸入。证明完毕。
\end{proof}

\begin{lemma}
\label{more-morphisms-lemma-qf-fp-flat-neighbourhood-dominates-fppf}
设 $f : X \to S$ 是平坦且局部有限表示的概形态射。设 $s \in S$
是 $f$ 的像中的一点。则存在交换图
$$
\xymatrix{
S' \ar[rr] \ar[rd]_g & & X \ar[ld]^f \\
& S
}
$$
，其中 $g : S' \to S$ 平坦、局部有限表示、局部拟有限，且
$s \in g(S')$。
\end{lemma}

\begin{proof}
由假设，纤维 $X_s$ 非空。因此存在闭点 $x \in X_s$，使 $f$ 在该点处
Cohen--Macaulay；见引理 \ref{more-morphisms-lemma-flat-finite-presentation-CM-open}。
应用引理 \ref{more-morphisms-lemma-slice-CM}，并令 $S' = Z$。
\end{proof}

\noindent
下一个引理说明，fppf 拓扑的层与“拟有限、平坦、有限表示”拓扑的层是
同一回事。

\begin{lemma}
\label{more-morphisms-lemma-qf-fp-flat-dominates-fppf}
设 $S$ 是概形。设 $\mathcal{U} = \{S_i \to S\}_{i \in I}$ 是 $S$
的 fppf 覆盖；见 拓扑章定义
\ref{topologies-definition-fppf-covering}。则存在加细 $\mathcal{U}$
（见 位点章定义 \ref{sites-definition-morphism-coverings}）的 fppf
覆盖 $\mathcal{V} = \{T_j \to S\}_{j \in J}$，使每个
$T_j \to S$ 都局部拟有限。
\end{lemma}

\begin{proof}
对每个 $s \in S$，存在 $i \in I$，使 $s$ 属于 $S_i \to S$ 的像。
由引理 \ref{more-morphisms-lemma-qf-fp-flat-neighbourhood-dominates-fppf}，可以找到态射
$g_s : T_s \to S$，使 $s \in g_s(T_s)$；该态射平坦、局部有限表示、
局部拟有限，并且 $g_s$ 经由 $S_i \to S$ 分解。因此
$\{T_s \to S\}$ 即为所求加细 $\mathcal{U}$ 的 $S$ 的覆盖。
\end{proof}



\section{泛纤维}
\label{more-morphisms-section-generic}

\noindent
这里给出若干关于泛纤维与邻近纤维之间关系的结果。

\begin{lemma}
\label{more-morphisms-lemma-empty-generic-fibre}
设 $f : X \to Y$ 是有限型概形态射。假设 $Y$ 不可约，泛点为 $\eta$。
若 $X_\eta = \emptyset$，则存在非空开集 $V \subset Y$，使
$X_V = V \times_Y X = \emptyset$。
\end{lemma}

\begin{proof}
这立即由更一般的 态射章引理
\ref{morphisms-lemma-quasi-compact-generic-point-not-in-image} 得出。
\end{proof}

\begin{lemma}
\label{more-morphisms-lemma-nonempty-generic-fibre}
设 $f : X \to Y$ 是有限型概形态射。假设 $Y$ 不可约，泛点为 $\eta$。
若 $X_\eta \not = \emptyset$，则存在非空开集 $V \subset Y$，使
$X_V = V \times_Y X \to V$ 满射。
\end{lemma}

\begin{proof}
取仿射开集后，这由 代数章引理
\ref{algebra-lemma-characterize-image-finite-type} 得出。（当然，这也可由
泛平坦性得出。）
\end{proof}

\begin{lemma}
\label{more-morphisms-lemma-nowhere-dense-generic-fibre}
设 $f : X \to Y$ 是有限型概形态射。假设 $Y$ 不可约，泛点为 $\eta$。
若 $Z \subset X$ 是闭子集，且 $Z_\eta$ 在 $X_\eta$ 中处处不稠密，
则存在非空开集 $V \subset Y$，使对所有 $y \in V$，$Z_y$ 在 $X_y$
中处处不稠密。
\end{lemma}

\begin{proof}
设 $Y' \subset Y$ 是 $Y$ 的约化，并令 $X' = Y' \times_Y X$、
$Z' = Y' \times_Y Z$。由 态射章引理
\ref{morphisms-lemma-reduction-universal-homeomorphism}，$Y' \to Y$
是泛同胚，故只需对 $Z' \subset X' \to Y'$ 证明本引理。因此可假设
$Y$ 是整概形；见 性质章引理
\ref{properties-lemma-characterize-integral}。由 态射章命题
\ref{morphisms-proposition-generic-flatness}，存在非空仿射开集
$V \subset Y$，使 $X_V \to V$ 与 $Z_V \to V$ 都平坦且有限表示。
我们断言这个 $V$ 可用。取 $y \in V$。若 $Z_y$ 有非空内部，则 $Z_y$
包含 $X_y$ 某个不可约分支的泛点 $\xi$。注意
$\eta \leadsto f(\xi)$。由于 $Z_V \to V$ 平坦，可以选择特化
$\xi' \leadsto \xi$，其中 $\xi' \in Z$ 且 $f(\xi') = \eta$；见
态射章引理 \ref{morphisms-lemma-generalizations-lift-flat}。由引理
\ref{more-morphisms-lemma-flat-finite-presentation-specialization-dimension}，有
$$
\dim_{\xi'}(Z_\eta) = \dim_{\xi}(Z_y) = \dim_{\xi}(X_y) = \dim_{\xi'}(X_\eta).
$$
因此，$Z_\eta$ 中某个经过 $\xi'$ 的不可约分支具有维数
$\dim_{\xi'}(X_\eta)$，这与 $Z_\eta$ 在 $X_\eta$ 中处处不稠密的假设
矛盾，结论得证。
\end{proof}

\begin{lemma}
\label{more-morphisms-lemma-scheme-theoretically-dense-generic-fibre}
设 $f : X \to Y$ 是有限型概形态射。假设 $Y$ 不可约，泛点为 $\eta$。
设 $U \subset X$ 是开子概形，且 $U_\eta$ 在 $X_\eta$ 中概形论稠密。
则存在非空开集 $V \subset Y$，使对所有 $y \in V$，$U_y$ 在 $X_y$
中概形论稠密。
\end{lemma}

\begin{proof}
设 $Y' \subset Y$ 是 $Y$ 的约化，并令 $X' = Y' \times_Y X$、
$U' = Y' \times_Y U$。由于 $Y' \to Y$ 在点上诱导双射，而
$U' \to U$ 与 $X' \to X$ 在概形论纤维上诱导同构，可以用 $Y'$ 替换
$Y$、用 $X'$ 替换 $X$。因此可假设 $Y$ 是整概形；见 性质章引理
\ref{properties-lemma-characterize-integral}。还可以用非空仿射开集替换
$Y$。换言之，可假设 $Y = \Spec(A)$，其中 $A$ 是分式域为 $K$ 的整环。

\medskip\noindent
由于 $f$ 有限型，$X$ 拟紧。对某些仿射开集 $X_i$，写成
$X = X_1 \cup \ldots \cup X_n$。由 态射章定义
\ref{morphisms-definition-scheme-theoretically-dense}，
$U_i = X_i \cap U$ 是 $X_i$ 的开子概形，且 $U_{i, \eta}$ 在
$X_{i, \eta}$ 中概形论稠密。因此只需对各对 $(X_i, U_i)$ 证明结论；
换言之，可假设 $X$ 仿射。

\medskip\noindent
写 $X = \Spec(B)$。注意，$B_K$ 是有限型 $K$-代数，故是诺特的。
因此 $U_\eta$ 拟紧。于是可找到有限多个 $g_1, \ldots, g_m \in B$，
使 $D(g_j) \subset U$，且
$U_\eta = D(g_1)_\eta \cup \ldots \cup D(g_m)_\eta$。$U_\eta$ 在
$X_\eta$ 中概形论稠密，意味着
$B_K \to \bigoplus_j (B_K)_{g_j}$ 是单射；见 态射章例
\ref{morphisms-example-scheme-theoretic-closure}。由 代数章引理
\ref{algebra-lemma-when-injective-covering}，这等价于
$B_K \to \bigoplus\nolimits_{j = 1, \ldots, m} B_K$，
$b \mapsto (g_1b, \ldots, g_mb)$ 是单射。令 $M$ 为此映射在 $A$
上的余核；也就是说，有正合列
$$
0 \to I \to B \xrightarrow{(g_1, \ldots, g_m)}
\bigoplus\nolimits_{j = 1, \ldots, m} B \to M \to 0
$$
对某个非零 $h$，用 $A_h$ 替换 $A$ 后，可以假设 $B$ 是平坦有限表示
$A$-代数，且 $M$ 在 $A$ 上平坦；见 代数章引理
\ref{algebra-lemma-generic-flatness}。$B$ 在 $A$ 上平坦，蕴含 $B$
作为 $A$-模无挠；见 代数进阶章引理
\ref{more-algebra-lemma-flat-torsion-free}。故 $B \subset B_K$。
由假设 $I_K = 0$，这蕴含 $I = 0$（因为
$I \subset B \subset B_K$ 是 $I_K$ 的子集）。于是现在有短正合列
$$
0 \to B \xrightarrow{(g_1, \ldots, g_m)}
\bigoplus\nolimits_{j = 1, \ldots, m} B \to M \to 0
$$
，其中 $M$ 在 $A$ 上平坦。因此，对每个到域 $\kappa$ 的同态
$A \to \kappa$，得到短正合列
$$
0 \to B \otimes_A \kappa \xrightarrow{(g_1 \otimes 1, \ldots, g_m \otimes 1)}
\bigoplus\nolimits_{j = 1, \ldots, m} B \otimes_A \kappa \to
M \otimes_A \kappa \to 0
$$
；见 代数章引理 \ref{algebra-lemma-flat-tor-zero}。反向应用上述论证，
这意味着 $\bigcup D(g_j \otimes 1)$ 在
$\Spec(B \otimes_A \kappa)$ 中概形论稠密。由于
$\bigcup D(g_j \otimes 1) = \bigcup D(g_j)_\kappa \subset U_\kappa$，
可知 $U_\kappa$ 在 $X_\kappa$ 中概形论稠密，这正是所要证明的。
\end{proof}

\noindent
设有概形态射 $f : X \to Y$ 与点 $y \in Y$。回忆纤维 $X_y$ 同胚于
$X$ 的子集 $f^{-1}(\{y\})$，后者赋予诱导拓扑；见 概形章引理
\ref{schemes-lemma-fibre-topological}。设有闭子集 $T(y) \subset X_y$。
令 $T$ 为 $T(y)$ 在 $X$ 中的闭包，并赋予 $T$ 诱导的约化概形结构。
则 $T$ 是 $X$ 的闭子概形，并且集合论意义下 $T_y = T(y)$。事实上，
$T$ 是具有此性质的 $X$ 的最小闭子概形。因此，若愿意，把 $X_y$
的闭子集记作 $T_y$ 是“无害的”。在下一个引理中，我们把这一点应用于
$f$ 的泛纤维。

\begin{lemma}
\label{more-morphisms-lemma-cover-generic-fibre-neighbourhood}
设 $f : X \to Y$ 是有限型概形态射。假设 $Y$ 不可约，泛点为 $\eta$。
设 $X_\eta = Z_{1, \eta} \cup \ldots \cup Z_{n, \eta}$ 是泛纤维的覆盖，
其中各项都是 $X_\eta$ 的闭子集。令 $Z_i$ 为 $Z_{i, \eta}$ 在 $X$
中的闭包（见上面的讨论）。则存在非空开集 $V \subset Y$，使对所有
$y \in V$，都有 $X_y = Z_{1, y} \cup \ldots \cup Z_{n, y}$。
\end{lemma}

\begin{proof}
若 $Y$ 诺特，则 $U = X \setminus (Z_1 \cup \ldots \cup Z_n)$ 在 $Y$
上有限型，可以直接应用引理 \ref{more-morphisms-lemma-empty-generic-fibre}，得到对某个
非空开集 $V \subset Y$ 有 $U_V = \emptyset$。一般情形下论证如下。
由于问题是拓扑性的，可以用其约化替换 $Y$。于是 $Y$ 是整概形；
见 性质章引理 \ref{properties-lemma-characterize-integral}。缩小 $Y$
后，可以假设 $X \to Y$ 平坦；见 态射章命题
\ref{morphisms-proposition-generic-flatness}。此时，由 态射章引理
\ref{morphisms-lemma-generalizations-lift-flat}，$X_y$ 中每个点 $x$
都是某个点 $x' \in X_\eta$ 的特化。由于各 $Z_i$ 在 $X$ 中闭，并覆盖
泛纤维，这蕴含对 $y \in Y$ 有 $X_y = \bigcup Z_{i, y}$，正合所求。
\end{proof}

\noindent
下一个引理说明，源约化的态射具有约化泛纤维。

\begin{lemma}
\label{more-morphisms-lemma-reduction-generic-fibre}
设 $f : X \to Y$ 是概形态射。设 $\eta \in Y$ 是 $Y$ 某个不可约分支
的泛点。则
$(X_\eta)_{red} = (X_{red})_\eta$.
\end{lemma}

\begin{proof}
选择 $\eta$ 的仿射邻域 $\Spec(A) \subset Y$，并选择经态射 $f$ 映入
$\Spec(A)$ 的仿射开集 $\Spec(B) \subset X$。设
$\mathfrak p \subset A$ 是 $\eta$ 对应的极小素理想。设 $B_{red}$
为 $B$ 对幂零根 $\sqrt{(0)}$ 的商。本引理的代数内容是
$C = B_{red} \otimes_A \kappa(\mathfrak p)$ 约化。以 $I \subset A$
表示幂零根，故 $A_{red} = A/I$。令
$\mathfrak p_{red} = \mathfrak p/I$；这是 $A_{red}$ 的极小素理想，且
$\kappa(\mathfrak p) = \kappa(\mathfrak p_{red})$。由于
$A \to B_{red}$ 与 $A \to \kappa(\mathfrak p)$ 都经由
$A \to A_{red}$ 分解，有
$C = B_{red} \otimes_{A_{red}} \kappa(\mathfrak p_{red})$。由 Algebra,
引理 \ref{algebra-lemma-minimal-prime-reduced-ring}，
$\kappa(\mathfrak p_{red}) = (A_{red})_{\mathfrak p_{red}}$ 是一个局部化。
故 $C$ 是 $B_{red}$ 的一个局部化（代数章引理
\ref{algebra-lemma-tensor-localization}），从而约化。
\end{proof}

\begin{lemma}
\label{more-morphisms-lemma-make-generic-fibre-geometrically-reduced}
设 $f : X \to Y$ 是概形态射。假设 $Y$ 不可约且 $f$ 有限型。
存在图
$$
\xymatrix{
X' \ar[d]_{f'} \ar[r]_{g'} & X_V \ar[r] \ar[d] & X \ar[d]^f \\
Y' \ar[r]^g & V \ar[r] & Y
}
$$
，其中
\begin{enumerate}
\item $V$ 是 $Y$ 的非空开集；
\item $X_V = V \times_Y X$；
\item $g : Y' \to V$ 是有限泛同胚；
\item $X' = (Y' \times_Y X)_{red} = (Y' \times_V X_V)_{red}$；
\item $g'$ 是有限泛同胚；
\item $Y'$ 是整仿射概形；
\item $f'$ 平坦且有限表示；
\item $f'$ 的泛纤维几何约化。
\end{enumerate}
\end{lemma}

\begin{proof}
设 $V = \Spec(A)$ 是 $Y$ 的非空仿射开集。由假设，$A$ 的幂零根
$\mathfrak p = \sqrt{(0)}$ 是素理想。令 $K = \kappa(\mathfrak p)$。
令 $p$ 在 $K$ 的特征为正时等于其特征，在特征为零时等于 $1$。
由 簇章引理
\ref{varieties-lemma-finite-extension-geometrically-reduced}，存在有限纯不可分
域扩张 $K'/K$，使 $(X_{K'})_{red}$ 在 $K'$ 上几何约化。选择元素
$x_1, \ldots, x_n \in K'$，使它们在 $K$ 上生成 $K'$，并且每个 $x_i$
的某个 $p$ 次幂属于 $A/\mathfrak p$。令 $A' \subset K'$ 为由
$x_1, \ldots, x_n$ 生成的 $K'$ 的有限 $A$-子代数。注意，$A'$ 是分式域为
$K'$ 的整环。由 代数章引理 \ref{algebra-lemma-p-ring-map}，
$A \to A'$ 在谱上诱导泛同胚。令 $Y' = \Spec(A')$，并令
$X' = (Y' \times_Y X)_{red}$。由引理
\ref{more-morphisms-lemma-reduction-generic-fibre}，$X' \to Y'$ 的泛纤维为
$(X_{K'})_{red}$，按构造它几何约化。注意，$X' \to X_V$ 是有限泛同胚，
因为它是约化态射 $X' \to Y' \times_Y X$（见 态射章引理
\ref{morphisms-lemma-reduction-universal-homeomorphism}）与 $g$ 的基变换的
复合。此时除可能的 (7) 外，本引理所有性质都成立。缩小 $Y'$，从而缩小
$V$，即可得到 (7)；见 态射章命题
\ref{morphisms-proposition-generic-flatness}。
\end{proof}

\begin{lemma}
\label{more-morphisms-lemma-make-components-generic-fibre-geometrically-irreducible}
设 $f : X \to Y$ 是概形态射。假设 $Y$ 不可约且 $f$ 有限型。
存在图
$$
\xymatrix{
X' \ar[d]_{f'} \ar[r]_{g'} & X_V \ar[r] \ar[d] & X \ar[d]^f \\
Y' \ar[r]^g & V \ar[r] & Y
}
$$
，其中
\begin{enumerate}
\item $V$ 是 $Y$ 的非空开集；
\item $X_V = V \times_Y X$；
\item $g : Y' \to V$ 是满有限平展态射；
\item $X' = Y' \times_Y X = Y' \times_V X_V$；
\item $g'$ 是满有限平展态射；
\item $Y'$ 是不可约仿射概形；
\item $f'$ 的泛纤维的所有不可约分支都几何不可约。
\end{enumerate}
\end{lemma}

\begin{proof}
设 $V = \Spec(A)$ 是 $Y$ 的非空仿射开集。由假设，$A$ 的幂零根
$\mathfrak p = \sqrt{(0)}$ 是素理想。令 $K = \kappa(\mathfrak p)$。
由 簇章引理
\ref{varieties-lemma-finite-extension-geometrically-irreducible-components}，
存在有限可分域扩张 $K'/K$，使 $X_{K'}$ 的所有不可约分支都在 $K'$
上几何不可约。选择在 $K$ 上生成 $K'$ 的元素 $\alpha \in K'$；见
域章引理 \ref{fields-lemma-primitive-element}。设
$P(T) \in K[T]$ 是 $\alpha$ 在 $K$ 上的极小多项式。对某个
$f \in A$、$f \not \in \mathfrak p$，用 $f \alpha$ 替换 $\alpha$ 后，
可以假设存在首一多项式
$T^d + a_1T^{d - 1} + \ldots + a_d \in A[T]$，它在映射
$A[T] \to K[T]$ 下映到 $P(T) \in K[T]$。
% P05-ERR-0038: frozen source quotients A[T] by P, previously defined in K[T].
令 $A' = A[T]/(P)$。则 $A \to A'$ 是有限自由环映射，并存在唯一的
位于 $\mathfrak p$ 上方的素理想 $\mathfrak q$，使
$K = \kappa(\mathfrak p) \subset \kappa(\mathfrak q) = K'$ 有限可分，
且 $\mathfrak pA'_{\mathfrak q}$ 是 $A'_{\mathfrak q}$ 的极大理想。
因此 $g : Y' = \Spec(A') \to V = \Spec(A)$ 在 $\mathfrak q$ 处平展；
见 代数章引理 \ref{algebra-lemma-characterize-etale}。这意味着存在开集
$W \subset \Spec(A')$，使 $g|_W : W \to \Spec(A)$ 平展。由于 $g$ 有限，
且 $\mathfrak q$ 是位于 $\mathfrak p$ 上方的唯一点，
$Z = g(Y' \setminus W)$ 是 $V$ 的不含 $\mathfrak p$ 的闭子集。
因此，用不与 $Z$ 相交的 $V$ 的一个主仿射开集替换 $V$ 后，$g$
成为有限平展态射。
\end{proof}







\section{相对伴随点}
\label{more-morphisms-section-assassin}

\begin{lemma}
\label{more-morphisms-lemma-relative-assassin-in-neighbourhood}
设 $f : X \to S$ 是概形态射，$\mathcal{F}$ 是拟凝聚
$\mathcal{O}_X$-模。设 $\xi \in \text{Ass}_{X/S}(\mathcal{F})$，并令
$Z = \overline{\{\xi\}} \subset X$。若 $f$ 局部有限型，且
$\mathcal{F}$ 是有限型 $\mathcal{O}_X$-模，则存在非空开集
$V \subset Z$，使对每个 $s \in f(V)$，$V_s$ 的泛点都属于
$\text{Ass}_{X/S}(\mathcal{F})$。
\end{lemma}

\begin{proof}
可以用 $f(\xi)$ 的一个仿射开邻域替换 $S$，并用 $\xi$ 的一个仿射开
邻域替换 $X$。故可假设 $S = \Spec(A)$、$X = \Spec(B)$，而 $f$ 由
有限型环映射 $A \to B$ 给出；见 态射章引理
\ref{morphisms-lemma-locally-finite-type-characterize}。此外，对某个有限
$B$-模 $M$，可写 $\mathcal{F} = \widetilde{M}$；见 性质章引理
\ref{properties-lemma-finite-type-module}。设 $\mathfrak q \subset B$
是 $\xi$ 对应的素理想，并设 $\mathfrak p \subset A$ 是 $A$ 中相应的
素理想。由假设，
$\mathfrak q \in \text{Ass}_B(M \otimes_A \kappa(\mathfrak p))$；见
代数章说明 \ref{algebra-remark-bourbaki} 与 除子章引理
\ref{divisors-lemma-associated-affine-open}。在此记号下，
$Z = V(\mathfrak q) \subset \Spec(B)$。特别地，
$f(Z) \subset V(\mathfrak p)$。因此，显然只需在分别用
$A/\mathfrak pA$、$B/\mathfrak pB$ 与 $M/\mathfrak pM$ 替换
$A$、$B$ 与 $M$ 后证明本引理。换言之，可假设 $A$ 是分式域为 $K$
的整环，且 $\mathfrak q \subset B$ 是 $M \otimes_A K$ 的伴随素理想。

\medskip\noindent
此时可以使用泛平坦性。确切地说，由 代数章引理
\ref{algebra-lemma-generic-flatness}，存在非零 $g \in A$，使 $M_g$
作为 $A_g$-模平坦。用 $A_g$ 替换 $A$ 后，可以假设 $M$ 作为 $A$-模
平坦。

\medskip\noindent
此时，由 代数章引理 \ref{algebra-lemma-post-bourbaki}，
$\mathfrak q$ 也是 $M$ 的伴随素理想。因此得到单射 $B$-模映射
$B/\mathfrak q \to M$。令 $Q$ 为其余核，得到有限 $B$-模的短正合列
$$
0 \to B/\mathfrak q \to M \to Q \to 0
$$
。再次应用泛平坦性 代数章引理
\ref{algebra-lemma-generic-flatness}，这次应用于 $B$-模 $Q$，可以假设
$Q$ 是平坦 $A$-模。特别地，可以假设上述短正合列泛单射；见 Algebra,
引理 \ref{algebra-lemma-flat-tor-zero}。在此情形下，对 $A$ 的任意素理想
$\mathfrak p'$，有
$(B/\mathfrak q) \otimes_A \kappa(\mathfrak p')
\subset M \otimes_A \kappa(\mathfrak p')$
。本引理随即得出，因为由 除子章引理
\ref{divisors-lemma-minimal-support-in-ass}，
$(B/\mathfrak q) \otimes_A \kappa(\mathfrak p')$ 的支集的极小素理想
$\mathfrak q'$ 是
$(B/\mathfrak q) \otimes_A \kappa(\mathfrak p')$ 的伴随素理想。
\end{proof}

\begin{lemma}
\label{more-morphisms-lemma-bad-case}
设 $f : X \to Y$ 是概形态射，$\mathcal{F}$ 是拟凝聚
$\mathcal{O}_X$-模，$U \subset X$ 是开子概形。假设
\begin{enumerate}
\item $f$ 有限型；
\item $\mathcal{F}$ 有限型；
\item $Y$ 不可约，泛点为 $\eta$；
\item $\text{Ass}_{X_\eta}(\mathcal{F}_\eta)$ 不包含于 $U_\eta$。
\end{enumerate}
则存在非空开子概形 $V \subset Y$，使对所有 $y \in V$，集合
$\text{Ass}_{X_y}(\mathcal{F}_y)$ 不包含于 $U_y$。
\end{lemma}

\begin{proof}
设 $Z \subset X$ 是 $\mathcal{F}$ 的概形论支集；见 态射章定义
\ref{morphisms-definition-scheme-theoretic-support}。则 $Z_\eta$ 是
$\mathcal{F}_\eta$ 的概形论支集（态射章引理
\ref{morphisms-lemma-flat-pullback-support}）。因此，由 除子章引理
\ref{divisors-lemma-minimal-support-in-ass}，$Z_\eta$ 各不可约分支的泛点
包含于 $\text{Ass}_{X_\eta}(\mathcal{F}_\eta)$。
% P05-ERR-0039: frozen "bad case" proof then asserts disjointness/empty T_eta;
% this conflicts with hypothesis (4) and appears to prove the opposite case.
由此可见 $Z_\eta \cap U_\eta = \emptyset$。因此
$T = Z \setminus U$ 是 $Z$ 的闭子集，且 $T_\eta = \emptyset$。若赋予
$T$ 诱导的约化概形结构，则 $T \to Y$ 是有限型态射。由引理
\ref{more-morphisms-lemma-empty-generic-fibre}，存在非空开集 $V \subset Y$，使
$T_V = \emptyset$。这个 $V$ 即为所求。
\end{proof}

\begin{lemma}
\label{more-morphisms-lemma-good-case}
设 $f : X \to Y$ 是概形态射，$\mathcal{F}$ 是拟凝聚
$\mathcal{O}_X$-模，$U \subset X$ 是开子概形。假设
\begin{enumerate}
\item $f$ 有限型；
\item $\mathcal{F}$ 有限型；
\item $Y$ 不可约，泛点为 $\eta$；
\item $\text{Ass}_{X_\eta}(\mathcal{F}_\eta) \subset U_\eta$。
\end{enumerate}
则存在非空开子概形 $V \subset Y$，使对所有 $y \in V$，都有
$\text{Ass}_{X_y}(\mathcal{F}_y) \subset U_y$。
\end{lemma}

\begin{proof}
（这个证明与引理 \ref{more-morphisms-lemma-scheme-theoretically-dense-generic-fibre}
的证明相同。我们建议读者先阅读那个证明。）由于陈述只涉及纤维，显然
可以用其约化替换 $Y$。因此可假设 $Y$ 是整概形；见 Properties,
引理 \ref{properties-lemma-characterize-integral}。还可假设
$Y = \Spec(A)$ 仿射。于是 $A$ 是分式域为 $K$ 的整环。

\medskip\noindent
由于 $f$ 有限型，$X$ 拟紧。对某些仿射开集 $X_i$，写成
$X = X_1 \cup \ldots \cup X_n$，并令
$\mathcal{F}_i = \mathcal{F}|_{X_i}$。由假设，
$U_i = X_i \cap U$ 的泛纤维包含
$\text{Ass}_{X_{i, \eta}}(\mathcal{F}_{i, \eta})$。因此只需对各三元组
$(X_i, \mathcal{F}_i, U_i)$ 证明结果；换言之，可假设 $X$ 仿射。

\medskip\noindent
写 $X = \Spec(B)$。设 $N$ 是满足 $\mathcal{F} = \widetilde{N}$ 的有限
$B$-模。注意，$B_K$ 是有限型 $K$-代数，故诺特。因此 $U_\eta$ 拟紧。
于是可找到有限多个 $g_1, \ldots, g_m \in B$，使 $D(g_j) \subset U$，
且 $U_\eta = D(g_1)_\eta \cup \ldots \cup D(g_m)_\eta$。由于
$\text{Ass}_{X_\eta}(\mathcal{F}_\eta) \subset U_\eta$，映射
$N_K \to \bigoplus_j (N_K)_{g_j}$ 是单射。由 代数章引理
\ref{algebra-lemma-when-injective-covering}，这等价于
$N_K \to \bigoplus\nolimits_{j = 1, \ldots, m} N_K$，
$n \mapsto (g_1n, \ldots, g_mn)$ 是单射。令 $I$ 与 $M$ 分别为此映射
在 $A$ 上的核与余核；也就是说，有正合列
$$
0 \to I \to N \xrightarrow{(g_1, \ldots, g_m)}
\bigoplus\nolimits_{j = 1, \ldots, m} N \to M \to 0
$$
对某个非零 $h$，用 $A_h$ 替换 $A$ 后，可以假设 $B$ 是平坦有限表示
$A$-代数，并且 $M$ 与 $N$ 都在 $A$ 上平坦；见 代数章引理
\ref{algebra-lemma-generic-flatness}。$N$ 在 $A$ 上平坦，蕴含 $N$
作为 $A$-模无挠；见 代数进阶章引理
\ref{more-algebra-lemma-flat-torsion-free}。故 $N \subset N_K$。
按构造 $I_K = 0$，这蕴含 $I = 0$（因为
$I \subset N \subset N_K$ 是 $I_K$ 的子集）。于是现在有短正合列
$$
0 \to N \xrightarrow{(g_1, \ldots, g_m)}
\bigoplus\nolimits_{j = 1, \ldots, m} N \to M \to 0
$$
，其中 $M$ 在 $A$ 上平坦。因此，对每个到域 $\kappa$ 的同态
$A \to \kappa$，得到短正合列
$$
0 \to N \otimes_A \kappa \xrightarrow{(g_1 \otimes 1, \ldots, g_m \otimes 1)}
\bigoplus\nolimits_{j = 1, \ldots, m} N \otimes_A \kappa \to
M \otimes_A \kappa \to 0
$$
；见 代数章引理 \ref{algebra-lemma-flat-tor-zero}。反向应用上述论证，
这意味着 $\bigcup D(g_j \otimes 1)$ 包含
$\text{Ass}_{B \otimes_A \kappa}(N \otimes_A \kappa)$。由于
$\bigcup D(g_j \otimes 1) = \bigcup D(g_j)_\kappa \subset U_\kappa$，
可知 $U_\kappa$ 包含
$\text{Ass}_{X \otimes \kappa}(\mathcal{F} \otimes \kappa)$，
这正是所要证明的。
\end{proof}

\begin{lemma}
\label{more-morphisms-lemma-base-change-assassin-in-U}
设 $f : X \to S$ 是局部有限型态射，$\mathcal{F}$ 是有限型拟凝聚
$\mathcal{O}_X$-模，而 $U \subset X$ 是开子概形。设 $g : S' \to S$
是概形态射，设 $f' : X' = X_{S'} \to S'$ 是 $f$ 的基变换，设
$g' : X' \to X$ 是投影，并令 $\mathcal{F}' = (g')^*\mathcal{F}$、
$U' = (g')^{-1}(U)$。最后，设 $s' \in S'$ 的像为 $s = g(s')$。
此时
$$
\text{Ass}_{X_s}(\mathcal{F}_s) \subset U_s
\Leftrightarrow
\text{Ass}_{X'_{s'}}(\mathcal{F}'_{s'}) \subset U'_{s'}.
$$
\end{lemma}

\begin{proof}
这立即由 除子章引理
\ref{divisors-lemma-base-change-relative-assassin} 得出。另见 Divisors,
说明 \ref{divisors-remark-base-change-relative-assassin}。
\end{proof}

\begin{lemma}
\label{more-morphisms-lemma-relative-assassin-constructible}
设 $f : X \to Y$ 是有限表示态射，$\mathcal{F}$ 是有限表示拟凝聚
$\mathcal{O}_X$-模。设 $U \subset X$ 是开子概形，且 $U \to Y$ 拟紧。
则集合
$$
E = \{y \in Y \mid \text{Ass}_{X_y}(\mathcal{F}_y) \subset U_y\}
$$
在 $Y$ 中局部可构造。
\end{lemma}

\begin{proof}
设 $y \in Y$。须证明存在 $y$ 在 $Y$ 中的开邻域 $V$，使
$E \cap V$ 在 $V$ 中可构造。因此可假设 $Y$ 仿射。写
$Y = \Spec(A)$，并把 $A = \colim A_i$ 写成有限型
$\mathbf{Z}$-代数的定向极限。由 极限章引理
\ref{limits-lemma-descend-finite-presentation}，可以找到 $i$ 以及有限表示态射
$f_i : X_i \to \Spec(A_i)$，其到 $Y$ 的基变换恢复 $f$。必要时增大 $i$，
可以假设存在有限表示拟凝聚 $\mathcal{O}_{X_i}$-模 $\mathcal{F}_i$，
其到 $X$ 的拉回同构于 $\mathcal{F}$；见 极限章引理
\ref{limits-lemma-descend-modules-finite-presentation}。必要时再增大一次 $i$，
可以假设存在开子概形 $U_i \subset X_i$，其在 $X$ 中的逆像为 $U$；
见 极限章引理 \ref{limits-lemma-descend-opens}。由引理
\ref{more-morphisms-lemma-base-change-assassin-in-U}，只需对 $f_i$ 证明本引理。因此问题
化到 $Y$ 是诺特环的谱的情形。

\medskip\noindent
在 $Y$ 为诺特概形时，我们用 拓扑章引理
\ref{topology-lemma-characterize-constructible-Noetherian} 的判别法证明 $E$
可构造。为此，设 $Z \subset Y$ 是不可约闭子概形。须证明
$E \cap Z$ 要么包含非空开子集，要么在 $Z$ 中不稠密。把引理
\ref{more-morphisms-lemma-bad-case} 与 \ref{more-morphisms-lemma-good-case} 应用于 $Z$ 上的基变换
$(X, \mathcal{F}, U) \times_Y Z$，即得结论。
\end{proof}











\section{约化纤维}
\label{more-morphisms-section-reduced}

\begin{lemma}
\label{more-morphisms-lemma-nonreduced-in-neighbourhood}
设 $f : X \to Y$ 是概形态射。假设 $Y$ 不可约，泛点为 $\eta$，且
$f$ 有限型。若 $X_\eta$ 非约化，则存在非空开集 $V \subset Y$，
使对所有 $y \in V$，纤维 $X_y$ 都非约化。
\end{lemma}

\begin{proof}
设 $Y' \subset Y$ 是 $Y$ 的约化，并设 $X' \to Y'$ 是 $f$ 的基变换。
注意，$Y' \to Y$ 在点上诱导双射，而 $X' \to X$ 识别各纤维。
% P05-ERR-0040: frozen source says assume Y' reduced after replacement, likely Y.
故可假设 $Y'$ 约化，也就是整；见 性质章引理
\ref{properties-lemma-characterize-integral}。还可用仿射开集替换 $Y$，
故可假设 $Y = \Spec(A)$，其中 $A$ 是整环。以 $K$ 表示 $A$ 的分式域。
选择仿射开集 $\Spec(B) = U \subset X$ 以及非零幂零截面
$h_\eta \in \Gamma(U_\eta, \mathcal{O}_{U_\eta}) = B_K$。缩小 $Y$ 后，
可以假设 $h$ 来自 $h \in \Gamma(U, \mathcal{O}_U) = B$。再稍微缩小
$Y$ 后，可以假设 $h$ 幂零。令
$I = \{b \in B \mid hb = 0\}$ 为 $h$ 的零化子。则 $C = B/I$ 是有限型
$A$-代数，其泛纤维 $(B/I)_K$ 非零（因为 $h_\eta \not = 0$）。
把泛平坦性应用于 $A \to C$ 与 $A \to B/hB$；见 代数章引理
\ref{algebra-lemma-generic-flatness}。得到 $g \in A$、$g \not = 0$，使
$C_g$ 作为 $A_g$-模自由，而 $(B/hB)_g$ 作为 $A_g$-模平坦。用
$D(g) \subset Y$ 替换 $Y$。现在有短正合列
$$
0 \to C \to B \to B/hB \to 0.
$$
，其中 $B/hB$ 在 $A$ 上平坦，而 $C$ 是非零自由 $A$-模。由此可知，
对到域的任意同态 $A \to \kappa$，环 $C \otimes_A \kappa$ 非零，且序列
$$
0 \to C \otimes_A \kappa \to B \otimes_A \kappa \to
B/hB \otimes_A \kappa \to 0
$$
正合；见 代数章引理 \ref{algebra-lemma-flat-tor-zero}。注意，由张量积的
右正合性，
$B/hB \otimes_A \kappa = (B \otimes_A \kappa) / h(B \otimes_A \kappa)$。
因此乘以 $h$ 在 $B \otimes_A \kappa$ 上不是零映射。这显然意味着，
对任意点 $y \in Y$，元素 $h$ 在 $U_y$ 上限制为非零元素，故 $X_y$
非约化。
\end{proof}

\begin{lemma}
\label{more-morphisms-lemma-base-change-fibres-geometrically-reduced}
设 $f : X \to Y$ 是概形态射。设 $g : Y' \to Y$ 是任意态射，并以
$f' : X' \to Y'$ 表示 $f$ 的基变换。则
\begin{align*}
\{y' \in Y' \mid X'_{y'}\text{ is geometrically reduced}\} \\
= g^{-1}(\{y \in Y \mid X_y\text{ is geometrically reduced}\}).
\end{align*}
\end{lemma}

\begin{proof}
这归结为如下陈述：若 $y' \in Y'$ 的像为 $y \in Y$，则纤维
$X'_{y'} = X_y \times_y y'$ 在 $\kappa(y')$ 上几何约化，当且仅当
$X_y$ 在 $\kappa(y)$ 上几何约化。这由 簇章引理
\ref{varieties-lemma-geometrically-reduced-upstairs} 得出。
\end{proof}

\begin{lemma}
\label{more-morphisms-lemma-not-geometrically-reduced-in-neighbourhood}
设 $f : X \to Y$ 是概形态射。假设 $Y$ 不可约，泛点为 $\eta$，且
$f$ 有限型。若 $X_\eta$ 非几何约化，则存在非空开集 $V \subset Y$，
使对所有 $y \in V$，纤维 $X_y$ 都非几何约化。
\end{lemma}

\begin{proof}
应用引理 \ref{more-morphisms-lemma-make-generic-fibre-geometrically-reduced}，得到
$$
\xymatrix{
X' \ar[d]_{f'} \ar[r]_{g'} & X_V \ar[r] \ar[d] & X \ar[d]^f \\
Y' \ar[r]^g & V \ar[r] & Y
}
$$
并具有该引理中列出的全部性质。设 $\eta'$ 是 $Y'$ 的泛点。考虑态射
$X' \to X_{Y'}$（即约化态射）以及由此得到的泛纤维态射
$X'_{\eta'} \to X_{\eta'}$。由于 $X'_{\eta'}$ 几何约化，而 $X_\eta$
不是几何约化，此态射不可能是同构；见 簇章引理
\ref{varieties-lemma-geometrically-reduced-upstairs}。故 $X_{\eta'}$ 非约化。
因此由引理 \ref{more-morphisms-lemma-nonreduced-in-neighbourhood}，对某个非空开集
$V' \subset Y'$ 的所有点 $y' \in V'$，$X_{Y'} \to Y'$ 的纤维都非约化。
由于 $g : Y' \to V$ 是同胚，引理
\ref{more-morphisms-lemma-base-change-fibres-geometrically-reduced} 说明 $g(V')$
就是所求开集。
\end{proof}

\begin{lemma}
\label{more-morphisms-lemma-geometrically-reduced-generic-fibre}
设 $f : X \to Y$ 是概形态射。假设
\begin{enumerate}
\item $Y$ 不可约，泛点为 $\eta$；
\item $X_\eta$ 几何约化；
\item $f$ 有限型。
\end{enumerate}
则存在非空开子概形 $V \subset Y$，使 $X_V \to V$ 具有几何约化纤维。
\end{lemma}

\begin{proof}
设 $Y' \subset Y$ 是 $Y$ 的约化，并设 $X' \to Y'$ 是 $f$ 的基变换。
注意，$Y' \to Y$ 在点上诱导双射，而 $X' \to X$ 识别各纤维。
% P05-ERR-0041: frozen source again says assume Y' reduced, likely Y.
故可假设 $Y'$ 约化，也就是整；见 性质章引理
\ref{properties-lemma-characterize-integral}。还可用仿射开集替换 $Y$，
故可假设 $Y = \Spec(A)$，其中 $A$ 是整环。以 $K$ 表示 $A$ 的分式域。
稍微缩小 $Y$ 后，还可假设 $X \to Y$ 平坦且有限表示；见 Morphisms,
命题 \ref{morphisms-proposition-generic-flatness}。

\medskip\noindent
由于 $X_\eta$ 几何约化，存在开稠密子集 $V \subset X_\eta$，使
$V \to \Spec(K)$ 光滑；见 簇章引理
\ref{varieties-lemma-geometrically-reduced-dense-smooth-open}。令
$U \subset X$ 为 $f$ 光滑的点所成的集合。由 态射章引理
\ref{morphisms-lemma-set-points-where-fibres-smooth}，有
$V \subset U_\eta$。因此 $U$ 的泛纤维在 $X$ 的泛纤维中稠密。由于
$X_\eta$ 约化，$U_\eta$ 在 $X_\eta$ 中概形论稠密；见 Morphisms,
引理 \ref{morphisms-lemma-reduced-scheme-theoretically-dense}。注意，由于
$U \to Y$ 光滑，$U \to Y$ 的所有纤维都几何约化。因此只需证明：缩小
$Y$ 后，对所有 $y \in Y$，概形 $U_y$ 在 $X_y$ 中概形论稠密；见
态射章引理 \ref{morphisms-lemma-reduced-subscheme-closure}。
这由引理 \ref{more-morphisms-lemma-scheme-theoretically-dense-generic-fibre} 得出。
\end{proof}

\begin{lemma}
\label{more-morphisms-lemma-geometrically-reduced-constructible}
设 $f : X \to Y$ 是拟紧且局部有限表示的态射。则集合
$$
E = \{y \in Y \mid X_y\text{ is geometrically reduced}\}
$$
在 $Y$ 中局部可构造。
\end{lemma}

\begin{proof}
设 $y \in Y$。须证明存在 $y$ 在 $Y$ 中的开邻域 $V$，使
$E \cap V$ 在 $V$ 中可构造。因此可假设 $Y$ 仿射。于是 $X$ 拟紧。
选择有限仿射开覆盖 $X = U_1 \cup \ldots \cup U_n$。各 $U_i \to Y$
在 $y$ 处的纤维构成 $X \to Y$ 在 $y$ 处纤维的仿射开覆盖。因此还可
假设 $X$ 仿射。写 $Y = \Spec(A)$，并把 $A = \colim A_i$ 写成有限型
$\mathbf{Z}$-代数的定向极限。由 极限章引理
\ref{limits-lemma-descend-finite-presentation}，可以找到 $i$ 与有限表示态射
$f_i : X_i \to \Spec(A_i)$，其到 $Y$ 的基变换恢复 $f$。由引理
\ref{more-morphisms-lemma-base-change-fibres-geometrically-reduced}，只需对 $f_i$
证明本引理。因此问题化到 $Y$ 是诺特环的谱的情形。

\medskip\noindent
在 $Y$ 为诺特概形时，我们用 拓扑章引理
\ref{topology-lemma-characterize-constructible-Noetherian} 的判别法证明 $E$
可构造。为此，设 $Z \subset Y$ 是泛点为 $\xi$ 的不可约闭子概形。
须证明 $E \cap Z$ 要么包含非空开子集，要么在 $Z$ 中不稠密。
若 $X_\xi$ 几何约化，则把引理
\ref{more-morphisms-lemma-geometrically-reduced-generic-fibre} 应用于态射
$X_Z \to Z$，可知对某个非空开集 $V \subset Z$，所有纤维 $X_y$
都几何约化。若 $X_\xi$ 非几何约化，则把引理
\ref{more-morphisms-lemma-not-geometrically-reduced-in-neighbourhood} 应用于态射
$X_Z \to Z$，可知对某个非空开集 $V \subset Z$，所有纤维 $X_y$
都非几何约化。结论得证。
\end{proof}

\begin{lemma}
\label{more-morphisms-lemma-proper-flat-over-dvr-reduced-fibre}
设 $f : X \to \Spec(R)$ 是固有态射，其中 $R$ 是离散赋值环。假设
$X$ 的每个不可约分支都支配 $\Spec(R)$（例如，当 $f$ 平坦时），
并假设特殊纤维约化。则 $X$ 与泛纤维 $X_\eta$ 都约化。
\end{lemma}

\begin{proof}
（平坦态射满足关于不可约分支的假设，例如可由 除子章引理
\ref{divisors-lemma-bourbaki} 或平坦环映射的下降性质得出。）假设特殊纤维
$X_s$ 约化。设 $x \in X$ 是任意点。我们证明 $\mathcal{O}_{X, x}$
约化；这将证明 $X$ 与 $X_\eta$ 约化。选择特化 $x \leadsto x'$，其中
$x'$ 位于特殊纤维中；由于固有态射是闭的，这样的特化存在。考虑局部环
$A = \mathcal{O}_{X, x'}$。由于 $\mathcal{O}_{X, x}$ 是 $A$ 的局部化，
只需证明 $A$ 约化。设 $\pi \in R$ 是一个素元。若 $a \in A$ 非零，
则存在 $n \geq 0$ 与元素 $a' \in A$，使 $a = \pi^n a'$ 且
$a' \not \in \pi A$。这由 Krull 交定理（代数章引理
\ref{algebra-lemma-intersect-powers-ideal-module-zero}）得出。
% P05-ERR-0042: frozen source has typo "If a is nilpotent the a is".
若 $a$ 幂零，则 $a$ 属于 $A$ 的所有极小素理想；而 $a'$ 也如此，
因为由关于 $X$ 的不可约分支的假设，$\pi$ 不属于 $A$ 的任何极小素理想。
故 $a'$ 幂零。但 $a'$ 在约化环
$A/\pi A = \mathcal{O}_{X_s, x'}$ 中的像非零。这产生矛盾，除非 $A$
约化，而这正是所要证明的。
\end{proof}

\begin{lemma}
\label{more-morphisms-lemma-geometrically-reduced-open}
设 $f : X \to Y$ 是有限表示平坦固有态射。则集合\linebreak[4]
$\{y \in Y \mid X_y\text{ is geometrically reduced}\}$ 在 $Y$ 中是开的。
\end{lemma}

\begin{proof}
可假设 $Y$ 仿射。于是 $Y$ 是有限型仿射 $\mathbf{Z}$-概形的余滤过极限。
故可假设 $X \to Y$ 是 $X_0 \to Y_0$ 的基变换，其中 $Y_0$ 是有限型
$\mathbf{Z}$-代数的谱，而 $X_0 \to Y_0$ 平坦且固有。见 极限章引理
\ref{limits-lemma-descend-finite-presentation}、
\ref{limits-lemma-descend-flat-finite-presentation} 与
\ref{limits-lemma-eventually-proper}。由于纤维几何约化的点集的形成与
基变换交换（引理 \ref{more-morphisms-lemma-base-change-fibres-geometrically-reduced}），
可以假设基底诺特。

\medskip\noindent
假设 $Y$ 诺特。由引理
\ref{more-morphisms-lemma-geometrically-reduced-constructible}，该集合可构造。因此只需
证明它在一般化下稳定（拓扑章引理
\ref{topology-lemma-characterize-closed-Noetherian}）。由 性质章引理
\ref{properties-lemma-locally-Noetherian-specialization-dvr}，问题化到如下
情形：$Y = \Spec(R)$，$R$ 是离散赋值环，且闭纤维 $X_y$ 几何约化。
须证泛纤维 $X_\eta$ 几何约化。

\medskip\noindent
若不然，则存在 $R$ 的分式域的有限扩张 $L$，使 $X_L$ 非约化；见
簇章引理 \ref{varieties-lemma-geometrically-reduced}。存在支配 $R$
的离散赋值环 $R' \subset L$，其分式域为 $L$；见 代数章引理
\ref{algebra-lemma-integral-closure-Dedekind}。用 $R'$ 替换 $R$ 后，
问题归结为引理 \ref{more-morphisms-lemma-proper-flat-over-dvr-reduced-fibre}。
\end{proof}







\section{纤维的不可约分支}
\label{more-morphisms-section-irreducible}

\begin{lemma}
\label{more-morphisms-lemma-irreducible-components-in-neighbourhood}
设 $f : X \to Y$ 是概形态射。假设 $Y$ 不可约，泛点为 $\eta$，且
$f$ 有限型。若 $X_\eta$ 有 $n$ 个不可约分支，则存在非空开集
$V \subset Y$，使对所有 $y \in V$，纤维 $X_y$ 至少有 $n$ 个不可约
分支。
\end{lemma}

\begin{proof}
由于问题纯粹是拓扑性的，可以用 $X$ 与 $Y$ 的约化分别替换它们。
特别地，这蕴含 $Y$ 是整概形；见 性质章引理
\ref{properties-lemma-characterize-integral}。设
$X_\eta = X_{1, \eta} \cup \ldots \cup X_{n, \eta}$ 是 $X_\eta$
的不可约分支分解。设 $X_i \subset X$ 是泛纤维为 $X_{i, \eta}$
的约化闭子概形。注意，$Z_{i, j} = X_i \cap X_j$ 是 $X_i$ 的闭子集，
其泛纤维 $Z_{i, j, \eta}$ 在 $X_{i, \eta}$ 中处处不稠密。因此缩小
$Y$ 后，可以假设对每个 $y \in Y$，$Z_{i, j, y}$ 在 $X_{i, y}$ 中
处处不稠密；见引理 \ref{more-morphisms-lemma-nowhere-dense-generic-fibre}。进一步缩小
$Y$ 后，可以假设对 $y \in Y$ 有 $X_y = \bigcup X_{i, y}$；见引理
\ref{more-morphisms-lemma-cover-generic-fibre-neighbourhood}。此外，缩小 $Y$ 后，可以
假设每个 $X_i \to Y$ 都平坦且有限表示；见 态射章命题
\ref{morphisms-proposition-generic-flatness}。态射 $X_i \to Y$ 是开的；见
态射章引理 \ref{morphisms-lemma-fppf-open}。因此存在 $\eta$ 的开邻域
$V$，对每个 $i$ 该邻域都包含于 $f(X_i)$。对每个 $y \in V$，概形
$X_{i, y}$ 是 $X_y$ 的非空闭子集，有 $X_y = \bigcup X_{i, y}$，且交
$Z_{i, j, y} = X_{i, y} \cap X_{j, y}$ 在 $X_{i, y}$ 中不稠密。
这显然蕴含 $X_y$ 至少有 $n$ 个不可约分支。
\end{proof}

\begin{lemma}
\label{more-morphisms-lemma-base-change-fibres-geometrically-irreducible}
设 $f : X \to Y$ 是概形态射。设 $g : Y' \to Y$ 是任意态射，并以
$f' : X' \to Y'$ 表示 $f$ 的基变换。则
\begin{align*}
\{y' \in Y' \mid X'_{y'}\text{ is geometrically irreducible}\} \\
= g^{-1}(\{y \in Y \mid X_y\text{ is geometrically irreducible}\}).
\end{align*}
\end{lemma}

\begin{proof}
这归结为如下陈述：若 $y' \in Y'$ 的像为 $y \in Y$，则纤维
$X'_{y'} = X_y \times_y y'$ 在 $\kappa(y')$ 上几何不可约，当且仅当
$X_y$ 在 $\kappa(y)$ 上几何不可约。这由 簇章引理
\ref{varieties-lemma-geometrically-irreducible-check-after-extension} 得出。
\end{proof}

\begin{lemma}
\label{more-morphisms-lemma-base-change-fibres-nr-geometrically-irreducible-components}
设 $f : X \to Y$ 是概形态射。设
$$
n_{X/Y} : Y \to \{0, 1, 2, 3, \ldots, \infty\}
$$
是如下函数：它把 $y \in Y$ 映到 $(X_y)_K$ 的不可约分支数，其中 $K$
是 $\kappa(y)$ 的可分闭扩张。此函数良定义；若 $g : Y' \to Y$ 是态射，
则
$$
n_{X'/Y'} = n_{X/Y} \circ g
$$
，其中 $X' \to Y'$ 是 $f$ 的基变换。
\end{lemma}

\begin{proof}
假设 $y' \in Y'$ 的像为 $y \in Y$。假设
$K \supset \kappa(y)$ 与 $K' \supset \kappa(y')$ 是可分闭扩张。
则可以选择域的交换图
$$
\xymatrix{
K \ar[r] & K'' & K' \ar[l] \\
\kappa(y) \ar[u] \ar[rr] & & \kappa(y') \ar[u]
}
$$
结论随即得出，因为概形态射
$$
\xymatrix{
(X'_{y'})_{K'} &
(X'_{y'})_{K''} = (X_y)_{K''} \ar[l] \ar[r] &
(X_y)_K
}
$$
在不可约分支之间诱导双射；见 簇章引理
\ref{varieties-lemma-separably-closed-field-irreducible-components}。
\end{proof}

\begin{lemma}
\label{more-morphisms-lemma-irreducible-polynomial-over-domain}
设 $A$ 是分式域为 $K$ 的整环，$P \in A[x_1, \ldots, x_n]$。
以 $\overline{K}$ 表示 $K$ 的代数闭包。假设 $P$ 在
$\overline{K}[x_1, \ldots, x_n]$ 中不可约。则存在 $f \in A$，使对所有
到域的同态 $\varphi : A_f \to \kappa$，多项式
$P^\varphi \in \kappa[x_1, \ldots, x_n]$ 都不可约。
\end{lemma}

\begin{proof}
存在 $A$ 上的 $A[x_1, \ldots, x_n]$ 的自同构 $\Psi$，使
$\Psi(P) = ax_n^d +$ 关于 $x_n$ 的低次项，其中 $a \not = 0$；见
代数章引理 \ref{algebra-lemma-helper-polynomial}。可以用 $\Psi(P)$
替换 $P$，并用 $A_a$ 替换 $A$。因此可假设 $P$ 是关于 $x_n$ 的次数
为 $d > 0$ 的首一多项式。对 $i = 1, \ldots, n - 1$，令 $d_i$
为 $P$ 关于 $x_i$ 的次数。注意，这蕴含：对每个到域 $\kappa$ 的同态
$\varphi : A \to \kappa$，$P^\varphi$ 是关于 $x_n$ 的次数为 $d$
的首一多项式，并且关于 $x_i$ 的次数为 $\leq d_i$。因此若
$P^\varphi$ 可约，则可以写成
$$
P^\varphi = Q_1 Q_2
$$
% P05-ERR-0044: frozen proof allows e_1 or e_2 = 0, hence trivial factors.
，其中 $Q_1, Q_2$ 关于 $x_n$ 分别为次数 $e_1, e_2 \geq 0$ 的首一
多项式，满足 $e_1 + e_2 = d$，并且对 $i = 1, \ldots, n - 1$，
关于 $x_i$ 的次数均为 $\leq d_i$。换言之，可以写成
\begin{equation}
\label{more-morphisms-equation-factors}
Q_j = x_n^{e_j} + \sum\nolimits_{0 \leq l < e_j}
\left( \sum\nolimits_{L \in \mathcal{L}} a_{j, l, L} x^L \right) x_n^l
\end{equation}
，其中求和遍历形如 $L = (l_1, \ldots, l_{n - 1})$ 且满足
$0 \leq l_i \leq d_i$ 的多重指标 $L$ 所成的集合 $\mathcal{L}$。
对任意满足 $e_1 + e_2 = d$ 的 $e_1, e_2 \geq 0$，考虑 $A$-代数
$$
B_{e_1, e_2} =
A[\{a_{1, l, L}\}_{0 \leq l < e_1, L \in \mathcal{L}},
\{a_{2, l, L}\}_{0 \leq l < e_2, L \in \mathcal{L}}]/(\text{relations})
$$
% P05-ERR-0045: frozen source has subject agreement "relations is".
，其中 $(\text{relations})$ 是由多项式
$$
P - Q_1Q_2 \in
A[\{a_{1, l, L}\}_{0 \leq l < e_1, L \in \mathcal{L}},
\{a_{2, l, L}\}_{0 \leq l < e_2, L \in \mathcal{L}}][x_1, \ldots, x_n]
$$
的系数生成的理想，这里 $Q_1$ 与 $Q_2$ 如
(\ref{more-morphisms-equation-factors}) 中定义。
% P05-ERR-0046: frozen source contains the editorial fragment "OK, and".
好，并且 $P$ 在 $\overline{K}$ 上不可约这一假设蕴含不存在任何
$A$-代数同态 $B_{e_1, e_2} \to \overline{K}$。由 Hilbert 零点定理，
见 代数章定理 \ref{algebra-theorem-nullstellensatz}，这意味着
$B_{e_1, e_2} \otimes_A K = 0$。由于 $B_{e_1, e_2}$ 是有限生成
$A$-代数，这表示可以找到 $f_{e_1, e_2} \in A$，使
$(B_{e_1, e_2})_{f_{e_1, e_2}} = 0$。按构造，这意味着：若
$\varphi : A_{f_{e_1, e_2}} \to \kappa$ 是到域的同态，则
$P^\varphi$ 不存在分解 $P^\varphi = Q_1 Q_2$，其中 $Q_1$ 关于
$x_n$ 的次数为 $e_1$，而 $Q_2$ 关于 $x_n$ 的次数为 $e_2$。
因此取 $f = \prod_{e_1, e_2 \geq 0, e_1 + e_2 = d} f_{e_1, e_2}$ 即得结论。
\end{proof}

\begin{lemma}
\label{more-morphisms-lemma-geom-irreducible-generic-fibre}
设 $f : X \to Y$ 是概形态射。假设
\begin{enumerate}
\item $Y$ 不可约，泛点为 $\eta$；
\item $X_\eta$ 几何不可约；
\item $f$ 有限型。
\end{enumerate}
则存在非空开子概形 $V \subset Y$，使 $X_V \to V$ 具有几何不可约纤维。
\end{lemma}

\begin{proof}[引理 \ref{more-morphisms-lemma-geom-irreducible-generic-fibre} 的第一种证明]
我们给出本引理的两种证明。二者本质等价；第二种更自足，但稍长。
选择图
$$
\xymatrix{
X' \ar[d]_{f'} \ar[r]_{g'} & X_V \ar[r] \ar[d] & X \ar[d]^f \\
Y' \ar[r]^g & V \ar[r] & Y
}
$$
，如引理 \ref{more-morphisms-lemma-make-generic-fibre-geometrically-reduced} 中所述。
注意，$f'$ 的泛纤维是 $f$ 的泛纤维的约化（见引理
\ref{more-morphisms-lemma-reduction-generic-fibre}），故几何不可约。假设本引理对态射
$f'$ 成立。则缩小 $V$ 后，$f'$ 的所有纤维都几何不可约。由于
$X' = (Y' \times_V X_V)_{red}$，这蕴含 $Y' \times_V X_V$ 的所有纤维
都几何不可约。因此由引理
\ref{more-morphisms-lemma-base-change-fibres-geometrically-irreducible}，
$X_V \to V$ 的所有纤维都几何不可约，结论得证。这样便可假设泛纤维
既几何约化又几何不可约，并可假设 $Y = \Spec(A)$，其中 $A$ 是整环。

\medskip\noindent
设 $x \in X_\eta$ 是泛点。由于 $X_\eta$ 几何不可约且约化，
$L = \kappa(x)$ 是 $K = \kappa(\eta)$ 的有限生成扩张，并且几何约化、
几何不可约；见 簇章引理
\ref{varieties-lemma-geometrically-reduced-at-point} 与
\ref{varieties-lemma-geometrically-irreducible-function-field}。特别地，
域扩张 $L/K$ 可分；见 代数章引理
\ref{algebra-lemma-characterize-separable-field-extensions}。因此可以找到
$x_1, \ldots, x_{r + 1} \in L$，使它们在 $K$ 上生成 $L$，并且
$x_1, \ldots, x_r$ 是 $L$ 在 $K$ 上的一组超越基；见 代数章引理
\ref{algebra-lemma-generating-finitely-generated-separable-field-extensions}。
设 $P \in K(x_1, \ldots, x_r)[T]$ 是 $x_{r + 1}$ 的极小多项式。
清除分母后，可以假设 $P$ 的系数属于 $A[x_1, \ldots, x_r]$。
注意，由于 $L$ 在 $K$ 上几何约化且几何不可约，多项式 $P$ 在
$\overline{K}[x_1, \ldots, x_r, T]$ 中不可约，其中 $\overline{K}$
是 $K$ 的代数闭包。记
$$
B' = A[x_1, \ldots, x_{r + 1}]/(P(x_{r + 1}))
$$
，并令 $X' = \Spec(B')$。按构造，$B'$ 的分式域作为 $K$-扩张同构于
$L = \kappa(x)$。因此存在开集 $U \subset X$、开集 $U' \subset X'$
以及 $Y$-同构 $U \to U'$；见 态射章引理
\ref{morphisms-lemma-common-open}。图示如下：
$$
\xymatrix{
X \ar[rd] &
U \ar[l] \ar@{=}[r] \ar[d] &
U' \ar[r] \ar[d] &
X' \ar[ld] \ar@{=}[r] & \Spec(B') \\
& Y \ar@{=}[r] & Y &
}
$$
注意，$U_\eta \subset X_\eta$ 与 $U'_\eta \subset X'_\eta$ 都是稠密开集。
因此，应用引理 \ref{more-morphisms-lemma-nowhere-dense-generic-fibre} 缩小 $Y$ 后，
对所有 $y \in Y$，$U_y$ 在 $X_y$ 中稠密，且 $U'_y$ 在 $X'_y$
中稠密。因此只需对 $X' \to Y$ 证明本引理，而这正是引理
\ref{more-morphisms-lemma-irreducible-polynomial-over-domain} 的内容。
\end{proof}

\begin{proof}[引理 \ref{more-morphisms-lemma-geom-irreducible-generic-fibre} 的第二种证明]
设 $Y' \subset Y$ 是 $Y$ 的约化，并设 $X' \to X$ 是 $X$ 的约化。
注意，$X' \to X  \to Y$ 经由 $Y'$ 分解；见 概形章引理
\ref{schemes-lemma-map-into-reduction}。由 态射章引理
\ref{morphisms-lemma-reduction-universal-homeomorphism}，$Y' \to Y$
与 $X' \to X$ 都是泛同胚，故只需对 $X' \to Y'$ 证明本引理。因此可
假设 $X$ 与 $Y$ 都约化。特别地，$Y$ 是整概形；见 性质章引理
\ref{properties-lemma-characterize-integral}。于是由 态射章命题
\ref{morphisms-proposition-generic-flatness}，存在非空仿射开集
$V \subset Y$，使 $X_V \to V$ 平坦且有限表示。用 $V$ 替换 $Y$ 后，
除 (1)、(2)、(3) 外，还可假设 $Y$ 是整仿射概形、$X$ 约化，而 $f$
平坦且有限表示。特别地，$f$ 泛开；见 态射章引理
\ref{morphisms-lemma-fppf-open}。

\medskip\noindent
选择非空仿射开集 $U \subset X$。则 $U \to Y$ 平坦且有限表示，泛纤维
几何不可约。补集 $X_\eta \setminus U_\eta$ 处处不稠密。因此缩小 $Y$
后，可以假设对所有 $y \in Y$，$U_y \subset X_y$ 是开稠密子集；见
引理 \ref{more-morphisms-lemma-nowhere-dense-generic-fibre}。故可以用 $U$ 替换 $X$，
把问题化到如下情形：$Y$ 是整仿射概形，$X$ 是在 $Y$ 上平坦且有限表示
的约化仿射概形，其泛纤维 $X_\eta$ 几何不可约。

\medskip\noindent
写 $X = \Spec(B)$、$Y = \Spec(A)$。则 $A$ 是整环，$B$ 约化，
$A \to B$ 是有限表示平坦映射，而 $B_K$ 在 $A$ 的分式域 $K$ 上几何
不可约。特别地，$B_K$ 是整环。设 $L$ 为 $B_K$ 的分式域。注意，由于
$B$ 是有限表示 $A$-代数，$L$ 是 $K$ 的有限生成域扩张。设 $K'/K$
是有限纯不可分扩张，使 $(L \otimes_K K')_{red}$ 是可分生成域扩张；
见 代数章引理 \ref{algebra-lemma-make-separable}。选择
$x_1, \ldots, x_n \in K'$，使它们在 $K$ 上生成域扩张 $K'$，并且对某个
素数幂 $q_i$ 有 $x_i^{q_i} \in A$（略去 $x_i$ 存在性的证明）。
设 $A'$ 为由 $x_1, \ldots, x_n$ 生成的 $K'$ 的 $A$-子代数。则 $A'$
是有限 $A$-子代数 $A' \subset K'$，其分式域为 $K'$。注意，
$\Spec(A') \to \Spec(A)$ 是泛同胚；见 代数章引理
\ref{algebra-lemma-p-ring-map}。因此只需在基变换到 $\Spec(A')$ 后证明
结论。我们将用 $A'$ 替换 $A$，并用 $(B \otimes_A A')_{red}$ 替换
$B$，从而达到 $L$ 是 $K$ 的可分生成域扩张的情形。当然，
$(B \otimes_A A')_{red}$ 可能不再平坦，或不再在 $A'$ 上有限表示；
但对合适的 $f \in A'$，用 $A'_f$ 替换 $A'$ 即可补救这一点；见
代数章引理 \ref{algebra-lemma-generic-flatness}。

\medskip\noindent
此时已知：$A$ 是整环，$B$ 约化，$A \to B$ 平坦且有限表示，$B_K$
是整环，其分式域 $L$ 是 $A$ 的分式域 $K$ 的可分生成域扩张。由
代数章引理
\ref{algebra-lemma-generating-finitely-generated-separable-field-extensions}，
可以写 $L = K(x_1, \ldots, x_{r + 1})$，其中
$x_1, \ldots, x_r$ 在 $K$ 上代数无关，而 $x_{r + 1}$ 在
$K(x_1, \ldots, x_r)$ 上可分。清除分母后，可以假设 $x_{r + 1}$
在 $K(x_1, \ldots, x_r)$ 上的极小多项式
$P \in K(x_1, \ldots, x_r)[T]$ 的系数属于
$A[x_1, \ldots, x_r]$。注意，由于 $L/K$ 可分，且 $L$ 在 $K$
上几何不可约，多项式 $P$ 在 $K$ 的代数闭包 $\overline{K}$ 上不可约。
记
$$
B' = A[x_1, \ldots, x_{r + 1}]/(P(x_{r + 1})).
$$
按构造，$B$ 与 $B'$ 的分式域作为 $K$-扩张同构。因此，对合适的
$h \in B$ 与 $h' \in B'$，存在 $A$-代数同构
$B_h \cong B'_{h'}$；见 态射章引理
\ref{morphisms-lemma-common-open}。换言之，$X$ 与
$X' = \Spec(B')$ 有一个共同的仿射开集 $U$。图示如下：
$$
\xymatrix{
X = \Spec(B) \ar[rd] &
U \ar[l] \ar[r] \ar[d] &
\Spec(B') = X' \ar[ld] \\
& Y = \Spec(A) &
}
$$
再次缩小 $Y$（把引理 \ref{more-morphisms-lemma-nowhere-dense-generic-fibre} 分别应用于
$X$ 中的 $Z = X \setminus U$ 与 $X'$ 中的 $Z' = X' \setminus U$）后，
对所有 $y \in Y$，$U_y$ 在 $X_y$ 中稠密，且 $U_y$ 在 $X'_y$
中稠密。因此只需对 $X' \to Y$ 证明本引理，而这正是引理
\ref{more-morphisms-lemma-irreducible-polynomial-over-domain} 的内容。
\end{proof}

\begin{lemma}
\label{more-morphisms-lemma-nr-geom-irreducible-components-good}
设 $f : X \to Y$ 是概形态射。设 $n_{X/Y}$ 是引理
\ref{more-morphisms-lemma-base-change-fibres-nr-geometrically-irreducible-components}
中引入的 $Y$ 上的函数，它计数 $f$ 的纤维的几何不可约分支数。
假设 $f$ 有限型。设 $y \in Y$ 是一点。则存在非空开集
$V \subset \overline{\{y\}}$，使 $n_{X/Y}|_V$ 为常数。
\end{lemma}

\begin{proof}
在 $\overline{\{y\}}$ 上赋予诱导的约化概形结构，并记为 $Z$。
设 $f_Z : X_Z \to Z$ 是 $f$ 的基变换。显然，只需对 $f_Z$ 与 $Z$
的泛点证明本引理。因此可假设 $Y$ 是整概形；见 性质章引理
\ref{properties-lemma-characterize-integral}。此时的目标是构造非空开集
$V \subset Y$，使 $n_{X/Y}|_V$ 为常数。

\medskip\noindent
把引理
\ref{more-morphisms-lemma-make-components-generic-fibre-geometrically-irreducible}
应用于 $f : X \to Y$，得到 $g : Y' \to V \subset Y$。由于
$g : Y' \to V$ 是满有限平展态射，特别地是开态射（见 态射章引理
\ref{morphisms-lemma-etale-open}），只需证明存在开集 $V' \subset Y'$，
使 $n_{X'/Y'}|_{V'}$ 为常数；见引理
\ref{more-morphisms-lemma-base-change-fibres-nr-geometrically-irreducible-components}。
因此可假设泛纤维 $X_\eta$ 的所有不可约分支都在 $\kappa(\eta)$ 上
几何不可约。

\medskip\noindent
此时，假设
$X_\eta = X_{1, \eta} \bigcup \ldots \bigcup X_{n, \eta}$
是泛纤维的（几何）不可约分支分解。特别地，
$n_{X/Y}(\eta) = n$。设 $X_i$ 是 $X_{i, \eta}$ 在 $X$ 中的闭包。
缩小 $Y$ 后，可以假设 $X = \bigcup X_i$；见引理
\ref{more-morphisms-lemma-cover-generic-fibre-neighbourhood}。进一步缩小 $Y$ 后，
$f$ 的每个纤维至少有 $n$ 个不可约分支；见引理
\ref{more-morphisms-lemma-irreducible-components-in-neighbourhood}。故对所有 $y \in Y$，
$n_{X/Y}(y) \geq n$。再次缩小 $Y$ 后，对每个 $i$ 与所有 $y \in Y$，
$X_{i, y}$ 几何不可约；见引理
\ref{more-morphisms-lemma-geom-irreducible-generic-fibre}。由于
$X_y = \bigcup X_{i, y}$，这说明 $n_{X/Y}(y) \leq n$，证明完成。
\end{proof}

\begin{lemma}
\label{more-morphisms-lemma-nr-geom-irreducible-components-constructible}
设 $f : X \to Y$ 是概形态射。设 $n_{X/Y}$ 是引理
\ref{more-morphisms-lemma-base-change-fibres-nr-geometrically-irreducible-components}
中引入的 $Y$ 上的函数，它计数 $f$ 的纤维的几何不可约分支数。
假设 $f$ 有限表示。则 $n_{X/Y}$ 的水平集
$$
E_n = \{y \in Y \mid n_{X/Y}(y) = n\}
$$
在 $Y$ 中局部可构造。
\end{lemma}

\begin{proof}
固定 $n$，并设 $y \in Y$。须证明存在 $y$ 在 $Y$ 中的开邻域 $V$，
使 $E_n \cap V$ 在 $V$ 中可构造。因此可假设 $Y$ 仿射。写
$Y = \Spec(A)$，并把 $A = \colim A_i$ 写成有限型
$\mathbf{Z}$-代数的定向极限。由 极限章引理
\ref{limits-lemma-descend-finite-presentation}，可以找到 $i$ 与有限表示态射
$f_i : X_i \to \Spec(A_i)$，其到 $Y$ 的基变换恢复 $f$。由引理
\ref{more-morphisms-lemma-base-change-fibres-nr-geometrically-irreducible-components}，
只需对 $f_i$ 证明本引理。因此问题化到 $Y$ 是诺特环的谱的情形。

\medskip\noindent
在 $Y$ 为诺特概形时，我们用 拓扑章引理
\ref{topology-lemma-characterize-constructible-Noetherian} 的判别法证明
$E_n$ 可构造。为此，设 $Z \subset Y$ 是不可约闭子概形。须证明
$E_n \cap Z$ 要么包含非空开子集，要么在 $Z$ 中不稠密。设
$\xi \in Z$ 是泛点。引理
\ref{more-morphisms-lemma-nr-geom-irreducible-components-good} 表明，$n_{X/Y}$ 在 $Z$
中 $\xi$ 的一个邻域上为常数。这显然蕴含所需结论。
\end{proof}













\section{纤维的连通分支}
\label{more-morphisms-section-connected}

\begin{lemma}
\label{more-morphisms-lemma-connected-components-in-neighbourhood}
设 $f : X \to Y$ 是概形态射。假设 $Y$ 不可约、泛点为 $\eta$，且
$f$ 是有限型的。若 $X_\eta$ 有 $n$ 个连通分支，则存在非空开子集
$V \subset Y$，使得对所有 $y \in V$，纤维 $X_y$ 至少有 $n$ 个连通分支。
\end{lemma}

\begin{proof}
由于问题纯粹是拓扑性的，可将 $X$ 和 $Y$ 换成其约化。特别地，这说明
$Y$ 是整概形，见 性质章引理
\ref{properties-lemma-characterize-integral}。设
$X_\eta = X_{1, \eta} \cup \ldots \cup X_{n, \eta}$ 是 $X_\eta$
的连通分支分解。令 $X_i \subset X$ 为泛纤维是 $X_{i, \eta}$ 的约化闭子概形。
注意，$Z_{i, j} = X_i \cap X_j$ 是 $X$ 的闭子集，且其泛纤维
$Z_{i, j, \eta}$ 为空。因此缩小 $Y$ 后，可假设
$Z_{i, j} = \emptyset$，见引理 \ref{more-morphisms-lemma-empty-generic-fibre}。进一步缩小
$Y$ 后，可假设对 $y \in Y$ 有 $X_y = \bigcup X_{i, y}$，见引理
\ref{more-morphisms-lemma-cover-generic-fibre-neighbourhood}。此外，缩小 $Y$ 后，可假设每个
$X_i \to Y$ 都平坦且有限表现，见 态射章命题
\ref{morphisms-proposition-generic-flatness}。态射 $X_i \to Y$ 是开态射，见
态射章引理 \ref{morphisms-lemma-fppf-open}。因此存在 $\eta$ 的开邻域
$V$，它对每个 $i$ 都包含于 $f(X_i)$。对每个 $y \in V$，概形
$X_{i, y}$ 是 $X_y$ 的非空闭子集，且 $X_y = \bigcup X_{i, y}$，而交
$Z_{i, j, y} = X_{i, y} \cap X_{j, y}$ 为空！显然，这说明 $X_y$ 至少有
$n$ 个连通分支。
\end{proof}

\begin{lemma}
\label{more-morphisms-lemma-base-change-fibres-geometrically-connected}
设 $f : X \to Y$ 是概形态射。任取态射 $g : Y' \to Y$，并将 $f$ 的基变换记为
$f' : X' \to Y'$。则
\begin{align*}
\{y' \in Y' \mid X'_{y'}\text{ 几何连通}\} \\
= g^{-1}(\{y \in Y \mid X_y\text{ 几何连通}\}).
\end{align*}
\end{lemma}

\begin{proof}
这归结为如下陈述：若 $y' \in Y'$ 的像为 $y \in Y$，则纤维
$X'_{y'} = X_y \times_y y'$ 在 $\kappa(y')$ 上几何连通，当且仅当 $X_y$
在 $\kappa(y)$ 上几何连通。这由 簇章引理
\ref{varieties-lemma-geometrically-connected-check-after-extension} 得出。
\end{proof}

\begin{lemma}
\label{more-morphisms-lemma-base-change-fibres-nr-geometrically-connected-components}
设 $f : X \to Y$ 是概形态射。令
$$
n_{X/Y} : Y \to \{0, 1, 2, 3, \ldots, \infty\}
$$
为如下函数：对 $y \in Y$，它取值为 $(X_y)_K$ 的连通分支数，其中 $K$
是 $\kappa(y)$ 的一个可分闭扩张。此定义良定；若 $g : Y' \to Y$ 是态射，则
$$
n_{X'/Y'} = n_{X/Y} \circ g
$$
其中 $X' \to Y'$ 是 $f$ 的基变换。
\end{lemma}

\begin{proof}
设 $y' \in Y'$ 的像为 $y \in Y$。设 $K \supset \kappa(y)$ 和
$K' \supset \kappa(y')$ 是可分闭扩张。于是可以选取域的交换图
$$
\xymatrix{
K \ar[r] & K'' & K' \ar[l] \\
\kappa(y) \ar[u] \ar[rr] & & \kappa(y') \ar[u]
}
$$
下列概形态射诱导连通分支之间的双射，因而得到结论：
$$
\xymatrix{
(X'_{y'})_{K'} &
(X'_{y'})_{K''} = (X_y)_{K''} \ar[l] \ar[r] &
(X_y)_K
}
$$
见 簇章引理
\ref{varieties-lemma-separably-closed-field-connected-components}。
\end{proof}

\begin{lemma}
\label{more-morphisms-lemma-geometrically-connected-generic-fibre}
设 $f : X \to Y$ 是概形态射。假设
\begin{enumerate}
\item $Y$ 不可约，泛点为 $\eta$；
\item $X_\eta$ 几何连通；并且
\item $f$ 是有限型的。
\end{enumerate}
则存在非空开子概形 $V \subset Y$，使得 $X_V \to V$ 的纤维几何连通。
\end{lemma}

\begin{proof}
按引理 \ref{more-morphisms-lemma-make-components-generic-fibre-geometrically-irreducible}
选取图
$$
\xymatrix{
X' \ar[d]_{f'} \ar[r]_{g'} & X_V \ar[r] \ar[d] & X \ar[d]^f \\
Y' \ar[r]^g & V \ar[r] & Y
}
$$
。注意 $f'$ 的泛纤维几何连通（例如由引理
\ref{more-morphisms-lemma-base-change-fibres-nr-geometrically-connected-components}）。假设本引理对态射
$f'$ 成立。这意味着存在非空开子集 $W \subset Y'$，使得 $X' \to Y'$ 在
$W$ 上的每个纤维都几何连通。由 态射章引理
\ref{morphisms-lemma-etale-open}，$g$ 是开态射，故 $f$ 在非空开子集
$V = g(W)$ 的各点处的纤维全都几何连通，见引理
\ref{more-morphisms-lemma-base-change-fibres-nr-geometrically-connected-components}。由此可见，可以假设泛纤维
$X_\eta$ 的不可约分支均几何不可约。

\medskip\noindent
令 $Y'$ 为 $Y$ 的约化，并置 $X' = Y' \times_Y X$。只需对态射
$X' \to Y'$ 证明本引理（例如再次使用引理
\ref{more-morphisms-lemma-base-change-fibres-nr-geometrically-connected-components}）。由于 $X' \to Y'$
的泛纤维与 $X \to Y$ 的泛纤维相同，可假设 $Y$ 不可约且约化（即为整概形，见
性质章引理 \ref{properties-lemma-characterize-integral}），并且泛纤维
$X_\eta$ 的不可约分支均几何不可约。

\medskip\noindent
现在设
$X_\eta = X_{1, \eta} \bigcup \ldots \bigcup X_{n, \eta}$
是泛纤维的（几何）不可约分支分解。令 $X_i$ 为 $X_{i, \eta}$ 在 $X$ 中的闭包。
缩小 $Y$ 后，由引理 \ref{more-morphisms-lemma-cover-generic-fibre-neighbourhood} 可假设
$X = \bigcup X_i$。令 $Z_{i, j} = X_i \cap X_j$。作分解
$$
\{1, \ldots, n\} \times \{1, \ldots, n\} = I \amalg J
$$
，其中当 $Z_{i, j, \eta} = \emptyset$ 时 $(i, j) \in I$，而当
$Z_{i, j, \eta} \not = \emptyset$ 时 $(i, j) \in J$。缩小 $Y$ 后，由引理
\ref{more-morphisms-lemma-empty-generic-fibre} 可假设对所有 $(i, j) \in I$ 都有
$Z_{i, j} = \emptyset$。再缩小 $Y$ 后，由引理
\ref{more-morphisms-lemma-geom-irreducible-generic-fibre}，对每个 $i$ 以及所有 $y \in Y$，
$X_{i, y}$ 都几何不可约。进一步缩小 $Y$，由 态射章命题
\ref{morphisms-proposition-generic-flatness}，可使每个 $Z_{i, j} \to Y$ 对所有
$(i, j) \in J$ 都平坦且有限表现。这意味着 $f(Z_{i, j}) \subset Y$ 是开集，见
态射章引理 \ref{morphisms-lemma-fppf-open}。我们断言
$$
V  = \bigcap\nolimits_{(i, j) \in J} f(Z_{i, j})
$$
满足要求，即对每个 $y \in V$，$X_y$ 都几何连通。事实上，$X_\eta$ 连通说明
由 $J$ 中各对生成的等价关系只有一个等价类。现在若 $y \in V$，且
$K \supset \kappa(y)$ 是可分闭扩张，则 $(X_y)_K$ 的不可约分支就是各纤维
$(X_{i, y})_K$。此外，由构造及 $y \in V$ 可见，$(X_{i, y})_K$ 与
$(X_{j, y})_K$ 相交，当且仅当 $(i, j) \in J$。因此关于等价类的上述说明表明
$(X_y)_K$ 连通，结论成立。
\end{proof}

\begin{lemma}
\label{more-morphisms-lemma-nr-geom-connected-components-good}
设 $f : X \to Y$ 是概形态射。令 $n_{X/Y}$ 为引理
\ref{more-morphisms-lemma-base-change-fibres-nr-geometrically-connected-components} 中引入的、在 $Y$
上计算 $f$ 的纤维之几何连通分支数的函数。假设 $f$ 是有限型的。设
$y \in Y$ 是一点。则存在非空开子集 $V \subset \overline{\{y\}}$，使得
$n_{X/Y}|_V$ 为常数。
\end{lemma}

\begin{proof}
令 $Z$ 为 $\overline{\{y\}}$ 上的诱导约化概形结构。令 $f_Z : X_Z \to Z$
为 $f$ 的基变换。显然，只需对 $f_Z$ 和 $Z$ 的泛点证明本引理。因此可以假设
$Y$ 是整概形，见 性质章引理
\ref{properties-lemma-characterize-integral}。此时目标是构造非空开子集
$V \subset Y$，使得 $n_{X/Y}|_V$ 为常数。

\medskip\noindent
对 $f : X \to Y$ 应用引理
\ref{more-morphisms-lemma-make-components-generic-fibre-geometrically-irreducible}，得到
$g : Y' \to V \subset Y$。由于 $g : Y' \to V$ 是满的有限 \'etale 态射，特别地是开态射
（见 态射章引理 \ref{morphisms-lemma-etale-open}），只需证明存在开子集
$V' \subset Y'$，使得 $n_{X'/Y'}|_{V'}$ 为常数，见
% P05-ERR-CH37-0047: Frozen source cites the irreducible-component base-change lemma here although the argument concerns connected components.
引理 \ref{more-morphisms-lemma-base-change-fibres-nr-geometrically-irreducible-components}.
因此可假设泛纤维 $X_\eta$ 的所有不可约分支在 $\kappa(\eta)$ 上几何不可约。由
簇章引理
\ref{varieties-lemma-irreducible-components-geometrically-irreducible}
可知，$X_\eta$ 的连通分支也几何连通。

\medskip\noindent
现在设
$X_\eta = X_{1, \eta} \bigcup \ldots \bigcup X_{n, \eta}$
是泛纤维的（几何）连通分支分解。特别地，$n_{X/Y}(\eta) = n$。令 $X_i$
为 $X_{i, \eta}$ 在 $X$ 中的闭包。缩小 $Y$ 后，由引理
\ref{more-morphisms-lemma-cover-generic-fibre-neighbourhood} 可假设 $X = \bigcup X_i$。进一步缩小
$Y$ 后，由引理 \ref{more-morphisms-lemma-connected-components-in-neighbourhood} 可知 $f$ 的每个纤维至少有
$n$ 个连通分支。因此对所有 $y \in Y$，均有 $n_{X/Y}(y) \geq n$。再次缩小
$Y$ 后，由引理 \ref{more-morphisms-lemma-geometrically-connected-generic-fibre}，对每个 $i$ 和所有
$y \in Y$，$X_{i, y}$ 都几何连通。由于 $X_y = \bigcup X_{i, y}$，这说明
$n_{X/Y}(y) \leq n$，证明完成。
\end{proof}

\begin{lemma}
\label{more-morphisms-lemma-nr-geom-connected-components-constructible}
设 $f : X \to Y$ 是概形态射。令 $n_{X/Y}$ 为引理
\ref{more-morphisms-lemma-base-change-fibres-nr-geometrically-connected-components} 中引入的、在 $Y$
上计算 $f$ 的纤维之几何连通分支数的函数。假设 $f$ 是有限表现的。则
$n_{X/Y}$ 的水平集
$$
E_n = \{y \in Y \mid n_{X/Y}(y) = n\}
$$
在 $Y$ 中局部可构造。
\end{lemma}

\begin{proof}
固定 $n$。设 $y \in Y$。须证明在 $Y$ 中存在 $y$ 的开邻域 $V$，使得
$E_n \cap V$ 在 $V$ 中可构造。因此可假设 $Y$ 仿射。写成
$Y = \Spec(A)$，并将 $A = \colim A_i$ 写成有限型 $\mathbf{Z}$-代数的有向极限。
由 极限章引理 \ref{limits-lemma-descend-finite-presentation}，可以找到 $i$ 及有限表现态射
$f_i : X_i \to \Spec(A_i)$，使其到 $Y$ 的基变换恢复 $f$。由引理
\ref{more-morphisms-lemma-base-change-fibres-nr-geometrically-connected-components}，只需对 $f_i$
证明本引理。因此问题化到 $Y$ 是诺特环的谱的情形。

\medskip\noindent
在 $Y$ 为诺特概形时，我们用 拓扑章引理
\ref{topology-lemma-characterize-constructible-Noetherian} 的判别法证明 $E_n$ 可构造。
为此，设 $Z \subset Y$ 是不可约闭子概形。须证明 $E_n \cap Z$ 要么包含非空开子集，
要么在 $Z$ 中不稠密。令 $\xi \in Z$ 为泛点。引理
\ref{more-morphisms-lemma-nr-geom-connected-components-good} 表明，$n_{X/Y}$ 在 $Z$ 中 $\xi$
的一个邻域上为常数。这显然蕴含所需结论。
\end{proof}

\begin{lemma}
\label{more-morphisms-lemma-connected-flat-over-dvr}
\begin{slogan}
不连通概形的平坦退化要么不连通，要么非约化。
\end{slogan}
设 $f : X \to S$ 是概形态射。假设
\begin{enumerate}
\item $S$ 是离散赋值环的谱；
\item $f$ 平坦；
\item $X$ 连通；
\item 闭纤维 $X_s$ 约化。
\end{enumerate}
则泛纤维 $X_\eta$ 连通。
\end{lemma}

\begin{proof}
写成 $S = \Spec(R)$，并令 $\pi \in R$ 为一致化元。反设 $X_\eta$ 不连通以导出矛盾。
这意味着存在非平凡幂等元 $e \in \Gamma(X_\eta, \mathcal{O}_{X_\eta})$。任取
$X$ 中的仿射开集 $U = \Spec(A)$。由于 $A$ 在 $R$ 上平坦，$\pi$ 是 $A$
上的非零因子，例如见 代数进阶章引理
\ref{more-algebra-lemma-flat-torsion-free}。于是 $e|_{U_\eta}$ 对应于一个元素
$e \in A[1/\pi]$。取 $z \in A$，使得 $e = z/\pi^n$，其中 $n \geq 0$ 最小。
注意 $z^2 = \pi^nz$。若 $n > 0$，这说明 $z \bmod \pi A$ 是幂零元。按假设，
$A/\pi A$ 约化，故 $n$ 的最小性蕴含 $n = 0$。因此 $e \in A$！换言之，
$e \in \Gamma(X, \mathcal{O}_X)$。由于 $X$ 连通，$e$ 必为平凡幂等元，矛盾。
\end{proof}







\section{与截面相交的连通分支}
\label{more-morphisms-section-connected-components}

\noindent
本节结果特别适用于群概形 $G \to S$ 及其单位截面 $e : S \to G$。

\begin{situation}
\label{more-morphisms-situation-connected-along-section}
% P05-ERR-CH37-0048: Frozen source has the ungrammatical phrase ``Here f ... be a morphism''.
这里 $f : X \to Y$ 是概形态射，而 $s : Y \to X$ 是 $f$ 的一个截面。对每个
$y \in Y$，以 $X^0_y$ 表示 $X_y$ 中包含 $s(y)$ 的连通分支。最后，置
$X^0 = \bigcup_{y \in Y} X^0_y$。
\end{situation}

\begin{lemma}
\label{more-morphisms-lemma-base-change-connected-along-section}
设 $f : X \to Y$、$s : Y \to X$ 如情形
\ref{more-morphisms-situation-connected-along-section}。若 $g : Y' \to Y$ 是任意态射，考虑基变换图
$$
\xymatrix{
X' \ar[r]_{g'} \ar[d]^{f'} & X \ar[d]_f \\
Y' \ar@/^1pc/[u]^{s'} \ar[r]^g & Y \ar@/_1pc/[u]_s
}
$$
，由此得到 $(X')^0 \subset X'$。则 $(X')^0 = (g')^{-1}(X^0)$。
\end{lemma}

\begin{proof}
设 $y' \in Y'$ 的像为 $y \in Y$。可将 $X^0_y$ 视为 $X_y$ 的闭子概形，例如见
态射章定义
\ref{morphisms-definition-scheme-structure-connected-component}。由于
$s(y) \in X^0_y$，簇章引理
\ref{varieties-lemma-geometrically-connected-if-connected-and-point} 表明，$X_y^0$
是 $\kappa(y)$ 上的几何连通概形。因此 $X_y^0 \times_y y' \to X'_{y'}$
是包含 $s'(y')$ 的连通闭子概形，从而
$X_y^0 \times_y y' \subset (X'_{y'})^0$。反向包含
$X_y^0 \times_y y' \supset (X'_{y'})^0$ 是显然的，因为 $(X'_{y'})^0$
在 $X_y$ 中的像是 $X_y$ 中包含 $s(y)$ 的连通子集。
\end{proof}

\begin{lemma}
\label{more-morphisms-lemma-connected-along-section-good}
设 $f : X \to Y$、$s : Y \to X$ 如情形
\ref{more-morphisms-situation-connected-along-section}。假设 $f$ 是有限型的，并设 $y \in Y$
为一点。则存在非空开子集 $V \subset \overline{\{y\}}$，使得 $X^0$
在基变换 $X_V$ 中的逆像在 $X_V$ 中既开又闭。
\end{lemma}

\begin{proof}
令 $Z \subset Y$ 为 $\overline{\{y\}}$ 上的诱导约化闭子概形结构。令
$f_Z : X_Z \to Z$ 和 $s_Z : Z \to X_Z$ 为 $f$ 和 $s$ 的基变换。由引理
\ref{more-morphisms-lemma-base-change-connected-along-section}，有 $(X_Z)^0 = (X^0)_Z$。因此只需对态射
% P05-ERR-CH37-0049: Frozen source introduces an undefined point x in X_Z although the lemma's distinguished point lies in the base Z.
$X_Z \to Z$ 以及映到 $Z$ 的泛点的点 $x \in X_Z$ 证明本引理。换言之，问题已化到
$Y$ 是泛点为 $\eta$ 的整概形的情形（见 性质章引理
\ref{properties-lemma-characterize-integral}）。目标是证明，缩小 $Y$ 后，子集 $X^0$
成为 $X$ 的既开又闭子集。

\medskip\noindent
注意，概形 $X_\eta$ 在一个域上有限型，因而是诺特概形。因此其连通分支既开又闭。
故可写成 $X_\eta = X_\eta^0 \amalg T_\eta$，其中 $T_\eta$ 是 $X_\eta$
的既开又闭子集。接着，令 $T \subset X$ 为 $T_\eta$ 的闭包，并令
$X^{00} \subset X$ 为 $X_\eta^0$ 的闭包。注意，$T_\eta$（相应地 $X^0_\eta$）
是 $T$（相应地 $X^{00}$）的泛纤维，见引理
\ref{more-morphisms-lemma-cover-generic-fibre-neighbourhood} 之前的讨论。此外，该引理表明，缩小 $Y$
后，可假设 $X = X^{00} \cup T$（在集合论意义下）。注意
$(T \cap X^{00})_\eta = T_\eta \cap X^0_\eta = \emptyset$。因此缩小 $Y$ 后，由引理
\ref{more-morphisms-lemma-empty-generic-fibre} 可假设 $T \cap X^{00} = \emptyset$。特别地，
$X^{00}$ 在 $X$ 中是开集。注意 $X^0_\eta$ 连通且有有理点，即 $s(\eta)$，故它几何连通，见
簇章引理
\ref{varieties-lemma-geometrically-connected-if-connected-and-point}。于是缩小 $Y$ 后，
由引理 \ref{more-morphisms-lemma-geometrically-connected-generic-fibre} 可假设 $X^{00} \to Y$
的所有纤维都几何连通。此时可知各纤维 $X^{00}_y$ 是 $X_y$ 中包含 $s(y)$
的开、闭且连通的子集。因此 $X^0 = X^{00}$，结论成立。
\end{proof}

\begin{lemma}
\label{more-morphisms-lemma-connected-along-section-locally-constructible}
设 $f : X \to Y$、$s : Y \to X$ 如情形
\ref{more-morphisms-situation-connected-along-section}。若 $f$ 是有限表现的，则 $X^0$ 在 $X$
中局部可构造。
\end{lemma}

\begin{proof}
设 $x \in X$。须证明存在 $x$ 的开邻域 $U$，使得 $X^0 \cap U$ 在 $U$
中可构造。这将问题化到 $Y$ 为仿射的情形。写成 $Y = \Spec(A)$，并将
$A = \colim A_i$ 写成有限型 $\mathbf{Z}$-代数的有向极限。由 极限章引理
\ref{limits-lemma-descend-finite-presentation}，可以找到 $i$ 及有限表现态射
$f_i : X_i \to \Spec(A_i)$，并配备截面 $s_i : \Spec(A_i) \to X_i$，使得到
$Y$ 的基变换恢复 $f$ 和截面 $s$。由引理
\ref{more-morphisms-lemma-base-change-connected-along-section}，只需对 $f_i, s_i$ 证明本引理。因此问题化到
$Y$ 是诺特环的谱的情形。

\medskip\noindent
假设 $Y$ 是诺特仿射概形。由于 $f$ 是有限表现的，因而是有限型的，故 $X$
也是诺特概形，见 态射章引理
\ref{morphisms-lemma-finite-type-noetherian}。为在此情形证明本引理，只需证明：对每个不可约闭子集
$Z \subset X$，交 $Z \cap X^0$ 要么包含 $Z$ 的一个非空开集，要么在 $Z$
中不稠密；见 拓扑章引理
\ref{topology-lemma-characterize-constructible-Noetherian}。令 $x \in Z$ 为泛点，并令
$y = f(x)$。由引理 \ref{more-morphisms-lemma-connected-along-section-good}，存在非空开子集
$V \subset \overline{\{y\}}$，使得 $X^0 \cap X_V$ 在 $X_V$ 中既开又闭。由于
$f(Z) \subset \overline{\{y\}}$ 且 $f(x) = y \in V$，可见
$W = f^{-1}(V) \cap Z$ 是 $Z$ 的非空开子集。因此 $X^0 \cap W$ 在 $W$
中既开又闭。由于 $W$ 不可约，$X^0 \cap W$ 要么为空，要么等于 $W$。本引理得证。
\end{proof}

\begin{lemma}
\label{more-morphisms-lemma-connected-along-section-open-neighbourhood}
设 $f : X \to Y$、$s : Y \to X$ 如情形
\ref{more-morphisms-situation-connected-along-section}。设 $y \in Y$ 为一点。假设
\begin{enumerate}
\item $f$ 有限表现且平坦；并且
\item 纤维 $X_y$ 几何约化。
\end{enumerate}
则 $X^0$ 是 $X$ 中 $X^0_y$ 的一个邻域。
\end{lemma}

\begin{proof}
可将 $Y$ 换成 $y$ 的一个仿射开邻域。写成 $Y = \Spec(A)$，并将
$A = \colim A_i$ 写成有限型 $\mathbf{Z}$-代数的有向极限。由 极限章引理
\ref{limits-lemma-descend-finite-presentation}，可以找到 $i$ 及有限表现态射
$f_i : X_i \to \Spec(A_i)$，并配备截面 $s_i : \Spec(A_i) \to X_i$，使得到
$Y$ 的基变换恢复 $f$ 和截面 $s$。必要时增大 $i$ 后，还可假设 $f_i$ 平坦，见
极限章引理 \ref{limits-lemma-descend-flat-finite-presentation}。令 $y_i$ 为 $y$
在 $Y_i$ 中的像。注意 $X_y = (X_{i, y_i}) \times_{y_i} y$。因此
$X_{i, y_i}$ 几何约化，见 簇章引理
\ref{varieties-lemma-geometrically-reduced-upstairs}。由引理
\ref{more-morphisms-lemma-base-change-connected-along-section}，只需对系统 $f_i, s_i, y_i \in Y_i$
证明本引理。因此问题化到 $Y$ 是诺特环的谱的情形。

\medskip\noindent
假设 $Y$ 是诺特环的谱。由于 $f$ 是有限表现的，因而是有限型的，故 $X$
也是诺特概形，见 态射章引理
\ref{morphisms-lemma-finite-type-noetherian}。设 $x \in X^0$ 是位于 $y$ 上方的一点。
由 拓扑章引理
\ref{topology-lemma-constructible-neighbourhood-Noetherian}，只需证明：对任何经过 $x$
的不可约闭子集 $Z \subset X$，交 $X^0 \cap Z$ 在 $Z$ 中稠密。特别地，只需证明
泛点 $x' \in Z$ 属于 $X^0$。由 性质章引理
\ref{properties-lemma-locally-Noetherian-specialization-dvr}，可以找到离散赋值环 $R$
和态射 $\Spec(R) \to X$，它把闭点映到 $x$，把泛点映到 $x'$。通过复合
$\Spec(R) \to X \to Y$，将 $\Spec(R)$ 视为 $Y$ 上的概形。由引理
\ref{more-morphisms-lemma-base-change-connected-along-section}，$(X_R)^0$ 是 $X^0$ 的逆像。按构造，结构态射
$X_R \to \Spec(R)$ 除 $s$ 的基变换 $s_R$ 外还有第二个截面
$t : \Spec(R) \to X_R$，使得 $t(\eta_R)$ 是 $X_R$ 中映到 $x'$ 的点，而
$t(0_R)$ 是 $X_R$ 中映到 $x$ 的点。注意 $t(0_R)$ 属于 $(X_R)^0$，并且
$t(\eta_R) \leadsto t(0_R)$。因此只需证明这蕴含
$t(\eta_R) \in (X_R)^0$。故只需在 $Y$ 是离散赋值环的谱、$y$ 是其闭点的情形证明本引理。

\medskip\noindent
假设 $Y$ 是离散赋值环的谱，而 $y$ 是其闭点。目标是证明 $X^0$ 是 $X^0_y$
的一个邻域。由于 $X_y$ 只有有限多个不可约分支，$X^0_y$ 在 $X_y$ 中既开又闭。
因此补集 $C = X_y \setminus X_y^0$ 在 $X$ 中是闭集。于是
$U = X \setminus C$ 是 $X^0_y$ 的开邻域，并且 $U^0 = X^0$。故只需对态射
$U \to Y$ 证明结论。换言之，可假设 $X_y$ 连通。假设 $X$ 不连通，并设
$X = X_1 \amalg \ldots \amalg X_n$ 为其连通分支分解。则 $s(Y)$ 完全包含在某个
$X_i$ 中。不妨设 $s(Y) \subset X_1$，于是 $X^0 \subset X_1$。故可将 $X$
换成 $X_1$，并假设 $X$ 连通。此时引理 \ref{more-morphisms-lemma-connected-flat-over-dvr}
表明 $X_\eta$ 连通，即 $X^0 = X$，结论成立。
\end{proof}

\begin{lemma}
\label{more-morphisms-lemma-connected-along-section-open}
设 $f : X \to Y$、$s : Y \to X$ 如情形
\ref{more-morphisms-situation-connected-along-section}。假设
\begin{enumerate}
\item $f$ 有限表现且平坦；并且
\item $f$ 的所有纤维均几何约化。
\end{enumerate}
则 $X^0$ 在 $X$ 中是开集。
\end{lemma}

\begin{proof}
这立即由引理 \ref{more-morphisms-lemma-connected-along-section-open-neighbourhood} 得出。
\end{proof}






\section{纤维的维数}
\label{more-morphisms-section-dimension}

\begin{lemma}
\label{more-morphisms-lemma-dimension-in-neighbourhood}
设 $f : X \to Y$ 是概形态射。假设 $Y$ 不可约、泛点为 $\eta$，且 $f$
是有限型的。若 $X_\eta$ 的维数为 $n$，则存在非空开子集 $V \subset Y$，使得对所有
$y \in V$，纤维 $X_y$ 的维数均为 $n$。
\end{lemma}

\begin{proof}
令 $Z = \{x \in X \mid \dim_x(X_{f(x)}) > n \}$。由 态射章引理
\ref{morphisms-lemma-openness-bounded-dimension-fibres}，这是 $X$ 的闭子集。按假设
$Z_\eta = \emptyset$。因此由引理 \ref{more-morphisms-lemma-empty-generic-fibre}，可缩小 $Y$
并假设 $Z = \emptyset$。令
$Z' = \{x \in X \mid \dim_x(X_{f(x)}) > n - 1 \} =
\{x \in X \mid \dim_x(X_{f(x)}) = n\}$。与前面一样，这是 $X$ 的闭子集。
按假设有 $Z'_\eta \not = \emptyset$。因此缩小
$Y$ 后，由引理 \ref{more-morphisms-lemma-nonempty-generic-fibre} 可假设 $Z' \to Y$ 是满射。
结论成立。
\end{proof}

\begin{lemma}
\label{more-morphisms-lemma-base-change-dimension-fibres}
设 $f : X \to Y$ 是有限型态射。令
$$
n_{X/Y} : Y \to \{0, 1, 2, 3, \ldots, \infty\}
$$
为把 $y \in Y$ 映到 $X_y$ 的维数的函数。若 $g : Y' \to Y$ 是态射，则
$$
n_{X'/Y'} = n_{X/Y} \circ g
$$
其中 $X' \to Y'$ 是 $f$ 的基变换。
\end{lemma}

\begin{proof}
这由 态射章引理
\ref{morphisms-lemma-dimension-fibre-after-base-change} 得出。
\end{proof}

\begin{lemma}
\label{more-morphisms-lemma-dimension-fibres-constructible}
设 $f : X \to Y$ 是概形态射。令 $n_{X/Y}$ 为引理
\ref{more-morphisms-lemma-base-change-dimension-fibres} 中引入的、在 $Y$ 上给出 $f$
之纤维维数的函数。假设 $f$ 是有限表现的。则 $n_{X/Y}$ 的水平集
$$
E_n = \{y \in Y \mid n_{X/Y}(y) = n\}
$$
在 $Y$ 中局部可构造。
\end{lemma}

\begin{proof}
固定 $n$。设 $y \in Y$。须证明在 $Y$ 中存在 $y$ 的开邻域 $V$，使得
$E_n \cap V$ 在 $V$ 中可构造。因此可假设 $Y$ 仿射。写成
$Y = \Spec(A)$，并将 $A = \colim A_i$ 写成有限型 $\mathbf{Z}$-代数的有向极限。
由 极限章引理 \ref{limits-lemma-descend-finite-presentation}，可以找到 $i$
及有限表现态射 $f_i : X_i \to \Spec(A_i)$，使其到 $Y$ 的基变换恢复 $f$。
由引理 \ref{more-morphisms-lemma-base-change-dimension-fibres}，只需对 $f_i$ 证明本引理。因此问题化到
$Y$ 是诺特环的谱的情形。

\medskip\noindent
在 $Y$ 为诺特概形时，我们用 拓扑章引理
\ref{topology-lemma-characterize-constructible-Noetherian} 的判别法证明 $E_n$ 可构造。
为此，设 $Z \subset Y$ 是不可约闭子概形。须证明 $E_n \cap Z$ 要么包含非空开子集，
要么在 $Z$ 中不稠密。令 $\xi \in Z$ 为泛点。引理
\ref{more-morphisms-lemma-dimension-in-neighbourhood} 表明，$n_{X/Y}$ 在 $Z$ 中 $\xi$
的一个邻域上为常数。这蕴含所需结论。
\end{proof}

\begin{lemma}
\label{more-morphisms-lemma-dimension-fibres-flat}
设 $f : X \to Y$ 是有限表现的平坦概形态射。令 $n_{X/Y}$ 为引理
\ref{more-morphisms-lemma-base-change-dimension-fibres} 中引入的、在 $Y$ 上给出 $f$
之纤维维数的函数。则 $n_{X/Y}$ 下半连续。
\end{lemma}

\begin{proof}
令 $W \subset X$、$W = \coprod_{d \geq 0} U_d$ 为引理
\ref{more-morphisms-lemma-flat-finite-presentation-CM-open} 和
\ref{more-morphisms-lemma-flat-finite-presentation-CM-pieces} 中构造的开集。设 $y \in Y$ 为一点。若
$n_{X/Y}(y) = \dim(X_y) = n$，则 $y$ 属于 $U_n \to Y$ 的像。由 Morphisms,
引理 \ref{morphisms-lemma-fppf-open}，$f(U_n)$ 在 $Y$ 中是开集。因此存在
$y$ 的开邻域，使得其上 $n_{X/Y}$ 均为 $\geq n$。
\end{proof}

\begin{lemma}
\label{more-morphisms-lemma-dimension-fibres-proper}
设 $f : X \to Y$ 是概形的固有态射。令 $n_{X/Y}$ 为引理
\ref{more-morphisms-lemma-base-change-dimension-fibres} 中引入的、在 $Y$ 上给出 $f$
之纤维维数的函数。则 $n_{X/Y}$ 上半连续。
\end{lemma}

\begin{proof}
令 $Z_d = \{x \in X \mid \dim_x(X_{f(x)}) > d\}$。由 态射章引理
\ref{morphisms-lemma-openness-bounded-dimension-fibres}，$Z_d$ 是 $X$ 的闭子集。
由于 $f$ 固有，$f(Z_d)$ 是闭集。又因
$y \in f(Z_d) \Leftrightarrow n_{X/Y}(y) > d$，本引理成立。
\end{proof}

\begin{lemma}
\label{more-morphisms-lemma-dimension-fibres-proper-flat}
设 $f : X \to Y$ 是有限表现的固有平坦概形态射。令 $n_{X/Y}$ 为引理
\ref{more-morphisms-lemma-base-change-dimension-fibres} 中引入的、在 $Y$ 上给出 $f$
之纤维维数的函数。则 $n_{X/Y}$ 局部常值。
\end{lemma}

\begin{proof}
这立即由引理 \ref{more-morphisms-lemma-dimension-fibres-flat} 和
\ref{more-morphisms-lemma-dimension-fibres-proper} 得出。
\end{proof}

















\section{弱相对 Noether 正规化}
\label{more-morphisms-section-relative-noether-normalization}

\noindent
本节目标是证明引理 \ref{more-morphisms-lemma-weak-relative-noether-normalization}。

\begin{lemma}
\label{more-morphisms-lemma-silly}
设 $R$ 是环。设 $\mathfrak p_1, \ldots, \mathfrak p_r$ 是 $R$ 的素理想，且当
$i \not = j$ 时有 $\mathfrak p_i \not \subset \mathfrak p_j$。设
$k_i \subset \kappa(\mathfrak p_i)$ 是子域，使得扩张
$\kappa(\mathfrak p_i)/k_i$ 不是代数扩张。设 $J \subset R$ 是不包含于任何
$\mathfrak p_i$ 的理想。则存在元素 $x \in J$，使得对 $i = 1, \ldots, r$，
$x$ 在 $\kappa(\mathfrak p_i)$ 中的像在 $k_i$ 上超越。
\end{lemma}

\begin{proof}
理想 $J_i = J \mathfrak p_1 \ldots \hat{\mathfrak p}_i \ldots \mathfrak p_r$
不包含于 $\mathfrak p_i$，见 代数章引理
\ref{algebra-lemma-product-ideals-in-prime}。于是
% P05-ERR-CH37-0050: Frozen source uses the undefined ring B in Frac(B/p_i); the ambient ring is R.
$\kappa(\mathfrak p_i) = \text{Frac}(B/\mathfrak p_i)$ 的每个元素 $\xi$
都可写成 $\xi = a/b$，其中 $a, b \in J_i$ 且
$b \not \in \mathfrak p_i$。选取在 $k_i$ 上超越的 $\xi$，可见 $a$ 或 $b$
在 $\kappa(\mathfrak p_i)$ 中的像在 $k_i$ 上超越。因此，对每个
$i = 1, \ldots, r$，可以找到元素
$x_i \in J_i = J \mathfrak p_1 \ldots \hat{\mathfrak p}_i \ldots \mathfrak p_r$，
使其在 $\kappa(\mathfrak p_i)$ 中的像在 $k_i$ 上超越。于是
$x = x_1 + \ldots + x_r$ 满足要求。
\end{proof}

\begin{lemma}
\label{more-morphisms-lemma-sum-is-ok}
设 $R \to S$ 是有限型环同态。设 $d \geq 0$，并设 $a, b \in S$。假设由把
$x$ 映到 $a$ 的 $R$-代数同态 $R[x] \to S$ 给出的态射
$$
f_a : \Spec(S) \longrightarrow \mathbf{A}^1_R
$$
的纤维维数为 $\leq d$。则存在 $n_0$，使得对 $n \geq n_0$，由把 $x$
映到 $a^n + b$ 的 $R$-代数同态 $R[x] \to S$ 给出的态射
$$
f_{a^n + b} : \Spec(S) \longrightarrow \mathbf{A}^1_R
$$
的纤维维数为 $\leq d$。
\end{lemma}

\begin{proof}
本段将问题化到 $R \to S$ 有限表现的情形。具体地，将
$S = R[A, B, x_1, \ldots, x_n]/J$ 写成这样的形式，其中
% P05-ERR-CH37-0051: Frozen source places J in R[x_1,...,x_n] although the presentation also uses indeterminates A and B.
$J \subset R[x_1, \ldots, x_n]$ 是某个理想，而 $A$ 和 $B$ 在 $S$ 中分别映到
$a$ 和 $b$。于是 $J$ 是其有限生成理想 $J_\lambda \subset J$ 的并。置
$S_\lambda = R[A, B, x_1, \ldots, x_n]/J_\lambda$，并以
$a_\lambda, b_\lambda \in S_\lambda$ 表示 $A$ 和 $B$ 的像。于是对某个 $\lambda$，
态射
$$
f_{a_\lambda} : \Spec(S_\lambda) \longrightarrow \mathbf{A}^1_R
$$
的纤维维数为 $\leq d$，见 极限章引理
\ref{limits-lemma-limit-dimension}。固定这样的 $\lambda$。若能找到适用于
$R \to S_\lambda$、$a_\lambda$、$b_\lambda$ 的 $n_0$，则该 $n_0$ 也适用于
$R \to S$。事实上，$f_{a_\lambda^n + b_\lambda} : \Spec(S_\lambda) \to \mathbf{A}^1_R$
的纤维包含 $f_{a^n + b} : \Spec(S) \to \mathbf{A}^1_R$ 的纤维。这将问题化到下一段所讨论的情形。

\medskip\noindent
假设 $R \to S$ 有限表现。本段将问题化到 $R$ 在 $\mathbf{Z}$ 上有限型的情形。
由 代数章引理 \ref{algebra-lemma-limit-module-finite-presentation}，可以找到有向集
$\Lambda$ 及定义在 $\Lambda$ 上的环同态系统 $R_\lambda \to S_\lambda$，其余极限为 $R \to S$，
使得当 $\mu \geq \lambda$ 时 $S_\mu = S_\lambda \otimes_{R_\lambda} R_\mu$，
并且每个 $R_\lambda$ 和 $S_\lambda$ 都在 $\mathbf{Z}$ 上有限型。选取
$\lambda_0 \in \Lambda$ 以及映到 $a, b \in S$ 的元素
$a_{\lambda_0}, b_{\lambda_0} \in S_{\lambda_0}$。对
$\lambda \geq \lambda_0$，以 $a_\lambda, b_\lambda \in S_\lambda$ 表示
$a_{\lambda_0}, b_{\lambda_0}$ 的像。于是当 $\lambda \geq \lambda_0$ 足够大时，态射
$$
f_{a_\lambda} : 
\Spec(S_\lambda) \longrightarrow \mathbf{A}^1_{R_\lambda}
$$
的纤维维数为 $\leq d$，见 极限章引理
\ref{limits-lemma-descend-dimension-d}。固定这样的 $\lambda$。若能找到适用于
$R_\lambda \to S_\lambda$、$a_\lambda$、$b_\lambda$ 的 $n_0$，则该 $n_0$
也适用于 $R \to S$。事实上，由 态射章引理
\ref{morphisms-lemma-dimension-fibre-after-base-change}，
$f_{a^n + b} : \Spec(S) \to \mathbf{A}^1_R$ 的任一纤维都与
$f_{a_\lambda^n + b_\lambda} : \Spec(S_\lambda) \to \mathbf{A}^1_{R_\lambda}$
的某个纤维具有相同维数。
% P05-ERR-CH37-0052: Frozen source has the duplicated phrase ``the the'' in this reduction sentence.
这将问题化到下一段所讨论的情形。

\medskip\noindent
假设 $R$ 和 $S$ 在 $\mathbf{Z}$ 上有限型。特别地，$R$ 的维数有限，因而可对
$\dim(R)$ 使用归纳法。因此可假设结论对引理中所有如下的 $R' \to S'$、$a$、$b$
成立：$R'$ 和 $S'$ 在 $\mathbf{Z}$ 上有限型，但
$\dim(R') < \dim(R)$。

\medskip\noindent
由于此陈述只涉及 $S$ 的谱的拓扑，可假设 $S$ 约化。令 $S^\nu$ 为 $S$
的正规化。由于 $S$ 优秀，$S \subset S^\nu$ 是有限扩张，见 Algebra,
命题 \ref{algebra-proposition-ubiquity-nagata} 和 态射章引理
\ref{morphisms-lemma-nagata-normalization}。因此 $\Spec(S^\nu) \to \Spec(S)$
是满的有限态射（代数章引理
\ref{algebra-lemma-integral-overring-surjective}）。由此可知，若结论对
$R \to S^\nu$ 以及 $a$、$b$ 在 $S^\nu$ 中的像成立，则它对
$R \to S$、$a$、$b$ 也成立。（略去一个小细节。）这将问题化到下一段所讨论的情形。

\medskip\noindent
假设 $R$ 和 $S$ 在 $\mathbf{Z}$ 上有限型，并且 $S$ 正规。于是对某些正规整环
$S_i$，有 $S = S_1 \times \ldots \times S_r$。若结论对每个 $R \to S_i$
及 $a$、$b$ 在 $S_i$ 中的像成立，则它对 $R \to S$、$a$、$b$
也成立。（略去一个小细节。）这将问题化到下一段所讨论的情形。

\medskip\noindent
假设 $R$ 和 $S$ 在 $\mathbf{Z}$ 上有限型，并且 $S$ 是正规整环。可将 $R$
换成 $R$ 在 $S$ 中的像（这不会增加 $R$ 的维数）。这将问题化到下一段所讨论的情形。

\medskip\noindent
假设 $R \subset S$ 在 $\mathbf{Z}$ 上有限型，且 $S$ 是正规整环。考虑态射
$$
f_a : \Spec(S) \to \mathbf{A}^1_R
$$
假设表明 $f_a$ 的纤维维数为 $\leq d$。因此
$f : \Spec(S) \to \Spec(R)$ 的纤维维数为 $\leq d + 1$（Morphisms,
引理 \ref{morphisms-lemma-dimension-fibre-at-a-point-additive}）。考虑整概形间的态射
$$
\phi : \Spec(S) \to \mathbf{A}^2_R = \Spec(R[x, y])
$$
，它对应于把 $x$ 映到 $a$、把 $y$ 映到 $b$ 的 $R$-代数同态
$R[x, y] \to S$。须考虑两种情形：
\begin{enumerate}
\item $\phi$ 是支配态射；
\item $\phi$ 不是支配态射。
\end{enumerate}
我们断言，在两种情形下，都存在整数 $n_0$ 和非空开子集
$V \subset \Spec(R)$，使得当 $n \geq n_0$ 时，$f_{a^n + b}$ 在各点
$q \in \mathbf{A}^1_V$ 处的纤维维数为 $\leq d$。

\medskip\noindent
在情形 (1) 中证明该断言。有 $f_{a^n + b} = \pi_n \circ \phi$，其中
$$
\pi_n :  \mathbf{A}^2_R \to \mathbf{A}^1_R
$$
是对应于把 $x$ 映到 $x^n + y$ 的 $R$-代数同态 $R[x] \to R[x, y]$
的平坦态射。由于 $\phi$ 支配，存在稠密开子集 $U \subset \Spec(S)$，使得
$\phi|_U : U \to \mathbf{A}^2_R$ 平坦（例如由泛平坦性得出，见 Morphisms,
命题 \ref{morphisms-proposition-generic-flatness}）。于是复合
$$
f_{a^n + b}|_U :
U \xrightarrow{\phi|_U}
\mathbf{A}^2_R
\xrightarrow{\pi_n}
\mathbf{A}^1_R
$$
也平坦。因此由 态射章引理
\ref{morphisms-lemma-dimension-fibre-at-a-point-additive}，此态射的纤维在
$f|_U : U \to \Spec(R)$ 的纤维中至少有余维 $1$。换言之，
$f_{a^n + b}|_U$ 的纤维维数为 $\leq d$。另一方面，由于 $U$ 在
$\Spec(S)$ 中稠密，可以找到非空开子集 $V \subset \Spec(R)$，使得对所有
$p \in V$，$U \cap f^{-1}(p) \subset f^{-1}(p)$ 都稠密（例如见引理
\ref{more-morphisms-lemma-nowhere-dense-generic-fibre}）。因此
$\dim(f^{-1}(p) \setminus U \cap f^{-1}(p)) \leq d$，从而该断言成立（因为
$f_{a^n + b} : \Spec(S) \to \mathbf{A}^1_R$ 的任一纤维都包含于 $f$ 的一个纤维中）。

\medskip\noindent
情形 (2)。此时可在 $R[x, y]$ 中找到非零元 $g = \sum c_{ij} x^i y^j$，使得
$\Im(\phi) \subset V(g)$。事实上，可假设 $g$ 在 $\text{Frac}(R)$ 上不可约。
若 $g \in R[x]$，设其首项系数为 $c$，则在
$V = D(c) \subset \Spec(R)$ 上，$f$ 的纤维维数已经为 $\leq d$（因为 $f_a$
的像包含于 $V(g) \subset \mathbf{A}^1_R$，而后者在 $V$ 上的纤维有限）。因此可假设
$g$ 不属于 $R[x]$。令 $s \geq 1$ 为 $g$ 作为 $y$ 的多项式的次数，并令 $t$
为 $\sum c_{is} x^i$ 作为 $x$ 的多项式的次数。于是 $c_{ts}$ 非零，并且
$$
g(x, -x^n) = (-1)^s c_{ts} x^{t + sn} + l.o.t.
$$
，只要 $n$ 大于 $g$ 作为 $x$ 的多项式的次数（略去一个小细节）。对这样的
$n$，多项式 $g(x, -x^n)$ 是 $x$ 的非零多项式，而且对所有
% P05-ERR-CH37-0053: Frozen source switches c_{ts} to c_{st} in the nonvanishing condition.
$\mathfrak p \subset R$，只要 $c_{st} \not \in \mathfrak p$，它在
$\kappa(\mathfrak p)[x]$ 中的像也是非零多项式。因此当 $V$ 取为
$\Spec(R)$ 的主开集 $D(c_{ts})$ 时，我们的断言成立。

\medskip\noindent
现在可以对 $\dim(R)$ 使用归纳法。具体地，令 $I \subset R$ 为满足
$V(I) = \Spec(R) \setminus V$ 的理想。由于 $R$ 是整环，注意有
$\dim(R/I) < \dim(R)$。令 $n'_0$ 为对 $\dim(R)$ 归纳应用于
$R/I \to S/IS$ 以及 $a$、$b$ 在 $S/IS$ 中的像而得到的整数。于是
$\max(n_0, n'_0)$ 满足要求。
\end{proof}

\begin{lemma}
\label{more-morphisms-lemma-weak-relative-noether-normalization}
设 $R \to S$ 是有限型环同态。令 $d$ 为 $\Spec(S) \to \Spec(R)$
的各纤维维数的最大值。则存在拟有限环同态 $R[t_1, \ldots, t_d] \to S$。
\end{lemma}

\begin{proof}
本段将问题化到 $R \to S$ 有限表现的情形。具体地，将
$S = R[x_1, \ldots, x_n]/J$ 写成这样的形式，其中
$J \subset R[x_1, \ldots, x_n]$ 是某个理想。于是 $J$ 是其有限生成理想
$J_\lambda \subset J$ 的并。置 $S_\lambda = R[x_1, \ldots, x_n]/J_\lambda$。
于是对某个 $\lambda$，$\Spec(S_\lambda) \to \Spec(R)$ 的纤维维数为
$\leq d$，见 极限章引理 \ref{limits-lemma-limit-dimension}。固定这样的
$\lambda$。若能找到拟有限同态 $R[t_1, \ldots, t_d] \to S_\lambda$，则复合
$R[t_1, \ldots, t_d] \to S$ 当然拟有限。这将问题化到下一段所讨论的情形。

\medskip\noindent
假设 $R \to S$ 有限表现。本段将问题化到 $R$ 在 $\mathbf{Z}$ 上有限型的情形。
由 代数章引理 \ref{algebra-lemma-limit-module-finite-presentation}，可以找到有向集
$\Lambda$ 及定义在 $\Lambda$ 上的环同态系统 $R_\lambda \to S_\lambda$，其余极限为 $R \to S$，
使得当 $\mu \geq \lambda$ 时 $S_\mu = S_\lambda \otimes_{R_\lambda} R_\mu$，
并且每个 $R_\lambda$ 和 $S_\lambda$ 都在 $\mathbf{Z}$ 上有限型。于是当
$\lambda$ 足够大时，$\Spec(S_\lambda) \to \Spec(R_\lambda)$ 的纤维维数为
$\leq d$，见 极限章引理 \ref{limits-lemma-descend-dimension-d}。固定这样的
$\lambda$。若能找到拟有限环同态
$R_\lambda[t_1, \ldots, t_d] \to S_\lambda$，则其基变换
$R[t_1, \ldots, t_d] \to S$ 也拟有限（代数章引理
\ref{algebra-lemma-quasi-finite-base-change}）。
% P05-ERR-CH37-0054: Frozen source again has the duplicated phrase ``the the'' in this reduction sentence.
这将问题化到下一段所讨论的情形。

\medskip\noindent
假设 $R$ 和 $S$ 在 $\mathbf{Z}$ 上有限型。若 $d = 0$，则环同态拟有限，证明完成。
假设 $d > 0$。我们将找到元素 $a \in S$，使得 $R$-代数同态
$R[x] \to S$、$x \mapsto a$ 的纤维维数为 $< d$。这将通过归纳完成证明。

\medskip\noindent
我们对 $\dim(R)$ 归纳证明 $a$ 的存在性。

\medskip\noindent
在 $S$ 的极小素理想中，令 $\mathfrak q_1, \ldots, \mathfrak q_r \subset S$
为满足 $\dim_{\mathfrak q_i}(S/R) = d$ 的那些。记号见 代数章定义
\ref{algebra-definition-relative-dimension}。设 $\mathfrak q_i$ 位于素理想
$\mathfrak p_i \subset R$ 上方。由于 $\mathfrak q_i$ 是其纤维的泛点，有
$\text{trdeg}_{\kappa(\mathfrak p_i)}(\kappa(\mathfrak q_i)) = d$；例如可对
$S \otimes_R \kappa(\mathfrak p_i)$ 应用 代数章引理
\ref{algebra-lemma-dimension-at-a-point-finite-type-field}。因此由引理
\ref{more-morphisms-lemma-silly}，可以找到元素 $a \in S$，使得对 $i = 1, \ldots, r$，
$a$ 在 $\kappa(\mathfrak q_i)$ 中的像在 $\kappa(\mathfrak p_i)$ 上超越。考虑态射
$$
f_a : \Spec(S) \longrightarrow \mathbf{A}^1_R
$$
% P05-ERR-CH37-0055: Frozen source has the malformed phrase ``corresponding the ... to mapping''.
，它对应于把 $x$ 映到 $a$ 的 $R$-代数同态 $R[x] \to S$。令
$U \subset \Spec(S)$ 为纤维维数为 $\leq d - 1$ 的开子集，见 Morphisms,
引理 \ref{morphisms-lemma-openness-bounded-dimension-fibres}。按构造，$U$ 包含
$\Spec(S)$ 的所有泛点。特别地，$U$ 包含 $\Spec(S) \to \Spec(R)$
的所有泛纤维的全部泛点，因为这样的点必为 $\Spec(S)$ 的泛点。置
$T = \Spec(S) \setminus U$，并将其视为 $\Spec(S)$ 的约化闭子概形。由刚才所述及假设
$\dim(S/R) \leq d$，$T \to \Spec(R)$ 的泛纤维维数为 $\leq d - 1$。因此多次应用引理
\ref{more-morphisms-lemma-dimension-in-neighbourhood}，以构造 $\Spec(R)$ 各泛点的开邻域，可以找到稠密开子集
$V \subset \Spec(R)$，使得 $T_V \to V$ 的纤维维数为 $\leq d - 1$。由此可知，对
$q \in \mathbf{A}^1_V$，$f_a$ 在 $q$ 上的纤维维数为 $\leq d - 1$（因为已经分别界定了
$f_a|_U$ 与 $f_a|_T$ 的纤维维数）。

\medskip\noindent
由素理想回避，可假设对某个 $f \in R$ 有 $V = D(f)$。于是环同态
$R_f[x] \to S_f$、$x \mapsto a$ 的纤维维数为 $\leq d - 1$。可将 $a$
换成 $fa$，并假设 $a \in (f)$。对 $\dim(R)$ 使用归纳法，可以找到元素
$\overline{b} \in S/fS$，使得
% P05-ERR-CH37-0056: Frozen source writes Spec(R/fR[x]), an ambiguous/ill-typed polynomial-ring quotient; likely Spec((R/fR)[x]).
$\Spec(S/fS) \to \Spec(R/fR[x])$、$x \mapsto \overline{b}$ 的纤维维数为
$\leq d - 1$。令 $b \in S$ 为 $\overline{b}$ 的一个提升。由引理
\ref{more-morphisms-lemma-sum-is-ok}，存在 $n > 0$，使得 $a^n + b$ 对 $R_f \to S_f$
仍然满足要求。另一方面，$a^n + b$ 在 $S/fS$ 中的像是 $\overline{b}$，证明完成。
\end{proof}













\section{Bertini 定理}
\label{more-morphisms-section-bertini}

\noindent
我们继续 簇章第 \ref{varieties-section-bertini} 节 中开始的讨论。本节遵循
\cite{Jou} 给出的精彩论证，证明几何不可约簇的一般超平面截面几何不可约。

\begin{lemma}
\label{more-morphisms-lemma-amazing-bertini-lemma}
\begin{reference}
见 \cite{Jou} 第 71、72 页
\end{reference}
设 $K/k$ 是几何不可约的有限生成域扩张。设 $n \geq 1$。设
$g_1, \ldots, g_n \in K$ 是元素，并且存在 $c_1, \ldots, c_n \in k$，使得元素
$$
x_1, \ldots, x_n, \sum g_ix_i, \sum c_ig_i \in K(x_1, \ldots, x_n)
$$
在 $k$ 上代数无关。则 $K(x_1, \ldots, x_n)$ 在
$k(x_1, \ldots, x_n, \sum g_ix_i)$ 上几何不可约。
\end{lemma}

\begin{proof}
取引理陈述中的 $c_1, \ldots, c_n \in k$。记
$\xi = \sum g_ix_i$、$\delta = \sum c_ig_i$。对 $a \in k$，考虑
$K(x_1, \ldots, x_n)$ 的自同构 $\sigma_a$；它在 $K$ 上为恒等，并由规则
$$
\sigma_a(x_i) = x_i + a c_i
$$
给出。注意 $\sigma_a(\xi) = \xi + a \delta$ 且
$\sigma_a(\delta) = \delta$。考虑域塔
$$
K_0 = k(x_1, \ldots, x_n) \subset
K_1 = K_0(\xi) \subset
K_2 = K_0(\xi, \delta) \subset K(x_1, \ldots, x_n) = \Omega
$$
注意 $\sigma_a(K_0) = K_0$ 且 $\sigma_a(K_2) = K_2$。设
$\theta \in \Omega$ 在 $K_1$ 上可分代数。须证明 $\theta \in K_1$，见
代数章引理
\ref{algebra-lemma-geometrically-irreducible-separable-elements}。

\medskip\noindent
以 $K'_2$ 表示 $K_2$ 在 $\Omega$ 中的可分代数闭包。由于 $K'_2/K_2$
有限（代数章引理 \ref{algebra-lemma-make-geometrically-irreducible}）且可分，
$K'_2$ 与 $K_2$ 之间只有有限多个中间域（域章引理
\ref{fields-lemma-primitive-element}）。若 $k$ 无限\footnote{有限域情形将在证明的最后一段处理。}，
则可在 $k$ 中找到不同的元素 $a_1, a_2$，使得
$$
K_2(\sigma_{a_1}(\theta)) = K_2(\sigma_{a_2}(\theta))
$$
作为 $\Omega$ 的子域相等。记 $\theta_i = \sigma_{a_i}(\theta)$ 以及
$\xi_i = \sigma_{a_i}(\xi) = \xi + a_i \delta$。注意
$$
K_2 = K_0(\xi_1, \xi_2)
$$
，因为 $\xi_i = \xi + a_i \delta$、
$\xi = (a_2 \xi_1 - a_1 \xi_2)/(a_2 - a_1)$，且
$\delta = (\xi_1 - \xi_2)/(a_1 - a_2)$。由于 $K_2/K_0$ 是超越次数为
$2$ 的纯超越扩张，$\xi_1$ 和 $\xi_2$ 在 $K_0$ 上代数无关。由于
$\theta_1$ 在 $K_0(\xi_1)$ 上代数，$\xi_2$ 在
$K_0(\xi_1, \theta_1)$ 上超越。

\medskip\noindent
按假设，$K/k$ 几何不可约。这蕴含 $K(x_1, \ldots, x_n)/K_0$ 几何不可约
（代数章引理
\ref{algebra-lemma-geometrically-irreducible-base-change-transcendental}）。进而，作为子扩张，
$K_0(\xi_1, \theta_1)/K_0$ 几何不可约（代数章引理
\ref{algebra-lemma-subalgebra-geometrically-irreducible}）。由于 $\xi_2$ 在
$K_0(\xi_1, \theta_1)$ 上超越，$K_0(\xi_1, \xi_2, \theta_1)/K_0(\xi_2)$
几何不可约（代数章引理
\ref{algebra-lemma-geometrically-irreducible-add-transcendental}）。按上面对
$a_1, a_2$ 的选取，有
$$
K_0(\xi_1, \xi_2, \theta_1) =
K_2(\sigma_{a_1}(\theta)) =
K_2(\sigma_{a_2}(\theta)) =
K_0(\xi_1, \xi_2, \theta_2)
$$
由于 $\theta_2$ 在 $K_0(\xi_2)$ 上可分代数，再次由 代数章引理
\ref{algebra-lemma-geometrically-irreducible-separable-elements} 得出
$\theta_2 \in K_0(\xi_2)$。对这一关系施以 $\sigma_{a_2}^{-1}$，便得到
% P05-ERR-CH37-0057: Frozen source has the typo ``givens''.
$\theta \in K_0(\xi) = K_1$，正是所需。

\medskip\noindent
这完成了 $k$ 无限时的证明。若 $k$ 有限，则可选取变量 $t$，并考虑扩张
$K(t)/k(t)$；由 代数章引理
\ref{algebra-lemma-geometrically-irreducible-base-change-transcendental}，它几何不可约。
% P05-ERR-CH37-0058: Frozen source has the ungrammatical phrase ``it is still be true''.
仍有
$x_1, \ldots, x_n, \sum g_ix_i, \sum c_ig_i$
在 $K(t, x_1, \ldots, x_n)$ 中于 $k(t)$ 上代数无关，故由上述论证，
$K(t, x_1, \ldots, x_n)$ 在
$k(t, x_1, \ldots, x_n, \sum g_ix_i)$ 上几何不可约。最后再用一次 Algebra,
引理
\ref{algebra-lemma-geometrically-irreducible-base-change-transcendental}
即可完成证明。
\end{proof}

\begin{lemma}
\label{more-morphisms-lemma-algebraically-independent}
设 $A$ 是域 $k$ 上的有限型整环。设 $n \geq 2$。设
$g_1, \ldots, g_n \in A$ 是使得 $V(g_1, g_2)$ 有一个维数为
$\dim(A) - 2$ 的不可约分支的元素。则存在 $c_1, \ldots, c_n \in k$，使得元素
$$
x_1, \ldots, x_n, \sum g_ix_i, \sum c_ig_i \in
\text{Frac}(A)(x_1, \ldots, x_n)
$$
在 $k$ 上代数无关。
\end{lemma}

\begin{proof}
在 $k$ 上代数无关意味着由 $y = \sum g_ix_i$ 和 $z = \sum c_ig_i$ 给出的态射
$$
T = \Spec(A[x_1, \ldots, x_n])
\longrightarrow
\Spec(k[x_1, \ldots, x_n, y, z]) = S
$$
为支配态射。置 $d = \dim(A)$。若 $T \to S$ 不支配，则其像的维数
$< n + 2$，因而每个纤维的每个不可约分支的维数
$> d + n - (n + 2) = d - 2$，见 簇章引理
\ref{varieties-lemma-dimension-fibres-locally-algebraic}。选取闭点
$u \in V(g_1, g_2)$，使其包含在一个维数为 $d - 2$ 的不可约分支中，而不属于
$V(g_1, g_2)$ 的任何其他分支。考虑 $T$ 中位于 $u$ 上方的闭点
$t = (u, 1, 0, \ldots 0)$。置
$(c_1, \ldots, c_n) = (0, 1, 0, \ldots, 0)$。于是 $t$ 映到 $S$ 中的点
$s = (1, 0, \ldots, 0)$。$T \to S$ 在 $s$ 上的纤维由下列元素截出：
$$
x_1 - 1, x_2, \ldots, x_n, \sum x_ig_i, g_2
$$
，因而等价地由下列元素截出：
$$
x_1 - 1, x_2, \ldots, x_n, g_1, g_2
$$
。由对 $g_1, g_2$ 的条件，此子概形有一个维数为 $d - 2$ 的不可约分支。
\end{proof}

\begin{lemma}
\label{more-morphisms-lemma-bertini-irreducible}
\begin{reference}
\cite[Theorem 6.3 part 4)]{Jou}
\end{reference}
在 Varieties, Situation \ref{varieties-situation-family-divisors} 中，假设
\begin{enumerate}
\item $X$ 在 $k$ 上有限型；
\item $X$ 在 $k$ 上几何不可约；
\item 存在 $v_1, v_2, v_3 \in V$ 以及 $H_{v_2} \cap H_{v_3}$ 的一个不可约分支
$Z$，使得 $Z \not \subset H_{v_1}$ 且 $\text{codim}(Z, X) = 2$；并且
\item $\bigcap_{v \in V} H_v$ 的每个不可约分支 $Y$ 都满足
$\text{codim}(Y, X) \geq 2$。
\end{enumerate}
则对一般的 $v \in V \otimes_k k'$，概形 $H_v$ 在 $k'$ 上几何不可约。
\end{lemma}

\begin{proof}
要使假设 (3) 成立，元素 $v_1, v_2, v_3$ 必须在 $V$ 中于 $k$ 上线性无关
（略去一个小细节）。因此可选取 $V$ 的基 $v_1, \ldots, v_r$，以前 $3$
个基向量包含这些元素。回忆 $H_{univ} \subset \mathbf{A}^r_k \times_k X$
是“泛除子”。考虑投影 $q : H_{univ} \to \mathbf{A}^r_k$，其概形论纤维为除子
$H_v$。由引理 \ref{more-morphisms-lemma-geom-irreducible-generic-fibre}，只需证明 $q$
的泛纤维几何不可约。为此，可将 $X$ 换成其约化，因而可假设 $X$
是在 $k$ 上有限型的整概形。

\medskip\noindent
设 $U \subset X$ 是满足 $\mathcal{L}|_U \cong \mathcal{O}_U$ 的非空仿射开集。
写成 $U = \Spec(A)$。以 $f_i \in A$ 表示在同构
$\mathcal{L}|_U \cong \mathcal{O}_U$ 下对应于截面 $\psi(v_i)|_U$ 的元素。于是
$H_{univ} \cap (\mathbf{A}^r_k \times_k U)$ 由下式给出：
$$
H_U = \Spec(A[x_1, \ldots, x_r]/(x_1f_1 + \ldots + x_rf_r))
$$
由所选基可知 $f_1$ 不可能为零，因为否则 $v_1 = 0$，进而
$H_{v_1} = X$，与假设 (3) 矛盾。因此 $\sum x_if_i$ 是
$A[x_1, \ldots, x_r]$ 中的非零因子。由此，$H_U$ 的每个不可约分支的维数均为
$d + r - 1$，其中 $d = \dim(X) = \dim(A)$。若 $U' = U \cap D(f_1)$，则
$$
H_{U'} =
\Spec(A_{f_1}[x_1, \ldots, x_r]/(x_1f_1 + \ldots + x_rf_r)) \cong
% P05-ERR-CH37-0059: Frozen source omits the comma after \ldots in the polynomial-variable list below.
\Spec(A_{f_1}[x_2, \ldots x_r]) =
\mathbf{A}^{r - 1}_k \times_k U'
$$
不可约。另一方面，
$$
H_U \setminus H_{U'} =
\Spec(A/(f_1)[x_1, \ldots, x_r]/(x_2f_2 + \ldots + x_rf_r))
$$
的维数至多为 $d + r - 2$。事实上，对 $i \not = 1$，概形
$(H_U \setminus H_{U'}) \times_U D(f_i)$ 要么为空（若 $f_i = 0$），要么由与上面相同的论证，
同构于 $\Spec(A/(f_1))$ 的某个开集上的 $r - 1$ 维仿射空间，因而维数至多为
$d + r - 2$。另一方面，$(H_U \setminus H_{U'}) \times_U V(f_2, \ldots, f_r)$
是 $\Spec(A/(f_1, \ldots, f_r))$ 上的 $r$ 维仿射空间，故假设 (4) 表明其维数至多为
$d + r - 2$。于是对上述每个 $U$，$H_U$ 都不可约。因此 $H_{univ}$ 不可约。

\medskip\noindent
因此只需证明 $H_{univ}$ 的泛点在 $\mathbf{A}^r_k$ 的泛点上几何不可约，见
簇章引理
\ref{varieties-lemma-geometrically-irreducible-function-field}。在 $X$ 中选取包含于
$X \setminus H_{v_1}$ 的非空仿射开集 $U = \Spec(A)$，使其与假设 (3) 所断言存在的
$H_{v_2} \cap H_{v_3}$ 的不可约分支 $Z$ 相交。沿用上述记号，须证明域扩张
$$
\text{Frac}(A[x_1, \ldots, x_r]/(x_1f_1 + \ldots + x_rf_r))/
k(x_1, \ldots , x_r)
$$
几何不可约。由对 $U$ 的选取，注意 $f_1$ 在 $A$ 中可逆。置
$K = \text{Frac}(A)$ 为 $A$ 的分式域。像上面一样消去变量 $x_1$ 后，只须证明域扩张
$$
K(x_2, \ldots, x_r)/
k(x_2, \ldots, x_r, -\sum\nolimits_{i = 2, \ldots, r} f_1^{-1}f_i x_i)
$$
几何不可约。由引理 \ref{more-morphisms-lemma-amazing-bertini-lemma}，只需证明对某些
$c_2, \ldots, c_r \in k$，元素
$$
x_2, \ldots, x_r, \sum\nolimits_{i = 2, \ldots, r} f_1^{-1}f_i x_i,
\sum\nolimits_{i = 2, \ldots, r} c_if_1^{-1}f_i
$$
在 $A[x_2, \ldots, x_r]$ 的分式域中于 $k$ 上代数无关。这由引理
\ref{more-morphisms-lemma-algebraically-independent} 以及如下事实得出：$Z \cap U$ 是
$V(f_1^{-1}f_2, f_1^{-1}f_3) \subset U$ 的一个不可约分支。
\end{proof}

\begin{remark}
\label{more-morphisms-remark-geometric-proof-bertini-irreducible}
我们概述引理 \ref{more-morphisms-lemma-bertini-irreducible} 一个特殊情形的“几何”证明。设 $k$
是代数闭域，而 $X \subset \mathbf{P}^n_k$ 光滑、不可约且维数 $\geq 2$。我们断言存在超平面
$H \subset \mathbf{P}^n_k$，使得 $X \cap H$ 光滑且不可约。事实上，由 Varieties,
引理 \ref{varieties-lemma-bertini}，对一般的
$v \in V = kT_0 \oplus \ldots \oplus kT_n$，相应的超平面截面 $X \cap H_v$
是光滑的。另一方面，由 Enriques--Severi--Zariski，概形 $X \cap H_v$ 连通，见
簇章引理 \ref{varieties-lemma-connectedness-ample-divisor}。因此
$X \cap H_v$ 光滑且不可约。
\end{remark}













\section{立方体定理}
\label{more-morphisms-section-theorem-cube}

\noindent
下述引理说明，在引理所作假设下，Picard 函子的对角可由局部闭浸入表示。

\begin{lemma}
\label{more-morphisms-lemma-diagonal-picard-flat-proper}
设 $f : X \to S$ 是有限表现的平坦固有态射。设 $\mathcal{E}$ 是有限局部自由
$\mathcal{O}_X$-模。对态射 $g : T \to S$，考虑基变换图
$$
\xymatrix{
X_T \ar[d]_p \ar[r]_q & X \ar[d]^f \\
T \ar[r]^g & S
}
$$
。假设对所有 $g : T \to S$，$\mathcal{O}_T \to p_*\mathcal{O}_{X_T}$
均为同构。则存在有限表现浸入 $j : Z \to S$，使得态射 $g : T \to S$
经 $Z$ 分解，当且仅当存在有限局部自由 $\mathcal{O}_T$-模 $\mathcal{N}$，满足
$p^*\mathcal{N} \cong q^*\mathcal{E}$。
\end{lemma}

\begin{proof}
由假设 $H^0(X_s, \mathcal{O}_{X_s}) = \kappa(s)$，注意 $f$ 的纤维 $X_s$
连通。因此 $\mathcal{E}$ 的秩在纤维上为常数。由于 $f$ 是开态射（Morphisms,
引理 \ref{morphisms-lemma-fppf-open}）且是闭态射，$S$ 有分解
$S = \coprod S_r$ 为开闭子概形，使得 $\mathcal{E}$ 在 $S_r$
的逆像上具有常秩 $r$。因此可假设 $\mathcal{E}$ 具有常秩 $r$。以
$\mathcal{E}^\vee = \SheafHom(\mathcal{E}, \mathcal{O}_X)$ 表示其秩为 $r$ 的对偶模。

\medskip\noindent
由上同调与基变换（更确切地，由 概形的导出范畴章引理
\ref{perfect-lemma-flat-proper-perfect-direct-image-general}），
$E = Rf_*\mathcal{E}$ 是 $S$ 的导出范畴中的完美对象，且其构造与任意基变换交换。
$E' = Rf_*\mathcal{E}^\vee$ 同样如此。由于在次数 $< 0$ 中绝无上同调，可见对某个
$b$，$E$ 和 $E'$（局部）具有位于 $[0, b]$ 中的 tor-振幅。注意，对任意
$g : T \to S$，有 $p_*(q^*\mathcal{E}) = H^0(Lg^*E)$ 以及
$p_*(q^*\mathcal{E}^\vee) = H^0(Lg^*E')$。令 $j : Z \to S$ 和
$j' : Z' \to S$ 为 概形的导出范畴章引理
\ref{perfect-lemma-locally-closed-where-H0-locally-free} 对 $E$ 和 $E'$、取
$a = 0$ 和 $r = r$ 时构造的有限表现浸入；粗略地说，它们由如下性质刻画：
$H^0(Lj^*E)$ 和 $H^0((j')^*E')$ 是与拉回相容的有限局部自由模。

\medskip\noindent
设 $g : T \to S$ 是态射。若存在引理中的 $\mathcal{N}$，则使用投影公式
上同调章引理 \ref{cohomology-lemma-projection-formula}，可见各模
$$
p_*(q^*\mathcal{E}) \cong
p_*(p^*\mathcal{N}) \cong
\mathcal{N} \otimes_{\mathcal{O}_T} p_*\mathcal{O}_{X_T} \cong
\mathcal{N}\quad\text{and similarly }\quad
p_*(q^*\mathcal{E}^\vee) \cong \mathcal{N}^\vee
$$
是秩为 $r$ 的有限局部自由模，而且在任何进一步的基变换 $T' \to T$ 后仍然如此。
因此此时 $T \to S$ 经 $j$ 和 $j'$ 分解。故可将 $S$ 换成 $Z \times_S Z'$，
并假设 $f_*\mathcal{E}$ 和 $f_*\mathcal{E}^\vee$ 是秩为 $r$ 的有限局部自由
$\mathcal{O}_S$-模，其构造与任意基变换交换（略去一个小细节）。

\medskip\noindent
% P05-ERR-CH37-0060: Frozen source has the ungrammatical phrase ``if g ... be a morphism''.
在此情形，若 $g : T \to S$ 是态射并且存在引理中的 $\mathcal{N}$，则映射（次数
$0$ 中的杯积）
$$
p_*(q^*\mathcal{E})
\otimes_{\mathcal{O}_T}
p_*(q^*\mathcal{E}^\vee)
\longrightarrow \mathcal{O}_T
$$
是完美配对。反之，若此杯积映射是完美配对，则在 $T$ 上局部可以分别选取
$p_*(q^*\mathcal{E})$ 中的截面基 $\sigma_1, \ldots, \sigma_r$ 和
$p_*(q^*\mathcal{E}^\vee)$ 中的截面基 $\tau_1, \ldots, \tau_r$，使其乘积满足
$\sigma_i \tau_j = \delta_{ij}$。将 $\sigma_i$ 视为 $X_T$ 上 $q^*\mathcal{E}$
的截面，并将 $\tau_j$ 视为 $X_T$ 上 $q^*\mathcal{E}^\vee$ 的截面，可知
$$
\sigma_1, \ldots, \sigma_r :
\mathcal{O}_{X_T}^{\oplus r}
\longrightarrow
q^*\mathcal{E}
$$
是同构，其逆由下式给出：
$$
\tau_1, \ldots, \tau_r :
q^*\mathcal{E}
\longrightarrow
\mathcal{O}_{X_T}^{\oplus r}
$$
。换言之，有 $p^*p_*q^*\mathcal{E} \cong q^*\mathcal{E}$。而杯积非退化这一条件截出一个逆紧开子概形
（即某个适当行列式非零的轨迹），证明完成。
\end{proof}

\noindent
上述引理特别说明：若向量丛在一个固有空间的固有平坦族的各纤维上平凡，则它是某个向量丛的拉回。下面稍作展开。

\begin{lemma}
\label{more-morphisms-lemma-trivial-on-fibres}
设 $f : X \to S$ 是有限表现的平坦固有态射，满足
$f_*\mathcal{O}_X = \mathcal{O}_S$，并且此性质在任意基变换后仍成立。设
$\mathcal{E}$ 是有限局部自由 $\mathcal{O}_X$-模。假设
\begin{enumerate}
\item 对所有 $s \in S$，$\mathcal{E}|_{X_s}$ 同构于
$\mathcal{O}_{X_s}^{\oplus r_s}$；并且
\item $S$ 约化。
\end{enumerate}
则对某个有限局部自由 $\mathcal{O}_S$-模 $\mathcal{N}$，有
$\mathcal{E} = f^*\mathcal{N}$。
\end{lemma}

\begin{proof}
事实上，此时引理 \ref{more-morphisms-lemma-diagonal-picard-flat-proper} 中的局部闭浸入
$j : Z \to S$ 是双射，因而是闭浸入。但由于 $S$ 约化，$j$ 是同构。
\end{proof}

\begin{lemma}
\label{more-morphisms-lemma-triviality-generic-fibre-valuation-ring}
设 $f : X \to S$ 是有限表现的固有平坦态射。设 $\mathcal{L}$ 是可逆
$\mathcal{O}_X$-模。假设
\begin{enumerate}
\item $S$ 是赋值环的谱；
\item $\mathcal{L}$ 在 $f$ 的泛纤维 $X_\eta$ 上平凡；
\item $f$ 的闭纤维 $X_0$ 是整概形；
\item $H^0(X_\eta, \mathcal{O}_{X_\eta})$ 等于 $S$ 的函数域。
\end{enumerate}
则 $\mathcal{L}$ 平凡。
\end{lemma}

\begin{proof}
写成 $S = \Spec(A)$。先在 $A$ 为离散赋值环时证明本引理（这是实践中最常用的情形）。
令 $\pi \in A$ 为一致化元。取平凡化截面
$s \in \Gamma(X_\eta, \mathcal{L}_\eta)$。必要时以 $\pi^n s$ 代替 $s$ 后，
可假设 $s \in \Gamma(X, \mathcal{L})$。若 $s|_{X_0} = 0$，则 $s$ 可被
$\pi$ 整除（因为 $X_0$ 是概形论纤维且 $X$ 在 $A$ 上平坦）。因此可假设
$s|_{X_0}$ 非零。于是 $s$ 的零化轨迹 $Z(s)$ 包含于 $X_0$，但不包含 $X_0$
的泛点（因为 $X_0$ 是整概形）。这意味着
% P05-ERR-CH37-0061: Frozen source has the malformed phrase ``the Z(s)''.
$Z(s)$ 在 $X$ 中的余维 $\geq 2$；由 除子章引理
\ref{divisors-lemma-effective-Cartier-codimension-1}，这将导致矛盾，除非
$Z(s) = \emptyset$，正是所需。

\medskip\noindent
现在证明一般情形。赋值环 $A$ 是凝聚环（代数章例
\ref{algebra-example-valuation-ring-coherent}），故 $H^0(X, \mathcal{L})$ 是凝聚
$A$-模，见 概形的导出范畴章引理
\ref{perfect-lemma-cohomology-over-coherent-ring}。等价地，$H^0(X, \mathcal{L})$
是有限表现 $A$-模（代数章引理 \ref{algebra-lemma-coherent-ring}）。由于
$H^0(X, \mathcal{L})$ 无挠（由 $X$ 在 $A$ 上平坦），代数进阶章引理
\ref{more-algebra-lemma-generalized-valuation-ring-modules} 表明，对某个 $n$，有
$H^0(X, \mathcal{L}) = A^{\oplus n}$。由平坦基变换（Cohomology of Schemes,
引理 \ref{coherent-lemma-flat-base-change-cohomology}），有
$$
K = H^0(X_\eta, \mathcal{O}_{X_\eta}) \cong
H^0(X_\eta, \mathcal{L}_\eta) =
H^0(X, \mathcal{L}) \otimes_A K
$$
其中 $K$ 是 $A$ 的分式域。因此 $n = 1$。选取生成元
$s \in H^0(X, \mathcal{L})$。令 $\mathfrak m \subset A$ 为极大理想。则
$\kappa = A/\mathfrak m = \colim A/\pi$，这里取遍非零元
$\pi \in \mathfrak m$ 的滤过余极限（此处使用 $A$ 是赋值环）。因此
$X_0 = \lim X \times_S \Spec(A/\pi)$。若 $s|_{X_0}$ 为零，则对某个 $\pi$，
$s$ 在 $X \times_S \Spec(A/\pi)$ 上的限制为零，见 极限章引理
\ref{limits-lemma-descend-section}。但若如此，$\pi^{-1} s$ 就是 $\mathcal{L}$
的整体截面，这与 $s$ 是 $H^0(X, \mathcal{L})$ 的生成元矛盾。因此
$s|_{X_0}$ 非零。令 $Z(s) \subset X$ 为 $s$ 的零概形。由于 $s|_{X_0}$
非零且 $X_0$ 是整概形，$Z(s)_0 \subset X_0$ 是有效 Cartier 除子。由于 $f$
固有且 $S$ 局部，$Z(s)$ 的每一点都特殊化到 $Z(s)_0$ 中的一点。因此由 Divisors,
引理 \ref{divisors-lemma-fibre-Cartier} 的第 (3) 部分，$Z(s)$ 是相对有效
Cartier 除子；特别地，$Z(s) \to S$ 平坦。因此，若
% P05-ERR-CH37-0062: Frozen source has the typo ``nonemtpy''.
$Z(s)$ 非空，则 $Z(s)_\eta$ 也非空，这与 $s|_{X_\eta}$ 是
$\mathcal{L}_\eta$ 的平凡化矛盾。故 $Z(s) = \emptyset$，正是所需。
\end{proof}

\begin{lemma}
\label{more-morphisms-lemma-get-a-closed}
设 $f : X \to S$ 和 $\mathcal{E}$ 如引理
\ref{more-morphisms-lemma-diagonal-picard-flat-proper}，并且进一步假设 $\mathcal{E}$ 是可逆
$\mathcal{O}_X$-模。若 $f$ 的几何纤维还是整概形，则 $Z$ 在 $S$ 中闭。
\end{lemma}

\begin{proof}
由于 $j : Z \to S$ 有限表现，只需证明：对任意态射
$g : \Spec(A) \to S$，其中 $A$ 是分式域为 $K$ 的赋值环，若
$g(\Spec(K)) \in j(Z)$，则 $g(\Spec(A)) \subset j(Z)$。见 Morphisms,
引理 \ref{morphisms-lemma-reach-points-scheme-theoretic-image}。这由引理
\ref{more-morphisms-lemma-triviality-generic-fibre-valuation-ring} 以及引理
\ref{more-morphisms-lemma-diagonal-picard-flat-proper} 中对 $j : Z \to S$ 的刻画得出。
\end{proof}

\begin{lemma}
\label{more-morphisms-lemma-H1-O-picard-flat-proper}
考虑概形的交换图
$$
\xymatrix{
X' \ar[rr] \ar[dr]_{f'} & & X \ar[dl]^f \\
& S
}
$$
，其中 $f' : X' \to S$ 和 $f : X \to S$ 满足引理
\ref{more-morphisms-lemma-diagonal-picard-flat-proper} 的假设。设 $\mathcal{L}$ 是可逆
$\mathcal{O}_X$-模，并令 $\mathcal{L}'$ 为其在 $X'$ 上的拉回。分别令
$Z \subset S$ 和 $Z' \subset S$ 为引理
\ref{more-morphisms-lemma-diagonal-picard-flat-proper} 对 $(f, \mathcal{L})$ 和
$(f', \mathcal{L}')$ 构造的局部闭子概形，于是 $Z \subset Z'$。若 $s \in Z$ 且
$$
H^1(X_s, \mathcal{O}) \longrightarrow H^1(X'_s, \mathcal{O})
$$
为单射，则对 $s$ 的某个开邻域 $U$，有 $Z \cap U = Z' \cap U$。
\end{lemma}

\begin{proof}
可将 $S$ 换成 $Z'$。将 $S$ 缩小为 $s$ 的仿射开邻域后，可假设
$\mathcal{L}' = \mathcal{O}_{X'}$。令
$E = Rf_*\mathcal{L}$ 以及
$E' = Rf'_*\mathcal{L}' = Rf'_*\mathcal{O}_{X'}$。它们是完美复形，且其构造与任意基变换交换
（概形的导出范畴章引理
\ref{perfect-lemma-flat-proper-perfect-direct-image-general}）。特别地，
$$
E \otimes_{\mathcal{O}_S}^\mathbf{L} \kappa(s) =
R\Gamma(X_s, \mathcal{L}_s) = R\Gamma(X_s, \mathcal{O}_{X_s})
$$
，其中第二个等式源于 $s \in Z$。置
$h_i = \dim_{\kappa(s)} H^i(X_s, \mathcal{O}_{X_s})$。缩小 $S$ 后，可用复形
$$
\mathcal{O}_S \to \mathcal{O}_S^{\oplus h_1} \to
\mathcal{O}_S^{\oplus h_2} \to \ldots
$$
表示 $E$，见 代数进阶章引理
\ref{more-algebra-lemma-lift-perfect-from-residue-field}（严格地说，这还用到
概形的导出范畴章引理
\ref{perfect-lemma-affine-compare-bounded} 和
\ref{perfect-lemma-perfect-affine}）。类似地，可假设 $E'$ 由复形
$$
\mathcal{O}_S \to \mathcal{O}_S^{\oplus h'_1} \to
\mathcal{O}_S^{\oplus h'_2} \to \ldots
$$
表示，其中 $h'_i = \dim_{\kappa(s)} H^i(X'_s, \mathcal{O}_{X'_s})$。由上同调的函子性，有映射
$$
E \longrightarrow E'
$$
位于 $D(\mathcal{O}_S)$ 中，且其构造与基变换交换。由于表示 $E$ 的复形是有限自由模的有限复形，
而 $S$ 仿射，可以选取复形映射
$$
\xymatrix{
\mathcal{O}_S \ar[r]_d \ar[d]_a &
\mathcal{O}_S^{\oplus h_1} \ar[r] \ar[d]_b &
\mathcal{O}_S^{\oplus h_2} \ar[r] \ar[d]_c & \ldots \\
\mathcal{O}_S \ar[r]^{d'} &
\mathcal{O}_S^{\oplus h'_1} \ar[r] &
\mathcal{O}_S^{\oplus h'_2} \ar[r] & \ldots
}
$$
表示给定映射 $E \to E'$。由于 $s \in Z$，$\mathcal{L}_s$ 的平凡化截面拉回为
$\mathcal{L}'_s = \mathcal{O}_{X'_s}$ 的平凡化截面。因此
$a \otimes \kappa(s)$ 是同构，故缩小 $S$ 后，$a$ 是同构。最后，使用
$H^1(X_s, \mathcal{O}) \to H^1(X'_s, \mathcal{O})$ 为单射这一假设，可知定义
$b$ 的矩阵有一个 $h_1 \times h_1$ 子式映到 $\kappa(s)$ 中的非零元。因此缩小
$S$ 后，可假设 $b$ 为单射。然而，由于 $\mathcal{L}' = \mathcal{O}_{X'}$，有
$d' = 0$。于是 $d = 0$。由此可见，$\mathcal{L}_s$ 的平凡化截面提升为 $X$
上 $\mathcal{L}$ 的截面。一个直接的拓扑论证（略）表明，必要时进一步略微缩小
$S$ 后，$\mathcal{L}$ 平凡。
\end{proof}

\begin{lemma}
\label{more-morphisms-lemma-H1-O-multiple-picard-flat-proper}
考虑概形的 $n$ 个交换图
$$
\xymatrix{
X_i \ar[rr] \ar[dr]_{f_i} & & X \ar[dl]^f \\
& S
}
$$
，其中 $f_i : X_i \to S$ 和 $f : X \to S$ 满足引理
\ref{more-morphisms-lemma-diagonal-picard-flat-proper} 的假设。设 $\mathcal{L}$ 是可逆
$\mathcal{O}_X$-模，并令 $\mathcal{L}_i$ 为其在 $X_i$ 上的拉回。分别令
$Z \subset S$ 和 $Z_i \subset S$ 为引理
\ref{more-morphisms-lemma-diagonal-picard-flat-proper} 对 $(f, \mathcal{L})$ 和
$(f_i, \mathcal{L}_i)$ 构造的局部闭子概形，于是
$Z \subset \bigcap_{i = 1, \ldots, n} Z_i$。若 $s \in Z$ 且
$$
H^1(X_s, \mathcal{O}) \longrightarrow
\bigoplus\nolimits_{i = 1, \ldots, n} H^1(X_{i, s}, \mathcal{O})
$$
为单射，则对 $s$ 的某个开邻域 $U$，有
$Z \cap U = (\bigcap_{i = 1, \ldots, n} Z_i) \cap U$（概形论交）。
\end{lemma}

\begin{proof}
本引理是引理 \ref{more-morphisms-lemma-H1-O-picard-flat-proper} 的一个变体；我们强烈建议读者先阅读该证明。
这里的证明基本上是对该证明稍作修改后的复本。由 $Z$ 和各 $Z_i$ 的（概形值）点的描述，
有 $Z \subset \bigcap_{i = 1, \ldots, n} Z_i$，其中取概形论交。因此可将 $S$
换成概形论交 $\bigcap_{i = 1, \ldots, n} Z_i$。将 $S$ 缩小为 $s$ 的仿射开邻域后，
可假设对 $i = 1, \ldots, n$ 有 $\mathcal{L}_i = \mathcal{O}_{X_i}$。令
$E = Rf_*\mathcal{L}$ 以及
$E_i = Rf_{i, *}\mathcal{L}_i = Rf_{i, *}\mathcal{O}_{X_i}$。它们是完美复形，
且其构造与任意基变换交换（概形的导出范畴章引理
\ref{perfect-lemma-flat-proper-perfect-direct-image-general}）。特别地，
$$
E \otimes_{\mathcal{O}_S}^\mathbf{L} \kappa(s) =
R\Gamma(X_s, \mathcal{L}_s) = R\Gamma(X_s, \mathcal{O}_{X_s})
$$
，其中第二个等式源于 $s \in Z$。置
$h_j = \dim_{\kappa(s)} H^j(X_s, \mathcal{O}_{X_s})$。缩小 $S$ 后，可用复形
$$
\mathcal{O}_S \to \mathcal{O}_S^{\oplus h_1} \to
\mathcal{O}_S^{\oplus h_2} \to \ldots
$$
表示 $E$，见 代数进阶章引理
\ref{more-algebra-lemma-lift-perfect-from-residue-field}（严格地说，这还用到
概形的导出范畴章引理
\ref{perfect-lemma-affine-compare-bounded} 和
\ref{perfect-lemma-perfect-affine}）。类似地，可假设 $E_i$ 由复形
$$
\mathcal{O}_S \to \mathcal{O}_S^{\oplus h_{i, 1}} \to
\mathcal{O}_S^{\oplus h_{i, 2}} \to \ldots
$$
表示，其中 $h_{i, j} = \dim_{\kappa(s)} H^j(X_{i, s}, \mathcal{O}_{X_{i, s}})$。
由上同调的函子性，有映射
$$
E \longrightarrow E_i
$$
位于 $D(\mathcal{O}_S)$ 中，且其构造与基变换交换。由于表示 $E$ 的复形是有限自由模的有限复形，
而 $S$ 仿射，可以选取复形映射
$$
\xymatrix{
\mathcal{O}_S \ar[r]_d \ar[d]_{a_i} &
\mathcal{O}_S^{\oplus h_1} \ar[r] \ar[d]_{b_i} &
\mathcal{O}_S^{\oplus h_2} \ar[r] \ar[d]_{c_i} & \ldots \\
\mathcal{O}_S \ar[r]^{d_i} &
\mathcal{O}_S^{\oplus h_{i, 1}} \ar[r] &
\mathcal{O}_S^{\oplus h_{i, 2}} \ar[r] & \ldots
}
$$
表示给定映射 $E \to E_i$。由于 $s \in Z$，$\mathcal{L}_s$ 的平凡化截面拉回为
$\mathcal{L}_{i, s} = \mathcal{O}_{X_{i, s}}$ 的平凡化截面。因此
$a_i \otimes \kappa(s)$ 是同构，故缩小 $S$ 后，$a_i$ 是同构。最后，使用
$H^1(X_s, \mathcal{O}) \to
\bigoplus_{i = 1, \ldots, n} H^1(X_{i, s}, \mathcal{O})$ 为单射这一假设，
可知定义 $\oplus b_i$ 的矩阵有一个 $h_1 \times h_1$ 子式映到 $\kappa(s)$
中的非零元。因此缩小 $S$ 后，可假设
% P05-ERR-CH37-0063: Frozen source writes O_S^{h_1} here, unlike the direct-sum notation used for the surrounding free modules.
$(b_1, \ldots, b_n) : \mathcal{O}_S^{h_1}
\to \bigoplus_{i = 1, \ldots, n} \mathcal{O}_S^{h_{i, 1}}$
为单射。然而，由于 $\mathcal{L}_i = \mathcal{O}_{X_i}$，对
$i = 1, \ldots n$ 有 $d_i = 0$。又因
$(b_1, \ldots, b_n) \circ d = (\oplus d_i) \circ (a_1, \ldots, a_n)$，故
$d = 0$。由此可见，$\mathcal{L}_s$ 的平凡化截面提升为 $X$ 上 $\mathcal{L}$
的截面。一个直接的拓扑论证（略）表明，必要时进一步略微缩小 $S$ 后，
$\mathcal{L}$ 平凡。
\end{proof}

\begin{lemma}
\label{more-morphisms-lemma-pic-of-product}
设 $f : X \to S$ 和 $g : Y \to S$ 是满足引理
\ref{more-morphisms-lemma-diagonal-picard-flat-proper} 之假设的概形态射。设
$\sigma : S \to X$ 和 $\tau : S \to Y$ 分别为 $f$ 和 $g$ 的截面。设
$s \in S$。设 $\mathcal{L}$ 是 $X \times_S Y$ 上的可逆层。若 $X$ 上的
$(1 \times \tau)^*\mathcal{L}$、$Y$ 上的 $(\sigma \times 1)^*\mathcal{L}$ 以及
$\mathcal{L}|_{(X \times_S Y)_s}$ 均平凡，则存在 $s$ 的开邻域 $U$，使得
$\mathcal{L}$ 在 $(X \times_S Y)_U$ 上平凡。
\end{lemma}

\begin{proof}
由 K\"unneth 公式（簇章引理 \ref{varieties-lemma-kunneth}），映射
$$
H^1(X_s \times_{\Spec(\kappa(s))} Y_s, \mathcal{O}) \to
H^1(X_s, \mathcal{O}) \oplus H^1(Y_s, \mathcal{O})
$$
为单射。因此可将引理 \ref{more-morphisms-lemma-H1-O-multiple-picard-flat-proper} 应用于两个态射
$$
1 \times \tau : X \to X \times_S Y
\quad\text{且}\quad
\sigma \times 1 : Y \to X \times_S Y
$$
，从而得到结论。
\end{proof}

\begin{theorem}[立方体定理]
\label{more-morphisms-theorem-of-the-cube}
设 $S$ 是概形。设 $X$、$Y$ 和 $Z$ 是 $S$ 上的概形。设 $x : S \to X$ 和
$y : S \to Y$ 为结构态射的截面。设 $\mathcal{L}$ 是
$X \times_S Y \times_S Z$ 上的可逆模。若
\begin{enumerate}
\item $X \to S$ 和 $Y \to S$ 是具有几何整纤维的有限表现平坦固有态射；
\item $\mathcal{L}$ 经 $x \times \text{id}_Y \times \text{id}_Z$ 和
$\text{id}_X \times y \times \text{id}_Z$ 的拉回分别在 $Y \times_S Z$ 和
$X \times_S Z$ 上平凡；
\item 存在一点 $z \in Z$，使得 $\mathcal{L}$ 在
$X \times_S Y \times_S z$ 上的限制平凡；并且
\item $Z$ 连通；
\end{enumerate}
则 $\mathcal{L}$ 平凡。
\end{theorem}

\noindent
一个常用的特殊情形如下。设 $k$ 是域。设 $X, Y, Z$ 是带有 $k$-有理点
$x, y, z$ 的簇。设 $\mathcal{L}$ 是 $X \times Y \times Z$ 上的可逆模。若
\begin{enumerate}
\item $\mathcal{L}$ 在 $x \times Y \times Z$、$X \times y \times Z$ 和
$X \times Y \times z$ 上平凡；并且
\item $X$ 和 $Y$ 在 $k$ 上几何整且固有；
\end{enumerate}
则 $\mathcal{L}$ 平凡。

\begin{proof}
% P05-ERR-CH37-0064: Frozen source omits the verb in ``whose geometrically integral fibres''.
注意，态射 $X \times_S Y \to S$ 是具有几何整纤维的有限表现平坦固有态射（关于纤维的陈述，
见 簇章引理 \ref{varieties-lemma-geometrically-integral}、
\ref{varieties-lemma-bijection-irreducible-components} 和
\ref{varieties-lemma-geometrically-reduced-any-base-change}）。由 Derived Categories of Schemes,
引理 \ref{perfect-lemma-proper-flat-geom-red-connected}，结构层经 $X \to S$、
$Y \to S$ 或 $X \times_S Y \to S$ 的推前都是 $S$ 的结构层，并且这在任何基变换后仍成立。
因此可将引理 \ref{more-morphisms-lemma-diagonal-picard-flat-proper} 应用于态射
$$
p : X \times_S Y \times_S Z \longrightarrow Z
$$
和可逆模 $\mathcal{L}$，得到“泛”局部闭子概形 $Z' \subset Z$，使得
$\mathcal{L}|_{X \times_S Y \times_S Z'}$ 是 $Z'$ 上某个可逆模 $\mathcal{N}$ 的拉回。
$z$ 的存在表明 $Z'$ 非空。由引理 \ref{more-morphisms-lemma-get-a-closed}，$Z' \subset Z$
是闭子概形。设 $z' \in Z'$ 为一点。注意，可将 $p$ 写成积态射
$$
(X \times_S Z) \times_Z (Y \times_S Z) \longrightarrow Z
$$
因此可将引理 \ref{more-morphisms-lemma-pic-of-product} 应用于态射 $p$、点 $z'$ 以及由 $x$ 和 $y$
给出的截面 $\sigma : Z \to X \times_S Z$ 和 $\tau : Z \to Y \times_S Z$。
由此 $Z'$ 是开集。因此 $Z' = Z$，且对 $Z$ 上某个可逆模 $\mathcal{N}$，有
$\mathcal{L} = p^*\mathcal{N}$。经
$x \times y \times \text{id}_Z : Z  \to X \times_S Y \times_S Z$
拉回，一方面得到 $\mathcal{N}$，另一方面由假设 (2) 得到平凡可逆模。因此
$\mathcal{N} = \mathcal{O}_Z$，证明完成。
\end{proof}









\section{极限论证}
\label{more-morphisms-section-limits}

\noindent
本节给出一些涉及概形极限与 Noether 逼近的引理。我们主要限于仿射情形。其中一些引理是极限章中引理的特殊情形。

\begin{lemma}
\label{more-morphisms-lemma-Noetherian-approximation}
设 $f : X \to S$ 是有限表现的仿射概形态射。则存在 Cartesian 图
$$
\xymatrix{
X_0 \ar[d]_{f_0} & X \ar[l]^g \ar[d]^f \\
S_0 & S \ar[l]
}
$$
使得
\begin{enumerate}
\item $X_0$、$S_0$ 是仿射概形；
\item $S_0$ 在 $\mathbf{Z}$ 上有限型；
\item $f_0$ 是有限型的。
\end{enumerate}
\end{lemma}

\begin{proof}
写成 $S = \Spec(A)$ 和 $X = \Spec(B)$。由于 $f$ 有限表现，$B$
作为 $A$-代数有限表现，见 态射章引理
\ref{morphisms-lemma-locally-finite-presentation-characterize}。因此本引理由
代数章引理 \ref{algebra-lemma-limit-module-finite-presentation} 得出。
\end{proof}

\begin{lemma}
\label{more-morphisms-lemma-Noetherian-approximation-module}
设 $f : X \to S$ 是有限表现的仿射概形态射。设 $\mathcal{F}$ 是有限表现的拟凝聚
$\mathcal{O}_X$-模。则存在引理 \ref{more-morphisms-lemma-Noetherian-approximation} 中的图，并且存在凝聚
$\mathcal{O}_{X_0}$-模 $\mathcal{F}_0$，满足 $g^*\mathcal{F}_0 = \mathcal{F}$。
\end{lemma}

\begin{proof}
写成 $S = \Spec(A)$、$X = \Spec(B)$ 以及 $\mathcal{F} = \widetilde{M}$。
由于 $f$ 有限表现，$B$ 作为 $A$-代数有限表现，见 态射章引理
\ref{morphisms-lemma-locally-finite-presentation-characterize}。由于 $\mathcal{F}$
作为 $\mathcal{O}_X$-模有限表现，$M$ 作为 $B$-模有限表现，见 Properties,
引理 \ref{properties-lemma-finite-presentation-module}。因此本引理由 代数章引理
\ref{algebra-lemma-limit-module-finite-presentation} 得出。
\end{proof}

\begin{lemma}
\label{more-morphisms-lemma-Noetherian-approximation-flat-module}
设 $f : X \to S$ 是有限表现的仿射概形态射。设 $\mathcal{F}$ 是有限表现的拟凝聚
$\mathcal{O}_X$-模，并且在 $S$ 上平坦。则可选取引理
\ref{more-morphisms-lemma-Noetherian-approximation-module} 中的图和层 $\mathcal{F}_0$，使得此外
$\mathcal{F}_0$ 在 $S_0$ 上平坦。
\end{lemma}

\begin{proof}
写成 $S = \Spec(A)$、$X = \Spec(B)$ 以及 $\mathcal{F} = \widetilde{M}$。
由于 $f$ 有限表现，$B$ 作为 $A$-代数有限表现，见 态射章引理
\ref{morphisms-lemma-locally-finite-presentation-characterize}。由于 $\mathcal{F}$
作为 $\mathcal{O}_X$-模有限表现，$M$ 作为 $B$-模有限表现，见 Properties,
引理 \ref{properties-lemma-finite-presentation-module}。由于 $\mathcal{F}$ 在 $S$
上平坦，$M$ 在 $A$ 上平坦，见 态射章引理
\ref{morphisms-lemma-flat-module-characterize}。因此本引理由 代数章引理
\ref{algebra-lemma-flat-finite-presentation-limit-flat} 得出。
\end{proof}

\begin{lemma}
\label{more-morphisms-lemma-Noetherian-approximation-flat}
设 $f : X \to S$ 是有限表现且平坦的仿射概形态射。则存在引理
\ref{more-morphisms-lemma-Noetherian-approximation} 中的图，使得此外 $f_0$ 平坦。
\end{lemma}

\begin{proof}
这是引理 \ref{more-morphisms-lemma-Noetherian-approximation-flat-module} 的特殊情形。
\end{proof}

\begin{lemma}
\label{more-morphisms-lemma-Noetherian-approximation-smooth}
设 $f : X \to S$ 是仿射概形的光滑态射。则存在引理
\ref{more-morphisms-lemma-Noetherian-approximation} 中的图，使得此外 $f_0$ 光滑。
\end{lemma}

\begin{proof}
写成 $S = \Spec(A)$、$X = \Spec(B)$。由于 $f$ 光滑，$B$ 作为 $A$-代数光滑，
见 态射章引理 \ref{morphisms-lemma-smooth-characterize}。因此本引理由
代数章引理 \ref{algebra-lemma-finite-presentation-fs-Noetherian} 得出。
\end{proof}

\begin{lemma}
\label{more-morphisms-lemma-Noetherian-approximation-geometrically-reduced}
设 $f : X \to S$ 是具有几何约化纤维的有限表现仿射概形态射。则存在引理
\ref{more-morphisms-lemma-Noetherian-approximation} 中的图，使得此外 $f_0$ 具有几何约化纤维。
\end{lemma}

\begin{proof}
应用引理 \ref{more-morphisms-lemma-Noetherian-approximation} 得到仿射概形的 Cartesian 图
$$
\xymatrix{
X_0 \ar[d]_{f_0} & X \ar[l]^g \ar[d]^f \\
S_0 & S \ar[l]_h
}
$$
，其中 $X_0 \to S_0$ 是在 $\mathbf{Z}$ 上有限型的概形之间的有限型态射。由引理
\ref{more-morphisms-lemma-geometrically-reduced-constructible}，其中 $f_0$ 的纤维几何约化的点集
$E \subset S_0$ 是可构造子集。由引理
\ref{more-morphisms-lemma-base-change-fibres-geometrically-reduced}，有 $h(S) \subset E$。写成
$S_0 = \Spec(A_0)$ 和 $S = \Spec(A)$。将 $A = \colim_i A_i$ 写成有限型
$A_0$-代数的直余极限。由 极限章引理
\ref{limits-lemma-limit-contained-in-constructible}，对某个 $i$，
$\Spec(A_i) \to S_0$ 的像包含于 $E$。将 $S_0$ 换成 $\Spec(A_i)$，并将
$X_0$ 换成 $X_0 \times_{S_0} \Spec(A_i)$ 后，可见 $f_0$ 的所有纤维均几何约化。
\end{proof}

\begin{lemma}
\label{more-morphisms-lemma-Noetherian-approximation-geometrically-irreducible}
设 $f : X \to S$ 是具有几何不可约纤维的有限表现仿射概形态射。则存在引理
\ref{more-morphisms-lemma-Noetherian-approximation} 中的图，使得此外 $f_0$ 具有几何不可约纤维。
\end{lemma}

\begin{proof}
应用引理 \ref{more-morphisms-lemma-Noetherian-approximation} 得到仿射概形的 Cartesian 图
$$
\xymatrix{
X_0 \ar[d]_{f_0} & X \ar[l]^g \ar[d]^f \\
S_0 & S \ar[l]_h
}
$$
，其中 $X_0 \to S_0$ 是在 $\mathbf{Z}$ 上有限型的概形之间的有限型态射。由引理
\ref{more-morphisms-lemma-nr-geom-irreducible-components-constructible}，其中 $f_0$
的纤维几何不可约的点集 $E \subset S_0$ 是可构造子集。由引理
\ref{more-morphisms-lemma-base-change-fibres-geometrically-irreducible}，有 $h(S) \subset E$。写成
$S_0 = \Spec(A_0)$ 和 $S = \Spec(A)$。将 $A = \colim_i A_i$ 写成有限型
$A_0$-代数的直余极限。由 极限章引理
\ref{limits-lemma-limit-contained-in-constructible}，对某个 $i$，
$\Spec(A_i) \to S_0$ 的像包含于 $E$。将 $S_0$ 换成 $\Spec(A_i)$，并将
$X_0$ 换成 $X_0 \times_{S_0} \Spec(A_i)$ 后，可见 $f_0$ 的所有纤维均几何不可约。
\end{proof}

\begin{lemma}
\label{more-morphisms-lemma-Noetherian-approximation-geometrically-connected}
设 $f : X \to S$ 是具有几何连通纤维的有限表现仿射概形态射。则存在引理
\ref{more-morphisms-lemma-Noetherian-approximation} 中的图，使得此外 $f_0$ 具有几何连通纤维。
\end{lemma}

\begin{proof}
应用引理 \ref{more-morphisms-lemma-Noetherian-approximation} 得到仿射概形的 Cartesian 图
$$
\xymatrix{
X_0 \ar[d]_{f_0} & X \ar[l]^g \ar[d]^f \\
S_0 & S \ar[l]_h
}
$$
，其中 $X_0 \to S_0$ 是在 $\mathbf{Z}$ 上有限型的概形之间的有限型态射。由引理
\ref{more-morphisms-lemma-nr-geom-connected-components-constructible}，其中 $f_0$
的纤维几何连通的点集 $E \subset S_0$ 是可构造子集。由引理
\ref{more-morphisms-lemma-base-change-fibres-geometrically-connected}，有 $h(S) \subset E$。写成
$S_0 = \Spec(A_0)$ 和 $S = \Spec(A)$。将 $A = \colim_i A_i$ 写成有限型
$A_0$-代数的直余极限。由 极限章引理
\ref{limits-lemma-limit-contained-in-constructible}，对某个 $i$，
$\Spec(A_i) \to S_0$ 的像包含于 $E$。将 $S_0$ 换成 $\Spec(A_i)$，并将
$X_0$ 换成 $X_0 \times_{S_0} \Spec(A_i)$ 后，可见 $f_0$ 的所有纤维均几何连通。
\end{proof}

\begin{lemma}
\label{more-morphisms-lemma-Noetherian-approximation-dimension-d}
设 $d \geq 0$ 是整数。设 $f : X \to S$ 是有限表现的仿射概形态射，且其所有纤维的维数均为
$d$。则存在引理 \ref{more-morphisms-lemma-Noetherian-approximation} 中的图，使得此外 $f_0$
的所有纤维维数均为 $d$。
\end{lemma}

\begin{proof}
应用引理 \ref{more-morphisms-lemma-Noetherian-approximation} 得到仿射概形的 Cartesian 图
$$
\xymatrix{
X_0 \ar[d]_{f_0} & X \ar[l]^g \ar[d]^f \\
S_0 & S \ar[l]_h
}
$$
，其中 $X_0 \to S_0$ 是在 $\mathbf{Z}$ 上有限型的概形之间的有限型态射。由引理
\ref{more-morphisms-lemma-dimension-fibres-constructible}，其中 $f_0$ 的纤维维数为 $d$
的点集 $E \subset S_0$ 是可构造子集。由引理
\ref{more-morphisms-lemma-base-change-dimension-fibres}，有 $h(S) \subset E$。写成
$S_0 = \Spec(A_0)$ 和 $S = \Spec(A)$。将 $A = \colim_i A_i$ 写成有限型
$A_0$-代数的直余极限。由 极限章引理
\ref{limits-lemma-limit-contained-in-constructible}，对某个 $i$，
$\Spec(A_i) \to S_0$ 的像包含于 $E$。将 $S_0$ 换成 $\Spec(A_i)$，并将
$X_0$ 换成 $X_0 \times_{S_0} \Spec(A_i)$ 后，可见 $f_0$ 的所有纤维维数均为 $d$。
\end{proof}

\begin{lemma}
\label{more-morphisms-lemma-Noetherian-approximation-standard-syntomic}
设 $f : X \to S$ 是仿射概形的标准 syntomic 态射（见 态射章定义
\ref{morphisms-definition-syntomic}）。则存在引理
\ref{more-morphisms-lemma-Noetherian-approximation} 中的图，使得此外 $f_0$ 是标准 syntomic 态射。
\end{lemma}

\begin{proof}
本引理是 代数章引理
\ref{algebra-lemma-relative-global-complete-intersection-Noetherian} 的复本。
\end{proof}

\begin{lemma}
\label{more-morphisms-lemma-Noetherian-approximation-combine}
（Noether 逼近与性质合并。）设 $P$、$Q$ 是在基变换下稳定的概形态射性质。设
$f : X \to S$ 是有限表现的仿射概形态射。假设可以找到 Cartesian 图
$$
\vcenter{
\xymatrix{
X_1 \ar[d]_{f_1} & X \ar[l] \ar[d]^f \\
S_1 & S \ar[l]
}
}
\quad\text{且}\quad
\vcenter{
\xymatrix{
X_2 \ar[d]_{f_2} & X \ar[l] \ar[d]^f \\
S_2 & S \ar[l]
}
}
$$
，其中 $S_1$、$S_2$ 在 $\mathbf{Z}$ 上有限型，$f_1$、$f_2$ 是有限型的，且
$f_1$ 具有性质 $P$，$f_2$ 具有性质 $Q$。则可以找到 Cartesian 图
$$
\xymatrix{
X_0 \ar[d]_{f_0} & X \ar[l] \ar[d]^f \\
S_0 & S \ar[l]
}
$$
，其中 $S_0$ 在 $\mathbf{Z}$ 上有限型，$f_0$ 是有限型的，并且 $f_0$
同时具有性质 $P$ 与性质 $Q$。
\end{lemma}

\begin{proof}
给定的一对图对应于环的余 Cartesian 图
$$
\vcenter{
\xymatrix{
B_1 \ar[r] & B \\
A_1 \ar[u] \ar[r] & A \ar[u]
}
}
\quad\text{且}\quad
\vcenter{
\xymatrix{
B_2 \ar[r] & B \\
A_2 \ar[u] \ar[r] & A \ar[u]
}
}
$$
令 $A_0 \subset A$ 为 $A$ 的有限型 $\mathbf{Z}$-子代数，包含
$A_1 \to A$ 和 $A_2 \to A$ 的像。这样的子代数存在，因为按假设 $A_1$ 和
$A_2$ 均在 $\mathbf{Z}$ 上有限型。注意，各环
$B_{0, 1} = B_1 \otimes_{A_1} A_0$ 和 $B_{0, 2} = B_2 \otimes_{A_2} A_0$
是有限型 $A_0$-代数，并满足
$B_{0, 1} \otimes_{A_0} A \cong B \cong B_{0, 2} \otimes_{A_0} A$
作为 $A$-代数同构。由于 $A$ 是其有限型 $A_0$-子代数的有向余极限，由 Limits,
引理 \ref{limits-lemma-descend-finite-presentation}，增大 $A_0$ 后可假设存在
$A_0$-代数同构 $B_{0, 1} \cong B_{0, 2}$。由于假设性质 $P$ 和 $Q$
在基变换下稳定，置 $S_0 = \Spec(A_0)$ 以及
$$
X_0 = X_1 \times_{S_1} S_0 =
\Spec(B_{0, 1}) \cong \Spec(B_{0, 2}) = X_2 \times_{S_2} S_0
$$
即可满足要求。
\end{proof}












\section{\'Etale 邻域}
\label{more-morphisms-section-etale-neighbourhoods}

\noindent
事实表明，作 \'etale 基变换后，态射的某些性质更易研究。引入以下术语会很方便。

\begin{definition}
\label{more-morphisms-definition-etale-neighbourhood}
设 $S$ 是概形，$s \in S$ 是一点。
\begin{enumerate}
\item $(S, s)$ 的一个{\it \'etale 邻域}是一个二元组 $(U, u)$，连同概形的
\'etale 态射 $\varphi : U \to S$，满足 $\varphi(u) = s$。
\item $(S, s)$ 的 \'etale 邻域之间的一个{\it 态射}
$f : (V, v) \to (U, u)$，就是满足 $f(v) = u$ 的 $S$-概形态射 $f : V \to U$。
\item {\it 初等 \'etale 邻域}是满足 $\kappa(s) = \kappa(u)$ 的 \'etale 邻域
$\varphi : (U, u) \to (S, s)$。
\end{enumerate}
\end{definition}

\noindent
% P05-ERR-CH37-0065: Frozen source has mismatched parentheses around the two literature citations in this sentence.
初等 \'etale 邻域这一概念在文献中有许多不同名称，例如有时称为
“\'etale 邻域”（\cite[Page 36]{Milne}）或“强 \'etale”（\cite[Page 108]{KPR}）。
这里遵循论文 \cite{GruRay} 的约定，称其为初等 \'etale 邻域。

\medskip\noindent
若 $f : (V, v) \to (U, u)$ 是 \'etale 邻域之间的态射，则 $f$ 自动为
\'etale，见 态射章引理 \ref{morphisms-lemma-etale-permanence}。因此它使
$(V, v)$ 成为 $(U, u)$ 的 \'etale 邻域。当然，由于 \'etale 态射的复合仍为
\'etale（态射章引理 \ref{morphisms-lemma-composition-etale}），反过来，
$(U, u)$ 的任意 \'etale 邻域 $(V, v)$ 也为 $(S, s)$ 的 \'etale 邻域。还要指出，
若 $U \subset S$ 是 $s$ 的开邻域，则 $(U, s) \to (S, s)$ 是 \'etale 邻域。
这是因为开浸入为 \'etale 态射（态射章引理
\ref{morphisms-lemma-open-immersion-etale}）。本节将不再说明地使用这些评注。

\medskip\noindent
注意，若 $(U, u) \to (S, s)$ 是 \'etale 邻域，则
$\kappa(u)/\kappa(s)$ 是有限可分扩张，见 态射章引理
\ref{morphisms-lemma-etale-at-point}。

\begin{lemma}
\label{more-morphisms-lemma-realize-prescribed-residue-field-extension-etale}
设 $S$ 是概形，$s \in S$。设 $k/\kappa(s)$ 是有限可分域扩张。则存在 \'etale 邻域
$(U, u) \to (S, s)$，使得域扩张 $\kappa(u)/\kappa(s)$ 同构于 $k/\kappa(s)$。
\end{lemma}

\begin{proof}
可假设 $S$ 仿射。此时本引理由 代数章引理
\ref{algebra-lemma-make-etale-map-prescribed-residue-field} 得出。
\end{proof}

\begin{lemma}
\label{more-morphisms-lemma-etale-neighbourhoods-not-quite-filtered}
设 $S$ 是概形，并设 $s$ 是 $S$ 的一点。\'Etale 邻域范畴具有以下性质：
\begin{enumerate}
\item 设 $(U_i, u_i)_{i=1, 2}$ 是 $S$ 中 $s$ 的两个 \'etale 邻域。则存在第三个
\'etale 邻域 $(U, u)$ 以及态射
$(U, u) \to (U_i, u_i)$, $i = 1, 2$.
\item 设 $h_1, h_2: (U, u) \to (U', u')$ 是 $s$ 的 \'etale 邻域之间的两个态射。
假设 $h_1$、$h_2$ 诱导相同的剩余域映射 $\kappa(u') \to \kappa(u)$。则存在
\'etale 邻域 $(U'', u'')$ 以及态射 $h : (U'', u'') \to (U, u)$，使 $h_1$
与 $h_2$ 等化，即满足 $h_1 \circ h = h_2 \circ h$。
\end{enumerate}
\end{lemma}

\begin{proof}
对第 (1) 部分，考虑纤维积 $U = U_1 \times_S U_2$。由于 \'etale 态射在基变换下保持，
它在 $U_1$ 和 $U_2$ 上均为 \'etale 态射，见 态射章引理
\ref{morphisms-lemma-base-change-etale}。例如由 概形章引理
\ref{schemes-lemma-points-fibre-product} 对纤维积之点的描述，$U$ 中有一点同时映到
$u_1$ 和 $u_2$。对第 (2) 部分，将 $U''$ 定义为纤维积
$$
\xymatrix{
U'' \ar[r] \ar[d] & U \ar[d]^{(h_1, h_2)} \\
U' \ar[r]^-\Delta & U' \times_S U'.
}
$$
。由于 $h_1$ 和 $h_2$ 诱导相同的剩余域映射
$\kappa(u') \to \kappa(u)$，存在位于 $u'$ 上方的点 $u'' \in U''$，满足
$\kappa(u'') = \kappa(u')$。特别地，$U'' \not = \emptyset$。此外，由于
$U'$ 在 $S$ 上 \'etale，纤维积 $U'\times_S U'$ 也如此（见 态射章引理
\ref{morphisms-lemma-base-change-etale} 和
\ref{morphisms-lemma-composition-etale}）。因此由 态射章引理
\ref{morphisms-lemma-etale-permanence}，竖直箭头 $(h_1, h_2)$ 为 \'etale 态射。
故由基变换，$U''$ 在 $U'$ 上 \'etale，进而也在 $S$ 上 \'etale（因为 \'etale
态射的复合仍为 \'etale 态射）。因此 $(U'', u'')$ 满足要求。
\end{proof}

\begin{lemma}
\label{more-morphisms-lemma-elementary-etale-neighbourhoods}
设 $S$ 是概形，并设 $s$ 是 $S$ 的一点。$(S, s)$ 的初等 \'etale 邻域范畴是余滤过的
（见 范畴章定义 \ref{categories-definition-codirected}）。
\end{lemma}

\begin{proof}
这立即由定义和引理 \ref{more-morphisms-lemma-etale-neighbourhoods-not-quite-filtered} 得出。
\end{proof}

\begin{lemma}
\label{more-morphisms-lemma-describe-henselization}
设 $S$ 是概形，并设 $s \in S$。则
$$
\mathcal{O}_{S, s}^h =
\colim_{(U, u)} \mathcal{O}(U)
$$
，其中余极限取遍 $(S, s)$ 的初等 \'etale 邻域 $(U, u)$ 的范畴之反范畴；后者为滤过范畴。
\end{lemma}

\begin{proof}
设 $\Spec(A) \subset S$ 是 $s$ 的仿射邻域。令 $\mathfrak p \subset A$ 为对应于
$s$ 的素理想。作此选取后，有典范同构
$\mathcal{O}_{S, s} = A_{\mathfrak p}$ 和 $\kappa(s) = \kappa(\mathfrak p)$。
由满足如下条件的初等 \'etale 邻域 $(U, u)$ 给出一个余终系统：$U$ 仿射且
$U \to S$ 经 $\Spec(A)$ 分解。换言之，右端等于
$\colim_{(B, \mathfrak q)} B$，其中余极限取遍配备一个位于 $\mathfrak p$
上方且满足 $\kappa(\mathfrak p) = \kappa(\mathfrak q)$ 的素理想 $\mathfrak q$
的 \'etale $A$-代数 $B$。因此本引理由 代数章引理
\ref{algebra-lemma-henselization-different} 得出。
\end{proof}

\noindent
可以把纤维上各点的 \'etale 邻域提升到全空间。

\begin{lemma}
\label{more-morphisms-lemma-lift-etale-neighbourhood-fibre}
\begin{slogan}
把纤维上一点的 \'etale 邻域提升到全空间。
\end{slogan}
设 $X \to S$ 是概形态射。设 $x \in X$ 的像为 $s \in S$。设
$(V, v) \to (X_s, x)$ 是 \'etale 邻域。则存在 \'etale 邻域
$(U, u) \to (X, x)$，并且存在 $(X_s, x)$ 的 \'etale 邻域之间的态射
$(U_s, u) \to (V, v)$，它是开浸入。
\end{lemma}

\begin{proof}
可假设 $X$、$V$ 和 $S$ 仿射。设态射 $X \to S$ 由 $A \to B$ 给出，点
$x$ 由素理想 $\mathfrak q \subset B$ 给出，点 $s$ 由
$\mathfrak p = A \cap \mathfrak q$ 给出，而态射 $V \to X_s$ 由
$B \otimes_A \kappa(\mathfrak p) \to C$ 给出。由于 $\kappa(\mathfrak p)$
是 $A/\mathfrak p$ 的一个局部化，存在 $f \in A$、$f \not \in \mathfrak p$
以及 \'etale 环同态 $B \otimes_A (A/\mathfrak p)_f \to D$，使得
$$
C = (B \otimes_A \kappa(\mathfrak p))
\otimes_{B \otimes_A (A/\mathfrak p)_f} D
$$
，见 代数章引理 \ref{algebra-lemma-etale} 的第 (9) 部分。将 $A$ 换成
$A_f$、将 $B$ 换成 $B_f$ 后，可假设 $D$ 在
$B \otimes_A A/\mathfrak p = B/\mathfrak p B$ 上 \'etale。于是可以应用
代数章引理 \ref{algebra-lemma-lift-etale}。本引理得证。
\end{proof}














\section{\'Etale 邻域与分支}
\label{more-morphisms-section-etale-nbhds-branches}

\noindent
概形在一点处的（几何）分支数已在 Properties, Section
\ref{properties-section-unibranch} 中定义。在 Varieties, Section
\ref{varieties-section-number-of-branches} 中，我们把它与正规化态射的纤维联系起来。
本节讨论用 \'etale 邻域刻画这一数目。

\begin{lemma}
\label{more-morphisms-lemma-nr-minimal-primes}
设 $R = \colim R_i$ 是一个环的有向系统的余极限，且其过渡同态忠实平坦。
则将 $R$ 的极小素理想数视为 $\{0, 1, 2, \ldots, \infty\}$ 中的元素时，
它等于各 $R_i$ 的极小素理想数的上确界。
\end{lemma}

\begin{proof}
若 $A \to B$ 是平坦环同态，则由下降定理，$\Spec(B) \to \Spec(A)$
把极小素理想映到极小素理想（代数章引理
\ref{algebra-lemma-flat-going-down}）。若 $A \to B$ 忠实平坦，则每个极小素理想都是某个极小素理想的像
（由 代数章引理 \ref{algebra-lemma-ff-rings} 和
\ref{algebra-lemma-minimal-prime-image-minimal-prime}）。因此，若 $i \leq i'$，则
% P05-ERR-CH37-0066: Frozen source reverses the monotonicity inequality for minimal-prime counts under a faithfully flat transition R_i -> R_i'.
$R_i$ 的极小素理想数 $\geq$ $R_{i'}$ 的极小素理想数。由 代数章引理
\ref{algebra-lemma-colimit-faithfully-flat}，每个同态 $R_i \to R$ 都忠实平坦，且还可见
$R$ 的极小素理想数 $\geq$ $R_i$ 的极小素理想数。最后，设
$\mathfrak q_1, \ldots, \mathfrak q_n$ 是 $R$ 中两两不同的极小素理想。则可以找到
$i$，使得 $R_i \cap \mathfrak q_1, \ldots, R_i \cap \mathfrak q_n$
两两不同（作为集合，因而也作为素理想）。这蕴含本引理。
\end{proof}

\begin{lemma}
\label{more-morphisms-lemma-nr-branches}
设 $X$ 是概形，$x \in X$ 是一点。则
\begin{enumerate}
\item $X$ 在 $x$ 处的分支数，等于当初等 \'etale 邻域
$(U, u) \to (X, x)$ 变化时，$U$ 中经过 $u$ 的不可约分支数的上确界；
\item $X$ 在 $x$ 处的几何分支数，等于当 \'etale 邻域
$(U, u) \to (X, x)$ 变化时，$U$ 中经过 $u$ 的不可约分支数的上确界；
\item $X$ 在 $x$ 处单分支，当且仅当对每个初等 \'etale 邻域
$(U, u) \to (X, x)$，$U$ 中恰有一个经过 $u$ 的不可约分支；并且
\item $X$ 在 $x$ 处几何单分支，当且仅当对每个 \'etale 邻域
$(U, u) \to (X, x)$，$U$ 中恰有一个经过 $u$ 的不可约分支。
\end{enumerate}
\end{lemma}

\begin{proof}
第 (3)、(4) 部分由第 (1)、(2) 部分和 性质章引理
\ref{properties-lemma-number-of-branches-1} 得出。

\medskip\noindent
证明 (1)。设 $\Spec(A)$ 是 $x$ 的仿射开邻域，并令 $\mathfrak p \subset A$
为对应于 $x$ 的素理想。可将 $X$ 换成 $\Spec(A)$；由于仿射初等 \'etale 邻域
$(U, u)$ 构成余终子系统，在取上确界时只需考虑它们。回忆 Hensel 化
$A_\mathfrak p^h$ 是各环 $B_\mathfrak q$ 在二元组 $(B, \mathfrak q)$
范畴上的余极限，其中 $B$ 是 \'etale $A$-代数，而 $\mathfrak q$ 是位于
$\mathfrak p$ 上方且满足 $\kappa(\mathfrak q) = \kappa(\mathfrak p)$ 的素理想，见
代数章引理 \ref{algebra-lemma-henselization-different}。由环与仿射概形之间的对应，
这些二元组 $(B, \mathfrak q)$ 恰好对应于仿射初等 \'etale 邻域 $(U, u)$。
注意，$\Spec(B)$ 中经过 $\mathfrak q$ 的不可约分支，恰为
$B_\mathfrak q$ 的极小素理想。由 性质章定义
\ref{properties-definition-number-of-branches}，$A_\mathfrak p^h$ 的极小素理想数就是
$X$ 在 $x$ 处的分支数。注意，系统中的过渡同态
$B_\mathfrak q \to B'_{\mathfrak q'}$ 全都平坦。由于平坦局部环同态忠实平坦
（代数章引理 \ref{algebra-lemma-local-flat-ff}），本引理由引理
\ref{more-morphisms-lemma-nr-minimal-primes} 得出。

\medskip\noindent
证明 (2)。除使用 代数章引理
\ref{algebra-lemma-strict-henselization-different} 外，证明与 (1) 相同。有一个细微差别：
给定 $\kappa(x)$ 的可分代数闭包 $\kappa^{sep}$，对每个 \'etale 邻域
$(U, u)$，都可以选取 $\kappa(x)$-嵌入
$\phi : \kappa(u) \to \kappa^{sep}$，因为 $\kappa(u)/\kappa(x)$ 有限可分
（态射章引理 \ref{morphisms-lemma-etale-at-point}）。因此可在所有三元组
$(U, u, \phi)$ 上取上确界，其中 $(U, u) \to (X, x)$ 是仿射 \'etale 邻域，
而 $\phi : \kappa(u) \to \kappa^{sep}$ 是 $\kappa(x)$-嵌入。这些三元组恰好对应于
代数章引理 \ref{algebra-lemma-strict-henselization-different} 中的三元组，
其余证明完全相同。
\end{proof}

\noindent
我们还需要上述引理的相对版本。

\begin{lemma}
\label{more-morphisms-lemma-nr-branches-fibre}
设 $X \to S$ 是概形态射，$x \in X$ 是像为 $s$ 的点。则
\begin{enumerate}
\item 纤维 $X_s$ 在 $x$ 处的分支数，等于当初等 \'etale 邻域
$(U, u) \to (X, x)$ 变化时，纤维 $U_s$ 中经过 $u$ 的不可约分支数的上确界；
\item 纤维 $X_s$ 在 $x$ 处的几何分支数，等于当 \'etale 邻域
$(U, u) \to (X, x)$ 变化时，纤维 $U_s$ 中经过 $u$ 的不可约分支数的上确界；
\item 纤维 $X_s$ 在 $x$ 处单分支，当且仅当对每个初等 \'etale 邻域
$(U, u) \to (X, x)$，纤维 $U_s$ 中恰有一个经过 $u$ 的不可约分支；并且
% P05-ERR-CH37-0067: Frozen source switches from the fibre X_s to X in item (4), unlike the relative statement and the three preceding items.
\item $X$ 在 $x$ 处几何单分支，当且仅当对每个 \'etale 邻域
$(U, u) \to (X, x)$，$U_s$ 中恰有一个经过 $u$ 的不可约分支。
\end{enumerate}
\end{lemma}

\begin{proof}
合并引理 \ref{more-morphisms-lemma-nr-branches} 与
\ref{more-morphisms-lemma-lift-etale-neighbourhood-fibre}。
\end{proof}

\begin{lemma}
\label{more-morphisms-lemma-number-of-branches-and-smooth}
设 $X \to S$ 是概形的光滑态射。设 $x \in X$ 的像为 $s \in S$。则
\begin{enumerate}
\item $X$ 在 $x$ 处的几何分支数等于 $S$ 在 $s$ 处的几何分支数。
\item 若 $\kappa(x)/\kappa(s)$ 是纯不可分\footnote{事实上，只需
$\kappa(x)$ 在 $\kappa(s)$ 上几何不可约。若日后需要，我们会补充详细证明。}
域扩张，则 $X$ 在 $x$ 处的分支数等于 $S$ 在 $s$ 处的分支数。
\end{enumerate}
\end{lemma}

\begin{proof}
这立即由 代数进阶章引理
\ref{more-algebra-lemma-invariance-number-branches-smooth} 和定义得出。
\end{proof}




\section{非分歧态射与 \'etale 态射}
\label{more-morphisms-section-unramified-is-etale}

\noindent
有时非分歧态射自动为 \'etale 态射。

\begin{lemma}
\label{more-morphisms-lemma-unramified-dominant-unibranch-is-etale}
设 $f : X \to Y$ 是概形态射。设 $x \in X$ 的像为 $y \in Y$。假设
\begin{enumerate}
\item $Y$ 是整概形，且在 $y$ 处几何单分支；
\item $f$ 局部有限型；
\item 存在特殊化 $x' \leadsto x$，使得 $f(x')$ 是 $Y$ 的泛点；
\item $f$ 在 $x$ 处非分歧。
\end{enumerate}
则 $f$ 在 $x$ 处为 \'etale 态射。
\end{lemma}

\begin{proof}
可将 $X$ 和 $Y$ 换成 $x$ 和 $y$ 的适当仿射开邻域。于是 $Y$ 是整环 $A$
的谱，而 $X$ 是有限型 $A$-代数 $B$ 的谱。令 $\mathfrak q \subset B$
为对应于 $x$ 的素理想，$\mathfrak p \subset A$ 为对应于 $y$ 的素理想。局部环
$A_\mathfrak p = \mathcal{O}_{Y, y}$ 几何单分支。环同态 $A \to B$ 在
$\mathfrak q$ 处非分歧。此外，(3) 中的点 $x'$ 对应于满足
$A \cap \mathfrak q' = (0)$ 的素理想 $\mathfrak q' \subset \mathfrak q$。
因此 $A_\mathfrak p \to B_\mathfrak q$ 为单射。结论由 代数进阶章引理
\ref{more-algebra-lemma-local-unramified-extension-unibranch-domain-is-etale} 得出。
\end{proof}

\begin{lemma}
\label{more-morphisms-lemma-global-unramified-dominant-unibranch-is-etale}
\begin{reference}
\cite[Expose I, Corollary 9.11]{SGA1}
\end{reference}
设 $f : X \to Y$ 是概形态射。假设
\begin{enumerate}
\item $Y$ 是整的且几何单分支；
\item $X$ 至少有一个不可约分支支配 $Y$；
\item $f$ 非分歧；并且
\item $X$ 连通。
\end{enumerate}
则 $f$ 为 \'etale 态射，且 $X$ 不可约。
\end{lemma}

\begin{proof}
令 $X' \subset X$ 为支配 $Y$ 的不可约分支。这意味着 $X'$ 的泛点映到 $Y$
的泛点（例如见 态射章引理
\ref{morphisms-lemma-dominant-finite-number-irreducible-components}）。由引理
\ref{more-morphisms-lemma-unramified-dominant-unibranch-is-etale}，$f$ 在 $X'$ 的每一点处为
\'etale 态射。特别地，$f$ 为 \'etale 的开子概形 $U \subset X$ 包含 $X'$。
注意，$U$ 的每个拟紧开集只有有限多个不可约分支，见 下降章引理
\ref{descent-lemma-locally-finite-nr-irred-local-fppf}。另一方面，由于 $Y$
几何单分支且 $U$ 在 $Y$ 上 \'etale，概形 $U$ 几何单分支。特别地，$U$
的每一点至多经过一个不可约分支。一个简单的拓扑论证表明，$X' \subset U$
既开又闭。于是 $X'$ 在 $X$ 中当然也既开又闭；由连通性得到
$X = U = X'$，正是所需。
\end{proof}

\begin{lemma}
\label{more-morphisms-lemma-weird-permanence-etale}
设 $f : X \to Y$ 和 $g : Y \to Z$ 是概形态射。设 $x \in X$ 的像为
$y \in Y$。假设
\begin{enumerate}
\item $Y$ 是整概形，且在 $y$ 处几何单分支；
\item $g \circ f$ 在 $x$ 处为 \'etale 态射；
\item 存在特殊化 $x' \leadsto x$，使得 $f(x')$ 是 $Y$ 的泛点。
\end{enumerate}
则 $f$ 在 $x$ 处为 \'etale 态射，且 $g$ 在 $y$ 处为 \'etale 态射。
\end{lemma}

\begin{proof}
将 $X$ 换成 $x$ 的一个开邻域后，可假设 $g \circ f$ 为 \'etale 态射。于是由
态射章引理 \ref{morphisms-lemma-unramified-permanence} 和
\ref{morphisms-lemma-etale-smooth-unramified}，$f$ 非分歧。因此由引理
\ref{more-morphisms-lemma-unramified-dominant-unibranch-is-etale}，$f$ 在 $x$ 处为 \'etale 态射。
再由 \'etale 下降，$g$ 在 $y$ 处为 \'etale 态射，例如见 下降章引理
\ref{descent-lemma-syntomic-smooth-etale-permanence}。
\end{proof}

\begin{lemma}
\label{more-morphisms-lemma-global-weird-permanence-etale}
设 $f : X \to Y$ 和 $g : Y \to Z$ 是概形态射。假设
\begin{enumerate}
\item $Y$ 是整的且几何单分支；
\item $g \circ f$ 为 \'etale 态射；
\item $X$ 的每个不可约分支都支配 $Y$。
\end{enumerate}
则 $f$ 为 \'etale 态射，且 $g$ 在 $f$ 的像中的每一点处为 \'etale 态射。
\end{lemma}

\begin{proof}
这立即由逐点版本引理 \ref{more-morphisms-lemma-weird-permanence-etale} 得出。
\end{proof}










\section{光滑态射的切片}
\label{more-morphisms-section-etale-over-smooth}

\noindent
本节阐释一个大意如下的结果：概形 $S$ 的光滑覆盖可由 \'etale 覆盖加细。其证明技术依赖于切片论证。

\begin{lemma}
\label{more-morphisms-lemma-slice-smooth-given-element}
设 $f : X \to S$ 是概形态射。设 $x \in X$ 的像为 $s \in S$。设
$h \in \mathfrak m_x \subset \mathcal{O}_{X, x}$。假设
\begin{enumerate}
\item $f$ 在 $x$ 处光滑；并且
\item $\text{d}h$ 在
$$
\Omega_{X_s/s, x} \otimes_{\mathcal{O}_{X_s, x}} \kappa(x) =
\Omega_{X/S, x} \otimes_{\mathcal{O}_{X, x}} \kappa(x)
$$
中的像 $\text{d}\overline{h}$ 非零。
\end{enumerate}
则存在 $x$ 的仿射开邻域 $U \subset X$，使得 $h$ 来自
$h \in \Gamma(U, \mathcal{O}_U)$，并且 $D = V(h)$ 是 $U$ 中的有效 Cartier
除子，满足 $x \in D$，且 $D \to S$ 光滑。
\end{lemma}

\begin{proof}
由于 $f$ 在 $x$ 处光滑，将 $X$ 替换为 $x$ 的一个开邻域后，可假设 $f$
光滑。特别地，$f$ 平坦且局部有限表现。由
引理 \ref{more-morphisms-lemma-slice-given-element}
已知存在 $x$ 的开邻域 $U \subset X$，使得 $h$ 来自
$h \in \Gamma(U, \mathcal{O}_U)$，并且 $D = V(h)$ 是 $U$ 中满足
$x \in D$ 的有效 Cartier 除子，而 $D \to S$ 平坦且有限表现。由
态射章引理 \ref{morphisms-lemma-differentials-relative-immersion}
有短正合列
$$
\mathcal{C}_{D/U} \to i^*\Omega_{U/S} \to \Omega_{D/S} \to 0
$$
其中 $i : D \to U$ 是闭浸入，而 $\mathcal{C}_{D/U}$ 是 $D$ 在 $U$ 中的余法层。
由于 $D$ 是由 $h \in \Gamma(U, \mathcal{O}_U)$ 截出的有效 Cartier 除子，可知
% P05-ERR-CH37-0068: Frozen source identifies the conormal sheaf with h · O_S, although the sheaf is on D and the usual invertible conormal is I/I^2.
$\mathcal{C}_{D/U} = h \cdot \mathcal{O}_S$。由于 $U \to S$ 光滑，层
$\Omega_{U/S}$ 有限局部自由，故其拉回 $i^*\Omega_{U/S}$ 也有限局部自由。
该序列的第一个箭头把自由生成元 $h$ 映到 $i^*\Omega_{U/S}$ 的截面
$\text{d}h|_D$；按假设，它在纤维
$\Omega_{U/S, x} \otimes \kappa(x)$ 中的值非零。由 $\otimes \kappa(x)$ 的右正合性可得
$$
\dim_{\kappa(x)} \left( \Omega_{D/S, x} \otimes \kappa(x) \right)
=
\dim_{\kappa(x)} \left( \Omega_{U/S, x} \otimes \kappa(x) \right) - 1.
$$
由
态射章引理 \ref{morphisms-lemma-smooth-at-point}
可知，$\Omega_{U/S, x} \otimes \kappa(x)$ 至多由 $\dim_x(U_s)$ 个元素生成。
由上式，$\Omega_{D/S, x} \otimes \kappa(x)$ 至多由
$\dim_x(U_s) - 1$ 个元素生成。注意到
$\dim_x(D_s) = \dim_x(U_s) - 1$；例如，由
代数章引理 \ref{algebra-lemma-one-equation}
有 $\dim(\mathcal{O}_{D_s, x}) = \dim(\mathcal{O}_{U_s, x}) - 1$
（$D_s \subset U_s$ 也是有效 Cartier 除子，见
除子章引理 \ref{divisors-lemma-relative-Cartier})
，再使用
态射章引理 \ref{morphisms-lemma-dimension-fibre-at-a-point}.
因此，$\Omega_{D/S, x} \otimes \kappa(x)$ 至多由 $\dim_x(D_s)$ 个元素生成，
再由
态射章引理 \ref{morphisms-lemma-smooth-at-point}
可知 $D \to S$ 在 $x$ 处光滑。缩小 $U$ 后，$D \to S$ 光滑，结论得证。
\end{proof}

\begin{lemma}
\label{more-morphisms-lemma-slice-smooth-once}
设 $f : X \to S$ 是概形态射。设 $x \in X$ 的像为 $s \in S$。假设
\begin{enumerate}
\item $f$ 在 $x$ 处光滑；并且
\item 映射
$$
\Omega_{X_s/s, x} \otimes_{\mathcal{O}_{X_s, x}} \kappa(x)
\longrightarrow
\Omega_{\kappa(x)/\kappa(s)}
$$
的核非零。
\end{enumerate}
则存在 $x$ 的仿射开邻域 $U \subset X$，以及包含 $x$ 的有效 Cartier 除子
$D \subset U$，使得 $D \to S$ 光滑。
\end{lemma}

\begin{proof}
记 $k = \kappa(s)$ 且 $R = \mathcal{O}_{X_s, x}$。以 $\mathfrak m$ 表示 $R$
的极大理想，并令 $\kappa = R/\mathfrak m$，于是 $\kappa = \kappa(x)$。
微分模的形成与局部化交换（见
代数章引理 \ref{algebra-lemma-differentials-localize})
，故 $\Omega_{X_s/s, x} = \Omega_{R/k}$。由
代数章引理 \ref{algebra-lemma-differential-seq}
有正合列
$$
\mathfrak m/\mathfrak m^2 \xrightarrow{\text{d}}
\Omega_{R/k} \otimes_R \kappa \to
\Omega_{\kappa/k} \to 0.
$$
因此，若 (2) 成立，则存在元素 $\overline{h} \in \mathfrak m$，使得
$\text{d}\overline{h}$ 非零。选取 $\overline{h}$ 的提升
$h \in \mathcal{O}_{X, x}$，并应用
引理 \ref{more-morphisms-lemma-slice-smooth-given-element}.
\end{proof}

\begin{remark}
\label{more-morphisms-remark-necessary-condition-slice-smooth}
即使 $x$ 是正维纤维的闭点，
引理 \ref{more-morphisms-lemma-slice-smooth-once}
中的第二个条件也是必要的。下面给出一个例子：设 $k$ 是特征为 $p > 0$
的非完美域。取不是 $p$ 次幂的元素 $a \in k$。令
$\mathfrak m = (x, y^p - a) \subset k[x, y]$。它对应于
$X = \mathbf{A}^2_k$ 的闭点 $w$。令 $S = \mathbf{A}^1_k$，并令
$f : X \to S$ 为对应于 $k[x] \to k[x, y]$ 的态射。此时不存在交换图
$$
\xymatrix{
S' \ar[rr]_h \ar[rd]_g & & X \ar[ld]^f \\
& S
}
$$
使得 $g$ 为 \'etale 且 $w$ 位于 $h$ 的像中。这是显然的，因为剩余域扩张
$\kappa(w)/\kappa(f(w))$ 是纯不可分的，而对任何满足
$g(s') = f(w)$ 的 $s' \in S'$，扩张 $\kappa(s')/\kappa(f(w))$ 都将是可分的。
\end{remark}

\noindent
若假设剩余域扩张可分，则
说明 \ref{more-morphisms-remark-necessary-condition-slice-smooth}
中的现象不会发生。准确结果如下。

\begin{lemma}
\label{more-morphisms-lemma-slice-smooth-once-separable-residue-field-extension}
设 $f : X \to S$ 是概形态射。设 $x \in X$ 的像为 $s \in S$。假设
\begin{enumerate}
\item $f$ 在 $x$ 处光滑；
\item 剩余域扩张 $\kappa(x)/\kappa(s)$ 可分；并且
\item $x$ 不是 $X_s$ 的泛点。
\end{enumerate}
则存在 $x$ 的仿射开邻域 $U \subset X$，以及包含 $x$ 的有效 Cartier 除子
$D \subset U$，使得 $D \to S$ 光滑。
\end{lemma}

\begin{proof}
记 $k = \kappa(s)$ 且 $R = \mathcal{O}_{X_s, x}$。以 $\mathfrak m$ 表示 $R$
的极大理想，并令 $\kappa = R/\mathfrak m$，于是 $\kappa = \kappa(x)$。
微分模的形成与局部化交换（见
代数章引理 \ref{algebra-lemma-differentials-localize})
，故 $\Omega_{X_s/s, x} = \Omega_{R/k}$。由假设 (2) 及
代数章引理 \ref{algebra-lemma-computation-differential}
，映射
$$
\text{d} :
\mathfrak m/\mathfrak m^2
\longrightarrow
\Omega_{R/k} \otimes_R \kappa(\mathfrak m)
$$
是单射。假设 (3) 蕴含 $\mathfrak m/\mathfrak m^2 \not = 0$。
因此存在元素 $\overline{h} \in \mathfrak m$，使得
$\text{d}\overline{h}$ 非零。选取 $\overline{h}$ 的提升
$h \in \mathcal{O}_{X, x}$，并应用
引理 \ref{more-morphisms-lemma-slice-smooth-given-element}.
\end{proof}

\noindent
下一引理构造的子概形 $Z$ 实际上是 $x$ 的某个仿射开邻域中的完全交。若以后需要这一事实，
我们会另行明确表述一个引理。

\begin{lemma}
\label{more-morphisms-lemma-slice-smooth}
设 $f : X \to S$ 是概形态射。设 $x \in X$ 的像为 $s \in S$。假设
\begin{enumerate}
\item $f$ 在 $x$ 处光滑；并且
\item $x$ 是 $X_s$ 的闭点，且 $\kappa(s) \subset \kappa(x)$ 可分。
\end{enumerate}
则存在包含 $x$ 的浸入 $Z \to X$，使得
\begin{enumerate}
\item $Z \to S$ 为 \'etale；并且
\item 集合论意义下 $Z_s = \{x\}$。
\end{enumerate}
\end{lemma}

\begin{proof}
我们可以并且确实把 $S$ 替换为 $s$ 的一个仿射开邻域。我们也可以并且确实把
$X$ 替换为 $x$ 的一个仿射开邻域，使得 $X \to S$ 光滑。下文对仿射概形间的
光滑态射按 $d = \dim_x(X_s)$ 归纳证明此引理。

\medskip\noindent
情形 $d = 0$。此时证明可取 $Z$ 为 $x$ 的一个开邻域。事实上，若 $d = 0$，则
$X \to S$ 在 $x$ 处拟有限，见
态射章引理 \ref{morphisms-lemma-locally-quasi-finite-rel-dimension-0}.
因此存在仿射开邻域 $U \subset X$，使得 $U \to S$ 拟有限，见
态射章引理 \ref{morphisms-lemma-quasi-finite-points-open}.
于是把 $X$ 替换为 $U$ 后，$X$ 在 $S$ 上拟有限且光滑，因而在 $S$ 上相对维数为
$0$ 且光滑，故在 $S$ 上为 \'etale。此外，纤维 $X_s$ 是有限离散集。因此，
% P05-ERR-CH37-0069: Frozen source says an affine open neighbourhood of X; the intended distinguished point is x.
把 $X$ 替换为 $X$ 的另一个仿射开邻域后，可使 $f^{-1}(\{s\}) = \{x\}$
（因为 $X_s$ 的拓扑由 $X$ 的拓扑诱导，见
概形章引理 \ref{schemes-lemma-fibre-topological}).
这就证明了此情形。

\medskip\noindent
接着假设 $d > 0$。注意到，由于 $x$ 是其纤维的闭点，扩张
$\kappa(x)/\kappa(s)$ 有限（由 Hilbert 零点定理，见
态射章引理 \ref{morphisms-lemma-closed-point-fibre-locally-finite-type}).
由于代数可分域扩张具有此性质，故 $\Omega_{\kappa(x)/\kappa(s)} = 0$。因此可以应用
引理 \ref{more-morphisms-lemma-slice-smooth-once}
得到图
$$
\xymatrix{
D \ar[r] \ar[rrd] & U \ar[r] \ar[rd] & X \ar[d] \\
& & S
}
$$
其中 $x \in D$。注意到 $\dim_x(D_s) = \dim_x(X_s) - 1$；例如，由
代数章引理 \ref{algebra-lemma-one-equation}
有 $\dim(\mathcal{O}_{D_s, x}) = \dim(\mathcal{O}_{X_s, x}) - 1$
（$D_s \subset X_s$ 也是有效 Cartier 除子，见
除子章引理 \ref{divisors-lemma-relative-Cartier})
，再使用
态射章引理 \ref{morphisms-lemma-dimension-fibre-at-a-point}.
因此态射 $D \to S$ 光滑，且
$\dim_x(D_s) = \dim_x(X_s) - 1 = d - 1$。由归纳假设，可找到所需的浸入
$Z \to D$，证明完成。
\end{proof}

\begin{lemma}
\label{more-morphisms-lemma-etale-nbhd-dominates-smooth}
\begin{slogan}
光滑态射具有 \'etale 局部截面。
\end{slogan}
\begin{reference}
见 \cite[Corollaire 17.16.3 (ii)]{EGA4}。
\end{reference}
设 $f : X \to S$ 是概形的光滑态射。设 $s \in S$ 是 $f$ 像中的点。
则存在 \'etale 邻域 $(S', s') \to (S, s)$ 以及 $S$-态射 $S' \to X$。
\end{lemma}

\begin{proof}[引理 \ref{more-morphisms-lemma-etale-nbhd-dominates-smooth} 的第一种证明]
按假设 $X_s \not = \emptyset$。由
簇章引理 \ref{varieties-lemma-smooth-separable-closed-points-dense}
，存在闭点 $x \in X_s$，使得 $\kappa(x)$ 是 $\kappa(s)$ 的有限可分域扩张。因此由
引理 \ref{more-morphisms-lemma-slice-smooth}
，存在浸入 $Z \to X$，使得 $Z \to S$ 为 \'etale 且 $x \in Z$。取
$(S' , s') = (Z, x)$。
\end{proof}

\begin{proof}[引理 \ref{more-morphisms-lemma-etale-nbhd-dominates-smooth} 的第二种证明]
选取满足 $f(x) = s$ 的点 $x \in X$。选择图
$$
\xymatrix{
X \ar[d] & U \ar[l] \ar[d] \ar[r]_-\pi & \mathbf{A}^d_V \ar[ld] \\
S & V \ar[l]
}
$$
使得 $\pi$ 为 \'etale，$x \in U$，且 $V = \Spec(R)$ 仿射，见
态射章引理 \ref{morphisms-lemma-smooth-etale-over-affine-space}.
特别地，$s \in V$。态射 $\pi : U \to \mathbf{A}^d_V$ 是开映射，见
态射章引理 \ref{morphisms-lemma-etale-open}.
因此 $W = \pi(U) \cap \mathbf{A}^d_s$ 是 $\mathbf{A}^d_s$ 的非空开子集。
取点 $w \in W$，使得 $\kappa(s) \subset \kappa(w)$ 有限可分，见
簇章引理 \ref{varieties-lemma-affine-space-over-field}.
由
代数章引理 \ref{algebra-lemma-dim-affine-space}
，存在 $d$ 个元素
$\overline{f}_1, \ldots, \overline{f}_d \in \kappa(s)[x_1, \ldots, x_d]$，
它们生成 $\kappa(s)[x_1, \ldots, x_d]$ 中对应于 $w$ 的极大理想。
把 $R$ 替换为一个主局部化后，可假设存在
$f_1, \ldots, f_d \in R[x_1, \ldots, x_d]$，其像为
$\overline{f}_1, \ldots, \overline{f}_d \in \kappa(s)[x_1, \ldots, x_d]$。
考虑 $R$-代数
$$
R' = R[x_1, \ldots, x_d]/(f_1, \ldots, f_d)
$$
并令 $S' = \Spec(R')$。按构造，有 $V$ 上的闭浸入
$j : S' \to \mathbf{A}^d_V$。按构造，$S' \to V$ 在 $s$ 上的纤维是单点
$s'$，其剩余域是 $\kappa(s)$ 的有限可分扩张。设
$\mathfrak q' \subset R'$ 为相应素理想。由
代数章引理 \ref{algebra-lemma-localize-relative-complete-intersection}
可知，对某个 $g \in R'$，$(R')_g$ 是 $R$ 上的相对全局完全交，且
% P05-ERR-CH37-0070: Frozen source uses undefined q here after defining q'.
$g \not \in \mathfrak q$。因此 $S' \to V$ 在 $s'$ 的一个邻域上平坦且有限表现，见
代数章引理 \ref{algebra-lemma-relative-global-complete-intersection}.
按构造，$S' \to V$ 在 $s$ 上的概形论纤维是 $\Spec(\kappa(s'))$。故由
态射章引理 \ref{morphisms-lemma-etale-at-point}
可知，$S' \to S$ 在 $s'$ 处为 \'etale。令
$$
S'' = U \times_{\pi, \mathbf{A}^d_V, j} S'.
$$
按构造，存在点 $s'' \in S''$，它经投影 $p : S'' \to S'$ 映到 $s'$。
注意到 $p$ 是 \'etale 态射 $\pi$ 的基变换，故为 \'etale，见
态射章引理 \ref{morphisms-lemma-base-change-etale}.
选取 $s''$ 的一个小仿射邻域 $S''' \subset S''$，使它映入 $s' \in S'$
的非空开邻域，而态射 $S' \to S$ 在该邻域上为 \'etale。于是 \'etale 邻域
$(S''', s'') \to (S, s)$ 即为本引理所求。
\end{proof}

\noindent
下一引理表明，光滑拓扑上的层与 \'etale 拓扑上的层是同一回事。

\begin{lemma}
\label{more-morphisms-lemma-etale-dominates-smooth}
设 $S$ 是概形。设 $\mathcal{U} = \{S_i \to S\}_{i \in I}$ 是 $S$ 的光滑覆盖，见
拓扑章定义 \ref{topologies-definition-smooth-covering}.
则存在 \'etale 覆盖 $\mathcal{V} = \{T_j \to S\}_{j \in J}$（见
拓扑章定义 \ref{topologies-definition-etale-covering})
，它加细 $\mathcal{U}$（见
位点章定义 \ref{sites-definition-morphism-coverings})
。
\end{lemma}

\begin{proof}
对每个 $s \in S$，存在 $i \in I$，使得 $s$ 位于 $S_i \to S$ 的像中。由
引理 \ref{more-morphisms-lemma-etale-nbhd-dominates-smooth}
可找到 \'etale 态射 $g_s : T_s \to S$，使得 $s \in g_s(T_s)$，且 $g_s$
经 $S_i \to S$ 分解。因此 $\{T_s \to S\}$ 是加细 $\mathcal{U}$ 的 $S$ 的
\'etale 覆盖。
\end{proof}

\begin{lemma}
\label{more-morphisms-lemma-cover-smooth-by-special}
设 $f : X \to S$ 是概形的光滑态射。则存在 \'etale 覆盖
$\{U_i \to X\}_{i \in I}$，使得 $U_i \to S$ 分解为
$U_i \to V_i \to S$，其中 $V_i \to S$ 为 \'etale，而 $U_i \to V_i$
是仿射概形间的光滑态射，具有截面且纤维几何连通。
\end{lemma}

\begin{proof}
设 $s \in S$。由
簇章引理 \ref{varieties-lemma-smooth-separable-closed-points-dense}
，使得 $\kappa(x)/\kappa(s)$ 可分的闭点 $x \in X_s$ 在 $X_s$ 中稠密。
因此，只需构造像包含 $x$ 的 \'etale 态射 $U \to X$，使得 $U \to S$
按引理所述方式分解。为此，选取经过 $x$ 的浸入 $Z \to X$，使得
$Z \to S$ 为 \'etale（引理 \ref{more-morphisms-lemma-slice-smooth}）。把 $S$ 替换为 $Z$，
把 $X$ 替换为 $Z \times_S X$ 后，可假设 $X \to S$ 具有截面
$\sigma : S \to X$，且 $\sigma(s) = x$。接着先把 $S$ 替换为 $s$ 的仿射开邻域，
再把 $X$ 替换为 $x$ 的仿射开邻域。最后，考虑 Section
\ref{more-morphisms-section-connected-components} 的子集 $X^0 \subset X$。由 Lemmas
\ref{more-morphisms-lemma-connected-along-section-open} 和
\ref{more-morphisms-lemma-connected-along-section-locally-constructible}
可知，这是一个包含 $\sigma$ 的逆紧开子概形，并且 $X^0 \to S$ 的纤维几何连通。
若 $X^0$ 不仿射，则选取包含 $x$ 的仿射开集 $U \subset X^0$。由于
$X^0 \to S$ 光滑，$U$ 的像是开集。选取 $s$ 的仿射开邻域 $V \subset S$，
使它包含于 $\sigma^{-1}(U)$ 及 $U \to S$ 的像中。最后，读者可看出
$U \cap f^{-1}(V) \to V$ 具有全部所需性质。例如，
$U \cap f^{-1}(V)$ 等于 $U \times_S V$，因而作为仿射概形的纤维积而仿射。此外，
% P05-ERR-CH37-0071: Frozen source says the fibres of X^0 -> S are irreducible although only geometric connectedness was established.
$U \cap f^{-1}(V) \to V$ 的几何纤维是 $X^0 \to S$ 的不可约纤维的非空开子概形，
故为连通的。略去若干细节。
\end{proof}











\section{\'Etale 邻域与 Artin 逼近}
\label{more-morphisms-section-etale-nbhds-artin}

\noindent
本节证明如下形式的结果：若两个带点概形具有同构的完备局部环，则它们具有同构的
\'etale 邻域。我们将依赖 Popescu 定理，见
环同态光滑化章定理 \ref{smoothing-theorem-popescu}.

\begin{lemma}
\label{more-morphisms-lemma-map-approximation}
设 $S$ 是局部 Noether 概形。设 $X$、$Y$ 是 $S$ 上局部有限型的概形。设
$x \in X$ 与 $y \in Y$ 位于同一点 $s \in S$ 之上。假设
$\mathcal{O}_{S, s}$ 是 G-环。再假设给定局部 $\mathcal{O}_{S, s}$-代数映射
$$
\varphi : \mathcal{O}_{Y, y} \longrightarrow \mathcal{O}_{X, x}^\wedge
$$
。对每个 $N \geq 1$，存在初等 \'etale 邻域
$(U, u) \to (X, x)$ 以及把 $u$ 映到 $y$ 的 $S$-态射 $f : U \to Y$，使得图
$$
\xymatrix{
\mathcal{O}_{X, x}^\wedge \ar[r] &
\mathcal{O}_{U, u}^\wedge \\
\mathcal{O}_{Y, y} \ar[r]^{f^\sharp_u} \ar[u]^\varphi &
\mathcal{O}_{U, u} \ar[u]
}
$$
模 $\mathfrak m_u^N$ 交换。
\end{lemma}

\begin{proof}
该问题在 $X$ 上是局部的，故可假设 $X$、$Y$、$S$ 仿射。记
$S = \Spec(R)$、$X = \Spec(A)$、$Y = \Spec(B)$。写成
$B = R[x_1, \ldots, x_n]/(f_1, \ldots, f_m)$。设
$\mathfrak p \subset A$ 是对应于 $x$ 的素理想。由
代数进阶章命题
\ref{more-algebra-proposition-finite-type-over-G-ring}，局部环
$\mathcal{O}_{X, x} = A_\mathfrak p$ 是 G-环。映射 $\varphi$ 是映射
$$
B_\mathfrak q^\wedge \longrightarrow A_\mathfrak p^\wedge
$$
，其中 $\mathfrak q \subset B$ 是对应于 $y$ 的素理想。令
$a_1, \ldots, a_n \in A_\mathfrak p^\wedge$ 为 $x_1, \ldots, x_n$ 经
$R[x_1, \ldots, x_n] \to B \to B_\mathfrak q^\wedge \to A_\mathfrak p^\wedge$
所得的像。于是可以应用 环同态光滑化章引理
\ref{smoothing-lemma-approximation-property-variant}
，得到 \'etale 环映射 $A \to A'$、素理想 $\mathfrak p' \subset A'$ 以及
$b_1, \ldots, b_n \in A'$，使得
$\kappa(\mathfrak p) = \kappa(\mathfrak p')$,
$a_i - b_i \in (\mathfrak p')^N(A'_{\mathfrak p'})^\wedge$, 及
% P05-ERR-CH37-0072: Frozen source indexes f_j by j = 1,...,n although the presentation defines f_1,...,f_m.
$f_j(b_1, \ldots, b_n) = 0$ for $j = 1, \ldots, n$.
把 $x_i$ 的类映到 $b_i \in A'$，便确定一个 $R$-代数映射 $B \to A'$。
取 $U = \Spec(A') \to \Spec(B)$ 为态射 $f$，并取 $u = \mathfrak p'$，证明完成。
\end{proof}

\begin{lemma}
\label{more-morphisms-lemma-isomorphism-approximation}
设 $S$ 是局部 Noether 概形。设 $X$、$Y$ 是 $S$ 上局部有限型的概形。设
$x \in X$ 与 $y \in Y$ 位于同一点 $s \in S$ 之上。假设
$\mathcal{O}_{S, s}$ 是 G-环。假设有完备局部环之间的
$\mathcal{O}_{S, s}$-代数同构
$$
\varphi : \mathcal{O}_{Y, y}^\wedge \longrightarrow \mathcal{O}_{X, x}^\wedge
$$
。则对每个 $N \geq 1$，存在 $S$ 上带点概形的态射
$$
(X, x) \leftarrow (U, u) \rightarrow (Y, y)
$$
，使得两个箭头都定义初等 \'etale 邻域，且图
$$
\xymatrix{
& \mathcal{O}_{U, u}^\wedge \\
\mathcal{O}_{Y, y}^\wedge \ar[rr]^\varphi \ar[ru] & &
\mathcal{O}_{X, x}^\wedge \ar[lu]
}
$$
模 $\mathfrak m_u^N$ 交换。
\end{lemma}

\begin{proof}
可假设 $N \geq 2$。应用 引理 \ref{more-morphisms-lemma-map-approximation} 得到
$(U, u) \to (X, x)$ 与 $f : (U, u) \to (Y, y)$。我们断言 $f$ 在 $u$
处为 \'etale，这将完成证明。事实上，将证明诱导映射
$\mathcal{O}_{Y, y}^\wedge \to \mathcal{O}_{U, u}^\wedge$
是同构。证明这一点后，
引理 \ref{more-morphisms-lemma-lifting-along-artinian-at-point}
将表明 $f$ 在 $u$ 处光滑；当然 $f$ 在 $u$ 处也非分歧，故
态射章引理 \ref{morphisms-lemma-etale-smooth-unramified}
表明 $f$ 在 $u$ 处为 \'etale。对局部环 $(R, \mathfrak m)$，令
$\text{Gr}_\mathfrak m(R) =
\bigoplus_{n \geq 0} \mathfrak m^n/\mathfrak m^{n + 1}$.
为证明该断言，考察诱导的分次环图
$$
\xymatrix{
& \text{Gr}_{\mathfrak m_u}(\mathcal{O}_{U, u}) \\
\text{Gr}_{\mathfrak m_y}(\mathcal{O}_{Y, y}) \ar[rr]^\varphi \ar[ru] & &
\text{Gr}_{\mathfrak m_x}(\mathcal{O}_{X, x}) \ar[lu]
}
$$
由于 $N \geq 2$，且图中的分次代数由次数 $1$ 的部分生成，所以此图确实交换！
按假设，下方箭头是同构。由 代数进阶章引理
\ref{more-algebra-lemma-flat-unramified}（例如），映射
$\mathcal{O}_{X, x}^\wedge \to \mathcal{O}_{U, u}^\wedge$
是同构，因而图中西北方向的箭头是同构。于是 $f$ 诱导同构
% P05-ERR-CH37-0073: Frozen source gives f an X-based graded source and an O_U target graded by m_y; f has source Y and the target maximal ideal is m_u.
$\text{Gr}_{\mathfrak m_x}(\mathcal{O}_{X, x}) \to
\text{Gr}_{\mathfrak m_y}(\mathcal{O}_{U, u})$.
对两个局部环使用归纳法及短正合列
$$
0 \to \text{Gr}^n_{\mathfrak m}(R) \to R/\mathfrak m^{n + 1} \to
R/\mathfrak m^n \to 0
$$
，由蛇引理可得 $f$ 对所有 $n$ 均诱导同构
$\mathcal{O}_{Y, y}/\mathfrak m_y^n \to \mathcal{O}_{U, u}/\mathfrak m_u^n$
，这正是所要证明的。
\end{proof}

\begin{lemma}
\label{more-morphisms-lemma-relative-map-approximation-pre}
设 $X \to S$、$Y \to T$、$x$、$s$、$y$、$t$、$\sigma$、$y_\sigma$ 与
$\varphi$ 如下给定：有概形态射
$$
\vcenter{
\xymatrix{
X \ar[d] & Y \ar[d] \\
S & T
}
}
\quad\text{带标记点}\quad
\vcenter{
\xymatrix{
x \ar[d] & y \ar[d] \\
s & t
}
}
$$
其中 $S$ 局部 Noether，$T$ 在 $\mathbf{Z}$ 上有限型。态射 $X \to S$ 与
$Y \to T$ 局部有限型。局部环 $\mathcal{O}_{S, s}$ 是 G-环。映射
$$
\sigma : \mathcal{O}_{T, t} \longrightarrow \mathcal{O}_{S, s}^\wedge
$$
是局部同态。令
$Y_\sigma = Y \times_{T, \sigma} \Spec(\mathcal{O}_{S, s}^\wedge)$.
再设 $y_\sigma$ 是 $Y_\sigma$ 中的点，它映到 $y$ 以及
$\Spec(\mathcal{O}_{S, s}^\wedge)$ 的闭点。最后，
$$
\varphi :
\mathcal{O}_{X, x}^\wedge
\longrightarrow
\mathcal{O}_{Y_\sigma, y_\sigma}^\wedge
$$
是 $\mathcal{O}_{S, s}^\wedge$-代数的同构。此时存在交换图
$$
\xymatrix{
X \ar[d] &
W \ar[l] \ar[rd] \ar[rr] & &
Y \times_{T, \tau} V \ar[r] \ar[ld] & Y \ar[d] \\
S & &
V \ar[ll] \ar[rr]^\tau & &
T
}
$$
以及点 $w \in W$、$v \in V$，使得
\begin{enumerate}
\item $(V, v) \to (S, s)$ 是初等 \'etale 邻域；
\item $(W, w) \to (X, x)$ 是初等 \'etale 邻域；并且
\item $\tau(v) = t$。
\end{enumerate}
令 $y_\tau \in Y \times_T V$ 为在恒等式
$(Y_\sigma)_s = (Y \times_T V)_v$ 下对应于 $y_\sigma$ 的点。则
\begin{enumerate}
\item[(4)] $(W, w) \to (Y \times_{T, \tau} V, y_\tau)$ 是初等 \'etale 邻域。
\end{enumerate}
\end{lemma}

\begin{proof}
记 $X_\sigma = X \times_S \Spec(\mathcal{O}_{S, s}^\wedge)$，并以
$x_\sigma \in X_\sigma$ 表示位于 $x$ 上方的唯一点。由
代数进阶章命题
\ref{more-algebra-proposition-Noetherian-complete-G-ring}，
$\mathcal{O}_{S, s}^\wedge$ 是 G-环。由 Lemma
\ref{more-morphisms-lemma-isomorphism-approximation}，可以选择
$$
(X_\sigma, x_\sigma) \leftarrow (U, u) \rightarrow (Y_\sigma, y_\sigma)
$$
，其中两个箭头都是初等 \'etale 邻域。

\medskip\noindent
把 $S$ 替换为 $s$ 的一个开邻域后，可假设 $S = \Spec(R)$ 仿射。由于
$\mathcal{O}_{S, s}$ 是 G-环，由 环同态光滑化章定理
\ref{smoothing-theorem-popescu}，环 $\mathcal{O}_{S, s}^\wedge$ 是光滑
$R$-代数的滤过余极限。因此可写成
$$
\Spec(\mathcal{O}_{S, s}^\wedge) = \lim S_i
$$
，这是在 $S$ 上光滑的仿射概形 $S_i$ 的有向极限。以 $s_i \in S_i$ 表示
$\Spec(\mathcal{O}_{S, s}^\wedge)$ 的闭点的像。注意
$\kappa(s) = \kappa(s_i)$。令 $X_i = X \times_S S_i$，并以
$x_i \in X_i$ 表示映到 $x$ 的唯一点。注意 $\kappa(x) = \kappa(x_i)$。
由于 $T$ 在 $\mathbf{Z}$ 上有限型，由 极限章命题
\ref{limits-proposition-characterize-locally-finite-presentation}
，可选择一个 $i$ 以及带点概形的态射
$\sigma_i : (S_i, s_i) \to (T, t)$，使它与
$\Spec(\mathcal{O}_{S, s}^\wedge) \to S_i$ 的复合等于 $\sigma$。令
$Y_i = Y \times_T S_i$，并以 $y_i$ 表示 $y_\sigma$ 的像。注意
$\kappa(y_i) = \kappa(y_\sigma)$。由 极限章引理
\ref{limits-lemma-descend-finite-presentation}，可选择一个 $i$ 以及图
$$
\xymatrix{
X_i \ar[rd] &
U_i \ar[l] \ar[d] \ar[r] &
Y_i \ar[ld] \\
& S_i
}
$$
，它到 $\Spec(\mathcal{O}_{S, s}^\wedge)$ 的基变换恢复
$X_\sigma \leftarrow U \rightarrow Y_\sigma$。由 极限章引理
\ref{limits-lemma-descend-etale}，增大 $i$ 后可假设态射
$X_i \leftarrow U_i \rightarrow Y_i$ 为 \'etale。设 $u_i \in U_i$ 是 $u$ 的像。
由于 $u_i \mapsto x_i$，故
$\kappa(x) = \kappa(x_\sigma) = \kappa(u) \supset \kappa(u_i) \supset
\kappa(x_i) = \kappa(x)$，从而 $\kappa(u_i) = \kappa(x_i)$。因此
$(X_i, x_i) \leftarrow (U_i, u_i)$ 是初等 \'etale 邻域。又由于
$\kappa(y_i) = \kappa(y_\sigma) = \kappa(u)$，可知
$(U_i, u_i) \to (Y_i, y_i)$ 也是初等 \'etale 邻域。

\medskip\noindent
至此已构造图
$$
\xymatrix{
X \ar[d] &
X \times_S S_i \ar[l] \ar[rd] &
U_i \ar[l] \ar[r] \ar[d] &
Y \times_T S_i \ar[r] \ar[ld] &
Y \ar[d] \\
S & &
S_i \ar[ll] \ar[rr] & &
T
}
$$
，它与引理陈述中的图相同，但 $S_i \to S$ 是光滑的。由 Lemma
\ref{more-morphisms-lemma-slice-smooth}，并缩小 $S_i$ 后，可假设存在经过 $s_i$ 的闭子概形
$V \subset S_i$，使得 $V \to S$ 为 \'etale。令 $W$ 等于 $V$ 在 $U_i$
中的概形论逆像，即得结论。
\end{proof}

\noindent
我们强烈建议读者跳过本节余下部分。

\begin{lemma}
\label{more-morphisms-lemma-relative-map-approximation}
考虑图
$$
\vcenter{
\xymatrix{
X \ar[d] & Y \ar[d] \\
S & T \ar[l]
}
}
\quad\text{带标记点}\quad
\vcenter{
\xymatrix{
x \ar[d] & y \ar[d] \\
s & t \ar[l]
}
}
$$
% P05-ERR-CH37-0074: Frozen source has the ungrammatical phrase “where S be”.
，其中 $S$ 是局部 Noether 概形，且各态射局部有限型。假设
$\mathcal{O}_{S, s}$ 是 G-环。再假设给定局部 $\mathcal{O}_{S, s}$-代数映射
$$
\sigma : \mathcal{O}_{T, t} \longrightarrow \mathcal{O}_{S, s}^\wedge
$$
以及局部 $\mathcal{O}_{S, s}$-代数映射
$$
\varphi :
\mathcal{O}_{X, x}
\longrightarrow
\mathcal{O}_{Y_\sigma, y_\sigma}^\wedge
$$
，其中 $Y_\sigma = Y \times_{T, \sigma} \Spec(\mathcal{O}_{S, s}^\wedge)$，
而 $y_\sigma$ 是 $Y_\sigma$ 中位于 $y$ 上方的唯一点。对每个 $N \geq 1$，
存在 $S$ 上概形的交换图
$$
\xymatrix{
X \ar[d] & X \times_S V \ar[l] \ar[rd] &
W \ar[l]^-f \ar[r] \ar[d] &
Y \times_{T, \tau} V \ar[r] \ar[ld] & Y \ar[d] \\
S & & V \ar[ll] \ar[rr]^\tau & & T
}
$$
以及点 $w \in W$、$v \in V$，使得
\begin{enumerate}
\item $v \mapsto s$、$\tau(v) = t$、$f(w) = (x, v)$，且 $w \mapsto (y, v)$；
\item $(V, v) \to (S, s)$ 是初等 \'etale 邻域；
\item 图
$$
\xymatrix{
\mathcal{O}_{S, s}^\wedge \ar[r] & \mathcal{O}_{V, v}^\wedge \\
\mathcal{O}_{T, t} \ar[r]^{\tau^\sharp_v} \ar[u]_\sigma &
\mathcal{O}_{V, v} \ar[u]
}
$$
% P05-ERR-CH37-0075: Frozen source says “commutes module” rather than “commutes modulo”.
模 $\mathfrak m_v^N$ 交换；
\item $(W, w) \to (Y \times_{T, \tau} V, (y, v))$ 是初等 \'etale 邻域；
\item 图
$$
\xymatrix{
\mathcal{O}_{X, x} \ar[r]_\varphi &
\mathcal{O}_{Y_\sigma, y_\sigma}^\wedge \ar[r] &
\mathcal{O}_{Y_\sigma, y_\sigma}/\mathfrak m_{y_\sigma}^N \ar@{=}[r] &
\mathcal{O}_{Y \times_{T, \tau} V, (y, v)}/\mathfrak m_{(y, v)}^N
\ar[d]_{\cong} \\
\mathcal{O}_{X, x} \ar[r] \ar@{=}[u] &
\mathcal{O}_{X \times_S V, (x, v)} \ar[r]^{f^\sharp_w} &
\mathcal{O}_{W, w} \ar[r] &
\mathcal{O}_{W, w}/\mathfrak m_w^N
}
$$
交换。该等号源于以下事实：由 (2) 与 (3)，$Y_\sigma$ 与
$Y \times_{T, \tau} V$ 在
$\mathcal{O}_{V, v}/\mathfrak m_v^N = \mathcal{O}_{S, s}/\mathfrak m_s^N$
上典范同构。
\end{enumerate}
\end{lemma}

\begin{proof}
把 $X$、$S$、$T$、$Y$ 替换为仿射开子概形后，可假设引理陈述中的图来自对下图
施行 $\Spec$：
$$
\vcenter{
\xymatrix{
A & B \\
R \ar[u] \ar[r] & C \ar[u]
}
}
\quad\text{with primes}\quad
\vcenter{
\xymatrix{
\mathfrak p_A & \mathfrak p_B \\
\mathfrak p_R \ar@{-}[u] \ar@{-}[r] & \mathfrak p_C \ar@{-}[u]
}
}
$$
，其中各环均 Noether，各环映射均有限型。在本证明中，每个环 $E$ 都是 Noether
$R$-代数，并带有位于 $\mathfrak p_R$ 上方的素理想 $\mathfrak p_E$；所有环映射
都是与给定素理想相容的 $R$-代数映射。此外，写 $E^\wedge$ 时，是指 $E$ 在
$\mathfrak p_E$ 处局部化后的完备化。我们还将不再另作说明地使用如下事实：若
\'etale 环映射 $E_1 \to E_2$ 满足
$\kappa(\mathfrak p_{E_1}) = \kappa(\mathfrak p_{E_2})$，则由
代数进阶章引理 \ref{more-algebra-lemma-flat-unramified}，它诱导同构
$E_1^\wedge = E_2^\wedge$。

\medskip\noindent
在此记号下，$\sigma$ 与 $\varphi$ 对应于环映射
$$
\sigma : C \to R^\wedge
\quad\text{且}\quad
\varphi : A \longrightarrow (B \otimes_{C, \sigma} R^\wedge)^\wedge
$$
。图示如下：
$$
\xymatrix{
A \ar@/^1em/[rrr]^\varphi &
B \ar[r] &
B \otimes_{C, \sigma} R^\wedge \ar[r] &
(B \otimes_{C, \sigma} R^\wedge)^\wedge \\
R \ar[r] \ar[u] &
C \ar[r]^\sigma \ar[u] & R^\wedge \ar[u] \ar[ru]
}
$$
由
代数进阶章命题
\ref{more-algebra-proposition-Noetherian-complete-G-ring}，$R^\wedge$ 是 G-环。因此由
代数进阶章命题
\ref{more-algebra-proposition-finite-type-over-G-ring}，
$B \otimes_{C, \sigma} R^\wedge$ 是 G-环。由 Lemma
\ref{more-morphisms-lemma-map-approximation}（转写成代数形式），存在 \'etale 环映射
$B \otimes_{C, \sigma} R^\wedge \to B'$，它诱导同构
$\kappa(\mathfrak p_{B \otimes_{C, \sigma} R^\wedge})
\to \kappa(\mathfrak p_{B'})$
；还存在 $R$-代数映射 $A \to B'$，使得复合
$$
A \to B' \to (B')^\wedge = (B \otimes_{C, \sigma} R^\wedge)^\wedge
$$
模 $(\mathfrak p_{(B \otimes_{C, \sigma} R^\wedge)^\wedge})^N$ 与
$\varphi$ 相同。因此可以用这一复合替换 $\varphi$，因为 $\varphi$ 在结论中
只通过引理陈述第 (5) 项的交换性要求出现。图示为：
$$
\xymatrix{
& & B' \ar[r] & (B')^\wedge \ar@{=}[d] \\
A \ar[rru] &
B \ar[r] &
B \otimes_{C, \sigma} R^\wedge \ar[r] \ar[u] &
(B \otimes_{C, \sigma} R^\wedge)^\wedge \\
R \ar[r] \ar[u] &
C \ar[r]^\sigma \ar[u] & R^\wedge \ar[u] \ar[ru]
}
$$
接着使用如下事实：由于按假设 $R_{\mathfrak p_R}$ 是 G-环，$R^\wedge$ 是光滑
$R$-代数的滤过余极限（环同态光滑化章定理
\ref{smoothing-theorem-popescu}）。由于 $C$ 在 $R$ 上有限表现，得到分解
$$
C \to R' \to R^\wedge
$$
，其中某个 $R \to R'$ 光滑，见
代数章引理 \ref{algebra-lemma-characterize-finite-presentation}.
增大 $R'$ 后，可假设存在 \'etale 的 $B \otimes_C R'$-代数 $B''$，其到
$B \otimes_{C, \sigma} R^\wedge$ 的基变换是 $B'$，见
代数章引理 \ref{algebra-lemma-etale}.
于是 $B'$ 是这些 $B''$ 的滤过余极限，从而增大 $R'$ 后，可假设存在
$R$-代数映射 $A \to B''$，使得 $A \to B'' \to B'$ 是先前构造的映射
（仍用上述引理）。图示为
$$
\xymatrix{
& & B'' \ar[r] & B' \ar[r] & (B')^\wedge \ar@{=}[d] \\
A \ar[rru] &
B \ar[r] &
B \otimes_C R' \ar[r] \ar[u] &
B \otimes_{C, \sigma} R^\wedge \ar[r] \ar[u] &
(B \otimes_{C, \sigma} R^\wedge)^\wedge \\
R \ar[r] \ar[u] &
C \ar[r] \ar[u] &
R' \ar[r] \ar[u] &
R^\wedge \ar[u] \ar[ru]
}
$$
并且
$$
B' = B'' \otimes_{(B \otimes_C R')} (B \otimes_{C, \sigma} R^\wedge)
$$
。这意味着可用 $R'$ 替换 $C$，用 $R' \to R^\wedge$ 替换
$\sigma : C \to R^\wedge$，并用 $B''$ 替换 $B$，从而化简为图
$$
\xymatrix{
A \ar[r] &
B \ar[r] &
B \otimes_{C, \sigma} R^\wedge \\
R \ar[r] \ar[u] &
C \ar[r]^\sigma \ar[u] & R^\wedge \ar[u]
}
$$
，其中 $\varphi$ 等于各水平箭头的复合再接上从
$B \otimes_{C, \sigma} R^\wedge$ 到其完备化的典范映射。
证明的最后一步，是把 引理 \ref{more-morphisms-lemma-map-approximation}（或其证明）再应用一次于
$\Spec(R)$ 上的 $\Spec(C)$ 和 $\Spec(R)$，以及映射 $C \to R^\wedge$。
该引理给出环映射 $C \to D$，使得 $R \to D$ 为 \'etale，
$\kappa(\mathfrak p_R) = \kappa(\mathfrak p_D)$，并且
$$
C \to D \to D^\wedge = R^\wedge
$$
模 $(\mathfrak p_{R^\wedge})^N$ 等于 $\sigma : C \to R^\wedge$。于是可取
$$
V = \Spec(D)
\quad\text{且}\quad
W = \Spec(B \otimes_C D)
$$
作为引理所求。事实上，图
$$
\xymatrix{
A \ar[r] &
B \otimes_{C, \sigma} R^\wedge \ar[r] &
B \otimes_{C, \sigma} R^\wedge/(\mathfrak p_{R^\wedge})^N \ar@{=}[r] &
B \otimes_C D/(\mathfrak p_D)^N \\
A \ar@{=}[u] \ar[r] &
A \otimes_R D \ar[r] &
B \otimes_R D \ar[r] &
B \otimes_C D/(\mathfrak p_D)^N \ar@{=}[u]
}
$$
交换，因为 $C \to D \to D^\wedge = R^\wedge$ 模
$(\mathfrak p_{R^\wedge})^N$ 等于 $\sigma$。这证明了第 (5) 项，其余性质由构造立即得到。
\end{proof}

\begin{lemma}
\label{more-morphisms-lemma-control-agreement}
% P05-ERR-CH37-0076: Frozen source uses plural “morphisms” for the single map T -> S.
设 $T \to S$ 是 Noether 概形的有限型态射。设 $t \in T$ 映到 $s \in S$，并设
$\sigma : \mathcal{O}_{T, t} \to \mathcal{O}_{S, s}^\wedge$
是局部 $\mathcal{O}_{S, s}$-代数映射。对每个 $N \geq 1$，存在有限型态射
$(T', t') \to (T, t)$，使得 $\sigma$ 经
$\mathcal{O}_{T, t} \to \mathcal{O}_{T', t'}$
分解，并且对每个经 $\mathcal{O}_{T, t} \to \mathcal{O}_{T', t'}$ 分解的局部
$\mathcal{O}_{S, s}$-代数映射
$\sigma' : \mathcal{O}_{T, t} \to \mathcal{O}_{S, s}^\wedge$
，映射 $\sigma$ 与 $\sigma'$ 模 $\mathfrak m_s^N$ 相同。
\end{lemma}

\begin{proof}
可假设 $S$ 与 $T$ 仿射。记 $S = \Spec(R)$ 且 $T = \Spec(C)$。设
$c_1, \ldots, c_n \in C$ 是 $C$ 作为 $R$-代数的一组生成元。设
$\mathfrak p \subset R$ 是对应于 $s$ 的素理想。记
$\mathfrak p = (f_1, \ldots, f_m)$。把 $R$ 替换为一个主局部化（以消去
$R_\mathfrak p$ 中的分母）后，可假设存在
$r_1, \ldots, r_n \in R$ 与 $a_{i, I} \in \mathcal{O}_{S, s}^\wedge$，
其中 $I = (i_1, \ldots, i_m)$ 且 $\sum i_j = N$，使得
$$
\sigma(c_i) = r_i + \sum\nolimits_I a_{i, I} f_1^{i_1} \ldots f_m^{i_m}
$$
在 $\mathcal{O}_{S, s}^\wedge$ 中成立。于是考虑
$$
C' = C[t_{i, I}]/
\left(c_i - r_i - \sum\nolimits_I t_{i, I} f_1^{i_1} \ldots f_m^{i_m}\right)
$$
，取 $\mathfrak p' = \mathfrak pC' + (t_{i, I})$；把 $t_{i, I}$ 映到
$a_{i, I}$，便给出 $\sigma : C \to \mathcal{O}_{S, s}^\wedge$ 经 $C'$ 的分解。
取 $T' = \Spec(C')$ 即可，因为引理陈述中的任何 $\sigma'$ 都把 $c_i$ 模极大理想的
$N$ 次幂映到 $r_i$。
\end{proof}

\begin{lemma}
\label{more-morphisms-lemma-control-graded}
设 $Y \to T \to S$ 是 Noether 概形的有限型态射。设 $t \in T$ 映到
$s \in S$，并设
$\sigma : \mathcal{O}_{T, t} \to \mathcal{O}_{S, s}^\wedge$
是局部 $\mathcal{O}_{S, s}$-代数映射。存在有限型态射
$(T', t') \to (T, t)$，使得 $\sigma$ 经
$\mathcal{O}_{T, t} \to \mathcal{O}_{T', t'}$
分解，并且对每个经 $\mathcal{O}_{T, t} \to \mathcal{O}_{T', t'}$ 分解的局部
$\mathcal{O}_{S, s}$-代数映射
$\sigma' : \mathcal{O}_{T, t} \to \mathcal{O}_{S, s}^\wedge$
，闭浸入
$$
Y \times_{T, \sigma} \Spec(\mathcal{O}_{S, s}^\wedge) = Y_\sigma
\longleftarrow Y_t \longrightarrow
Y_{\sigma'} =
Y \times_{T, \sigma'} \Spec(\mathcal{O}_{S, s}^\wedge)
$$
具有同构的余法代数。
\end{lemma}

\begin{proof}
一个有用的观察是，$\sigma$ 的存在蕴含 $\kappa(s) = \kappa(t)$。陈述确有意义，
因为 $Y_\sigma$ 与 $Y_{\sigma'}$ 在
$s \in \Spec(\mathcal{O}_{S, s}^\wedge)$ 上的纤维都典范同构于 $Y_t$。
我们把“$\sigma'$ 经 $\mathcal{O}_{T, t} \to \mathcal{O}_{T', t'}$ 分解”这一性质
视为对 $\sigma'$ 的约束。若有若干这样的约束，例如由
$(T'_i, t'_i) \to (T, t)$、$i = 1, \ldots, n$ 给出，则可考虑
% P05-ERR-CH37-0077: Frozen source has the ungrammatical phrase “can combined”.
$(T'_1 \times_T \ldots \times_T T'_n, (t'_1, \ldots, t'_n)) \to
(T, t)$ 把它们合并。下文将不再另作说明地使用这一点。

\medskip\noindent
由 引理 \ref{more-morphisms-lemma-control-agreement}，可假设引理陈述中的任何 $\sigma'$ 模
$\mathfrak m_s^2$ 都与 $\sigma$ 相同。注意，$Y_t$ 在 $Y_\sigma$ 中的余法代数
正是拟凝聚分次 $\mathcal{O}_{Y_t}$-代数
$$
\bigoplus\nolimits_{n \geq 0}
\mathfrak m_s^n\mathcal{O}_{Y_\sigma}/
\mathfrak m_s^{n + 1}\mathcal{O}_{Y_\sigma}
$$
；$Y_{\sigma'}$ 的情形类似。由于 $\sigma$ 与 $\sigma'$ 模
$\mathfrak m_s^2$ 相同，这两个代数在次数 $0$ 与 $1$ 上相同。另一方面，这些余法代数
在次数 $0$ 上由次数 $1$ 的部分生成。因此，若存在延拓刚在次数 $0$ 与 $1$ 上构造的
同构的同构，则它是唯一的。

\medskip\noindent
可假设 $S$ 与 $T$ 仿射。设 $Y = Y_1 \cup \ldots \cup Y_n$ 是仿射开覆盖。
若能按引理构造 $(T_i', t'_i) \to (T, t)$，使得对 $Y_i \to T \to S$ 与
$\sigma$ 存在所需同构（见上一段），则由唯一性，这些同构可以黏合，从而证明
$Y \to T$ 的结果。因此可假设 $Y$ 仿射。

\medskip\noindent
写成 $S = \Spec(R)$、$T = \Spec(C)$ 与 $Y = \Spec(B)$。选择表现
$B = C[x_1, \ldots, x_n]/(f_1, \ldots, f_m)$。记
$R^\wedge = \mathcal{O}_{S, s}^\wedge$。取多项式
$a_{kj} \in R^\wedge[x_1, \ldots, x_n]$，使得
$$
\sum\nolimits_{j = 1, \ldots, m} a_{kj}\sigma(f_j) = 0,\quad
\text{对 }k = 1, \ldots, K
$$
给出 $\sigma(f_j) \in R^\wedge[x_1, \ldots, x_n]$ 之间的关系模的一组生成元。
于是有正合列
\begin{equation}
\label{more-morphisms-equation-resolution}
R^\wedge[x_1, \ldots, x_n]^{\oplus K} \to
R^\wedge[x_1, \ldots, x_n]^{\oplus m} \to
R^\wedge[x_1, \ldots, x_n] \to B \otimes_{C, \sigma} R^\wedge \to 0
\end{equation}
按
代数进阶章第 \ref{more-algebra-section-artin-rees} 节.
中的定义，取整数 $c$，使得 Artin--Rees 引理对该序列的第一、第二个映射以及理想
$\mathfrak m_{R^\wedge}R^\wedge[x_1, \ldots, x_n]$ 均适用。
用多重指标记号写成
$$
a_{kj} = \sum\nolimits_{I \in \Omega} a_{kj, I} x^I
\quad\text{且}\quad
f_j = \sum\nolimits_{I \in \Omega} f_{j, I} x^I
$$
，其中 $a_{kj, I} \in R^\wedge$、$f_{j, I} \in C$，而 $\Omega$ 是有限多重指标集。
% P05-ERR-CH37-0079: Frozen source omits “is” in “and Omega a finite set of multiindices”.
于是可知
$$
\sum\nolimits_{j = 1, \ldots, m,\ I, I' \in \Omega,\ I + I' = I''}
a_{kj, I} \sigma(f_{j, I'}) = 0,\quad
I''\text{ a multiindex}
$$
在 $R^\wedge$ 中成立。因此取
$$
% P05-ERR-CH37-0078: Frozen source defines C' with t_{jk,I}, uses t_{kj,I} in the relations, and then maps t_{kj,I} to a_{jk,I}; the k,j indices are inconsistent.
C' = C[t_{jk, I}]/
\left(
\sum\nolimits_{j = 1, \ldots, m,\ I, I' \in \Omega,\ I + I' = I''}
t_{kj, I} f_{j, I'},\ I''\text{ a multiindex}\right)
$$
。于是 $\sigma$ 经映射 $\tilde\sigma : C' \to R^\wedge$ 分解，该映射把
$t_{kj, I}$ 映到 $a_{jk, I}$。因此 $T' = \Spec(C')$ 带有点 $t' \in T'$，
使得 $\sigma$ 经 $\mathcal{O}_{T, t} \to \mathcal{O}_{T', t'}$ 分解。在
$C'[x_1, \ldots, x_n]$ 中令 $t_{kj} = \sum t_{kj, I} x^I$。于是有复形
\begin{equation}
\label{more-morphisms-equation-resolution-new}
C'[x_1, \ldots, x_n]^{\oplus K} \to
C'[x_1, \ldots, x_n]^{\oplus m} \to
C'[x_1, \ldots, x_n] \to B \otimes_C C' \to 0
\end{equation}
，它在 $C'[x_1, \ldots, x_n]$ 处正合，且经 $\tilde\sigma$ 的基变换给出
(\ref{more-morphisms-equation-resolution})。

\medskip\noindent
由 引理 \ref{more-morphisms-lemma-control-agreement}，可找到进一步的态射
$(T'', t'') \to (T', t')$，使得 $\tilde\sigma$ 经
$\mathcal{O}_{T', t'} \to \mathcal{O}_{T'', t''}$ 分解，并且若
$\sigma' : C \to R^\wedge$ 经 $\mathcal{O}_{T'', t''}$ 分解，则诱导映射
$\tilde \sigma' : C' \to R^\wedge$ 模 $\mathfrak m_s^{c + 1}$ 与
$\tilde \sigma$ 相同。因此，若 $\sigma'$ 是这样的映射，就得到复形
$$
R^\wedge[x_1, \ldots, x_n]^{\oplus K} \to
R^\wedge[x_1, \ldots, x_n]^{\oplus m} \to
R^\wedge[x_1, \ldots, x_n] \to B \otimes_{C, \sigma'} R^\wedge \to 0
$$
：这是把 $\tilde\sigma'$ 应用于多项式 $t_{kj}$ 与 $f_j$ 而得到的
$R^\wedge[x_1, \ldots, x_n]$ 上复形。换言之，它是复形
(\ref{more-morphisms-equation-resolution-new}) 经 $\tilde\sigma'$ 的基变换。由于
$\tilde \sigma$ 与 $\tilde \sigma'$ 模 $\mathfrak m_s^{c + 1}$ 同余，定义此复形的
矩阵与定义复形 (\ref{more-morphisms-equation-resolution}) 的矩阵模
$\mathfrak m_s^{c + 1}$ 同余。由于 (\ref{more-morphisms-equation-resolution}) 正合，可以应用
代数进阶章引理 \ref{more-algebra-lemma-approximate-complex-graded}
得出
$$
\text{Gr}_{\mathfrak m_s}(B \otimes_{C, \sigma'} R^\wedge)
\cong
\text{Gr}_{\mathfrak m_s}(B \otimes_{C, \sigma} R^\wedge)
$$
，如所求。
\end{proof}

\begin{lemma}
\label{more-morphisms-lemma-relative-isomorphism-approximation}
% P05-ERR-CH37-0080: Frozen source has “notation an assumptions” instead of “notation and assumptions”.
沿用 引理 \ref{more-morphisms-lemma-relative-map-approximation} 的记号与假设，并假设 $\varphi$
在完备化上诱导同构。则可以选择图，使得 $f$ 为 \'etale。
\end{lemma}

\begin{proof}
可假设 $N \geq 2$，并可按 引理 \ref{more-morphisms-lemma-control-graded} 用 $(T', t')$ 替换
$(T, t)$。由于 $(V, v) \to (S, s)$ 是初等 \'etale 邻域，
$(X \times_S V, (x, v)) \to (X, x)$ 也是如此。因此，由
代数进阶章引理 \ref{more-algebra-lemma-flat-unramified}.
，$\mathcal{O}_{X, x} \to \mathcal{O}_{X \times_S V, (x, v)}$ 在完备化上诱导同构。
我们断言 $\mathcal{O}_{X, x} \to \mathcal{O}_{W, w}$ 在完备化上诱导同构。
证明这一点后，
引理 \ref{more-morphisms-lemma-lifting-along-artinian-at-point}
将表明 $f$ 在 $w$ 处光滑；当然，
% P05-ERR-CH37-0081: Frozen source says f is unramified at u although the distinguished point in this lemma is w.
$f$ 在 $u$ 处也非分歧，故
态射章引理 \ref{morphisms-lemma-etale-smooth-unramified}
表明 $f$ 在 $w$ 处为 \'etale。

\medskip\noindent
首先使用
引理 \ref{more-morphisms-lemma-relative-map-approximation}
第 (5) 项的交换性，看到对 $i \leq N$ 有交换图
$$
\xymatrix{
\text{Gr}^i_{\mathfrak m_x}(\mathcal{O}_{X, x})
\ar[r]_-\varphi &
\text{Gr}^i_{\mathfrak m_{y_\sigma}}(\mathcal{O}_{Y_\sigma, y_\sigma}^\wedge)
\ar@{=}[r] &
\text{Gr}^i_{\mathfrak m_{(y, v)}}(\mathcal{O}_{Y \times_{T, \tau} V, (y, v)})
\ar[d]_{\cong} \\
\text{Gr}^i_{\mathfrak m_x}(\mathcal{O}_{X, x})
\ar[r]^-{\cong} \ar@{=}[u] &
\text{Gr}^i_{\mathfrak m_{(x, v)}}(\mathcal{O}_{X \times_S V, (x, v)})
\ar[r]^{f^\sharp_w} &
\text{Gr}^i_{\mathfrak m_w}(\mathcal{O}_{W, w})
}
$$
这蕴含 $f^\sharp_w$ 在剩余域上给出同构 $\kappa(x) \to \kappa(w)$，在余切空间上
给出同构 $\mathfrak m_x/\mathfrak m_x^2 \to \mathfrak m_w/\mathfrak m_w^2$。
因此 $f^\sharp_w$ 在完备局部环上给出满射
$\mathcal{O}_{X, x}^\wedge \to \mathcal{O}_{W, w}^\wedge$
。

\medskip\noindent
由 引理 \ref{more-morphisms-lemma-control-graded}，有
$\text{Gr}_{\mathfrak m_s}(\mathcal{O}_{(Y \times_{T, \tau} V, (y, v)})$
与
$\text{Gr}_{\mathfrak m_s}(\mathcal{O}_{Y_\sigma, y_\sigma})$.
之间的同构。这可通过在点 $y$ 处取余法层同构的茎得到。由于所涉局部环 Noether，
关于 $\mathfrak m_s$ 取相伴分次与完备化交换，因为关于理想的完备化是 Noether 环上
有限模范畴的正合函子。这给出下列分次环图中的右侧竖直同构：
$$
\xymatrix{
\text{Gr}_{\mathfrak m_s}(\mathcal{O}_{W, w}^\wedge) &
% P05-ERR-CH37-0082: Frozen source omits a closing parenthesis in the point subscript of the upper-right local ring.
\text{Gr}_{\mathfrak m_s}
(\mathcal{O}_{(Y \times_{T, \tau} V, (y, v)}^\wedge) \ar[l] \\
\text{Gr}_{\mathfrak m_s}(\mathcal{O}_{X, x}^\wedge)
\ar[r]^\varphi \ar[u] &
\text{Gr}_{\mathfrak m_s}(\mathcal{O}_{Y_\sigma, y_\sigma}^\wedge)
\ar[u]_{\cong}
}
$$
我们不声称此图交换。由上一段的结果，左侧箭头满射；另外三个箭头是同构。因此左侧
箭头是两个同构 Noether 环之间的满射。故由 代数章引理
\ref{algebra-lemma-surjective-endo-noetherian-ring-is-iso}，它是同构
（也可用 Hilbert 函数直接论证）。特别地，
$\mathcal{O}_{X, x}^\wedge \to \mathcal{O}_{W, w}^\wedge$ 既单射又满射，证明完成。
\end{proof}















\section{有限自由覆盖在局部上支配 \'etale 覆盖}
\label{more-morphisms-section-finite-free-over-etale}

\noindent
本节说明一个大意如下的结果：概形 $S$ 的 \'etale 覆盖可由 $S$ 的有限局部自由覆盖的
Zariski 覆盖加细。

\begin{lemma}
\label{more-morphisms-lemma-dominate-etale-neighbourhood-finite-flat}
设 $S$ 是概形，$s \in S$。设 $f : (U, u) \to (S, s)$ 是 \'etale 邻域。
则存在仿射开邻域 $s \in V \subset S$ 以及满的有限局部自由态射
$\pi : T \to V$，使得对每个 $t \in \pi^{-1}(s)$，都存在开邻域
$t \in W_t \subset T$ 及交换图
$$
\xymatrix{
T \ar[d]^\pi & W_t \ar[l] \ar[rr]_{h_t} \ar[rd] & & U \ar[dl] \\
V \ar[rr] & & S
}
$$
，其中 $h_t(t) = u$。
\end{lemma}

\begin{proof}
该问题在 $S$ 上是局部的，故可把 $S$ 替换为 $s$ 的任意开邻域。也可把 $U$ 替换为
$u$ 的一个开邻域。因此，由 态射章引理
\ref{morphisms-lemma-etale-locally-standard-etale}，可假设 $U \to S$ 是仿射概形间的
标准 \'etale 态射。在此情形，引理（取 $V = S$）由
代数章引理 \ref{algebra-lemma-standard-etale-finite-flat-Zariski}.
得到。
\end{proof}

\begin{lemma}
\label{more-morphisms-lemma-dominate-etale-affine-finite-flat}
设 $f : U \to S$ 是仿射概形间的满 \'etale 态射。则存在满的有限局部自由态射
$\pi : T \to S$ 以及有限开覆盖 $T = T_1 \cup \ldots \cup T_n$，使得每个
$T_i \to S$ 都经 $U \to S$ 分解。图示为：
$$
\xymatrix{
& \coprod T_i  \ar[rd] \ar[ld] & \\
T \ar[rd]^\pi & & U \ar[ld]_f \\
& S &
}
$$
其中西南方向的箭头是 Zariski 覆盖。
\end{lemma}

\begin{proof}
这正是
代数章引理 \ref{algebra-lemma-etale-finite-flat-zariski}.
的改述。
\end{proof}

\begin{remark}
\label{more-morphisms-remark-topologies}
用拓扑的语言，引理 \ref{more-morphisms-lemma-dominate-etale-neighbourhood-finite-flat} 与
\ref{more-morphisms-lemma-dominate-etale-affine-finite-flat} 表示如下结论。设 $S$ 是任意概形，
$\{f_i : U_i \to S\}$ 是 $S$ 的 \'etale 覆盖。存在 Zariski 开覆盖
$S = \bigcup V_j$；对每个 $j$，存在有限局部自由满态射 $W_j \to V_j$；并且对每个
$j$，存在 Zariski 开覆盖 $\{W_{j, k} \to W_j\}$，使得族
$\{W_{j, k} \to S\}$ 加细给定的 \'etale 覆盖
$\{f_i : U_i \to S\}$。这在实践中意味着什么？例如，假设有一个下降问题，我们知道
如何对 Zariski 覆盖以及形如 $\{\pi : T \to S\}$ 的 fppf 覆盖解决它，其中
$\pi$ 有限局部自由且满。那么，这个下降问题对 \'etale 覆盖也有肯定答案。Gabber
在证明仿射概形 $X$ 满足 $\text{Br}(X) = \text{Br}'(X)$ 时使用了这一技巧，见
\cite{Hoobler}。
\end{remark}




\section{拟有限态射的 \'etale 局部化}
\label{more-morphisms-section-etale-localization}

\noindent
现在给出围绕“拟有限态射经 \'etale 局部化后变为有限态射”这一主题的一系列引理。
一般思路如下。设给定概形态射 $f : X \to S$ 以及点 $s \in S$。设
$\varphi : (U, u) \to (S, s)$ 是 $S$ 中 $s$ 的 \'etale 邻域。考虑纤维积
$X_U = U \times_S X$ 与基本图
\begin{equation}
\label{more-morphisms-equation-basic-diagram}
\vcenter{
\xymatrix{
V \ar[r] \ar[dr] & X_U \ar[d] \ar[r] & X \ar[d]^f \\
& U \ar[r]^\varphi & S
}
}
\end{equation}
，其中 $V \subset X_U$ 是开集。对态射 $f_U : X_U \to U$，或者对适当开集 $V$
的态射 $V \to U$，是否存在某种标准模型？一般而言答案当然是否定的，但对拟有限态射
可以得到一些结论。

\begin{lemma}
\label{more-morphisms-lemma-etale-makes-quasi-finite-finite-at-point}
设 $f : X \to S$ 是概形态射。设 $x \in X$，令 $s = f(x)$。假设
\begin{enumerate}
\item $f$ 局部有限型；并且
\item $x \in X_s$ 是孤立点\footnote{在 (1) 成立时，这意味着 $f$ 在 $x$ 处拟有限，
见 态射章引理 \ref{morphisms-lemma-quasi-finite-at-point-characterize}。}。
\end{enumerate}
则存在
\begin{enumerate}
\item[(a)] 初等 \'etale 邻域 $(U, u) \to (S, s)$；
\item[(b)] 开子概形 $V \subset X_U$（见 \ref{more-morphisms-equation-basic-diagram}）
\end{enumerate}
，使得
\begin{enumerate}
\item[(\romannumeral1)] $V \to U$ 是有限态射；
\item[(\romannumeral2)] $V$ 中恰有一个映到 $U$ 中 $u$ 的点 $v$；并且
\item[(\romannumeral3)] 点 $v$ 在态射 $X_U \to X$ 下映到 $x$，并诱导
$\kappa(x) = \kappa(v)$。
\end{enumerate}
此外，对任意初等 \'etale 邻域 $(U', u') \to (U, u)$，令
$V' = U' \times_U V \subset X_{U'}$，则三元组 $(U', u', V')$ 仍满足性质
(\romannumeral1)、(\romannumeral2) 与 (\romannumeral3)。
\end{lemma}

\begin{proof}
取仿射开集 $Y \subset X$、$W \subset S$，使得 $f(Y) \subset W$ 且
$x \in Y$。注意，$x$ 也是态射 $f|_Y : Y \to W$ 的纤维的孤立点。若能对
$f|_Y : Y \to W$ 证明该定理，则可得到 $f$ 的结论。因此归约到 $f$ 是仿射概形间
态射的情形；这正是
代数章引理 \ref{algebra-lemma-etale-makes-quasi-finite-finite-one-prime}.
。
\end{proof}

\noindent
在上一引理和下一引理中，我们不假设态射 $f$ 分离。这意味着其中构造的开集
$V$、$V_i$ 未必在 $X_U$ 中闭。此外，若选择仿射邻域 $U$，则每个 $V_i$ 都仿射，
但在非分离情形，交 $V_i \cap V_j$ 未必仿射。

\begin{lemma}
\label{more-morphisms-lemma-etale-makes-quasi-finite-finite-multiple-points}
设 $f : X \to S$ 是概形态射。设 $x_1, \ldots, x_n \in X$ 是在 $S$ 中具有
同一像 $s$ 的点。假设
\begin{enumerate}
\item $f$ 局部有限型；并且
\item 对 $i = 1, \ldots, n$，$x_i \in X_s$ 是孤立点。
\end{enumerate}
则存在
\begin{enumerate}
\item[(a)] 初等 \'etale 邻域 $(U, u) \to (S, s)$；
\item[(b)] 对每个 $i$，开子概形 $V_i \subset X_U$，
\end{enumerate}
使得对每个 $i$ 均有
\begin{enumerate}
\item[(\romannumeral1)] $V_i \to U$ 是有限态射；
\item[(\romannumeral2)] $V_i$ 中恰有一个映到 $U$ 中 $u$ 的点 $v_i$；并且
\item[(\romannumeral3)] 点 $v_i$ 映到 $X$ 中的 $x_i$，且
$\kappa(x_i) = \kappa(v_i)$。
\end{enumerate}
\end{lemma}

\begin{proof}
对 $n$ 使用归纳法。具体而言，假设 $(U, u) \to (S, s)$ 与
$V_i \subset X_U$（$i = 1, \ldots, n - 1$）对 $x_1, \ldots, x_{n - 1}$
适用。由于 $\kappa(s) = \kappa(u)$，纤维 $(X_U)_u = X_s$。因此存在唯一的点
$x'_n \in X_u \subset X_U$ 对应于 $x_n \in X_s$。此外，$x'_n$ 在 $X_u$
中是孤立的。故由 引理 \ref{more-morphisms-lemma-etale-makes-quasi-finite-finite-at-point}，
存在初等 \'etale 邻域 $(U', u') \to (U, u)$ 以及对 $x'_n$、因而对 $x_n$
适用的开集 $V_n \subset X_{U'}$。由
引理 \ref{more-morphisms-lemma-etale-makes-quasi-finite-finite-at-point}
的最后一个断言，对 $i = 1, \ldots, n - 1$，开子概形
$V'_i = U'\times_U V_i$ 对 $x_1, \ldots, x_{n - 1}$ 仍适用。结论得证。
\end{proof}

\noindent
若允许非平凡域扩张 $\kappa(u)/\kappa(s)$，即允许一般的 \'etale 邻域，则可如下分裂这些点。

\begin{lemma}
\label{more-morphisms-lemma-etale-makes-quasi-finite-finite-multiple-points-var}
设 $f : X \to S$ 是概形态射。设 $x_1, \ldots, x_n \in X$ 是在 $S$ 中具有
同一像 $s$ 的点。假设
\begin{enumerate}
\item $f$ 局部有限型；并且
\item 对 $i = 1, \ldots, n$，$x_i \in X_s$ 是孤立点。
\end{enumerate}
则存在
\begin{enumerate}
\item[(a)] \'etale 邻域 $(U, u) \to (S, s)$；
\item[(b)] 对每个 $i$，整数 $m_i$ 及开子概形
$V_{i, j} \subset X_U$，$j = 1, \ldots, m_i$
\end{enumerate}
，使得
\begin{enumerate}
\item[(\romannumeral1)] 每个 $V_{i, j} \to U$ 都是有限态射；
\item[(\romannumeral2)] $V_{i, j}$ 中恰有一个映到 $U$ 中 $u$ 的点 $v_{i, j}$，且
$\kappa(u) \subset \kappa(v_{i, j})$ 是有限纯不可分扩张；
% P05-ERR-CH37-0083: Frozen source orders item (iv) before item (iii).
\item[(\romannumeral4)] 若 $v_{i, j} = v_{i', j'}$，则 $i = i'$ 且 $j = j'$；并且
\item[(\romannumeral3)] 各点 $v_{i, j}$ 映到 $X$ 中的 $x_i$，且
$(X_U)_u$ 中没有其他点映到 $x_i$。
\end{enumerate}
\end{lemma}

\begin{proof}
本证明是
代数章引理 \ref{algebra-lemma-etale-makes-quasi-finite-finite-variant}
证明的概形语言变体。由 态射章引理
\ref{morphisms-lemma-quasi-finite-at-point-characterize}，态射 $f$ 在每个点 $x_i$
处拟有限。因此对每个 $i$，$\kappa(s) \subset \kappa(x_i)$ 有限
（态射章引理 \ref{morphisms-lemma-residue-field-quasi-finite}）。
对每个 $i$，令 $\kappa(s) \subset L_i \subset \kappa(x_i)$ 为满足如下性质的子域：
$L_i/\kappa(s)$ 可分，而 $\kappa(x_i)/L_i$ 纯不可分。选取有限 Galois 扩张
$L/\kappa(s)$，使得对 $i = 1, \ldots, n$ 存在 $\kappa(s)$-嵌入
$L_i \to L$。选择 \'etale 邻域 $(U, u) \to (S, s)$，使得作为
$\kappa(s)$-扩张有 $L \cong \kappa(u)$（Lemma
\ref{more-morphisms-lemma-realize-prescribed-residue-field-extension-etale}）。

\medskip\noindent
设 $y_{i, j}$（$j = 1, \ldots, m_i$）为 $X_U$ 中位于 $x_i \in X$ 与
$u \in U$ 上方的点。由
概形章引理 \ref{schemes-lemma-points-fibre-product}
，这些点 $y_{i, j}$ 恰好对应于环
$\kappa(u) \otimes_{\kappa(s)} \kappa(x_i)$ 中的素理想。这也说明它们为何只有有限个；
事实上 $m_i = [L_i : \kappa(s)]$，但我们不需要这一点。由 $L$ 的选择
（以及初等域论），对每对 $i, j$，$\kappa(u) \subset \kappa(y_{i, j})$
都是有限纯不可分扩张。此外，例如由 态射章引理
\ref{morphisms-lemma-base-change-quasi-finite}，态射 $X_U \to U$ 在所有点
$y_{i, j}$（所有 $i, j$）处拟有限。

\medskip\noindent
把 引理 \ref{more-morphisms-lemma-etale-makes-quasi-finite-finite-multiple-points} 应用于态射
$X_U \to U$、点 $u \in U$ 及各点 $y_{i, j} \in (X_U)_u$。这给出 \'etale 邻域
$(U', u') \to (U, u)$，满足 $\kappa(u) = \kappa(u')$；还给出开集
$V_{i, j} \subset X_{U'}$，满足该引理的性质 (\romannumeral1)、
(\romannumeral2) 与 (\romannumeral3)。我们断言 \'etale 邻域
$(U', u') \to (S, s)$ 及开集 $V_{i, j} \subset X_{U'}$ 即为本引理所求。
验证从略。
\end{proof}

\begin{lemma}
\label{more-morphisms-lemma-etale-splits-off-quasi-finite-part-technical}
设 $f : X \to S$ 是概形态射。设 $s \in S$，并设
$x_1, \ldots, x_n \in X_s$。假设
\begin{enumerate}
\item $f$ 局部有限型；
\item $f$ 分离；并且
\item $x_1, \ldots, x_n$ 是 $X_s$ 中两两不同的孤立点。
\end{enumerate}
则存在初等 \'etale 邻域 $(U, u) \to (S, s)$ 以及分解
$$
U \times_S X = W \amalg V_1 \amalg \ldots \amalg V_n
$$
为开闭子概形，使得态射 $V_i \to U$ 有限，$V_i \to U$ 在 $u$ 上的纤维是单点集
$\{v_i\}$，每个 $v_i$ 映到 $x_i$ 且 $\kappa(x_i) = \kappa(v_i)$，而
$W \to U$ 在 $u$ 上的纤维不含映到任何 $x_i$ 的点。
\end{lemma}

\begin{proof}
按
引理 \ref{more-morphisms-lemma-etale-makes-quasi-finite-finite-multiple-points}.
选择 $(U, u) \to (S, s)$ 与 $V_i \subset X_U$。由于 $X_U \to U$ 分离
(概形章引理 \ref{schemes-lemma-separated-permanence})
，而 $V_i \to U$ 有限、因而固有
(态射章引理 \ref{morphisms-lemma-finite-proper})
，由
态射章引理 \ref{morphisms-lemma-image-proper-scheme-closed}.
可知 $V_i \subset X_U$ 闭。因此 $V_i \cap V_j$ 是 $V_i$ 中不含 $v_i$ 的闭子集。
故 $V_i \cap V_j$ 在 $U$ 中的像是不含 $u$ 的闭集（因为 $V_i \to U$ 固有）。
缩小 $U$ 后，因而可假设
% P05-ERR-CH37-0084: Frozen source says V_i cap V_j is empty for all i,j, which is impossible when i=j; pairwise disjointness requires i != j.
$V_i \cap V_j = \emptyset$ 对所有 $i, j$ 成立。这给出引理中的分解。
\end{proof}

\noindent
下面给出归约到纯不可分域扩张的变体。

\begin{lemma}
\label{more-morphisms-lemma-etale-splits-off-quasi-finite-part-technical-variant}
设 $f : X \to S$ 是概形态射。设 $s \in S$，并设
$x_1, \ldots, x_n \in X_s$。假设
\begin{enumerate}
\item $f$ 局部有限型；
\item $f$ 分离；并且
\item $x_1, \ldots, x_n$ 是 $X_s$ 中两两不同的孤立点。
\end{enumerate}
则存在 \'etale 邻域 $(U, u) \to (S, s)$ 以及分解
$$
U \times_S X =
W \amalg
\ \coprod\nolimits_{i = 1, \ldots, n}
\ \coprod\nolimits_{j = 1, \ldots, m_i}
V_{i, j}
$$
为开闭子概形，使得态射 $V_{i, j} \to U$ 有限，$V_{i, j} \to U$ 在 $u$ 上的纤维
是单点集 $\{v_{i, j}\}$，每个 $v_{i, j}$ 映到 $x_i$，
$\kappa(u) \subset \kappa(v_{i, j})$ 纯不可分，而 $W \to U$ 在 $u$ 上的纤维
不含映到任何 $x_i$ 的点。
\end{lemma}

\begin{proof}
证明与 引理 \ref{more-morphisms-lemma-etale-splits-off-quasi-finite-part-technical}
完全相同，只是使用 Lemma
\ref{more-morphisms-lemma-etale-makes-quasi-finite-finite-multiple-points-var} 代替 Lemma
\ref{more-morphisms-lemma-etale-makes-quasi-finite-finite-multiple-points}。
\end{proof}

\noindent
下面的版本可能稍易理解。

\begin{lemma}
\label{more-morphisms-lemma-etale-splits-off-quasi-finite-part}
设 $f : X \to S$ 是概形态射。设 $s \in S$。假设
\begin{enumerate}
\item $f$ 局部有限型；
\item $f$ 分离；并且
\item $X_s$ 至多有有限多个孤立点。
\end{enumerate}
则存在初等 \'etale 邻域 $(U, u) \to (S, s)$ 以及分解
$$
U \times_S X = W \amalg V
$$
为开闭子概形，使得态射 $V \to U$ 有限，并且态射 $W \to U$ 的纤维 $W_u$ 不含
孤立点。特别地，若 $f^{-1}(s)$ 是有限集，则 $W_u = \emptyset$。
\end{lemma}

\begin{proof}
这由
引理 \ref{more-morphisms-lemma-etale-splits-off-quasi-finite-part-technical}
立即可得：取 $x_1, \ldots, x_n$ 为 $X_s$ 的全部孤立点，并令
$V = \bigcup V_i$。
\end{proof}






\section{整态射的 \'etale 局部化}
\label{more-morphisms-section-etale-localization-integral}

\noindent
下面给出 第 \ref{more-morphisms-section-etale-localization} 节 中结果在整态射情形的若干变体。

\begin{lemma}
\label{more-morphisms-lemma-etale-makes-integral-split}
设 $R \to S$ 是整环映射。设 $\mathfrak p \subset R$ 是素理想。假设
\begin{enumerate}
\item 位于 $\mathfrak p$ 上方的素理想 $\mathfrak q_1, \ldots, \mathfrak q_n$ 只有有限个；并且
\item 对每个 $i$，极大可分子扩张
% P05-ERR-CH37-0085: Frozen source uses undefined q in the extension tower; the surrounding indexed field is kappa(q_i).
$\kappa(\mathfrak q)/\kappa(\mathfrak q_i)_{sep}/\kappa(\mathfrak p)$
(域章引理 \ref{fields-lemma-separable-first})
在 $\kappa(\mathfrak p)$ 上有限。
\end{enumerate}
则存在 \'etale 环映射 $R \to R'$ 以及位于 $\mathfrak p$ 上方的素理想
$\mathfrak p'$，使得
$$
S \otimes_R R' = A_1 \times \ldots \times A_m
$$
，其中整映射 $R' \to A_j$ 恰有一个位于 $\mathfrak p'$ 上方的素理想
$\mathfrak r_j$，且 $\kappa(\mathfrak r_j)/\kappa(\mathfrak p')$ 纯不可分。
\end{lemma}

\begin{proof}[第一种证明]
本证明使用
代数章引理 \ref{algebra-lemma-etale-makes-quasi-finite-finite-variant}.
具体而言，选取该域在 $\kappa(\mathfrak p)$ 上的生成元
$\theta_i \in \kappa(\mathfrak q_i)_{sep}$（域章引理
\ref{fields-lemma-primitive-element}）。纤维环 $S \otimes_R \kappa(\mathfrak p)$
的谱是有限离散空间，其点对应于 $\mathfrak q_1, \ldots, \mathfrak q_n$。由中国剩余定理
(代数章引理 \ref{algebra-lemma-chinese-remainder})
可知 $S \otimes_R \kappa(\mathfrak p) \to \prod \kappa(\mathfrak q_i)$ 满射。
因此，对某个 $g \in R$、$g \not \in \mathfrak p$，用 $R_g$ 替换 $R$ 后，
可假设
$(0, \ldots, 0, \theta_i, 0, \ldots, 0) \in \prod \kappa(\mathfrak q_i)$
是某个 $x_i \in S$ 的像。设 $S' \subset S$ 是由这些 $x_i$ 生成的 $R$-子代数。
由于 $\Spec(S) \to \Spec(S')$ 满射
(代数章引理 \ref{algebra-lemma-integral-overring-surjective})
，可知 $\mathfrak q_i' = S' \cap \mathfrak q_i$ 是 $S'$ 中位于
$\mathfrak p$ 上方的各素理想。由 $x_i$ 的选择可知这些素理想彼此不同，并且
% P05-ERR-CH37-0086: Frozen source has the malformed phrase “distinct that and”.
$\kappa(\mathfrak q'_i)_{sep} = \kappa(\mathfrak q_i)_{sep}$.
特别地，域扩张 $\kappa(\mathfrak q_i)/\kappa(\mathfrak q'_i)$ 纯不可分。
由于 $R \to S'$ 有限，可以应用
代数章引理 \ref{algebra-lemma-etale-makes-quasi-finite-finite-variant}.
% P05-ERR-CH37-0087: Frozen source places a period before “and we get”, breaking the sentence.
从而得到 $R \to R'$、$\mathfrak p'$ 以及分解
$$
S' \otimes_R R' = A'_1 \times \ldots \times A'_m \times B'
$$
，其中整映射 $R' \to A'_j$ 恰有一个位于 $\mathfrak p'$ 上方的素理想
$\mathfrak r'_j$，使得 $\kappa(\mathfrak r'_j)/\kappa(\mathfrak p')$ 纯不可分；
而 $B'$ 没有位于 $\mathfrak p'$ 上方的素理想。由于 $R' \to B'$ 有限
（因为 $R \to S'$ 有限），在某个 $g' \in R'$、$g' \not \in \mathfrak p'$
处局部化 $R'$ 后，可假设 $B' = 0$。经映射
$S' \otimes_R R' \to S \otimes_R R'$，得到 $S$ 的相应分解。
\end{proof}

\begin{proof}[第二种证明]
本证明使用严格 Hensel 化。首先，
% P05-ERR-CH37-0088: Frozen source says “R is strictly henselization”; the intended condition is that R is strictly henselian/a strict henselization.
假设 $R$ 是以 $\mathfrak p$ 为极大理想的严格 Hensel 化。则
$S/\mathfrak p S$ 有有限多个对应于 $\mathfrak q_1, \ldots, \mathfrak q_n$ 的素理想；
它们均为极大理想，且其剩余域在 $\kappa(\mathfrak p)$ 上纯不可分。因此
% P05-ERR-CH37-0089: Frozen source localizes at undefined p_i after naming the relevant primes q_i.
$S/\mathfrak p S$ 等于 $\prod (S/\mathfrak p S)_{\mathfrak p_i}$。由
代数进阶章引理 \ref{more-algebra-lemma-characterize-henselian-pair}，可把这一乘积分解
提升为 $S$ 的如下所述的
% P05-ERR-CH37-0090: Frozen source says “product composition” instead of “product decomposition”.
乘积复合。

\medskip\noindent
在一般情形，设 $R^{sh}$ 是 $R_\mathfrak p$ 的严格 Hensel 化。于是可把第一段的结果
应用于 $R^{sh} \to S \otimes_R R^{sh}$。考虑 $S \otimes_R R^{sh}$ 中对应于
该乘积分解的 $m$ 个两两正交幂等元。由
代数章引理 \ref{algebra-lemma-strict-henselization-different}
，$R^{sh}$ 是 \'etale 环映射 $(R, \mathfrak p) \to (R', \mathfrak p')$ 的滤过余极限，
故这些幂等元按要求下降到某个 $R'$。
\end{proof}







\section{Zariski 主定理}
\label{more-morphisms-section-application-etale-neighbourhoods}

\noindent
本节证明 Grothendieck 重述的 Zariski 主定理。在本文中说“Zariski 主定理”时，通常指
% P05-ERR-CH37-0091: Frozen source says “in this content”; the idiom requires “in this context”.
引理 \ref{more-morphisms-lemma-finite-type-separated},
引理 \ref{more-morphisms-lemma-quasi-finite-separated-quasi-affine} 或
引理 \ref{more-morphisms-lemma-quasi-finite-separated-pass-through-finite} 中的任一个。在多数文献中，
最后一个结果被称为 Zariski 主定理。

\medskip\noindent
我们已经在
代数章定理 \ref{algebra-theorem-main-theorem}
中证明了代数版本，也已经用概形语言重述了该版本，见
态射章定理 \ref{morphisms-theorem-main-theorem}.
本节的版本更为精细；为得到完整结果，我们使用 Section
\ref{more-morphisms-section-etale-localization} 的 \'etale 局部化技术归约到代数情形。

\begin{lemma}
\label{more-morphisms-lemma-finite-type-separated}
设 $f : X \to S$ 是概形态射。假设 $f$ 有限型且分离。设 $S'$ 是 $S$ 在 $X$
中的正规化，见
态射章定义 \ref{morphisms-definition-normalization-X-in-Y}.
图示为：
$$
\xymatrix{
X \ar[rd]_f \ar[rr]_{f'} & & S' \ar[ld]^\nu \\
& S &
}
$$
则存在开子概形 $U' \subset S'$，使得
\begin{enumerate}
\item $(f')^{-1}(U') \to U'$ 是同构；并且
\item $(f')^{-1}(U') \subset X$ 是 $f$ 在其处拟有限的点集。
\end{enumerate}
\end{lemma}

\begin{proof}
由 态射章引理 \ref{morphisms-lemma-quasi-finite-points-open}，$f$ 在其处拟有限的
点组成开子集 $U \subset X$。本引理等价于
\begin{enumerate}
\item[(a)] $U' = f'(U) \subset S'$ 是开集；
\item[(b)] $U = (f')^{-1}(U')$；并且
\item[(c)] $U \to U'$ 是同构。
\end{enumerate}
任取 $x \in U$。我们断言存在开邻域 $f'(x) \in V \subset S'$，使得
$(f')^{-1}V \to V$ 是同构。先证明该断言蕴含本引理。事实上，
$(f')^{-1}V \cong V$ 在 $S$ 上局部有限型（因为它是 $X$ 的开子概形）；而对
$v \in V$，剩余域扩张 $\kappa(v)/\kappa(\nu(v))$ 是代数的（因为
$V \subset S'$ 且 $S'$ 在 $S$ 上整）。因此 $V \to S$ 的纤维离散
（态射章引理
\ref{morphisms-lemma-algebraic-residue-field-extension-closed-point-fibre})
，而 $(f')^{-1}V \to S$ 局部拟有限
(态射章引理 \ref{morphisms-lemma-locally-quasi-finite-fibres}).
这蕴含 $(f')^{-1}V \subset U$ 与 $V \subset U'$。由于 $x$ 任意，(a)、(b)、(c) 成立。

\medskip\noindent
令 $s = f(x)$。设 $(T, t) \to (S, s)$ 是初等 \'etale 邻域。用下标 ${}_T$
表示到 $T$ 的基变换。设 $y = (x, t) \in X_T$ 是纤维 $X_t$ 中位于 $x$
上方的唯一点。注意，$U_T \subset X_T$ 是 $f_T$ 在其处拟有限的点集，见
态射章引理 \ref{morphisms-lemma-base-change-quasi-finite}.
还要注意
$$
X_T \xrightarrow{f'_T} S'_T \xrightarrow{\nu_T} T
$$
是 $T$ 在 $X_T$ 中的正规化，见
引理 \ref{more-morphisms-lemma-normalization-smooth-localization}.
假设该断言对 $y \in U_T \subset X_T \to S'_T \to T$ 成立，即假设可找到开邻域
$f'_T(y) \in V' \subset S'_T$，使得 $(f'_T)^{-1}V' \to V'$ 是同构。态射
$S'_T \to S'$ 为 \'etale，故 $V'$ 的像 $V \subset S'$ 是开集。由于
$f'_T(y) \in V'$，有 $f'(x) \in V$。注意图
$$
\xymatrix{
(f'_T)^{-1}V' \ar[r] \ar[d] & (f')^{-1}(V) \ar[d] \\
V' \ar[r] & V
}
$$
是纤维方块（因为 $S'_T \times_{S'} X = X_T$）。由于左侧竖直箭头是同构，且
$\{V' \to V\}$ 是 \'etale 覆盖，由
下降章引理 \ref{descent-lemma-descending-property-isomorphism}.
可知右侧竖直箭头是同构。换言之，该断言对 $x \in U \subset X \to S' \to S$ 成立。

\medskip\noindent
由上一段的结果，为证明该断言，可以把 $S$ 替换为 $s = f(x)$ 的初等 \'etale 邻域。
因此可假设有分解
$$
X = V \amalg W
$$
为开闭子概形，其中 $V \to S$ 有限且 $x \in V$，见 Lemma
\ref{more-morphisms-lemma-etale-splits-off-quasi-finite-part-technical}。由于 $X$ 是 $S$ 上 $V$ 与
$W$ 的不交并，且 $V \to S$ 有限，可知 $S$ 在 $X$ 中的正规化是态射
$$
X = V \amalg W \longrightarrow V \amalg W' \longrightarrow S
$$
，其中 $W'$ 是 $S$ 在 $W$ 中的正规化，见
态射章引理 \ref{morphisms-lemma-normalization-in-disjoint-union},
\ref{morphisms-lemma-finite-integral} 与
\ref{morphisms-lemma-normalization-in-integral}.
该断言随即成立，证明完成。
\end{proof}

\begin{lemma}
\label{more-morphisms-lemma-quasi-finite-separated-quasi-affine}
\begin{slogan}
拟有限分离态射是拟仿射态射
\end{slogan}
设 $f : X \to S$ 是概形态射。假设 $f$ 拟有限且分离。设 $S'$ 是 $S$ 在 $X$
中的正规化，见
态射章定义 \ref{morphisms-definition-normalization-X-in-Y}.
图示为：
$$
\xymatrix{
X \ar[rd]_f \ar[rr]_{f'} & & S' \ar[ld]^\nu \\
& S &
}
$$
则 $f'$ 是拟紧开浸入，而 $\nu$ 是整态射。特别地，$f$ 拟仿射。
\end{lemma}

\begin{proof}
这由 引理 \ref{more-morphisms-lemma-finite-type-separated} 得到。事实上，由该引理，存在开子概形
$U' \subset S'$，使得 $(f')^{-1}(U') = X$，且 $X \to U'$ 是同构。换言之，
$f'$ 是开浸入。注意，$f'$ 拟紧，因为 $f$ 拟紧且 $\nu : S' \to S$ 分离
（概形章引理 \ref{schemes-lemma-quasi-compact-permanence}）。于是由
态射章引理 \ref{morphisms-lemma-characterize-quasi-affine}.
可知 $f$ 拟仿射。
\end{proof}

\begin{lemma}[Zariski's Main Theorem]
\label{more-morphisms-lemma-quasi-finite-separated-pass-through-finite}
\begin{reference}
\cite[IV Corollary 18.12.13]{EGA}
\end{reference}
设 $f : X \to S$ 是概形态射。假设 $f$ 拟有限且分离，并假设 $S$ 拟紧且拟分离。
则存在分解
$$
\xymatrix{
X \ar[rd]_f \ar[rr]_j & & T \ar[ld]^\pi \\
& S &
}
$$
，其中 $j$ 是拟紧开浸入，而 $\pi$ 有限。
\end{lemma}

\begin{proof}
设 $X \to S' \to S$ 如 Lemma
\ref{more-morphisms-lemma-quasi-finite-separated-quasi-affine} 的结论。由
性质章引理
\ref{properties-lemma-integral-algebra-directed-colimit-finite}
，可写成 $\nu_*\mathcal{O}_{S'} = \colim_{i \in I} \mathcal{A}_i$，即有限拟凝聚
$\mathcal{O}_S$-代数
$\mathcal{A}_i \subset \nu_*\mathcal{O}_{S'}$ 的有向余极限。于是对每个 $i$，
$\pi_i : T_i = \underline{\Spec}_S(\mathcal{A}_i) \to S$
是有限态射。注意，转移态射 $T_{i'} \to T_i$ 仿射，且 $S' = \lim T_i$。

\medskip\noindent
由 极限章引理 \ref{limits-lemma-descend-opens}，存在一个 $i$ 及拟紧开集
$U_i \subset T_i$，其在 $S'$ 中的逆像等于 $f'(X)$。对 $i' \geq i$，令
$U_{i'}$ 为 $U_i$ 在 $T_{i'}$ 中的逆像。于是
$X \cong f'(X) = \lim_{i' \geq i} U_{i'}$，见
极限章引理 \ref{limits-lemma-directed-inverse-system-has-limit}.
由 极限章引理 \ref{limits-lemma-finite-type-eventually-closed} 可知，对某个
$i' \geq i$，$X \to U_{i'}$ 是闭浸入。（事实上，对充分大的 $i'$，
$X \cong U_{i'}$，但我们不需要这一点。）因此 $X \to T_{i'}$ 是浸入。由
态射章引理 \ref{morphisms-lemma-factor-quasi-compact-immersion}
可将其分解为 $X \to T \to T_{i'}$，其中第一个箭头是开浸入，第二个是闭浸入。
结论得证。
\end{proof}

\begin{lemma}
\label{more-morphisms-lemma-quasi-finite-separated-pass-through-finite-addendum}
沿用 引理 \ref{more-morphisms-lemma-quasi-finite-separated-pass-through-finite} 的记号与假设。
再假设 $f$ 局部有限表现。则可选择该分解，使得 $T$ 在 $S$ 上有限且有限表现。
\end{lemma}

\begin{proof}
由 极限章引理
\ref{limits-lemma-finite-in-finite-and-finite-presentation}，可写
$T = \lim T_i$，其中所有 $T_i$ 在 $S$ 上有限且有限表现，转移态射
$T_{i'} \to T_i$ 是闭浸入。由
极限章引理 \ref{limits-lemma-descend-opens}
，存在一个 $i$ 及开子概形 $U_i \subset T_i$，其在 $T$ 中的逆像是 $X$。由
极限章引理
\ref{limits-lemma-finite-type-eventually-closed}
可知对充分大的 $i$ 有 $X \cong U_i$。用 $T_i$ 替换 $T$ 即完成证明。
\end{proof}








\section{Zariski 主定理的应用，I}
\label{more-morphisms-section-applications-zmt}

\noindent
第一个应用是把有限态射刻画为具有有限纤维的固有态射。

\begin{lemma}
\label{more-morphisms-lemma-characterize-finite}
设 $f : X \to S$ 是概形态射。下列条件等价：
\begin{enumerate}
\item $f$ 有限；
\item $f$ 固有且纤维有限；
\item $f$ 固有且局部拟有限；
\item $f$ 普遍闭、分离、局部有限型且纤维有限。
\end{enumerate}
\end{lemma}

\begin{proof}
由 态射章引理 \ref{morphisms-lemma-finite-proper},
\ref{morphisms-lemma-quasi-finite},
以及 \ref{morphisms-lemma-finite-quasi-finite}，(1) 蕴含 (2)。由 态射章引理
\ref{morphisms-lemma-finite-fibre}，(2) 蕴含 (3)。由固有态射的定义以及 Morphisms,
引理 \ref{morphisms-lemma-quasi-finite-locally-quasi-compact} 和
\ref{morphisms-lemma-quasi-finite}，(3) 蕴含 (4)。

\medskip\noindent
假设 (4)。取 $s \in S$。由 态射章引理
\ref{morphisms-lemma-finite-fibre}，$X_s$ 的有限多个点在 $X_s$ 中全都是孤立点。
按 引理 \ref{more-morphisms-lemma-etale-splits-off-quasi-finite-part}，选择初等 \'etale 邻域
$(U, u) \to (S, s)$ 以及分解 $X_U = V \amalg W$。
注意，由于 $X_s$ 的所有点都是孤立点，$W_u = \emptyset$。由于 $f$ 普遍闭，$W$
在 $U$ 中的像是不含 $u$ 的闭集。缩小 $U$ 后，可假设 $W = \emptyset$。换言之，
$X_U = V$ 在 $U$ 上有限。由于 $s \in S$ 任意，这意味着存在像覆盖 $S$ 的
\'etale 态射族 $\{U_i \to S\}$，使得基变换 $X_{U_i} \to U_i$ 有限。注意，
$\{U_i \to S\}$ 是 \'etale 覆盖，见 拓扑章定义
\ref{topologies-definition-etale-covering}。因此它是 fpqc 覆盖，见
拓扑章
引理 \ref{topologies-lemma-zariski-etale-smooth-syntomic-fppf-fpqc}。故由
下降章引理 \ref{descent-lemma-descending-property-finite} 可知 $f$ 有限。
\end{proof}

\noindent
由此得到以下有用结果。

\begin{lemma}
\label{more-morphisms-lemma-proper-finite-fibre-finite-in-neighbourhood}
设 $f : X \to S$ 是概形态射。设 $s \in S$。假设 $f$ 固有且
$f^{-1}(\{s\})$ 是有限集。则存在 $s$ 的开邻域 $V \subset S$，使得
$f|_{f^{-1}(V)} : f^{-1}(V) \to V$ 有限。
\end{lemma}

\begin{proof}
由 态射章引理 \ref{morphisms-lemma-finite-fibre}，态射 $f$ 在
$f^{-1}(\{s\})$ 的所有点处拟有限。由 态射章引理
\ref{morphisms-lemma-quasi-finite-points-open}，$f$ 在其处拟有限的点集是开集
$U \subset X$。令 $Z = X \setminus U$。则 $s \not \in f(Z)$。由于 $f$ 固有，
集合 $f(Z) \subset S$ 闭。选取 $s$ 的任意开邻域 $V \subset S$，使得
$f(Z) \cap V = \emptyset$。则 $f^{-1}(V) \to V$ 局部拟有限且固有，故拟有限
（态射章引理 \ref{morphisms-lemma-quasi-finite-locally-quasi-compact}），
从而纤维有限（态射章引理 \ref{morphisms-lemma-quasi-finite}），进而由
引理 \ref{more-morphisms-lemma-characterize-finite} 可知它有限。
\end{proof}

\begin{lemma}
\label{more-morphisms-lemma-flat-proper-family-cannot-collapse-fibre}
考虑概形的交换图
$$
\xymatrix{
X \ar[rr]_h \ar[rd]_f & & Y \ar[ld]^g \\
& S
}
$$
设 $s \in S$。假设
\begin{enumerate}
\item $X \to S$ 是固有态射；
\item $Y \to S$ 分离且局部有限型；并且
\item $X_s \to Y_s$ 的像有限。
\end{enumerate}
则存在包含 $s$ 的开子空间 $U \subset S$，使得 $X_U \to Y_U$ 经在 $U$ 上有限的
闭子概形 $Z \subset Y_U$ 分解。
% P05-ERR-CH37-0092: Frozen source says “open subspace” although S is a scheme and the intended object is an open subscheme.
\end{lemma}

\begin{proof}
设 $Z \subset Y$ 是 $h$ 的概形论像，见
态射章第 \ref{morphisms-section-scheme-theoretic-image} 节.
由 态射章引理 \ref{morphisms-lemma-scheme-theoretic-image-is-proper}，态射
$X \to Z$ 满，而 $Z \to S$ 固有。因此 $X_s \to Z_s$ 满。可知
% P05-ERR-CH37-0093: Frozen source has the stray word “either” in “We see that either (3) implies...”.
(3) 蕴含 $Z_s$ 有限。故由
引理 \ref{more-morphisms-lemma-proper-finite-fibre-finite-in-neighbourhood}.
$Z \to S$ 在 $s$ 的一个开邻域上有限。
\end{proof}





\section{Zariski 主定理的应用，II}
\label{more-morphisms-section-applications-zmt-II}

\noindent
本节给出 Zariski 主定理关于拟有限态射结构的另外若干推论。

\begin{lemma}
\label{more-morphisms-lemma-separated-locally-quasi-finite-over-affine}
设 $f : X \to Y$ 是分离且局部拟有限的态射，并且 $Y$ 仿射。则 $X$ 的每个有限点集
都包含在 $X$ 的某个仿射开集内。
\end{lemma}

\begin{proof}
设 $x_1, \ldots, x_n \in X$。选择拟紧开集 $U \subset X$，使得 $x_i \in U$。
由 引理 \ref{more-morphisms-lemma-quasi-finite-separated-quasi-affine}，$U \to Y$ 拟仿射。
故由
性质章引理 \ref{properties-lemma-ample-finite-set-in-affine}.
，存在包含 $x_1, \ldots, x_n$ 的仿射开集 $V \subset U$。
\end{proof}

\begin{lemma}
\label{more-morphisms-lemma-quasi-finite-finite-over-dense-open}
设 $f : Y \to X$ 是拟有限态射。则存在稠密开集 $U \subset X$，使得
$f|_{f^{-1}(U)} : f^{-1}(U) \to U$ 有限。
\end{lemma}

\begin{proof}
若 $U_i \subset X$（$i \in I$）是一族开集，使得限制
$f|_{f^{-1}(U_i)} : f^{-1}(U_i) \to U_i$ 有限，则令 $U = \bigcup U_i$，限制
$f|_{f^{-1}(U)} : f^{-1}(U) \to U$ 也有限，见
态射章引理 \ref{morphisms-lemma-finite-local}.
因此该问题在 $X$ 上是局部的，可假设 $X$ 仿射。

\medskip\noindent
假设 $X$ 仿射。写成 $Y = \bigcup_{j = 1, \ldots, m} V_j$，其中 $V_j$ 仿射。
这是可能的，因为 $f$ 拟有限，特别地拟紧。每个 $V_j \to X$ 拟有限且分离。
设 $\eta \in X$ 是 $X$ 的某个不可约分支的泛点。由 态射章引理
\ref{morphisms-lemma-quasi-finite} 与 \ref{morphisms-lemma-generically-finite}
可知，存在开邻域 $\eta \in U_\eta$，使得 $f^{-1}(U_\eta) \cap V_j \to U_\eta$
有限。可以选择对每个 $j = 1, \ldots, m$ 都适用的 $U_\eta$。注意，$X$ 的泛点集
在 $X$ 中稠密。因此存在稠密开集 $W = \bigcup_\eta U_\eta$，使得每个
$f^{-1}(W) \cap V_j \to W$ 都有限。只需证明存在稠密开集 $U \subset W$，使得
$f|_{f^{-1}(U)} : f^{-1}(U) \to U$ 有限。因此可把 $X$ 替换为 $W$ 的一个仿射开子概形，
并假设每个 $V_j \to X$ 有限。

\medskip\noindent
假设 $X$ 仿射，$Y = \bigcup_{j = 1, \ldots, m} V_j$ 且 $V_j$ 仿射，并且限制
$f|_{V_j} : V_j \to X$ 有限。令
$$
\Delta_{ij} =
\Big(\overline{V_i \cap V_j} \setminus V_i \cap V_j\Big) \cap V_j.
$$
这是 $V_j$ 的无处稠密闭子集，因为它是开子集 $V_i \cap V_j$ 在 $V_j$ 中的边界。由
态射章引理 \ref{morphisms-lemma-image-nowhere-dense-finite}
，像 $f(\Delta_{ij})$ 是 $X$ 的无处稠密闭子集。由
拓扑章引理 \ref{topology-lemma-nowhere-dense}
，并集 $T = \bigcup f(\Delta_{ij})$ 是 $X$ 的无处稠密闭子集。因此
$U = X \setminus T$ 是 $X$ 的稠密开子集。我们断言
$f|_{f^{-1}(U)} : f^{-1}(U) \to U$ 有限。为此，设 $U' \subset U$ 是仿射开集。
令 $Y' = f^{-1}(U') = U' \times_X Y$、
$V_j' = Y' \cap V_j = U' \times_X V_j$。考虑 $f$ 的限制
$$
f' = f|_{Y'} : Y' \longrightarrow U'
$$
。此态射现在具有如下性质：$Y' = \bigcup_{j = 1, \ldots, m} V'_j$ 是仿射开覆盖，
每个 $V'_j \to U'$ 有限，而 $V_i' \cap V_j'$ 在 $V'_i$ 与 $V'_j$ 中都（开且）闭。
因此 $V_i' \cap V_j'$ 仿射，且映射
$$
\mathcal{O}(V'_i) \otimes_{\mathbf{Z}} \mathcal{O}(V'_j)
\longrightarrow
\mathcal{O}(V'_i \cap V'_j)
$$
满射。这蕴含 $Y'$ 分离，见
概形章引理 \ref{schemes-lemma-characterize-separated}.
最后考虑交换图
$$
\xymatrix{
\coprod_{j = 1, \ldots, m} V'_j \ar[rd] \ar[rr] & & Y' \ar[ld] \\
& U' &
}
$$
东南方向的箭头有限、因而固有，水平箭头满，而西南方向的箭头分离。因此由
态射章引理 \ref{morphisms-lemma-image-proper-is-proper}
可知 $Y' \to U'$ 固有。由于它也拟有限，由 Lemma
\ref{more-morphisms-lemma-characterize-finite} 可知它有限，结论得证。
\end{proof}

\begin{lemma}
\label{more-morphisms-lemma-stratify-flat-fp-lqf-universally-bounded}
设 $f : X \to S$ 平坦、局部有限表现、分离、局部拟有限，并且纤维普遍有界。
则存在闭子集
$$
\emptyset = Z_{-1} \subset Z_0 \subset Z_1 \subset Z_2 \subset
\ldots \subset Z_n = S
$$
，使得令 $S_r = Z_r \setminus Z_{r - 1}$ 后，分层
$S = \coprod_{r = 0, \ldots, n} S_r$ 由如下泛性质刻画：给定 $g : T \to S$，
投影 $X \times_S T \to T$ 是次数为 $r$ 的有限局部自由态射，当且仅当集合论意义下
$g(T) \subset S_r$。
\end{lemma}

\begin{proof}
设整数 $n$ 是 $X \to S$ 各纤维次数的上界。由 态射章引理
\ref{morphisms-lemma-base-change-universally-bounded}，任意基变换的纤维次数也以
$n$ 为上界。特别地，引理陈述中出现的所有整数 $r$ 都满足 $\leq n$。对 $n$
归纳证明本引理。基础情形 $n = 0$ 显然。

\medskip\noindent
我们断言，满足 $\deg_{\kappa(s)}(X_s) = n$ 的点 $s \in S$ 组成开子集
$S_n \subset S$，而 $X \times_S S_n \to S_n$ 是次数为 $n$ 的有限局部自由态射。
事实上，设 $s \in S$ 是这样的点。按
引理 \ref{more-morphisms-lemma-etale-splits-off-quasi-finite-part}.
选择初等 \'etale 态射 $(U, u) \to (S, s)$ 及分解
$U \times_S X = W \amalg V$。由于 $V \to U$ 有限、平坦且局部有限表现，
可知 $V \to U$ 有限局部自由，见
态射章引理 \ref{morphisms-lemma-finite-flat}.
把 $U$ 缩小为 $u$ 的更小邻域后，可假设 $V \to U$ 是某个次数 $d$ 的有限局部自由态射，见
态射章引理 \ref{morphisms-lemma-finite-locally-free}.
由于 $u \mapsto s$ 且 $W_u = \emptyset$，有 $d = n$。由于 $n$ 是纤维次数的最大值，
可知 $W = \emptyset$！因此 $U \times_S X \to U$ 是次数为 $n$ 的有限局部自由态射。
由 下降章引理 \ref{descent-lemma-descending-property-finite-locally-free}，
$X \to S$ 在 $s$ 的开邻域 $\Im(U \to S)$ 上是次数为 $n$ 的有限局部自由态射
（态射章引理 \ref{morphisms-lemma-etale-open}）。该断言得证。

\medskip\noindent
赋予 $S' = S \setminus S_n$ 约化诱导概形结构，并令 $X' = X \times_S S'$。
注意，$X' \to S'$ 的纤维次数普遍以 $n - 1$ 为上界。由归纳假设，得到适合于态射
$X' \to S'$ 的分层 $S' = S_0 \amalg \ldots \amalg S_{n - 1}$。我们断言
$S = \coprod_{r = 0, \ldots, n} S_r$ 对态射 $X \to S$ 适用。设
$g : T \to S$ 是概形态射，并假设 $X \times_S T \to T$ 是次数为 $r$ 的有限局部自由
态射。如上所述，这蕴含 $r \leq n$。若 $r = n$，则 $T \to S$ 显然经 $S_n$
分解。若 $r < n$，则
% P05-ERR-CH37-0094: Frozen source uses undefined S_d here; the complement fixed in the preceding sentence is S \ S_n.
$g(T) \subset S' = S \setminus S_d$ (集合论意义下) hence
集合论意义下成立，故 $T_{red} \to S$ 经 $S'$ 分解，见
概形章引理 \ref{schemes-lemma-map-into-reduction}.
注意，作为基变换，$X \times_S T_{red} \to T_{red}$ 也是次数为 $r$ 的有限局部自由态射。
由分层 $S' = \coprod_{r = 0, \ldots, n - 1} S_r$ 的泛性质，
$g(T) = g(T_{red})$ 包含于 $S_r$。反过来，假设有 $g : T \to S$，使得集合论意义下
$g(T) \subset S_r$。若 $r = n$，则 $g$ 经 $S_n$ 分解，而作为基变换，
$X \times_S T \to T$ 显然是次数为 $n$ 的有限局部自由态射。若 $r < n$，则
$X \times_S T \to T$ 分离、平坦且局部有限表现，并且它在 $T_{red}$ 上的限制是次数为
$r$ 的有限局部自由态射。由于 $T_{red} \to T$ 是普遍同胚，可知
$X \times_S T_{red} \to X \times_S T$ 也是普遍同胚，因而
$X \times_S T \to T$ 普遍闭（因为有限态射
$X \times_S T_{red} \to T_{red}$ 具有这一性质）。例如由 Lemma
\ref{more-morphisms-lemma-characterize-finite}，可知 $X \times_S T \to T$ 有限。再由 Morphisms,
引理 \ref{morphisms-lemma-finite-flat}，可知 $X \times_S T \to T$ 有限局部自由。最后，由于所有纤维的次数都是
$r$，其次数为 $r$。
\end{proof}

\begin{lemma}
\label{more-morphisms-lemma-stratify-flat-fp-qf}
设 $f : X \to S$ 是平坦、局部有限表现、分离且拟有限的概形态射。则存在闭子集
$$
\emptyset = Z_{-1} \subset Z_0 \subset Z_1 \subset Z_2 \subset
\ldots \subset S
$$
，使得令 $S_r = Z_r \setminus Z_{r - 1}$ 后，分层 $S = \coprod S_r$ 由如下泛性质
刻画：给定态射 $g : T \to S$，投影 $X \times_S T \to T$ 是次数为 $r$ 的有限局部自由
态射，当且仅当集合论意义下 $g(T) \subset S_r$。此外，包含映射 $S_r \to S$ 拟紧。
\end{lemma}

\begin{proof}
该问题在 $S$ 上是局部的，故可假设 $S$ 仿射。由 态射章引理
\ref{morphisms-lemma-locally-quasi-finite-qc-source-universally-bounded}
，在此情形 $f$ 的纤维普遍有界。因此该分层的存在性由
引理 \ref{more-morphisms-lemma-stratify-flat-fp-lqf-universally-bounded}.
得到。

\medskip\noindent
我们将证明对每个 $r \geq 0$，$U_r = S \setminus Z_r \to S$ 拟紧。由初等拓扑，
这将证明最后一个断言。由于拟紧映射的复合拟紧，只需证明 $U_r \to U_{r - 1}$ 拟紧。
选择仿射开集 $W \subset U_{r - 1}$。写成 $W = \Spec(A)$。则对某个理想
$I \subset A$ 有 $Z_r \cap W = V(I)$，且
$X \times_S \Spec(A/I) \to \Spec(A/I)$ 是次数为 $r$ 的有限局部自由态射。注意
$A/I = \colim A/I_i$，其中 $I_i \subset I$ 遍历有限生成理想。由
极限章引理 \ref{limits-lemma-descend-finite-locally-free}
可知，对某个 $i$，$X \times_S \Spec(A/I_i) \to \Spec(A/I_i)$ 是次数为 $r$
的有限局部自由态射。（这里使用了 $X \to S$ 有限表现；事实上，它局部有限表现、分离且
拟紧。）因此，集合论意义下，$\Spec(A/I_i) \to \Spec(A) = W$ 经
$Z_r \cap W$ 分解。于是 $Z_r \cap W = V(I_i)$ 是 $A$ 中某个有限元素集的零点集。
这意味着 $W \setminus Z_r$ 是有限多个标准开集之并，因而拟紧，如所求。
\end{proof}

\begin{lemma}
\label{more-morphisms-lemma-stratify-flat-fp-lqf}
设 $f : X \to S$ 是平坦、局部有限表现、分离且局部拟有限的概形态射。则存在开子概形
$$
S = U_0 \supset U_1 \supset U_2 \supset \ldots
$$
，使得当 $k$ 是域时，态射 $\Spec(k) \to S$ 经 $U_d$ 分解，当且仅当
$X \times_S \Spec(k)$ 在 $k$ 上的次数 $\geq d$。
\end{lemma}

\begin{proof}
该陈述仅仅是说，纤维次数 $\geq d$ 的点集是开集。因此可在 $S$ 上局部工作，并假设
$S$ 仿射。在此情形，对每个拟紧开集 $W \subset X$，由 Lemma
\ref{more-morphisms-lemma-stratify-flat-fp-qf}，$W \to S$ 的纤维次数 $\geq d$ 的点集 $U_d(W)$
是开集。由于 $U_d = \bigcup_W U_d(W)$，结论成立。
\end{proof}

\begin{lemma}
\label{more-morphisms-lemma-go-down-with-annihilators}
设 $f : X \to S$ 是平坦、局部有限表现且局部拟有限的概形态射。设
$g \in \Gamma(X, \mathcal{O}_X)$ 非零。则存在满足 $g|_V \not = 0$ 的开集
$V \subset X$、开集 $U \subset S$，它们构成交换图
$$
\xymatrix{
V \ar[r] \ar[d]_\pi & X \ar[d]^f \\
U \ar[r] & S,
}
$$
；还存在拟凝聚子层 $\mathcal{F} \subset \mathcal{O}_U$、整数 $r > 0$ 以及单射的
$\mathcal{O}_U$-模映射 $\mathcal{F}^{\oplus r} \to \pi_*\mathcal{O}_V$，其像包含
$g|_V$。
\end{lemma}

\begin{proof}
可假设 $X$ 与 $S$ 仿射。由 Lemmas
\ref{more-morphisms-lemma-stratify-flat-fp-lqf-universally-bounded} 与
\ref{more-morphisms-lemma-stratify-flat-fp-qf}，得到滤过
$\emptyset = Z_{-1} \subset Z_0 \subset Z_1 \subset Z_2 \subset \ldots
\subset Z_n = S$。设 $T \subset X$ 是有限 $\mathcal{O}_X$-模
$\Im(g : \mathcal{O}_X \to \mathcal{O}_X)$ 的概形论支撑。注意，$T$ 是 $g$ 作为
$\mathcal{O}_X$ 截面的支撑（模章定义 \ref{modules-definition-support}），
且对任意开集 $V \subset X$，$g|_V \not = 0$ 当且仅当
$V \cap T \not = \emptyset$。设 $r$ 是使得集合论意义下 $f(T) \subset Z_r$ 的最小整数。
取 $\xi \in T$ 为 $T$ 的某个不可约分支的泛点，使得
$f(\xi) \not \in Z_{r - 1}$（因而 $f(\xi) \in Z_r$）。可把 $S$ 替换为 $f(\xi)$
的一个包含于 $S \setminus Z_{r - 1}$ 的仿射邻域。写成 $S = \Spec(A)$，并设
$I = (a_1, \ldots, a_m) \subset A$ 是有限生成理想，使得集合论意义下
$V(I) = Z_r$（见
代数章引理 \ref{algebra-lemma-qc-open}).
由 $r$ 的选择，$g$ 的支撑包含于 $f^{-1}V(I)$，故存在整数 $N$，使得对
$j = 1, \ldots, m$ 有 $a_j^N g = 0$。
% P05-ERR-CH37-0095: Frozen source replaces a_j by a_j^r although the annihilating exponent just constructed is N.
用 $a_j^r$ 替换 $a_j$ 后，可假设 $Ig = 0$。对任意 $A$-模 $M$，以 $M[I]$
表示 $M$ 的 $I$-挠部分，即 $M[I] = \{m \in M \mid Im = 0\}$。写成
$X = \Spec(B)$，于是 $g \in B[I]$。由于 $A \to B$ 平坦，可知
$$
B[I] = A[I] \otimes_A B \cong A[I] \otimes_{A/I} B/IB
$$
由 $Z_r$ 的选择，$A/I$-模 $B/IB$ 是秩为 $r$ 的有限局部自由模。因此，把 $S$
替换为 $f(\xi)$ 的更小仿射开邻域后，可假设作为 $A/I$-模有
$B/IB \cong (A/IA)^{\oplus r}$。选择映射 $\psi : A^{\oplus r} \to B$，使其模
$I$ 约化为上一句的同构。于是诱导映射
$$
A[I]^{\oplus r} \longrightarrow B[I]
$$
是同构。取 $\mathcal{F}$ 为与 $A$-模 $A[I]$ 相伴的拟凝聚层，并取映射
$\mathcal{F}^{\oplus r} \to \pi_*\mathcal{O}_V$ 为对应于
$A[I]^{\oplus r} \subset A^{\oplus r} \to B$ 的映射，即得本引理。
\end{proof}

\begin{lemma}
\label{more-morphisms-lemma-there-is-a-scheme-integral-over}
设 $U \to X$ 是概形的满 \'etale 态射。假设 $X$
拟紧且拟分离。则存在满整态射 $Y \to X$，使得对
每个 $y \in Y$，都有一个开邻域 $V \subset Y$，
且 $V \to X$ 经由 $U$ 分解。事实上，可以假设
$Y \to X$ 是有限且有限表示的。
\end{lemma}

\begin{proof}
由于 $X$ 拟紧，存在有限多个仿射开集
$U_i \subset U$，使得 $U' = \coprod U_i \to X$ 是满射。
用 $U'$ 替换 $U$ 后，可假设 $U$ 是仿射的。
特别地，$U \to X$ 是分离的
（概形章引理 \ref{schemes-lemma-affine-separated}）。
于是存在一个整数 $d$，界定 $U \to X$ 的几何纤维的次数
（见 态射章引理
\ref{morphisms-lemma-locally-quasi-finite-qc-source-universally-bounded}）。
我们对 $d$ 归纳，为所有满 \'etale 映到 $X$ 的拟紧分离概形
$U$ 证明本引理。若 $d = 1$，则 $U = X$，取 $Y = U$
即可得到结论。以下假设 $d > 1$。

\medskip\noindent
应用引理 \ref{more-morphisms-lemma-quasi-finite-separated-quasi-affine}，
得到一个分解
$$
\xymatrix{
U \ar[rr]_j \ar[rd] & & Y \ar[ld]^\pi \\
& X
}
$$
其中 $\pi$ 是整态射，$j$ 是拟紧开浸入。我们可以且确实假设
$j(U)$ 在 $Y$ 中概形论稠密。注意
$$
U \times_X Y = U \amalg W
$$
其中第一个直和项是 $U \to U \times_X Y$ 的像
（由 概形章引理 \ref{schemes-lemma-semi-diagonal}，它是闭的；
又因为它作为两个在 $Y$ 上 \'etale 的概形之间的态射是 \'etale 的，
所以也是开的），第二个直和项是其开闭补集。
$W$ 的像 $V \subset Y$ 是包含 $Y \setminus U$ 的开子概形。

\medskip\noindent
\'Etale 态射 $W \to Y$ 的几何纤维的基数均为 $< d$。
事实上，对 $U \subset Y$ 的几何点直接考察即可看出这一点。
由于 $U \subset Y$ 稠密，例如由引理
\ref{more-morphisms-lemma-stratify-flat-fp-lqf-universally-bounded}，
它对 $Y$ 的所有几何点均成立
（拟紧分离 \'etale 态射的纤维次数在特化下不会增大）。因此，可以将
归纳假设应用于 $W \to V$，得到一个满整态射
$Z \to V$，其中 $Z$ 是概形，且该态射在 Zariski 局部经由 $W$ 分解。
选取分解 $Z \to Z' \to Y$，使得 $Z' \to Y$ 是整态射，
$Z \to Z'$ 是开浸入
（引理 \ref{more-morphisms-lemma-quasi-finite-separated-quasi-affine}）。
用 $Z$ 在 $Z'$ 中的概形论闭包替换 $Z'$ 后，
可假设 $Z$ 在 $Z'$ 中概形论稠密。这样便有
$Z' \times_Y V = Z$。最后，令 $T \subset Y$ 为赋予
$Y \setminus V$ 的诱导既约闭子概形结构。考虑态射
$$
Z' \amalg T \longrightarrow X
$$
按构造，这是满整态射。由于 $T \subset U$，态射 $T \to X$
显然经由 $U$ 分解。另一方面，取一点 $z \in Z'$。
若 $z \not \in Z$，则 $z$ 映到 $Y \setminus V \subset U$
中的一点，因而可找到 $z$ 的一个邻域，使该邻域上的态射经由 $U$ 分解。
若 $z \in Z$，则存在邻域 $\Omega \subset Z$，它经由
$W \subset U \times_X Y$ 分解，因而也经由 $U$ 分解。
这证明了所需对象的存在性。

\medskip\noindent
假设已经找到整且满的 $Y \to X$，并且它在 Zariski 局部经由
$U$ 分解。选取有限仿射开覆盖 $Y = \bigcup V_j$，使得
$V_j \to X$ 经由 $U$ 分解。可以写成 $Y = \lim Y_i$，
其中 $Y_i \to X$ 是有限且有限表示的；见 极限章引理
\ref{limits-lemma-integral-limit-finite-and-finite-presentation}。
当 $i$ 充分大时，可以找到仿射开集 $V_{i, j} \subset Y_i$，
其在 $Y$ 中的逆像恰为 $V_j$；见
极限章引理 \ref{limits-lemma-descend-opens}。
进一步增大 $i$，则 $X$ 上的态射 $V_j \to U$ 来自
$X$ 上的态射 $V_{i, j} \to U$；见 极限章命题
\ref{limits-proposition-characterize-locally-finite-presentation}。
证明完毕。
\end{proof}







\section{连通纤维态射的应用}
\label{more-morphisms-section-application-geometrically-connected}

\noindent
本节证明若干引理，用以构造所有纤维均几何连通或几何整的态射。
这些结果稍后在研究有限型态射的局部结构时会很有用。

% P05-EN-0096: the frozen source reads "an a point"; Chinese follows the evident intended syntax.
\begin{lemma}
\label{more-morphisms-lemma-descent-connected-fibres}
考虑概形态射的图表
$$
\xymatrix{
Z \ar[r]_{\sigma} \ar[rd] & X \ar[d] \\
& Y
}
$$
以及一点 $y \in Y$。假设
\begin{enumerate}
\item $X \to Y$ 是有限表示且平坦的，
\item $Z \to Y$ 是有限局部自由的，
\item $Z_y \not = \emptyset$,
\item $X \to Y$ 的所有纤维均几何既约，并且
\item $X_y$ 在 $\kappa(y)$ 上几何连通。
\end{enumerate}
则存在拟紧开集 $X^0 \subset X$，使得 $X^0_y = X_y$，
并且 $X^0 \to Y$ 的所有非空纤维均几何连通。
\end{lemma}

\begin{proof}
本证明将用到平坦、有限表示和有限局部自由这些性质在基变换与复合下保持，
还将用到有限局部自由态射既是开态射又是闭态射。这些事实见
态射章引理
\ref{morphisms-lemma-base-change-flat},
\ref{morphisms-lemma-base-change-finite-presentation},
\ref{morphisms-lemma-base-change-finite-locally-free},
\ref{morphisms-lemma-composition-flat},
\ref{morphisms-lemma-composition-finite-presentation},
\ref{morphisms-lemma-composition-finite-locally-free},
\ref{morphisms-lemma-fppf-open}, 及
\ref{morphisms-lemma-finite-proper}.

\medskip\noindent
注意，$X_Z \to Z$ 是有限表示的平坦态射，并具有由 $\sigma$
给出的截面 $s$。以 $X_Z^0$ 表示
Situation \ref{more-morphisms-situation-connected-along-section}
所定义的 $X_Z$ 的子集。由
引理 \ref{more-morphisms-lemma-connected-along-section-open}，
它是 $X_Z$ 的开子集。

\medskip\noindent
$X$ 到 $Z \times_Y Z$ 的拉回 $X_{Z \times_Y Z}$ 带有两个截面
$s_0, s_1$，即 $s$ 沿
$\text{pr}_0, \text{pr}_1 : Z \times_Y Z \to Z$ 的基变换。
Situation \ref{more-morphisms-situation-connected-along-section} 的构造给出两个子集
$(X_{Z \times_Y Z})_{s_0}^0$ 和
$(X_{Z \times_Y Z})_{s_1}^0$。由
引理 \ref{more-morphisms-lemma-base-change-connected-along-section}，
它们分别是 $X_Z^0$ 在态射
$1_X \times \text{pr}_0, 1_X \times \text{pr}_1 : X_{Z \times_Y Z} \to X_Z$
下的逆像。特别地，这些子集是开的。

\medskip\noindent
令 $(Z \times_Y Z)_y = \{z_1, \ldots, z_n\}$。
由于 $X_y$ 几何连通，可见
$(X_{Z \times_Y Z})_{s_0}^0$ 和 $(X_{Z \times_Y Z})_{s_1}^0$
在每个 $z_i$ 上的纤维相同（都等于整个纤维）。换言之，
$$
s_0(z_i) \in (X_{Z \times_Y Z})_{s_1}^0
\quad\text{且}\quad
s_1(z_i) \in (X_{Z \times_Y Z})_{s_0}^0.
$$
由于子集 $(X_{Z \times_Y Z})_{s_0}^0$ 和 $(X_{Z \times_Y Z})_{s_1}^0$
在 $X_{Z \times_Y Z}$ 中是开的，存在 $(Z \times_Y Z)_y$ 的一个开邻域
$W \subset Z \times_Y Z$，使得
$$
s_0(W) \subset (X_{Z \times_Y Z})_{s_1}^0
\quad\text{且}\quad
s_1(W) \subset (X_{Z \times_Y Z})_{s_0}^0.
$$
于是由 Situation \ref{more-morphisms-situation-connected-along-section}
中的构造直接可得
$$
p^{-1}(W) \cap (X_{Z \times_Y Z})_{s_0}^0
=
p^{-1}(W) \cap (X_{Z \times_Y Z})_{s_1}^0
$$
% P05-EN-0097: the frozen source gives the projection target as Z \times_W Z;
% the evident ambient base is Z \times_Y Z. Formula retained unchanged.
其中 $p : X_{Z \times_Y Z} \to Z \times_W Z$ 是投影。
由于 $Z \times_Y Z \to Y$ 有限局部自由，因而既开又闭，
存在 $y$ 的一个仿射开邻域 $V \subset Y$，使得
$q^{-1}(V) \subset W$，其中 $q : Z \times_Y Z \to Y$ 是结构态射。
为证明本引理，可以用 $V$ 替换 $Y$。这样便看出
% P05-EN-0098: the frozen source says X_Z^0 \subset Y_Z; expected ambient is X_Z.
$X_Z^0 \subset Y_Z$ 是开集，并且
$$
(1_X \times \text{pr}_0)^{-1}(X_Z^0) =
(1_X \times \text{pr}_1)^{-1}(X_Z^0).
$$
这意味着 $X_Z^0$ 的像 $X^0 \subset X$ 是开集，且
$(X_Z \to X)^{-1}(X^0) = X_Z^0$；见
下降章引理 \ref{descent-lemma-open-fpqc-covering}。
最后，由 引理 \ref{more-morphisms-lemma-connected-along-section-locally-constructible}，
$X_Z^0$ 是拟紧的，因而 $X^0$ 也是拟紧的
（这里用到此时 $Y$ 仿射，故 $X$ 拟紧且拟分离；于是局部可构造等同于
可构造，特别是拟紧的；略去细节）。由此可见，$X^0$ 具有全部所需性质。
\end{proof}

\begin{lemma}
\label{more-morphisms-lemma-fibre-geometrically-connected-reduced}
设 $h : Y \to S$ 是概形态射，$s \in S$ 是一点，
$T \subset Y_s$ 是开子概形。假设
\begin{enumerate}
\item $h$ 平坦且有限表示，
\item $h$ 的所有纤维均几何既约，并且
\item $T$ 在 $\kappa(s)$ 上几何连通。
\end{enumerate}
则可以找到一个仿射初等 \'etale 邻域
$(S', s') \to (S, s)$ 和一个拟紧开集 $V \subset Y_{S'}$，使得
\begin{enumerate}
\item[(a)] $V \to S'$ 的所有纤维均几何连通；
\item[(b)] $V_{s'} = T \times_s s'$.
\end{enumerate}
\end{lemma}

\begin{proof}
问题显然在 $S$ 上是局部的，故可以用 $s$ 的一个仿射开邻域替换 $S$。
$Y_s$ 上的拓扑由 $Y$ 上的拓扑诱导；见
概形章引理 \ref{schemes-lemma-fibre-topological}。
因此，可以找到拟紧开集 $V \subset Y$，使得 $V_s = T$。
$h$ 在 $V$ 上的限制是拟紧的（因为 $S$ 仿射而 $V$ 拟紧）、
拟分离的、局部有限表示的且平坦的，因而是有限表示的平坦态射。
所以用 $V$ 替换 $Y$ 后，除 (1) 和 (2) 外，还可假设
$Y_s = T$ 且 $S$ 仿射。

\medskip\noindent
选取闭点 $y \in Y_s$，使得 $h$ 在 $y$ 处是 Cohen--Macaulay 的；见
引理 \ref{more-morphisms-lemma-flat-finite-presentation-CM-open}。
由 引理 \ref{more-morphisms-lemma-slice-CM}，存在图表
$$
\xymatrix{
Z \ar[r] \ar[rd] & Y \ar[d] \\
& S
}
$$
使得 $Z \to S$ 平坦、局部有限表示且局部拟有限，并且
$Z_s = \{y\}$。应用
引理 \ref{more-morphisms-lemma-etale-makes-quasi-finite-finite-at-point}，
找到一个初等邻域 $(S', s') \to (S, s)$ 和一个开集
$Z' \subset Z_{S'} = S' \times_S Z$，使得 $Z' \to S'$ 有限，
并且其中恰有一点 $z' \in Z'$ 位于 $s$ 上方。注意，
$Z' \to S'$ 还是局部有限表示且平坦的（因为它是 $Z \to S$
的基变换中的开集），故 $Z' \to S'$ 有限局部自由；见
态射章引理 \ref{morphisms-lemma-finite-flat}。
注意，作为 $h$ 的基变换，$Y_{S'} \to S'$ 平坦且有限表示，
并具有几何既约纤维。此外，$Y_{s'} = Y_s$ 几何连通。
将 引理 \ref{more-morphisms-lemma-descent-connected-fibres}
应用于 $S'$ 上的 $Z' \to Y_{S'}$，得到拟紧开集
$V \subset Y_{S'}$，它满足 (2)，且在 $S'$ 上的纤维或为空，
或几何连通。由于 $V \to S'$ 是开态射
（态射章引理 \ref{morphisms-lemma-fppf-open}），用 $s'$
的一个仿射开邻域替换 $S'$ 后，可以假设 $V \to S'$ 是满射，
从而 (1) 成立。
\end{proof}

\begin{lemma}
\label{more-morphisms-lemma-cover-by-geometrically-connected}
设 $f : X \to S$ 是局部有限表示且平坦、并具有几何既约纤维的概形态射。
则存在一个 \'etale 覆盖 $\{X_i \to X\}_{i \in I}$，
使得 $X_i \to S$ 分解为 $X_i \to S_i \to S$，其中
$S_i \to S$ 是 \'etale 态射，而 $X_i \to S_i$ 是有限表示的平坦态射，
并具有几何连通且几何既约的纤维。
\end{lemma}

\begin{proof}
取一点 $x \in X$，其像为 $s \in S$。我们将构造图表
$$
\xymatrix{
X' \ar[r] \ar[rd] & S' \times_S X \ar[r] \ar[d] & X \ar[d] \\
& S' \ar[r] & S
}
$$
以及点 $s' \in S'$、$x' \in X'$、$y \in S' \times_S X$，使得
$x'$ 映到 $x$，$(S', s') \to (S, s)$ 是 \'etale 邻域，
$(X', x') \to (S' \times_S X, y)$ 是 \'etale 邻域\footnote{证明实际上给出
一个开集 $X' \subset S' \times_S X$。}，并且
$X' \to S'$ 具有几何连通纤维。如果对每个 $x \in X$ 都能如此构造，
则本引理成立（覆盖的成员由这样得到的诸 \'etale 态射
$X' \to X$ 给出）。
% P05-EN-0099: the frozen source reads "is the replace"; Chinese follows the intended infinitive.
第一步是用 $x$ 与 $s$ 的仿射开邻域分别替换 $X$ 与 $S$，
从而化归到 $X$ 和 $S$ 均仿射的情形（此时 $f$ 为有限表示）。

\medskip\noindent
选取 $\kappa(s)$ 的一个可分代数扩张 $\overline{k}$，并以
$X_{\overline{k}}$ 表示 $X_s$ 的相应基变换。选取
$X_{\overline{k}}$ 中映到 $x \in X_s$ 的一点 $\overline{x}$，
以及 $\overline{x}$ 的一个连通拟紧开邻域
$\overline{V} \subset X_{\overline{k}}$。
（这是可能的，因为域上局部有限型的任意概形作为局部 Noether 拓扑空间
是局部连通的。）由
簇章引理 \ref{varieties-lemma-Galois-action-quasi-compact-open}，
可以找到有限可分扩张 $k'/\kappa(s)$ 和拟紧开集
$V' \subset X_{k'}$，使其基变换为 $\overline{V}$。特别地，
$V'$ 在 $k'$ 上几何连通；见
簇章引理 \ref{varieties-lemma-characterize-geometrically-connected}。
由 引理 \ref{more-morphisms-lemma-realize-prescribed-residue-field-extension-etale}，
可以找到 \'etale 邻域 $(S', s') \to (S, s)$，使得
$\kappa(s')$ 作为 $\kappa(s)$ 的扩张同构于 $k'$。
以 $x' \in (S' \times_S X)_{s'} = X_{k'}$ 表示 $\overline{x}$ 的像。
于是，用 $(S', s')$ 替换 $(S, s)$，并用
$(S' \times_S X, x')$ 替换 $(X, x)$ 后，就化归到下一段处理的情形。

\medskip\noindent
假设存在包含 $x$ 且几何连通的拟紧开集 $V \subset X_s$。
于是可以应用 引理 \ref{more-morphisms-lemma-fibre-geometrically-connected-reduced}，
找到一个仿射 \'etale 邻域 $(S', s') \to (S, s)$ 和一个拟紧开集
$X' \subset S' \times_S X$，使得 $X' \to S'$ 具有几何连通纤维，
并且 $X'$ 含有一个映到 $x$ 的点。证明完毕。
\end{proof}

\begin{lemma}
\label{more-morphisms-lemma-normal-morphism-irreducible}
设 $h : Y \to S$ 是概形态射，$s \in S$ 是一点，
$T \subset Y_s$ 是开子概形。假设
\begin{enumerate}
\item $h$ 是有限表示的，
\item $h$ 是正规态射，并且
\item $T$ 在 $\kappa(s)$ 上几何不可约。
\end{enumerate}
则可以找到一个仿射初等 \'etale 邻域
$(S', s') \to (S, s)$ 和一个拟紧开集 $V \subset Y_{S'}$，使得
\begin{enumerate}
\item[(a)] $V \to S'$ 的所有纤维均几何整；
\item[(b)] $V_{s'} = T \times_s s'$.
\end{enumerate}
\end{lemma}

\begin{proof}
应用 引理 \ref{more-morphisms-lemma-fibre-geometrically-connected-reduced}，
找到一个仿射初等 \'etale 邻域 $(S', s') \to (S, s)$ 和一个拟紧开集
$V \subset Y_{S'}$，使得 $V \to S'$ 的所有纤维均几何连通，并且
$V_{s'} = T \times_s s'$。由于 $V$ 是 $h$ 的基变换中的开集，
$V \to S'$ 的所有纤维均几何正规；见 引理 \ref{more-morphisms-lemma-normal}。
特别地，它们几何既约。为完成证明，只需说明它们几何不可约。
若 $t \in S'$，则 $V_t$ 在 $\kappa(t)$ 上有限型，因此
$V_t \times_{\kappa(t)} \overline{\kappa(t)}$ 在
$\overline{\kappa(t)}$ 上有限型，从而是 Noether 的。
由 $S' \to S$ 的选取，概形
$V_t \times_{\kappa(t)} \overline{\kappa(t)}$ 是连通的。
故由 性质章引理 \ref{properties-lemma-normal-Noetherian}，
$V_t \times_{\kappa(t)} \overline{\kappa(t)}$ 不可约，结论得证。
\end{proof}








\section{有限型态射结构的应用}
\label{more-morphisms-section-application-finite-type}

\noindent
本节的结果可见于 \cite{GruRay}。粗略地说，它断言：在源与目标上
\'etale 局部地，一个有限型态射是一个有限态射与一个具有几何连通纤维的
光滑态射之复合，而该光滑态射的相对维数等于原态射的纤维维数。

\begin{lemma}
\label{more-morphisms-lemma-local-structure-finite-type}
设 $f : X \to S$ 是态射。取 $x \in X$，并令 $s = f(x)$。
假设 $f$ 局部有限型，且 $n = \dim_x(X_s)$。则存在交换图表
$$
\xymatrix{
X \ar[dd] & X' \ar[l]^g \ar[d]^\pi & x \ar@{|->}[dd] &
x' \ar@{|->}[l]  \ar@{|->}[d] \\
& Y \ar[d]^h & & y \ar@{|->}[d] \\
S \ar@{=}[r] & S & s & s \ar@{=}[l]
}
$$
以及满足 $g(x') = x$ 的一点 $x' \in X'$，使得令 $y = \pi(x')$ 后，
\begin{enumerate}
\item $h : Y \to S$ 是相对维数为 $n$ 的光滑态射；
\item $g : (X', x') \to (X, x)$ 是初等 \'etale 邻域；
\item $\pi$ 有限，且 $\pi^{-1}(\{y\}) = \{x'\}$；
\item $\kappa(y)$ 是 $\kappa(s)$ 的纯超越扩张。
\end{enumerate}
此外，若 $f$ 局部有限表示，则 $\pi$ 是有限表示的。
\end{lemma}

\begin{proof}
问题在 $X$ 和 $S$ 上是局部的，故可假设 $X$ 与 $S$ 均仿射。
由 代数章引理 \ref{algebra-lemma-refined-quasi-finite-over-polynomial-algebra}，
用 $x$ 在 $X$ 中的一个标准开邻域替换 $X$ 后，可假设存在分解
$$
\xymatrix{
X \ar[r]^\pi & \mathbf{A}^n_S \ar[r] & S
}
$$
使得 $\pi$ 拟有限，并且 $\kappa(\pi(x))$ 是 $\kappa(s)$ 的纯超越扩张。
由 引理 \ref{more-morphisms-lemma-etale-makes-quasi-finite-finite-at-point}，
存在初等 \'etale 邻域
$$
(Y, y) \to (\mathbf{A}^n_S, \pi(x))
$$
以及开集 $X' \subset X \times_{\mathbf{A}^n_S} Y$；它包含唯一一个
位于 $y$ 上方的点 $x'$，并且 $X' \to Y$ 有限。
这证明了 (1) -- (4)。最后一个断言由
态射章引理 \ref{morphisms-lemma-finite-presentation-permanence} 得到。
\end{proof}

\begin{lemma}
\label{more-morphisms-lemma-local-local-structure-finite-type}
\begin{slogan}
在 \'etale 邻域中，有限型态射是光滑态射上的有限态射。
\end{slogan}
设 $f : X \to S$ 是态射。取 $x \in X$，并令 $s = f(x)$。
假设 $f$ 局部有限型，且 $n = \dim_x(X_s)$。则存在交换图表
$$
\xymatrix{
X \ar[dd] & X' \ar[l]^g \ar[d]^\pi & x \ar@{|->}[dd] &
x' \ar@{|->}[l] \ar@{|->}[d] \\
& Y' \ar[d]^h & & y' \ar@{|->}[d] \\
S & S' \ar[l]_e & s & s' \ar@{|->}[l]
}
$$
以及满足 $g(x') = x$ 的一点 $x' \in X'$，使得令 $y' = \pi(x')$、
$s' = h(y')$ 后，
\begin{enumerate}
\item $h : Y' \to S'$ 是相对维数为 $n$ 的光滑态射；
\item $Y' \to S'$ 的所有纤维均几何整；
\item $g : (X', x') \to (X, x)$ 是初等 \'etale 邻域；
\item $\pi$ 有限，且 $\pi^{-1}(\{y'\}) = \{x'\}$；
\item $\kappa(y')$ 是 $\kappa(s')$ 的纯超越扩张；
\item $e : (S', s') \to (S, s)$ 是初等 \'etale 邻域。
\end{enumerate}
此外，若 $f$ 局部有限表示，则 $\pi$ 是有限表示的。
\end{lemma}

\begin{proof}
问题在 $S$ 上是局部的，故可用 $s$ 的一个仿射开邻域替换 $S$。
接着应用 引理 \ref{more-morphisms-lemma-local-structure-finite-type}，得到交换图表
$$
\xymatrix{
X \ar[dd] & X' \ar[l]^g \ar[d]^\pi & x \ar@{|->}[dd] &
x' \ar@{|->}[l] \ar@{|->}[d] \\
& Y \ar[d]^h & & y \ar@{|->}[d] \\
S \ar@{=}[r] & S & s & s \ar@{=}[l]
}
$$
其中 $h$ 是相对维数为 $n$ 的光滑态射，而 $\kappa(y)$ 是
$\kappa(s)$ 的纯超越扩张。由于问题在 $X$ 上也是局部的，
可以用 $y$ 的一个仿射邻域替换 $Y$（并用其在 $\pi$ 下的逆像替换
$X'$）。由于 $S$ 仿射，这保证 $Y \to S$ 拟紧、分离且光滑，
特别是有限表示的。令 $T$ 为 $Y_s$ 中含有 $y$ 的连通分支。
由于 $Y_s$ 是 Noether 的，可见 $T$ 是开的。由 Varieties,
引理 \ref{varieties-lemma-geometrically-connected-if-connected-and-point}，
$T$ 在 $\kappa(s)$ 上几何连通。又因为 $T$ 在 $\kappa(s)$ 上光滑，
所以它几何正规；见
簇章引理 \ref{varieties-lemma-smooth-geometrically-normal}。
由此断定 $T$ 在 $\kappa(s)$ 上几何不可约
（因为连通的 Noether 正规概形不可约；见
性质章引理 \ref{properties-lemma-normal-Noetherian}）。
最后，由 引理 \ref{more-morphisms-lemma-smooth-normal}，光滑态射 $h$ 是正规态射。
% P05-EN-0100: the frozen source has singular "all assumption"; Chinese follows the intended plural.
至此已经验证：对态射 $h : Y \to S$ 和开集 $T \subset Y_s$，
引理 \ref{more-morphisms-lemma-normal-morphism-irreducible} 的全部假设均成立。
应用 引理 \ref{more-morphisms-lemma-normal-morphism-irreducible} 后，得到
$e : S' \to S$、$s' \in S'$ 和 $Y'$，如交换图表
$$
\xymatrix{
X \ar[dd] & X' \ar[l]^g \ar[d]^\pi & X' \times_Y Y' \ar[l] \ar[d] &
x \ar@{|->}[dd] & x' \ar@{|->}[l] \ar@{|->}[d] &
(x', s') \ar@{|->}[l] \ar@{|->}[d] \\
& Y \ar[d]^h & Y' \ar[d] \ar[l] & & y \ar@{|->}[d] &
(y, s') \ar@{|->}[l] \ar@{|->}[d] \\
S \ar@{=}[r] & S & S' \ar[l]_e & s & s \ar@{=}[l] & s' \ar@{|->}[l]
}
$$
所示，其中 $e : (S', s') \to (S, s)$ 是初等 \'etale 邻域，
而 $Y' \subset Y_{S'}$ 是一个开邻域，其在 $S'$ 上的所有纤维均几何不可约，
并且在等同 $Y_s = Y_{S', s'}$ 下有 $Y'_{s'} = T$。
令 $(y, s') \in Y'$ 为对应于 $y \in T$ 的点；它也是
$Y \times_S S'$ 中位于 $y$ 上方、剩余域等于 $\kappa(y)$ 且映到
$S'$ 中 $s'$ 的唯一一点。类似地，令
$(x', s') \in X' \times_Y Y' \subset X' \times_S S'$ 为位于
$x'$ 上方、剩余域等于 $\kappa(x')$ 且位于 $s'$ 上方的唯一一点。
于是，该图表的外部部分给出了本引理所提问题的解。略去若干次要细节。
\end{proof}

\begin{lemma}
\label{more-morphisms-lemma-local-local-structure-finite-type-affine}
假设与记号同 引理 \ref{more-morphisms-lemma-local-local-structure-finite-type}。
除性质 (1) -- (6) 外，还可以安排使得
\begin{enumerate}
\item[(7)] $S'$、$Y'$、$X'$ 均仿射。
\end{enumerate}
\end{lemma}

\begin{proof}
注意，若 $Y'$ 仿射，则由于 $\pi$ 有限，$X'$ 也仿射。
选取 $s'$ 的仿射开邻域 $U' \subset S'$，再选取 $y'$ 的仿射开邻域
$V' \subset h^{-1}(U')$。令 $W' = h(V')$。由
态射章引理 \ref{morphisms-lemma-smooth-open}，
这是 $S'$ 中包含于 $U'$ 的 $s'$ 的开邻域。选取 $s'$ 的仿射开邻域
$U'' \subset W'$。于是 $h^{-1}(U'') \cap V'$ 仿射，因为它等于
$U'' \times_{U'} V'$。按构造，
$h^{-1}(U'') \cap V' \to U''$ 是满光滑态射，其纤维是
$Y' \to S'$ 的几何整纤维的（非空）开子概形，因而也是几何整的。
所以可以用 $U''$ 替换 $S'$，并用 $h^{-1}(U'') \cap V'$ 替换 $Y'$。
\end{proof}

\noindent
性质 $\pi^{-1}(\{y'\}) = \{x'\}$ 的意义由下述引理部分说明。

\begin{lemma}
\label{more-morphisms-lemma-finite-morphism-single-point-in-fibre}
设 $\pi : X \to Y$ 是有限态射。取 $x \in X$，令 $y = \pi(x)$，
并假设 $\pi^{-1}(\{y\}) = \{x\}$。则
\begin{enumerate}
\item 对 $x$ 在 $X$ 中的每个邻域 $U \subset X$，都存在 $y$
的邻域 $V \subset Y$，使得 $\pi^{-1}(V) \subset U$。
\item 环映射 $\mathcal{O}_{Y, y} \to \mathcal{O}_{X, x}$ 是有限的。
\item 若 $\pi$ 是有限表示的，则
$\mathcal{O}_{Y, y} \to \mathcal{O}_{X, x}$ 是有限表示的。
\item 对任意拟凝聚 $\mathcal{O}_X$-模 $\mathcal{F}$，作为
$\mathcal{O}_{Y, y}$-模，有 $\mathcal{F}_x = \pi_*\mathcal{F}_y$。
\end{enumerate}
\end{lemma}

\begin{proof}
第一个断言纯属拓扑性质；只需使用 $\pi$ 是连续闭映射并且
$\pi^{-1}(\{y\}) = \{x\}$。为证明第二、第三部分，可以假设
$X = \Spec(B)$ 且 $Y = \Spec(A)$。此时 $A \to B$ 是有限环映射，
而 $y$ 对应于 $A$ 的素理想 $\mathfrak p$，并且 $B$ 中恰有一个
位于 $\mathfrak p$ 上方的素理想 $\mathfrak q$。于是
$B_{\mathfrak q} = B_{\mathfrak p}$；见
代数章引理 \ref{algebra-lemma-unique-prime-over-localize-below}。
换言之，映射 $A_{\mathfrak p} \to B_{\mathfrak q}$ 等于将
$A \to B$ 在 $\mathfrak p$ 处局部化所得的映射
$A_{\mathfrak p} \to B_{\mathfrak p}$。因此，(2) 和 (3) 由局部化的
简单性质得到（略去若干细节）。对最后一个断言，设对某个 $B$-模 $M$
有 $\mathcal{F} = \widetilde M$。
% P05-EN-0101: frozen source has F = M_q here; the stalk F_x is evidently intended.
则 $\mathcal{F} = M_{\mathfrak q}$ 且
$\pi_*\mathcal{F}_y = M_{\mathfrak p}$。由上述讨论，这两个局部化相同。
也可以直接利用 (1) 和茎的定义看出
$\mathcal{F}_x = \pi_*\mathcal{F}_y$。
\end{proof}






\section{fppf 拓扑的应用}
\label{more-morphisms-section-application-fppf}

\noindent
可以利用上述 \'etale 局部化技术证明下述结果：fppf 拓扑等于由 Zariski
覆盖以及形如 $\{f : T \to S\}$ 的覆盖“生成”的拓扑，其中 $f$
是满的有限局部自由态射。

\begin{lemma}
\label{more-morphisms-lemma-dominate-fppf-etale-locally}
设 $S$ 是概形，$\{S_i \to S\}_{i \in I}$ 是 fppf 覆盖。则存在
\begin{enumerate}
\item 一个 \'etale 覆盖 $\{S'_a \to S\}$；
\item 满的有限局部自由态射 $V_a \to S'_a$；
\end{enumerate}
使得 fppf 覆盖 $\{V_a \to S\}$ 加细给定覆盖 $\{S_i \to S\}$。
\end{lemma}

\begin{proof}
可以假设每个 $S_i \to S$ 都局部拟有限；见
引理 \ref{more-morphisms-lemma-qf-fp-flat-dominates-fppf}。

\medskip\noindent
固定一点 $s \in S$。选取 $i \in I$ 以及映到 $s$ 的一点
$s_i \in S_i$。选取初等 \'etale 邻域
$(S', s) \to (S, s)$，使得存在开集
$$
S_i \times_S S'  \supset V
$$
它包含唯一一个映到 $s \in S'$ 的点 $v \in V$，并且
$V \to S'$ 有限；见
引理 \ref{more-morphisms-lemma-etale-makes-quasi-finite-finite-at-point}。
于是 $V \to S'$ 有限局部自由：它是有限的，而
$S_i \times_S S' \to S'$ 作为 $S_i \to S$ 的基变换平坦且局部有限表示；见
态射章引理 \ref{morphisms-lemma-base-change-finite-presentation},
\ref{morphisms-lemma-base-change-flat}, 及
\ref{morphisms-lemma-finite-flat}。
因此 $V \to S'$ 是开态射；缩小 $S'$ 后，可以假设 $V \to S'$
是满的有限局部自由态射。由于对 $S$ 的每一点都可以这样做，断定
$\{S_i \to S\}$ 可由形如 $\{V_a \to S\}_{a \in A}$ 的覆盖加细，
其中每个 $V_a \to S$ 分解为 $V_a \to S'_a \to S$，
$S'_a \to S$ 为 \'etale 态射，而 $V_a \to S'_a$ 为满的有限局部自由态射。
\end{proof}

\begin{lemma}
\label{more-morphisms-lemma-dominate-fppf}
设 $S$ 是概形，$\{S_i \to S\}_{i \in I}$ 是 fppf 覆盖。则存在
\begin{enumerate}
\item Zariski 开覆盖 $S = \bigcup U_j$；
\item 满的有限局部自由态射 $W_j \to U_j$；
\item Zariski 开覆盖 $W_j = \bigcup_k W_{j, k}$；
\item 满的有限局部自由态射 $T_{j, k} \to W_{j, k}$；
\end{enumerate}
使得 fppf 覆盖 $\{T_{j, k} \to S\}$ 加细给定覆盖
$\{S_i \to S\}$。
\end{lemma}

\begin{proof}
令 $\{V_a \to S\}_{a \in A}$ 为
引理 \ref{more-morphisms-lemma-dominate-fppf-etale-locally} 中得到的 fppf 覆盖。
换言之，该覆盖加细 $\{S_i \to S\}$，并且每个 $V_a \to S$
分解为 $V_a \to S'_a \to S$，其中 $S'_a \to S$ 是 \'etale 态射，
$V_a \to S'_a$ 是满的有限局部自由态射。

\medskip\noindent
由 说明 \ref{more-morphisms-remark-topologies}，存在 Zariski 开覆盖
$S = \bigcup U_j$；对每个 $j$，存在有限局部自由满态射
$W_j \to U_j$；并且对每个 $j$，存在 Zariski 开覆盖
$\{W_{j, k} \to W_j\}$，使得族 $\{W_{j, k} \to S\}$ 加细
\'etale 覆盖 $\{S'_a \to S\}$。也就是说，对每一对 $j, k$，
存在 $a(j, k)$ 以及 $W_{j, k} \to S$ 的分解
$W_{j, k} \to S'_a \to S$。令
$T_{j, k} = W_{j, k} \times_{S'_a} V_a$，一切即清楚。
\end{proof}

\begin{lemma}
\label{more-morphisms-lemma-extend-integral-surjective-morphisms}
设 $S$ 是概形。若 $U \subset S$ 是开集且 $V \to U$ 是满整态射，
则存在满整态射 $\overline{V} \to S$，使得作为 $U$ 上的概形，
$\overline{V} \times_S U$ 同构于 $V$。
\end{lemma}

\begin{proof}
% P05-EN-0102: the frozen source says normalization of S in U, but the claimed inverse image is V;
% the construction evidently requires normalization/integral closure of S in V. Rendered formula/variables retained.
令 $V' \to S$ 为 $S$ 在 $U$ 中的正规化；见
态射章第 \ref{morphisms-section-normalization-X-in-Y} 节。
按构造，$V' \to S$ 是整态射。由 态射章引理
\ref{morphisms-lemma-normalization-localization} 及
\ref{morphisms-lemma-normalization-in-integral}，可见 $U$ 在 $V'$ 中的逆像是
$V$。令 $Z$ 为赋予 $S \setminus U$ 的既约诱导概形结构。于是
$\overline{V} = V' \amalg Z$ 即满足要求。
\end{proof}

\begin{lemma}
\label{more-morphisms-lemma-extend-finite-surjective-morphisms}
设 $S$ 是拟紧拟分离概形。若 $U \subset S$ 是拟紧开集，且
$V \to U$ 是满有限态射，则存在满有限态射
$\overline{V} \to S$，使得作为 $U$ 上的概形，
$\overline{V} \times_S U$ 同构于 $V$。
\end{lemma}

\begin{proof}
由 Zariski 主定理
（引理 \ref{more-morphisms-lemma-quasi-finite-separated-pass-through-finite}），
可以假设 $V$ 是某个在 $S$ 上有限的概形 $V'$ 中的拟紧开集。
用 $V$ 的概形论像替换 $V'$ 后，可以假设 $V$ 在 $V'$ 中稠密。
于是 $V' \times_S U = V$，因为 $V$ 在 $U$ 上有限，故
$V \to V' \times_S U$ 是闭的。令 $Z$ 为赋予
$S \setminus U$ 的既约诱导概形结构。则
$\overline{V} = V' \amalg Z$ 即满足要求。
\end{proof}

\begin{lemma}
\label{more-morphisms-lemma-dominate-fppf-integral}
设 $S$ 是概形，$\{S_i \to S\}_{i \in I}$ 是 fppf 覆盖。
则存在满整态射 $S' \to S$ 和开覆盖 $S' = \bigcup U'_\alpha$，
使得对每个 $\alpha$，态射 $U'_\alpha \to S$ 经由某个 $i$ 所对应的
$S_i \to S$ 分解。
\end{lemma}

\begin{proof}
按 引理 \ref{more-morphisms-lemma-dominate-fppf} 选取
$S = \bigcup U_j$、$W_j \to U_j$、$W_j = \bigcup W_{j, k}$ 以及
$T_{j, k} \to W_{j, k}$。由
引理 \ref{more-morphisms-lemma-extend-integral-surjective-morphisms}，可以将
$W_j \to U_j$ 延拓为满整态射 $\overline{W}_j \to S$。
之后，可以将 $T_{j, k} \to W_{j, k}$ 延拓为满整态射
$\overline{T}_{j, k} \to \overline{W}_j$。
令 $\overline{T}_j$ 等于所有概形 $\overline{T}_{j, k}$ 在
$\overline{W}_j$ 上的积
（极限章引理 \ref{limits-lemma-infinite-product}）。
再令 $S'$ 等于所有概形 $\overline{T}_j$ 在 $S$ 上的积。
若 $x \in S'$，则存在 $j$，使得 $x$ 在 $S$ 中的像属于 $U_j$。
因此存在 $k$，使得 $x$ 在投影 $S' \to \overline{W}_j$ 下的像属于
$W_{j, k}$。从而在投影
$S' \to \overline{T}_j \to \overline{T}_{j, k}$ 下，点 $x$ 落在
$T_{j, k}$ 中，而 $T_{j, k} \to S$ 经由某个 $i$ 所对应的 $S_i$ 分解。
最后，由 极限章引理 \ref{limits-lemma-infinite-product-integral} 及
\ref{limits-lemma-infinite-product-surjective}，态射 $S' \to S$ 整且满。
\end{proof}

\begin{lemma}
\label{more-morphisms-lemma-dominate-fppf-finite}
设 $S$ 是拟紧拟分离概形，$\{S_i \to S\}_{i \in I}$ 是 fppf 覆盖。
则存在有限表示的满有限态射 $S' \to S$ 以及开覆盖
$S' = \bigcup U'_\alpha$，使得对每个 $\alpha$，态射
$U'_\alpha \to S$ 经由某个 $i$ 所对应的 $S_i \to S$ 分解。
\end{lemma}

\begin{proof}
% P05-EN-0103: the frozen proof switches from the statement's base S to undefined X throughout.
% Variables are retained; the evident intended base is S.
令 $Y \to X$ 为 引理 \ref{more-morphisms-lemma-dominate-fppf-integral} 中得到的整满态射。
选取有限仿射开覆盖 $Y = \bigcup V_j$，使得
$V_j \to X$ 经由 $S_{i(j)}$ 分解。可以写成 $Y = \lim Y_\lambda$，
其中 $Y_\lambda \to X$ 有限且有限表示；见
极限章引理 \ref{limits-lemma-integral-limit-finite-and-finite-presentation}。
当 $\lambda$ 充分大时，可以找到仿射开集
$V_{\lambda, j} \subset Y_\lambda$，其在 $Y$ 中的逆像恰为 $V_j$；见
极限章引理 \ref{limits-lemma-descend-opens}。
进一步增大 $\lambda$，则 $X$ 上的态射 $V_j \to S_{i(j)}$ 来自
$X$ 上的态射 $V_{\lambda, j} \to S_{i(j)}$；见 极限章命题
\ref{limits-proposition-characterize-locally-finite-presentation}。
对这个 $\lambda$ 令 $S' = Y_\lambda$，证明完毕。
\end{proof}

\begin{lemma}
\label{more-morphisms-lemma-fppf-ph}
概形的 fppf 覆盖是 ph 覆盖。
\end{lemma}

\begin{proof}
设 $\{T_i \to T\}$ 是概形的 fppf 覆盖；见
拓扑章定义 \ref{topologies-definition-fppf-covering}。
注意，$T_i \to T$ 局部有限型。设 $U \subset T$ 是仿射开集。
只需说明 $\{T_i \times_T U \to U\}$ 可由一个标准 ph 覆盖加细；见
拓扑章定义 \ref{topologies-definition-ph-covering}。
这立即由 引理 \ref{more-morphisms-lemma-dominate-fppf-finite} 以及有限态射为固有态射这一事实得到
（态射章引理 \ref{morphisms-lemma-finite-proper}）。
\end{proof}

\begin{remark}
\label{more-morphisms-remark-change-topologies}
作为 引理 \ref{more-morphisms-lemma-fppf-ph} 的推论，得到比较态射
$$
\epsilon : (\Sch/S)_{ph} \longrightarrow (\Sch/S)_{fppf}
$$
这是由底层范畴上的恒等函子给出的位点态射
（按 拓扑章说明 \ref{topologies-remark-choice-sites} 适当选取位点）。
函子 $\epsilon_*$ 在底层预层上是恒等函子，而函子
% P05-EN-0104: frozen source omits the verb in "epsilon^{-1} associated ... its ph sheafification".
$\epsilon^{-1}$ 将一个 fppf 层对应到其 ph 层化。
通过复合，还可以将 ph 拓扑与 syntomic、光滑、\'etale 及 Zariski
拓扑进行比较。
\end{remark}




\section{拟投射概形}
\label{more-morphisms-section-quasi-projective}

\noindent
术语“拟投射概形”尚未定义。一种可能的定义是：具有充足可逆层的概形。
然而，若 $X$ 是基概形 $S$ 上的概形，则当态射 $X \to S$ 拟投射时，
我们称{\it $X$ 在 $S$ 上拟投射}
（态射章定义 \ref{morphisms-definition-quasi-projective}）。
由于任意概形的恒等态射都是拟投射的，可见在 $S$ 上拟投射的概形未必具有
充足可逆层。因此，最好仍不定义术语“拟投射概形”。

\begin{lemma}
\label{more-morphisms-lemma-quasi-projective}
设 $S$ 是具有充足可逆层的概形，$f : X \to S$ 是概形态射。
下列条件等价：
\begin{enumerate}
\item $X \to S$ 拟投射；
\item $X \to S$ 是 H-拟投射的；
\item 存在 $S$ 上概形的拟紧开浸入 $X \to X'$，并且
$X' \to S$ 投射；
\item $X \to S$ 有限型，且 $X$ 具有充足可逆层；
\item $X \to S$ 有限型，且存在一个 $f$-极充足可逆层。
\end{enumerate}
\end{lemma}

\begin{proof}
蕴含 (2) $\Rightarrow$ (1) 是
态射章引理 \ref{morphisms-lemma-H-quasi-projective-quasi-projective}。
蕴含 (1) $\Rightarrow$ (2) 是 态射章引理
\ref{morphisms-lemma-projective-over-quasi-projective-is-H-projective}。
蕴含 (2) $\Rightarrow$ (3) 是
态射章引理 \ref{morphisms-lemma-H-quasi-projective-open-H-projective}

\medskip\noindent
假设 $X \subset X'$ 如 (3) 所述。特别地，$X \to S$ 有限型。
% P05-EN-0105: frozen source calls X -> S H-projective here, while the cited lemma and ensuing immersion
% characterize H-quasi-projective; rendered terminology follows the frozen source.
由 态射章引理 \ref{morphisms-lemma-H-quasi-projective-open-H-projective}，
态射 $X \to S$ 是 H-投射的。因此存在拟紧浸入
$i : X \to \mathbf{P}^n_S$。于是
$\mathcal{L} = i^*\mathcal{O}_{\mathbf{P}^n_S}(1)$ 是 $f$-极充足的。
由于 $X \to S$ 拟紧，由
态射章引理 \ref{morphisms-lemma-ample-very-ample}，
$\mathcal{L}$ 是 $f$-充足的。因此按定义 $X \to S$ 拟投射。

\medskip\noindent
蕴含 (4) $\Rightarrow$ (2) 是
态射章引理 \ref{morphisms-lemma-quasi-projective-finite-type-over-S}.

\medskip\noindent
蕴含 (5) $\Rightarrow$ (4) 由
态射章引理 \ref{morphisms-lemma-pullback-ample-tensor-relatively-ample}.

\medskip\noindent
条件 (1) 与 (5) 的等价性由
态射章引理 \ref{morphisms-lemma-ample-very-ample} 及
\ref{morphisms-lemma-finite-type-ample-very-ample}.
\end{proof}

\begin{lemma}
\label{more-morphisms-lemma-category-quasi-projective}
设 $S$ 是具有充足可逆层的概形。令 $\text{QP}_S$ 为 $S$ 上概形范畴的
全子范畴，其对象满足 引理 \ref{more-morphisms-lemma-quasi-projective} 中的等价条件。
\begin{enumerate}
\item 若 $S' \to S$ 是概形态射且 $S'$ 具有充足可逆层，则基变换确定函子
$\text{QP}_S \to \text{QP}_{S'}$；
\item 若 $X \in \text{QP}_S$ 且 $Y \in \text{QP}_X$，则 $Y \in \text{QP}_S$；
\item 范畴 $\text{QP}_S$ 对纤维积封闭；
\item 范畴 $\text{QP}_S$ 对有限不交并封闭；
\item 若 $X \to S$ 投射，则 $X \in \text{QP}_S$；
\item 若 $X \to S$ 是有限型拟仿射态射，则 $X$ 属于 $\text{QP}_S$；
\item 若 $X \to S$ 拟有限且分离，则 $X \in \text{QP}_S$；
\item 若 $X \to S$ 是拟紧浸入，则 $X \in \text{QP}_S$；
% P05-EN-0106: the frozen source contains this editorial placeholder as a numbered mathematical item.
\item 此处添加更多内容。
\end{enumerate}
\end{lemma}

\begin{proof}
第 (1) 部分由 态射章引理
\ref{morphisms-lemma-base-change-quasi-projective}.

\medskip\noindent
第 (2) 部分由 引理 \ref{more-morphisms-lemma-quasi-projective} 的第四种刻画得到。

\medskip\noindent
若 $X \to S$ 与 $Y \to S$ 拟投射，则由
态射章引理 \ref{morphisms-lemma-base-change-quasi-projective}，
$X \times_S Y \to Y$ 拟投射。因此 (3) 由 (2) 得到。

\medskip\noindent
若 $X = Y \amalg Z$ 是概形的不交并，而 $\mathcal{L}$ 是可逆
$\mathcal{O}_X$-模，并且 $\mathcal{L}|_Y$ 与 $\mathcal{L}|_Z$
均充足，则 $\mathcal{L}$ 充足（略去细节）。因此，第 (4) 部分由
引理 \ref{more-morphisms-lemma-quasi-projective} 的第四种刻画得到。

\medskip\noindent
第 (5) 部分由
态射章引理 \ref{morphisms-lemma-projective-quasi-projective}.

\medskip\noindent
第 (6) 部分由
态射章引理
\ref{morphisms-lemma-quasi-affine-finite-type-quasi-projective}.

\medskip\noindent
第 (7) 部分由第 (6) 部分及
引理 \ref{more-morphisms-lemma-quasi-finite-separated-quasi-affine}.

\medskip\noindent
第 (8) 部分由第 (7) 部分及
态射章引理 \ref{morphisms-lemma-immersion-locally-quasi-finite}.
\end{proof}

\noindent
下述引理其实并不属于本节，但在别处似乎也没有合适的位置。

\begin{lemma}
\label{more-morphisms-lemma-integral-over-quasi-affine}
设 $X$ 是拟仿射概形，$f : U \to X$ 是整态射。
则 $U$ 拟仿射，并且图表
$$
\xymatrix{
U \ar[r] \ar[d] & \Spec(\Gamma(U, \mathcal{O}_U)) \ar[d] \\
X \ar[r] & \Spec(\Gamma(X, \mathcal{O}_X))
}
$$
是笛卡尔的。
\end{lemma}

\begin{proof}
概形 $U$ 拟仿射，因为整态射是仿射的，仿射态射是拟仿射的；一个概形
拟仿射，当且仅当它到 $\Spec(\mathbf{Z})$ 的结构态射拟仿射；并且
拟仿射态射的复合拟仿射。前两个陈述立即由定义得到，第三个是
态射章引理 \ref{morphisms-lemma-composition-quasi-affine}。
令 $U' =
X \times_{\Spec(\Gamma(X, \mathcal{O}_X))} \Spec(\Gamma(U, \mathcal{O}_U))$
，并考虑扩展图表
$$
\xymatrix{
U \ar[r]_j \ar[rd] & U' \ar[d] \ar[r] &
\Spec(\Gamma(U, \mathcal{O}_U)) \ar[d] \\
& X \ar[r] & \Spec(\Gamma(X, \mathcal{O}_X))
}
$$
由 态射章引理 \ref{morphisms-lemma-image-proper-scheme-closed}，
结合整态射普遍闭这一事实
（态射章引理 \ref{morphisms-lemma-integral-universally-closed}）
% P05-EN-0107: frozen source reads "vertical arrows are in the diagram are separated".
以及图表中的竖直箭头分离这一事实，态射 $j$ 是闭的。
另一方面，由 性质章引理 \ref{properties-lemma-quasi-affine}，
引理图表中的水平箭头是开的，故 $j$ 是开的。因此，$j$ 将 $U$
等同于 $U'$ 的一个开闭子概形。若 $U \not = U'$，则 $U$ 在 $U'$
中不稠密，更不可能在 $\Gamma(U, \mathcal{O}_U)$ 的谱中稠密。
然而，$U$ 在 $\Spec(\Gamma(U, \mathcal{O}_U))$ 中的概形论像是
$\Spec(\Gamma(U, \mathcal{O}_U))$，因为
$\Gamma(U, \mathcal{O}_U)$ 中任何切出一个使 $U$ 的态射经其分解的
闭子概形的理想都必须为零。因此，例如由
态射章引理 \ref{morphisms-lemma-quasi-compact-scheme-theoretic-image}，
$U$ 在 $\Spec(\Gamma(U, \mathcal{O}_U))$ 中稠密。
所以 $U = U'$，结论得证。
\end{proof}









\section{投射概形}
\label{more-morphisms-section-projective}

\noindent
本节是 第 \ref{more-morphisms-section-quasi-projective} 节 针对投射态射的对应版本。

\begin{lemma}
\label{more-morphisms-lemma-projective}
设 $S$ 是具有充足可逆层的概形，$f : X \to S$ 是概形态射。
下列条件等价：
\begin{enumerate}
\item $X \to S$ 投射；
\item $X \to S$ 是 H-投射的；
\item $X \to S$ 拟投射且固有；
\item $X \to S$ 是 H-拟投射且固有的；
\item $X \to S$ 固有，且 $X$ 具有充足可逆层；
\item $X \to S$ 固有，且存在一个 $f$-充足可逆层；
\item $X \to S$ 固有，且存在一个 $f$-极充足可逆层；
\item 存在拟凝聚分次 $\mathcal{O}_S$-代数 $\mathcal{A}$，它由
$\mathcal{A}_1$ 在 $\mathcal{A}_0$ 上生成，其中 $\mathcal{A}_1$
是有限型 $\mathcal{O}_S$-模，并且
$X = \underline{\text{Proj}}_S(\mathcal{A})$. 
\end{enumerate}
\end{lemma}

\begin{proof}
首先注意，在每种情形下态射 $f$ 都是固有的；见
态射章引理 \ref{morphisms-lemma-H-projective} 及
\ref{morphisms-lemma-locally-projective-proper}。
因此，只需在 $f$ 为固有态射的情形证明这些概念等价。
以下将直接使用这一点而不再说明。

\medskip\noindent
等价关系 (1) $\Leftrightarrow$ (3) 和 (2) $\Leftrightarrow$ (4) 是
态射章引理 \ref{morphisms-lemma-projective-is-quasi-projective-proper}.

\medskip\noindent
蕴含 (2) $\Rightarrow$ (1) 是
态射章引理 \ref{morphisms-lemma-H-projective}.

\medskip\noindent
蕴含 (1) $\Rightarrow$ (2) 和 (3) $\Rightarrow$ (4) 是
态射章引理
\ref{morphisms-lemma-projective-over-quasi-projective-is-H-projective}.

\medskip\noindent
蕴含 (1) $\Rightarrow$ (7) 立即由
态射章定义 \ref{morphisms-definition-projective} 及
\ref{morphisms-definition-very-ample}.

\medskip\noindent
由 态射章定义 \ref{morphisms-definition-quasi-projective}，
条件 (3) 和 (6) 等价。

\medskip\noindent
因此，(1) -- (4)、(6) 等价且蕴含 (7)。由
引理 \ref{more-morphisms-lemma-quasi-projective}，条件 (3)、(5) 和 (7) 等价。
由此可见，(1) -- (7) 等价。

\medskip\noindent
由 除子章引理 \ref{divisors-lemma-relative-proj-projective}，
(8) 蕴含 (1)。反之，若 (2) 成立，则可选取闭浸入
$$
i :
X
\longrightarrow
\mathbf{P}^n_S = \underline{\text{Proj}}_S(\mathcal{O}_S[T_0, \ldots, T_n]).
$$
该等式见 构造章引理 \ref{constructions-lemma-projective-space-bundle}。
由 除子章引理 \ref{divisors-lemma-closed-subscheme-proj}，
$X$ 是 $\mathcal{O}_S[T_0, \ldots, T_n]$ 的某个拟凝聚分次商代数
$\mathcal{A}$ 的相对 Proj。于是 $\mathcal{A}$ 满足 (8) 的条件。
\end{proof}

\begin{lemma}
\label{more-morphisms-lemma-category-projective}
设 $S$ 是具有充足可逆层的概形。令 $\text{P}_S$ 为 $S$ 上概形范畴的
全子范畴，其对象满足 引理 \ref{more-morphisms-lemma-projective} 中的等价条件。
\begin{enumerate}
\item 若 $S' \to S$ 是概形态射且 $S'$ 具有充足可逆层，则基变换确定函子
$\text{P}_S \to \text{P}_{S'}$；
\item 若 $X \in \text{P}_S$ 且 $Y \in \text{P}_X$，则 $Y \in \text{P}_S$；
\item 范畴 $\text{P}_S$ 对纤维积封闭；
\item 范畴 $\text{P}_S$ 对有限不交并封闭；
\item 若 $X \to S$ 有限，则 $X$ 属于 $\text{P}_S$；
% P05-EN-0108: the frozen source contains this editorial placeholder as a numbered mathematical item.
\item 此处添加更多内容。
\end{enumerate}
\end{lemma}

\begin{proof}
第 (1) 部分由 态射章引理
\ref{morphisms-lemma-base-change-projective}.

\medskip\noindent
第 (2) 部分由 引理 \ref{more-morphisms-lemma-projective} 的第五种刻画，
以及固有态射的复合仍固有这一事实得到
（态射章引理 \ref{morphisms-lemma-composition-proper}）。

\medskip\noindent
若 $X \to S$ 与 $Y \to S$ 投射，则由
态射章引理 \ref{morphisms-lemma-base-change-projective}，
$X \times_S Y \to Y$ 投射。因此 (3) 由 (2) 得到。

\medskip\noindent
若 $X = Y \amalg Z$ 是概形的不交并，而 $\mathcal{L}$ 是可逆
$\mathcal{O}_X$-模，并且 $\mathcal{L}|_Y$ 与 $\mathcal{L}|_Z$
均充足，则 $\mathcal{L}$ 充足（略去细节）。因此，第 (4) 部分由
引理 \ref{more-morphisms-lemma-projective} 的第五种刻画得到。

\medskip\noindent
第 (5) 部分由
态射章引理 \ref{morphisms-lemma-finite-projective}.
\end{proof}

\noindent
下面给出一个类型稍有不同的结果。

\begin{lemma}
\label{more-morphisms-lemma-ample-in-neighbourhood}
\begin{reference}
\cite[IV Corollary 9.6.4]{EGA}
\end{reference}
设 $f : X \to Y$ 是概形的固有态射，$\mathcal{L}$ 是可逆
$\mathcal{O}_X$-模。设 $y \in Y$ 是一点，并且 $\mathcal{L}_y$
在 $X_y$ 上充足。则存在 $y$ 的开邻域 $V \subset Y$，使得
$\mathcal{L}|_{f^{-1}(V)}$ 在 $f^{-1}(V)/V$ 上充足。
\end{lemma}

\begin{proof}
可以假设 $Y$ 仿射。于是找到有向集 $I$ 以及概形态射的逆系统
$X_i \to Y_i$，其中 $Y_i$ 在 $\mathbf{Z}$ 上有限型，转移态射
$X_i \to X_{i'}$ 与 $Y_i \to Y_{i'}$ 仿射，$X_i \to Y_i$ 固有，
并且 $X \to Y = \lim (X_i \to Y_i)$。见 极限章引理
\ref{limits-lemma-proper-limit-of-proper-finite-presentation-noetherian}。
缩小 $I$ 后，可以假设有一族相容的可逆 $\mathcal{O}_{X_i}$-模
$\mathcal{L}_i$，其拉回为 $\mathcal{L}$；见
极限章引理 \ref{limits-lemma-descend-invertible-modules}。
令 $y_i \in Y_i$ 为 $y$ 的像。则
$\kappa(y) = \colim \kappa(y_i)$。因此，由
极限章引理 \ref{limits-lemma-limit-ample}，对某个 $i$，
$\mathcal{L}_{i, y_i}$ 在 $X_{i, y_i}$ 上充足。
由 概形上同调章引理 \ref{coherent-lemma-ample-in-neighbourhood}，
可找到 $y_i$ 的开邻域 $V_i \subset Y_i$，使得 $\mathcal{L}_i$
在 $f_i^{-1}(V_i)$ 上的限制相对于 $V_i$ 充足。
令 $V \subset Y$ 为 $V_i$ 的逆像，即完成证明
（提示：使用 态射章引理 \ref{morphisms-lemma-ample-base-change}、
$X \to Y \times_{Y_i} X_i$ 为仿射态射这一事实，以及由
态射章引理 \ref{morphisms-lemma-pullback-ample-tensor-relatively-ample}
可知充足可逆层沿仿射态射的拉回仍充足这一事实）。
\end{proof}







\section{Proj 与 Spec}
\label{more-morphisms-section-proj-spec}

\noindent
本节阐明分次环的 Proj 与谱之间的关系。

\medskip\noindent
设 $R$ 是环，$A$ 是分次 $R$-代数；见
代数章第 \ref{algebra-section-graded} 节。
对 $m \geq 0$，记 $A_{\geq m} = \bigoplus_{d \geq m} A_d$。
考虑分次环
$$
B = \bigoplus\nolimits_{d \geq 0} A_{\geq d}
$$
对 $d' \geq d$ 及 $a \in A_{d'}$，以 $a^{(d)} \in B$ 表示
$B_d$ 中与 $a$ 对应的元素。以 $\sigma : A \to B$ 和
$\psi : A \to B$ 表示两个显然的环映射：若 $a \in A_d$，则
$\sigma(a) = a^{(0)}$ 且 $\psi(a) = a^{(d)}$。
于是 $\psi$ 是分次环映射，而 $\sigma$ 使 $B$ 成为 $A$ 上的分次代数。
还有一个满分次环映射 $\tau : B \to A$：对 $d' \geq d$ 及
$a \in A_{d'}$，若 $d' > d$，它将 $a^{(d)}$ 映到 $0$；
若 $d' = d$，则映到 $a$。

\medskip\noindent
仿射概形与谱。令 $X = \Spec(A)$。无关理想 $A_+$ 切出闭子概形
$Z = V(A_+) = \Spec(A/A_+) = \Spec(A_0)$。令 $U = X \setminus Z$。
$$
U \longrightarrow X \longrightarrow Z
$$
投射概形与 Proj。令 $P = \text{Proj}(A)$。我们可以且确实将 $P$
视为 $\Spec(A_0) = Z$ 上的概形。令 $L = \text{Proj}(B)$。
我们可以且确实将 $L$ 视为 $\Spec(B_0) = \Spec(A) = X$ 上的概形；
注意，$B_0$ 与 $A$ 的等同由 $\sigma$ 给出。满射 $\tau$ 定义闭浸入
$0 : P \to L$。由于 $A \xrightarrow{\sigma} B \to A$ 等于映射
$A \to A_0 \to A$，断定
$$
\xymatrix{
P \ar[d] \ar[r]_0 & L \ar[d] \\
Z \ar[r] & X
}
$$
交换。

\medskip\noindent
我们断言 $\psi$ 定义态射 $L \to P$。由 构造章引理
\ref{constructions-lemma-morphism-proj}，只需对每个满足
$B_+ \not \subset \mathfrak p$ 的齐次素理想 $\mathfrak p \subset B$
验证 $\psi(A_+) \not \subset \mathfrak p$。首先，选取齐次元素
$g \in B_+$，使得 $g \not \in \mathfrak p$。可以将 $g$ 写成有限和
$g = \sum a_i^{(d)}$，其中对某些 $d_i \geq d$ 有
$a_i \in A_{d_i}$。从而存在 $d' \geq d$ 及 $a \in A_{d'}$，
使得 $a^{(d)} \not \in \mathfrak p$。于是
$$
(a^{(d)})^{d'} =
(a^{d'})^{(d'd)} =
a^{(d)} (a^{d' - 1})^{(d(d' - 1))} =
\psi(a) (a^{d' - 1})^{(d(d' - 1))}
$$
（这个记号并不十分理想）不属于 $\mathfrak p$。因此
$\psi(a) \not \in \mathfrak p$，断言得证。于是可以将上图扩展为交换图表
$$
\xymatrix{
P \ar[d] \ar[r]_0 & L \ar[d] \ar[r]_\pi & P \ar[d] \\
Z \ar[r] & X \ar[r] & Z
}
$$
其中 $X \to Z$ 由 $A_0 \to A$ 给出。由于
$\tau \circ \psi = \text{id}_A$，可见 $\pi \circ 0 = \text{id}_P$。

\medskip\noindent
注意，$\pi$ 是仿射态射。这由 构造章引理
\ref{constructions-lemma-morphism-proj} 中的构造显然可见。
事实上，若对某个 $d > 0$ 有 $f \in A_d$，则令 $g = \psi(f)$ 后，有
$\pi^{-1}(D_+(f)) = D_+(g)$。在此情形，齐次部分有如下等式
$$
(B[1/g])_{m'} = \bigoplus\nolimits_{m \geq m'} (A[1/f])_m
$$
该同构与进一步的局部化相容。取 $m' = 0$，可见
$\pi_*\mathcal{O}_L$ 是所有 $m \geq 0$ 的 $\mathcal{O}_P(m)$ 的直和\footnote{类似地可得
$\pi_*\mathcal{O}_L(i) = \bigoplus_{m \geq -i} \mathcal{O}_P(m)$。}。
于是 $L$ 等同于相对谱：
$$
L = \underline{\Spec}_P
\left(
\bigoplus\nolimits_{m \geq 0} \mathcal{O}_P(m)
\right)
$$
特别地，$L \to P$ 是锥\footnote{$L$ 往往是 $P$ 上的线丛；见下文。}；见
构造章第 \ref{constructions-section-cone} 节。
此外，$0 : P \to L$ 显然是该锥的顶点。

\medskip\noindent
如上一段，取某个 $d > 0$ 的 $f \in A_d$，并令
$g = \psi(f) \in B_d$。考察环映射的结构
$$
\xymatrix{
A_0 \ar[r] \ar[d] & A \ar[d]^\sigma \ar[r] & A_0 \ar[d] \\
(A[1/f])_0 \ar[r]^-\psi &
(B[1/g])_0 = \bigoplus\nolimits_{m \geq 0} (A[1/f])_m
\ar[r]^-\tau &
(A[1/f])_0
}
$$
% P05-EN-0109: the frozen source misspells "computations" as "compuations".
分次环中的一些计算\footnote{(1) 和 (2) 是清楚的。
为看出 (3)，注意若 $a \in A_d$，则
$\sigma(a) = \sigma(f) \psi(a/f)$。对 (4)，注意
$b/g^m$ 属于 $\tau$ 的核，当且仅当
$b \in A_{\geq md}$ 在 $A_{md}$ 中的像为零。
因此，只需说明：若 $m' > md$ 且 $a \in A_{m'}$，则
$a^{(md)}/g^m$ 的某个幂属于由 $\sigma(f)$ 生成的理想。
取 $e$ 使得 $em' - emd \geq d$。于是
$$
(a^{(md)}/g^m)^e = (a^e)^{(emd)}/g^{em} =
(fa^e)^{(emd + d)}/g^{em + 1} = \sigma(f) \cdot (a^e)^{(emd + d)}/g^{em + 1}
$$
正如所需（抱歉使用了如此糟糕的记号）。为看出 (5)，如前论证，并注意当
$d = 1$ 时，$a^{(md)}/g^m = \sigma(f) \cdot a^{(md + 1)}/g^{m + 1}$。}
表明
\begin{enumerate}
\item $\sigma(A_+)(B[1/g])_0 \subset \Ker(\tau : (B[1/g])_0 \to (A[1/f])_0)$,
\item $\sigma(f) \in (B[1/g])_0$ 是非零因子；
\item 作为理想，$\sigma(f) (B[1/g])_0 = \sigma(A_d) (B[1/g])_0$；
\item $\sigma(f) (B[1/g])_0$ 与
$\Ker(\tau : (B[1/g])_0 \to (A[1/f])_0)$ 具有相同的根；
\item 若 $d = 1$，则
$\sigma(f) (B[1/g])_0 = \Ker(\tau : (B[1/g])_0 \to (A[1/f])_0)$.
\end{enumerate}
特别地，可见
$$
0(D_+(f)) = V(\sigma(f)) \subset D_+(g) = \Spec((B[1/g])_0)
$$
在集合论意义下成立。换言之，由 $\sigma(A_d)$ 生成的理想在
$D_+(g)$ 上切出一个有效 Cartier 除子，它在集合论意义下等于闭浸入
$0 : P \to L$ 的像。

\medskip\noindent
我们断言 $L \to X$ 在 $U$ 上是同构。事实上，若对某个 $d > 0$
有 $f \in A_d$，则
$$
\Spec(A_f) \times_X L = \text{Proj}(A_f \otimes_A B) =
\text{Proj}(B_{\sigma(f)})
$$
对每个 $e$，有
$(B_{\sigma(f)})_e = A_f \otimes_B B_e = A_f \otimes_A A_{\geq e} = A_f$,
其中最后一个等式由单射 $A_{\geq e} \subset A$ 诱导。
因此 $B_{\sigma(f)} \cong A_f[T]$，其中 $T$ 的次数为 $1$。
由于 $\text{Proj}(A_f[T]) \to \Spec(A_f)$ 是同构，断言得证。
从现在起，将 $U$ 与 $L$ 中相应的开集等同。

\medskip\noindent
上一段所作的等同使我们可以考虑限制 $\pi|_U : U \to P$。
像上面多次所做的那样，选取某个 $d > 0$ 的 $f \in A_d$，并令
$g = \psi(f) \in B_d$。则
$$
U \cap \pi^{-1}(D_+(f)) = U \cap D_+(g)
$$
在将 $D_+(g)$ 与 $(B[1/g])_0$ 的谱等同后，它是
$\sigma(f) \in (B[1/g])_0$ 的零点轨迹的补集。这正是上面的断言 (4)。
因此 $U \cap D_+(g)$ 仿射，并且
$$
\mathcal{O}_L(U \cap D_+(g)) = (B[1/g])_0[1/\sigma(f)] =
\bigoplus\nolimits_{m \in \mathbf{Z}} (A[1/f])_m
$$
其中最后一个等号是上面等同
$(B[1/g])_0 = \bigoplus_{m \geq 0} (A[1/f])_m$ 的自然延拓。
完全如先前对 $\pi : L \to P$ 所作的那样，断定
$\pi|_U : U \to P$ 仿射，并且作为 $P$ 上的概形，
$$
U = \underline{\Spec}_P
\left(
\bigoplus\nolimits_{m \in \mathbf{Z}} \mathcal{O}_P(m)
\right)
$$
成立。

\medskip\noindent
综上，我们的构造给出概形的交换图表
\begin{equation}
\label{more-morphisms-equation-proj-and-spec}
\vcenter{
\xymatrix{
\underline{\Spec}_P \left(
\bigoplus\nolimits_{m \in \mathbf{Z}} \mathcal{O}_P(m)
\right)
\ar[r] \ar@{=}[d] &
L =
\underline{\Spec}_P \left(
\bigoplus\nolimits_{m \geq 0} \mathcal{O}_P(m)
\right) \ar[d]^\sigma \ar[r]_-\pi & P \ar[d] \\
U \ar[r] & X \ar[r] & Z
}
}
\end{equation}
，其中 $\pi$ 是一个锥，其零截面 $0 : P \to L$ 在集合论意义下映到
$Z$ 在 $L$ 中的逆像上。

\medskip\noindent
令 $W \subset P$ 为满足下列条件的最大开集：$\mathcal{O}_P(1)|_W$
可逆，并且对所有 $m \in \mathbf{Z}$，自然映射诱导同构
$\mathcal{O}_P(m)|_W \cong \mathcal{O}_P(1)^{\otimes m}|_W$；
也就是说，它是 构造章引理
\ref{constructions-lemma-where-invertible} 中 $d = 1$ 所对应的开集。
于是可见，$L|_W = \pi^{-1}(W) \to W$ 是秩为 $1$ 的向量丛
（构造章第 \ref{constructions-section-vector-bundle} 节），即
$$
L|_W = \mathbf{V}(\mathcal{O}_P(1)|_W)
$$
，这里采用 Grothendieck 式记号。上文已表明 $L|_W$ 等于
$\mathcal{O}_P(1)|_W$ 在 $\mathcal{O}_W$ 上的对称代数的相对谱，
故这一点立即成立。显然，态射 $0|_W : W \to L|_W$ 是该向量丛的
零截面。特别地，$0(W)$ 是 $L|_W$ 上的有效 Cartier 除子。
此外，开集 $U|_W = (\pi|_U)^{-1}(W)$ 是零截面的补集。

\medskip\noindent
若 $A$ 由 $f_1, \ldots, f_r \in A_1$ 在 $A_0$ 上生成，则对所有
$m \geq 0$，有 $(f_1, \ldots, f_r)^m = A_{\geq m}$；因此上述 $B$
是 $A_+ = (f_1, \ldots, f_r)$ 的 Rees 代数。在此情形，
$L \to X$ 是沿 $Z$ 的爆破，并且上一段中的 $W$ 满足 $W = P$。

\medskip\noindent
若 $P$ 拟紧，则当 $d$ 充分可除时，由
$\sigma(A_d)\mathcal{O}_L$ 切出的闭子概形 $D \subset L$ 是有效
Cartier 除子，$0 : P \to L$ 经由 $D$ 分解，并且在集合论意义下
$0(P) = D$。这由 构造章引理
\ref{constructions-lemma-proj-quasi-compact} 以及上面证明的 (1)、(2)、
(3)、(4) 得到。（取任意一个能被该引理中所得元素的次数之最小公倍数整除的
$d$。）

\medskip\noindent
继续假设 $P$ 拟紧。设 $\mathcal{F}$ 是拟凝聚 $\mathcal{O}_P$-模，
并令 $\mathcal{F}_U = \pi^*\mathcal{F}|_U$。则有
\begin{equation}
\label{more-morphisms-equation-cohomology-torsor}
R\Gamma(U, \mathcal{F}_U) =
\bigoplus\nolimits_{m \in \mathbf{Z}}
R\Gamma(P, \mathcal{F} \otimes_{\mathcal{O}_P} \mathcal{O}_P(m))
\end{equation}
此外，该直和分解关于 $\mathcal{F}$ 是函子性的；右侧由此诱导的
$A$-模结构与左侧由 $U \subset X$ 给出的 $A$-模结构相同。
为证明该公式，由于 $\pi|_U$ 仿射，并且
$(\pi|_U)_*\mathcal{O}_U = \bigoplus_{m \in \mathbf{Z}} \mathcal{O}_P(m)$，
得到
\begin{align*}
R(\pi|_U)_*\mathcal{F}_U
& =
(\pi|_U)_*\mathcal{F}_U \\
& =
(\pi|_U)_*(\pi|_U)^*\mathcal{F} \\
& =
\mathcal{F} \otimes_{\mathcal{O}_P} 
\bigoplus\nolimits_{m \in \mathbf{Z}} \mathcal{O}_P(m) \\
& =
\bigoplus\nolimits_{m \in \mathbf{Z}}
\mathcal{F} \otimes_{\mathcal{O}_P} \mathcal{O}_P(m)
\end{align*}
由 Leray，得到
$R\Gamma(U, \mathcal{F}_U) =
R\Gamma(P, R(\pi|_U)_*\mathcal{F}_U)$, see
上同调章引理 \ref{cohomology-lemma-apply-Leray}。
在此情形，取上同调与直和交换，故证明完成；见
概形的导出范畴章引理
\ref{perfect-lemma-quasi-coherence-pushforward-direct-sums}。
此处正用到了 $P$ 拟紧；由 构造章引理
\ref{constructions-lemma-proj-separated}，$P$ 是分离的。

\begin{lemma}
\label{more-morphisms-lemma-apply-proj-spec}
设 $R$ 是环，$P$ 是 $R$ 上的固有概形，$\mathcal{L}$ 是充足可逆
$\mathcal{O}_P$-模。令
$A = \bigoplus_{m \geq 0} \Gamma(P, \mathcal{L}^{\otimes m})$。
则 $P = \text{Proj}(A)$，且图表 (\ref{more-morphisms-equation-proj-and-spec}) 变为
$$
\xymatrix{
\underline{\Spec}_P \left(
\bigoplus\nolimits_{m \in \mathbf{Z}} \mathcal{L}^{\otimes m}
\right)
\ar[r] \ar@{=}[d] &
L =
\underline{\Spec}_P \left(
\bigoplus\nolimits_{m \geq 0} \mathcal{L}^{\otimes m}
\right) \ar[d]^\sigma \ar[r]_-\pi & P \ar[d] \\
U \ar[r] & X \ar[r] & Z
}
$$
并具有上述性质。
\end{lemma}

\begin{proof}
由 态射章引理 \ref{morphisms-lemma-proper-ample-is-proj}，
有 $P = \text{Proj}(A)$。此外，由
性质章引理 \ref{properties-lemma-ample-gcd-is-one}，在该等同下，
对所有 $m \in \mathbf{Z}$ 均有
$\mathcal{O}_P(m) = \mathcal{L}^{\otimes m}$。
\end{proof}









\section{纤维中的闭点}
\label{more-morphisms-section-closed-points-fibres}

\noindent
本节部分材料取自预印本 \cite{Osserman-Payne}。

\begin{lemma}
\label{more-morphisms-lemma-locally-principal-vertical}
设 $f : X \to S$ 是概形态射，$Z \subset X$ 是闭子概形，
$s \in S$。假设
\begin{enumerate}
\item $S$ 不可约，其泛点为 $\eta$；
\item $X$ 不可约；
\item $f$ 优势；
\item $f$ 局部有限型；
\item $\dim(X_s) \leq \dim(X_\eta)$,
\item $Z$ 在 $X$ 中局部主；
\item $Z_\eta = \emptyset$.
\end{enumerate}
则在集合论意义下，纤维 $Z_s$ 是 $X_s$ 的若干不可约分支之并。
\end{lemma}

\begin{proof}
以 $X_{red}$ 表示 $X$ 的既约化。则 $Z \cap X_{red}$ 是
$X_{red}$ 的局部主闭子概形；见
除子章引理 \ref{divisors-lemma-pullback-locally-principal}。
因此可以假设 $X$ 既约。换言之，$X$ 是整概形；见
性质章引理 \ref{properties-lemma-characterize-integral}。
在此情形，态射 $X \to S$ 经由 $S_{red}$ 分解；见
概形章引理 \ref{schemes-lemma-map-into-reduction}。
所以可以用 $S_{red}$ 替换 $S$，并假设 $S$ 也是整概形。

\medskip\noindent
$f$ 为优势态射这一断言意味着 $X$ 的泛点映到 $\eta$；见 Morphisms,
引理 \ref{morphisms-lemma-dominant-finite-number-irreducible-components}。
此外，概形 $X_\eta$ 是域 $\kappa(\eta)$ 上局部有限型的整概形。
因此，对 $X_\eta$ 的每一点 $\xi$，都有
$d = \dim(X_\eta) \geq 0$ 等于 $\dim_\xi(X_\eta)$；见
代数章引理 \ref{algebra-lemma-dimension-spell-it-out} 及
\ref{algebra-lemma-dimension-at-a-point-finite-type-over-field}。
结合 态射章引理 \ref{morphisms-lemma-openness-bounded-dimension-fibres}
与条件 (5)，断定对每个 $x \in X_s$ 都有 $\dim_x(X_s) = d$。

\medskip\noindent
在 Noether 情形，该断言可如下证明。
% P05-EN-0110: frozen source has grammatical "does not holds"; Chinese follows intended syntax.
若引理不成立，则存在 $x \in Z_s$，它是 $Z_s$ 的某个不可约分支的泛点，
却不是 $X_s$ 的任何不可约分支的泛点。于是
$\dim_x(Z_s) \leq d - 1$，因为 $\dim_x(X_s) = d$，并且在 $X_s$
中 $x$ 的一个邻域内，闭子概形 $Z_s$ 不包含 $X_s$ 的任何不可约分支。
所以用 $x$ 的一个开邻域替换 $X$ 后，可以假设对所有 $z \in Z$ 都有
$\dim_z(Z_{f(z)}) \leq d - 1$；见
态射章引理 \ref{morphisms-lemma-openness-bounded-dimension-fibres}。
令 $\xi' \in Z$ 为 $Z$ 的某个不可约分支的泛点，并令
% P05-EN-0111: frozen source sets s'=f(xi) before xi is introduced; expected f(xi').
$s' = f(\xi)$。由于 $Z \not = X$ 局部主，可见
% P05-EN-0112: frozen source uses O_{X,xi}; the codimension-one point is evidently xi'.
$\dim(\mathcal{O}_{X, \xi}) = 1$；见
代数章引理 \ref{algebra-lemma-minimal-over-1}
（此处用到了 $X$ 为 Noether 概形）。令 $\xi \in X$ 为 $X$ 的泛点，
并令 $\xi_1$ 为 $X_{s'}$ 的任意一个包含 $\xi'$ 的不可约分支的泛点。
于是有特化
$$
\xi \leadsto \xi_1 \leadsto \xi'.
$$
由于 $\dim(\mathcal{O}_{X, \xi}) = 1$，两个特化之一必为等式。
按假设 $s' \not = \eta$，故第一个特化不是等式。因此
$\xi' = \xi_1$ 是 $X_{s'}$ 的某个不可约分支的泛点。
再次应用 态射章引理
\ref{morphisms-lemma-openness-bounded-dimension-fibres}，可得
$\dim_{\xi'}(Z_{s'}) = \dim_{\xi'}(X_{s'}) \geq \dim(X_\eta) = d$
，这给出所需矛盾。

\medskip\noindent
在一般情形，按如下方式化归到 Noether 情形。若引理不成立，则存在位于
$s$ 上方的一点 $x \in X$，$x$ 是 $Z_s$ 的某个不可约分支的泛点，
却不是 $X_s$ 的任何不可约分支的泛点。令 $U \subset S$ 为 $s$
的仿射邻域，并令 $V \subset X$ 为满足 $f(V) \subset U$ 的 $x$
的仿射邻域。写成 $U = \Spec(A)$、$V = \Spec(B)$，使得 $f|_V$
由环映射 $A \to B$ 给出。令 $\mathfrak q \subset B$、相应地令
$\mathfrak p \subset A$ 为分别对应于 $x$、$s$ 的素理想。
必要时缩小 $V$ 后，可假设 $Z \cap V$ 由某个元素 $g \in B$ 切出。
以 $K$ 表示 $A$ 的分式域。此时已知如下事实：
\begin{enumerate}
\item $A \subset B$ 是整环的有限生成扩张；
\item 元素 $g \otimes 1$ 在 $B \otimes_A K$ 中可逆；
\item $d = \dim(B \otimes_A K) = \dim(B \otimes_A \kappa(\mathfrak p))$,
\item $g \otimes 1$ 不是 $B \otimes_A \kappa(\mathfrak p)$ 的单位；
\item $g \otimes 1$ 不属于 $B \otimes_A \kappa(\mathfrak p)$ 的任何极小素理想。
\end{enumerate}
我们要导出矛盾。

\medskip\noindent
选取在 $A$ 上生成 $B$ 的元素 $x_1, \ldots, x_n \in B$。
对有限生成的 $\mathbf{Z}$-代数 $A_0 \subset A$，令
$B_0 \subset B$ 为由 $x_1, \ldots, x_n$ 生成的 $A_0$-子代数，
以 $K_0$ 表示 $A_0$ 的分式域，并令
$\mathfrak p_0 = A_0 \cap \mathfrak p$。我们断言，当 $A_0$ 充分大时，
(1) -- (5) 对系统 $(A_0 \subset B_0, g, \mathfrak p_0)$ 也成立。

\medskip\noindent
逐一证明这些条件。第 (1) 部分按构造成立。对第 (2) 部分，
对某个 $h \otimes 1/a \in B \otimes_A K$ 写成
$(g \otimes 1) h = 1$。对某些系数 $a_I, a'_I \in A$，写成
$g = \sum a_I x^I$、$h = \sum a'_I x^I$（多重指标记号）。
只要 $A_0$ 包含 $a$ 及所有 $a_I, a'_I$，则 (2) 成立，因为
$B_0 \otimes_{A_0} K_0 \subset B \otimes_A K$
（它们是单射 $B_0 \to B$ 的局部化）。为得到 (3)，考虑正合列
$$
0 \to I \to A[X_1, \ldots, X_n] \to B \to 0
$$
它定义 $I$，其中第二个映射将 $X_i$ 映到 $x_i$。由于 $\otimes$
右正合，可见 $I \otimes_A K$、相应地
$I \otimes_A \kappa(\mathfrak p)$ 分别是满射
$K[X_1, \ldots, X_n] \to B \otimes_A K$、相应地
$\kappa(\mathfrak p)[X_1, \ldots, X_n] \to B \otimes_A \kappa(\mathfrak p)$
的核。由于域上的多项式环是 Noether 环，存在有限多个元素
$h_j \in I$，$j = 1, \ldots, m$，它们生成
$I \otimes_A K$ 和 $I \otimes_A \kappa(\mathfrak p)$。
写成 $h_j = \sum a_{j, I}X^I$。只要 $A_0$ 包含所有 $a_{j, I}$，
便得到
$$
B_0 \otimes_{A_0} K_0 \otimes_{K_0} K = B \otimes_A K
\quad\text{且}\quad
B_0 \otimes_{A_0} \kappa(\mathfrak p_0)
\otimes_{\kappa(\mathfrak p_0)} \kappa(\mathfrak p)
=
B \otimes_A \kappa(\mathfrak p).
$$
由 态射章引理 \ref{morphisms-lemma-dimension-fibre-after-base-change}
或 代数章引理 \ref{algebra-lemma-dimension-preserved-field-extension}，
可见 (3) 中的维数等式成立。第 (4) 部分立即成立。由于
$B_0 \otimes_{A_0} \kappa(\mathfrak p_0) \subset
B \otimes_A \kappa(\mathfrak p)$，
$B_0 \otimes_{A_0} \kappa(\mathfrak p_0)$ 的每个极小素理想，都位于
$B \otimes_A \kappa(\mathfrak p)$ 的某个极小素理想之下；见
代数章引理 \ref{algebra-lemma-image-dense-generic-points}。
这蕴含 (5) 成立。由此将问题化归到上面已经处理的 Noether 情形。
\end{proof}

\noindent
下面是上述引理的一个代数应用。若 $A \to B$ 平坦，则该引理的第四个假设
成立；见 引理 \ref{more-morphisms-lemma-equality-dimensions}。

\begin{lemma}
\label{more-morphisms-lemma-horizontal}
设 $A \to B$ 是局部环的局部同态，且 $g \in \mathfrak m_B$。假设
\begin{enumerate}
\item $A$ 与 $B$ 是整环，且 $A \subset B$；
\item $B$ 在 $A$ 上本质有限型；
\item $g$ 不属于 $\mathfrak m_AB$ 上方的任何极小素理想；
\item $\dim(B/\mathfrak m_AB) +
\text{trdeg}_{\kappa(\mathfrak m_A)}(\kappa(\mathfrak m_B)) =
\text{trdeg}_A(B)$.
\end{enumerate}
则 $A \subset B/gB$，即 $\Spec(A)$ 的泛点属于态射
$\Spec(B/gB) \to \Spec(A)$ 的像。
\end{lemma}

\begin{proof}
由 代数章引理 \ref{algebra-lemma-image-dense-generic-points}，
注意这两个断言等价。为开始证明，令 $C$ 为有限型 $A$-代数，
$\mathfrak q$ 为 $C$ 的素理想，使得 $B = C_{\mathfrak q}$。
当然可以假设 $C$ 是整环且 $g \in C$。用一个局部化替换 $C$ 后，有
$\dim(C/\mathfrak m_AC) = \dim(B/\mathfrak m_AB) +
\text{trdeg}_{\kappa(\mathfrak m_A)}(\kappa(\mathfrak m_B))$；见
态射章引理 \ref{morphisms-lemma-dimension-fibre-at-a-point}。
令 $K$ 等于 $A$ 的分式域，由同一引理可见
$\dim(C \otimes_A K) = \text{trdeg}_A(B)$。因此，假设 (4) 意味着态射
$\Spec(C) \to \Spec(A)$ 的泛纤维与闭纤维具有相同维数。

\medskip\noindent
假设引理不成立。则 $(B/gB) \otimes_A K = 0$。这意味着
$g \otimes 1$ 在 $B \otimes_A K
= C_{\mathfrak q} \otimes_A K$
中可逆。由于 $C_{\mathfrak q}$ 是主局部化的极限，断定对某个
$h \in C$、$h \not \in \mathfrak q$，$g \otimes 1$ 在
$C_h \otimes_A K$ 中可逆。因此，用 $C_h$ 替换 $C$ 后，可以假设
$(C/gC) \otimes_A K = 0$。再替换一次 $C$，以保证
$C/\mathfrak m_AC$ 的极小素理想与 $B/\mathfrak m_AB$ 的极小素理想
一一对应。此时，将 引理 \ref{more-morphisms-lemma-locally-principal-vertical}
应用于 $X = \Spec(C) \to \Spec(A) = S$ 以及局部闭子概形
$Z = \Spec(C/gC)$。由于 $Z_K = \emptyset$，可见
$Z \otimes \kappa(\mathfrak m_A)$ 必须包含
$X \otimes \kappa(\mathfrak m_A) = \Spec(C/\mathfrak m_AC)$ 的一个
不可约分支。但这与 $g$ 不属于 $\mathfrak m_AB$ 上方任何极小素理想的
假设矛盾。引理得证。
\end{proof}

\begin{lemma}
\label{more-morphisms-lemma-equality-dimensions}
设 $A \to B$ 是局部环的局部同态。假设
\begin{enumerate}
\item $A$ 与 $B$ 是整环，且 $A \subset B$；
\item $B$ 在 $A$ 上本质有限型；
\item $B$ 在 $A$ 上平坦。
\end{enumerate}
则有
$$
\dim(B/\mathfrak m_AB) +
\text{trdeg}_{\kappa(\mathfrak m_A)}(\kappa(\mathfrak m_B)) =
\text{trdeg}_A(B).
$$
\end{lemma}

\begin{proof}
令 $C$ 为有限型 $A$-代数，$\mathfrak q$ 为 $C$ 的素理想，使得
$B = C_{\mathfrak q}$。可以假设 $C$ 是整环。由 态射章引理
\ref{morphisms-lemma-dimension-fibre-at-a-point}，有
$\dim_{\mathfrak q}(C/\mathfrak m_AC) = \dim(B/\mathfrak m_AB) +
\text{trdeg}_{\kappa(\mathfrak m_A)}(\kappa(\mathfrak m_B))$。
令 $K$ 等于 $A$ 的分式域，由同一引理可见
$\dim(C \otimes_A K) = \text{trdeg}_A(B)$。
因此，实际要证明的是
$\dim_{\mathfrak q}(C/\mathfrak m_AC) = \dim(C \otimes_A K)$。
选取 $K$ 中支配 $A$ 的赋值环 $A'$；见
代数章引理 \ref{algebra-lemma-dominate}。令 $C' = C \otimes_A A'$。
选取 $C'$ 中位于 $\mathfrak q$ 上方的素理想 $\mathfrak q'$；
这样的素理想存在，因为
$$
C'/\mathfrak m_{A'}C' =
C/\mathfrak m_AC \otimes_{\kappa(\mathfrak m_A)} \kappa(\mathfrak m_{A'})
$$
这还证明 $C/\mathfrak m_AC \to C'/\mathfrak m_{A'}C'$ 忠实平坦。
此外，这证明
$\dim_{\mathfrak q}(C/\mathfrak m_AC) =
\dim_{\mathfrak q'}(C'/\mathfrak m_{A'}C')$；见
代数章引理 \ref{algebra-lemma-dimension-at-a-point-preserved-field-extension}。
注意，$B' = C'_{\mathfrak q'}$ 是 $B \otimes_A A'$ 的一个局部化。
因此 $B'$ 在 $A'$ 上平坦。泛纤维 $B' \otimes_{A'} K$ 是
$B \otimes_A K$ 的一个局部化，故 $B'$ 是整环。若对
$A' \subset B'$ 证明本引理，则得到等式
$\dim_{\mathfrak q'}(C'/\mathfrak m_{A'}C') = \dim(C' \otimes_{A'} K)$
，结合上述讨论，它蕴含所需等式
$\dim_{\mathfrak q}(C/\mathfrak m_AC) = \dim(C \otimes_A K)$
。由此将本引理化归到 $A$ 为赋值环的情形。

\medskip\noindent
设 $A \subset B$ 如引理所述，且 $A$ 是赋值环。像先前一样，
对某个有限型 $A$-整代数 $C$ 写成 $B = C_{\mathfrak q}$。
由 代数章引理
\ref{algebra-lemma-finite-type-domain-over-valuation-ring-dim-fibres}，
得到 $\dim(C/\mathfrak m_AC) = \dim(C \otimes_A K)$，结论得证。
\end{proof}

\begin{lemma}
\label{more-morphisms-lemma-closed-point-nearby-fibre}
设 $f : X \to S$ 是概形态射，$x \leadsto x'$ 是 $X$ 中点的特化。
令 $s = f(x)$ 且 $s' = f(x')$。假设
\begin{enumerate}
\item $x'$ 是 $X_{s'}$ 的闭点；
\item $f$ 局部有限型。
\end{enumerate}
则集合
$$
\{x_1 \in X
\text{ 使得 }
f(x_1) = s
\text{ 且 }
x_1\text{ is closed in }X_s
\text{ 且 }
x \leadsto x_1 \leadsto x'
\}
$$
在 $x$ 于 $X_s$ 中的闭包内稠密。
\end{lemma}

\begin{proof}
将 概形章引理 \ref{schemes-lemma-points-specialize}
应用于特化 $x \leadsto x'$。它给出态射
$\varphi : \Spec(B) \to X$，其中 $B$ 是赋值环，且 $\varphi$
将泛点映到 $x$、闭点映到 $x'$。还可以假设 $\kappa(x)$ 是 $B$
的分式域。令 $A = B \cap \kappa(s)$。注意，这是一个赋值环
（见 代数章引理 \ref{algebra-lemma-valuation-ring-cap-field}），
它支配 $\mathcal{O}_{S, s'} \to \kappa(s)$ 的像。考虑交换图表
$$
\xymatrix{
\Spec(B) \ar[rd] \ar[r] &
X_A \ar[d] \ar[r] & X \ar[d] \\
& \Spec(A) \ar[r] & S
}
$$
$B$ 的泛点（相应地，闭点）映到 $X_A$ 中位于 $\Spec(A)$ 泛点
（相应地，闭点）上方的一点 $x_A$（相应地，$x'_A$）。由 Morphisms,
引理 \ref{morphisms-lemma-base-change-closed-point-fibre-locally-finite-type}，
$x'_A$ 是 $X_A$ 的特殊纤维的闭点。注意，$X_A \to \Spec(A)$ 的泛纤维
同构于 $X_s$。因此，已经将引理化归到如下情形：$S$ 是一个赋值环的谱，
$s = \eta \in S$ 是泛点，而 $s' \in S$ 是闭点。

\medskip\noindent
对 $\dim_x(X_\eta)$ 归纳证明本引理。若 $\dim_x(X_\eta) = 0$，
则 $X_\eta$ 中没有其他点特化到 $x$，且 $x$ 在其纤维中闭；见
态射章引理 \ref{morphisms-lemma-quasi-finite-at-point-characterize}。
此时结论成立。以下假设 $\dim_x(X_\eta) > 0$。

\medskip\noindent
令 $X' \subset X$ 为 $X$ 的不可约闭子概形 $\overline{\{x\}}$
上赋予的既约诱导概形结构；见
概形章定义 \ref{schemes-definition-reduced-induced-scheme}。
为证明本引理，可以用 $X'$ 替换 $X$，因为这只会减小
$\dim_x(X_\eta)$。所以还可以假设 $X$ 是整概形，而 $x$ 是其泛点。
此外，可以用 $x'$ 的一个仿射邻域替换 $X$。于是
$X = \Spec(B)$，其中 $A \subset B$ 是整环的有限型扩张。
注意，在此情形
$\dim_x(X_\eta) = \dim(X_\eta) = \dim(X_{s'})$，并且事实上
$X_{s'}$ 等维；见 代数章引理
\ref{algebra-lemma-finite-type-domain-over-valuation-ring-dim-fibres}。

\medskip\noindent
令 $W \subset X_\eta$ 为真闭子集（这是我们想要“避开”的子集）。
由于 $X_s$ 在一个域上有限型，可见 $W$ 只有有限多个不可约分支，写成
$W = W_1 \cup \ldots \cup W_n$。令
% P05-EN-0113: frozen source indexes the n corresponding components by j=1,...,r.
$\mathfrak q_j \subset B$，$j = 1, \ldots, r$，为相应的素理想。
令 $\mathfrak q \subset B$ 为对应于点 $x'$ 的极大理想。
令 $\mathfrak p_1, \ldots, \mathfrak p_s \subset B$ 为位于
$\mathfrak m_AB$ 上方的极小素理想。这样的素理想只有有限多个，
因为它们对应于 Noether 概形 $X_{s'}$ 的不可约分支。此外，这些不可约分支
的维数都为 $> 0$（见上文），故对所有 $i$ 都有
$\mathfrak p_i \not = \mathfrak q$。现在选取元素
$g \in \mathfrak q$，使得对所有 $j$ 有 $g \not \in \mathfrak q_j$，
且对所有 $i$ 有 $g \not \in \mathfrak p_i$；见
代数章引理 \ref{algebra-lemma-silly}。
以 $Z \subset X$ 表示由 $g$ 定义的局部主闭子概形。令
$Z_\eta = Z_{1, \eta} \cup \ldots \cup Z_{n, \eta}$，$n \geq 0$，
为 $Z$ 的泛纤维的不可约分支分解（分支有限，因为泛纤维是 Noether 的）。
以 $Z_i \subset X$ 表示 $Z_{i, \eta}$ 的闭包。用一个更小的仿射邻域替换
$X$ 后，可以假设对每个 $i = 1, \ldots, n$ 都有 $x' \in Z_i$。
按构造，$Z \cap X_{s'}$ 不包含 $X_{s'}$ 的任何不可约分支。
因此，由 引理 \ref{more-morphisms-lemma-locally-principal-vertical}，断定
$Z_\eta \not = \emptyset$！换言之，$n \geq 1$。令
$x_1 \in Z_1$ 为泛点，则 $x_1 \leadsto x'$ 且 $f(x_1) = \eta$。
此外，按构造，$Z_{1, \eta} \cap W_j \subset W_j$ 是真闭子集。
所以 $Z_{1, \eta} \cap W_j$ 的每个不可约分支在 $X_\eta$ 中的余维
均为 $\geq 2$，而由 代数章引理 \ref{algebra-lemma-minimal-over-1}，
$\text{codim}(Z_{1, \eta}, X_\eta) = 1$。因此
$W \cap Z_{1, \eta}$ 是真闭子集。此时，归纳假设可应用于
$Z_1 \to S$ 及特化 $x_1 \leadsto x'$。由此得到
$Z_{1, \eta}$ 的一个不属于 $W$ 且特化到 $x'$ 的闭点 $x_2$。
于是得到 $x \leadsto x_2 \leadsto x'$；点 $x_2$ 在 $X_\eta$ 中闭，
且如所需 $x_2 \not \in W$。
\end{proof}

\begin{remark}
\label{more-morphisms-remark-full-specialization-sequence}
引理 \ref{more-morphisms-lemma-closed-point-nearby-fibre} 的证明实际上表明，
存在特化序列
$$
x \leadsto x_1 \leadsto x_2 \leadsto \ldots \leadsto x_d \leadsto x'
$$
其中所有 $x_i$ 都属于纤维 $X_s$，每个特化都是直接的，而
$x_d$ 是 $X_s$ 的闭点。整数
$d = \text{trdeg}_{\kappa(s)}(\kappa(x)) = \dim(\overline{\{x\}})$
，其中闭包取于 $X_s$ 中。此外，可以选取这些点 $x_i$，使其避开
$X_s$ 中任何不包含点 $x$ 的闭子集。
\end{remark}

\noindent
Examples, 第 \ref{examples-section-topology-finite-type} 节 表明，
若不假设 $A$ 为 Noether 环，则下述引理不成立。

\begin{lemma}
\label{more-morphisms-lemma-quasi-finite-quasi-section-meeting-nearby-open}
设 $\varphi : A \to B$ 是局部环的局部环映射。设
$V \subset \Spec(B)$ 是一个开子概形，且至少包含一个不位于
$\mathfrak m_A$ 上方的素理想。假设 $A$ 是 Noether 环，$\varphi$
本质有限型，并且 $A/\mathfrak m_A \subset B/\mathfrak m_B$ 有限。
则存在 $\mathfrak q \in V$，满足
$\mathfrak m_A \not = \mathfrak q \cap A$，使得
$A \to B/\mathfrak q$ 是某个拟有限环映射的局部化。
\end{lemma}

\begin{proof}
由于 $A$ 是 Noether 环，而 $A \to B$ 本质有限型，可知 $B$ 也是
Noether 环。由 性质章引理
\ref{properties-lemma-complement-closed-point-Jacobson}，拓扑空间
$\Spec(B) \setminus \{\mathfrak m_B\}$ 是 Jacobson 的。因此可以在非空开集
$$
V \setminus \{\mathfrak q \subset B \mid \mathfrak m_A = \mathfrak q \cap A\}.
$$
中选取闭点 $\mathfrak q$。（该集按假设非空；它是开的，因为
$\{\mathfrak m_A\}$ 是 $\Spec(A)$ 的闭子集。）
于是 $\Spec(B/\mathfrak q)$ 有两个点，即 $\mathfrak m_B$ 和
$\mathfrak q$，而 $\mathfrak q$ 不位于 $\mathfrak m_A$ 上方。
对某个有限型 $A$-代数 $C$ 及素理想 $\mathfrak m$，写成
$B/\mathfrak q = C_{\mathfrak m}$。由
代数章引理 \ref{algebra-lemma-isolated-point-fibre} (2)，
$A \to C$ 在 $\mathfrak m$ 处拟有限。因此，由
代数章引理 \ref{algebra-lemma-quasi-finite-open}，
用一个主局部化替换 $C$ 后，环映射 $A \to C$ 拟有限。
\end{proof}

\begin{lemma}
\label{more-morphisms-lemma-quasi-finite-quasi-section-meeting-nearby-open-X}
设 $f : X \to S$ 是概形态射。取 $x \in X$，其像为 $s \in S$。
设 $U \subset X$ 是开子概形。假设 $f$ 局部有限型，$S$ 局部
Noether，$x$ 是 $X_s$ 的闭点，并假设存在一点 $x' \in U$，满足
$x' \leadsto x$ 且 $f(x') \not = s$。则存在闭子概形 $Z \subset X$，
使得 (a) $x \in Z$；(b) $f|_Z : Z \to S$ 在 $x$ 处拟有限；
(c) 存在 $z \in Z$，$z \in U$、$z \leadsto x$ 且 $f(z) \not = s$。
\end{lemma}

\begin{proof}
这是 引理 \ref{more-morphisms-lemma-quasi-finite-quasi-section-meeting-nearby-open}
的改述。具体地，令 $A = \mathcal{O}_{S, s}$、
$B = \mathcal{O}_{X, x}$。以 $V \subset \Spec(B)$ 表示 $U$ 的逆像。
环映射 $f^\sharp : A \to B$ 本质有限型。按假设，$V$ 中至少有一点
不映到 $\Spec(A)$ 的闭点。因此，
引理 \ref{more-morphisms-lemma-quasi-finite-quasi-section-meeting-nearby-open}
的所有假设均成立，故得到素理想 $\mathfrak q \subset B$，它不位于
$\mathfrak m_A$ 上方，并且 $A \to B/\mathfrak q$ 是某个拟有限环映射的
局部化。令 $z \in X$ 为点 $\mathfrak q$ 在典范态射
$\Spec(B) \to X$ 下的像。令 $Z = \overline{\{z\}}$，并赋予诱导既约
概形结构。由于 $z \leadsto x$，可见 $x \in Z$ 且
$\mathcal{O}_{Z, x} = B/\mathfrak q$。按构造，$Z \to S$ 在 $x$ 处拟有限。
\end{proof}

\begin{remark}
\label{more-morphisms-remark-alternative-closed-point-nearby-fibre}
当 $S$ 局部 Noether 时，可以利用
引理 \ref{more-morphisms-lemma-quasi-finite-quasi-section-meeting-nearby-open}
或其变体
引理 \ref{more-morphisms-lemma-quasi-finite-quasi-section-meeting-nearby-open-X}，
给出 引理 \ref{more-morphisms-lemma-closed-point-nearby-fibre} 的另一证明。
下面是一个粗略纲要。首先，用 $s'$ 处局部环的谱替换 $S$，然后对
$\dim(S)$ 归纳。当 $\dim(S) = 0$ 时，由于此时 $s' = s$，结论平凡。
用 $\overline{\{x\}}$ 上的既约诱导概形结构替换 $X$。
将 引理 \ref{more-morphisms-lemma-quasi-finite-quasi-section-meeting-nearby-open-X}
应用于 $X \to S$、$x' \mapsto s'$ 以及任意含有 $x$ 的非空开集
$U \subset X$。由此得到闭子概形
% P05-EN-0114: frozen source omits conjunction/punctuation between x' in Z subset X and a point z in Z.
$x' \in Z \subset X$ 及一点 $z \in Z$，使得 $Z \to S$ 在 $x'$ 处
拟有限，并且 $f(z) \not = s'$。于是 $z$ 是 $X_{f(z)}$ 的闭点，
且 $z \leadsto x'$。由于 $f(z) \not = s'$，可见
$\dim(\mathcal{O}_{S, f(z)}) < \dim(S)$。由于 $x$ 是 $X$ 的泛点，
有 $x \leadsto z$，因而 $s = f(x) \leadsto f(z)$。
将归纳假设应用于 $s \leadsto f(z)$ 及 $z \mapsto f(z)$，即得结论。
\end{remark}

\begin{lemma}
\label{more-morphisms-lemma-change-hypotheses}
假设 $f : X \to S$ 局部有限型，$S$ 局部 Noether，$x \in X$ 是其纤维
$X_s$ 的闭点，而 $U \subset X$ 是满足 $U \cap X_s = \emptyset$
及 $x \in \overline{U}$ 的开子概形。则
引理 \ref{more-morphisms-lemma-quasi-finite-quasi-section-meeting-nearby-open-X}
的结论成立。
\end{lemma}

\begin{proof}
可以按如下方式化归到所引引理：首先分别用 $x$ 与 $s$ 的仿射邻域替换
$X$ 与 $S$。于是 $X$ 是 Noether 的，特别地，$U$ 拟紧（见
态射章引理 \ref{morphisms-lemma-finite-type-noetherian} 及
拓扑章引理 \ref{topology-lemma-Noetherian} 及
\ref{topology-lemma-Noetherian-quasi-compact}）。
因此存在满足 $x' \in U$ 的特化 $x' \leadsto x$（见
态射章引理 \ref{morphisms-lemma-reach-points-scheme-theoretic-image}）。
注意，$f(x') \not = s$。由此可见，该引理的全部假设均成立，结论得证。
\end{proof}



\section{Stein 分解}
\label{more-morphisms-section-stein-factorization}

\noindent
Stein 分解断言：若固有态射 $f : X \to S$ 满足
$f_*\mathcal{O}_X = \mathcal{O}_S$，则它具有连通纤维。

\begin{lemma}
\label{more-morphisms-lemma-stein-universally-closed}
设 $S$ 是概形，$f : X \to S$ 是普遍闭且拟分离的态射。
则存在分解
$$
\xymatrix{
X \ar[rr]_{f'} \ar[rd]_f & & S' \ar[dl]^\pi \\
& S &
}
$$
具有下列性质：
\begin{enumerate}
\item 态射 $f'$ 普遍闭、拟紧、拟分离且满；
\item 态射 $\pi : S' \to S$ 是整态射；
\item 有 $f'_*\mathcal{O}_X = \mathcal{O}_{S'}$；
\item 有 $S' = \underline{\Spec}_S(f_*\mathcal{O}_X)$；
\item $S'$ 是 $S$ 在 $X$ 中的正规化；见
态射章定义 \ref{morphisms-definition-normalization-X-in-Y}.
\end{enumerate}
分解 $f = \pi \circ f'$ 的形成与平坦基变换交换。
\end{lemma}

\begin{proof}
由 态射章引理 \ref{morphisms-lemma-universally-closed-quasi-compact}，
态射 $f$ 拟紧。因此，$S$ 在 $X$ 中的正规化 $S'$ 有定义
（态射章定义 \ref{morphisms-definition-normalization-X-in-Y}），
并且有分解 $X \to S' \to S$。由 态射章引理
\ref{morphisms-lemma-normalization-in-universally-closed}，得到 (2)、(4)、
(5)。由 态射章引理 \ref{morphisms-lemma-image-proper-scheme-closed}，
态射 $f'$ 普遍闭。由 概形章引理
\ref{schemes-lemma-quasi-compact-permanence}，它拟紧；由
概形章引理 \ref{schemes-lemma-compose-after-separated}，它拟分离。

\medskip\noindent
为证明其余陈述，可以假设基概形 $S$ 仿射，写成 $S = \Spec(R)$。
于是 $S' = \Spec(A)$，其中 $A = \Gamma(X, \mathcal{O}_X)$ 是整
$R$-代数。因而 $f'_*\mathcal{O}_X$ 显然就是 $\mathcal{O}_{S'}$
（由 概形章引理 \ref{schemes-lemma-push-forward-quasi-coherent}，
$f'_*\mathcal{O}_X$ 拟凝聚，故等于 $\widetilde{A}$）。这证明了 (3)。

\medskip\noindent
下面证明 $f'$ 是满射。由于 $f'$ 普遍闭（见上文），$f'$ 的像是闭子集
$V(I) \subset S' = \Spec(A)$。取 $h \in I$。则
$h|_X = f^\sharp(h)$ 是 $X$ 的结构层的一个整体截面，并在每一点消失。
由于 $X$ 拟紧，这意味着 $h|_X$ 是幂零截面，即对某个 $n > 0$ 有
% P05-EN-0115: frozen source writes h^n|X without the expected restriction underscore.
$h^n|X = 0$。但 $A = \Gamma(X, \mathcal{O}_X)$，故 $h^n = 0$。
由于 $I$ 的每个元素都幂零，断定 $V(I) = S'$，正如所需。

\medskip\noindent
由 概形上同调章引理
\ref{coherent-lemma-flat-base-change-cohomology}，$f_*\mathcal{O}_X$ 的形成
与平坦基变换交换。由 构造章引理
\ref{constructions-lemma-spec-properties}，相对谱的形成与任意基变换交换。
因此，该分解的形成与平坦基变换交换。
\end{proof}

\begin{lemma}
\label{more-morphisms-lemma-stein-universally-closed-residue-fields}
在 引理 \ref{more-morphisms-lemma-stein-universally-closed} 中进一步假设 $f$ 局部有限型。
则对 $s \in S$，纤维 $\pi^{-1}(\{s\}) = \{s_1, \ldots, s_n\}$
是有限的，并且域扩张 $\kappa(s_i)/\kappa(s)$ 有限。
\end{lemma}

\begin{proof}
回忆，$\pi^{-1}(\{s\})$ 的点之间不存在特化；见
代数章引理 \ref{algebra-lemma-integral-no-inclusion}。由于 $f'$ 满，
可知 $|X_s| \to \pi^{-1}(\{s\})$ 满。注意，$X_s$ 是域上的有限型
拟分离概形（拟紧性已在所引引理的证明中说明）。因此 $X_s$ 是
Noether 的（态射章引理 \ref{morphisms-lemma-finite-type-noetherian}）。
现在，一个拓扑论证（略去）表明 $\pi^{-1}(\{s\})$ 有限。
对每个 $i$，可以选取映到 $s_i$ 的有限型点 $x_i \in X_s$
（态射章引理 \ref{morphisms-lemma-enough-finite-type-points}）。
由此断定 $\kappa(s_i)/\kappa(s)$ 有限：$x_i$ 可由有限型态射
$\Spec(k_i) \to X_s$ 表示（按有限型点的定义），因此
$\Spec(k_i) \to s = \Spec(\kappa(s))$ 是有限型态射
（作为有限型态射的复合），故 $k_i/\kappa(s)$ 有限
（态射章引理 \ref{morphisms-lemma-point-finite-type}）。
\end{proof}

\begin{lemma}
\label{more-morphisms-lemma-characterize-geometrically-connected-fibres}
设 $f : X \to S$ 是概形态射，$s \in S$。则 $X_s$ 几何连通，
当且仅当对每个 \'etale 邻域 $(U, u) \to (S, s)$，基变换
$X_U \to U$ 的纤维 $X_u$ 连通。
\end{lemma}

\begin{proof}
若 $X_s$ 几何连通，则它的任意基变换都连通。反之，假设 $X_s$
不几何连通。由 簇章引理
\ref{varieties-lemma-characterize-geometrically-disconnected}，
对某个有限可分域扩张 $k/\kappa(s)$，
$X_s \times_{\Spec(\kappa(s))} \Spec(k)$ 不连通。
由 引理 \ref{more-morphisms-lemma-realize-prescribed-residue-field-extension-etale}，
存在仿射 \'etale 邻域 $(U, u) \to (S, s)$，使得
$\kappa(u)/\kappa(s)$ 与 $k/\kappa(s)$ 等同。在此情形，$X_u$ 不连通。
\end{proof}

\begin{theorem}[Stein 分解；Noether 情形]
\label{more-morphisms-theorem-stein-factorization-Noetherian}
设 $S$ 是局部 Noether 概形，$f : X \to S$ 是固有态射。
则存在分解
$$
\xymatrix{
X \ar[rr]_{f'} \ar[rd]_f & & S' \ar[dl]^\pi \\
& S &
}
$$
具有下列性质：
\begin{enumerate}
\item 态射 $f'$ 固有，并具有几何连通纤维；
\item 态射 $\pi : S' \to S$ 有限；
\item 有 $f'_*\mathcal{O}_X = \mathcal{O}_{S'}$；
\item 有 $S' = \underline{\Spec}_S(f_*\mathcal{O}_X)$；
\item $S'$ 是 $S$ 在 $X$ 中的正规化；见
态射章定义 \ref{morphisms-definition-normalization-X-in-Y}.
\end{enumerate}
\end{theorem}

\begin{proof}
令 $f = \pi \circ f'$ 为 引理 \ref{more-morphisms-lemma-stein-universally-closed}
中的分解。注意，除 引理 \ref{more-morphisms-lemma-stein-universally-closed} 的结论外，
$f'$ 还是分离的
（概形章引理 \ref{schemes-lemma-compose-after-separated}）且有限型
（态射章引理 \ref{morphisms-lemma-permanence-finite-type}）。
因此 $f'$ 固有。由 概形上同调章命题
\ref{coherent-proposition-proper-pushforward-coherent}，
$f_*\mathcal{O}_X$ 是凝聚 $\mathcal{O}_S$-模。因此 $\pi$ 有限，
即 (2) 成立。

\medskip\noindent
这证明了除最有趣的断言外的所有结论；剩下要证的是 $f'$ 的所有纤维均
几何连通。由上述讨论，显然可以用 $S'$ 替换 $S$，因而可假设 $S$
是 Noether 仿射概形，$f : X \to S$ 固有，并且
$f_*\mathcal{O}_X = \mathcal{O}_S$。取 $S$ 的一点 $s \in S$。
需要证明 $X_s$ 几何连通。由
引理 \ref{more-morphisms-lemma-characterize-geometrically-connected-fibres}，
只需对每个 \'etale 邻域 $(U, u) \to (S, s)$ 证明 $X_u$ 连通。
可以假设 $U$ 仿射。于是 $U$ 是 Noether 的
（态射章引理 \ref{morphisms-lemma-finite-type-noetherian}），
基变换 $f_U : X_U \to U$ 固有
（态射章引理 \ref{morphisms-lemma-base-change-proper}），并且
$(f_U)_*\mathcal{O}_{X_U} = \mathcal{O}_U$
（概形上同调章引理
\ref{coherent-lemma-flat-base-change-cohomology}）。
因此，用基变换 $(f_U : X_U \to U, u)$ 替换
$(f : X \to S, s)$ 后，只需在 $f_*\mathcal{O}_X = \mathcal{O}_S$
时证明纤维 $X_s$ 连通。可以从 概形的导出范畴章引理
\ref{perfect-lemma-proper-idempotent-on-fibre} 推出这一点
（考察 $X_s$ 结构层中的幂等元），不过下面还将给出直接论证。

\medskip\noindent
具体地，应用形式函数定理，更确切地说，应用 概形上同调章引理
\ref{coherent-lemma-formal-functions-stalk}。它给出
$$
\mathcal{O}^\wedge_{S, s} = (f_*\mathcal{O}_X)_s^\wedge =
\lim_n H^0(X_n, \mathcal{O}_{X_n})
$$
其中 $X_n$ 是 $X_s$ 的第 $n$ 个无穷小邻域。由于 $X_n$ 的底层拓扑空间
等于 $X_s$ 的底层拓扑空间，若 $X_s = T_1 \amalg T_2$ 是非空开闭子概形的
不交并，则对所有 $n$，同样有
$X_n = T_{1, n} \amalg T_{2, n}$。这又意味着
$H^0(X_n, \mathcal{O}_{X_n})$ 含有非平凡幂等元 $e_{1, n}$，即在
$T_{1, n}$ 上恒等于 $1$、在 $T_{2, n}$ 上恒等于 $0$ 的函数。
显然，$e_{1, n + 1}$ 在 $X_n$ 上限制为 $e_{1, n}$。所以
$e_1 = \lim e_{1, n}$ 是该极限的非平凡幂等元。这与
$\mathcal{O}^\wedge_{S, s}$ 为局部环矛盾。因此假设错误，即 $X_s$
连通，结论得证。
\end{proof}

\begin{theorem}[Stein 分解；一般情形]
\label{more-morphisms-theorem-stein-factorization-general}
设 $S$ 是概形，$f : X \to S$ 是固有态射。则存在分解
$$
\xymatrix{
X \ar[rr]_{f'} \ar[rd]_f & & S' \ar[dl]^\pi \\
& S &
}
$$
具有下列性质：
\begin{enumerate}
\item 态射 $f'$ 固有，并具有几何连通纤维；
\item 态射 $\pi : S' \to S$ 是整态射；
\item 有 $f'_*\mathcal{O}_X = \mathcal{O}_{S'}$；
\item 有 $S' = \underline{\Spec}_S(f_*\mathcal{O}_X)$；
\item $S'$ 是 $S$ 在 $X$ 中的正规化；见
态射章定义 \ref{morphisms-definition-normalization-X-in-Y}.
\end{enumerate}
\end{theorem}

\begin{proof}
应用 引理 \ref{more-morphisms-lemma-stein-universally-closed}，得到态射
$f' : X \to S'$。注意，除 引理 \ref{more-morphisms-lemma-stein-universally-closed}
的结论外，$f'$ 还是分离的
（概形章引理 \ref{schemes-lemma-compose-after-separated}）且有限型
（态射章引理 \ref{morphisms-lemma-permanence-finite-type}）。
因此 $f'$ 固有。至此，除 $f'$ 具有几何连通纤维这一陈述外，
已经证明了所有陈述。

\medskip\noindent
可以假设 $S = \Spec(R)$ 仿射。令
$R' = \Gamma(X, \mathcal{O}_X)$。则 $S' = \Spec(R')$。
所以可以用 $S'$ 替换 $S$，并假设
% P05-EN-0116: frozen source omits conjunction/punctuation between affine and R=Gamma(X,O_X).
$S = \Spec(R)$ 仿射且 $R = \Gamma(X, \mathcal{O}_X)$。
接着，取一点 $s \in S$。设 $U \to S$ 是仿射概形的 \'etale 态射，
$u \in U$ 是映到 $s$ 的一点。令 $X_U \to U$ 为 $X$ 的基变换。
由 引理 \ref{more-morphisms-lemma-characterize-geometrically-connected-fibres}，
只需证明 $X_U \to U$ 在 $u$ 上的纤维连通。由
概形上同调章引理 \ref{coherent-lemma-flat-base-change-cohomology}，
有 $\Gamma(X_U, \mathcal{O}_{X_U}) = \Gamma(U, \mathcal{O}_U)$。
因此需要证明：给定仿射概形 $S = \Spec(R)$、满足
$\Gamma(X, \mathcal{O}_X) = R$ 的固有态射 $X \to S$，以及一点
$s \in S$，纤维 $X_s$ 连通。

\medskip\noindent
为此，只需说明幂等元 $e \in H^0(X_s, \mathcal{O}_{X_s})$ 只有
$0$ 和 $1$（由 引理 \ref{more-morphisms-lemma-stein-universally-closed}，已经知道
$X_s$ 非空）。由 概形的导出范畴章引理
\ref{perfect-lemma-proper-idempotent-on-fibre}，用一个主局部化替换 $R$ 后，
可以假设 $e$ 是 $R$ 中某个元素的像。由于
$R \to H^0(X_s, \mathcal{O}_{X_s})$ 经由 $\kappa(s)$ 分解，结论即得。
\end{proof}

\noindent
下面给出一个应用。

\begin{lemma}
\label{more-morphisms-lemma-geometrically-connected-fibres-towards-normal}
设 $f : X \to S$ 是概形态射。假设
\begin{enumerate}
\item $f$ 固有；
\item $S$ 是泛点为 $\xi$ 的整概形；
\item $S$ 正规；
\item $X$ 既约；
\item $X$ 的每个不可约分支的泛点均映到 $\xi$；
\item 有 $H^0(X_\xi, \mathcal{O}) = \kappa(\xi)$。
\end{enumerate}
则 $f_*\mathcal{O}_X = \mathcal{O}_S$，并且 $f$ 具有几何连通纤维。
\end{lemma}

\begin{proof}
应用 定理 \ref{more-morphisms-theorem-stein-factorization-general}，得到分解
$X \to S' \to S$。只需证明 $S' = S$。这将由 态射章引理
\ref{morphisms-lemma-finite-birational-over-normal} 得到。
具体地，由于 $X$ 既约，$S'$ 既约
（态射章引理 \ref{morphisms-lemma-normalization-in-reduced}）。
由上述定理，态射 $S' \to S$ 是整态射。由 态射章引理
\ref{morphisms-lemma-normalization-generic} 及假设 (5)，$S'$ 的每个泛点
都位于 $\xi$ 上方。另一方面，由于 $S'$ 是 $f_*\mathcal{O}_X$ 的相对谱，
概形论纤维 $S'_\xi$ 是 $H^0(X_\xi, \mathcal{O})$ 的谱，而按假设后者
等于 $\kappa(\xi)$。因此 $S'$ 是整概形，其函数域等于 $S$ 的函数域。
证明完毕。
\end{proof}

\noindent
下面是另一个应用。

\begin{lemma}
\label{more-morphisms-lemma-proper-flat-nr-geom-conn-comps-lower-semicontinuous}
设 $X \to S$ 是有限表示的平坦固有态射。令 $n_{X/S}$ 为 $S$ 上计数
% P05-EN-0117: frozen source refers to fibres of undefined f and says "the numbers";
% the named morphism is X -> S and the function counts the number.
$f$ 的纤维的几何连通分支个数的函数，它引入于
引理 \ref{more-morphisms-lemma-base-change-fibres-nr-geometrically-connected-components}.
则 $n_{X/S}$ 下半连续。
\end{lemma}

\begin{proof}
取 $s \in S$，令 $n = n_{X/S}(s)$。注意 $n < \infty$：
$X \to S$ 在 $s$ 处的几何纤维是域上的固有概形，因而是 Noether 的，
故只有有限多个连通分支。需要找到 $s$ 的开邻域 $V$，使得
$n_{X/S}|_V \geq n$。令 $X \to S' \to S$ 为
定理 \ref{more-morphisms-theorem-stein-factorization-general} 中的 Stein 分解。
由 引理 \ref{more-morphisms-lemma-stein-universally-closed-residue-fields}，位于 $s$
上方的点只有有限多个 $s'_1, \ldots, s'_m \in S'$，并且扩张
$\kappa(s'_i)/\kappa(s)$ 有限。引理 \ref{more-morphisms-lemma-etale-makes-integral-split}
于是表明：用 $s$ 的一个 \'etale 邻域替换 $S$ 后，可以假设作为概形
$S' = V_1 \amalg \ldots \amalg V_m$，且 $s'_i \in V_i$、
$\kappa(s'_i)/\kappa(s)$ 纯不可分。于是概形 $X_{s_i'}$ 在
$\kappa(s)$ 上几何连通，故 $m = n$。概形
$X_i = (f')^{-1}(V_i)$，$i = 1, \ldots, n$，在 $S$ 上平坦且有限表示。
因此 $X_i \to S$ 的像是开的
（态射章引理 \ref{morphisms-lemma-fppf-open}）。
由此可见，在 $s$ 的一个邻域内，$n_{X/S}$ 至少为 $n$。
\end{proof}

\begin{lemma}
\label{more-morphisms-lemma-proper-flat-geom-red}
设 $f : X \to S$ 是概形态射。假设
\begin{enumerate}
\item $f$ 固有、平坦且有限表示；
\item $f$ 的几何纤维既约。
\end{enumerate}
则计数 $f$ 的纤维的几何连通分支个数的函数
$n_{X/S} : S \to \mathbf{Z}$ 局部常值。
\end{lemma}

\begin{proof}
由 引理 \ref{more-morphisms-lemma-proper-flat-nr-geom-conn-comps-lower-semicontinuous}，
% P05-EN-0118: frozen source misspells "semicontinuous" as "semincontinuous".
函数 $n_{X/S}$ 下半连续。对 $s \in S$，考虑 $\kappa(s)$-代数
$$
A_s = H^0(X_s, \mathcal{O}_{X_s})
$$
% P05-EN-0119: frozen source defines beta_0 on X but then sends s in S and compares functions on S;
% the evident domain is S. Formula retained.
定义函数 $\beta_0 : X \to \mathbf{Z}$，将 $s$ 映到
$\dim_{\kappa(s)} A_s$。由 簇章引理
\ref{varieties-lemma-proper-geometrically-reduced-global-sections} 以及
$X_s$ 几何既约这一事实，$A_s$ 是 $\kappa(s)$ 的有限可分扩张的有限积。
因此 $A_s \otimes_{\kappa(s)} \kappa(\overline{s})$ 是
$\beta_0(s)$ 个 $\kappa(\overline{s})$ 的积。故 $X_{\overline{s}}$
有 $\beta_0(s)$ 个连通分支。换言之，作为 $S$ 上的函数，有
$n_{X/S} = \beta_0$。所以由 概形的导出范畴章引理
\ref{perfect-lemma-jump-loci-geometric}，$n_{X/S}$ 上半连续。证明完毕。
\end{proof}

\noindent
最后一个应用。

\begin{lemma}
\label{more-morphisms-lemma-split-off-proper-part-henselian}
\begin{reference}
Adic Noether 基的情形可参见
\cite[III, Proposition 5.5.1]{EGA}。
\end{reference}
设 $(A, I)$ 是 Hensel 对。设 $X \to \Spec(A)$ 分离且有限型。
令 $X_0 = X \times_{\Spec(A)} \Spec(A/I)$。设 $Y \subset X_0$
是开闭子概形，并且 $Y \to \Spec(A/I)$ 固有。则存在在 $A$ 上固有的
开闭子概形 $W \subset X$，使得
$W \times_{\Spec(A)} \Spec(A/I) = Y$。
\end{lemma}

\begin{proof}
以 $T \mapsto T_0$ 表示沿 $\Spec(A/I) \to \Spec(A)$ 的基变换。
由 Chow 引理（采用 极限章引理 \ref{limits-lemma-chow-finite-type}
的形式），存在满固有态射 $\varphi : X' \to X$，使得 $X'$ 可浸入
$\mathbf{P}^n_A$。令 $Y' = \varphi^{-1}(Y)$。这是 $X'_0$ 的开闭子概形。
假设本引理对 $(X', Y')$ 成立。令 $W' \subset X'$ 为在 $A$ 上固有且满足
$Y' = W'_0$ 的开闭子概形。由 态射章引理
\ref{morphisms-lemma-image-proper-scheme-closed}，
$W = \varphi(W') \subset X$ 和 $Q = \varphi(X' \setminus W') \subset X$
是闭子集；由 态射章引理
\ref{morphisms-lemma-image-proper-is-proper}，$W$ 在 $A$ 上固有。
$W \cap Q$ 在 $\Spec(A)$ 中的像是闭的。由于 $(A, I)$ 是 Hensel 对，
若 $W \cap Q$ 非空，则 $W \cap Q$ 有一个位于 $\Spec(A/I)$ 上方的点。
但这是不可能的，因为 $W'_0 = Y' = \varphi^{-1}(Y)$。
所以 $W$ 是 $X$ 的开闭子概形，在 $A$ 上固有，并满足 $W_0 = Y$。
由此化归到下一段所述的情形。

\medskip\noindent
假设存在 $A$ 上的浸入 $j : X \to \mathbf{P}^n_A$。令
$\overline{X}$ 为 $j$ 的概形论像。由于 $j$ 是拟紧态射
（概形章引理 \ref{schemes-lemma-quasi-compact-permanence}），
可见 $j : X \to \overline{X}$ 是开浸入
（态射章引理 \ref{morphisms-lemma-quasi-compact-immersion}）。
因此，基变换 $j_0 : X_0 \to \overline{X}_0$ 也是开浸入。
所以 $j_0(Y) \subset \overline{X}_0$ 是开的。由 态射章引理
\ref{morphisms-lemma-image-proper-scheme-closed}，它也是闭的。
假设本引理对 $(\overline{X}, j_0(Y))$ 成立。令
$\overline{W} \subset \overline{X}$ 为相应的、在 $A$ 上固有的开闭子概形，
使得 $j_0(Y) = \overline{W}_0$。于是
$T = \overline{W} \setminus j(X)$ 在 $\overline{W}$ 中闭；由于
$\overline{W}$ 在 $A$ 上固有，该子集在 $\Spec(A)$ 中的像是闭的。
由于 $(A, I)$ 是 Hensel 对，若 $T$ 非空，则 $T$ 中有一点映到
$\Spec(A/I)$。这是不可能的，因为
$j_0(Y) = \overline{W}_0$ 包含于 $j(X)$。故 $\overline{W}$ 包含于
$j(X)$，并可令 $W \subset X$ 等于唯一一个经 $j$ 同构映到
$\overline{W}$ 的开闭子概形。由此化归到下一段所述的情形。

\medskip\noindent
假设 $X \subset \mathbf{P}^n_A$ 是闭子概形。则
$X \to \Spec(A)$ 是固有态射。令 $Z = X_0 \setminus Y$。
这是 $X_0$ 的开闭子概形，并且 $X_0 = Y \amalg Z$。
令 $X \to X' \to \Spec(A)$ 为
定理 \ref{more-morphisms-theorem-stein-factorization-general} 中的 Stein 分解。
令 $Y' \subset X'_0$ 与 $Z' \subset X'_0$ 分别为 $Y$ 与 $Z$ 的像。
% P05-EN-0120: frozen source says fibres of X -> Z; the Stein map with connected fibres is X -> X'.
由于 $X \to Z$ 的纤维几何连通，可见 $Y' \cap Z' = \emptyset$。
又因为 $X \to X'$ 满，所以 $X'_0 = Y' \amalg Z'$。
由于 $X' \to \Spec(A)$ 是整态射，可见 $X'$ 是某个在 $A$ 上整的
$A$-代数的谱。回忆，谱的开闭子集与相应环中的幂等元 $1$-到-$1$ 对应；见
代数章引理 \ref{algebra-lemma-disjoint-decomposition}。
因此，由 代数进阶章引理
\ref{more-algebra-lemma-characterize-henselian-pair}，可以写成
$X' = W' \amalg V'$，其中 $W'$ 与 $V'$ 开且闭，并且
$Y' = W'_0$、$Z' = V'_0$。令 $W$ 为其在 $X$ 中的逆像，即完成证明。
\end{proof}









\section{泛平坦性分层}
\label{more-morphisms-section-generic-flatness-stratification}

\noindent
可以利用泛平坦性构造基的一个分层，使给定模在各层上成为平坦模。

\begin{lemma}[泛平坦性分层]
\label{more-morphisms-lemma-generic-flatness-stratification}
设 $f : X \to S$ 是拟紧拟分离概形之间的有限表示态射，
$\mathcal{F}$ 是有限表示 $\mathcal{O}_X$-模。则存在 $t \geq 0$
及闭子概形
$$
S \supset S_0 \supset S_1 \supset \ldots \supset S_t = \emptyset
$$
使得 $S_i \to S$ 由有限型理想层定义，$S_0 \subset S$ 是增厚，
并且 $\mathcal{F}$ 到 $X \times_S (S_i \setminus S_{i + 1})$ 的拉回
在 $S_i \setminus S_{i + 1}$ 上平坦。
\end{lemma}

\begin{proof}
可以找到笛卡尔图表
$$
\xymatrix{
X \ar[d] \ar[r] & X_0 \ar[d] \\
S \ar[r] & S_0
}
$$
以及有限表示 $\mathcal{O}_{X_0}$-模 $\mathcal{F}_0$，其拉回为
$\mathcal{F}$，并且 $X_0$ 与 $S_0$ 在 $\mathbf{Z}$ 上有限型。见
极限章命题 \ref{limits-proposition-approximate} 及
引理 \ref{limits-lemma-descend-finite-presentation} 及
\ref{limits-lemma-descend-modules-finite-presentation}。
因此，可以假设 $X$ 与 $S$ 在 $\mathbf{Z}$ 上有限型，而
$\mathcal{F}$ 是凝聚 $\mathcal{O}_X$-模。

\medskip\noindent
假设 $X$ 与 $S$ 在 $\mathbf{Z}$ 上有限型，而 $\mathcal{F}$ 是凝聚
$\mathcal{O}_X$-模。在此情形，每个拟凝聚理想都是有限型的，故无需检查
$S_i$ 由有限型理想切出这一条件。令 $S_0 = S_{red}$ 为 $S$ 的既约化。
由 态射章命题
\ref{morphisms-proposition-generic-flatness-reduced} 所述的泛平坦性，
存在稠密开集 $U_0 \subset S_0$，使得 $\mathcal{F}$ 到
$X \times_S U_0$ 的拉回在 $U_0$ 上平坦。令 $S_1 \subset S_0$
为底层闭子集等于 $S \setminus U_0$ 的既约闭子概形。若
$S_1 \not = \emptyset$，便依此继续，得到
$S_0 \supset S_1 \supset \ldots$。由于 $S$ 是 Noether 的，
任意闭子集降链均稳定，所以对某个 $t \geq 0$ 有
$S_t = \emptyset$。
\end{proof}

\begin{lemma}
\label{more-morphisms-lemma-generic-flatness-stratification-scheme}
设 $f : X \to S$ 是拟紧拟分离概形之间的有限表示态射。
则存在 $t \geq 0$ 及闭子概形
$$
S \supset S_0 \supset S_1 \supset \ldots \supset S_t = \emptyset
$$
使得 $S_i \to S$ 由有限型理想层定义，$S_0 \subset S$ 是增厚，并且
$X \times_S (S_i \setminus S_{i + 1})$ 在
$S_i \setminus S_{i + 1}$ 上平坦。
\end{lemma}

\begin{proof}
在 $\mathcal{F} = \mathcal{O}_X$ 的情形应用
引理 \ref{more-morphisms-lemma-generic-flatness-stratification}。
\end{proof}

\begin{lemma}[Longke Tang]
\label{more-morphisms-lemma-flatness-criterion}
设 $f : X \to S$ 是有限表示的满态射。设 $M$ 是
$D_\QCoh(\mathcal{O}_S)$ 的对象，并且对 $i > 0$ 有
$H^i(M) = 0$。下列条件等价：
\begin{enumerate}
\item $M$ 同构于置于次数 $0$ 的平坦 $\mathcal{O}_S$-模；
\item $Lf^*M$ 同构于置于次数 $0$ 的平坦 $\mathcal{O}_X$-模。
\end{enumerate}
\end{lemma}

\begin{proof}
略去 (1) 蕴含 (2) 的证明。假设 (2)。证明 (1) 是局部问题，
故可假设 $S$ 仿射。按
引理 \ref{more-morphisms-lemma-generic-flatness-stratification-scheme} 选取分层
$S \supset S_0 \supset S_1 \supset \ldots \supset S_t = \emptyset$
。
% P05-EN-0121: frozen source defines T_i=S_i\S_{i-1}; the strata above are S_i\S_{i+1},
% and i=0 would require undefined S_{-1}. Formula retained.
令 $T_i = S_i \setminus S_{i - 1}$、$Y_i = X \times_S T_i$。
以 $g_i : T_i \to S$ 和 $h_i : Y_i \to X$ 表示包含态射。
% P05-EN-0122: frozen source uses undefined X_i as source of f_i after defining Y_i.
由于 $f$ 满，态射 $f_i : X_i \to T_i$ 满且平坦。由于
$Lh_i^* \circ Lf^* = Lf_i^* \circ Lg_i^*$，由 (2) 断定
$Lf_i^* Lg_i^*M$ 同构于置于次数 $0$ 的平坦 $\mathcal{O}_{Y_i}$-模。
由于 $f_i$ 平坦且满，可见 $Lg_i^*M$ 同构于置于次数 $0$ 的平坦
$\mathcal{O}_{T_i}$-模。

\medskip\noindent
我们对 $t$ 归纳证明这蕴含所需结论。采用如下记号：
$S = \Spec(A)$，并且 $M$ 是 $D_\QCoh(\mathcal{O}_S)$ 中与
$D(A)$ 中的 $N$ 对应的对象；见 概形的导出范畴章引理
\ref{perfect-lemma-affine-compare-bounded}。

\medskip\noindent
基础情形：$t = 1$。写成 $S_0 = \Spec(A/I)$，其中 $I$ 是有限生成理想。
由于在集合论意义下 $S_0 = S$，$I$ 幂零。假设表明
$N \otimes_A^\mathbf{L} A/I$ 的 Tor 振幅在 $[0, 0]$ 中。
由 代数进阶章引理
\ref{more-algebra-lemma-check-tor-dimension-modulo-nilpotent}，
$N$ 亦然。

\medskip\noindent
归纳步骤。假设 $t > 1$。写成 $S_{t - 1} = \Spec(A/I)$，其中
$I = (f_1, \ldots, f_r)$ 是有限生成理想。对 $r$ 归纳论证。
注意，由归纳假设，$D(A_{f_r})$ 中的
$N_{f_r} = N \otimes_A^\mathbf{L} A_{f_r}$ 的 Tor 振幅在
$[0, 0]$ 中（因为主开集 $D(f_r)$ 上的分层长度至多为 $t - 1$）。
另一方面，考虑 $D(A/f_rA)$ 中的
$N' = N \otimes_A^\mathbf{L} A/f_rA$。令
$S' = \Spec(A/f_rA)$，得到分层
$$
S' \supset S' \cap S_0 \supset S' \cap S_1 \supset \ldots \supset
S' \cap S_{t - 1} \supset
S' \cap S_t = \emptyset
$$
其中 $S' \cap S_{t - 1}$ 在 $S'$ 中由 $r - 1$ 个方程切出。
与先前类似，$M$ 到各部分
$S' \cap S_i \setminus S' \cap S_{i + 1}$ 的导出拉回的 Tor 振幅均在
$[0, 0]$ 中。因此，由对 $r$ 的归纳，$N'$ 的 Tor 振幅在
$[0, 0]$ 中。最后，由 代数进阶章引理
\ref{more-algebra-lemma-f-flat}，断定 $N$ 的 Tor 振幅在 $[0, 0]$ 中。
\end{proof}

\begin{lemma}
\label{more-morphisms-lemma-cokernel-flat}
设 $R$ 是 Noether 整环，$R \to A \to B$ 是有限型环映射。
设 $M$ 是有限 $A$-模，$N$ 是有限 $B$-模，$M \to N$ 是
$A$-线性映射。则存在非零的 $f \in R$，使得
$M_f \to N_f$ 的余核是平坦 $R_f$-模。
\end{lemma}

\begin{proof}
用 $M \to N$ 的像替换 $M$ 后，可以假设 $M \subset N$。
选取滤过 $0 = N_0 \subset N_1 \subset \ldots \subset N_t = N$，
使得对某个素理想 $\mathfrak q_i \subset B$，有
$N_i/N_{i - 1} = B/\mathfrak q_i$；见
代数章引理 \ref{algebra-lemma-filter-Noetherian-module}。
令 $M_i = M \cap N_i$。则 $Q = N/M$ 具有由子模
$Q_i = N_i/M_i$ 给出的滤过。只需证明在某个非零的 $f$ 的元素处局部化后，
% P05-EN-0123: frozen source says "a nonzero element of f"; the element should be of R.
$Q_i/Q_{i - 1}$ 成为平坦模（因为由 代数章引理
\ref{algebra-lemma-flat-ses}，平坦模的扩张平坦）。由于
$Q_i/Q_{i - 1}$ 同构于映射
$M_i/M_{i - 1} \to N_i/N_{i - 1}$ 的余核，故化归到下一段讨论的情形。

\medskip\noindent
假设 $B$ 是整环，且 $M \subset N = B$。用 $A$ 在 $B$ 中的像替换
$A$ 后，可以假设 $A \subset B$。由泛平坦性，可假设 $A$ 与 $B$
在 $R$ 上平坦
（代数章引理 \ref{algebra-lemma-generic-flatness-Noetherian}）。
现在，只需说明：用一个主局部化替换 $R$ 后，$M \to B$ 成为
$R$-普遍单射
（代数章引理 \ref{algebra-lemma-ui-flat-domain}）。
由泛自由性，可以找到非零的 $g \in A$，使得 $B_g$ 是自由
$A_g$-模（代数章引理
\ref{algebra-lemma-generic-flatness-Noetherian}）。
因此，可以选取作为 $A_g$-模的直和项 $M' \subset B_g$，使它是有限自由
$A_g$-模，并且 $M \to B \to B_g$ 经由 $M'$ 分解。
显然，只需说明：用一个主局部化替换 $R$ 后，$M \to M'$ 成为
$R$-普遍单射。

\medskip\noindent
设 $M' = A_g^{\oplus n}$。由于 $M \subset M'$ 是有限 $A$-模，
可见对某个 $m \geq 0$，$M$ 包含于 $(1/g^m)A^{\oplus n}$。
改变 $M'$ 的基后，可以假设 $M \subset A^{\oplus n}$。
于是只需说明：用一个主局部化替换 $R$ 后，$A^{\oplus n}/M$ 与
$A_g/A$ 成为 $R$-平坦模。事实上，此时由 代数章引理
\ref{algebra-lemma-flat-tor-zero}，$M \to A^{\oplus n}$ 与
$A^{\oplus n} \to A_g^{\oplus n}$ 普遍单射，因而它们的复合
$M \to M' = A_g^{\oplus n}$ 也普遍单射。

\medskip\noindent
由泛平坦性（见上述引文），可以假设模 $A^{\oplus n}/M$ 是
$R$-平坦的。对商 $A_g/A$，使用事实
$$
A_g/A = \colim (1/g^m)A/A \cong \colim A/g^mA
$$
以及模 $A/g^mA$ 具有长度为 $m$ 的滤过，其相继商均同构于 $A/gA$。
再次由泛平坦性，可以假设 $A/gA$ 是 $R$-平坦的，因而每个
$A/g^mA$ 都是 $R$-平坦的，从而 $A_g/A$ 亦然。
\end{proof}

\noindent
设 $f : X \to Y$ 是基概形 $S$ 上的概形态射。令 $Z \subset Y$
为 $f$ 的概形论像；见
态射章第 \ref{morphisms-section-scheme-theoretic-image} 节。
设 $g : S' \to S$ 是概形态射，并令
$f' : X \times_S S' \to Y \times_S S'$ 为 $f$ 沿 $g$ 的基变换。
$Z \times_S S' \subset Y \times_S S'$ 未必是 $f'$ 的概形论像。
若对每个上述 $g$，$f'$ 的概形论像都等于 $Z \times_S S'$，
则称{\it $f/S$ 的概形论像的形成与任意基变换交换}。

\begin{lemma}
\label{more-morphisms-lemma-base-change-scheme-theoretic-image}
设 $S$ 是拟紧拟分离概形。设 $f : X \to Y$ 是 $S$ 上的概形态射，
且 $X$ 与 $Y$ 均在 $S$ 上有限表示。则存在 $t \geq 0$ 及闭子概形
$$
S \supset S_0 \supset S_1 \supset \ldots \supset S_t = \emptyset
$$
具有下列性质：
\begin{enumerate}
\item $S_i \to S$ 由有限型理想层定义；
\item $S_0 \subset S$ 是增厚；
\item 令 $T_i = S_i \setminus S_{i + 1}$，并令 $f_i$ 为 $f$ 到
$T_i$ 的基变换，则 $f_i/T_i$ 的概形论像的形成与任意基变换交换
（见引理上方的讨论）。
\end{enumerate}
\end{lemma}

\begin{proof}
可以找到交换图表
$$
\xymatrix{
X \ar[d] \ar[r] & Y \ar[d] \ar[r] & S \ar[d] \\
U \ar[r] & V \ar[r] & W
}
$$
，其中两个方块均笛卡尔，且 $U$、$V$、$W$ 在 $\mathbf{Z}$ 上有限型。
具体地，首先利用 极限章命题 \ref{limits-proposition-approximate}，
将 $S$ 写成转移态射仿射的 $\mathbf{Z}$ 上有限型概形的余滤极限，
再利用 极限章引理 \ref{limits-lemma-descend-finite-presentation}
下降态射 $X \to Y$。由此化归到下一段讨论的情形。

\medskip\noindent
假设 $S$ 是 Noether 的。在此情形，每个拟凝聚理想都是有限型的，
故无需检查 $S_i$ 由有限型理想切出这一条件。令
$S_0 = S_{red}$ 为 $S$ 的既约化。令 $\eta \in S_0$ 为 $S_0$
的某个不可约分支的泛点。对 $S_0$ 的底层拓扑空间作 Noether 归纳，
可以假设结论对 $S_0$ 的任何不包含 $\eta$ 的闭子概形均成立。
所以只需说明存在开邻域 $U_0 \subset S_0$，使得 $f$ 到 $U_0$
的基变换 $f_0$ 具有性质 (3)。

\medskip\noindent
设 $R$ 是 Noether 整环，$f : X \to Y$ 是 $R$ 上有限型概形的态射。
由上一段的讨论，只需说明：对某个非零的 $g \in R$，用 $R_g$ 替换
$R$，并用到 $R_g$ 的基变换替换 $X$、$Y$ 后，$f/R$ 的概形论像的形成
与任意基变换交换。

\medskip\noindent
令 $Y = V_1 \cup \ldots V_n$ 为仿射开覆盖，并令
$U_i = f^{-1}(V_i)$。若该陈述对每个 $R$ 上的态射
$U_i \to V_i$ 均成立，则它对 $f$ 成立。事实上，由 态射章引理
\ref{morphisms-lemma-quasi-compact-scheme-theoretic-image}，
$U_i \to V_i$ 的概形论像是 $V_i$ 与 $f : X \to Y$ 的概形论像之交。
所以可以假设 $Y$ 仿射。

\medskip\noindent
令 $X = U_1 \cup \ldots U_n$ 为仿射开覆盖。则 $X \to Y$ 的概形论像
与 $\coprod U_i \to Y$ 的概形论像相同。
% P05-EN-0124: frozen source misspells "image" as "imge".
因此可以假设 $X$ 仿射。

\medskip\noindent
写成 $X = \Spec(A)$、$Y = \Spec(B)$，并设 $f$ 对应于 $R$-代数映射
% P05-EN-0125: frozen source reverses the affine contravariance: X=Spec(A)->Y=Spec(B)
% should correspond to B->A, and the image should be Spec(B/Ker). Formulas retained.
$\varphi : A \to B$。则 $f$ 的概形论像是
$\Spec(A/\Ker(\varphi))$；基变换后亦同（这里基变换采用仿射态射，
但只需对此检查）。因此，若对所有环映射 $R \to R'$ 都有
$\Ker(\varphi \otimes_R R') = \Ker(\varphi) \otimes_R R'$，
则概形论像的形成与基变换交换。

\medskip\noindent
对 $R$ 中适当的非零元素 $g$，分别用 $R_g$、$A_g$、$B_g$ 替换
$R$、$A$、$B$ 后，可以假设 $A$ 与 $B$ 在 $R$ 上平坦。
由 引理 \ref{more-morphisms-lemma-cokernel-flat}，还可假设 $B/A$ 是平坦 $R$-模。
于是 $0 \to \Ker(\varphi) \to A \to B \to B/A \to 0$
是平坦 $R$-模的正合列，这蕴含所需的基变换陈述。
\end{proof}






\section{对态射分层}
\label{more-morphisms-section-stratifying-morphisms}

\noindent
设 $f : X \to S$ 是拟紧拟分离概形之间的有限表现态射。
在 第 \ref{more-morphisms-section-generic-flatness-stratification} 节 中，
我们已经看到，可以对 $S$ 分层，使得 $X$ 在各层上平坦。
本节尝试同时对 $S$ 和 $X$ 分层，以得到光滑的各层；
这不能完全直接实现，我们还需要沿有限局部自由态射作基变换。

\begin{lemma}
\label{more-morphisms-lemma-smoothness-stratification-at-generic-point}
设 $f : X \to S$ 是有限表现的概形态射。
设 $\eta \in S$ 是 $S$ 的一个不可约分支的泛点。
假设 $S$ 既约。则存在
\begin{enumerate}
\item 包含 $\eta$ 的开子概形 $U \subset S$，
\item 满射、普遍单射的有限局部自由态射 $V \to U$，
\item 某个 $t \geq 0$ 以及闭子概形
$$
X \times_S V \supset Z_0 \supset Z_1 \supset \ldots \supset Z_t = \emptyset
$$
使得 $Z_i \to X \times_S V$ 由有限型理想层定义，
$Z_0 \subset X \times_S V$ 是一个增厚，并且态射
$Z_i \setminus Z_{i + 1} \to V$ 光滑。
\end{enumerate}
\end{lemma}

\begin{proof}
显然，我们可以用 $\eta$ 的一个开邻域替换 $S$，并用 $X$
在此开集上的限制替换之。因此可以假设 $S = \Spec(A)$，
其中 $A$ 是既约环，而 $\eta$ 对应于一个极小素理想 $\mathfrak p$。
回忆在这种情形下，局部环 $\mathcal{O}_{S, \eta} = A_\mathfrak p$
等于 $\kappa(\mathfrak p)$，见
代数章引理 \ref{algebra-lemma-minimal-prime-reduced-ring}.

\medskip\noindent
对 $k = \kappa(\eta)$ 上的概形 $X_\eta$ 应用
簇章引理 \ref{varieties-lemma-smooth-stratification}。
将由此得到的纯不可分域扩张记为 $k'/k$。
下一段把问题约化到 $k' = k$ 的情形。
（这一步对应于寻找引理陈述中的态射 $V \to U$；特别地，
若 $\kappa(\mathfrak p)$ 的特征为零，则可以取 $V = U$。）

\medskip\noindent
若 $k = \kappa(\mathfrak p)$ 的特征为零，则 $k' = k$。
若 $k = \kappa(\mathfrak p)$ 的特征为 $p > 0$，则 $p$
在 $A_\mathfrak p = \kappa(\mathfrak p)$ 中的像为零。
因此，用一个主局部化替换 $A$（即缩小 $S$）后，
可以假设在 $A$ 中有 $p = 0$。若 $k' \not = k$，则存在
$\beta \in k'$、$\beta \not \in k$，使得 $\beta^p \in k$。
用一个主局部化替换 $A$ 后，可以假设存在 $a \in A$，
使得 $\beta^p = a$。置 $A' = A[x]/(x^p - a)$。
于是 $S' = \Spec(A') \to \Spec(A) = S$ 是有限局部自由、满射且普遍单射的。
此外，若 $\mathfrak p' \subset A'$ 表示位于 $\mathfrak p$ 上方的唯一素理想，
则 $A'_{\mathfrak p'} = k(\beta)$，且 $k'/k(\beta)$ 的次数更小。
因此，用 $S'$ 替换 $S$，并用对应于 $\mathfrak p'$ 的点 $\eta'$
替换 $\eta$ 后，可见 $k'$ 在 $\eta$ 的剩余域上的次数已经减小。
如此继续并作归纳，就约化到
$k' = \kappa(\mathfrak p) = \kappa(\eta)$ 的情形。

\medskip\noindent
因此，可以假设 $S$ 仿射且既约，并且存在 $t \geq 0$ 及闭子概形
$$
X_\eta \supset Z_{\eta, 0} \supset Z_{\eta, 1}
\supset \ldots \supset Z_{\eta, t} = \emptyset
$$
使得 $Z_{\eta, 0} = (X_\eta)_{red}$，并且对所有 $i$，
$Z_{\eta, i} \setminus Z_{\eta, i + 1}$ 在 $\eta$ 上光滑。
回忆 $\kappa(\eta) = \kappa(\mathfrak p) = A_\mathfrak p$
是所有 $A_a$（其中 $a \in A$、$a \not \in \mathfrak p$）的滤过余极限。
见 代数章引理 \ref{algebra-lemma-localization-colimit}。
因此，对某个 $a \in A$、$a \not \in \mathfrak p$，
可以把上图下降为 $\Spec(A_a)$ 上的相应图。
更确切地，用 $\Spec(A_a)$ 替换 $S$ 后，可以假设存在
$t \geq 0$ 及闭子概形
$$
X \supset Z_0 \supset Z_1 \supset \ldots \supset Z_t = \emptyset
$$
使得 $Z_i \to X$ 是有限表现的闭浸入，$Z_0 \to X$ 是一个增厚，
并且 $Z_i \setminus Z_{i + 1}$ 在 $S$ 上光滑。换言之，引理成立。
更确切地，我们先用
极限章引理 \ref{limits-lemma-descend-finite-presentation}
得到态射
$$
Z_t \to Z_{t - 1} \to \ldots \to Z_0 \to X
$$
于 $S$ 上；每个态射都是有限表现的，并且它们到 $\eta$ 的基变换给出
上述给定闭子概形之间的包含。进一步缩小 $S$ 后，可以假设
每个态射都是闭浸入，见
极限章引理
\ref{limits-lemma-descend-closed-immersion-finite-presentation}.
缩小 $S$ 后，可以假设 $Z_0 \to X$ 满射，因而是一个增厚，见
极限章引理 \ref{limits-lemma-descend-surjective}.
再缩小一次 $S$，可以假设 $Z_i \setminus Z_{i + 1} \to S$ 光滑，见
极限章引理 \ref{limits-lemma-descend-smooth}.
证明完毕。
\end{proof}

\begin{lemma}
\label{more-morphisms-lemma-smoothness-stratification}
设 $f : X \to S$ 是拟紧拟分离概形之间的有限表现态射。
则存在 $t \geq 0$ 及闭子概形
$$
S \supset S_0 \supset S_1 \supset \ldots \supset S_t = \emptyset
$$
使得
\begin{enumerate}
\item $S_i \to S$ 由有限型理想层定义，
\item $S_0 \subset S$ 是一个增厚，
\item 对每个 $i$，存在满射的有限局部自由态射
$T_i \to S_i \setminus S_{i + 1}$，
\item 对每个 $i$，存在 $t_i \geq 0$ 及闭子概形
$$
X_i = X \times_S T_i \supset Z_{i, 0}
\supset Z_{i, 1} \supset \ldots \supset Z_{i, t_i} = \emptyset
$$
使得 $Z_{i, j} \to X_i$ 由有限型理想层定义，
$Z_{i, 0} \subset X_i$ 是一个增厚，并且态射
$Z_{i, j} \setminus Z_{i, j + 1} \to T_i$ 光滑。
\end{enumerate}
\end{lemma}

\begin{proof}
可以找到笛卡尔图
$$
\xymatrix{
X \ar[d] \ar[r] & X_0 \ar[d] \\
S \ar[r] & S_0
}
$$
使得 $X_0$ 和 $S_0$ 在 $\mathbf{Z}$ 上为有限型。见
极限章命题 \ref{limits-proposition-approximate} 及
引理 \ref{limits-lemma-descend-finite-presentation}.
因此，可以假设 $X$ 和 $S$ 在 $\mathbf{Z}$ 上为有限型。
事实上，引理对 $X_0 \to S_0$ 所提出问题的一个解，
经基变换后就给出 $S$ 上的一个解；细节略。

\medskip\noindent
假设 $X$ 和 $S$ 在 $\mathbf{Z}$ 上为有限型。
在这种情形下，每个拟凝聚理想都是有限型的，因而无须检验
$S_i$ 由有限型理想截出的条件。置 $S_0 = S_{red}$，即 $S$ 的既约化。
设 $\eta \in S_0$ 是一个不可约分支的泛点。
由 引理 \ref{more-morphisms-lemma-smoothness-stratification-at-generic-point}，
可以找到开子概形 $U \subset S_0$、满射且普遍单射的有限局部自由态射
$V \to U$、某个 $t_0 \geq 0$ 以及闭子概形
$$
X \times_S V \supset Z_{0, 0} \supset Z_{0, 1} \supset \ldots \supset
Z_{0, t_0} = \emptyset
$$
使得 $Z_{0, i} \to X \times_S V$ 由有限型理想层定义，
$Z_{0, 0} \subset X \times_S V$ 是一个增厚，并且态射
$Z_{0, i} \setminus Z_{0, i + 1} \to V$ 光滑。
然后，令 $S_1 \subset S_0$ 是 $S_0 \setminus U$ 上的既约诱导子概形结构。
对 $S$ 的底拓扑空间作诺特归纳，可以假设引理对
$X \times_S S_1 \to S_1$ 成立。这给出 $t \geq 1$ 以及
$$
S_1 = S_1 \supset S_2 \supset \ldots \supset S_t = \emptyset
$$
以及引理陈述中的 $t_i$ 和 $Z_{i, j}$。
这就证明了引理。
\end{proof}


















\section{改进相对维数为一的态射}
\label{more-morphisms-section-make-good-curves}

\noindent
通过扩张基域、紧化及正规化，可以使任意曲线成为光滑射影曲线。
这还蕴含关于泛纤维维数为 $1$ 的有限型态射的结果。

\begin{lemma}
\label{more-morphisms-lemma-make-good-curves}
设 $f : X \to S$ 是概形态射。设 $\eta \in S$ 是 $S$
的一个不可约分支的泛点。假设 $f$ 分离、有限表现，且
$\dim(X_\eta) \leq 1$。则存在交换图
$$
\xymatrix{
\overline{Y}_1 \amalg \ldots \amalg \overline{Y}_n \ar[rd] &
Y_1 \amalg \ldots \amalg Y_n \ar[r]_-\nu \ar[d] \ar[l]^j &
X_V \ar[r] \ar[d] &
X_U \ar[r] \ar[d] &
X \ar[d]^f \\
& T_1 \amalg \ldots \amalg T_n \ar[r] &
V \ar[r] &
U \ar[r] &
S
}
$$
满足下列性质：
\begin{enumerate}
\item $U \subset S$ 是 $\eta$ 的一个开邻域，
\item $V \to U$ 是有限、满射且普遍单射的态射，
\item $X_U = U \times_S X$ 和 $X_V = V \times_S X$ 是相应的基变换，
\item $\nu$ 有限且满射，并且存在开集 $W \subset X_V$，使得
\begin{enumerate}
\item $W$ 在 $X_V \to V$ 的每个纤维中稠密，
\item $\nu^{-1}(W) \cap Y_i$ 在 $Y_i \to T_i$ 的每个纤维中稠密，并且
\item $\nu^{-1}(W) \to W$ 是一个增厚，
\end{enumerate}
\item $j$ 是开浸入，
\item $T_i \to V$ 是有限 \'etale 态射，
\item $Y_i \to T_i$ 满射且光滑，
\item $\overline{Y}_i \to T_i$ 光滑且固有，其纤维几何连通且维数 $\leq 1$。
\end{enumerate}
\end{lemma}

\begin{proof}
显然，我们可以用 $\eta$ 的一个开邻域替换 $S$，并用 $X$
在此开集上的限制替换之。此外，可以用 $S$ 的既约化替换它，
并用 $X$ 到此既约化的基变换替换之。
因此可以假设 $S = \Spec(A)$，其中 $A$ 是既约环，
而 $\eta$ 对应于一个极小素理想 $\mathfrak p$。
回忆在这种情形下，局部环 $\mathcal{O}_{S, \eta} = A_\mathfrak p$
等于 $\kappa(\mathfrak p)$，见
代数章引理 \ref{algebra-lemma-minimal-prime-reduced-ring}.

\medskip\noindent
对 $k = \kappa(\eta)$ 上的概形 $X_\eta$ 应用 簇章引理
\ref{varieties-lemma-dim-1-projective-completion-after-insep}
。将由此得到的纯不可分域扩张记为 $k'/k$。
下一段把问题约化到 $k' = k$ 的情形。
（这一步对应于寻找引理陈述中的态射 $V \to U$；特别地，
若 $\kappa(\mathfrak p)$ 的特征为零，则可以取 $V = U$。）

\medskip\noindent
若 $k = \kappa(\mathfrak p)$ 的特征为零，则 $k' = k$。
若 $k = \kappa(\mathfrak p)$ 的特征为 $p > 0$，则 $p$
在 $A_\mathfrak p = \kappa(\mathfrak p)$ 中的像为零。
因此，用一个主局部化替换 $A$（即缩小 $S$）后，
可以假设在 $A$ 中有 $p = 0$。若 $k' \not = k$，则存在
$\beta \in k'$、$\beta \not \in k$，使得 $\beta^p \in k$。
用一个主局部化替换 $A$ 后，可以假设存在 $a \in A$，
使得 $\beta^p = a$。置 $A' = A[x]/(x^p - a)$。
于是 $S' = \Spec(A') \to \Spec(A) = S$ 有限、满射且普遍单射。
此外，若 $\mathfrak p' \subset A'$ 表示位于 $\mathfrak p$ 上方的唯一素理想，
则 $A'_{\mathfrak p'} = k(\beta)$，且 $k'/k(\beta)$ 的次数更小。
因此，用 $S'$ 替换 $S$，并用对应于 $\mathfrak p'$ 的点 $\eta'$
替换 $\eta$ 后，可见 $k'$ 在 $\eta$ 的剩余域上的次数已经减小。
如此继续并作归纳，就约化到
$k' = \kappa(\mathfrak p) = \kappa(\eta)$ 的情形。

\medskip\noindent
因此，可以假设 $S$ 仿射且既约，并且有图
$$
\xymatrix{
\overline{Y}_{1, \eta} \amalg \ldots \amalg \overline{Y}_{n, \eta} \ar[rd] &
Y_{1, \eta} \amalg \ldots \amalg Y_{n, \eta} \ar[r]_-\nu \ar[d] \ar[l]^j &
X_\eta \ar[d] \\
& \Spec(k_1) \amalg \ldots \amalg \Spec(k_n) \ar[r] &
\eta
}
$$
满足下列性质：
\begin{enumerate}
\item $\nu$ 是 $X_\eta$ 的正规化，
\item $j$ 是像稠密的开浸入，
\item 对 $i = 1, \ldots, n$，$k_i/\kappa(\eta)$ 是有限可分扩张，
\item $\overline{Y}_{i, \eta}$ 在 $k_i$ 上光滑、射影、几何不可约，
且维数 $\leq 1$。
\end{enumerate}
回忆 $\kappa(\eta) = \kappa(\mathfrak p) = A_\mathfrak p$
是所有 $A_a$（其中 $a \in A$、$a \not \in \mathfrak p$）的滤过余极限。
见 代数章引理 \ref{algebra-lemma-localization-colimit}。
因此，对某个 $a \in A$、$a \not \in \mathfrak p$，
可以把上图下降为 $\Spec(A_a)$ 上的相应图。
更确切地，用 $\Spec(A_a)$ 替换 $S$ 后，可以假设有交换图
$$
\xymatrix{
\overline{Y}_1 \amalg \ldots \amalg \overline{Y}_n \ar[rd] &
Y_1 \amalg \ldots \amalg Y_n \ar[r]_-\nu \ar[d] \ar[l]^j &
X \ar[d] \\
& T_1 \amalg \ldots \amalg T_n \ar[r] & S
}
$$
其到 $\eta$ 的基变换是上图，并满足下列性质：
\begin{enumerate}
\item $\nu$ 是有限满射，
\item $j$ 是开浸入，
\item 对 $i = 1, \ldots, n$，$T_i \to S$ 是有限 \'etale 态射，
\item $Y_i \to T_i$ 光滑且满射，
\item $\overline{Y}_i \to T_i$ 光滑且固有，并具有维数 $\leq 1$
的几何连通纤维。
\end{enumerate}
为此，先用
极限章引理 \ref{limits-lemma-descend-finite-presentation}
得到一个经基变换后为前图的图。然后使用 极限章引理
\ref{limits-lemma-descend-etale},
\ref{limits-lemma-descend-smooth},
\ref{limits-lemma-descend-finite-finite-presentation},
\ref{limits-lemma-limit-affine},
\ref{limits-lemma-descend-open-immersion},
\ref{limits-lemma-eventually-proper}, 及
\ref{limits-lemma-descend-surjective}
，得到 $\nu$ 有限且满射、$j$ 为开浸入、$T_i \to S$ 为有限
\'etale 态射、$Y_i \to T$ 光滑，以及 $\overline{Y}_i \to T_i$
固有且光滑。由于 $Y_i$ 不可能为空，光滑态射是开的，且
$T_i \to S$ 为有限 \'etale 态射，缩小 $S$ 后可以假设
$Y_i \to T_i$ 满射。最后，$\overline{Y}_i \to T_i$ 在 $T_i$
中位于 $\eta$ 上方的唯一点 $\eta_i = \Spec(k_i)$ 处的纤维几何连通。
因此，再次缩小后，可以假设所有纤维都具有同一性质，见
引理 \ref{more-morphisms-lemma-proper-flat-geom-red}.

\medskip\noindent
还需证明存在满足 (a)、(b)、(c) 的开集 $W \subset X$。
由于 $\nu_\eta : \coprod Y_{i, \eta} \to X_\eta$ 是正规化态射，
由
簇章引理 \ref{varieties-lemma-normalization-locally-algebraic}
可知，存在稠密开集 $W_\eta \subset X_\eta$，使得
$\nu^{-1}(W_\eta) \to W_\eta$ 等于 $W_\eta$ 的既约化到
$W_\eta$ 的包含。取拟紧开集 $W \subset X$，使其在 $\eta$
上的纤维就是刚找到的开集 $W_\eta$。
用另一个局部化替换 $A = \Gamma(S, \mathcal{O}_S)$ 后，
可以假设 $\nu^{-1}(W) \to W$ 是闭浸入，见 极限章引理
\ref{limits-lemma-descend-closed-immersion-finite-presentation}.
又因为 $\nu$ 满射，所以 $\nu^{-1}(W) \to W$ 是一个增厚。
置 $W_i = \nu^{-1}(W) \cap Y_i$。再次缩小 $S$，可以假设
对所有 $i$，$W_i \to T_i$ 满射（论证同上）。
由于 $Y_i \to T_i$ 具有几何不可约纤维，于是
$W_i \subset Y_i$ 在 $Y_i \to T_i$ 的每个纤维中稠密。
又因 $\nu$ 有限且满射，遂有 $W = \nu(\nu^{-1}(W))$
在 $X \to S$ 的每个纤维中也稠密。
\end{proof}











\section{分离且局部拟有限态射的下降}
\label{more-morphisms-section-separated-locally-quasi-finite}

\noindent
本节证明“分离且局部拟有限的态射对 fppf 覆盖满足下降”。
术语见 下降章定义
\ref{descent-definition-descending-types-morphisms}。
这个结果隐藏在 Raynaud 与 Gruson 那篇（在许多方面都）极精彩的论文
\cite[Lemma 5.7.1]{GruRay} 的证明中。
它也见于 \cite{Murre-representation} 以及
\cite[Expos\'e X, Lemma 5.4]{SGA3}
；后两者附加假设该态射局部有限表现。下面是正式陈述。

\begin{lemma}
\label{more-morphisms-lemma-separated-locally-quasi-finite-morphisms-fppf-descend}
设 $S$ 是概形。设 $\{X_i \to S\}_{i\in I}$ 是一个 fppf 覆盖，见
拓扑章定义 \ref{topologies-definition-fppf-covering}.
设 $(V_i/X_i, \varphi_{ij})$ 是相对于 $\{X_i \to S\}$ 的下降数据。
若每个态射 $V_i \to X_i$ 都分离且局部拟有限，则此下降数据有效。
\end{lemma}

\begin{proof}
分离性与局部拟有限性都是在任意基变换下保持的概形态射性质，见
概形章引理 \ref{schemes-lemma-separated-permanence} 及
态射章引理 \ref{morphisms-lemma-base-change-quasi-finite}.
因此可以应用 下降章引理 \ref{descent-lemma-descending-types-morphisms}，
只需证明 fppf 覆盖由仿射概形之间的单个平坦、满射、有限表现态射
$\{X \to S\}$ 给出的情形。写 $X = \Spec(A)$、$S = \Spec(R)$，
于是 $R \to A$ 是忠实平坦环映射。设 $(V, \varphi)$ 是相对于
$S$ 上 $X$ 的下降数据，并假设 $\pi : V \to X$ 分离且局部拟有限。

\medskip\noindent
任取仿射开集 $W^1 \subset V$。考虑
$W = \text{pr}_1(\varphi(W^1 \times_S X)) \subset V$。图示如下：
$$
\xymatrix{
W^1 \times_S X \ar[rrrrr] \ar[ddd] \ar[rd]
& & & & &
\varphi(W^1 \times_S X) \ar[ddd] \ar[ld] \\
& V \times_S X \ar[rrr]^\varphi \ar[rd] \ar[dd]
& & &
X \times_S V \ar[ld] \ar[dd] & \\
& &
X \times_S X \ar[r]^1 \ar[d]_{\text{pr}_0}
&
X \times_S X \ar[d]^{\text{pr}_1}
& & \\
W^1 \ar[r] &
V \ar[r] &
X &
X &
V \ar[l] &
W \ar[l]
}
$$
因为 $X \to S$ 平坦且有限表现，所以它普遍开
（态射章引理 \ref{morphisms-lemma-fppf-open}）。
因此 $W$ 是开集。显然，$W$ 还拟紧，并且 $W^1 \subset W$。
另外，由 $\varphi$ 的上闭链条件，
$\varphi(W \times_S X) = X \times_S W$。
因此，把 $\varphi$ 限制到 $W \times_S X$，得到新的下降数据
$(W, \varphi')$。注意态射 $W \to X$ 拟紧、分离且局部拟有限。
按定义，这蕴含它分离且拟有限。于是由
引理 \ref{more-morphisms-lemma-quasi-finite-separated-quasi-affine}
可知它是拟仿射的。故由
下降章引理 \ref{descent-lemma-quasi-affine}
可知下降数据 $(W, \varphi')$ 有效。

\medskip\noindent
换言之，存在由拟紧开集 $W_i$ 构成的开覆盖
$V = \bigcup W_i$，其中各开集在下降态射 $\varphi$ 下稳定。
此外，对每个这样的拟紧开集 $W \subset V$，相应的下降数据
$(W, \varphi')$ 都有效。把由各开集 $W_i$ 下降所得的概形粘合起来，
便知原下降数据有效，见
下降章引理 \ref{descent-lemma-effective-for-fpqc-is-local-upstairs}.
\end{proof}






\section{相对有限表现}
\label{more-morphisms-section-finite-type-finite-presentation}

\noindent
设 $R \to A$ 是有限型环映射，$M$ 是 $A$-模。在
代数进阶章
第 \ref{more-algebra-section-relative-finite-presentation} 节
中，我们定义了 $M$ 相对于 $R$ 有限表现的含义。
还证明了这一概念具有良好的局部化性质，并且可以粘合。
因此可以如下定义相应的整体概念。

\begin{definition}
\label{more-morphisms-definition-relatively-finitely-presented-sheaf}
设 $f : X \to S$ 是局部有限型的概形态射，$\mathcal{F}$ 是拟凝聚
$\mathcal{O}_X$-模。若存在仿射开覆盖 $S = \bigcup V_i$，并且对每个
$i$ 存在仿射开覆盖 $f^{-1}(V_i) = \bigcup_j U_{ij}$，使得
$\mathcal{F}(U_{ij})$ 作为 $\mathcal{O}_X(U_{ij})$-模相对于
$\mathcal{O}_S(V_i)$ 有限表现，则称 $\mathcal{F}$
{\it 相对于 $S$ 有限表现}，或称它{\it 在 $S$ 上相对有限表现}。
\end{definition}

\noindent
注意，这蕴含 $\mathcal{F}$ 是有限型 $\mathcal{O}_X$-模。
若 $X \to S$ 仅仅局部有限型，则即使 $X \to S$ 不是局部有限表现，
$\mathcal{F}$ 仍可能相对于 $S$ 有限表现。我们将看到，
$X \to S$ 局部有限表现，当且仅当 $\mathcal{O}_X$ 相对于 $S$ 有限表现。

\begin{lemma}
\label{more-morphisms-lemma-relative-finite-presentation-characterize}
设 $f : X \to S$ 是局部有限型的概形态射，$\mathcal{F}$ 是拟凝聚
$\mathcal{O}_X$-模。下列条件等价：
\begin{enumerate}
\item $\mathcal{F}$ 相对于 $S$ 有限表现；
\item 对每一对满足 $f(U) \subset V$ 的仿射开集 $U \subset X$、
$V \subset S$，$\mathcal{O}_X(U)$-模 $\mathcal{F}(U)$
相对于 $\mathcal{O}_S(V)$ 有限表现。
\end{enumerate}
此外，若这些条件成立，则对每一对满足 $f(U) \subset V$ 的开子概形
$U \subset X$、$V \subset S$，限制 $\mathcal{F}|_U$ 相对于 $V$ 有限表现。
\end{lemma}

\begin{proof}
最后一个陈述由 (1) 与 (2) 的等价性显然成立。(2) 蕴含 (1) 也显然。
假设 (1) 成立。取
$S = \bigcup V_i$ 和 $f^{-1}(V_i) = \bigcup U_{ij}$ 为
定义 \ref{more-morphisms-definition-relatively-finitely-presented-sheaf}.
中的仿射开覆盖。设 $U \subset X$ 和 $V \subset S$ 如 (2) 所述。
由 代数进阶章引理
\ref{more-algebra-lemma-glue-relative-finite-presentation}
，只需找到 $U$ 的标准开覆盖 $U = \bigcup U_k$，使得
$\mathcal{F}(U_k)$ 相对于 $\mathcal{O}_S(V)$ 有限表现。
换言之，对每个 $u \in U$，只需找到标准仿射开集
$u \in U' \subset U$，使得 $\mathcal{F}(U')$ 相对于
$\mathcal{O}_S(V)$ 有限表现。选取 $i$ 使 $f(u) \in V_i$，
再选取 $j$ 使 $u \in U_{ij}$。由
概形章引理 \ref{schemes-lemma-standard-open-two-affines}
，可以找到 $v \in V' \subset V \cap V_i$，使它在 $V'$ 与
$V_i$ 中都是标准仿射开集。于是 $f^{-1}V'  \cap U$ 与
$f^{-1}V' \cap U_{ij}$ 分别是 $U$ 与 $U_{ij}$ 的标准仿射开集。
再次应用该引理，可以找到
$u \in U' \subset f^{-1}V' \cap U \cap U_{ij}$，使它在
$f^{-1}V'  \cap U$ 与 $f^{-1}V' \cap U_{ij}$ 中都是标准仿射开集。
因此 $U'$ 也是 $U$ 与 $U_{ij}$ 的标准仿射开集。由 代数进阶章引理
\ref{more-algebra-lemma-localize-relative-finite-presentation}
，$\mathcal{F}(U_{ij})$ 相对于 $\mathcal{O}_S(V_i)$ 有限表现这一假设蕴含
$\mathcal{F}(U')$ 相对于 $\mathcal{O}_S(V_i)$ 有限表现。由于
$\mathcal{O}_X(U') =
\mathcal{O}_X(U') \otimes_{\mathcal{O}_S(V_i)} \mathcal{O}_S(V')$
，由 代数进阶章引理
\ref{more-algebra-lemma-base-change-relative-finite-presentation}
可知 $\mathcal{F}(U')$ 相对于 $\mathcal{O}_S(V')$ 有限表现。
再次应用 代数进阶章引理
\ref{more-algebra-lemma-localize-relative-finite-presentation}
，得到 $\mathcal{F}(U')$ 相对于 $\mathcal{O}_S(V)$ 有限表现。
证明完毕。
\end{proof}

\begin{lemma}
\label{more-morphisms-lemma-relative-finite-presentation}
设 $f : X \to S$ 是局部有限型的概形态射，$\mathcal{F}$ 是拟凝聚
$\mathcal{O}_X$-模。
\begin{enumerate}
\item 若 $f$ 局部有限表现，则 $\mathcal{F}$ 相对于 $S$ 有限表现，
当且仅当 $\mathcal{F}$ 有限表现。
\item 态射 $f$ 局部有限表现，当且仅当 $\mathcal{O}_X$
相对于 $S$ 有限表现。
\end{enumerate}
\end{lemma}

\begin{proof}
这直接由定义得到，见
代数进阶章定义
\ref{more-algebra-definition-relatively-finitely-presented} 后的讨论。
\end{proof}

\begin{lemma}
\label{more-morphisms-lemma-finite-morphism-relative-finite-presentation}
设 $\pi : X \to Y$ 是在基概形 $S$ 上局部有限型的概形之间的有限态射，
$\mathcal{F}$ 是拟凝聚 $\mathcal{O}_X$-模。则 $\mathcal{F}$
相对于 $S$ 有限表现，当且仅当 $\pi_*\mathcal{F}$ 相对于 $S$ 有限表现。
\end{lemma}

\begin{proof}
这是把
代数进阶章引理 \ref{more-algebra-lemma-finite-extension}
的结果翻译成概形语言。
\end{proof}

\begin{lemma}
\label{more-morphisms-lemma-base-change-relative-finite-presentation}
设 $f : X \to S$ 是局部有限型的概形态射，$\mathcal{F}$ 是拟凝聚
$\mathcal{O}_X$-模。设 $S' \to S$ 是概形态射，置 $X' = X \times_S S'$，
并以 $\mathcal{F}'$ 表示 $\mathcal{F}$ 到 $X'$ 的拉回。
若 $\mathcal{F}$ 相对于 $S$ 有限表现，则 $\mathcal{F}'$
相对于 $S'$ 有限表现。
\end{lemma}

\begin{proof}
这是把
代数进阶章引理
\ref{more-algebra-lemma-base-change-relative-finite-presentation}
的结果翻译成概形语言。
\end{proof}

\begin{lemma}
\label{more-morphisms-lemma-pull-relative-finite-presentation}
设 $X \to Y \to S$ 是局部有限型的概形态射，$\mathcal{G}$ 是拟凝聚
$\mathcal{O}_Y$-模。若 $f : X \to Y$ 局部有限表现，且
$\mathcal{G}$ 相对于 $S$ 有限表现，则 $f^*\mathcal{G}$
相对于 $S$ 有限表现。
\end{lemma}

\begin{proof}
这是把
代数进阶章引理
\ref{more-algebra-lemma-pull-relative-finite-presentation}
的结果翻译成概形语言。
\end{proof}

\begin{lemma}
\label{more-morphisms-lemma-composition-relative-finite-presentation}
设 $X \to Y \to S$ 是局部有限型的概形态射，$\mathcal{F}$ 是拟凝聚
$\mathcal{O}_X$-模。若 $Y \to S$ 局部有限表现，且 $\mathcal{F}$
相对于 $Y$ 有限表现，则 $\mathcal{F}$ 相对于 $S$ 有限表现。
\end{lemma}

\begin{proof}
这是把
代数进阶章引理
\ref{more-algebra-lemma-composition-relative-finite-presentation}
的结果翻译成概形语言。
\end{proof}

\begin{lemma}
\label{more-morphisms-lemma-ses-relatively-finite-presentation}
设 $X \to S$ 是局部有限型的概形态射。设
$0 \to \mathcal{F}' \to \mathcal{F} \to \mathcal{F}'' \to 0$
是拟凝聚 $\mathcal{O}_X$-模的短正合列。
\begin{enumerate}
\item 若 $\mathcal{F}', \mathcal{F}''$ 相对于 $S$ 有限表现，
则 $\mathcal{F}$ 也如此。
\item 若 $\mathcal{F}'$ 是有限型 $\mathcal{O}_X$-模，且
$\mathcal{F}$ 相对于 $S$ 有限表现，则 $\mathcal{F}''$
相对于 $S$ 有限表现。
\end{enumerate}
\end{lemma}

\begin{proof}
这是把
代数进阶章引理
\ref{more-algebra-lemma-ses-relatively-finite-presentation}
的结果翻译成概形语言。
\end{proof}

\begin{lemma}
\label{more-morphisms-lemma-sum-relatively-finite-presentation}
\begin{slogan}
模的直和项继承相对于基有限表现这一性质。
\end{slogan}
设 $X \to S$ 是局部有限型的概形态射，$\mathcal{F}, \mathcal{F}'$
是拟凝聚 $\mathcal{O}_X$-模。若 $\mathcal{F} \oplus \mathcal{F}'$
相对于 $S$ 有限表现，则 $\mathcal{F}$ 和 $\mathcal{F}'$ 也如此。
\end{lemma}

\begin{proof}
这是把
代数进阶章引理
\ref{more-algebra-lemma-sum-relatively-finite-presentation}
的结果翻译成概形语言。
\end{proof}










\section{相对伪凝聚性}
\label{more-morphisms-section-relative-pseudo-coherence}

\noindent
本节是
代数进阶章第 \ref{more-algebra-section-relative-pseudo-coherent} 节
在概形情形下的对应版本。我们强烈建议读者先阅读该节。
尽管本节的内容以及
上同调章第 \ref{cohomology-section-strictly-perfect} 节,
\ref{cohomology-section-pseudo-coherent},
\ref{cohomology-section-tor}, 及
\ref{cohomology-section-perfect}
中关于伪凝聚复形的内容，都是针对任意 $\mathcal{O}_X$-模复形建立的，
但若 $X$ 是概形，则只处理 $D_\QCoh(\mathcal{O}_X)$ 中的对象，
会极大简化许多引理和论证，常常能立即把手头的问题归结为其代数对应物。
此外，我们首先证明的结果之一是：相对伪凝聚性蕴含上同调层拟凝聚，
见 引理 \ref{more-morphisms-lemma-relative-pseudo-coherence}。
因此，初次阅读时，建议读者只处理 $D_\QCoh(\mathcal{O}_X)$ 中的对象。

\begin{lemma}
\label{more-morphisms-lemma-relatively-pseudo-coherent}
设 $X \to S$ 是仿射概形之间的有限型态射，$E$ 是
$D(\mathcal{O}_X)$ 的对象，并设 $m \in \mathbf{Z}$。
下列条件等价：
\begin{enumerate}
\item 对某个闭浸入 $i : X \to \mathbf{A}^n_S$，
$D(\mathcal{O}_{\mathbf{A}^n_S})$ 的对象 $Ri_*E$ 是 $m$-伪凝聚的；
\item 对所有闭浸入 $i : X \to \mathbf{A}^n_S$，
$D(\mathcal{O}_{\mathbf{A}^n_S})$ 的对象 $Ri_*E$ 都是 $m$-伪凝聚的。
\end{enumerate}
\end{lemma}

\begin{proof}
写 $S = \Spec(R)$、$X = \Spec(A)$。设 $i$ 对应于满射
$\alpha : R[x_1, \ldots, x_n] \to A$，而 $X \to \mathbf{A}^m_S$
对应于 $\beta : R[y_1, \ldots, y_m] \to A$。
选取满足 $\alpha(f_j) = \beta(y_j)$ 的
$f_j \in R[x_1, \ldots, x_n]$，以及满足
$\beta(g_i) = \alpha(x_i)$ 的 $g_i \in R[y_1, \ldots, y_m]$。
于是得到交换图
$$
\xymatrix{
R[x_1, \ldots, x_n, y_1, \ldots, y_m]
\ar[d]^{x_i \mapsto g_i} \ar[rr]_-{y_j \mapsto f_j} & &
R[x_1, \ldots, x_n] \ar[d] \\
R[y_1, \ldots, y_m] \ar[rr] & & A
}
$$
它对应于闭浸入的交换图
$$
\xymatrix{
\mathbf{A}^{n + m}_S & \mathbf{A}^n_S \ar[l] \\
\mathbf{A}^m_S \ar[u] & X \ar[u] \ar[l]
}
$$
因此，只需证明：对闭浸入
$$
f : \mathbf{A}^m_S \to \mathbf{A}^{n + m}_S
$$
，$D(\mathcal{O}_{\mathbf{A}^m_S})$ 的对象 $E$ 是 $m$-伪凝聚的，
当且仅当 $Rf_*E$ 是 $m$-伪凝聚的。这由
概形的导出范畴章引理
\ref{perfect-lemma-closed-push-pseudo-coherent}
以及 $f_*\mathcal{O}_{\mathbf{A}^m_S}$ 是伪凝聚
$\mathcal{O}_{\mathbf{A}^{n + m}_S}$-模这一事实得到。
$f_*\mathcal{O}_{\mathbf{A}^m_S}$ 的伪凝聚性很容易直接证明，
但它也可由
概形的导出范畴章引理 \ref{perfect-lemma-pseudo-coherent-affine}
和
代数进阶章引理 \ref{more-algebra-lemma-relatively-pseudo-coherent}.
得到。
\end{proof}

\noindent
回忆，若 $f : X \to S$ 是局部有限型的概形态射，则对每一对满足
$f(U) \subset V$ 的仿射开集 $U \subset X$、$V \subset S$，
环映射 $\mathcal{O}_S(V) \to \mathcal{O}_X(U)$ 都是有限型的
（态射章引理 \ref{morphisms-lemma-locally-finite-type-characterize}）。
因此总存在闭浸入 $U \to \mathbf{A}^n_V$，下面的定义是有意义的。

\begin{definition}
\label{more-morphisms-definition-relative-pseudo-coherence}
设 $f : X \to S$ 是局部有限型的概形态射。设 $E$ 是
$D(\mathcal{O}_X)$ 的对象，$\mathcal{F}$ 是 $\mathcal{O}_X$-模。
固定 $m \in \mathbf{Z}$。
\begin{enumerate}
\item 若存在仿射开覆盖 $S = \bigcup V_i$，且对每个 $i$
存在仿射开覆盖 $f^{-1}(V_i) = \bigcup U_{ij}$，使得对每一对
$(U_{ij} \to V_i, E|_{U_{ij}})$，
引理 \ref{more-morphisms-lemma-relatively-pseudo-coherent}
中的等价条件都成立，则称 $E${\it 相对于 $S$ 是 $m$-伪凝聚的}。
\item 若对所有 $m \in \mathbf{Z}$，$E$ 相对于 $S$ 都是
$m$-伪凝聚的，则称 $E${\it 相对于 $S$ 是伪凝聚的}。
\item 若把 $\mathcal{F}$ 看作 $D(\mathcal{O}_X)$ 的对象时，
它相对于 $S$ 是 $m$-伪凝聚的，则称 $\mathcal{F}$
{\it 相对于 $S$ 是 $m$-伪凝聚的}。
\item 若把 $\mathcal{F}$ 看作 $D(\mathcal{O}_X)$ 的对象时，
它相对于 $S$ 是伪凝聚的，则称 $\mathcal{F}$
{\it 相对于 $S$ 是伪凝聚的}。
\end{enumerate}
\end{definition}

\noindent
若 $X$ 拟紧，并且对某个 $m$，$E$ 相对于 $S$ 是 $m$-伪凝聚的，
则 $E$ 上有界。若 $E$ 相对于 $S$ 是伪凝聚的，
则 $E$ 具有拟凝聚上同调层。

\begin{lemma}
\label{more-morphisms-lemma-relative-pseudo-coherence}
设 $f : X \to S$ 是局部有限型的概形态射。若
$D(\mathcal{O}_X)$ 中的 $E$ 相对于 $S$ 是 $m$-伪凝聚的，
则对 $i > m$，$H^i(E)$ 是拟凝聚 $\mathcal{O}_X$-模。
若 $E$ 相对于 $S$ 是伪凝聚的，则 $E$ 是
$D_\QCoh(\mathcal{O}_X)$ 的对象。
\end{lemma}

\begin{proof}
选取仿射开覆盖 $S = \bigcup V_i$，并对每个 $i$ 选取仿射开覆盖
$f^{-1}(V_i) = \bigcup U_{ij}$，使得对每一对
$(U_{ij} \to V_i, E|_{U_{ij}})$，
引理 \ref{more-morphisms-lemma-relatively-pseudo-coherent}
中的等价条件都成立。由于拟凝聚性在 $X$ 上是局部的，
可以假设存在闭浸入 $i : X \to \mathbf{A}^n_S$，使得
$Ri_*E$ 在 $\mathbf{A}^n_S$ 上是 $m$-伪凝聚的。
由 概形的导出范畴章引理 \ref{perfect-lemma-pseudo-coherent}，
这意味着当 $q > m$ 时，$H^q(Ri_*E)$ 拟凝聚。
因为 $i_*$ 是正合函子，故
$i_*H^q(E) = H^q(Ri_*E)$ 在 $\mathbf{A}^n_S$ 上拟凝聚。
由 态射章引理 \ref{morphisms-lemma-i-star-equivalence}，
这蕴含 $H^q(E)$ 如所需是拟凝聚的（严格说来，它蕴含存在某个拟凝聚
$\mathcal{O}_X$-模 $\mathcal{F}$，使得
$i_*\mathcal{F} = i_*H^q(E)$；于是
模章引理 \ref{modules-lemma-i-star-equivalence}
告诉我们 $\mathcal{F} \cong H^q(E)$，从而得到结论）。
\end{proof}

\noindent
接下来证明，相对伪凝聚条件具有良好的局部化性质。

\begin{lemma}
\label{more-morphisms-lemma-localize-relative-pseudo-coherent}
设 $S$ 是仿射概形，$V \subset S$ 是标准开集，
$X \to V$ 是仿射概形之间的有限型态射，$U \subset X$ 是仿射开集，
$E$ 是 $D(\mathcal{O}_X)$ 的对象。若对 $(X \to V, E)$，
引理 \ref{more-morphisms-lemma-relatively-pseudo-coherent}
中的等价条件成立，则对 $(U \to S, E|_U)$，
引理 \ref{more-morphisms-lemma-relatively-pseudo-coherent}
中的等价条件也成立。
\end{lemma}

\begin{proof}
写 $S = \Spec(R)$、$V = D(f)$、$X = \Spec(A)$、$U = D(g)$。
假设对 $(X \to V, E)$，
引理 \ref{more-morphisms-lemma-relatively-pseudo-coherent}
中的等价条件成立。

\medskip\noindent
选取满射 $R_f[x_1, \ldots, x_n] \to A$。写
$R_f = R[x_0]/(fx_0 - 1)$。于是 $R[x_0, x_1, \ldots, x_n] \to A$
是满射，且 $R_f[x_1, \ldots, x_n]$ 作为
$R[x_0, \ldots, x_n]$-模是伪凝聚的。因此有
$$
X \to \mathbf{A}^n_V \to \mathbf{A}^{n + 1}_S
$$
，并可应用
概形的导出范畴章
引理 \ref{perfect-lemma-closed-push-pseudo-coherent}
断定 $E$ 到 $\mathbf{A}^{n + 1}_S$ 的推前 $E'$ 是 $m$-伪凝聚的。

\medskip\noindent
选取映到 $g \in A$ 的元素 $g' \in R[x_0, x_1, \ldots, x_n]$。
考虑满射
$R[x_0, \ldots, x_{n + 1}] \to R[x_0, \ldots, x_n, 1/g']$.
得到
$$
\xymatrix{
X \ar[d] & U \ar[d] \ar[l] \ar[dr] \\
\mathbf{A}^{n + 1}_S & D(g')\ar[l] \ar[r] & \mathbf{A}^{n + 2}_S
}
$$
其中左下箭头是开浸入，右下箭头是闭浸入。与前面一样可断定，
$E'|_{D(g')}$ 到 $\mathbf{A}^{n + 2}_S$ 的推前是 $m$-伪凝聚的。
由于这也正是 $E|_U$ 到 $\mathbf{A}^{n + 2}_S$ 的推前，引理得证。
\end{proof}

\begin{lemma}
\label{more-morphisms-lemma-glue-relative-pseudo-coherent}
设 $X \to S$ 是仿射概形之间的有限型态射，$E$ 是
$D(\mathcal{O}_X)$ 的对象，$m \in \mathbf{Z}$。
设 $X = \bigcup U_i$ 是标准仿射开覆盖。下列条件等价：
\begin{enumerate}
\item 对所有 $(U_i \to S, E|_{U_i})$，
引理 \ref{more-morphisms-lemma-relatively-pseudo-coherent}
中的等价条件成立；
\item 对 $(X \to S, E)$，
引理 \ref{more-morphisms-lemma-relatively-pseudo-coherent}
中的等价条件成立。
\end{enumerate}
\end{lemma}

\begin{proof}
蕴含关系 (2) $\Rightarrow$ (1) 就是
引理 \ref{more-morphisms-lemma-localize-relative-pseudo-coherent}.
假设 (1)。写 $S = \Spec(R)$、$X = \Spec(A)$、$U_i = D(f_i)$，
并在 $A$ 中写 $1 = \sum f_ig_i$。考虑满射
$$
R[x_i, y_i, z_i] \to R[x_i, y_i, z_i]/(\sum y_iz_i - 1) \to A.
$$
它把 $y_i$ 映到 $f_i$，把 $z_i$ 映到 $g_i$。注意
$R[x_i, y_i, z_i]/(\sum y_iz_i - 1)$ 作为
$R[x_i, y_i, z_i]$-模是伪凝聚的。因此，只需证明 $E$ 到
$T = \Spec(R[x_i, y_i, z_i]/(\sum y_iz_i - 1))$ 的推前是
$m$-伪凝聚的，见
概形的导出范畴章
引理 \ref{perfect-lemma-closed-push-pseudo-coherent}.
对每个 $i_0$，只需证明此推前到
$W_{i_0} = \Spec(R[x_i, y_i, z_i, 1/y_{i_0}]/(\sum y_iz_i - 1))$
的限制是 $m$-伪凝聚的。注意有交换图
$$
\xymatrix{
X \ar[d] & U_{i_0} \ar[l] \ar[d] \\
T & W_{i_0} \ar[l]
}
$$
它蕴含 $E$ 到 $T$ 的推前在 $W_{i_0}$ 上的限制，正是
$E|_{U_{i_0}}$ 到 $W_{i_0}$ 的推前。由于\linebreak[4]
$R[x_i, y_i, z_i, 1/y_{i_0}]/(\sum y_iz_i - 1)$
同构于 $R$ 上的多项式环，这就证明了所需结论。
\end{proof}

\begin{lemma}
\label{more-morphisms-lemma-relative-pseudo-coherence-characterize}
设 $f : X \to S$ 是局部有限型的概形态射，$E$ 是
$D(\mathcal{O}_X)$ 的对象。固定 $m \in \mathbf{Z}$。
下列条件等价：
\begin{enumerate}
\item $E$ 相对于 $S$ 是 $m$-伪凝聚的；
\item 对每一对满足 $f(U) \subset V$ 的仿射开集 $U \subset X$、
$V \subset S$，对 $(U \to V, E|_U)$，
引理 \ref{more-morphisms-lemma-relatively-pseudo-coherent}
中的等价条件成立。
\end{enumerate}
此外，若这些条件成立，则对每一对满足 $f(U) \subset V$
的开子概形 $U \subset X$、$V \subset S$，限制 $E|_U$
相对于 $V$ 是 $m$-伪凝聚的。
\end{lemma}

\begin{proof}
最后一个陈述由 (1) 与 (2) 的等价性显然成立。(2) 蕴含 (1) 也显然。
假设 (1) 成立。取 $S = \bigcup V_i$ 和
$f^{-1}(V_i) = \bigcup U_{ij}$ 为
定义 \ref{more-morphisms-definition-relative-pseudo-coherence}.
中的仿射开覆盖。设 $U \subset X$ 和 $V \subset S$ 如 (2) 所述。
由 引理 \ref{more-morphisms-lemma-glue-relative-pseudo-coherent}，
只需找到 $U$ 的标准开覆盖 $U = \bigcup U_k$，使得对每一对
$(U_k \to V, E|_{U_k})$，
引理 \ref{more-morphisms-lemma-relatively-pseudo-coherent}
中的等价条件成立。换言之，对每个 $u \in U$，只需找到标准仿射开集
$u \in U' \subset U$，使得对 $(U' \to V, E|_{U'})$，
引理 \ref{more-morphisms-lemma-relatively-pseudo-coherent}
中的等价条件成立。选取 $i$ 使 $f(u) \in V_i$，再选取 $j$
使 $u \in U_{ij}$。由
概形章引理 \ref{schemes-lemma-standard-open-two-affines}
，可以找到 $v \in V' \subset V \cap V_i$，使它在 $V'$ 与 $V_i$
中都是标准仿射开集。于是 $f^{-1}V'  \cap U$ 与
$f^{-1}V' \cap U_{ij}$ 分别是 $U$ 与 $U_{ij}$ 的标准仿射开集。
再次应用该引理，可以找到
$u \in U' \subset f^{-1}V' \cap U \cap U_{ij}$，使它在
$f^{-1}V'  \cap U$ 与 $f^{-1}V' \cap U_{ij}$ 中都是标准仿射开集。
因此 $U'$ 也是 $U$ 与 $U_{ij}$ 的标准仿射开集。
由 引理 \ref{more-morphisms-lemma-localize-relative-pseudo-coherent}，
对 $(U_{ij} \to V_i, E|_{U_{ij}})$，
引理 \ref{more-morphisms-lemma-relatively-pseudo-coherent}
中的等价条件成立这一假设，蕴含对 $(U' \to V, E|_{U'})$，
引理 \ref{more-morphisms-lemma-relatively-pseudo-coherent}
中的等价条件成立。
\end{proof}

\noindent
对于上同调层拟凝聚的导出范畴对象，可以把相对 $m$-伪凝聚性
与 代数进阶章定义
\ref{more-algebra-definition-relatively-pseudo-coherent}.
中定义的概念联系起来。我们将使用如下事实：对仿射概形
$U = \Spec(A)$，函子 $R\Gamma(U, -)$ 诱导
$D_\QCoh(\mathcal{O}_U)$ 与 $D(A)$ 之间的等价，见
概形的导出范畴章引理
\ref{perfect-lemma-affine-compare-bounded}.
此函子与拉回相容：若 $E$ 是 $D_\QCoh(\mathcal{O}_U)$ 的对象，
而 $A \to B$ 是对应于仿射概形态射
$g : V = \Spec(B) \to \Spec(A) = U$ 的环映射，则
$R\Gamma(V, Lg^*E) = R\Gamma(U, E) \otimes_A^\mathbf{L} B$.
见 概形的导出范畴章引理
\ref{perfect-lemma-quasi-coherence-pullback}.

\begin{lemma}
\label{more-morphisms-lemma-qcoh-relative-pseudo-coherence-characterize}
设 $f : X \to S$ 是局部有限型的概形态射，$E$ 是
$D_\QCoh(\mathcal{O}_X)$ 的对象。固定 $m \in \mathbf{Z}$。
下列条件等价：
\begin{enumerate}
\item $E$ 相对于 $S$ 是 $m$-伪凝聚的；
\item 存在仿射开覆盖 $S = \bigcup V_i$，且对每个 $i$
存在仿射开覆盖 $f^{-1}(V_i) = \bigcup U_{ij}$，使
$\mathcal{O}_X(U_{ij})$-模复形 $R\Gamma(U_{ij}, E)$
相对于 $\mathcal{O}_S(V_i)$ 是 $m$-伪凝聚的；
\item 对每一对满足 $f(U) \subset V$ 的仿射开集 $U \subset X$、
$V \subset S$，$\mathcal{O}_X(U)$-模复形 $R\Gamma(U, E)$
相对于 $\mathcal{O}_S(V)$ 是 $m$-伪凝聚的。
\end{enumerate}
\end{lemma}

\begin{proof}
设 $U$ 和 $V$ 如 (2) 所述，并选取闭浸入
$i : U \to \mathbf{A}^n_V$。利用
引理 \ref{more-morphisms-lemma-relative-pseudo-coherence-characterize} 作形式论证，
可见只需证明：$Ri_*(E|_U)$ 是 $m$-伪凝聚的，当且仅当
$R\Gamma(U, E)$ 相对于 $\mathcal{O}_S(V)$ 是 $m$-伪凝聚的。
写 $U = \Spec(A)$、$V = \Spec(R)$，并且
% P05-EN-0127: frozen source omits the closing parenthesis in this Spec expression.
$\mathbf{A}^n_V = \Spec(R[x_1, \ldots, x_n]$. 由上述说明
由引理前的说明，$E|_U$ 拟同构于 $U$ 上由 $D(A)$ 的对象
$R\Gamma(U, E)$ 所伴随的拟凝聚层复形。注意，因为 $i$ 是闭浸入
（因而 $i_*$ 正合），有
$R\Gamma(U, E) = R\Gamma(\mathbf{A}^n_V, Ri_*(E|_U))$。
于是 $Ri_*E$ 由视为 $D(R[x_1, \ldots, x_n])$ 对象的
$R\Gamma(U, E)$ 所伴随。由
概形的导出范畴章引理 \ref{perfect-lemma-pseudo-coherent-affine}
，$Ri_*(E|_U)$ 的 $m$-伪凝聚性等价于 $R\Gamma(U, E)$
在 $D(R[x_1, \ldots, x_n])$ 中的 $m$-伪凝聚性；按定义，
后者又等价于 $R\Gamma(U, E)$ 相对于 $R = \mathcal{O}_S(V)$
是 $m$-伪凝聚的。结论得证。
\end{proof}

\begin{lemma}
\label{more-morphisms-lemma-closed-morphism-relative-pseudo-coherence}
设 $i : X \to Y$ 是基概形 $S$ 上局部有限型的概形之间的态射。
假设 $i$ 诱导 $X$ 到 $Y$ 的一个闭子集的同胚。
设 $E$ 是 $D(\mathcal{O}_X)$ 的对象。则 $E$ 相对于 $S$
是 $m$-伪凝聚的，当且仅当 $Ri_*E$ 相对于 $S$ 是 $m$-伪凝聚的。
\end{lemma}

\begin{proof}
由 态射章引理 \ref{morphisms-lemma-homeomorphism-affine}，
态射 $i$ 是仿射的。因此可以假设 $S$、$Y$、$X$ 都仿射。
写 $S = \Spec(R)$、$Y = \Spec(A)$、$X = \Spec(B)$。
该条件意味着 $A/\text{rad}(A) \to B/\text{rad}(B)$ 满射；
% P05-EN-0128: frozen source calls rad(A), rad(B) Jacobson radicals although the nilpotence argument uses nilradicals.
这里 $\text{rad}(A)$ 和 $\text{rad}(B)$ 分别表示 $A$ 和 $B$
的 Jacobson 根。由于 $B$ 在 $A$ 上为有限型，可以找到
$b_1, \ldots, b_m \in \text{rad}(B)$，它们生成 $B$ 作为
$A$-代数。设对所有 $j$ 都有 $b_j^N = 0$。考虑环的图
$$
\xymatrix{
B & R[x_i, y_j]/(y_j^N) \ar[l] & R[x_i, y_j] \ar[l] \\
A \ar[u] & R[x_i] \ar[l] \ar[u] \ar[ru]
}
$$
它转化为仿射概形的图
$$
\xymatrix{
X \ar[d] \ar[r] & T \ar[d] \ar[r] & \mathbf{A}^{n + m}_S \ar[ld] \\
Y \ar[r] & \mathbf{A}^n_S
}
$$
。由 引理 \ref{more-morphisms-lemma-relative-pseudo-coherence-characterize}，
$E$ 相对于 $S$ 是 $m$-伪凝聚的，当且仅当它到
$\mathbf{A}^{n + m}_S$ 的推前是 $m$-伪凝聚的。由
概形的导出范畴章引理
\ref{perfect-lemma-closed-push-pseudo-coherent}
，这又等价于它到 $T$ 的推前是 $m$-伪凝聚的。同一引理表明，
后者等价于它到 $\mathbf{A}^n_S$ 的推前是 $m$-伪凝聚的。
再次由
引理 \ref{more-morphisms-lemma-relative-pseudo-coherence-characterize}
，这等价于 $Ri_*E$ 相对于 $S$ 是 $m$-伪凝聚的。
\end{proof}

\begin{lemma}
\label{more-morphisms-lemma-finite-morphism-relative-pseudo-coherence}
设 $\pi : X \to Y$ 是基概形 $S$ 上局部有限型的概形之间的有限态射。
设 $E$ 是 $D_\QCoh(\mathcal{O}_X)$ 的对象。则 $E$ 相对于 $S$
是 $m$-伪凝聚的，当且仅当 $R\pi_*E$ 相对于 $S$ 是 $m$-伪凝聚的。
\end{lemma}

\begin{proof}
这是把
代数进阶章引理
\ref{more-algebra-lemma-finite-extension-pseudo-coherent}
的结果翻译成概形语言。注意，由 概形的导出范畴章引理
\ref{perfect-lemma-quasi-coherence-direct-image}.
，$R\pi_*$ 确实把 $D_\QCoh(\mathcal{O}_X)$ 映入
$D_\QCoh(\mathcal{O}_Y)$。在这一翻译中使用
引理 \ref{more-morphisms-lemma-relative-pseudo-coherence-characterize}.
\end{proof}

\begin{lemma}
\label{more-morphisms-lemma-cone-relatively-pseudo-coherent}
设 $f : X \to S$ 是局部有限型的概形态射。
设 $(E, E', E'')$ 是 $D(\mathcal{O}_X)$ 的可区别三角，
并设 $m \in \mathbf{Z}$。
\begin{enumerate}
\item 若 $E$ 相对于 $S$ 是 $(m + 1)$-伪凝聚的，且 $E'$
相对于 $S$ 是 $m$-伪凝聚的，则 $E''$ 相对于 $S$ 是 $m$-伪凝聚的。
\item 若 $E, E''$ 相对于 $S$ 是 $m$-伪凝聚的，则 $E'$
相对于 $S$ 是 $m$-伪凝聚的。
\item 若 $E'$ 相对于 $S$ 是 $(m + 1)$-伪凝聚的，且 $E''$
相对于 $S$ 是 $m$-伪凝聚的，则 $E$ 相对于 $S$
是 $(m + 1)$-伪凝聚的。
\end{enumerate}
此外，若 $E, E', E''$ 中有两个相对于 $S$ 是伪凝聚的，
则第三个也如此。
\end{lemma}

\begin{proof}
这直接由 引理 \ref{more-morphisms-lemma-relative-pseudo-coherence-characterize} 和
上同调章引理 \ref{cohomology-lemma-cone-pseudo-coherent}.
得到。
\end{proof}

\begin{lemma}
\label{more-morphisms-lemma-rel-n-pseudo-module}
设 $X \to S$ 是局部有限型的概形态射，$\mathcal{F}$ 是
$\mathcal{O}_X$-模。则
\begin{enumerate}
\item 对所有 $m > 0$，$\mathcal{F}$ 相对于 $S$ 都是 $m$-伪凝聚的；
\item $\mathcal{F}$ 相对于 $S$ 是 $0$-伪凝聚的，当且仅当
$\mathcal{F}$ 是有限型 $\mathcal{O}_X$-模；
\item $\mathcal{F}$ 相对于 $S$ 是 $(-1)$-伪凝聚的，当且仅当
$\mathcal{F}$ 拟凝聚且相对于 $S$ 有限表现。
\end{enumerate}
\end{lemma}

\begin{proof}
第 (1) 部分直接由定义得到。为证明第 (3) 部分，可以在 $X$ 上局部工作
（两个性质都是局部的）。因此可以假设 $X$ 和 $S$ 仿射。
选取闭浸入 $i : X \to \mathbf{A}^n_S$。于是 $\mathcal{F}$
相对于 $S$ 是 $(-1)$-伪凝聚的，当且仅当 $i_*\mathcal{F}$
是 $(-1)$-伪凝聚的；而这又当且仅当 $i_*\mathcal{F}$
是有限表现的 $\mathcal{O}_{\mathbf{A}^n_S}$-模，见
上同调章引理 \ref{cohomology-lemma-finite-cohomology}.
有限表现的模拟凝聚，见
模章引理 \ref{modules-lemma-finite-presentation-quasi-coherent}.
由 态射章引理 \ref{morphisms-lemma-i-star-equivalence}，
$\mathcal{F}$ 拟凝聚，当且仅当 $i_*\mathcal{F}$ 拟凝聚。
至此第 (3) 部分得证。第 (2) 部分的证明类似但更简单。
\end{proof}

\begin{lemma}
\label{more-morphisms-lemma-summands-relative-pseudo-coherent}
设 $X \to S$ 是局部有限型的概形态射，$m \in \mathbf{Z}$，
$E, K$ 是 $D(\mathcal{O}_X)$ 的对象。若 $E \oplus K$
相对于 $S$ 是 $m$-伪凝聚的，则 $E$ 和 $K$ 也如此。
\end{lemma}

\begin{proof}
这由
上同调章引理 \ref{cohomology-lemma-summands-pseudo-coherent}
及定义得到。
\end{proof}

\begin{lemma}
\label{more-morphisms-lemma-complex-relative-pseudo-coherent-modules}
设 $X \to S$ 是局部有限型的概形态射，$m \in \mathbf{Z}$。
设 $\mathcal{F}^\bullet$ 是（局部）上有界的 $\mathcal{O}_X$-模复形，
使得对所有 $i$，$\mathcal{F}^i$ 相对于 $S$ 都是
$(m - i)$-伪凝聚的。则 $\mathcal{F}^\bullet$ 相对于 $S$
是 $m$-伪凝聚的。
\end{lemma}

\begin{proof}
这由
上同调章引理 \ref{cohomology-lemma-complex-pseudo-coherent-modules}
及定义得到。
\end{proof}

\begin{lemma}
\label{more-morphisms-lemma-cohomology-relative-pseudo-coherent}
设 $X \to S$ 是局部有限型的概形态射，$m \in \mathbf{Z}$，
$E$ 是 $D(\mathcal{O}_X)$ 的对象。若 $E$（局部）上有界，
并且对所有 $i$，$H^i(E)$ 相对于 $S$ 都是 $(m - i)$-伪凝聚的，
则 $E$ 相对于 $S$ 是 $m$-伪凝聚的。
\end{lemma}

\begin{proof}
这由
上同调章引理 \ref{cohomology-lemma-cohomology-pseudo-coherent}
及定义得到。
\end{proof}

\begin{lemma}
\label{more-morphisms-lemma-base-change-relative-pseudo-coherent}
设 $X \to S$ 是局部有限型的概形态射，$m \in \mathbf{Z}$。
设 $E$ 是 $D(\mathcal{O}_X)$ 的对象，且相对于 $S$ 是 $m$-伪凝聚的。
设 $S' \to S$ 是概形态射。置 $X' = X \times_S S'$，并以 $E'$
表示 $E$ 到 $X'$ 的导出拉回。若 $S'$ 与 $X$ 在 $S$ 上 Tor 独立，
则 $E'$ 相对于 $S'$ 是 $m$-伪凝聚的。
\end{lemma}

\begin{proof}
问题在 $X$ 与 $X'$ 上是局部的，故可假设 $X$、$S$、$S'$、$X'$
都是仿射的。选取闭浸入 $i : X \to \mathbf{A}^n_S$，并以
$i' : X' \to \mathbf{A}^n_{S'}$ 表示它到 $S'$ 的基变换。
以 $g : X' \to X$ 和 $g' : \mathbf{A}^n_{S'} \to \mathbf{A}^n_S$
表示投影，于是 $E' = Lg^*E$。由于 $X$ 与 $S'$ 在 $S$ 上 Tor 独立，
基变换映射
(上同调章说明 \ref{cohomology-remark-base-change})
诱导同构
$$
Ri'_*(Lg^*E) = L(g')^*Ri_*E
$$
。具体而言，对位于 $x \in X$ 上方的点 $x' \in X'$，
基变换映射在 $x'$ 处的茎上的映射为
$$
E_x \otimes_{\mathcal{O}_{\mathbf{A}^n_S, x}}^\mathbf{L}
\mathcal{O}_{\mathbf{A}^n_{S'}, x'}
\longrightarrow
E_x \otimes_{\mathcal{O}_{X, x}}^\mathbf{L}
\mathcal{O}_{X', x'}
$$
，它来自闭浸入 $i$ 与 $i'$。注意，源拟同构于
$E_x \otimes_{\mathcal{O}_{S, s}}^\mathbf{L} \mathcal{O}_{S', s'}$
的一个局部化；它同构于目标，因为根据假设，
$\mathcal{O}_{X', x'}$ 同构于
$\mathcal{O}_{X, x} \otimes_{\mathcal{O}_{S, s}}^\mathbf{L}
\mathcal{O}_{S', s'}$ 的（同一个）局部化。应用
上同调章引理 \ref{cohomology-lemma-pseudo-coherent-pullback}.
即得引理。
\end{proof}

\begin{lemma}
\label{more-morphisms-lemma-pull-relative-pseudo-coherent}
设 $f : X \to Y$ 是基 $S$ 上局部有限型的概形态射。
设 $m \in \mathbf{Z}$，$E$ 是 $D(\mathcal{O}_Y)$ 的对象。假设
\begin{enumerate}
\item $\mathcal{O}_X$ 相对于 $Y$ 是伪凝聚的\footnote{这意味着
$f$ 是伪凝聚态射，见 定义 \ref{more-morphisms-definition-pseudo-coherent}。}；
\item $E$ 相对于 $S$ 是 $m$-伪凝聚的。
\end{enumerate}
则 $Lf^*E$ 相对于 $S$ 是 $m$-伪凝聚的。
\end{lemma}

\begin{proof}
问题在 $X$ 上是局部的。因此可以假设 $X$、$Y$、$S$ 都仿射。
仿照 代数进阶章引理
\ref{more-algebra-lemma-pull-relative-pseudo-coherent}
的证明进行论证，可以找到交换图
$$
\xymatrix{
X \ar[r]_i \ar[d]_f &
\mathbf{A}^d_Y \ar[r]_j \ar[ld]^p &
\mathbf{A}^{n + d}_S \ar[ld] \\
Y \ar[r] &
\mathbf{A}^n_S
}
$$
注意
$$
Ri_* Lf^*E = Ri_* Li^* Lp^*E =
Lp^*E \otimes_{\mathcal{O}_{\mathbf{A}_Y^n}}^\mathbf{L} Ri_*\mathcal{O}_X
$$
，这是由 上同调章引理
\ref{cohomology-lemma-projection-formula-closed-immersion}.
得到的。由假设以及 $Y$ 仿射这一事实，可以用有限自由
% P05-EN-0129: frozen source switches the diagram's relative affine dimension d to n in coefficient sheaves.
$\mathcal{O}_{\mathbf{A}_Y^n}$-模复形 $\mathcal{F}^\bullet$
表示 $Ri_*\mathcal{O}_X = i_*\mathcal{O}_X$，并且
$\mathcal{F}^q = 0$ 对 $q > 0$ 成立
（细节略；使用 概形的导出范畴章引理
\ref{perfect-lemma-pseudo-coherent-affine}
和
代数进阶章引理
\ref{more-algebra-lemma-rel-n-pseudo-module}）。
由假设，$E$ 上有界；设对 $q > a$ 有 $H^q(E) = 0$。
用满足对 $q > a$ 有 $\mathcal{E}^q = 0$ 的 $\mathcal{O}_Y$-模复形
$\mathcal{E}^\bullet$ 表示 $E$。于是上述导出张量积由
$\text{Tot}(p^*\mathcal{E}^\bullet
\otimes_{\mathcal{O}_{\mathbf{A}_Y^n}} \mathcal{F}^\bullet)$ 表示。

\medskip\noindent
因为 $j$ 是闭浸入，函子 $j_*$ 正合，而且可以通过对任一代表层复形
应用 $j_*$ 来计算 $Rj_*$。因此必须证明
% P05-EN-0130: frozen source uses ambient dimension n+m here although the diagram has n+d.
$j_*\text{Tot}(p^*\mathcal{E}^\bullet \otimes_{\mathcal{O}_{\mathbf{A}_Y^n}}
\mathcal{F}^\bullet)$ 作为 $\mathcal{O}_{\mathbf{A}^{n + m}_S}$-模复形
是 $m$-伪凝聚的。注意
$\text{Tot}(p^*\mathcal{E}^\bullet \otimes_{\mathcal{O}_{\mathbf{A}_Y^n}}
\mathcal{F}^\bullet)$ 具有由子复形构成的滤过，其相继商为复形
$p^*\mathcal{E}^\bullet
\otimes_{\mathcal{O}_{\mathbf{A}_Y^n}} \mathcal{F}^q[-q]$。
注意，当 $q \ll 0$ 时，复形
$p^*\mathcal{E}^\bullet \otimes_{\mathcal{O}_{\mathbf{A}_Y^n}}
\mathcal{F}^q[-q]$
在次数 $\leq m$ 上的上同调为零，因而是 $m$-伪凝聚的。因此，应用
引理 \ref{more-morphisms-lemma-cone-relatively-pseudo-coherent}
并作归纳，只需证明
% P05-EN-0131: frozen source asks for full pseudo-coherence here but the next reduction proves only m-pseudo-coherence.
$p^*\mathcal{E}^\bullet \otimes_{\mathcal{O}_{\mathbf{A}_Y^n}}
\mathcal{F}^q[-q]$
对所有 $q$ 都相对于 $S$ 是伪凝聚的。注意，对 $q > 0$
有 $\mathcal{F}^q = 0$。又因为 $\mathcal{F}^q$ 有限自由，
这就归结为证明 $p^*\mathcal{E}^\bullet$ 相对于 $S$
是 $m$-伪凝聚的；例如，这由
引理 \ref{more-morphisms-lemma-base-change-relative-pseudo-coherent}
得到。
\end{proof}

\begin{lemma}
\label{more-morphisms-lemma-composition-relative-pseudo-coherent}
设 $f : X \to Y$ 是基 $S$ 上局部有限型的概形态射。
设 $m \in \mathbf{Z}$，$E$ 是 $D(\mathcal{O}_X)$ 的对象。
假设 $\mathcal{O}_Y$ 相对于 $S$ 是伪凝聚的\footnote{这意味着
$Y \to S$ 是伪凝聚态射，见 定义 \ref{more-morphisms-definition-pseudo-coherent}。}。
则下列条件等价：
\begin{enumerate}
\item $E$ 相对于 $Y$ 是 $m$-伪凝聚的；
\item $E$ 相对于 $S$ 是 $m$-伪凝聚的。
\end{enumerate}
\end{lemma}

\begin{proof}
问题在 $X$ 上是局部的，故可假设 $X$、$Y$、$S$ 都仿射。
仿照 代数进阶章引理
\ref{more-algebra-lemma-pull-relative-pseudo-coherent}
的证明进行论证，可以找到交换图
$$
\xymatrix{
X \ar[r]_i \ar[d]_f &
\mathbf{A}^m_Y \ar[r]_j \ar[ld]^p &
\mathbf{A}^{n + m}_S \ar[ld] \\
Y \ar[r] &
\mathbf{A}^n_S
}
$$
假设 $\mathcal{O}_Y$ 相对于 $S$ 是伪凝聚的，这蕴含
$\mathcal{O}_{\mathbf{A}^m_Y}$ 相对于 $\mathbf{A}^m_S$ 是伪凝聚的
（通过平坦基变换；例如可用
引理 \ref{more-morphisms-lemma-base-change-relative-pseudo-coherent} 看出）。
% P05-EN-0132: frozen source writes A^n_Y here although j has domain A^m_Y in the displayed diagram.
这又蕴含 $j_*\mathcal{O}_{\mathbf{A}^n_Y}$ 作为
$\mathcal{O}_{\mathbf{A}^{n + m}_S}$-模是伪凝聚的。于是引理中的等价性由
概形的导出范畴章引理
\ref{perfect-lemma-closed-push-pseudo-coherent}.
得到。
\end{proof}

\begin{lemma}
\label{more-morphisms-lemma-check-relative-pseudo-coherence-on-charts}
设
$$
\xymatrix{
X \ar[rd] \ar[rr]_i & & P \ar[ld] \\
& S
}
$$
是概形的交换图。假设 $i$ 是闭浸入，且 $P \to S$ 平坦并局部有限表现。
设 $E$ 是 $D(\mathcal{O}_X)$ 的对象。则下列条件等价：
\begin{enumerate}
\item $E$ 相对于 $S$ 是 $m$-伪凝聚的；
\item $Ri_*E$ 相对于 $S$ 是 $m$-伪凝聚的；
\item $Ri_*E$ 在 $P$ 上是 $m$-伪凝聚的。
\end{enumerate}
\end{lemma}

\begin{proof}
条件 (1) 与 (2) 的等价性就是
引理 \ref{more-morphisms-lemma-finite-morphism-relative-pseudo-coherence}.
条件 (2) 与 (3) 的等价性可由
引理
\ref{more-morphisms-lemma-composition-relative-pseudo-coherent}
应用于 $\text{id} : P \to P$ 得到，只须证明 $\mathcal{O}_P$
相对于 $S$ 是伪凝聚的。这由
代数进阶章引理
\ref{more-algebra-lemma-flat-finite-presentation-perfect}
及定义得到。
\end{proof}









\section{伪凝聚态射}
\label{more-morphisms-section-pseudo-coherent}

\noindent
务必避免阅读本节。如果确实需要其中的一些材料，请先阅读相应的代数章节，见
代数进阶章第 \ref{more-algebra-section-pseudo-coherent} 节,
\ref{more-algebra-section-relative-pseudo-coherent}, 及
\ref{more-algebra-section-pseudo-coherent-perfect-ring-map}.
目前只需知道：环映射 $A \to B$ 伪凝聚，当且仅当
$B = A[x_1, \ldots, x_n]/I$，并且 $B$ 作为
$A[x_1, \ldots, x_n]$-模具有由有限自由
$A[x_1, \ldots, x_n]$-模构成的分解。

\begin{lemma}
\label{more-morphisms-lemma-pseudo-coherent}
设 $f : X \to S$ 是概形态射。下列条件等价：
\begin{enumerate}
\item 存在仿射开覆盖 $S = \bigcup V_j$，且对每个 $j$
存在仿射开覆盖 $f^{-1}(V_j) = \bigcup U_{ji}$，使得
% P05-EN-0133: frozen source reverses the covering index U_ji to U_ij in the ring-map clause.
$\mathcal{O}_S(V_j) \to \mathcal{O}_X(U_{ij})$ 是伪凝聚环映射；
\item 对每一对满足 $f(U) \subset V$ 的仿射开集 $U \subset X$、
$V \subset S$，环映射 $\mathcal{O}_S(V) \to \mathcal{O}_X(U)$ 伪凝聚；
\item $f$ 局部有限型，且 $\mathcal{O}_X$ 相对于 $S$ 是伪凝聚的。
\end{enumerate}
\end{lemma}

\begin{proof}
为证明 (1) 与 (2) 等价，只需对伪凝聚环映射这一性质检验
态射章定义 \ref{morphisms-definition-property-local}
的条件 (1)(a)、(b)、(c)。这些性质由下列引理得到
（其中用到局部化是平坦的）：
代数进阶章引理
\ref{more-algebra-lemma-base-change-relative-pseudo-coherent},
\ref{more-algebra-lemma-localize-relative-pseudo-coherent}, 及
\ref{more-algebra-lemma-glue-relative-pseudo-coherent}.

\medskip\noindent
若 (1) 成立，则 $f$ 局部有限型，因为伪凝聚环映射按定义是有限型的。
此外，由 引理 \ref{more-morphisms-lemma-qcoh-relative-pseudo-coherence-characterize}
及定义，(1) 蕴含 $\mathcal{O}_X$ 相对于 $S$ 是伪凝聚的。
反之，若 (3) 成立，则对 (2) 中的每一对 $U$ 和 $V$，
环 $\mathcal{O}_X(U)$ 在 $\mathcal{O}_S(V)$ 上为有限型，
并且 $\mathcal{O}_X(U)$ 作为模相对于 $\mathcal{O}_S(V)$ 是伪凝聚的，见
引理 \ref{more-morphisms-lemma-relative-pseudo-coherence-characterize} 及
\ref{more-morphisms-lemma-qcoh-relative-pseudo-coherence-characterize}.
这正是伪凝聚环映射的定义，故 (2) 和 (1) 成立。
\end{proof}

\begin{definition}
\label{more-morphisms-definition-pseudo-coherent}
若概形态射 $f : X \to S$ 满足
引理 \ref{more-morphisms-lemma-pseudo-coherent}
中的等价条件，则称它是{\it 伪凝聚的}。在这种情形下，
也称 $X$ 在 $S$ 上伪凝聚。
\end{definition}

\noindent
注意，伪凝聚态射的基变换一般不伪凝聚。

\begin{lemma}
\label{more-morphisms-lemma-flat-base-change-pseudo-coherent}
伪凝聚态射的平坦基变换伪凝聚。
\end{lemma}

\begin{proof}
这转化为如下代数结果：设 $A \to B$ 是伪凝聚环映射，
$A \to A'$ 平坦。则 $A' \to B \otimes_A A'$ 伪凝聚。
这由更一般的
代数进阶章
引理 \ref{more-algebra-lemma-base-change-relative-pseudo-coherent}.
得到。
\end{proof}

\begin{lemma}
\label{more-morphisms-lemma-composition-pseudo-coherent}
伪凝聚概形态射的复合伪凝聚。
\end{lemma}

\begin{proof}
这转化为如下代数结果：若 $A \to B \to C$ 是可复合的伪凝聚环映射，
则 $A \to C$ 伪凝聚。这可由
代数进阶章
引理 \ref{more-algebra-lemma-pull-relative-pseudo-coherent}
或
代数进阶章
引理 \ref{more-algebra-lemma-composition-relative-pseudo-coherent}.
中的任一个得到。
\end{proof}

\begin{lemma}
\label{more-morphisms-lemma-pseudo-coherent-finite-presentation}
伪凝聚态射局部有限表现。
\end{lemma}

\begin{proof}
直接由定义得到。
\end{proof}

\begin{lemma}
\label{more-morphisms-lemma-flat-finite-presentation-pseudo-coherent}
平坦且局部有限表现的态射伪凝聚。
\end{lemma}

\begin{proof}
这是因为有限表现的平坦环映射伪凝聚（甚至完美），见
代数进阶章
引理 \ref{more-algebra-lemma-flat-finite-presentation-perfect}.
\end{proof}

\begin{lemma}
\label{more-morphisms-lemma-permanence-pseudo-coherent}
设 $f : X \to Y$ 是基概形 $S$ 上伪凝聚的概形之间的态射。
则 $f$ 伪凝聚。
\end{lemma}

\begin{proof}
这转化为如下代数结果：若 $R \to A \to B$ 是可复合的环映射，
$R \to A$、$R \to B$ 伪凝聚，则
% P05-EN-0134: frozen source concludes R->B, already assumed; the lemma requires the relative map A->B.
$R \to B$ 伪凝聚。这由
代数进阶章
引理 \ref{more-algebra-lemma-composition-relative-pseudo-coherent}.
得到。
\end{proof}

\begin{lemma}
\label{more-morphisms-lemma-finite-pseudo-coherent}
设 $f : X \to S$ 是有限概形态射。则 $f$ 伪凝聚，
当且仅当 $f_*\mathcal{O}_X$ 作为 $\mathcal{O}_S$-模伪凝聚。
\end{lemma}

\begin{proof}
把此引理翻译成代数语言，就是如下陈述：若 $R \to A$ 是有限环映射，
则 $R \to A$ 作为环映射伪凝聚（按定义，这意味着 $A$ 作为
$A$-模相对于 $R$ 伪凝聚），当且仅当 $A$ 作为 $R$-模伪凝聚。
这由更一般的
代数进阶章引理
\ref{more-algebra-lemma-finite-extension-pseudo-coherent}.
得到。
\end{proof}

\begin{lemma}
\label{more-morphisms-lemma-Noetherian-pseudo-coherent}
设 $f : X \to S$ 是概形态射。若 $S$ 局部诺特，
则 $f$ 伪凝聚，当且仅当 $f$ 局部有限型。
\end{lemma}

\begin{proof}
这转化为如下代数结果：若 $R \to A$ 是有限型环映射，且 $R$ 诺特，
则 $R \to A$ 伪凝聚，当且仅当 $R \to A$ 有限型。
为此，注意伪凝聚环映射按定义是有限型的。反之，若 $R \to A$
有限型，则可以写成 $A = R[x_1, \ldots, x_n]/I$；由
代数进阶章
引理 \ref{more-algebra-lemma-Noetherian-pseudo-coherent}
可知 $A$ 作为 $R[x_1, \ldots, x_n]$-模伪凝聚，
即 $R \to A$ 是伪凝聚环映射。
\end{proof}

\begin{lemma}
\label{more-morphisms-lemma-descending-property-pseudo-coherent}
性质 $\mathcal{P}(f) =$“$f$ 伪凝聚”在基上是 fpqc 局部的。
\end{lemma}

\begin{proof}
我们使用
下降章引理 \ref{descent-lemma-descending-properties-morphisms}
的判据来证明。由
定义 \ref{more-morphisms-definition-pseudo-coherent}
，伪凝聚性在基上是 Zariski 局部的。由
引理 \ref{more-morphisms-lemma-flat-base-change-pseudo-coherent}
，伪凝聚性在平坦基变换下保持。
下降章引理 \ref{descent-lemma-descending-properties-morphisms}
的最后一个假设 (3) 转化为如下代数陈述：设 $A \to B$ 是忠实平坦环映射，
$C = A[x_1, \ldots, x_n]/I$ 是 $A$-代数。若 $C \otimes_A B$
作为 $B[x_1, \ldots, x_n]$-模伪凝聚，则 $C$ 作为
$A[x_1, \ldots, x_n]$-模伪凝聚。这就是
代数进阶章引理 \ref{more-algebra-lemma-flat-descent-pseudo-coherent}.
\end{proof}

\begin{lemma}
\label{more-morphisms-lemma-quotient-of-flat-finitely-presented}
设 $A \to B$ 是有限表现的平坦环映射，$I \subset B$ 是理想。
则 $A \to B/I$ 伪凝聚，当且仅当 $I$ 作为 $B$-模伪凝聚。
\end{lemma}

\begin{proof}
选取表现 $B = A[x_1, \ldots, x_n]/J$。注意 $B$ 作为
$A[x_1, \ldots, x_n]$-模伪凝聚，因为由
引理 \ref{more-morphisms-lemma-flat-finite-presentation-pseudo-coherent}.
，$A \to B$ 是伪凝聚环映射。还要注意，$A \to B/I$ 伪凝聚，
当且仅当 $B/I$ 作为 $A[x_1, \ldots, x_n]$-模伪凝聚。由
代数进阶章引理 \ref{more-algebra-lemma-finite-push-pseudo-coherent}
，这等价于 $B/I$ 作为 $B$-模伪凝聚。短正合列
$0 \to I \to B \to B/I \to 0$ 表明，$I$ 伪凝聚当且仅当
$B/I$ 伪凝聚（见
代数进阶章
引理 \ref{more-algebra-lemma-two-out-of-three-pseudo-coherent}），
故引理得证。
\end{proof}

\noindent
下面的引理将被更强的
引理 \ref{more-morphisms-lemma-pseudo-coherent-fppf-local-source}.
取代。

\begin{lemma}
\label{more-morphisms-lemma-pseudo-coherent-syntomic-local-source}
性质 $\mathcal{P}(f) =$“$f$ 伪凝聚”在源上是 syntomic 局部的。
\end{lemma}

\begin{proof}
我们使用
下降章引理 \ref{descent-lemma-properties-morphisms-local-source}
的判据来证明。由
引理 \ref{more-morphisms-lemma-flat-finite-presentation-pseudo-coherent} 及
\ref{more-morphisms-lemma-composition-pseudo-coherent}
可知，伪凝聚性在与局部有限表现的平坦态射前复合时保持，
特别是在与 syntomic 态射前复合时保持（见
态射章引理 \ref{morphisms-lemma-syntomic-flat} 及
\ref{morphisms-lemma-syntomic-locally-finite-presentation}）。由
定义 \ref{more-morphisms-definition-pseudo-coherent}
显然可知，伪凝聚性在源和目标上是 Zariski 局部的。因此，按照上述
下降章引理 \ref{descent-lemma-properties-morphisms-local-source}
，只需证明如下陈述：设 $X' \to X \to Y$ 是仿射概形的态射，
$X' \to X$ syntomic，且 $X' \to Y$ 伪凝聚，则 $X \to Y$ 伪凝聚。
首先注意，无论如何，由
下降章引理 \ref{descent-lemma-flat-finitely-presented-permanence-algebra}.
$X \to Y$ 都有限表现。选取闭浸入 $X \to \mathbf{A}^n_Y$。由
代数章引理 \ref{algebra-lemma-lift-syntomic}
，可以找到仿射开覆盖 $X' = \bigcup_{i = 1, \ldots, n} X'_i$，
以及提升态射 $X'_i \to X$ 的 syntomic 态射
$W_i \to \mathbf{A}^n_Y$；也就是说，有纤维积图
$$
\xymatrix{
X'_i \ar[d] \ar[r] & W_i \ar[d] \\
X \ar[r] & \mathbf{A}^n_Y
}
$$
用 $\coprod X'_i$ 替换 $X'$ 并置 $W = \coprod W_i$，得到纤维积图
$$
\xymatrix{
X' \ar[d] \ar[r] & W \ar[d]^h \\
X \ar[r] & \mathbf{A}^n_Y
}
$$
，其中 $W \to \mathbf{A}^n_Y$ 平坦且有限表现，而 $X' \to Y$
仍伪凝聚。由于 $W \to \mathbf{A}^n_Y$ 是开的（见
态射章引理 \ref{morphisms-lemma-fppf-open})
），且 $X' \to X$ 满射，可以找到
$f \in \Gamma(\mathbf{A}^n_Y, \mathcal{O})$，使得
$X \subset D(f) \subset \Im(h)$。写
$Y = \Spec(R)$、$X = \Spec(A)$、$X' = \Spec(A')$、
$W = \Spec(B)$、$A = R[x_1, \ldots, x_n]/I$、$A' = B/IB$。
于是 $R \to A'$ 伪凝聚。图示为
$$
\xymatrix{
A' = B/IB & B \ar[l] \\
A = R[x_1, \ldots, x_n]/I \ar[u] & R[x_1, \ldots, x_n] \ar[l] \ar[u]
}
$$
由
引理 \ref{more-morphisms-lemma-quotient-of-flat-finitely-presented}
可知 $IB$ 作为 $B$-模伪凝聚。由上面对 $f$ 的选择，环映射
$R[x_1, \ldots, x_n]_f \to B_f$ 忠实平坦。这蕴含
$I_f \subset R[x_1, \ldots, x_n]_f$ 伪凝聚，见
代数进阶章引理 \ref{more-algebra-lemma-flat-descent-pseudo-coherent}.
再应用一次
引理 \ref{more-morphisms-lemma-quotient-of-flat-finitely-presented}
，可知 $R \to A$ 伪凝聚。
\end{proof}

\begin{lemma}
\label{more-morphisms-lemma-pseudo-coherent-fppf-local-source}
性质 $\mathcal{P}(f) =$“$f$ 伪凝聚”在源上是 fppf 局部的。
\end{lemma}

\begin{proof}
设 $f : X \to S$ 是概形态射。设 $\{g_i : X_i \to X\}$
是 fppf 覆盖，使得每个复合 $f \circ g_i$ 都伪凝聚。由
引理 \ref{more-morphisms-lemma-dominate-fppf}
，存在
\begin{enumerate}
\item Zariski 开覆盖 $X = \bigcup U_j$；
\item 满射的有限局部自由态射 $W_j \to U_j$；
\item Zariski 开覆盖 $W_j = \bigcup_k W_{j, k}$；
\item 满射的有限局部自由态射 $T_{j, k} \to W_{j, k}$，
\end{enumerate}
使得 fppf 覆盖 $\{h_{j, k} : T_{j, k} \to X\}$ 加细给定覆盖
$\{X_i \to X\}$。以 $\psi_{j, k} : T_{j, k} \to X_{\alpha(j, k)}$
表示见证 $\{T_{j, k} \to X\}$ 加细给定覆盖 $\{X_i \to X\}$ 的态射。
注意 $T_{j, k} \to X$ 平坦且局部有限表现，故由
引理 \ref{more-morphisms-lemma-flat-finite-presentation-pseudo-coherent}.
，$X_i$ 和 $T_{j, k}$ 都在 $X$ 上伪凝聚。因此
$\psi_{j, k} : T_{j, k} \to X_i$ 伪凝聚，见
引理 \ref{more-morphisms-lemma-permanence-pseudo-coherent}.
从而 $T_{j, k} \to S$ 作为 $\psi_{j, k}$ 与
$f \circ g_{\alpha(j, k)}$ 的复合是伪凝聚的，见
引理 \ref{more-morphisms-lemma-composition-pseudo-coherent}.
由此可见，已经把引理归结为 Zariski 开覆盖的情形（这没有问题），
以及覆盖由单个满射的有限局部自由态射给出的情形；后者在下一段处理。

\medskip\noindent
假设 $X' \to X \to S$ 是一列概形态射，其中 $X' \to X$
满射且有限局部自由，并且
% P05-EN-0135: frozen source says X'->Y although this paragraph has base S and no Y.
$X' \to Y$ 伪凝聚。目标是证明 $X \to S$ 伪凝聚。注意，由
下降章引理 \ref{descent-lemma-flat-finitely-presented-permanence}
，态射 $X \to S$ 局部有限表现。显然，问题可归结为 $X'$、$X$、$S$
都仿射，且 $X' \to X$ 为某个秩 $r > 0$ 的自由态射的情形。
相应的代数问题如下：设 $R \to A \to A'$ 是环映射，
$R \to A'$ 伪凝聚，$R \to A$ 有限表现，并且 $A' \cong A^{\oplus r}$
作为 $A$-模。目标是证明 $R \to A$ 伪凝聚。
假设 $R \to A'$ 伪凝聚，意味着 $A'$ 作为 $A'$-模相对于 $R$
伪凝聚。由
代数进阶章
引理 \ref{more-algebra-lemma-finite-extension-pseudo-coherent}
，这蕴含 $A'$ 作为 $A$-模相对于 $R$ 伪凝聚。
由于 $A' \cong A^{\oplus r}$ 作为 $A$-模，故 $A$ 作为 $A$-模
相对于 $R$ 伪凝聚，见
代数进阶章
引理 \ref{more-algebra-lemma-summands-relative-pseudo-coherent}.
按定义，这意味着 $R \to A$ 伪凝聚，证明完成。
\end{proof}








\section{完美态射}
\label{more-morphisms-section-perfect}

\noindent
要理解本节内容，必须稍微了解伪凝聚态射一节的内容。
目前只需知道：环映射 $A \to B$ 完美，当且仅当它伪凝聚，
且 $B$ 作为 $A$-模具有有限 Tor 维数。

\begin{lemma}
\label{more-morphisms-lemma-perfect}
设 $f : X \to S$ 是局部有限型的概形态射。下列条件等价：
\begin{enumerate}
\item 存在仿射开覆盖 $S = \bigcup V_j$，且对每个 $j$
存在仿射开覆盖 $f^{-1}(V_j) = \bigcup U_{ji}$，使得
% P05-EN-0136: frozen source reverses the covering index U_ji to U_ij in the ring-map clause.
$\mathcal{O}_S(V_j) \to \mathcal{O}_X(U_{ij})$ 是完美环映射；
\item 对每一对满足 $f(U) \subset V$ 的仿射开集 $U \subset X$、
$V \subset S$，环映射 $\mathcal{O}_S(V) \to \mathcal{O}_X(U)$ 完美。
\end{enumerate}
\end{lemma}

\begin{proof}
假设 (1)，并设 $U, V$ 如 (2) 所述。由
引理 \ref{more-morphisms-lemma-pseudo-coherent}
可知 $\mathcal{O}_S(V) \to \mathcal{O}_X(U)$ 伪凝聚。
因此，只需证明环映射“具有有限 Tor 维数”这一性质满足
态射章定义 \ref{morphisms-definition-property-local}.
的条件 (1)(a)、(b)、(c)。这些性质由
代数进阶章
引理 \ref{more-algebra-lemma-flat-push-tor-amplitude},
\ref{more-algebra-lemma-flat-base-change-finite-tor-dimension}, 及
\ref{more-algebra-lemma-glue-tor-amplitude}.
得到。略去一些细节。
\end{proof}

\begin{definition}
\label{more-morphisms-definition-perfect}
若概形态射 $f : X \to S$ 满足
引理 \ref{more-morphisms-lemma-perfect}
中的等价条件，则称它是{\it 完美的}。在这种情形下，
也称 $X$ 在 $S$ 上完美。
\end{definition}

\noindent
注意，完美态射特别地伪凝聚，因而局部有限表现。
还要注意，完美态射的基变换一般不完美。

\begin{lemma}
\label{more-morphisms-lemma-flat-base-change-perfect}
完美态射的平坦基变换完美。
\end{lemma}

\begin{proof}
这转化为如下代数结果：设 $A \to B$ 是完美环映射，
$A \to A'$ 平坦。则 $A' \to B \otimes_A A'$ 完美。
关于伪凝聚环映射的相应结果已在
引理 \ref{more-morphisms-lemma-flat-base-change-pseudo-coherent}.
中见过。有限 Tor 维数的相应事实由
代数进阶章
引理 \ref{more-algebra-lemma-flat-base-change-finite-tor-dimension}.
得到。
\end{proof}

\begin{lemma}
\label{more-morphisms-lemma-composition-perfect}
完美概形态射的复合完美。
\end{lemma}

\begin{proof}
这转化为如下代数结果：若 $A \to B \to C$ 是可复合的完美环映射，
则 $A \to C$ 完美。关于伪凝聚性的相应结论已见于
引理 \ref{more-morphisms-lemma-composition-pseudo-coherent}
及其证明。由假设，存在整数 $n$、$m$，使得 $B$ 在 $A$ 上的
Tor 维数 $\leq n$，而 $C$ 在 $B$ 上的 Tor 维数 $\leq m$。
于是对任意 $A$-模 $M$，有
$$
M \otimes_A^{\mathbf{L}} C =
(M \otimes_A^{\mathbf{L}} B) \otimes_B^{\mathbf{L}} C
$$
而
代数进阶章例 \ref{more-algebra-example-tor}
的谱序列表明，当 $p > n + m$ 时
$\text{Tor}^A_p(M, C) = 0$，正是所需结论。
\end{proof}

\begin{lemma}
\label{more-morphisms-lemma-flat-finite-presentation-perfect}
设 $f : X \to S$ 是概形态射。下列条件等价：
\begin{enumerate}
\item $f$ 平坦且完美；
\item $f$ 平坦且局部有限表现。
\end{enumerate}
\end{lemma}

\begin{proof}
蕴含关系 (2) $\Rightarrow$ (1) 就是
代数进阶章
引理 \ref{more-algebra-lemma-flat-finite-presentation-perfect}.
反向蕴含由伪凝聚态射局部有限表现这一事实得到，见
引理 \ref{more-morphisms-lemma-pseudo-coherent-finite-presentation}.
\end{proof}

\begin{lemma}
\label{more-morphisms-lemma-regular-target-perfect}
设 $f : X \to S$ 是概形态射。假设 $S$ 正则且 $f$ 局部有限型。
则 $f$ 完美。
\end{lemma}

\begin{proof}
见
代数进阶章引理 \ref{more-algebra-lemma-regular-perfect-ring-map}.
\end{proof}

\begin{lemma}
\label{more-morphisms-lemma-regular-immersion-perfect}
概形的正则浸入完美。概形的 Koszul-正则浸入完美。
\end{lemma}

\begin{proof}
由于正则浸入是 Koszul-正则浸入，见
除子章引理 \ref{divisors-lemma-regular-quasi-regular-immersion},
只需证明第二个陈述。这转化为如下代数陈述：假设 $I \subset A$
是由 $A$ 的 Koszul-正则列 $f_1, \ldots, f_r$ 生成的理想，
则 $A \to A/I$ 是完美环映射。由于 $A \to A/I$ 满射，
这给出 $A/I$ 由 $A$ 上多项式代数所得的一个表现。
因此，只需说明 $A/I$ 作为 $A$-模伪凝聚且具有有限 Tor 维数。
按 Koszul 列的定义，Koszul 复形 $K(A, f_1, \ldots, f_r)$
是 $A/I$ 的有限自由分解。因此 $A/I$ 是 $A$-模的完美复形，结论得证。
\end{proof}

\begin{lemma}
\label{more-morphisms-lemma-perfect-permanence}
设
$$
\xymatrix{
X \ar[rr]_f \ar[rd] & & Y \ar[ld] \\
& S
}
$$
是概形态射的交换图。假设 $Y \to S$ 光滑且 $X \to S$ 完美。
则 $f : X \to Y$ 完美。
\end{lemma}

\begin{proof}
可以把 $f$ 分解为复合
$$
X \longrightarrow X \times_S Y \longrightarrow Y
$$
，其中第一个态射是映射 $i = (1, f)$，第二个态射是投影。
由于 $Y \to S$ 平坦，见
态射章引理 \ref{morphisms-lemma-smooth-flat},
由 引理 \ref{more-morphisms-lemma-flat-base-change-perfect} 可知
$X \times_S Y \to Y$ 完美。
由于 $Y \to S$ 光滑，$X \times_S Y \to X$ 也光滑，见
态射章引理 \ref{morphisms-lemma-base-change-smooth}.
因此 $i$ 是光滑态射的截面，从而 $i$ 是正则浸入，见
除子章引理 \ref{divisors-lemma-section-smooth-regular-immersion}.
这蕴含 $i$ 完美，见
引理 \ref{more-morphisms-lemma-regular-immersion-perfect}.
因为完美态射的复合完美，故 $f$ 完美，见
引理 \ref{more-morphisms-lemma-composition-perfect}.
\end{proof}

\begin{remark}
\label{more-morphisms-remark-perfect-permanence}
基 $S$ 上完美概形 $X, Y$ 之间的态射未必完美。例如，取
$S = \Spec(k)$、$X = \Spec(k)$，
% P05-EN-0137: frozen source omits the closing parenthesis in this Spec expression.
$Y = \Spec(k[x]/(x^2)$，并取 $X \to Y$ 为唯一的 $S$-态射。
\end{remark}

\begin{lemma}
\label{more-morphisms-lemma-descending-property-perfect}
性质 $\mathcal{P}(f) =$“$f$ 完美”在基上是 fpqc 局部的。
\end{lemma}

\begin{proof}
我们使用
下降章引理 \ref{descent-lemma-descending-properties-morphisms}
的判据来证明。由
定义 \ref{more-morphisms-definition-perfect}
，完美性在基上是 Zariski 局部的。由
引理 \ref{more-morphisms-lemma-flat-base-change-perfect}
，完美性在平坦基变换下保持。
下降章引理 \ref{descent-lemma-descending-properties-morphisms}
的最后一个假设 (3) 转化为如下代数陈述：设 $A \to B$ 是忠实平坦环映射，
$C = A[x_1, \ldots, x_n]/I$ 是 $A$-代数。若 $C \otimes_A B$
作为 $B[x_1, \ldots, x_n]$-模完美，则 $C$ 作为
$A[x_1, \ldots, x_n]$-模完美。这就是
代数进阶章引理 \ref{more-algebra-lemma-flat-descent-perfect}.
\end{proof}

\begin{lemma}
\label{more-morphisms-lemma-check-perfect-stalks}
设 $f : X \to S$ 是伪凝聚概形态射。下列条件等价：
\begin{enumerate}
\item $f$ 完美；
\item $\mathcal{O}_X$ 作为 $f^{-1}\mathcal{O}_S$-模层局部具有有限 Tor 维数；
\item 对所有 $x \in X$，环 $\mathcal{O}_{X, x}$ 作为
$\mathcal{O}_{S, f(x)}$-模具有有限 Tor 维数。
\end{enumerate}
\end{lemma}

\begin{proof}
问题在 $X$ 和 $S$ 上是局部的。因此可以假设
$X = \Spec(B)$、$S = \Spec(A)$，且 $f$ 对应于伪凝聚环映射
$A \to B$。

\medskip\noindent
若 (1) 成立，则 $B$ 作为 $A$-模具有有限 Tor 维数 $d$。
于是对每个满足 $\mathfrak p = A \cap \mathfrak q$ 的素理想
$\mathfrak q \subset B$，$B_\mathfrak q$ 作为 $A_\mathfrak p$-模
具有 Tor 维数 $d$，见
代数进阶章引理 \ref{more-algebra-lemma-tor-amplitude-localization}.
于是由
上同调章引理 \ref{cohomology-lemma-tor-amplitude-stalk}.
，$\mathcal{O}_X$ 作为 $f^{-1}\mathcal{O}_S$-模层具有 Tor 维数 $d$。
故 (1) 蕴含 (2)。

\medskip\noindent
由 上同调章引理 \ref{cohomology-lemma-tor-amplitude-stalk}，
(2) 蕴含 (3)。

\medskip\noindent
假设 (3)。不能使用
代数进阶章引理 \ref{more-algebra-lemma-tor-amplitude-localization}
直接得出结论，因为并未给定 $B_\mathfrak q$ 在 $A_\mathfrak p$
上的 Tor 维数具有与 $\mathfrak q$ 无关的界。选取表现
$A[x_1, \ldots, x_n] \to B$。则 $B$ 作为
$A[x_1, \ldots, x_n]$-模伪凝聚。设
$\mathfrak q \subset A[x_1, \ldots, x_n]$ 是位于
$\mathfrak p \subset A$ 上方的素理想。则 $B_\mathfrak q$ 或者为零，
或者由假设，它作为 $A_\mathfrak p$-模具有有限 Tor 维数。
由于 $A \to A[x_1, \ldots, x_n]$ 的纤维具有有限整体维数，可以把
代数进阶章引理 \ref{more-algebra-lemma-perfect-over-polynomial-ring}
应用于 $A_\mathfrak p \to A[x_1, \ldots, x_n]_\mathfrak q$，
从而看出 $B_\mathfrak q$ 是完美
$A[x_1, \ldots, x_n]_\mathfrak q$-模。因此，由
代数进阶章引理 \ref{more-algebra-lemma-check-perfect-stalks}.
$B$ 是完美 $A[x_1, \ldots, x_n]$-模。故按定义，
$A \to B$ 是完美环映射。
\end{proof}

\begin{lemma}
\label{more-morphisms-lemma-perfect-closed-immersion-perfect-direct-image}
设 $i : Z \to X$ 是完美的概形闭浸入。则 $i_*\mathcal{O}_Z$
是完美 $\mathcal{O}_X$-模，即它是 $D(\mathcal{O}_X)$ 的完美对象。
\end{lemma}

\begin{proof}
这基本上直接由定义得到。具体地，设 $U = \Spec(A)$ 是 $X$
的仿射开集。则对某个理想 $I \subset A$，有
$i^{-1}(U) = \Spec(A/I)$；并且由
代数进阶章引理 \ref{more-algebra-lemma-perfect-ring-map}.
，$A/I$ 具有由有限投射 $A$-模构成的有限分解。因此
$i_*\mathcal{O}_Z|_U$ 可由有限局部自由 $\mathcal{O}_U$-模的
有限长复形表示。这正是所需结论，见
上同调章第 \ref{cohomology-section-perfect} 节.
\end{proof}

\begin{lemma}
\label{more-morphisms-lemma-perfect-proper-perfect-direct-image}
设 $S$ 是诺特概形，$f : X \to S$ 是完美固有概形态射。
设 $E \in D(\mathcal{O}_X)$ 完美。则 $Rf_*E$ 是
$D(\mathcal{O}_S)$ 的完美对象。
\end{lemma}

\begin{proof}
我们断言 概形的导出范畴章引理
\ref{perfect-lemma-perfect-direct-image} 适用。
适用。条件 (1) 和 (2) 是直接的。条件 (3) 在 $X$ 上是局部的。
因此可以假设 $X$ 和 $S$ 仿射，且 $E$ 由严格完美的
$\mathcal{O}_X$-模复形表示。于是只需证明 $\mathcal{O}_X$
作为 $f^{-1}\mathcal{O}_S$-模层具有有限 Tor 维数。
由 引理 \ref{more-morphisms-lemma-check-perfect-stalks}，这等价于完美性。
\end{proof}

\begin{lemma}
\label{more-morphisms-lemma-perfect-fppf-local-source}
性质 $\mathcal{P}(f) =$“$f$ 完美”在源上是 fppf 局部的。
\end{lemma}

\begin{proof}
设 $\{g_i : X_i \to X\}_{i \in I}$ 是概形的 fppf 覆盖，
$f : X \to S$ 是态射，使得每个 $f \circ g_i$ 都完美。由
引理 \ref{more-morphisms-lemma-pseudo-coherent-fppf-local-source}
可知 $f$ 伪凝聚。因此由
引理 \ref{more-morphisms-lemma-check-perfect-stalks}
，只需检验对所有 $x \in X$，$\mathcal{O}_{X, x}$ 作为
$\mathcal{O}_{S, f(x)}$-模具有有限 Tor 维数。
选取映到 $x$ 的 $i \in I$ 和 $x_i \in X_i$。于是
$\mathcal{O}_{X_i, x_i}$ 在 $\mathcal{O}_{S, f(x)}$ 上具有有限 Tor 维数，
且 $\mathcal{O}_{X, x} \to \mathcal{O}_{X_i, x_i}$ 忠实平坦。
所需结论由
代数进阶章引理 \ref{more-algebra-lemma-flat-descent-tor-amplitude}.
得到。
\end{proof}

\begin{lemma}
\label{more-morphisms-lemma-factor-regular-immersion}
设 $i : Z \to Y$ 和 $j : Y \to X$ 是概形浸入。假设
\begin{enumerate}
\item $X$ 局部诺特；
\item $j \circ i$ 是正则浸入；
\item $i$ 完美。
\end{enumerate}
则 $i$ 和 $j$ 都是正则浸入。
\end{lemma}

\begin{proof}
因为 $X$（从而 $Y$）局部诺特，四种正则浸入概念全部一致；
而且可以在局部环层面检验态射是否为正则浸入，见
除子章引理 \ref{divisors-lemma-Noetherian-scheme-regular-ideal}.
因此结论由
除幂代数章引理 \ref{dpa-lemma-regular-sequence}.
得到。
\end{proof}









\section{局部完全交态射}
\label{more-morphisms-section-lci}

\noindent
在
除子章第 \ref{divisors-section-regular-immersions} 节
中，我们定义了四种不同类型的正则浸入：正则、Koszul-正则、
$H_1$-正则及拟正则。本节考虑在 $X$ 上局部具有如下分解的态射
$f : X \to S$：
$$
\xymatrix{
X \ar[rr]_i \ar[rd] & & \mathbf{A}^n_S \ar[ld] \\
& S
}
$$
其中对 $* \in \{\emptyset, Koszul, H_1, quasi\}$，$i$ 是 $*$-正则浸入。
然而，当 $* = \emptyset$，即要求 $i$ 是正则浸入时，
我们不知道如何证明这一条件与分解的选择无关。另一方面，
我们希望局部完全交态射是完美的，而这只有在 $* = Koszul$
或 $* = \emptyset$ 时才成立。因此，我们把在 $X$ 上局部具有上述分解，
且其中 $i$ 是 Koszul-正则浸入的概形态射 $f : X \to S$，
定义为{\it 局部完全交态射}或{\it Koszul 态射}。
为说明这样定义可行，先证明此性质与所选分解无关。

\begin{lemma}
\label{more-morphisms-lemma-koszul-independence-factorization}
设 $S$ 是概形，$U$、$P$、$P'$ 是 $S$-概形，$u \in U$。
设 $i : U \to P$、$i' : U \to P'$ 是 $S$ 上的浸入。
假设 $P$ 和 $P'$ 在 $S$ 上光滑。则下列条件等价：
\begin{enumerate}
\item $i$ 在 $u$ 的一个邻域中是 Koszul-正则浸入；
\item $i'$ 在 $u$ 的一个邻域中是 Koszul-正则浸入。
\end{enumerate}
\end{lemma}

\begin{proof}
假设 $i$ 在 $u$ 的一个邻域中是 Koszul-正则浸入。
考虑态射 $j = (i, i') : U \to P \times_S P' = P''$。
因为 $P'' = P \times_S P' \to P$ 光滑，由
除子章引理 \ref{divisors-lemma-lift-regular-immersion-to-smooth}
可知 $j$ 是 Koszul-正则浸入；进而由
除子章引理 \ref{divisors-lemma-push-regular-immersion-thru-smooth}
可知 $i'$ 是 Koszul-正则浸入。
\end{proof}

\noindent
在给出定义之前，先作如下简单说明。设 $f : X \to S$
是局部有限型的概形态射，$x \in X$。则存在开邻域
$U \subset X$，以及 $f|_U$ 的一个分解：先作浸入
$i : U \to \mathbf{A}^n_S$，再作光滑投影
$\mathbf{A}^n_S \to S$。图示为
$$
\xymatrix{
X \ar[rd] & U \ar[l] \ar[d] \ar[r]_-i & \mathbf{A}^n_S = P \ar[ld]^\pi \\
& S
}
$$
事实上，对 $x$ 在 $X$ 中的任意仿射开邻域 $U$ 都可以这样做，见
态射章引理 \ref{morphisms-lemma-quasi-affine-finite-type-over-S}.

\begin{definition}
\label{more-morphisms-definition-lci}
设 $f : X \to S$ 是概形态射。
\begin{enumerate}
\item 设 $x \in X$。若 $f$ 在 $x$ 处为有限型，并且存在开邻域
及 $f|_U$ 的分解 $\pi \circ i$，其中 $i : U \to P$
是 Koszul-正则浸入，$\pi : P \to S$ 光滑，
则称 $f${\it 在 $x$ 处为 Koszul}。
\item 若 $f$ 在每一点处都为 Koszul，则称 $f$ 是
{\it Koszul 态射}，或称 $f$ 是{\it 局部完全交态射}。
\end{enumerate}
\end{definition}

\noindent
上面已经看到，分解 $f|_U = \pi \circ i$ 的选择无关紧要。
也就是说，给定 $f|_U$ 的分解——先作浸入 $i$，再作光滑态射 $\pi$——
$i$ 在 $x$ 的一个邻域中是否 Koszul-正则，是 $f$ 在 $x$ 处的内蕴性质。
为便于引用，把这一点明确记录为引理。

\begin{lemma}
\label{more-morphisms-lemma-lci}
设 $f : X \to S$ 是局部完全交态射，$P$ 是在 $S$ 上光滑的概形。
设 $U \subset X$ 是开子概形，$i : U \to P$ 是 $S$ 上的概形浸入。
则 $i$ 是 Koszul-正则浸入。
\end{lemma}

\begin{proof}
这就是局部完全交态射的定义性质。见上面的讨论。
\end{proof}

\noindent
下面汇集所有 Koszul 态射共有的一些性质。

\begin{lemma}
\label{more-morphisms-lemma-lci-properties}
设 $f : X \to S$ 是局部完全交态射。则
\begin{enumerate}
\item $f$ 局部有限表现；
\item $f$ 伪凝聚；
\item $f$ 完美。
\end{enumerate}
\end{lemma}

\begin{proof}
完美态射伪凝聚（因为完美环映射伪凝聚），而伪凝聚态射局部有限表现
（因为伪凝聚环映射有限表现），所以只需证明最后一个陈述。
完美性是局部性质，故可假设 $f$ 分解为 $\pi \circ i$，
其中 $\pi$ 光滑，而 $i$ 是 Koszul-正则浸入。
Koszul-正则浸入完美，见
引理 \ref{more-morphisms-lemma-regular-immersion-perfect}.
光滑态射平坦且局部有限表现，因而完美，见
引理 \ref{more-morphisms-lemma-flat-finite-presentation-perfect}.
最后，完美态射的复合完美，见
引理 \ref{more-morphisms-lemma-composition-perfect}.
\end{proof}

\begin{lemma}
\label{more-morphisms-lemma-affine-lci}
设 $f : X = \Spec(B) \to S = \Spec(A)$ 是仿射概形态射。
则 $f$ 是局部完全交态射，当且仅当 $A \to B$ 是局部完全交同态，见
代数进阶章定义
\ref{more-algebra-definition-local-complete-intersection}.
\end{lemma}

\begin{proof}
直接由定义得到。
\end{proof}

\noindent
注意，Koszul 态射的基变换一般不是 Koszul 态射。

\begin{lemma}
\label{more-morphisms-lemma-flat-base-change-lci}
局部完全交态射的平坦基变换仍是局部完全交态射。
\end{lemma}

\begin{proof}
略。提示：光滑态射的基变换光滑，而 Koszul-正则浸入的平坦基变换
仍是 Koszul-正则浸入，故结论成立，见
除子章引理 \ref{divisors-lemma-regular-immersion-noetherian}.
\end{proof}

\begin{lemma}
\label{more-morphisms-lemma-composition-lci}
局部完全交态射的复合仍是局部完全交态射。
\end{lemma}

\begin{proof}
设 $g : Y \to S$ 和 $f : X \to Y$ 是局部完全交态射。
设 $x \in X$，并置 $y = f(x)$。选取 $y$ 的开邻域 $V \subset Y$，
以及分解 $g|_V = \pi \circ i$，其中 $i : V \to P$
是 Koszul-正则浸入，$\pi : P \to S$ 光滑。
接着选取 $x \in X$ 的开邻域 $U$，以及分解
$f|_U = \pi' \circ i'$，其中 $i' : U \to P'$ 是 Koszul-正则浸入，
$\pi' : P' \to Y$ 光滑。事实上，可以假设
$P' = \mathbf{A}^n_V$，见 定义 \ref{more-morphisms-definition-lci}
前后的讨论。图示为：
$$
\xymatrix{
X \ar[d] & U \ar[l] \ar[r]_-{i'} & P' = \mathbf{A}^n_V \ar[d] \\
Y \ar[d] &  & V \ar[ll] \ar[r]_i & P \ar[d] \\
S & & & S \ar[lll]
}
$$
置 $P'' = \mathbf{A}^n_P$。则 $U \to P' \to P''$
是 Koszul-正则浸入的复合，即 $i'$ 与 $i$ 沿 $P'' \to P$
的平坦基变换之复合，因而是 Koszul-正则浸入，见
除子章
引理 \ref{divisors-lemma-regular-immersion-noetherian}
及
除子章引理 \ref{divisors-lemma-composition-regular-immersion}.
此外，$P'' \to P \to S$ 作为光滑态射的复合是光滑的，见
态射章引理 \ref{morphisms-lemma-composition-smooth}.
因此，$X \to S$ 如所需在 $x$ 处为 Koszul。
\end{proof}

\begin{lemma}
\label{more-morphisms-lemma-flat-lci}
\begin{slogan}
态射平坦且为局部完全交，当且仅当它是 syntomic 的。
\end{slogan}
设 $f : X \to S$ 是概形态射。下列条件等价：
\begin{enumerate}
\item $f$ 平坦且为局部完全交态射；
\item $f$ syntomic。
\end{enumerate}
\end{lemma}

\begin{proof}
在仿射局部工作时，这就是
代数进阶章引理 \ref{more-algebra-lemma-syntomic-lci}.
下面也给出一个更几何的证明。

\medskip\noindent
假设 (2)。由
态射章引理 \ref{morphisms-lemma-syntomic-locally-standard-syntomic}
，对 $X$ 的每个点 $x$，存在 $x$ 的仿射开邻域 $U$ 和 $f(x)$
的仿射开邻域 $V$，使得 $f|_U : U \to V$ 是标准 syntomic 态射。
这意味着 $U = \Spec(R[x_1, \ldots, x_n]/(f_1, \ldots, f_c))
\to V = \Spec(R)$，其中 $R[x_1, \ldots, x_n]/(f_1, \ldots, f_c)$
是 $R$ 上的相对整体完全交。由
代数章
引理 \ref{algebra-lemma-relative-global-complete-intersection-conormal}
，对每个素理想 $\mathfrak q \supset (f_1, \ldots, f_c)$，列
$f_1, \ldots, f_c$ 在局部环
$R[x_1, \ldots, x_n]_{\mathfrak q}$ 中都是正则列。
考虑 Koszul 复形
$K_\bullet = K_\bullet(R[x_1, \ldots, x_n], f_1, \ldots, f_c)$，
其同调群为 $H_i = H_i(K_\bullet)$。由
代数进阶章引理 \ref{more-algebra-lemma-regular-koszul-regular}
可知，对上述每个 $\mathfrak q$，当 $i > 0$ 时有
$(H_i)_{\mathfrak q} = 0$。另一方面，由
代数进阶章引理 \ref{more-algebra-lemma-homotopy-koszul}
可知 $(f_1, \ldots, f_c)$ 零化 $H_i$。因此，当 $i > 0$ 时
$H_i = 0$，并且 $f_1, \ldots, f_c$ 是 Koszul-正则列。
这证明 $U \to V$ 如所需分解为 Koszul-正则浸入
$U \to \mathbf{A}^n_V$ 后接一个光滑态射。

\medskip\noindent
假设 (1)。则 $f$ 平坦且局部有限表现
（引理 \ref{more-morphisms-lemma-lci-properties}）。因此，根据
态射章引理 \ref{morphisms-lemma-syntomic-locally-standard-syntomic}
，只需证明局部环 $\mathcal{O}_{X_s, x}$ 是局部完全交环。
在 $X$ 上局部选取分解 $f = \pi \circ i$，其中
$i : X \to P$ 是 Koszul-正则浸入，$\pi : P \to S$ 光滑。
注意 $X \to P$ 是 $S$ 上的相对拟正则浸入，见
除子章定义 \ref{divisors-definition-relative-H1-regular-immersion}.
因此，由
除子章
引理 \ref{divisors-lemma-relative-regular-immersion-flat-in-neighbourhood}
可知 $X \to P$ 是正则浸入，且这一性质在任意基变换后仍成立。
因此每个纤维都是正则浸入，从而 $X$ 的所有纤维的所有局部环
都是局部完全交。
\end{proof}

\begin{lemma}
\label{more-morphisms-lemma-regular-immersion-lci}
概形的正则浸入是局部完全交态射。
概形的 Koszul-正则浸入是局部完全交态射。
\end{lemma}

\begin{proof}
由于正则浸入是 Koszul-正则浸入，见
除子章引理 \ref{divisors-lemma-regular-quasi-regular-immersion},
只需证明第二个陈述，而它直接由定义得到。
\end{proof}

\begin{lemma}
\label{more-morphisms-lemma-lci-permanence}
设
$$
\xymatrix{
X \ar[rr]_f \ar[rd] & & Y \ar[ld] \\
& S
}
$$
是概形态射的交换图。假设 $Y \to S$ 光滑，且 $X \to S$
是局部完全交态射。则 $f : X \to Y$ 是局部完全交态射。
\end{lemma}

\begin{proof}
直接由定义得到。
\end{proof}

\begin{lemma}
\label{more-morphisms-lemma-morphism-regular-schemes-is-lci}
设 $f : X \to Y$ 是概形态射。若 $f$ 局部有限型，
且 $X$ 和 $Y$ 正则，则 $f$ 是局部完全交态射。
\end{lemma}

\begin{proof}
可以假设存在分解 $X \to \mathbf{A}^n_Y \to Y$，其中第一个箭头是浸入。
由于 $Y$ 正则，由
代数章引理 \ref{algebra-lemma-regular-goes-up}.
可知 $\mathbf{A}^n_Y$ 也正则。因此由
除子章引理 \ref{divisors-lemma-immersion-regular-regular-immersion}.
$X \to \mathbf{A}^n_Y$ 是正则浸入。
\end{proof}

\noindent
下面的引理具有不同的性质。

\begin{lemma}
\label{more-morphisms-lemma-lci-avramov}
设
$$
\xymatrix{
X \ar[rr]_f \ar[rd] & & Y \ar[ld] \\
& S
}
$$
是概形态射的交换图。假设
\begin{enumerate}
\item $S$ 局部诺特；
\item $Y \to S$ 局部有限型；
\item $f : X \to Y$ 完美；
\item $X \to S$ 是局部完全交态射。
\end{enumerate}
则 $X \to Y$ 是局部完全交态射，并且对所有 $x \in X$，
$Y \to S$ 在 $f(x)$ 处为 Koszul。
\end{lemma}

\begin{proof}
在此证明中，所有概形都局部诺特，所有环都诺特。我们将不再另行说明地使用：
在这种情形下，正则列与 Koszul-正则列一致，见 代数进阶章引理
\ref{more-algebra-lemma-noetherian-finite-all-equivalent}.
此外，理想（相应地，理想层）的正则性可以在局部环（相应地，茎）上检验，见
代数章引理 \ref{algebra-lemma-regular-sequence-in-neighbourhood}
（相应地，除子章引理 \ref{divisors-lemma-Noetherian-scheme-regular-ideal}）。

\medskip\noindent
问题是局部的。因此可以假设 $S$、$X$、$Y$ 都仿射。
在这种情形下，可以选取交换图
$$
\xymatrix{
\mathbf{A}^{n + m}_S \ar[d] & X \ar[l] \ar[d] \\
\mathbf{A}^n_S \ar[d] & Y \ar[l] \ar[ld] \\
S
}
$$
，其中水平箭头是闭浸入。设 $x \in X$ 是一点，并考虑相应的局部环交换图
$$
\xymatrix{
J \ar[r] &
\mathcal{O}_{\mathbf{A}^{n + m}_S, x} \ar[r] &
\mathcal{O}_{X, x} \\
I \ar[r] \ar[u] &
\mathcal{O}_{\mathbf{A}^n_S, f(x)} \ar[r] \ar[u] &
\mathcal{O}_{Y, f(x)} \ar[u]
}
$$
，其中 $J$ 和 $I$ 是水平箭头的核。由于 $X \to S$ 是局部完全交态射，
理想 $J$ 由正则列生成。由于 $X \to Y$ 完美，环
$\mathcal{O}_{X, x}$ 在 $\mathcal{O}_{Y, f(x)}$ 上具有有限 Tor 维数。
因此可以应用
除幂代数章引理 \ref{dpa-lemma-perfect-map-ci}
断定 $I$ 和 $J/I$ 都由正则列生成。由开头的说明，证明完毕。
\end{proof}

\begin{lemma}
\label{more-morphisms-lemma-lci-to-regular}
设
$$
\xymatrix{
X \ar[rr]_f \ar[rd] & & Y \ar[ld] \\
& S
}
$$
是概形态射的交换图。假设 $S$ 局部诺特，$Y \to S$ 局部有限型，
$Y$ 正则，且 $X \to S$ 是局部完全交态射。
则 $f : X \to Y$ 是局部完全交态射，并且对所有 $x \in X$，
$Y \to S$ 在 $f(x)$ 处为 Koszul。
\end{lemma}

\begin{proof}
这是 引理 \ref{more-morphisms-lemma-lci-avramov} 的特殊情形；这一点由
引理 \ref{more-morphisms-lemma-regular-target-perfect}
（以及 态射章引理 \ref{morphisms-lemma-permanence-finite-type}）得到。
\end{proof}

\begin{lemma}
\label{more-morphisms-lemma-perfect-conormal-free-lci}
设 $i : X \to Y$ 是浸入。若
\begin{enumerate}
\item $i$ 完美；
\item $Y$ 局部诺特；
\item 余法层 $\mathcal{C}_{X/Y}$ 有限局部自由，
\end{enumerate}
则 $i$ 是正则浸入。
\end{lemma}

\begin{proof}
翻译成代数语言，这就是
除幂代数章命题 \ref{dpa-proposition-regular-ideal}.
\end{proof}

\begin{lemma}
\label{more-morphisms-lemma-lci-NL}
% P05-EN-0138: frozen source calls the scheme morphism a homomorphism.
设 $f : X \to Y$ 是局部完全交同态。则朴素余切复形 $\NL_{X/Y}$
是 $D(\mathcal{O}_X)$ 中 Tor 振幅位于 $[-1, 0]$ 的完美对象。
\end{lemma}

\begin{proof}
翻译成代数语言，这就是
代数进阶章引理 \ref{more-algebra-lemma-lci-NL}。
在这一翻译中使用 引理 \ref{more-morphisms-lemma-affine-lci} 和
\ref{more-morphisms-lemma-NL-affine}，以及
概形的导出范畴章引理
\ref{perfect-lemma-affine-compare-bounded},
\ref{perfect-lemma-tor-dimension-affine} 及
\ref{perfect-lemma-perfect-affine}.
\end{proof}

\begin{lemma}
\label{more-morphisms-lemma-perfect-NL-lci}
设 $f : X \to Y$ 是局部诺特概形之间的完美态射。下列条件等价：
\begin{enumerate}
\item $f$ 是局部完全交态射；
\item $\NL_{X/Y}$ 的 Tor 振幅位于 $[-1, 0]$；
\item $\NL_{X/Y}$ 完美且 Tor 振幅位于 $[-1, 0]$。
\end{enumerate}
\end{lemma}

\begin{proof}
翻译成代数语言，这就是
除幂代数章引理 \ref{dpa-lemma-perfect-NL-lci}。
在这一翻译中使用 引理 \ref{more-morphisms-lemma-affine-lci} 和
\ref{more-morphisms-lemma-NL-affine}，以及
概形的导出范畴章引理
\ref{perfect-lemma-affine-compare-bounded},
\ref{perfect-lemma-tor-dimension-affine} 及
\ref{perfect-lemma-perfect-affine}.
\end{proof}

\begin{lemma}
\label{more-morphisms-lemma-flat-fp-NL-lci}
设 $f : X \to Y$ 是有限表现的平坦态射。下列条件等价：
\begin{enumerate}
\item $f$ 是局部完全交态射；
\item $f$ syntomic；
\item $\NL_{X/Y}$ 的 Tor 振幅位于 $[-1, 0]$；
\item $\NL_{X/Y}$ 完美且 Tor 振幅位于 $[-1, 0]$。
\end{enumerate}
\end{lemma}

\begin{proof}
翻译成代数语言，这就是
除幂代数章引理 \ref{dpa-lemma-flat-fp-NL-lci}。
在这一翻译中使用 引理 \ref{more-morphisms-lemma-affine-lci} 和
\ref{more-morphisms-lemma-NL-affine}，以及
概形的导出范畴章引理
\ref{perfect-lemma-affine-compare-bounded},
\ref{perfect-lemma-tor-dimension-affine} 及
\ref{perfect-lemma-perfect-affine}.
\end{proof}

\noindent
下面的引理把光滑态射刻画为对角态射完美的平坦态射。

\begin{lemma}
\label{more-morphisms-lemma-smooth-diagonal-perfect}
设 $f : X \to Y$ 是局部诺特概形之间的有限型态射。
以 $\Delta : X \to X \times_Y X$ 表示对角态射。下列条件等价：
\begin{enumerate}
\item $f$ 光滑；
\item $f$ 平坦，且 $\Delta : X \to X \times_Y X$ 是正则浸入；
\item $f$ 平坦，且 $\Delta : X \to X \times_Y X$ 是局部完全交态射；
\item $f$ 平坦，且 $\Delta : X \to X \times_Y X$ 完美。
\end{enumerate}
\end{lemma}

\begin{proof}
假设 (1)。由 态射章引理 \ref{morphisms-lemma-smooth-flat}，
$f$ 平坦。由 态射章引理 \ref{morphisms-lemma-base-change-smooth}，
投影 $X \times_Y X \to X$ 光滑。因此对角态射是光滑态射的截面，
从而是正则浸入，见
除子章引理 \ref{divisors-lemma-section-smooth-regular-immersion}.
故 (1) $\Rightarrow$ (2)。蕴含关系 (2) $\Rightarrow$ (3) 就是
引理 \ref{more-morphisms-lemma-regular-immersion-lci}.
蕴含关系 (3) $\Rightarrow$ (4) 就是
引理 \ref{more-morphisms-lemma-lci-properties}.
较有趣的蕴含关系 (4) $\Rightarrow$ (1) 直接由
除幂代数章引理 \ref{dpa-lemma-perfect-diagonal} 得到。
\end{proof}

\begin{lemma}
\label{more-morphisms-lemma-descending-property-lci}
性质 $\mathcal{P}(f) =$“$f$ 是局部完全交态射”在基上是 fpqc 局部的。
\end{lemma}

\begin{proof}
设 $f : X \to S$ 是概形态射，$\{S_i \to S\}$ 是 $S$ 的 fpqc 覆盖。
假设 $f$ 的每个基变换 $f_i : X_i \to S_i$ 都是局部完全交态射。
注意，这特别地蕴含 $f$ 局部有限型，见
引理 \ref{more-morphisms-lemma-lci-properties}
和
下降章引理 \ref{descent-lemma-descending-property-locally-finite-type}.
设 $x \in X$。选取 $x$ 的开邻域 $U$，以及 $S$ 上的浸入
$j : U \to \mathbf{A}^n_S$（见 定义 \ref{more-morphisms-definition-lci}
之前的讨论）。必须证明 $j$ 是 Koszul-正则浸入。
由于 $f_i$ 是局部完全交态射，基变换
$j_i : U \times_S S_i \to \mathbf{A}^n_{S_i}$
是 Koszul-正则浸入，见
引理 \ref{more-morphisms-lemma-lci}.
因为 $\{\mathbf{A}^n_{S_i} \to \mathbf{A}^n_S\}$ 是 fpqc 覆盖，由
下降章引理 \ref{descent-lemma-descending-property-regular-immersion}
可知 $j$ 如所需是 Koszul-正则浸入。
\end{proof}

\begin{lemma}
\label{more-morphisms-lemma-lci-syntomic-local-source}
性质 $\mathcal{P}(f) =$“$f$ 是局部完全交态射”在源上是 syntomic 局部的。
\end{lemma}

\begin{proof}
我们使用
下降章引理 \ref{descent-lemma-properties-morphisms-local-source}
的判据来证明。由
引理 \ref{more-morphisms-lemma-flat-lci} 及
\ref{more-morphisms-lemma-composition-lci}
可知，局部完全交性质在与 syntomic 态射前复合时保持。由
定义 \ref{more-morphisms-definition-lci}
显然可知，局部完全交性质在源和目标上是 Zariski 局部的。
因此，按照上述
下降章引理 \ref{descent-lemma-properties-morphisms-local-source}
，只需证明如下陈述：设 $X' \to X \to Y$ 是仿射概形的态射，
$X' \to X$ syntomic，且 $X' \to Y$ 是局部完全交态射，
则 $X \to Y$ 是局部完全交态射。首先注意，无论如何，由
下降章引理 \ref{descent-lemma-flat-finitely-presented-permanence-algebra}.
$X \to Y$ 都有限表现。选取闭浸入 $X \to \mathbf{A}^n_Y$。由
代数章引理 \ref{algebra-lemma-lift-syntomic}
，可以找到仿射开覆盖 $X' = \bigcup_{i = 1, \ldots, n} X'_i$，
以及提升态射 $X'_i \to X$ 的 syntomic 态射
$W_i \to \mathbf{A}^n_Y$；也就是说，有纤维积图
$$
\xymatrix{
X'_i \ar[d] \ar[r] & W_i \ar[d] \\
X \ar[r] & \mathbf{A}^n_Y
}
$$
用 $\coprod X'_i$ 替换 $X'$ 并置 $W = \coprod W_i$，
得到仿射概形的纤维积图
$$
\xymatrix{
X' \ar[d] \ar[r] & W \ar[d]^h \\
X \ar[r] & \mathbf{A}^n_Y
}
$$
，其中 $h : W \to \mathbf{A}^n_Y$ syntomic，而 $X' \to Y$
仍是局部完全交态射。由于 $W \to \mathbf{A}^n_Y$ 是开的（见
态射章引理 \ref{morphisms-lemma-fppf-open})
），且 $X' \to X$ 满射，可见 $X$ 包含于 $W \to \mathbf{A}^n_Y$ 的像中。
选取 $\mathbf{A}^n_Y$ 上的闭浸入
$W \to \mathbf{A}^{n + m}_Y$。现在图形如下：
$$
\xymatrix{
X' \ar[d] \ar[r] & W \ar[d]^h \ar[r] & \mathbf{A}^{n + m}_Y \ar[ld] \\
X \ar[r] & \mathbf{A}^n_Y
}
$$
因为 $h$ syntomic，因而是局部完全交态射（见上文），所以态射
$W \to \mathbf{A}^{n + m}_Y$ 是 Koszul-正则浸入。
因为 $X' \to Y$ 是局部完全交态射，所以态射
$X' \to \mathbf{A}^{n + m}_Y$ 是 Koszul-正则浸入。由
除子章引理 \ref{divisors-lemma-permanence-regular-immersion}
可知 $X' \to W$ 是 Koszul-正则浸入。又因为 Koszul-正则浸入
在目标上是 fpqc 局部的（见
下降章引理 \ref{descent-lemma-descending-property-regular-immersion})
），所以 $X \to \mathbf{A}^n_Y$ 是 Koszul-正则浸入，
这正是必须证明的。
\end{proof}

\begin{lemma}
\label{more-morphisms-lemma-base-change-lci-fibres}
设 $S$ 是概形，$f : X \to Y$ 是 $S$ 上的概形态射。
假设 $X$ 和 $Y$ 在 $S$ 上都平坦且局部有限表现。则集合
$$
\{x \in X \mid f\text{ Koszul at }x\}.
$$
在 $X$ 中是开的，且其构造与任意基变换 $S' \to S$ 可交换。
\end{lemma}

\begin{proof}
该集合按定义是开的（见 定义 \ref{more-morphisms-definition-lci}）。
设 $S' \to S$ 是概形态射。置 $X' = S' \times_S X$、
$Y' = S' \times_S Y$，并以 $f' : X' \to Y'$ 表示 $f$ 的基变换。
设 $x' \in X'$ 是一点，使得 $f'$ 在 $x'$ 处为 Koszul。
以 $s' \in S'$、$x \in X$、$y' \in Y'$、$y \in Y$、$s \in S$
表示 $x'$ 的像。注意 $f$ 局部有限表现，见
态射章引理 \ref{morphisms-lemma-finite-presentation-permanence}.
因此可以选取 $x$ 的仿射邻域 $U \subset X$ 及浸入
$i : U \to \mathbf{A}^n_Y$。以 $U' = S' \times_S U$ 和
$i' : U' \to \mathbf{A}^n_{Y'}$ 表示 $i$ 的基变换。
$f'$ 在 $x'$ 处为 Koszul 这一假设蕴含 $i'$ 在 $x'$ 的一个邻域中
是 Koszul-正则浸入，见
引理 \ref{more-morphisms-lemma-lci}.
概形 $X'$ 作为 $X$ 的基变换，在 $S'$ 上平坦且局部有限表现（见
态射章引理 \ref{morphisms-lemma-base-change-flat} 及
\ref{morphisms-lemma-base-change-finite-presentation}）。
因此 $i'$ 在 $x'$ 的一个邻域中是 $S'$ 上的相对
$H_1$-正则浸入（见
除子章定义 \ref{divisors-definition-relative-H1-regular-immersion}).
所以基变换 $i'_{s'} : U'_{s'} \to \mathbf{A}^n_{Y'_{s'}}$
在 $x'$ 的一个开邻域中是 $H_1$-正则浸入，见
除子章引理 \ref{divisors-lemma-relative-regular-immersion}
及 除子章定义 \ref{divisors-definition-relative-H1-regular-immersion}
之后的讨论。由于 $s' = \Spec(\kappa(s')) \to \Spec(\kappa(s)) = s$
是满射、平坦且普遍开的态射（见
态射章引理 \ref{morphisms-lemma-scheme-over-field-universally-open})
），可知基变换 $i_s : U_s \to \mathbf{A}^n_{Y_s}$ 在 $x$
的一个邻域中是 $H_1$-正则浸入，见
下降章引理 \ref{descent-lemma-descending-property-regular-immersion}.
最后，注意 $\mathbf{A}^n_Y$ 在 $S$ 上平坦且局部有限表现，故
除子章引理 \ref{divisors-lemma-fibre-quasi-regular}
蕴含 $i$ 在 $x$ 的一个邻域中如所需是（Koszul-）正则浸入。
\end{proof}

\begin{lemma}
\label{more-morphisms-lemma-unramified-lci}
设 $f : X \to Y$ 是局部完全交概形态射。则 $f$ 非分歧，
当且仅当 $f$ 形式非分歧；在这种情形下，余法层
$\mathcal{C}_{X/Y}$ 在 $X$ 上有限局部自由。
\end{lemma}

\begin{proof}
第一个断言直接由
引理 \ref{more-morphisms-lemma-unramified-formally-unramified}
及局部完全交态射局部有限型这一事实得到。为计算 $f$ 的余法层，
在 $X$ 上局部选取 $f$ 的分解 $f = p \circ i$，其中
$i : X \to V$ 是 Koszul-正则浸入，而 $V \to Y$ 光滑。由
引理 \ref{more-morphisms-lemma-two-unramified-morphisms-formally-smooth}
可知 $\mathcal{C}_{X/Y}$ 是 $\mathcal{C}_{X/V}$ 的局部直和项。
由于 $i$ 是 Koszul-正则（因而拟正则）浸入，后者有限局部自由，见
除子章引理 \ref{divisors-lemma-quasi-regular-immersion}.
\end{proof}

\begin{lemma}
\label{more-morphisms-lemma-transitivity-conormal-lci}
设 $Z \to Y \to X$ 是形式非分歧的概形态射。
假设 $Z \to Y$ 是局部完全交态射。则
$$
0 \to i^*\mathcal{C}_{Y/X} \to
\mathcal{C}_{Z/X} \to
\mathcal{C}_{Z/Y} \to 0
$$
这一来自
引理 \ref{more-morphisms-lemma-transitivity-conormal}
的正合列是短正合列。
\end{lemma}

\begin{proof}
问题在 $Z$ 上是局部的，故可假设态射 $Z \to Y$ 存在分解
$Z \to \mathbf{A}^n_Y \to Y$。于是得到交换图
$$
\xymatrix{
Z \ar[r]_{i'} \ar@{=}[d] &
\mathbf{A}^n_Y \ar[r] \ar[d] &
\mathbf{A}^n_X \ar[d] \\
Z \ar[r]^i & Y \ar[r] & X
}
$$
由于 $Z \to Y$ 是局部完全交态射，$Z \to \mathbf{A}^n_Y$
是 Koszul-正则浸入。因此由
除子章引理 \ref{divisors-lemma-transitivity-conormal-quasi-regular}
，正合列
$$
0 \to (i')^*\mathcal{C}_{\mathbf{A}^n_Y/\mathbf{A}^n_X} \to
\mathcal{C}_{Z/\mathbf{A}^n_X} \to
\mathcal{C}_{Z/\mathbf{A}^n_Y} \to 0
$$
是局部分裂的。注意，有
$i^*\mathcal{C}_{Y/X} = (i')^*\mathcal{C}_{\mathbf{A}^n_Y/\mathbf{A}^n_X}$
；这是由
引理 \ref{more-morphisms-lemma-universal-thickening-fibre-product-flat}
得到的。还要注意图
$$
\xymatrix{
(i')^*\mathcal{C}_{\mathbf{A}^n_Y/\mathbf{A}^n_X} \ar[r] &
\mathcal{C}_{Z/\mathbf{A}^n_X} \\
i^*\mathcal{C}_{Y/X} \ar[u]^{\cong} \ar[r] & \mathcal{C}_{Z/X} \ar[u]
}
$$
交换。因此下方的水平箭头是局部分裂的单射。引理得证。
\end{proof}













\section{微分层与余法层的正合列}
\label{more-morphisms-section-exact}

\noindent
本节汇集余法层与微分层的正合列的一些结果。在某种意义下，
它们都是与一对可复合概形态射相伴的余切复形三角的具体实现。

\medskip\noindent
设 $g : Z \to Y$ 和 $f : Y \to X$ 是概形态射。
\begin{enumerate}
\item 有典范正合列
$$
g^*\Omega_{Y/X} \to \Omega_{Z/X} \to \Omega_{Z/Y} \to 0,
$$
，见
态射章引理 \ref{morphisms-lemma-triangle-differentials}.
若 $g : Z \to Y$ 光滑，或更一般地形式光滑，则此列是短正合列，见
态射章引理 \ref{morphisms-lemma-triangle-differentials-smooth}
或 引理 \ref{more-morphisms-lemma-triangle-differentials-formally-smooth}。
\item 若 $g$ 是浸入，或更一般地形式非分歧，则有典范正合列
$$
\mathcal{C}_{Z/Y} \to g^*\Omega_{Y/X} \to \Omega_{Z/X} \to 0,
$$
，见 态射章引理 \ref{morphisms-lemma-differentials-relative-immersion}
或 引理 \ref{more-morphisms-lemma-universally-unramified-differentials-sequence}。
若 $f \circ g : Z \to X$ 光滑，或更一般地形式光滑，
则此列是短正合列，见
态射章引理 \ref{morphisms-lemma-differentials-relative-immersion-smooth}
或 引理 \ref{more-morphisms-lemma-differentials-formally-unramified-formally-smooth}。
\item 若 $g$ 与 $f \circ g$ 是浸入，或更一般地形式非分歧，
则有典范正合列
$$
\mathcal{C}_{Z/X} \to \mathcal{C}_{Z/Y} \to g^*\Omega_{Y/X} \to 0,
$$
，见 态射章引理 \ref{morphisms-lemma-two-immersions} 或
引理 \ref{more-morphisms-lemma-two-unramified-morphisms}。若 $f : Y \to X$ 光滑，
或更一般地形式光滑，则此列是短正合列，见
态射章引理 \ref{morphisms-lemma-two-immersions-smooth} 或
引理 \ref{more-morphisms-lemma-two-unramified-morphisms-formally-smooth}。
\item 若 $g$ 和 $f$ 是浸入，或更一般地形式非分歧，则有典范正合列
$$
g^*\mathcal{C}_{Y/X} \to \mathcal{C}_{Z/X} \to \mathcal{C}_{Z/Y} \to 0.
$$
见 态射章引理 \ref{morphisms-lemma-transitivity-conormal} 或
引理 \ref{more-morphisms-lemma-transitivity-conormal}。
若 $g : Z \to Y$ 是正则浸入\footnote{$g$ 是 $H_1$-正则浸入即已足够。
注意，作为局部完全交态射的浸入是 Koszul-正则的。}，
或更一般地是局部完全交态射，则此列是短正合列，见
除子章引理 \ref{divisors-lemma-transitivity-conormal-quasi-regular}
或 引理 \ref{more-morphisms-lemma-transitivity-conormal-lci}。
\end{enumerate}










\section{弱 \'etale 态射}
\label{more-morphisms-section-weakly-etale}

\noindent
若 $A \to B$ 和 $B \otimes_A B \to B$ 都平坦，
则称环同态 $A \to B$ 弱 \'etale，见
代数进阶章定义 \ref{more-algebra-definition-weakly-etale}.
概形态射的对应概念如下。

\begin{definition}
\label{more-morphisms-definition-weakly-etale}
若 $X \to Y$ 与对角态射 $X \to X \times_Y X$ 都平坦，
则称概形态射 $X \to Y$ 是{\it 弱 \'etale 的}或{\it 绝对平坦的}。
\end{definition}

\noindent
\'etale 态射是弱 \'etale 的；反过来，弱 \'etale 态射确实在某种程度上
类似 \'etale 态射。例如，若 $X \to Y$ 弱 \'etale，则
$L_{X/Y} = 0$，这是由
余切复形章引理 \ref{cotangent-lemma-when-zero}.
得到的。稍后将证明一个非常精确的结果，把弱 \'etale 态射与
\'etale 态射联系起来（见
Pro-\'etale 上同调章第 \ref{proetale-section-weakly-etale} 节).
本节只讨论基础内容。

\begin{lemma}
\label{more-morphisms-lemma-check-weakly-etale-stalks}
设 $f : X \to Y$ 是概形态射。下列条件等价：
\begin{enumerate}
\item $X \to Y$ 弱 \'etale；
\item 对每个 $x \in X$，环映射
$\mathcal{O}_{Y, f(x)} \to \mathcal{O}_{X, x}$ 弱 \'etale。
\end{enumerate}
\end{lemma}

\begin{proof}
注意，在假设 (1) 和 (2) 中，态射 $f$ 都平坦。因此可以假设 $f$ 平坦。
设 $x \in X$ 在 $Y$ 中的像为 $y = f(x)$。有典范环映射
$$
\mathcal{O}_{X, x} \otimes_{\mathcal{O}_{Y, y}} \mathcal{O}_{X, x}
\longrightarrow
\mathcal{O}_{X \times_Y X, \Delta_{X/Y}(x)}
\longrightarrow
\mathcal{O}_{X, x}
$$
，其中第一个映射是局部化（因而平坦），第二个映射是满射。
条件 (1) 意味着对所有 $x$，第二个箭头平坦；条件 (2) 意味着对所有
$x$，复合平坦。因此等价性由
代数章引理 \ref{algebra-lemma-flat-localization} 的第 (2) 部分得到。
\end{proof}

\begin{lemma}
\label{more-morphisms-lemma-key}
设 $X \to Y$ 是概形态射，使得 $X \to X \times_Y X$ 平坦。
设 $\mathcal{F}$ 是 $\mathcal{O}_X$-模。若 $\mathcal{F}$
在 $Y$ 上平坦，则 $\mathcal{F}$ 在 $X$ 上平坦。
\end{lemma}

\begin{proof}
% P05-EN-0139: frozen source uses f(x) although the morphism in the lemma is unnamed.
设 $x \in X$ 在 $Y$ 中的像为 $y = f(x)$。
由于 $X \to X \times_Y X$ 平坦，可知
$\mathcal{O}_{X, x} \otimes_{\mathcal{O}_{Y, y}} \mathcal{O}_{X, x} \to
\mathcal{O}_{X, x}$ 平坦。因此结论由
代数进阶章引理 \ref{more-algebra-lemma-key}
及定义得到。
\end{proof}

\begin{lemma}
\label{more-morphisms-lemma-weakly-etale-characterize}
设 $f : X \to S$ 是概形态射。下列条件等价：
\begin{enumerate}
\item 态射 $f$ 弱 \'etale。
\item 对每一对满足 $f(U) \subset V$ 的仿射开集 $U \subset X$、
$V \subset S$，环映射 $\mathcal{O}_S(V) \to \mathcal{O}_X(U)$ 弱 \'etale。
\item 存在开覆盖 $S = \bigcup_{j \in J} V_j$，以及开覆盖
$f^{-1}(V_j) = \bigcup_{i \in I_j} U_i$，使得对
$j\in J, i\in I_j$，每个态射 $U_i \to V_j$ 都弱 \'etale。
\item 存在仿射开覆盖 $S = \bigcup_{j \in J} V_j$，以及仿射开覆盖
$f^{-1}(V_j) = \bigcup_{i \in I_j} U_i$，使得对所有
$j\in J, i\in I_j$，环映射
$\mathcal{O}_S(V_j) \to \mathcal{O}_X(U_i)$ 都弱 \'etale。
\end{enumerate}
此外，若 $f$ 弱 \'etale，则对任意满足 $f(U) \subset V$
的开子概形 $U \subset X$、$V \subset S$，限制
$f|_U : U \to V$ 弱 \'etale。
\end{lemma}

\begin{proof}
设给定满足 $f(U) \subset V$ 的开子概形 $U \subset X$、$V \subset S$。
% P05-EN-0140: frozen source switches the base S to an undefined Y in the ambient fibre product and diagonal.
则 $U \times_V U \subset X \times_Y X$ 是开集
（概形章引理 \ref{schemes-lemma-open-fibre-product}），而
$f|_U : U \to V$ 的对角态射 $\Delta_{U/V}$ 是限制
$\Delta_{X/Y}|_U : U \to U \times_V U$。由于平坦性是概形态射的局部性质
（态射章引理 \ref{morphisms-lemma-flat-characterize}），
引理的最后一个陈述以及 (1) 与 (3) 的等价性都随之成立。
若 $X$ 和 $Y$ 仿射，则 $X \to Y$ 弱 \'etale，当且仅当
$\mathcal{O}_Y(Y) \to \mathcal{O}_X(X)$ 弱 \'etale（再次使用
态射章引理 \ref{morphisms-lemma-flat-characterize}).
因此，(1) 与 (3) 也等价于 (2) 与 (4)。
\end{proof}

\begin{lemma}
\label{more-morphisms-lemma-composition-weakly-etale}
设 $X \to Y \to Z$ 是概形态射。
\begin{enumerate}
\item 若 $X \to X \times_Y X$ 和 $Y \to Y \times_Z Y$ 平坦，
则 $X \to X \times_Z X$ 平坦。
\item 若 $X \to Y$ 和 $Y \to Z$ 弱 \'etale，则 $X \to Z$ 弱 \'etale。
\end{enumerate}
\end{lemma}

\begin{proof}
第 (1) 部分由 $X$ 在 $Z$ 上的对角态射的分解
$$
X \to X \times_Y X \to X \times_Z X
$$
、等式
$$
X \times_Y X = (X \times_Z X) \times_{(Y \times_Z Y)} Y,
$$
、平坦态射的基变换平坦这一事实，以及平坦态射的复合平坦这一事实得到
（态射章引理 \ref{morphisms-lemma-base-change-flat} 和
\ref{morphisms-lemma-composition-flat}）。
第 (2) 部分由第 (1) 部分以及刚刚使用的平坦态射复合仍平坦这一事实得到。
\end{proof}

\begin{lemma}
\label{more-morphisms-lemma-base-change-weakly-etale}
设 $X \to Y$ 和 $Y' \to Y$ 是概形态射，并设
$X' = Y' \times_Y X$ 是 $X$ 的基变换。
\begin{enumerate}
\item 若 $X \to X \times_Y X$ 平坦，则
$X' \to X' \times_{Y'} X'$ 平坦。
\item 若 $X \to Y$ 弱 \'etale，则 $X' \to Y'$ 弱 \'etale。
\end{enumerate}
\end{lemma}

\begin{proof}
假设 $X \to X \times_Y X$ 平坦。态射 $X' \to X' \times_{Y'} X'$
是 $X \to X \times_Y X$ 沿 $Y' \to Y$ 的基变换。因此由
态射章引理 \ref{morphisms-lemma-base-change-flat}，它平坦。
这证明了 (1)。第 (2) 部分由 (1) 及刚刚使用的平坦态射的基变换平坦得到。
\end{proof}

\begin{lemma}
\label{more-morphisms-lemma-go-down}
设 $X \to Y \to Z$ 是概形态射。假设 $X \to Y$ 平坦且满射，
并且 $X \to X \times_Z X$ 平坦。则 $Y \to Y \times_Z Y$ 平坦。
\end{lemma}

\begin{proof}
考虑交换图
$$
\xymatrix{
X \ar[r] \ar[d] & X \times_Z X \ar[d] \\
Y \ar[r] & Y \times_Z Y
}
$$
上方的水平箭头平坦，竖直箭头也平坦。因此 $X$ 在
$Y \times_Z Y$ 上平坦。由
态射章引理 \ref{morphisms-lemma-flat-permanence}
可知 $Y$ 在 $Y \times_Z Y$ 上平坦。
\end{proof}

\begin{lemma}
\label{more-morphisms-lemma-weakly-etale-formally-unramified}
设 $f : X \to Y$ 是弱 \'etale 概形态射。
则 $f$ 形式非分歧，即 $\Omega_{X/Y} = 0$。
\end{lemma}

\begin{proof}
回忆，由
引理 \ref{more-morphisms-lemma-formally-unramified-differentials}.
，$f$ 形式非分歧当且仅当 $\Omega_{X/Y} = 0$。通过
引理 \ref{more-morphisms-lemma-weakly-etale-characterize} 和
态射章引理 \ref{morphisms-lemma-differentials-affine}
，这归结为环的情形，即
代数进阶章引理 \ref{more-algebra-lemma-formally-unramified}.
\end{proof}

\begin{lemma}
\label{more-morphisms-lemma-when-weakly-etale}
设 $f : X \to Y$ 是概形态射。在下列每种情形中，$X \to Y$ 都弱 \'etale：
\begin{enumerate}
\item $X \to Y$ 是平坦单态射；
\item $X \to Y$ 是开浸入；
\item $X \to Y$ 平坦且非分歧；
\item $X \to Y$ 是 \'etale 态射。
\end{enumerate}
\end{lemma}

\begin{proof}
若 (1) 成立，则 $\Delta_{X/Y}$ 是同构
（概形章引理 \ref{schemes-lemma-monomorphism}），故 $f$ 当然弱 \'etale。
情形 (2) 是 (1) 的特殊情形。非分歧态射的对角态射是开浸入
（态射章引理 \ref{morphisms-lemma-diagonal-unramified-morphism}），
因而平坦。所以平坦非分歧态射弱 \'etale。
\'etale 态射平坦且非分歧
（态射章引理 \ref{morphisms-lemma-etale-smooth-unramified}），
故 (4) 由 (3) 得到。
\end{proof}

\begin{lemma}
\label{more-morphisms-lemma-reduced-goes-up}
设 $f : X \to Y$ 是概形态射。若 $Y$ 既约且 $f$ 弱 \'etale，
则 $X$ 既约。
\end{lemma}

\begin{proof}
通过 引理 \ref{more-morphisms-lemma-weakly-etale-characterize}，
这归结为环的情形，即
代数进阶章引理
\ref{more-algebra-lemma-absolutely-flat-over-absolutely-flat}.
\end{proof}

\noindent
下面的引理使用了关于弱 \'etale 环映射的一个非平凡结果。

\begin{lemma}
\label{more-morphisms-lemma-weakly-etale-strictly-henselian-local-rings}
设 $f : X \to Y$ 是概形态射。下列条件等价：
\begin{enumerate}
\item $f$ 弱 \'etale；
\item 对 $x \in X$，局部环映射
$\mathcal{O}_{Y, f(x)} \to \mathcal{O}_{X, x}$
在严格 Hensel 化上诱导同构。
\end{enumerate}
\end{lemma}

\begin{proof}
设 $x \in X$ 在 $Y$ 中的像为 $y = f(x)$。
选取 $\kappa(x)$ 的可分代数闭包 $\kappa^{sep}$。
设 $\mathcal{O}_{X, x}^{sh}$ 是对应于 $\kappa^{sep}$ 的严格 Hensel 化，
而 $\mathcal{O}_{Y, y}^{sh}$ 是相对于 $\kappa(y)$ 在
$\kappa^{sep}$ 中的可分代数闭包的严格 Hensel 化。考虑交换图
$$
\xymatrix{
\mathcal{O}_{X, x} \ar[r] & \mathcal{O}_{X, x}^{sh} \\
\mathcal{O}_{Y, y} \ar[u] \ar[r] & \mathcal{O}_{Y, y}^{sh} \ar[u]
}
$$
，其中各映射为局部环的局部同态，见
代数章引理 \ref{algebra-lemma-strictly-henselian-functorial}.
由于严格 Hensel 化是 \'etale 环映射的滤过余极限，
代数进阶章引理 \ref{more-algebra-lemma-when-weakly-etale}
表明水平映射弱 \'etale。此外，由
代数进阶章引理 \ref{more-algebra-lemma-dumb-properties-henselization}.
，水平映射忠实平坦。

\medskip\noindent
假设 $f$ 弱 \'etale。由 引理 \ref{more-morphisms-lemma-check-weakly-etale-stalks}，
左侧竖直箭头弱 \'etale。由
代数进阶章引理 \ref{more-algebra-lemma-composition-weakly-etale} 及
\ref{more-algebra-lemma-weakly-etale-permanence}
可知右侧竖直箭头弱 \'etale。由
代数进阶章定理 \ref{more-algebra-theorem-olivier}
可知右侧竖直映射是同构。

\medskip\noindent
假设 $\mathcal{O}_{Y, y}^{sh} \to \mathcal{O}_{X, x}^{sh}$ 是同构。
则 $\mathcal{O}_{Y, y} \to \mathcal{O}_{X, x}^{sh}$ 弱 \'etale。
由于 $\mathcal{O}_{X, x} \to \mathcal{O}_{X, x}^{sh}$ 忠实平坦，
由
代数进阶章引理 \ref{more-algebra-lemma-go-down}.
可知 $\mathcal{O}_{Y, y} \to \mathcal{O}_{X, x}$ 弱 \'etale。
因此由 引理 \ref{more-morphisms-lemma-check-weakly-etale-stalks}，(2) 蕴含 (1)。
\end{proof}

\begin{lemma}
\label{more-morphisms-lemma-normal-goes-up}
设 $f : X \to Y$ 是概形态射。若 $Y$ 是正规概形且 $f$ 弱 \'etale，
则 $X$ 是正规概形。
\end{lemma}

\begin{proof}
由 代数进阶章引理 \ref{more-algebra-lemma-henselization-normal}，
概形 $S$ 正规，当且仅当对所有 $s \in S$，
$\mathcal{O}_{S, s}$ 的严格 Hensel 化都是正规整环。因此引理由
引理 \ref{more-morphisms-lemma-weakly-etale-strictly-henselian-local-rings}.
得到。
\end{proof}

\begin{lemma}
\label{more-morphisms-lemma-weakly-etale-permanence}
设 $S$ 是概形，$f : X \to Y$ 是 $S$ 上的概形态射。
若 $X$、$Y$ 在 $S$ 上弱 \'etale，则 $f$ 弱 \'etale。
\end{lemma}

\begin{proof}
我们将不再另行说明地使用 态射章引理
\ref{morphisms-lemma-base-change-flat} 和
\ref{morphisms-lemma-composition-flat}。把 $X \to Y$ 写成复合
$X \to X \times_S Y \to Y$。第二个态射是平坦态射 $X \to S$
的基变换，因而平坦。第一个态射是平坦态射
$Y \to Y \times_S Y$ 沿态射 $X \times_S Y \to Y \times_S Y$
的基变换，因而也平坦。所以 $X \to Y$ 平坦。
态射 $X \times_Y X \to X \times_S X$ 是浸入。
因此 引理 \ref{more-morphisms-lemma-key} 蕴含：由于 $X$ 在 $X \times_S X$
上平坦，$X$ 在 $X \times_Y X$ 上也平坦。
\end{proof}

\noindent
下面是这样一个观察的概形论推广：同时可分且纯不可分的域扩张必为同构。

\begin{lemma}
\label{more-morphisms-lemma-weakly-etale-universal-homeomorphism}
设 $f : X \to Y$ 是概形态射。若 $f$ 弱 \'etale 且为普遍同胚，
则它是同构。
\end{lemma}

\begin{proof}
因为 $f$ 是普遍同胚，由
态射章引理 \ref{morphisms-lemma-homeomorphism-affine} 及
\ref{morphisms-lemma-universally-injective}，对角态射
$\Delta : X \to X \times_Y X$ 是满射闭浸入。由于 $\Delta$ 也平坦，
由
态射章引理 \ref{morphisms-lemma-characterize-flat-closed-immersions}.
可知 $\Delta$ 必为同构。换言之，$f$ 是单态射
（概形章引理 \ref{schemes-lemma-monomorphism}）。
由于 $f$ 是普遍同胚，它当然拟紧。因此由 下降章引理
\ref{descent-lemma-flat-surjective-quasi-compact-monomorphism-isomorphism}
可知 $f$ 是同构。
\end{proof}

\noindent
下面是
\'Etale 态射章引理 \ref{etale-lemma-relative-frobenius-etale}.
的弱 \'etale 推广。

\begin{lemma}
\label{more-morphisms-lemma-relative-frobenius-weakly-etale}
设 $U \to X$ 是弱 \'etale 概形态射，其中 $X$ 是特征 $p$ 的概形。
则相对 Frobenius 态射 $F_{U/X} : U \to U \times_{X, F_X} X$ 是同构。
\end{lemma}

\begin{proof}
由 簇章引理 \ref{varieties-lemma-relative-frobenius}，
$F_{U/X}$ 是普遍同胚。由于它是 $X$ 上弱 \'etale 概形之间的态射，
由 引理 \ref{more-morphisms-lemma-weakly-etale-permanence}
可知 $F_{U/X}$ 弱 \'etale。因此由
引理 \ref{more-morphisms-lemma-weakly-etale-universal-homeomorphism}，$F_{U/X}$ 是同构。
\end{proof}






\section{约化纤维定理}
\label{more-morphisms-section-reduced-fibre-theorem}

\noindent
本节讨论题名所预告的此类定理中最简单的一种。虽然该结果的证明
稍显繁复，但从本质上说，它直接由 Epp 关于消除分歧的结果推出，见
代数进阶，定理 \ref{more-algebra-theorem-epp}。

\medskip\noindent
设 $A$ 是分式域为 $K$ 的戴德金整环。
设 $X$ 是 $A$ 上平坦且有限型的概形。
设 $L$ 是 $K$ 的有限扩张，并设 $B$ 是 $A$ 在 $L$ 中的整闭包。
于是 $B$ 是戴德金整环
（代数，引理 \ref{algebra-lemma-integral-closure-Dedekind}）。
令 $X_B = X \times_{\Spec(A)} \Spec(B)$ 为相应的基变换。
则 $X_B \to \Spec(B)$ 是有限型的
（态射，引理 \ref{morphisms-lemma-base-change-finite-type}）。
因此 $X_B$ 是诺特概形
（态射，引理 \ref{morphisms-lemma-finite-type-noetherian}）。
从而正规化 $\nu : Y \to X_B$ 存在（见
态射，定义 \ref{morphisms-definition-normalization}
及其后的讨论）。图示为
\begin{equation}
\label{more-morphisms-equation-normalized-base-change}
\xymatrix{
Y \ar[rd] \ar[r]_\nu & X_B \ar[r] \ar[d] & X \ar[d] \\
 & \Spec(B) \ar[r] & \Spec(A)
}
\end{equation}
我们有时称 $Y$ 为 $X$ 的{\it 正规化基变换}。
一般而言，态射 $\nu$ 未必有限。但若 $A$ 是纳伽塔环（这是一项在实践中
几乎总能满足的条件），则 $\nu$ 有限，且 $Y$ 在 $B$ 上有限型，见
态射，引理 \ref{morphisms-lemma-nagata-normalization} 和
\ref{morphisms-lemma-finite-type-nagata}。

\medskip\noindent
取正规化基变换与复合相容。更确切地说，若 $M/L/K$ 是域的有限扩张，
相应的整闭包为 $A \subset B \subset C$，则 $Y \to \Spec(B)$
关于 $M/L$ 的正规化基变换 $Z$，等于 $X \to \Spec(A)$
关于 $M/K$ 的正规化基变换。

\begin{theorem}
\label{more-morphisms-theorem-normalized-base-change-with-reduced-fibre}
设 $A$ 是分式域为 $K$ 的戴德金环。
设 $X$ 是 $A$ 上平坦且有限型的概形，并假设 $A$ 是纳伽塔环。
存在有限扩张 $L/K$，使得正规化基变换 $Y$ 在所有纤维的全部泛点处
都在 $\Spec(B)$ 上光滑。
\end{theorem}

\begin{proof}
在证明中，我们将反复使用如下事实：对于 $X \to \Spec(A)$ 这样的
（平坦、有限表示）态射，其光滑点集的形成与基变换相容，见
态射，引理 \ref{morphisms-lemma-set-points-where-fibres-smooth}。

\medskip\noindent
先选择有限扩张 $L/K$，使得 $(X_L)_{red}$ 在 $L$ 上几何约化，见
簇，引理 \ref{varieties-lemma-finite-extension-geometrically-reduced}。
由于 $Y \to (X_B)_{red}$ 是双有理的，应用
簇，引理 \ref{varieties-lemma-generic-points-geometrically-reduced}
可知 $Y_L$ 也在 $L$ 上几何约化。因此由
簇，引理 \ref{varieties-lemma-geometrically-reduced-dense-smooth-open}，
$Y_L \to \Spec(L)$ 在某个稠密开子集 $V \subset Y_L$ 上光滑。
于是态射 $Y \to \Spec(B)$ 的光滑轨迹 $U \subset Y$ 是开集
（由态射，定义 \ref{morphisms-definition-smooth}），
并且在泛纤维中稠密。以 $B$ 代替 $A$、以 $Y$ 代替 $X$，
便归约到下一段所处理的情形。

\medskip\noindent
假设 $X$ 正规，且 $X \to \Spec(A)$ 的光滑轨迹 $U \subset X$
在泛纤维中稠密。这蕴含 $U$ 在除有限多个之外的所有纤维中都稠密，
见 引理 \ref{more-morphisms-lemma-nowhere-dense-generic-fibre}。
令 $x_1, \ldots, x_r \in X \setminus U$ 为如下有限多个点：它们既是
$X \setminus U$ 的不可约分支的泛点，又是 $X \to \Spec(A)$ 的纤维的
不可约分支的泛点。置 $\mathcal{O}_i = \mathcal{O}_{X, x_i}$。
令 $A_i$ 为 $A$ 在如下极大理想处的局部化：该极大理想对应于 $x_i$
在 $\Spec(A)$ 中的像。由
代数进阶，Proposition
\ref{more-algebra-proposition-epp-essentially-finite-type}，
存在有限扩张 $K_i/K$，它们是离散赋值环扩张
$A_i \to \mathcal{O}_i$ 的解。取一个支配所有扩张 $K_i/K$ 的有限扩张
$L/K$。由代数进阶，引理 \ref{more-algebra-lemma-solution-goes-up}，
$L/K$ 仍是 $A_i \to \mathcal{O}_i$ 的解。

\medskip\noindent
考虑刚构造的扩张 $L/K$ 所对应的图
(\ref{more-morphisms-equation-normalized-base-change})。注意 $U_B \subset X_B$ 在 $B$
上光滑，因而正规（例如使用代数，Lemma
\ref{algebra-lemma-normal-goes-up}）。故 $Y \to X_B$ 在 $U_B$ 上是同构。
设 $y \in Y$ 是 $Y \to \Spec(B)$ 的某个纤维的一个不可约分支的泛点，
该纤维位于极大理想 $\mathfrak m \subset B$ 之上。假设
$y \not \in U_B$。则 $y$ 映到某个点 $x_i$。由此可知
$\mathcal{O}_{Y, y}$ 是 $\mathcal{O}_i$ 在 $R(X) \otimes_K L$
中的整闭包的一个局部环（略去细节）。由于 $L/K$ 是
$A_i \to \mathcal{O}_i$ 的解，故
$B_\mathfrak m \to \mathcal{O}_{Y, y}$ 对 $\mathfrak m_y$-进拓扑形式光滑
（这正是成为“解”的定义）。换言之，
$\mathfrak m\mathcal{O}_{Y, y} = \mathfrak m_y$，并且剩余域扩张可分，见
代数进阶，引理 \ref{more-algebra-lemma-extension-dvrs-formally-smooth}。
因此纤维在 $y$ 处的局部环是 $\kappa(y)$。例如由代数，Lemma
\ref{algebra-lemma-separable-smooth}，这蕴含该纤维在 $y$ 处于
$\kappa(\mathfrak m)$ 上光滑。证明完成。
\end{proof}

\begin{lemma}[曲线上的变体]
\label{more-morphisms-lemma-normalized-base-change-with-reduced-fibre-over-curve}
设 $f : X \to S$ 是平坦有限型的概形态射。假设 $S$ 是纳伽塔整概形，
其函数域为 $K$，并且是维数为 $1$ 的正则概形。
则存在有限扩张 $L/K$，使得在图
$$
\xymatrix{
Y \ar[rd]_g \ar[r]_-\nu & X \times_S T \ar[d] \ar[r] & X \ar[d]_f \\
& T \ar[r] & S
}
$$
中，态射 $g$ 在纤维的所有泛点处光滑。这里，$T$ 是 $S$ 在
$\Spec(L)$ 中的正规化，而 $\nu : Y \to X \times_S T$ 是正规化。
\end{lemma}

\begin{proof}
选择有限仿射开覆盖 $S = \bigcup \Spec(A_i)$。于是对每个 $i$，
$K$ 都等于 $A_i$ 的分式域。令 $X_i = X \times_S \Spec(A_i)$。
对态射 $X_i \to \Spec(A_i)$，按照
定理 \ref{more-morphisms-theorem-normalized-base-change-with-reduced-fibre}
选择 $L_i/K$。令 $B_i \subset L_i$ 为 $A_i$ 的整闭包，并令 $Y_i$
为 $X$ 到 $B_i$ 的正规化基变换。取支配每个 $L_i$ 的有限扩张
$L/K$。令 $T_i \subset T$ 为 $\Spec(A_i)$ 的原像。对每个 $i$，
得到交换图
$$
\xymatrix{
g^{-1}(T_i) \ar[r] \ar[d] & Y_i \ar[r] \ar[d] & X \times_S \Spec(A_i) \ar[d] \\
T_i \ar[r] & \Spec(B_i) \ar[r] & \Spec(A_i)
}
$$
而且左侧方块正是本节开头所讨论的正规化基变换。在
定理 \ref{more-morphisms-theorem-normalized-base-change-with-reduced-fibre}
的证明中，我们已经看到，$Y \to T$ 的光滑轨迹包含如下点集在
$g^{-1}(T_i)$ 中的原像：$Y_i$ 在这些点处于 $B_i$ 上光滑。
这就证明了引理。
\end{proof}

\begin{lemma}[可分扩张情形的变体]
\label{more-morphisms-lemma-normalized-base-change-with-reduced-fibre-separable}
设 $A$ 是分式域为 $K$ 的戴德金环。设 $X$ 是 $A$ 上平坦且有限型的概形。
假设 $A$ 是纳伽塔环，并且对 $X$ 的每个不可约分支的每个泛点
$\eta$，域扩张 $\kappa(\eta)/K$ 都可分。则存在有限可分扩张 $L/K$，
使得正规化基变换 $Y$ 在所有纤维的全部泛点处都在 $\Spec(B)$ 上光滑。
\end{lemma}

\begin{proof}
证明方法与 定理 \ref{more-morphisms-theorem-normalized-base-change-with-reduced-fibre}
完全相同，只需作几处小修改。最重要的改动是使用代数进阶，Lemma
\ref{more-algebra-lemma-epp-essentially-finite-type-separable}，
而不是代数进阶，Proposition
\ref{more-algebra-proposition-epp-essentially-finite-type}。
在证明中，我们将反复使用如下事实：对于 $X \to \Spec(A)$ 这样的
（平坦、有限表示）态射，其光滑点集的形成与基变换相容，见
态射，引理 \ref{morphisms-lemma-set-points-where-fibres-smooth}。

\medskip\noindent
由于 $X$ 在 $A$ 上平坦，$X$ 的每个泛点 $\eta$ 都映到 $\Spec(A)$
的泛点。以 $X$ 的约化概形代替它后，可以假设 $X$ 约化。
在这种情形下，由簇，Lemma
\ref{varieties-lemma-generic-points-geometrically-reduced}，$X_K$ 在 $K$
上几何约化。因此由簇，Lemma
\ref{varieties-lemma-geometrically-reduced-dense-smooth-open}，
$X_K \to \Spec(K)$ 在某个稠密开集上光滑。于是态射
$X \to \Spec(A)$ 的光滑轨迹 $U \subset X$ 是开集
（由态射，定义 \ref{morphisms-definition-smooth}），
并且在泛纤维中稠密。这就把问题归约到下一段的情形。

\medskip\noindent
假设 $X$ 正规，且 $X \to \Spec(A)$ 的光滑轨迹 $U \subset X$
在泛纤维中稠密。这蕴含 $U$ 在除有限多个之外的所有纤维中都稠密，
见 引理 \ref{more-morphisms-lemma-nowhere-dense-generic-fibre}。
令 $x_1, \ldots, x_r \in X \setminus U$ 为如下有限多个点：它们既是
$X \setminus U$ 的不可约分支的泛点，又是 $X \to \Spec(A)$ 的纤维的
不可约分支的泛点。置 $\mathcal{O}_i = \mathcal{O}_{X, x_i}$。
注意，$\mathcal{O}_i$ 的分式域是 $X$ 的某个泛点的剩余域。
令 $A_i$ 为 $A$ 在如下极大理想处的局部化：该极大理想对应于 $x_i$
在 $\Spec(A)$ 中的像。我们可以应用代数进阶，Lemma
\ref{more-algebra-lemma-epp-essentially-finite-type-separable}，
得到有限可分扩张 $K_i/K$，它们是 $A_i \to \mathcal{O}_i$ 的解。
取一个支配所有扩张 $K_i/K$ 的有限可分扩张 $L/K$。由代数进阶，
引理 \ref{more-algebra-lemma-solution-goes-up}，$L/K$ 仍是
$A_i \to \mathcal{O}_i$ 的解。

\medskip\noindent
考虑刚构造的扩张 $L/K$ 所对应的图
(\ref{more-morphisms-equation-normalized-base-change})。注意 $U_B \subset X_B$ 在 $B$
上光滑，因而正规（例如使用代数，Lemma
\ref{algebra-lemma-normal-goes-up}）。故 $Y \to X_B$ 在 $U_B$ 上是同构。
设 $y \in Y$ 是 $Y \to \Spec(B)$ 的某个纤维的一个不可约分支的泛点，
该纤维位于极大理想 $\mathfrak m \subset B$ 之上。假设
$y \not \in U_B$。则 $y$ 映到某个点 $x_i$。由此可知
$\mathcal{O}_{Y, y}$ 是 $\mathcal{O}_i$ 在 $R(X) \otimes_K L$
中的整闭包的一个局部环（略去细节）。由于 $L/K$ 是
$A_i \to \mathcal{O}_i$ 的解，故
$B_\mathfrak m \to \mathcal{O}_{Y, y}$ 形式光滑
（这正是成为“解”的定义）。换言之，
$\mathfrak m\mathcal{O}_{Y, y} = \mathfrak m_y$，并且剩余域扩张可分。
因此纤维在 $y$ 处的局部环是 $\kappa(y)$。例如由代数，Lemma
\ref{algebra-lemma-separable-smooth}，这蕴含该纤维在 $y$ 处于
$\kappa(\mathfrak m)$ 上光滑。证明完成。
\end{proof}

\begin{lemma}[曲线上的可分扩张变体]
\label{more-morphisms-lemma-normalized-base-change-with-reduced-fibre-over-curve-separable}
设 $f : X \to S$ 是平坦有限型的概形态射。假设 $S$ 是纳伽塔整概形，
其函数域为 $K$，并且是维数为 $1$ 的正则概形。假设对 $X$
的每个不可约分支的每个泛点 $\eta$，域扩张 $\kappa(\eta)/K$
都可分。则存在有限可分扩张 $L/K$，使得在图
$$
\xymatrix{
Y \ar[rd]_g \ar[r]_-\nu & X \times_S T \ar[d] \ar[r] & X \ar[d]_f \\
& T \ar[r] & S
}
$$
中，态射 $g$ 在纤维的所有泛点处光滑。这里，$T$ 是 $S$ 在
$\Spec(L)$ 中的正规化，而 $\nu : Y \to X \times_S T$ 是正规化。
\end{lemma}

\begin{proof}
这由 Lemma
\ref{more-morphisms-lemma-normalized-base-change-with-reduced-fibre-separable}
推出，其方式与 Lemma
\ref{more-morphisms-lemma-normalized-base-change-with-reduced-fibre-over-curve}
由 定理 \ref{more-morphisms-theorem-normalized-base-change-with-reduced-fibre}
推出的方式完全相同。
\end{proof}








\section{Ind-拟仿射态射}
\label{more-morphisms-section-ind-quasi-affine}

\noindent
下面给出稍后要用到的一点理论。

\begin{definition}
\label{more-morphisms-definition-ind-quasi-affine}
如果 $X$ 的每个拟紧开子概形都是拟仿射的，则称概形 $X$ 是
{\it ind-拟仿射的}。类似地，如果对 $Y$ 中的每个仿射开集 $V$，
$f^{-1}(V)$ 都是 ind-拟仿射的，则称概形态射 $X \to Y$ 是
{\it ind-拟仿射的}。
\end{definition}

\noindent
仿射概形的任一开子概形都是 ind-拟仿射概形。若
$X = \bigcup_{i \in I} U_i$ 是拟仿射开集的并，且任意两个 $U_i$
都包含在第三个这样的开集中，则 $X$ 是 ind-拟仿射的。
Ind-拟仿射概形 $X$ 是分离的，因为任意两个仿射开集 $U, V$
都包含在 $X$ 的一个分离开子概形中，即 $U \cup V$ 中。
类似地，ind-拟仿射态射是分离的。

\begin{lemma}
\label{more-morphisms-lemma-ind-quasi-affine-alternative-definition}
对概形态射 $f : X \to Y$，下列条件等价：
\begin{enumerate}
\item $f$ 是 ind-拟仿射的；
\item 对每个仿射开子概形 $V \subset Y$ 以及每个拟紧开子概形
$U \subset f^{-1}(V)$，所诱导的态射 $U \to V$ 是拟仿射的；
\item
存在 $Y$ 的一个由拟紧拟分离开子概形 $V_j \subset Y$ 组成的覆盖
$\{ V_j \}_{j \in J}$，使得对每个 $j \in J$ 以及每个拟紧开子概形
$U \subset f^{-1}(V_j)$，所诱导的态射 $U \to V_j$ 是拟仿射的；
\item 对每个拟紧拟分离开子概形 $V \subset Y$ 以及每个拟紧开子概形
$U \subset f^{-1}(V)$，所诱导的态射 $U \to V$ 是拟仿射的。
\end{enumerate}
特别地，成为 ind-拟仿射态射这一性质在基底上是 Zariski 局部的。
\end{lemma}

\begin{proof}
(1) $\Leftrightarrow$ (2) 的等价性由定义和
态射，引理 \ref{morphisms-lemma-characterize-quasi-affine} 得出。
为证 (2) $\Rightarrow$ (4)，令 $U$ 和 $V$ 如 (4) 所述。由
概形，引理 \ref{schemes-lemma-quasi-compact-permanence}，所诱导的态射
$U \to V$ 是拟紧的。因此对每个仿射开集 $V' \subset V$，纤维积
$V' \times_V U$ 是拟紧的，所以由 (2)，所诱导的映射
$V' \times_V U \to V'$ 是拟仿射的。于是由
态射，引理 \ref{morphisms-lemma-characterize-quasi-affine}，
$U \to V$ 也是拟仿射的。
这一论证也给出 (3) $\Rightarrow$ (4)：沿用相同记号，包含在某个
$V_j$ 中的那些仿射开集 $V' \subset V$ 覆盖 $V$，因此只需证明拟紧映射
$V' \times_V U \to V'$ 是拟仿射的。然而由 (3)，复合
$V' \times_V U \to V' \to V_j$ 是拟仿射的；而由
概形，引理 \ref{schemes-lemma-compose-after-separated}，映射
$V' \to V_j$ 是拟分离的。因此由
态射，引理 \ref{morphisms-lemma-quasi-affine-permanence}，
$V' \times_V U \to V'$ 是拟仿射的。最后的蕴含
(4) $\Rightarrow$ (2) 和 (4) $\Rightarrow$ (3) 是显然的。
\end{proof}

\begin{lemma}
\label{more-morphisms-lemma-ind-quasi-affine-composition}
成为 ind-拟仿射态射这一性质在复合下稳定。
\end{lemma}

\begin{proof}
设 $f : X \to Y$ 和 $g : Y \to Z$ 是 ind-拟仿射态射。设
$V \subset Z$ 和 $U \subset f^{-1}(g^{-1}(V))$ 是拟紧开集，并且
$V$ 还是拟分离的。像 $f(U)$ 是 $g^{-1}(V)$ 的拟紧子集，故包含在某个
拟紧开集 $W \subset g^{-1}(V)$ 中（它是有限多个仿射开集的并）。
我们得到分解 $U \to W \to V$。由
引理 \ref{more-morphisms-lemma-ind-quasi-affine-alternative-definition}，映射
$W \to V$ 是拟仿射的，特别地，$W$ 是拟分离的。再次应用
引理 \ref{more-morphisms-lemma-ind-quasi-affine-alternative-definition}，可知
$U \to W$ 也是拟仿射的。因此由态射，Lemma
\ref{morphisms-lemma-composition-quasi-affine}，复合 $U \to V$
也是拟仿射的；最后再应用一次
引理 \ref{more-morphisms-lemma-ind-quasi-affine-alternative-definition} 即可。
\end{proof}

\begin{lemma}
\label{more-morphisms-lemma-ind-quasi-affine-examples}
任一拟仿射态射都是 ind-拟仿射的。
任一浸入都是 ind-拟仿射的。
\end{lemma}

\begin{proof}
第一项断言直接由定义得出。特别地，仿射态射（例如闭浸入）
是 ind-拟仿射的。因此由 Lemma
\ref{more-morphisms-lemma-ind-quasi-affine-composition}，只需证明开浸入是
ind-拟仿射的，而这也直接由定义得出。
\end{proof}

\begin{lemma}
\label{more-morphisms-lemma-ind-quasi-affine-permanence}
设 $f : X \to Y$ 和 $g : Y \to Z$ 是概形态射。若 $g \circ f$
是 ind-拟仿射的，则 $f$ 是 ind-拟仿射的。
\end{lemma}

\begin{proof}
由 引理 \ref{more-morphisms-lemma-ind-quasi-affine-alternative-definition}，可以先在
$Z$ 上、再在 $Y$ 上 Zariski 局部地工作，故不妨假设 $Z$ 仿射，
进而也假设 $Y$ 仿射。此时 $X$ 的任一拟紧开集都是拟仿射的，
故结论由 引理 \ref{more-morphisms-lemma-ind-quasi-affine-alternative-definition} 得出。
\end{proof}

\begin{lemma}
\label{more-morphisms-lemma-base-change-ind-quasi-affine}
成为 ind-拟仿射态射这一性质在基变换下稳定。
\end{lemma}

\begin{proof}
设 $f : X \to Y$ 是 ind-拟仿射态射。为检验 $f$ 的每个基变换都是
ind-拟仿射的，由 Lemma
\ref{more-morphisms-lemma-ind-quasi-affine-alternative-definition}，可以在 $Y$ 上
Zariski 局部地工作，故假设 $Y$ 仿射。进一步，也可以假设基变换态射
$Z \to Y$ 中的概形 $Z$ 仿射。基变换 $X \times_Y Z \to X$ 是仿射态射，
故由 引理 \ref{more-morphisms-lemma-ind-quasi-affine-composition} 和
\ref{more-morphisms-lemma-ind-quasi-affine-examples}，映射 $X \times_Y Z \to Y$
是 ind-拟仿射的。再由 引理 \ref{more-morphisms-lemma-ind-quasi-affine-permanence}，
基变换 $X \times_Y Z \to Z$ 是 ind-拟仿射的，正合所需。
\end{proof}

\begin{lemma}
\label{more-morphisms-lemma-descending-property-ind-quasi-affine}
成为 ind-拟仿射态射这一性质在基底上是 fpqc 局部的。
\end{lemma}

\begin{proof}
引理 \ref{more-morphisms-lemma-base-change-ind-quasi-affine} 给出的 ind-拟仿射性在
基变换下的稳定性证明了一个方向。反过来，设 $f : X \to Y$ 是概形态射，
$\{g_i : Y_i \to Y\}$ 是一个 fpqc 覆盖，并且对每个 $i$，基变换
$f_i : X_i \to Y_i$ 都是 ind-拟仿射的。我们要证明 $f$ 是
ind-拟仿射的。

\medskip\noindent
由 引理 \ref{more-morphisms-lemma-ind-quasi-affine-alternative-definition}，可以在 $Y$
上 Zariski 局部地工作，故假设 $Y$ 仿射。再利用
引理 \ref{more-morphisms-lemma-base-change-ind-quasi-affine} 所保证的基变换稳定性细化覆盖，
并假设该覆盖由单个仿射忠实平坦态射 $g : Y' \to Y$ 给出。
对任一拟紧开集 $U \subset X$，其 $Y'$-基变换
$U \times_Y Y' \subset X \times_Y Y'$ 也是拟紧的。最后注意，由
下降，引理 \ref{descent-lemma-descending-property-quasi-affine}，
映射 $U \to Y$ 是拟仿射的，当且仅当 $U \times_Y Y' \to Y'$ 是拟仿射的。
\end{proof}

\begin{lemma}
\label{more-morphisms-lemma-etale-separated-ind-quasi-affine}
分离局部拟有限的概形态射是 ind-拟仿射的。
\end{lemma}

\begin{proof}
设 $f : X \to Y$ 是分离局部拟有限的概形态射。设 $V \subset Y$
仿射，而 $U \subset f^{-1}(V)$ 是拟紧开集。我们要证明 $U$ 拟仿射。
由于 $U \to V$ 是分离拟有限的概形态射，结论由 Zariski 主定理得出。
见 引理 \ref{more-morphisms-lemma-quasi-finite-separated-quasi-affine}。
\end{proof}




\section{概形范畴中的推出，II}
\label{more-morphisms-section-pushouts-II}

\noindent
本节是 第 \ref{more-morphisms-section-pushouts} 节 的延续。我们将在 $Z \to X$
为闭浸入、$Z \to Y$ 为整态射且满足一项附加条件时，构造
$Y \leftarrow Z \rightarrow X$ 的推出。详细讨论见
\cite{Ferrand-Conducteur}。

\begin{situation}
\label{more-morphisms-situation-pushout-along-closed-immersion-and-integral}
这里，$S$ 是概形，$i : Z \to X$ 和 $j : Z \to Y$ 是 $S$-概形的态射。
我们假设
\begin{enumerate}
\item $i$ 是闭浸入；
\item $j$ 是整的概形态射；
\item 对 $y \in Y$，存在仿射开集 $U \subset X$，使得
$j^{-1}(\{y\}) \subset i^{-1}(U)$。
\end{enumerate}
\end{situation}

\begin{lemma}
\label{more-morphisms-lemma-prepare-pushout-along-closed-immersion-and-integral}
在 Situation \ref{more-morphisms-situation-pushout-along-closed-immersion-and-integral}
中，对 $y \in Y$，存在仿射开集 $U \subset X$ 和 $V \subset Y$，
使得 $i^{-1}(U) = j^{-1}(V)$ 且 $y \in V$。
\end{lemma}

\begin{proof}
令 $y \in Y$。选择仿射开集 $U \subset X$，使得
$j^{-1}(\{y\}) \subset i^{-1}(U)$（由假设可以做到）。选择 $y$ 的一个
仿射开邻域 $V \subset Y$，使得 $j^{-1}(V) \subset i^{-1}(U)$。
这是可能的，因为 $j : Z \to Y$ 是闭态射
（态射，引理 \ref{morphisms-lemma-integral-universally-closed}），
且 $i^{-1}(U)$ 包含 $y$ 上的纤维。由于 $j$ 是整的，概形论纤维 $Z_y$
是某个域上的整代数的谱。由极限，Lemma
\ref{limits-lemma-ample-profinite-set-in-principal-affine}，可以找到
$\overline{f} \in \Gamma(i^{-1}(U), \mathcal{O}_{i^{-1}(U)})$，使得
$Z_y \subset D(\overline{f}) \subset j^{-1}(V)$。由于
$i|_{i^{-1}(U)} : i^{-1}(U) \to U$ 是仿射概形之间的闭浸入，可以选择
$f \in \Gamma(U, \mathcal{O}_U)$，其在 $i^{-1}(U)$ 上的限制是
$\overline{f}$。以主开集 $D(f) \subset U$ 代替 $U$ 后，得到仿射开集
$y \in V \subset Y$ 和 $U \subset X$，满足
$$
j^{-1}(\{y\}) \subset i^{-1}(U) \subset j^{-1}(V)
$$
现在我们（在某种意义下）重复这一论证。具体来说，选择
$g \in \Gamma(V, \mathcal{O}_V)$，使得 $y \in D(g)$ 且
$j^{-1}(D(g)) \subset i^{-1}(U)$（与上面相同的论证保证这可以做到）。
然后可以取 $f \in \Gamma(U, \mathcal{O}_U)$，使其在 $i^{-1}(U)$
上的限制是 $g$ 沿 $i^{-1}(U) \to V$ 的拉回（仍由上述同一理由）。
于是最终得到仿射开集 $y \in V' = D(g) \subset V \subset Y$ 和
% P05-EN-0141: frozen source writes i^{-1}(V') although V' is an open of Y; the defined open of X is U'.
$U' = D(f) \subset U \subset X$，满足 $j^{-1}(V') = i^{-1}(V')$。
\end{proof}

\begin{proposition}
\label{more-morphisms-proposition-pushout-along-closed-immersion-and-integral}
\begin{reference}
\cite[Theorem 7.1 part iii]{Ferrand-Conducteur}
\end{reference}
在 Situation \ref{more-morphisms-situation-pushout-along-closed-immersion-and-integral}
中，推出 $Y \amalg_Z X$ 在概形范畴中存在。图示为
$$
\xymatrix{
Z \ar[r]_i \ar[d]_j & X \ar[d]^a \\
Y \ar[r]^-b & Y \amalg_Z X
}
$$
该图是纤维方块，态射 $a$ 是整的，态射 $b$ 是闭浸入，并且
$$
\mathcal{O}_{Y \amalg_Z X} =
b_*\mathcal{O}_Y \times_{c_*\mathcal{O}_Z} a_*\mathcal{O}_X
$$
作为环层成立，其中 $c = a \circ i = b \circ j$。
\end{proposition}

\begin{proof}
作为拓扑空间，我们把 $Y \amalg_Z X$ 定义为该图在拓扑空间范畴中的推出
（拓扑，第 \ref{topology-section-colimits} 节）。它就是底层集合的推出
（拓扑，引理 \ref{topology-lemma-colimits}）赋予商拓扑。
在 $Y \amalg_Z X$ 上，有环层的映射
$$
b_*\mathcal{O}_Y \longrightarrow c_*\mathcal{O}_Z
\longleftarrow a_*\mathcal{O}_X
$$
并且可以把
$$
\mathcal{O}_{Y \amalg_Z X} =
b_*\mathcal{O}_Y \times_{c_*\mathcal{O}_Z} a_*\mathcal{O}_X
$$
定义为环层范畴中的纤维积。为证明所得对象是概形，需要证明每个点都有
仿射开邻域。对不在 $c$ 的像中的点，这是清楚的，因为 $c$ 的像是闭子集，
其补集作为环化空间同构于
$(Y \setminus j(Z)) \amalg (X \setminus i(Z))$。

\medskip\noindent
$c$ 的像中的一个点对应于 $j$ 的像中的唯一一点 $y \in Y$。由
引理 \ref{more-morphisms-lemma-prepare-pushout-along-closed-immersion-and-integral}，
可以找到仿射开集 $U \subset X$ 和 $V \subset Y$，使得
$y \in V$ 且 $i^{-1}(U) = j^{-1}(V)$。第一段的构造显然与限制到
相容的开子概形相容，所以为证明它给出一个概形，可以假设
$X$、$Y$ 和 $Z$ 都仿射。

\medskip\noindent
若 $X = \Spec(A)$、$Y = \Spec(B)$ 和 $Z = \Spec(C)$ 都仿射，则
代数进阶，引理 \ref{more-algebra-lemma-points-of-fibre-product} 表明，
作为拓扑空间有 $Y \amalg_Z X = \Spec(B \times_C A)$。
% P05-EN-0142: frozen source switches the pushout to the undefined fibre product Y times_Z X.
为完成 $Y \times_Z X$ 是概形的证明，只需说明在
$\Spec(B \times_C A)$ 上，结构层是各直像的纤维积。这可由把
代数进阶，引理 \ref{more-algebra-lemma-diagram-localize}
应用于 $\Spec(B \times_C A)$ 的主仿射开集得到。

\medskip\noindent
上述讨论表明，概形 $Y \amalg_Z X$ 有仿射开覆盖
$Y \amalg_Z X = \bigcup W_i$，使得 $U_i = a^{-1}(W_i)$、
$V_i = b^{-1}(W_i)$ 和 $\Omega_i = c^{-1}(W_i)$ 分别是 $X$、$Y$ 和 $Z$
中的仿射开集。因此 $a$ 和 $b$ 都是仿射的。进一步，若 $A_i$、$B_i$、
$C_i$ 分别是对应于 $U_i$、$V_i$、$\Omega_i$ 的环，则
$A_i \to C_i$ 满射，而 $W_i$ 对应于 $A_i \times_{C_i} B_i$，
后者满射到 $B_i$。故 $b$ 是闭浸入。由代数进阶，Lemma
\ref{more-algebra-lemma-fibre-product-integral}，环映射
$A_i \times_{C_i} B_i \to A_i$ 是整的，故 $a$ 是整的。
该图是笛卡尔的，因为
$$
C_i \cong B_i \otimes_{B_i \times_{C_i} A_i} A_i
$$
这是因为 $B_i \times_{C_i} A_i \to B_i$ 和 $A_i \to C_i$
是具有相同核的满射。

\medskip\noindent
最后，应用 引理 \ref{more-morphisms-lemma-basic-example-pushout} 和
\ref{more-morphisms-lemma-pushout-fpqc-local}，可知我们的构造是概形范畴中的推出。
\end{proof}

\begin{lemma}
\label{more-morphisms-lemma-pushout-separated}
在 Situation \ref{more-morphisms-situation-pushout-along-closed-immersion-and-integral}
中，若 $X$ 和 $Y$ 分离，则推出 $Y \amalg_Z X$
（命题 \ref{more-morphisms-proposition-pushout-along-closed-immersion-and-integral}）
分离。把“分离”替换为“在 $S$ 上分离”、“拟分离”或“在 $S$ 上拟分离”，
结论同样成立。
\end{lemma}

\begin{proof}
态射 $Y \amalg X \to Y \amalg_Z X$ 满射且普遍闭。因此可以应用
态射，引理 \ref{morphisms-lemma-image-universally-closed-separated}。
\end{proof}

\begin{lemma}
\label{more-morphisms-lemma-pushout-finite-type}
在 Situation \ref{more-morphisms-situation-pushout-along-closed-immersion-and-integral}
中，假设 $S$ 是局部诺特概形，而 $X$、$Y$ 和 $Z$ 都在 $S$ 上局部有限型。
则推出 $Y \amalg_Z X$
（命题 \ref{more-morphisms-proposition-pushout-along-closed-immersion-and-integral}）
在 $S$ 上局部有限型。
\end{lemma}

\begin{proof}
在仿射开集上考察，结论化为代数进阶，Lemma
\ref{more-algebra-lemma-fibre-product-finite-type} 的结果。
\end{proof}

\begin{lemma}
\label{more-morphisms-lemma-pushout-functor}
在 Situation \ref{more-morphisms-situation-pushout-along-closed-immersion-and-integral}
中，假设给定交换图
$$
\xymatrix{
Y' \ar[d]^g & Z' \ar[l]^{j'} \ar[r]_{i'} \ar[d]^h & X' \ar[d]^f \\
Y & Z \ar[l] \ar[r] & X
}
$$
其中两个方块均为笛卡尔方块，且 $f, g, h$ 分离并局部拟有限。则
\begin{enumerate}
\item 推出 $Y \amalg_Z X$ 和 $Y' \amalg_{Z'} X'$ 存在；
\item $Y' \amalg_{Z'} X' \to Y \amalg_Z X$ 分离且局部拟有限；
\item 图中的两个方块
$$
\xymatrix{
Y' \ar[r] \ar[d] & Y' \amalg_{Z'} X' \ar[d] & X' \ar[l] \ar[d] \\
Y \ar[r] & Y \amalg_Z X & X \ar[l]
}
$$
均为笛卡尔方块。
\end{enumerate}
\end{lemma}

\begin{proof}
由 命题 \ref{more-morphisms-proposition-pushout-along-closed-immersion-and-integral}，
推出 $Y \amalg_Z X$ 存在。为证明推出 $Y' \amalg_{Z'} X'$ 存在，
只需检验 Situation \ref{more-morphisms-situation-pushout-along-closed-immersion-and-integral}
的条件 (3) 对 $(X', Y', Z', i', j')$ 成立。具体地，令
$y' \in Y'$，并把它的像记作 $y \in Y$。选择仿射开集 $U \subset X$，
使得 $i(j^{-1}(y)) \subset U$。选择包含在 $f^{-1}(U)$ 中的拟紧开集
$U' \subset X'$，使其包含拟紧子集 $i'((j')^{-1}(\{y'\}))$。由
引理 \ref{more-morphisms-lemma-etale-separated-ind-quasi-affine}，$U'$ 是拟仿射的。
由于 $Z'_{y'}$ 是某个域上的整代数的谱，可以应用极限，Lemma
\ref{limits-lemma-ample-profinite-set-in-principal-affine}，从而找到
$U'$ 的一个包含 $i'((j')^{-1}(\{y'\}))$ 的仿射开子概形，正合所需。

\medskip\noindent
验证存在性后，我们检验其余断言。在仿射局部，这恰好是代数进阶，Lemma
\ref{more-algebra-lemma-properties-algebras-over-fibre-product}
中 $B \to D$ 和 $A' \to C'$ 局部拟有限的情形\footnote{确切地说，
$X, Y, Z, Y \amalg_Z X, X', Y', Z', Y' \amalg_{Z'} X'$
分别对应于 $A', B, A, B', C', D, C, D'$。}。特别地，态射
$Y' \amalg_{Z'} X' \to Y \amalg_Z X$ 局部有限型。由代数进阶，Lemma
\ref{more-algebra-lemma-module-over-fibre-product}，图中的方块是笛卡尔的。
由于局部拟有限性可以在纤维上检验
（态射，引理 \ref{morphisms-lemma-quasi-finite-at-point-characterize}），
可知 $Y' \amalg_{Z'} X' \to Y \amalg_Z X$ 局部拟有限。

\medskip\noindent
还需检验 $Y' \amalg_{Z'} X' \to Y \amalg_Z X$ 分离。由
命题 \ref{more-morphisms-proposition-pushout-along-closed-immersion-and-integral}，
$Y' \amalg X' \to Y' \amalg_{Z'} X'$ 普遍闭且满射。又因为态射
$Y' \amalg X' \to Y \amalg_Z X$ 分离（它分解为
$Y' \amalg X' \to Y \amalg X \to Y \amalg_Z X$），结论由
态射，引理 \ref{morphisms-lemma-image-universally-closed-separated} 得出。
\end{proof}

\begin{lemma}
\label{more-morphisms-lemma-pushout-functor-equivalence-flat}
在 Situation \ref{more-morphisms-situation-pushout-along-closed-immersion-and-integral}
中，在推出 $Y \amalg_Z X$ 上平坦、分离且局部拟有限的概形所成的范畴，
等价于 引理 \ref{more-morphisms-lemma-pushout-functor} 中满足 $f, g, h$ 平坦的
$(X', Y', Z', i', j', f, g, h)$ 所成的范畴。把“平坦”替换为
“\'etale”，同样成立。
\end{lemma}

\begin{proof}
若从 引理 \ref{more-morphisms-lemma-pushout-functor} 中的
$(X', Y', Z', i', j', f, g, h)$ 出发，并且 $f, g, h$ 平坦或 \'etale，
则由代数进阶，Lemma
\ref{more-algebra-lemma-properties-algebras-over-fibre-product}，
$Y' \amalg_{Z'} X' \to Y \amalg_Z X$ 平坦或 \'etale。

\medskip\noindent
反过来，设 $W \to Y \amalg_Z X$ 是分离且局部拟有限的态射。置
$X' = W \times_{Y \amalg_Z X} X$、$Y' = W \times_{Y \amalg_Z X} Y$ 和
$Z' = W \times_{Y \amalg_Z X} Z$，并赋予它们显然的态射
$i', j', f, g, h$。形成推出 $Y' \amalg_{Z'} X'$，便得到态射
$$
Y' \amalg_{Z'} X' \longrightarrow W
$$
% P05-EN-0143: frozen source gives the base as Y amalg_X Z; the constructed pushout throughout is Y amalg_Z X.
它由推出的泛性质成为 $Y \amalg_X Z$ 上的概形态射。若不假设
$W \to Y \amalg_Z X$ 平坦，则一般而言该态射不是同构。
（事实上，代数进阶，Lemma
\ref{more-algebra-lemma-module-over-fibre-product-bis}
表明所显示的箭头是闭浸入，但一般不是同构。）然而，如果
% P05-EN-0144: frozen source switches the pushout to the undefined fibre product Y times_Z X in the flatness condition.
$W \to Y \times_Z X$ 平坦，则由代数进阶，Lemma
\ref{more-algebra-lemma-properties-algebras-over-fibre-product}，它是同构。
\end{proof}

\noindent
下面讨论两个态射均为闭浸入时的存在性。

\begin{lemma}
\label{more-morphisms-lemma-pushout-along-closed-immersions}
设 $i : Z \to X$ 和 $j : Z \to Y$ 是概形的闭浸入。则推出
$Y \amalg_Z X$ 在概形范畴中存在。图示为
$$
\xymatrix{
Z \ar[r]_i \ar[d]_j & X \ar[d]^a \\
Y \ar[r]^-b & Y \amalg_Z X
}
$$
该图是纤维方块，态射 $a$ 和 $b$ 都是闭浸入，并且有短正合列
$$
0 \to \mathcal{O}_{Y \amalg_Z X} \to
a_*\mathcal{O}_X \oplus b_*\mathcal{O}_Y \to
c_*\mathcal{O}_Z \to 0
$$
其中 $c = a \circ i = b \circ j$。
\end{lemma}

\begin{proof}
这是 Proposition
\ref{more-morphisms-proposition-pushout-along-closed-immersion-and-integral} 的特殊情形。
注意，Situation \ref{more-morphisms-situation-pushout-along-closed-immersion-and-integral}
中的假设 (3) 直接成立，因为 $j$ 的各纤维都是单点集。最后交换两个箭头的
角色，即可得出 $a$ 和 $b$ 都是闭浸入。
\end{proof}

\begin{lemma}
\label{more-morphisms-lemma-pushout-along-closed-immersions-properties-above}
设 $i : Z \to X$ 和 $j : Z \to Y$ 是概形的闭浸入。设
$f : X' \to X$ 和 $g : Y' \to Y$ 是概形态射，并设
$\varphi : X' \times_{X, i} Z \to Y' \times_{Y, j} Z$ 是 $Z$-概形的同构。
考虑态射
$$
h :
X' \amalg_{X' \times_{X, i} Z, \varphi} Y'
\longrightarrow
X \amalg_Z Y
$$
则有
\begin{enumerate}
\item $h$ 局部有限型，当且仅当 $f$ 和 $g$ 局部有限型；
\item $h$ 平坦，当且仅当 $f$ 和 $g$ 平坦；
\item $h$ 平坦且局部有限表示，当且仅当 $f$ 和 $g$ 平坦且局部有限表示；
\item $h$ 光滑，当且仅当 $f$ 和 $g$ 光滑；
\item $h$ 是 \'etale 的，当且仅当 $f$ 和 $g$ 是 \'etale 的；
% P05-EN-0145: frozen source contains the editorial placeholder "add more here as needed" as a numbered mathematical assertion.
\item 可按需要在此添加更多性质。
\end{enumerate}
\end{lemma}

\begin{proof}
由 引理 \ref{more-morphisms-lemma-pushout-along-closed-immersions}，我们知道这些推出存在，
特别地，得到态射 $h$。因此可以把出现的所有概形都替换为仿射概形。
在这种情形下，本引理的各项断言等价于代数进阶，Lemma
\ref{more-algebra-lemma-properties-algebras-over-fibre-product}
的相应断言。
\end{proof}





\section{相对态射}
\label{more-morphisms-section-relative-morphisms}

\noindent
本节证明一个可表性结果；我们将在基本群，Section
\ref{pione-section-finite-etale} 中用它证明关于一个概形的有限 \'etale
覆盖范畴的结果。本节材料在可表性判据，Section
\ref{criteria-section-relative-morphisms} 中以适当的一般性加以讨论。

\medskip\noindent
设 $S$ 是概形，$Z$ 和 $X$ 是 $S$ 上的概形。给定 $S$ 上的概形 $T$，
可以考虑 $S$ 上的态射 $b : T \times_S Z \to T \times_S X$。图示为
\begin{equation}
\label{more-morphisms-equation-hom}
\vcenter{
\xymatrix{
T \times_S Z \ar[rd] \ar[rr]_b & &
T \times_S X \ar[ld] & Z \ar[rd] & & X \ar[ld] \\
& T \ar[rrr] & & & S
}
}
\end{equation}
当然，也可以把 $b$ 看作态射 $b : T \times_S Z \to X$，使得
$$
\xymatrix{
T \times_S Z \ar[r] \ar[d] \ar@/^1pc/[rrr]_-b &
Z \ar[rd] & & X \ar[ld] \\
T \ar[rr] & & S
}
$$
交换。在这种情形下，可以定义函子
\begin{equation}
\label{more-morphisms-equation-hom-functor}
\mathit{Mor}_S(Z, X) : (\Sch/S)^{opp} \longrightarrow \textit{Sets},
\quad
T \longmapsto \{b\text{ 如上}\}
\end{equation}
下面是一个基本的可表性结果。

\begin{lemma}
\label{more-morphisms-lemma-hom-from-finite-free-into-affine}
设 $Z \to S$ 和 $X \to S$ 是仿射概形的态射。假设
$\Gamma(Z, \mathcal{O}_Z)$ 是有限自由的
$\Gamma(S, \mathcal{O}_S)$-模。则 $\mathit{Mor}_S(Z, X)$
可由 $S$ 上的一个仿射概形表示。
\end{lemma}

\begin{proof}
写成 $S = \Spec(R)$。为 $\Gamma(Z, \mathcal{O}_Z)$ 选择一个 $R$-基
$\{e_1, \ldots, e_m\}$。选择一个表示
$$
\Gamma(X, \mathcal{O}_X) = R[\{x_i\}_{i \in I}]/(\{f_k\}_{k \in K}).
$$
用 $\overline{x}_i$ 表示 $x_i$ 在该商中的像。写成
$$
P = R[\{a_{ij}\}_{i \in I, 1 \leq j \leq m}].
$$
考虑 $R$-代数映射
$$
\Psi :
R[\{x_i\}_{i \in I}]
\longrightarrow
P \otimes_R \Gamma(Z, \mathcal{O}_Z), \quad
x_i \longmapsto \sum\nolimits_j a_{ij} \otimes e_j.
$$
写成 $\Psi(f_k) = \sum c_{kj} \otimes e_j$，其中 $c_{kj} \in P$。
最后，用 $J \subset P$ 表示由元素 $c_{kj}$（$k \in K$，
$1 \leq j \leq m$）生成的理想。我们断言 $W = \Spec(P/J)$
表示函子 $\mathit{Mor}_S(Z, X)$。

\medskip\noindent
首先注意，按构造 $P/J$ 是 $R$-代数，因而给出态射 $W \to S$。
其次，按构造，映射 $\Psi$ 分解经过 $\Gamma(X, \mathcal{O}_X)$，
所以得到 $P/J$-代数同态
$$
P/J \otimes_R \Gamma(X, \mathcal{O}_X)
\longrightarrow
P/J \otimes_R \Gamma(Z, \mathcal{O}_Z)
$$
它确定一个态射 $b_{univ} : W \times_S Z \to W \times_S X$。
由 Yoneda 引理，$b_{univ}$ 确定函子变换
$W \to \mathit{Mor}_S(Z, X)$；我们断言它是同构。为证明这一点，
只需证明对任一仿射概形 $T$，它诱导集合的双射
$W(T) \to \mathit{Mor}_S(Z, X)(T)$。

\medskip\noindent
假设 $T = \Spec(R')$ 是 $S$ 上的仿射概形，并且
$b \in \mathit{Mor}_S(Z, X)(T)$。结构态射 $T \to S$ 在 $R'$ 上定义
$R$-代数结构，而 $b$ 定义一个 $R'$-代数映射
$$
b^\sharp :
R' \otimes_R \Gamma(X, \mathcal{O}_X)
\longrightarrow
R' \otimes_R \Gamma(Z, \mathcal{O}_Z).
$$
特别地，可以写成
$b^\sharp(1 \otimes \overline{x}_i) = \sum \alpha_{ij} \otimes e_j$，
其中 $\alpha_{ij} \in R'$。这对应于由规则
$a_{ij} \mapsto \alpha_{ij}$ 确定的 $R$-代数映射 $P \to R'$。
由理想 $J$ 的构造，该映射分解经过商 $P/J$，给出映射
$P/J \to R'$。这又对应于态射 $T \to W$，使得 $b$ 是 $b_{univ}$
的拉回。略去一些细节。
\end{proof}

\begin{lemma}
\label{more-morphisms-lemma-hom-from-finite-locally-free-into-affine}
设 $Z \to S$ 和 $X \to S$ 是概形态射。若 $Z \to S$ 有限局部自由，
且 $X \to S$ 仿射，则 $\mathit{Mor}_S(Z, X)$ 可由一个在 $S$
上仿射的概形表示。
\end{lemma}

\begin{proof}
选择仿射开覆盖 $S = \bigcup U_i$，使得
$\Gamma(Z \times_S U_i, \mathcal{O}_{Z \times_S U_i})$ 是有限自由的
$\mathcal{O}_S(U_i)$-模。令 $F_i \subset \mathit{Mor}_S(Z, X)$
为如下子函子：若 $T \to S$ 不分解经过 $U_i$，则它把 $T/S$
映到空集；否则映到 $\mathit{Mor}_S(Z, X)(T)$。这些子函子的集合满足
概形，引理 \ref{schemes-lemma-glue-functors} 的条件
(2)(a)、(2)(b)、(2)(c)，从而证明引理。条件 (2)(a) 由
引理 \ref{more-morphisms-lemma-hom-from-finite-free-into-affine} 得出，
其余两个条件由直接的论证得出。
\end{proof}

\noindent
下一个引理中关于态射 $f : X \to S$ 的条件，对证明这类断言非常有用。
在下列任一情形中，该条件成立：$X$ 拟仿射；$f$ 拟仿射；$f$ 拟射影；
$f$ 局部射影；$X$ 上存在充足可逆层；$X$ 上存在 $f$-充足可逆层；
或者 $X$ 上存在 $f$-很充足可逆层。

\begin{lemma}
\label{more-morphisms-lemma-hom-from-finite-locally-free-representable}
设 $Z \to S$ 和 $X \to S$ 是概形态射。假设
\begin{enumerate}
\item $Z \to S$ 有限局部自由；
\item 对所有满足 $s \in S$ 且 $x_1, \ldots, x_d \in X_s$ 的
$(s, x_1, \ldots, x_d)$，存在仿射开集 $U \subset X$，使得
$x_1, \ldots, x_d \in U$。
\end{enumerate}
则 $\mathit{Mor}_S(Z, X)$ 可由概形表示。
\end{lemma}

\begin{proof}
考虑由偶对 $(U, V)$ 组成的集合 $I$，其中 $U \subset X$ 和 $V \subset S$
是仿射开集，且 $U \to S$ 分解经过 $V$。对 $i \in I$，用
$(U_i, V_i)$ 表示相应的偶对。置
$F_i = \mathit{Mor}_{V_i}(Z_{V_i}, U_i)$。显然，$F_i$ 是
$\mathit{Mor}_S(Z, X)$ 的子函子。于是我们断言概形，Lemma
% P05-EN-0146: frozen source omits the predicate "hold" after claiming the three gluing conditions.
\ref{schemes-lemma-glue-functors} 的条件 (2)(a)、(2)(b)、(2)(c) 成立，
这就证明了引理。

\medskip\noindent
条件 (2)(a) 由
引理 \ref{more-morphisms-lemma-hom-from-finite-locally-free-into-affine} 得出。

\medskip\noindent
为检验条件 (2)(b)，考虑 $T/S$ 和 $b \in \mathit{Mor}_S(Z, X)(T)$。
把 $b$ 看作态射 $T \times_S Z \to X$，便得到开集
$b^{-1}(U_i) \subset T \times_S Z$。显然，$b \in F_i(T)$ 当且仅当
$b^{-1}(U_i) = T \times_S Z$。由于投影 $p : T \times_S Z \to T$
有限，因而是闭态射，满足 $p^{-1}(\{t\}) \subset b^{-1}(U_i)$
的点 $t \in T$ 所成的集合 $U_{i, b} \subset T$ 是开集。
于是 $f : T' \to T$ 分解经过 $U_{i, b}$，当且仅当
% P05-EN-0147: frozen source writes b composed with f although b has domain T times_S Z and f has codomain T; the base-changed map is required.
$b \circ f \in F_i(T')$；条件 (2)(b) 得证。

\medskip\noindent
最后检验条件 (2)(c)，此处要使用我们对 $X \to S$ 所加的条件。
具体地，考虑 $T/S$ 和 $b \in \mathit{Mor}_S(Z, X)(T)$。
只需证明每个 $t \in T$ 都包含在上一段定义的某个开集 $U_{i, b}$ 中。
这等价于存在某个 $i$，使得 $b(p^{-1}(\{t\})) \subset U_i$，
其中 $p : T \times_S Z \to T$ 是投影，而 $b : T \times_S Z \to X$
是给定的态射。由于 $p$ 有限，集合 $b(p^{-1}(\{t\})) \subset X$
是有限集，并且包含在 $X \to S$ 的如下纤维中：该纤维位于 $t$
在 $S$ 中的像 $s$ 之上。因此，我们对 $X \to S$ 的条件恰好表明
存在合适的偶对。
\end{proof}

\begin{lemma}
\label{more-morphisms-lemma-hom-from-finite-locally-free-separated-lqf}
设 $Z \to S$ 和 $X \to S$ 是概形态射。假设 $Z \to S$ 有限局部自由，
而 $X \to S$ 分离且局部拟有限。则 $\mathit{Mor}_S(Z, X)$
可由概形表示。
\end{lemma}

\begin{proof}
这由 引理 \ref{more-morphisms-lemma-hom-from-finite-locally-free-representable} 和
\ref{more-morphisms-lemma-separated-locally-quasi-finite-over-affine} 得出。
\end{proof}









\section{伪凝聚复形的刻画，III}
\label{more-morphisms-section-characterize-pseudo-coherent}

\noindent
本节讨论以其上同调刻画伪凝聚复形的方法。这是概形的导出范畴，Section
\ref{perfect-section-pseudo-coherent} 的延续。基本工具是使用 Chow 引理的
一个导出版本，把问题归约到射影空间的情形，见
引理 \ref{more-morphisms-lemma-derived-chow}。

\begin{lemma}
\label{more-morphisms-lemma-case-of-tor-independence}
考虑概形的交换图
$$
\xymatrix{
Z' \ar[d] \ar[r] & Y' \ar[d] \\
X' \ar[r] & S'
}
$$
设 $S \to S'$ 是态射。用 $X$ 和 $Y$ 分别表示 $X'$ 和 $Y'$ 到 $S$
的基变换。假设 $Y' \to S'$ 和 $Z' \to X'$ 平坦。则
$X \times_S Y$ 和 $Z'$ 在 $X' \times_{S'} Y'$ 上 Tor 独立。
\end{lemma}

\begin{proof}
问题是局部的，故可以假设所有概形都仿射（略去一些细节）。注意
$$
\xymatrix{
X \times_S Y \ar[r] \ar[d] & X' \times_{S'} Y' \ar[d] \\
X \ar[r] & X'
}
$$
是竖直箭头平坦的笛卡尔方块。写成 $X = \Spec(A)$、
$X' = \Spec(A')$、$X' \times_{S'} Y' = \Spec(B')$。则
$X \times_S Y = \Spec(A \otimes_{A'} B')$。再写成
$Z' = \Spec(C')$。我们要证明
$$
\text{Tor}_p^{B'}(A \otimes_{A'} B', C') = 0,
\quad\text{对 } p > 0
$$
由于 $A' \to B'$ 平坦，有
$A \otimes_{A'} B' = A \otimes_{A'}^\mathbf{L} B'$。因此
$$
(A \otimes_{A'} B') \otimes_{B'}^\mathbf{L} C' =
(A \otimes_{A'}^\mathbf{L} B') \otimes_{B'}^\mathbf{L} C' =
A \otimes_{A'}^\mathbf{L} C' =
A \otimes_{A'} C'
$$
第二个等号由代数进阶，Lemma
\ref{more-algebra-lemma-double-base-change} 得出；最后一个等号则因为
$A' \to C'$ 平坦。这就证明了引理。
\end{proof}

\begin{lemma}[导出 Chow 引理]
\label{more-morphisms-lemma-derived-chow}
设 $A$ 是环，$X$ 是 $A$ 上分离且有限表示的概形，并设 $x \in X$。
则存在 $x$ 的开邻域 $U \subset X$、整数 $n \geq 0$、
开集 $V \subset \mathbf{P}^n_A$、闭子概形
$Z \subset X \times_A \mathbf{P}^n_A$、点 $z \in Z$ 以及
$D(\mathcal{O}_{X \times_A \mathbf{P}^n_A})$ 中的对象 $E$，使得
\begin{enumerate}
\item $Z \to X \times_A \mathbf{P}^n_A$ 是有限表示的；
\item $b : Z \to X$ 在 $U$ 上是同构，且 $b(z) = x$；
\item $c : Z \to \mathbf{P}^n_A$ 在 $V$ 上是闭浸入；
\item $b^{-1}(U) = c^{-1}(V)$，特别地，$c(z) \in V$；
\item $E|_{X \times_A V} \cong
(b, c)_*\mathcal{O}_Z|_{X \times_A V}$；
\item $E$ 伪凝聚，且支撑在 $Z$ 上。
\end{enumerate}
\end{lemma}

\begin{proof}
可以找到有限型 $\mathbf{Z}$-子代数 $A' \subset A$，以及 $A'$ 上
分离且有限表示的概形 $X'$，使其到 $A$ 的基变换是 $X$。见极限，
引理 \ref{limits-lemma-descend-finite-presentation} 和
\ref{limits-lemma-descend-separated-finite-presentation}。
令 $x' \in X'$ 为 $x$ 的像。若能对 $x' \in X'/A'$ 证明本引理，
则对 $x \in X/A$ 的引理随即成立。具体地，若
$U', n', V', Z', z', E'$ 给出了 $x' \in X'/A'$ 的解，则可以令
$U \subset X$ 为 $U'$ 的原像，令 $n = n'$，令
$V \subset \mathbf{P}^n_A$ 为 $V'$ 的原像，令
$Z \subset X \times \mathbf{P}^n$ 为 $Z'$ 的概形论原像，
令 $z \in Z$ 为映到 $x$ 的唯一一点，并令 $E$ 为 $E'$ 的导出拉回。
由上同调，引理 \ref{cohomology-lemma-pseudo-coherent-pullback}，
$E$ 是伪凝聚的。只需再检验 (5)。为此，置
$W = b^{-1}(U) = c^{-1}(V)$ 和
% P05-EN-0148: frozen source uses U instead of U' in the definition of W' over X'.
$W' = (b')^{-1}(U) = (c')^{-1}(V')$，并考虑笛卡尔方块
$$
\xymatrix{
W \ar[d]_{(b, c)} \ar[r] & W' \ar[d]^{(b', c')} \\
X \times_A V \ar[r] & X' \times_{A'} V'
}
$$
由 引理 \ref{more-morphisms-lemma-case-of-tor-independence}，概形 $X \times_A V$ 和
$W'$ 在 $X' \times_{A'} V'$ 上 Tor 独立。因此由概形的导出范畴，
引理 \ref{perfect-lemma-compare-base-change}，
$(b', c')_*\mathcal{O}_{W'}$ 到 $X \times_A V$ 的导出拉回是
$(b, c)_*\mathcal{O}_W$。这里还用到
$R(b', c')_*\mathcal{O}_{Z'} = (b', c')_*\mathcal{O}_{Z'}$，
因为 $(b', c')$ 是闭浸入；对 $(b, c)_*\mathcal{O}_Z$ 也类似。
由于 $E'|_{U' \times_{A'} V'} =
(b', c')_*\mathcal{O}_{W'}$，我们得到
$E|_{U \times_A V} = (b, c)_*\mathcal{O}_W$，故 (5) 成立。
这把我们归约到下一段所描述的情形。

\medskip\noindent
假设 $A$ 在 $\mathbf{Z}$ 上有限型。选择 $x$ 的仿射开邻域
$U \subset X$。则 $U$ 在 $A$ 上有限型。选择闭浸入
$U \to \mathbf{A}^n_A$，并用 $j : U \to \mathbf{P}^n_A$ 表示它与开浸入
$\mathbf{A}^n_A \to \mathbf{P}^n_A$ 复合所得的浸入。令 $Z$ 为
$$
(\text{id}_U, j) : U \longrightarrow X \times_A \mathbf{P}^n_A
$$
的概形论闭包。由于投影 $X \times \mathbf{P}^n \to X$ 分离，
由态射，Lemma
\ref{morphisms-lemma-scheme-theoretic-image-of-partial-section} 可知
$b : Z \to X$ 在 $U$ 上是同构。令 $z \in Z$ 为位于 $x$ 上的唯一一点。

\medskip\noindent
令 $Y \subset \mathbf{P}^n_A$ 为 $j$ 的概形论闭包。于是显然，
$Z \subset X \times_A Y$ 是
$(\text{id}_U, j) : U \to X \times_A Y$ 的概形论闭包。
由于 $X$ 分离，态射 $X \times_A Y \to Y$ 也分离。因此由上面使用的
同一引理，$Z \to Y$ 在开子概形 $j(U) \subset Y$ 上是同构。
选择开集 $V \subset \mathbf{P}^n_A$，使得 $V \cap Y = j(U)$。
于是 (3) 和 (4) 成立。

\medskip\noindent
由于 $A$ 是诺特环，$X$ 和 $X \times_A \mathbf{P}^n_A$ 都是诺特概形。
因此在这种情形下可以取 $E = (b, c)_*\mathcal{O}_Z$，见概形的导出范畴，
引理 \ref{perfect-lemma-identify-pseudo-coherent-noetherian}。
证明完成。
\end{proof}

\begin{lemma}
\label{more-morphisms-lemma-compute-Fourier-Mukai-for-derived-chow}
设 $A$、$x \in X$ 以及 $U, n, V, Z, z, E$ 如
引理 \ref{more-morphisms-lemma-derived-chow} 所述。对任意
$K \in D_\QCoh(\mathcal{O}_X)$，有
$$
Rq_*(Lp^*K \otimes^\mathbf{L} E)|_V = R(U \to V)_*K|_U
$$
其中 $p : X \times_A \mathbf{P}^n_A \to X$ 和
$q : X \times_A \mathbf{P}^n_A \to \mathbf{P}^n_A$ 是投影，
而态射 $U \to V$ 是有限表示的闭浸入 $c \circ (b|_U)^{-1}$。
\end{lemma}

\begin{proof}
由于 $b^{-1}(U) = c^{-1}(V)$，且 $c$ 在 $V$ 上是闭浸入，
$c \circ (b|_U)^{-1}$ 是闭浸入。它是有限表示的，因为 $U$ 和 $V$
都在 $A$ 上有限表示，见态射，Lemma
\ref{morphisms-lemma-finite-presentation-permanence}。首先有
$$
Rq_*(Lp^*K \otimes^\mathbf{L} E)|_V =
Rq'_*\left((Lp^*K \otimes^\mathbf{L} E)|_{X \times_A V}\right)
$$
其中 $q' : X \times_A V \to V$ 是投影；这里使用了总直像的形成与局部化相容。
置 $W = b^{-1}(U) = c^{-1}(V)$，并用 $i : W \to X \times_A V$
表示闭浸入 $i = (b, c)|_W$。则
$$
Rq'_*\left((Lp^*K \otimes^\mathbf{L} E)|_{X \times_A V}\right) =
Rq'_*(Lp^*K|_{X \times_A V} \otimes^\mathbf{L} i_*\mathcal{O}_W)
$$
由性质 (5) 得到该等式。由于 $i$ 是闭浸入，有
$i_*\mathcal{O}_W = Ri_*\mathcal{O}_W$。使用概形的导出范畴，Lemma
\ref{perfect-lemma-cohomology-base-change}，可以把它改写为
$$
Rq'_* Ri_* Li^* Lp^*K|_{X \times_A V} =
R(q' \circ i)_* Lb^*K|_W =
R(U \to V)_* K|_U
$$
这正是所需结论。
\end{proof}

\begin{lemma}
\label{more-morphisms-lemma-characterize-pseudo-coherent}
设 $A$ 是环，$X$ 是 $A$ 上分离且有限表示的概形，并设
$K \in D_\QCoh(\mathcal{O}_X)$。若对 $D(\mathcal{O}_X)$ 中的每个
伪凝聚对象 $E$，$R\Gamma(X, E \otimes^\mathbf{L} K)$ 在 $D(A)$
中都是伪凝聚的，则 $K$ 相对于 $A$ 伪凝聚。
\end{lemma}

\begin{proof}
假设 $K \in D_\QCoh(\mathcal{O}_X)$，并且对 $D(\mathcal{O}_X)$
中的每个伪凝聚对象 $E$，$R\Gamma(X, E \otimes^\mathbf{L} K)$
在 $D(A)$ 中都是伪凝聚的。令 $x \in X$。我们将证明 $K$
在 $x$ 的某个邻域中相对于 $A$ 伪凝聚；这便证明了引理。

\medskip\noindent
选择 引理 \ref{more-morphisms-lemma-derived-chow} 中的 $U, n, V, Z, z, E$。
用 $p : X \times \mathbf{P}^n \to X$ 和
$q : X \times \mathbf{P}^n \to \mathbf{P}^n_A$ 表示投影。
则对任意 $i \in \mathbf{Z}$，有
\begin{align*}
& R\Gamma(\mathbf{P}^n_A,
Rq_*(Lp^*K \otimes^\mathbf{L} E)
\otimes^\mathbf{L}
\mathcal{O}_{\mathbf{P}^n_A}(i)) \\
& =
R\Gamma(X \times \mathbf{P}^n,
Lp^*K \otimes^\mathbf{L} E \otimes^\mathbf{L}
Lq^*\mathcal{O}_{\mathbf{P}^n_A}(i)) \\
& =
R\Gamma(X,
K \otimes^\mathbf{L} Rp_*(E \otimes^\mathbf{L}
Lq^*\mathcal{O}_{\mathbf{P}^n_A}(i)))
\end{align*}
这是由概形的导出范畴，Lemma
\ref{perfect-lemma-cohomology-base-change} 得到的。由概形的导出范畴，
引理 \ref{perfect-lemma-flat-proper-pseudo-coherent-direct-image-general}，
复形 $Rp_*(E \otimes^\mathbf{L} Lq^*\mathcal{O}_{\mathbf{P}^n_A}(i))$
在 $X$ 上伪凝聚。因此假设表明，所显示公式中的表达式是 $D(A)$
中的伪凝聚对象。由概形的导出范畴，Lemma
\ref{perfect-lemma-pseudo-coherent-on-projective-space}，可知
$Rq_*(Lp^*K \otimes^\mathbf{L} E)$ 在 $\mathbf{P}^n_A$ 上伪凝聚。
由 引理 \ref{more-morphisms-lemma-compute-Fourier-Mukai-for-derived-chow}，有
$$
% P05-EN-0149: frozen source restricts a complex on projective space to X times_A V; Lemma compute-Fourier-Mukai states the restriction is to V.
Rq_*(Lp^*K \otimes^\mathbf{L} E)|_{X \times_A V} =
R(U \to V)_*K|_U
$$
由于 $U \to V$ 是到 $\mathbf{P}^n_A$ 的一个开子概形中的闭浸入，
由 引理 \ref{more-morphisms-lemma-check-relative-pseudo-coherence-on-charts}，
这表明 $K|_U$ 相对于 $A$ 伪凝聚。
\end{proof}

\begin{lemma}
\label{more-morphisms-lemma-characterize-pseudo-coh-improved}
设 $A$ 是环，$X$ 是 $A$ 上分离且有限表示的概形。设
$K \in D_\QCoh(\mathcal{O}_X).$ 若对每个完美对象
$E \in D(\mathcal{O}_X)$，$R \Gamma (X, E \otimes ^{\mathbf{L}} K)$
在 $D(A)$ 中都伪凝聚，则 $K$ 相对于 $A$ 伪凝聚。
\end{lemma}

\begin{proof}
鉴于 引理 \ref{more-morphisms-lemma-characterize-pseudo-coherent}，只需证明对每个
伪凝聚对象 $E \in D(\mathcal{O}_X)$，
$R \Gamma (X, E \otimes ^{\mathbf{L}} K)$ 在 $D(A)$ 中都伪凝聚。
由概形的导出范畴，Proposition
\ref{perfect-proposition-detecting-bounded-above}，可知
$K \in D^-_\QCoh (\mathcal{O}_X)$。现在结论由概形的导出范畴，Lemma
\ref{perfect-lemma-perfect-enough} 得出。
\end{proof}

\begin{lemma}
\label{more-morphisms-lemma-characterize-relatively-perfect}
设 $A$ 是环，$X$ 是 $A$ 上分离、有限表示且平坦的概形。设
$K \in D_\QCoh(\mathcal{O}_X).$ 若对每个完美对象
$E \in D(\mathcal{O}_X)$，$R \Gamma (X, E \otimes^\mathbf{L} K)$
在 $D(A)$ 中都完美，则 $K$ 是 $\Spec(A)$-完美的。
\end{lemma}

\begin{proof}
由 引理 \ref{more-morphisms-lemma-characterize-pseudo-coh-improved}，$K$ 相对于 $A$
伪凝聚。由 引理 \ref{more-morphisms-lemma-check-relative-pseudo-coherence-on-charts}，
$K$ 在 $D( \mathcal{O}_X)$ 中伪凝聚。由概形的导出范畴，Proposition
\ref{perfect-proposition-detecting-bounded-below}，可知 $K$ 属于
$D^-(\mathcal{O}_X)$。设 $\mathfrak{p}$ 是 $A$ 的素理想，并用
$i : Y \to X$ 表示 $\mathfrak{p}$ 上的概形论纤维的嵌入，即 $Y$
是 $\kappa(\mathfrak p)$ 上的概形。由概形的导出范畴，Lemma
\ref{perfect-lemma-bounded-on-fibres}，只要证明 $Li^*(K)$ 下有界即可。
令 $G \in D_{perf} (\mathcal{O}_X)$ 为生成 $D_\QCoh (\mathcal{O}_X)$
的完美复形，见概形的导出范畴，Theorem
\ref{perfect-theorem-bondal-van-den-Bergh}。我们有
\begin{align*}
R\Hom _{\mathcal{O}_Y}(Li^*(G), Li^*(K))
& =
R\Gamma(Y, Li^*(G ^\vee \otimes ^\mathbf{L} K)) \\
& =
R\Gamma(X, G^\vee \otimes ^{\mathbf{L}} K)
\otimes^\mathbf{L}_A \kappa(\mathfrak{p})
\end{align*}
第一个等号使用了 $Li^*$ 保持完美对象及其对偶这一事实，以及上同调，
引理 \ref{cohomology-lemma-dual-perfect-complex}；略去一些细节。
由于 $X$ 在 $A$ 上平坦，第二个等号由概形的导出范畴，Lemma
\ref{perfect-lemma-compare-base-change} 得出。由假设，它是
$D(\kappa(\mathfrak{p}))$ 中的完美对象。由概形的导出范畴，Remark
\ref{perfect-remark-pullback-generator}，对象
$Li^*(G) \in D_{perf}(\mathcal{O}_Y)$ 生成 $D_\QCoh(\mathcal{O}_Y)$。
因此概形的导出范畴，Proposition
\ref{perfect-proposition-detecting-bounded-below} 现在表明 $Li^*(K)$
下有界，证明完成。
\end{proof}




\section{复形有限性性质的下降}
\label{more-morphisms-section-descent-finiteness}

\noindent
本节是概形的导出范畴，Section
\ref{perfect-section-descent-finiteness} 的延续。

\begin{lemma}
\label{more-morphisms-lemma-relative-pseudo-coherent-descends-fppf}
设 $X \to S$ 局部有限型，$\{f_i : X_i \to X\}$ 是概形的 fppf 覆盖。
设 $E \in D_\QCoh(\mathcal{O}_X)$，并设 $m \in \mathbf{Z}$。
则 $E$ 相对于 $S$ 是 $m$-伪凝聚的，当且仅当每个 $Lf_i^*E$
相对于 $S$ 都是 $m$-伪凝聚的。
\end{lemma}

\begin{proof}
假设 $E$ 相对于 $S$ 是 $m$-伪凝聚的。由
引理 \ref{more-morphisms-lemma-flat-finite-presentation-pseudo-coherent}，态射 $f_i$
是伪凝聚的。因此由 引理 \ref{more-morphisms-lemma-pull-relative-pseudo-coherent}，
$Lf_i^*E$ 相对于 $S$ 是 $m$-伪凝聚的。

\medskip\noindent
反过来，假设对每个 $i$，$Lf_i^*E$ 相对于 $S$ 都是 $m$-伪凝聚的。
按照 引理 \ref{more-morphisms-lemma-dominate-fppf} 选择 $S = \bigcup U_j$、
$W_j \to U_j$、$W_j = \bigcup W_{j, k}$、$T_{j, k} \to W_{j, k}$，
以及 $S$ 上的态射
$\alpha_{j, k} : T_{j, k} \to X_{i(j, k)}$。由于态射
% P05-EN-0150: frozen source switches the lower index k to an uppercase K in T_{j,K}.
$T_{j, K} \to S$ 平坦且有限表示，由
引理 \ref{more-morphisms-lemma-permanence-pseudo-coherent} 可知 $\alpha_{j, k}$
是伪凝聚的。因此
$$
% P05-EN-0151: frozen source switches T_{j,k} to T_{i,k} on the right side of the pullback identity.
L\alpha_{j, k}^*Lf_{i(j, k)}^*E = L(T_{i, k} \to S)^*E
$$
由 引理 \ref{more-morphisms-lemma-pull-relative-pseudo-coherent}，它相对于 $S$
是 $m$-伪凝聚的。现在要沿覆盖
$\{T_{j, k} \to W_{j, k}\}$、$W_j = \bigcup W_{j, k}$、
$\{W_j \to U_j\}$ 和 $S = \bigcup U_j$ 下降这一性质。
对 Zariski 覆盖，结论由相对于 $S$ 的 $m$-伪凝聚性的定义成立，
故可以假设有单个满射有限局部自由态射 $\pi : Y \to X$，使得
% P05-EN-0152: frozen source drops the qualifier m- at the decisive descent step and in the following direct-image claim.
$L\pi^*E$ 相对于 $S$ 伪凝聚。此时由
引理 \ref{more-morphisms-lemma-finite-morphism-relative-pseudo-coherence}，
$R\pi_*L\pi^*E$ 相对于 $S$ 伪凝聚（这是我们第一次使用 $E$
具有拟凝聚上同调层）。我们有
$R\pi_*L\pi^*E = E \otimes^\mathbf{L}_{\mathcal{O}_X} \pi_*\mathcal{O}_Y$
例如由概形的导出范畴，Lemma
\ref{perfect-lemma-cohomology-base-change} 得到该等式；并且在 $X$
上局部地，映射 $\mathcal{O}_X \to \pi_*\mathcal{O}_Y$ 是一个直和项的嵌入。
因此结论由 引理 \ref{more-morphisms-lemma-summands-relative-pseudo-coherent} 得出。
\end{proof}

\begin{lemma}
\label{more-morphisms-lemma-relative-pseudo-coherent-post-compose}
设 $X \to T \to S$ 是概形态射。假设 $T \to S$ 平坦且局部有限表示，
而 $X \to T$ 局部有限型。设 $E \in D(\mathcal{O}_X)$，并设
$m \in \mathbf{Z}$。则 $E$ 相对于 $S$ 是 $m$-伪凝聚的，当且仅当
$E$ 相对于 $T$ 是 $m$-伪凝聚的。
\end{lemma}

\begin{proof}
在 $X$ 上局部地，可以选择闭浸入 $i : X \to \mathbf{A}^n_T$。
于是 $\mathbf{A}^n_T \to S$ 平坦且局部有限表示。因此应用
引理 \ref{more-morphisms-lemma-composition-relative-pseudo-coherent} 即得所述等价性。
\end{proof}

\begin{lemma}
\label{more-morphisms-lemma-relative-pseudo-coherent-descends-fppf-base}
设 $f : X \to S$ 局部有限型，$\{S_i \to S\}$ 是概形的 fppf 覆盖。
用 $f_i : X_i \to S_i$ 表示 $f$ 的基变换，并用 $g_i : X_i \to X$
表示投影。设 $E \in D_\QCoh(\mathcal{O}_X)$，并设 $m \in \mathbf{Z}$。
则 $E$ 相对于 $S$ 是 $m$-伪凝聚的，当且仅当每个 $Lg_i^*E$
相对于 $S_i$ 都是 $m$-伪凝聚的。
\end{lemma}

\begin{proof}
这形式地由 引理 \ref{more-morphisms-lemma-relative-pseudo-coherent-descends-fppf} 和
\ref{more-morphisms-lemma-relative-pseudo-coherent-post-compose} 得出。具体地，若
$E$ 相对于 $S$ 是 $m$-伪凝聚的，则 $Lg_i^*E$ 相对于 $S$
是 $m$-伪凝聚的（由第一个引理），因而 $Lg_i^*E$ 相对于 $S_i$ 是
$m$-伪凝聚的（由第二个引理）。反过来，若 $Lg_i^*E$ 相对于 $S_i$
是 $m$-伪凝聚的，则 $Lg_i^*E$ 相对于 $S$ 是 $m$-伪凝聚的（由第二个引理），
因而 $E$ 相对于 $S$ 是 $m$-伪凝聚的（由第一个引理）。
\end{proof}






\section{相对完美对象}
\label{more-morphisms-section-relatively-perfect}

\noindent
本节是概形的导出范畴，Section
\ref{perfect-section-relatively-perfect} 中讨论的延续。

\begin{lemma}
\label{more-morphisms-lemma-thickening-pseudo-coherent}
设 $i : X \to X'$ 是概形的有限阶增厚。设
$K' \in D(\mathcal{O}_{X'})$ 是对象，并且 $K = Li^*K'$ 伪凝聚。
则 $K'$ 伪凝聚。
\end{lemma}

\begin{proof}
先证明 $K'$ 具有拟凝聚上同调层。为此，可以归约到一阶增厚的情形，见
第 \ref{more-morphisms-section-thickenings} 节。令
$\mathcal{I} \subset \mathcal{O}_{X'}$ 为定义 $X$ 的拟凝聚理想层。
把短正合列
$$
0 \to \mathcal{I} \to \mathcal{O}_{X'} \to i_*\mathcal{O}_X \to 0
$$
与 $K'$ 作张量积，得到优三角
$$
K' \otimes_{\mathcal{O}_{X'}}^\mathbf{L} \mathcal{I}
\to K' \to
K' \otimes_{\mathcal{O}_{X'}}^\mathbf{L} i_*\mathcal{O}_X
\to
(K' \otimes_{\mathcal{O}_{X'}}^\mathbf{L} \mathcal{I})[1]
$$
由于 $i_* = Ri_*$，并且可以把 $\mathcal{I}$ 看成拟凝聚
$\mathcal{O}_X$-模（因为这里是一阶增厚），可以把它改写为
$$
i_*(K \otimes_{\mathcal{O}_X}^\mathbf{L} \mathcal{I})
\to K' \to
i_*K \to
i_*(K \otimes_{\mathcal{O}_X}^\mathbf{L} \mathcal{I})[1]
$$
使用上同调，Lemma
\ref{cohomology-lemma-projection-formula-closed-immersion} 来识别各项。
由于 $K$ 属于 $D_\QCoh(\mathcal{O}_X)$，可知 $K'$ 属于
$D_\QCoh(\mathcal{O}_{X'})$；这里使用了概形的导出范畴，Lemmas
\ref{perfect-lemma-pseudo-coherent}、
\ref{perfect-lemma-quasi-coherence-tensor-product} 和
\ref{perfect-lemma-quasi-coherence-direct-image}。

\medskip\noindent
假设 $K'$ 属于 $D_\QCoh(\mathcal{O}_{X'})$。问题在 $X'$ 上是局部的，
故可以假设 $X'$ 仿射。设 $X' = \Spec(A')$、$X = \Spec(A)$，其中
$A = A'/I$ 且 $I$ 幂零。于是 $K'$ 来自某个对象 $M' \in D(A')$，
见概形的导出范畴，引理 \ref{perfect-lemma-affine-compare-bounded}。
因此由概形的导出范畴，Lemma
\ref{perfect-lemma-pseudo-coherent-affine} 以及对 $K$ 的假设，
$M = M' \otimes_{A'}^\mathbf{L} A$ 是 $D(A)$ 中的伪凝聚对象。
进而由代数进阶，Lemma
\ref{more-algebra-lemma-check-pseudo-coherent-modulo-nilpotent}，
$M'$ 伪凝聚。证明完成。
\end{proof}

\begin{lemma}
\label{more-morphisms-lemma-thickening-relatively-perfect}
考虑概形的笛卡尔图
$$
\xymatrix{
X \ar[r]_i \ar[d]_f & X' \ar[d]^{f'} \\
Y \ar[r]^j & Y'
}
$$
假设 $X' \to Y'$ 平坦且局部有限表示，而 $Y \to Y'$ 是有限阶增厚。
设 $E' \in D(\mathcal{O}_{X'})$。若 $E = Li^*(E')$ 是 $Y$-完美的，
则 $E'$ 是 $Y'$-完美的。
\end{lemma}

\begin{proof}
回忆，$E$ 是 $Y$-完美的，意味着 $E$ 伪凝聚，并且作为
$f^{-1}\mathcal{O}_Y$-模复形局部具有有限 Tor 维数
（概形的导出范畴，Definition
\ref{perfect-definition-relatively-perfect}）。由
引理 \ref{more-morphisms-lemma-thickening-pseudo-coherent}，$E'$ 伪凝聚。特别地，
$E'$ 属于 $D_\QCoh(\mathcal{O}_{X'})$，见概形的导出范畴，Lemma
\ref{perfect-lemma-pseudo-coherent}。为证明 $E'$ 局部具有有限 Tor 维数，
可以在 $X'$ 上局部地工作。因此可以假设
% P05-EN-0153: frozen source switches the diagram's bases Y',Y to undefined S',S in this affine reduction.
$X'$、$S'$、$X$、$S$ 都仿射，设它们分别由环 $A'$、$R'$、$A$、$R$
给出。于是由概形的导出范畴，Lemma
\ref{perfect-lemma-affine-locally-rel-perfect}，问题归约到交换代数版本；
该版本见代数进阶，Lemma
\ref{more-algebra-lemma-thickening-relatively-perfect}。
\end{proof}

\begin{lemma}
\label{more-morphisms-lemma-henselian-relatively-perfect}
设 $(R, I)$ 是由环和包含在 Jacobson 根中的理想 $I$ 组成的二元组。
置 $S = \Spec(R)$ 和 $S_0 = \Spec(R/I)$。设 $f : X \to S$
固有、平坦且有限表示。记 $X_0 = S_0 \times_S X$。设
$E \in D(\mathcal{O}_X)$ 伪凝聚。若 $E$ 到 $X_0$ 的导出限制 $E_0$
是 $S_0$-完美的，则 $E$ 是 $S$-完美的。
\end{lemma}

\begin{proof}
选择有限仿射开覆盖 $X = U_1 \cup \ldots \cup U_n$。对每个 $i$，
可以选择闭浸入 $U_i \to \mathbf{A}^{d_i}_S$。置
$U_{i, 0} = S_0 \times_S U_i$。对每个 $i$，复形
$E_0|_{U_{i, 0}}$ 的 Tor 振幅包含在某个 $[a_i, b_i]$ 中，其中
$a_i, b_i \in \mathbf{Z}$。令 $x \in X$。我们将证明对某个 $i$，
$E_x$ 在 $R$ 上的 Tor 振幅包含在 $[a_i - d_i, b_i]$ 中。
这将完成证明，因为由上同调，Lemma
\ref{cohomology-lemma-tor-amplitude-stalk}，Tor 振幅可以从茎读出。

\medskip\noindent
由于 $f$ 固有，$f(\overline{\{x\}})$ 是 $S$ 的闭子集。由于 $I$
包含在 Jacobson 根中，$f(\overline{\{x\}})$ 与闭子集
% P05-EN-0154: frozen source has the malformed finite verb "meeting" where the proof requires "meets".
$S_0 \subset S$ 相交。因此存在特化 $x \leadsto x_0$，其中
$x_0 \in X_0$。选择满足 $x_0 \in U_i$ 的 $i$，于是
$x_0 \in U_{i, 0}$。在余下的证明中固定该 $i$。
写成 $U_i = \Spec(A)$。则 $A$ 是平坦有限表示的 $R$-代数，
并且是一个含 $d_i$ 个变量的多项式 $R$-代数的商。限制
$E|_{U_i}$ 对应于 $D(A)$ 中的某个伪凝聚对象 $K$
（由概形的导出范畴，Lemmas
\ref{perfect-lemma-affine-compare-bounded} 和
\ref{perfect-lemma-pseudo-coherent-affine}）。注意，$E_0$ 对应于
$K \otimes_A^\mathbf{L} A/IA$。令
$\mathfrak q \subset \mathfrak q_0 \subset A$ 为对应于
$x \leadsto x_0$ 的素理想。则 $E_x = K_{\mathfrak q}$，而
$K_{\mathfrak q}$ 是 $K_{\mathfrak q_0}$ 的局部化。因此只需证明
$K_{\mathfrak q_0}$ 作为 $R$-模复形的 Tor 振幅包含在
$[a_i - d_i, b_i]$ 中。令 $I \subset \mathfrak p_0 \subset R$
为对应于 $f(x_0)$ 的素理想。于是有
\begin{align*}
K \otimes_R^\mathbf{L} \kappa(\mathfrak p_0)
& =
(K \otimes_R^\mathbf{L} R/I) \otimes_{R/I}^\mathbf{L}
\kappa(\mathfrak p_0) \\
& =
(K \otimes_A^\mathbf{L} A/IA) \otimes_{R/I}^\mathbf{L} \kappa(\mathfrak p_0)
\end{align*}
第二个等号成立是因为 $R \to A$ 平坦。由对 $a_i, b_i$ 的选择，
该复形仅在区间 $[a_i, b_i]$ 中的次数上有上同调。因此最后把
代数进阶，Lemma
\ref{more-algebra-lemma-lift-from-fibre-relatively-perfect}
应用于 $R \to A$、$\mathfrak q_0$、$\mathfrak p_0$ 和 $K$ 即得结论。
\end{proof}










\section{收缩有理曲线}
\label{more-morphisms-section-contracting}

\noindent
本节研究固有态射 $f : X \to Y$，其纤维维数 $\leq 1$，并且满足
$R^1f_*\mathcal{O}_X = 0$。为理解本节题名，请参阅代数曲线，Sections
\ref{curves-section-contracting-rational-tails}、
\ref{curves-section-contracting-rational-bridges} 和
\ref{curves-section-contracting-to-stable}。

\begin{lemma}
\label{more-morphisms-lemma-check-h1-fibre-zero}
设 $f : X \to Y$ 是概形的固有态射，$y \in Y$ 是满足
$\dim(X_y) \leq 1$ 的点。若
\begin{enumerate}
\item $R^1f_*\mathcal{O}_X = 0$；或者更一般地，
\item 存在态射 $g : Y' \to Y$，使得 $y$ 位于 $g$ 的像中，并且
$R^1f'_*\mathcal{O}_{X'} = 0$，其中 $f' : X' \to Y'$ 是 $f$
沿 $g$ 的基变换，
\end{enumerate}
则 $H^1(X_y, \mathcal{O}_{X_y}) = 0$。
\end{lemma}

\begin{proof}
为证明引理，可以用 $y$ 的开邻域代替 $Y$。因此可以假设 $Y$ 仿射，
并且 $f$ 的所有纤维维数均 $\leq 1$，见态射，Lemma
\ref{morphisms-lemma-openness-bounded-dimension-fibres}。在这种情形下，
$R^1f_*\mathcal{O}_X$ 是有限型拟凝聚 $\mathcal{O}_Y$-模，且其形成与
任意基变换相容，见极限，Lemmas
\ref{limits-lemma-proper-top-cohomology-finite-type} 和
\ref{limits-lemma-higher-direct-images-zero-above-dimension-fibre}。
引理立即得证。
\end{proof}

\begin{lemma}
\label{more-morphisms-lemma-h1-fibre-zero}
设 $f : X \to Y$ 是概形的固有态射，$y \in Y$ 是满足
$\dim(X_y) \leq 1$ 且 $H^1(X_y, \mathcal{O}_{X_y}) = 0$ 的点。
则存在 $y$ 的开邻域 $V \subset Y$，使得
$R^1f_*\mathcal{O}_X|_V = 0$，并且沿任意 $Y' \to V$ 基变换后仍成立。
\end{lemma}

\begin{proof}
为证明引理，可以用 $y$ 的开邻域代替 $Y$。因此可以假设 $Y$ 仿射，
并且 $f$ 的所有纤维维数均 $\leq 1$，见态射，Lemma
\ref{morphisms-lemma-openness-bounded-dimension-fibres}。在这种情形下，
$R^1f_*\mathcal{O}_X$ 是有限型拟凝聚 $\mathcal{O}_Y$-模，且其形成与
任意基变换相容，见极限，Lemmas
\ref{limits-lemma-proper-top-cohomology-finite-type} 和
\ref{limits-lemma-higher-direct-images-zero-above-dimension-fibre}。
设 $Y = \Spec(A)$，$y$ 对应于素理想 $\mathfrak p \subset A$，而
$R^1f_*\mathcal{O}_X$ 对应于有限 $A$-模 $M$。由上述基变换断言，
$H^1(X_y, \mathcal{O}_{X_y}) = 0$ 意味着
$\mathfrak pM_\mathfrak p = M_\mathfrak p$。由 Nakayama 引理，
$M_\mathfrak p = 0$。由于 $M$ 有限，可以找到 $f \in A$、
$f \not \in \mathfrak p$，使得 $M_f = 0$。因此取主开集
$V$ 为 $D(f)$，便得到所需结论。
\end{proof}

\begin{lemma}
\label{more-morphisms-lemma-globally-generated-vanishing}
设 $f : X \to Y$ 是概形的固有态射，并且对所有 $y \in Y$ 均有
$\dim(X_y) \leq 1$ 和 $H^1(X_y, \mathcal{O}_{X_y}) = 0$。
设 $\mathcal{F}$ 是 $X$ 上的拟凝聚层。则
\begin{enumerate}
\item 对 $p > 1$，有 $R^pf_*\mathcal{F} = 0$；
\item 若存在满射 $f^*\mathcal{G} \to \mathcal{F}$，其中
$\mathcal{G}$ 是 $Y$ 上的拟凝聚层，则 $R^1f_*\mathcal{F} = 0$。
\end{enumerate}
若 $Y$ 仿射，则还有
\begin{enumerate}
\item[(3)] 对 $p \not \in \{0, 1\}$，有 $H^p(X, \mathcal{F}) = 0$；
\item[(4)] 若 $\mathcal{F}$ 全局生成，则 $H^1(X, \mathcal{F}) = 0$。
\end{enumerate}
\end{lemma}

\begin{proof}
(1) 中的消失就是极限，Lemma
\ref{limits-lemma-higher-direct-images-zero-above-dimension-fibre}。
为证明 (2)，可以在 $Y$ 上局部地工作，并假设 $Y$ 仿射。此时
$R^1f_*\mathcal{F}$ 是与模 $H^1(X, \mathcal{F})$ 相伴的 $Y$ 上拟凝聚模。
这里使用了 $Y$ 仿射、高阶直像的拟凝聚性（概形的上同调，Lemma
\ref{coherent-lemma-quasi-coherence-higher-direct-images}），以及概形的上同调，
引理 \ref{coherent-lemma-quasi-coherence-higher-direct-images-application}。
由于 $Y$ 仿射，拟凝聚模 $\mathcal{G}$ 全局生成，因而
$f^*\mathcal{G}$ 和 $\mathcal{F}$ 也全局生成。这样便看出 (4) 蕴含 (2)。
(3) 由 (1) 以及刚才关于直像拟凝聚性的说明得出。因此只需证明 (4)。
若 $\mathcal{F}$ 全局生成，则存在满射
$\bigoplus_{i \in I} \mathcal{O}_X \to \mathcal{F}$。由 (1) 和上同调长正合列，
它诱导 $H^1$ 上的满射。由 引理 \ref{more-morphisms-lemma-h1-fibre-zero}，
$R^1f_*\mathcal{O}_X = 0$，故 $H^1(X, \mathcal{O}_X) = 0$；并且
$H^1(X, -)$ 与直和相容（上同调，Lemma
\ref{cohomology-lemma-quasi-separated-cohomology-colimit}）。结论随即得出。
\end{proof}

\begin{lemma}
\label{more-morphisms-lemma-h1-fibre-zero-check-h0-kappa}
设 $f : X \to Y$ 是概形的固有态射。假设
\begin{enumerate}
\item 对所有 $y \in Y$，均有 $\dim(X_y) \leq 1$ 和
$H^1(X_y, \mathcal{O}_{X_y}) = 0$；
\item $\mathcal{O}_Y \to f_*\mathcal{O}_X$ 满射。
\end{enumerate}
则对 $f$ 的任一基变换 $f' : X' \to Y'$，态射
$\mathcal{O}_{Y'} \to f'_*\mathcal{O}_{X'}$ 满射。
\end{lemma}

\begin{proof}
可以假设 $Y$ 和 $Y'$ 仿射。于是可以选择闭浸入 $Y' \to Y''$，
其中 $Y'' \to Y$ 是仿射概形之间的平坦态射。由平坦基变换
（概形的上同调，引理 \ref{coherent-lemma-flat-base-change-cohomology}），
结论对 $X'' \to Y''$ 成立。因此可以假设 $Y'$ 是 $Y$ 的闭子概形。
令 $\mathcal{I} \subset \mathcal{O}_Y$ 为定义 $Y'$ 的理想。于是有短正合列
$$
0 \to \mathcal{I}\mathcal{O}_X \to
\mathcal{O}_X \to \mathcal{O}_{X'} \to 0
$$
其中把 $\mathcal{O}_{X'}$ 看成 $X$ 上的拟凝聚模。由
引理 \ref{more-morphisms-lemma-globally-generated-vanishing}，有
$H^1(X, \mathcal{I}\mathcal{O}_X) = 0$。因此
$$
H^0(Y, \mathcal{O}_Y) \to
H^0(Y, f_*\mathcal{O}_X) = H^0(X, \mathcal{O}_X) \to
H^0(X, \mathcal{O}_{X'})
$$
如所需是满射。第一个箭头满射，是因为 $Y$ 仿射，并且我们假设
$\mathcal{O}_Y \to f_*\mathcal{O}_X$ 满射；第二个箭头满射，则由上述
短正合列所伴随的上同调长正合列以及刚证明的消失得到。
\end{proof}

\begin{lemma}
\label{more-morphisms-lemma-h1-fibre-zero-h0-kappa}
考虑交换图表
$$
\xymatrix{
X \ar[rr]_f \ar[rd] & & Y \ar[ld] \\
& S
}
$$
为概形态射的交换图。设 $s \in S$ 是一点。假设
\begin{enumerate}
\item $X \to S$ 局部有限表示，并且在 $X_s$ 的各点处平坦；
\item $f$ 固有；
\item $f_s : X_s \to Y_s$ 的各纤维维数 $\leq 1$，并且
$R^1f_{s, *}\mathcal{O}_{X_s} = 0$；
\item $\mathcal{O}_{Y_s} \to f_{s, *}\mathcal{O}_{X_s}$ 满射。
\end{enumerate}
则存在开集 $Y_s \subset V \subset Y$，使得
(a) $f^{-1}(V)$ 在 $S$ 上平坦，
(b) 对 $y \in V$ 有 $\dim(X_y) \leq 1$，
(c) $R^1f_*\mathcal{O}_X|_V = 0$，
(d) $\mathcal{O}_V \to f_*\mathcal{O}_X|_V$ 满射，
并且沿任意 $Y' \to V$ 基变换后，(b)、(c)、(d) 仍成立。
\end{lemma}

\begin{proof}
令 $y \in Y$ 为位于 $s$ 上的一点。只需找到 $y$ 的一个具有所需性质的
开邻域。第一步，用 引理 \ref{more-morphisms-lemma-h1-fibre-zero} 中找到的开集 $V$
代替 $Y$，使得 $R^1f_*\mathcal{O}_X$ 普遍为零（该引理的假设由
引理 \ref{more-morphisms-lemma-check-h1-fibre-zero} 成立）。还可缩小 $Y$，使得 $f$
的所有纤维维数均 $\leq 1$（使用态射，Lemma
\ref{morphisms-lemma-openness-bounded-dimension-fibres} 以及 $f$ 的固有性）。
因此可以假设在 $V = Y$ 时 (b)、(c) 成立，并且沿任意
$Y' \to Y$ 基变换后仍成立。于是由
引理 \ref{more-morphisms-lemma-h1-fibre-zero-check-h0-kappa}，只需在 $Y$ 上证明 (d)。
仍可进一步缩小 $Y$；例如，可以并且确实假设 $Y$ 和 $S$ 都仿射。

\medskip\noindent
由 定理 \ref{more-morphisms-theorem-openness-flatness}，存在开子集 $U \subset X$，
$X \to S$ 在其上平坦；按假设，该开集包含 $X_s$。于是
$f(X \setminus U)$ 是不含 $y$ 的闭子集。因此缩小 $Y$ 后，
可以假设 $X$ 在 $S$ 上平坦。

\medskip\noindent
设 $S = \Spec(R)$。选择闭浸入 $Y \to Y'$，其中 $Y'$ 是集合 $E$
上的多项式环 $R[x_e; e \in E]$ 的谱。用 $f' : X \to Y'$ 表示 $f$
与 $Y \to Y'$ 的复合。于是对 $f'$ 和 $s$，假设 (1)--(4) 以及
(b)、(c) 均成立。若能证明 $\mathcal{O}_{Y'} \to f'_*\mathcal{O}_X$
在 $y$ 的某个开邻域中满射，则 $\mathcal{O}_Y \to f_*\mathcal{O}_X$
也如此。因此可以假设 $Y$ 是 $R[x_e; e \in E]$ 的谱。

\medskip\noindent
此时 $X$ 和 $Y$ 都在 $S$ 上平坦。于是 $Y_s$ 和 $X$ 在 $Y$ 上
Tor 独立。我们鼓励读者自己寻找证明；它也可由
引理 \ref{more-morphisms-lemma-case-of-tor-independence} 推出，只需把该引理应用于
顶点为 $X, Y, S, S$ 的方块及其沿 $s \to S$ 的基变换。因此
$$
Rf_{s, *}\mathcal{O}_{X_s} = L(Y_s \to Y)^*Rf_*\mathcal{O}_X
$$
这是由概形的导出范畴，Lemma
\ref{perfect-lemma-compare-base-change} 得到的。由已经建立的消失，它蕴含
$f_{s, *}\mathcal{O}_{X_s} = (Y_s \to Y)^*f_*\mathcal{O}_X$。
因此 $\mathcal{O}_Y \to f_*\mathcal{O}_X$ 是拟凝聚
$\mathcal{O}_Y$-模的映射，其到 $Y_s$ 的拉回满射。我们断言
$f_*\mathcal{O}_X$ 是有限型 $\mathcal{O}_Y$-模。若此断言成立，则
$\mathcal{O}_Y \to f_*\mathcal{O}_X$ 的余核 $\mathcal{F}$ 是有限型拟凝聚
$\mathcal{O}_Y$-模，并且 $\mathcal{F}_y \otimes \kappa(y) = 0$。
由 Nakayama 引理（代数，引理 \ref{algebra-lemma-NAK}），
$\mathcal{F}_y = 0$。因此 $\mathcal{F}$ 在 $y$ 的某个开邻域中为零
（模，引理 \ref{modules-lemma-finite-type-stalk-zero}），证明完成。

\medskip\noindent
证明该断言。对有限子集 $E' \subset E$，置
$Y' = \Spec(R[x_e; e \in E'])$。当 $E'$ 足够大时，态射
$f' : X \to Y \to Y'$ 固有，见极限，Lemma
\ref{limits-lemma-finite-type-eventually-proper}。以下固定 $E'$ 和 $Y'$。
把 $R = \colim R_i$ 写成其有限型 $\mathbf{Z}$-子代数的余极限。置
$S_i = \Spec(R_i)$ 和 $Y'_i = \Spec(R_i[x_e; e \in E'])$。
当 $i$ 足够大时，可以找到图
$$
\xymatrix{
X \ar[d] \ar[r]_{f'} & Y' \ar[d] \ar[r] & S \ar[d] \\
X_i \ar[r]^{f'_i} & Y'_i \ar[r] & S_i
}
$$
其中两个方块均为笛卡尔方块，$X_i$ 在 $S_i$ 上平坦，且
$X_i \to Y'_i$ 固有。见极限，Lemmas
\ref{limits-lemma-descend-finite-presentation}、
\ref{limits-lemma-descend-flat-finite-presentation} 和
\ref{limits-lemma-eventually-proper}。与上面相同的论证表明，$Y'$ 和
$X_i$ 在 $Y'_i$ 上 Tor 独立，因此
$$
R\Gamma(X, \mathcal{O}_X) =
R\Gamma(X_i, \mathcal{O}_{X_i})
\otimes^\mathbf{L}_{R_i[x_e; e \in E']}
R[x_e; e \in E']
$$
这里使用与上面相同的参考结果。由概形的上同调，Lemma
\ref{coherent-lemma-proper-over-affine-cohomology-finite}，复形
$R\Gamma(X_i, \mathcal{O}_{X_i})$ 在诺特环 $R_i[x_e; e \in E']$
的导出范畴中伪凝聚（见代数进阶，Lemma
\ref{more-algebra-lemma-Noetherian-pseudo-coherent}）。因此
$R\Gamma(X, \mathcal{O}_X)$ 在 $R[x_e; e \in E']$ 的导出范畴中伪凝聚，
见代数进阶，引理 \ref{more-algebra-lemma-pull-pseudo-coherent}。
由于唯一不消失的上同调模是 $H^0(X, \mathcal{O}_X)$，可知它是有限
$R[x_e; e \in E']$-模，见代数进阶，Lemma
\ref{more-algebra-lemma-n-pseudo-module}。证明完成。
\end{proof}

\begin{lemma}
\label{more-morphisms-lemma-h1-fibre-zero-h0-flat}
考虑交换图表
$$
\xymatrix{
X \ar[rr]_f \ar[rd] & & Y \ar[ld] \\
& S
}
$$
为概形态射的交换图。假设 $X \to S$ 平坦，$f$ 固有，对 $y \in Y$
有 $\dim(X_y) \leq 1$，并且 $R^1f_*\mathcal{O}_X = 0$。
则 $f_*\mathcal{O}_X$ 在 $S$ 上平坦，且 $f_*\mathcal{O}_X$
的形成与任意基变换 $S' \to S$ 相容。
\end{lemma}

\begin{proof}
可以假设 $Y$ 和 $S$ 仿射，设 $S = \Spec(A)$。为证明拟凝聚
$\mathcal{O}_Y$-模 $f_*\mathcal{O}_X$ 相对于 $S$ 平坦，只需证明
$H^0(X, \mathcal{O}_X)$ 在 $A$ 上平坦（略去一些细节）。由
引理 \ref{more-morphisms-lemma-globally-generated-vanishing}，对每个 $A$-模 $M$ 均有
$H^1(X, \mathcal{O}_X \otimes_A M) = 0$。又因 $\mathcal{O}_X$ 在 $A$
上平坦，可知函子 $M \mapsto H^0(X, \mathcal{O}_X \otimes_A M)$ 正合。
此外，由上同调，Lemma
\ref{cohomology-lemma-quasi-separated-cohomology-colimit}，该函子与直和相容。
于是容易验证，关于 $M$ 函子性地有
$H^0(X, \mathcal{O}_X \otimes_A M) = M \otimes_A H^0(X, \mathcal{O}_X)$；
这就给出所需的平坦性。最后，若 $S' \to S$ 是由环映射
$A \to A'$ 给出的仿射概形态射，则在刚讨论的仿射情形中有
$$
H^0(X \times_S S', \mathcal{O}_{X \times_S S'}) =
H^0(X, \mathcal{O}_X \otimes_A A') = H^0(X, \mathcal{O}_X) \otimes_A A'
$$
这表明 $f_*\mathcal{O}_X$ 的形成与任意基变换 $S' \to S$ 相容。
略去一些细节。
\end{proof}

\begin{lemma}
\label{more-morphisms-lemma-h1-fibre-zero-isom}
考虑交换图表
$$
\xymatrix{
X \ar[rr]_f \ar[rd] & & Y \ar[ld] \\
& S
}
$$
为概形态射的交换图。设 $s \in S$ 是一点。假设
\begin{enumerate}
\item $X \to S$ 局部有限表示，并且在 $X_s$ 的各点处平坦；
\item $Y \to S$ 局部有限表示；
\item $f$ 固有；
\item $f_s : X_s \to Y_s$ 的各纤维维数 $\leq 1$，并且
$R^1f_{s, *}\mathcal{O}_{X_s} = 0$；
\item $\mathcal{O}_{Y_s} \to f_{s, *}\mathcal{O}_{X_s}$ 是同构。
\end{enumerate}
则存在开集 $Y_s \subset V \subset Y$，使得
(a) $V$ 在 $S$ 上平坦，
(b) $f^{-1}(V)$ 在 $S$ 上平坦，
(c) 对 $y \in V$ 有 $\dim(X_y) \leq 1$，
(d) $R^1f_*\mathcal{O}_X|_V = 0$，
(e) $\mathcal{O}_V \to f_*\mathcal{O}_X|_V$ 是同构，
并且沿任意 $S' \to S$ 对 $f^{-1}(V) \to V$ 作基变换后，
(a)--(e) 仍成立。
\end{lemma}

\begin{proof}
令 $y \in Y_s$。总可以用 $y$ 的开邻域代替 $Y$。因此可以假设
$Y$ 和 $S$ 仿射。还可以假设 $X$ 在 $S$ 上平坦，对 $y \in Y$ 有
$\dim(X_y) \leq 1$，$R^1f_*\mathcal{O}_X = 0$ 普遍成立，并且
$\mathcal{O}_Y \to f_*\mathcal{O}_X$ 满射，见
引理 \ref{more-morphisms-lemma-h1-fibre-zero-h0-kappa}。（我们不会用到所有这些条件。）

\medskip\noindent
假设 $S$ 和 $Y$ 仿射。把 $S = \lim S_i$ 写成仿射诺特概形 $S_i$
的余滤逆极限。由极限，Lemma
\ref{limits-lemma-descend-finite-presentation}，存在元素 $0 \in I$ 以及图
$$
\xymatrix{
X_0 \ar[rr]_{f_0} \ar[rd] & & Y_0 \ar[ld] \\
& S_0
}
$$
它由概形的有限型态射组成，并且到 $S$ 的基变换是本引理中的图。
增大 $0$ 后，可以假设 $Y_0$ 仿射，且 $X_0 \to S_0$ 固有，见极限，
引理 \ref{limits-lemma-eventually-proper} 和
\ref{limits-lemma-limit-affine}。令 $s_0 \in S_0$ 为 $s$ 的像。
由于 $Y_s$ 仿射，$R^1f_{s, *}\mathcal{O}_{X_s} = 0$ 等价于
$H^1(X_s, \mathcal{O}_{X_s}) = 0$。由于 $X_s$ 是 $X_{0, s_0}$
沿忠实平坦映射 $\kappa(s_0) \to \kappa(s)$ 的基变换，可知
$H^1(X_{0, s_0}, \mathcal{O}_{X_{0, s_0}}) = 0$，因而
% P05-EN-0155: frozen source applies f_0 directly to the fibre X_{0,s_0}; the restricted morphism should be indexed f_{0,s_0}.
$R^1f_{0, *}\mathcal{O}_{X_{0, s_0}} = 0$。类似地，由于
$\mathcal{O}_{Y_s} \to f_{s, *}\mathcal{O}_{X_s}$ 是同构，
$\mathcal{O}_{Y_{0, s_0}} \to f_{0, *}\mathcal{O}_{X_{0, s_0}}$ 也是同构。
由于 $X_s \to Y_s$ 的各纤维维数至多为 $1$，对态射
$X_{0, s_0} \to Y_{0, s_0}$ 也如此。最后，由于 $X \to S$ 平坦，
增大 $0$ 后，可以假设 $X_0$ 在 $S_0$ 上平坦，见极限，Lemma
\ref{limits-lemma-descend-flat-finite-presentation}。因此只需对
$X_0 \to Y_0 \to S_0$ 和点 $s_0$ 证明本引理。

\medskip\noindent
合并上述归约论证，我们归约到如下情形：$S$ 和 $Y$ 仿射，$S$ 诺特，
$f$ 的各纤维维数 $\leq 1$，并且 $R^1f_*\mathcal{O}_X = 0$ 普遍成立。
设 $y \in Y_s$ 是一点。断言：
$$
\mathcal{O}_{Y, y} \longrightarrow (f_*\mathcal{O}_X)_y
$$
是同构。该断言蕴含本引理。事实上，$f_*\mathcal{O}_X$ 是凝聚的
（概形的上同调，Proposition
\ref{coherent-proposition-proper-pushforward-coherent}），故该断言表明，
可以用 $y$ 的开邻域代替 $Y$，并得到同构
$\mathcal{O}_Y \to f_*\mathcal{O}_X$。于是由
引理 \ref{more-morphisms-lemma-h1-fibre-zero-h0-flat}，$Y$ 在 $S$ 上平坦。
最后，由 引理 \ref{more-morphisms-lemma-h1-fibre-zero-h0-flat} 的最后一个断言，
同构 $\mathcal{O}_Y \to f_*\mathcal{O}_X$ 沿任意基变换 $S' \to S$
后仍为同构。

\medskip\noindent
证明该断言。我们已经知道
$\mathcal{O}_{Y, y} \longrightarrow (f_*\mathcal{O}_X)_y$ 满射
（引理 \ref{more-morphisms-lemma-h1-fibre-zero-h0-kappa}），并且
$(f_*\mathcal{O}_X)_y$ 是 $\mathcal{O}_{S, s}$-平坦的
（引理 \ref{more-morphisms-lemma-h1-fibre-zero-h0-flat}），而且所诱导的映射
$$
\mathcal{O}_{Y_s, y} =  \mathcal{O}_{Y, y}/\mathfrak m_s\mathcal{O}_{Y, y}
\longrightarrow
(f_*\mathcal{O}_X)_y/\mathfrak m_s (f_*\mathcal{O}_X)_y
\to
(f_{s, *}\mathcal{O}_{X_s})_y
$$
由本引理的假设是单射。于是由代数，Lemma
\ref{algebra-lemma-mod-injective}，
$\mathcal{O}_{Y, y} \longrightarrow (f_*\mathcal{O}_X)_y$ 是单射，
正合所需。
\end{proof}

\begin{lemma}
\label{more-morphisms-lemma-bijection-on-Pic}
设 $f : X \to Y$ 是诺特概形的固有态射，满足
$f_*\mathcal{O}_X = \mathcal{O}_Y$，$f$ 的各纤维维数 $\leq 1$，
并且对 $y \in Y$ 有 $H^1(X_y, \mathcal{O}_{X_y}) = 0$。
则 $f^* : \Pic(Y) \to \Pic(X)$ 双射到由如下元素组成的子群：
$\mathcal{L} \in \Pic(X)$，并且对所有 $y \in Y$ 都有
$\mathcal{L}|_{X_y} \cong \mathcal{O}_{X_y}$。
\end{lemma}

\begin{proof}
由投影公式（上同调，引理 \ref{cohomology-lemma-projection-formula}），
对 $\mathcal{N} \in \Pic(Y)$ 有 $f_*f^*\mathcal{N} \cong \mathcal{N}$。
我们断言，若 $\mathcal{L} \in \Pic(X)$ 满足对所有 $y \in Y$ 都有
$\mathcal{L}|_{X_y} \cong \mathcal{O}_{X_y}$，则
$\mathcal{N} = f_*\mathcal{L}$ 可逆，并且
$\mathcal{L} \cong f^*\mathcal{N}$。这将完成证明。

\medskip\noindent
由概形的上同调，Proposition
\ref{coherent-proposition-proper-pushforward-coherent}，
$\mathcal{O}_Y$-模 $\mathcal{N} = f_*\mathcal{L}$ 是凝聚的。
因此为证明它是可逆 $\mathcal{O}_Y$-模，只需在茎上检验（代数，Lemma
\ref{algebra-lemma-finite-projective}）。由于诺特局部环到其完备化的映射
忠实平坦，只需检验完备化 $(f_*\mathcal{L})_y^\wedge$ 自由（见代数，
第 \ref{algebra-section-completion-noetherian} 节 和
引理 \ref{algebra-lemma-finite-projective-descends}）。为此，将使用概形的上同调，
引理 \ref{coherent-lemma-formal-functions-stalk} 所表述的形式函数定理。
由于 $f_*\mathcal{O}_X = \mathcal{O}_Y$，从而
$(f_*\mathcal{O}_X)_y^\wedge \cong \mathcal{O}_{Y, y}^\wedge$，
只需证明对每个 $n$ 都有 $\mathcal{L}|_{X_n} \cong \mathcal{O}_{X_n}$
% P05-EN-0156: frozen source opens but does not close the parenthetical "compatibly for varying n".
（并且关于变化的 $n$ 相容。
由 引理 \ref{more-morphisms-lemma-picard-group-first-order-thickening}，有正合列
$$
H^1(X_y, \mathfrak m_y^n\mathcal{O}_X/\mathfrak m_y^{n + 1}\mathcal{O}_X)
\to \Pic(X_{n + 1}) \to \Pic(X_n)
$$
记号与形式函数定理中相同。注意有满射
$$
\mathcal{O}_{X_y}^{\oplus r_n} \cong
\mathfrak m_y^n/\mathfrak m_y^{n + 1} \otimes_{\kappa(y)} \mathcal{O}_{X_y}
\longrightarrow
\mathfrak m_y^n\mathcal{O}_X/\mathfrak m_y^{n + 1}\mathcal{O}_X
$$
其中 $r_n \geq 0$ 是某个整数。由于 $\dim(X_y) \leq 1$，该满射诱导
第一上同调群之间的满射（由上同调，Proposition
\ref{cohomology-proposition-vanishing-Noetherian}，次数 $\geq 2$
的上同调消失）。因此正合列中的 $H^1$ 为零，而转移映射
$\Pic(X_{n + 1}) \to \Pic(X_n)$ 如所需是单射。

\medskip\noindent
还需证明 $f^*\mathcal{N} \cong \mathcal{L}$。这可用相同方法证明，
略去细节。
\end{proof}









\section{仿射分层}
\label{more-morphisms-section-affine-stratifications}

\noindent
本节材料取自 \cite{RV}。阅读本节之前，请先阅读拓扑，Section
\ref{topology-section-stratifications} 中关于分层的一些内容。

\medskip\noindent
若 $X$ 是概形，则 $X$ 的分层通常是指 $X$ 的底层拓扑空间的分层。
各层是局部闭子集。使用概形，Remark
\ref{schemes-remark-reduced-induced-locally-closed}，我们把这些层看作
$X$ 的约化局部闭子概形。

\begin{definition}
\label{more-morphisms-definition-affine-stratification}
设 $X$ 是概形。{\it 仿射分层}是局部有限分层
$X = \coprod_{i \in I} X_i$，其中各层 $X_i$ 仿射，并且嵌入态射
$X_i \to X$ 是仿射的。
\end{definition}

\noindent
在嵌入态射 $X_i \to X$ 拟紧的条件下，分层 $X = \coprod X_i$
局部有限，等价于各层是 $X$ 的局部可构造子集，见性质，Lemma
\ref{properties-lemma-stratification-locally-finite-constructible}。

\medskip\noindent
$X_i \to X$ 是仿射态射这一条件，与我们赋予局部闭子集 $X_i$
的概形结构无关，见 引理 \ref{more-morphisms-lemma-thicken-property-morphisms}。
此外，若 $X$ 分离（或更一般地，其对角态射仿射），且
$X = \coprod X_i$ 是具有仿射层的局部有限分层，则态射
$X_i \to X$ 仿射。见态射，Lemma
\ref{morphisms-lemma-affine-permanence}。因此在大多数关注的情形中，
可以不再单独考虑嵌入态射 $X_i \to X$ 的仿射性条件。

\medskip\noindent
我们通常关注分层的偏序指标集 $I$ 有限的情形。回忆，偏序集 $I$
的{\it 长度}是 $I$ 的元素链 $i_0 < i_1 < \ldots < i_p$ 的长度 $p$
的上确界。

\begin{lemma}
\label{more-morphisms-lemma-improve-stratification}
设 $X$ 是概形，$X = \coprod_{i \in I} X_i$ 是有限仿射分层。
则存在指标集为 $\{0, \ldots, n\}$ 的仿射分层，其中 $n$ 是 $I$ 的长度。
\end{lemma}

\begin{proof}
回忆，$I$ 上有偏序，使得对所有 $i \in I$，$X_i$ 的闭包包含在
$\bigcup_{j \leq i} X_j$ 中。令 $I' \subset I$ 为 $I$ 的极大指标所成的集合。
若 $i \in I'$，则 $X_i$ 在 $X$ 中开，因为其余各层的闭包之并正是
$X_i$ 的补集。把 $U = \bigcup_{i \in I'} X_i$ 看作 $X$ 的开子概形，
于是作为概形有 $U_{red} = \coprod_{i \in I'} X_i$。由概形，Lemma
\ref{schemes-lemma-disjoint-union-affines} 和
引理 \ref{more-morphisms-lemma-thickening-affine-scheme}，$U$ 是仿射概形。态射
$U \to X$ 是仿射的，因为使用 Lemma
\ref{more-morphisms-lemma-thicken-property-morphisms} 作同样的论证可知，每个
$X_i \to X$（$i \in I'$）都是仿射的。给补集 $Z = X \setminus U$
赋予约化诱导概形结构，则它有仿射分层
$Z = \bigcup_{i \in I \setminus I'} X_i$。这里使用了如下事实：
概形态射 $T \to Z$ 仿射，当且仅当复合 $T \to X$ 仿射；这由态射，
引理 \ref{morphisms-lemma-closed-immersion-affine}、
\ref{morphisms-lemma-composition-affine} 和
\ref{morphisms-lemma-affine-permanence} 得出。注意，偏序集
$I \setminus I'$ 的长度恰好比 $I$ 的长度小一。因此由归纳可知，$Z$
有仿射分层 $Z = Z_0 \amalg \ldots \amalg Z_{n - 1}$，其指标集为
% P05-EN-0157: frozen source labels strata Z_0 through Z_{n-1} but states the index set is {1,...,n}; it should be {0,...,n-1}.
$\{1, \ldots, n\}$。置 $Z_n = U$，即得 $X$ 的所需分层。
\end{proof}

\noindent
若概形 $X$ 有有限仿射分层，则 $X$ 当然拟紧。稍不明显的是，
这还迫使 $X$ 拟分离。

\begin{lemma}
\label{more-morphisms-lemma-qc-affine-stratification}
设 $X$ 是概形。下列条件等价：
\begin{enumerate}
\item $X$ 有有限仿射分层；
\item $X$ 拟紧且拟分离。
\end{enumerate}
\end{lemma}

\begin{proof}
设 $X = \bigcup X_i$ 是有限仿射分层。由于每个 $X_i$ 仿射，因而拟紧，
可知 $X$ 拟紧。设 $U, V \subset X$ 是仿射开集。由于 $X_i \to X$
是仿射态射，$U \cap X_i$ 和 $V \cap X_i$ 是 $X_i$ 中的仿射开集。
因此 $U \cap V \cap X_i$ 是仿射概形 $X_i$ 的仿射开集（例如见概形，
引理 \ref{schemes-lemma-characterize-separated}）。于是
$U \cap V = \coprod U \cap V \cap X_i$ 作为有限多个仿射层的并是拟紧的。
由概形，引理 \ref{schemes-lemma-characterize-quasi-separated}，
$X$ 拟分离。

\medskip\noindent
假设 $X$ 拟紧且拟分离。可以使用概形的上同调，Lemma
\ref{coherent-lemma-induction-principle} 的归纳原理，证明 $X$
有有限仿射分层。若 $X$ 为空，则它有空的仿射分层。若 $X$
是非空仿射概形，则它有仅含一层的仿射分层。接下来假设
$X = U \cup V$，其中 $U$ 是拟紧开集，$V$ 是仿射开集，并且已有有限
仿射分层 $U = \bigcup_{i \in I} U_i$ 和
$U \cap V = \coprod_{j \in J} W_j$。记 $Z = X \setminus V$ 和
$Z' = X \setminus U$。注意，$Z$ 在 $U$ 中闭，而 $Z'$ 在 $V$ 中闭。
注意，$U_i \cap Z$ 和 $U_i \cap W_j = U_i \times_U W_j$ 是在 $U$
上仿射的仿射概形。（提示：$U_i \times_U W_j \to W_j$ 是
$U_i \to U$ 的基变换，故仿射；从而 $U_i \cap W_j$ 仿射，进而
$U_i \cap W_j \to U_i$ 仿射，最后 $U_i \cap W_j \to U$ 仿射。）
因此
$$
U = \coprod\nolimits_{i \in I} (U_i \cap Z) \amalg
\coprod\nolimits_{(i, j) \in I \times J} (U_i \cap W_j)
$$
是有限仿射分层；$I \amalg I \times J$ 上的偏序由
$i' \leq (i, j) \Leftrightarrow i' \leq i$ 和
$(i', j') \leq (i, j) \Leftrightarrow i' \leq i$ 且 $j' \leq j$ 给出。
注意，$(U_i \cap Z) \times_X V = \emptyset$ 和
$(U_i \cap W_j) \times_X V = U_i \cap W_j$ 都仿射。因此态射
$U_i \cap Z \to X$ 和 $U_i \cap W_j \to X$ 仿射，因为态射的仿射性
可以在靶上局部地检验（态射，Lemma
\ref{morphisms-lemma-characterize-affine}），而我们已经有了在 $U$ 和 $V$
上的仿射性。为完成证明，在上述分层中再加入一层，即 $Z'$，
并把它的指标作为该偏序集的一个极小元。
\end{proof}

\begin{definition}
\label{more-morphisms-definition-affine-stratification-number}
设 $X$ 是非空拟紧拟分离概形。{\it 仿射分层数}是使下列等价条件成立的
最小整数 $n \geq 0$：
\begin{enumerate}
\item 存在有限仿射分层 $X = \coprod_{i \in I} X_i$，其中 $I$ 的长度为 $n$；
\item 存在仿射分层 $X = X_0 \amalg X_1 \amalg \ldots \amalg X_n$，
其指标集为 $\{0, \ldots, n\}$。
\end{enumerate}
\end{definition}

\noindent
这些条件的等价性由 引理 \ref{more-morphisms-lemma-improve-stratification} 成立。
有限仿射分层的存在性已在 Lemma
\ref{more-morphisms-lemma-qc-affine-stratification} 中证明。

\begin{lemma}
\label{more-morphisms-lemma-affine-stratification-number-bound}
设 $X$ 是分离概形，且有一个由 $n + 1$ 个仿射开集组成的开覆盖。
则 $X$ 的仿射分层数至多为 $n$。
\end{lemma}

\begin{proof}
设 $X = U_0 \cup \ldots \cup U_n$ 是仿射开覆盖。置
$$
X_i = (U_i \cup \ldots \cup U_n) \setminus
(U_{i + 1} \cup \ldots \cup U_n)
$$
于是 $X_i$ 作为 $U_i$ 的闭子概形是仿射的。由态射，Lemma
\ref{morphisms-lemma-affine-permanence}，态射 $X_i \to X$ 仿射。
最后有 $\overline{X_i} \subset X_i \cup X_{i - 1} \cup \ldots X_0$。
\end{proof}

\begin{lemma}
\label{more-morphisms-lemma-affine-stratification-number-bound-Noetherian}
设 $X$ 是维数为 $\infty > d \geq 0$ 的诺特概形。
则 $X$ 的仿射分层数至多为 $d$。
\end{lemma}

\begin{proof}
对 $d$ 归纳。若 $d = 0$，则 $X$ 仿射，见性质，Lemma
\ref{properties-lemma-locally-Noetherian-dimension-0}。假设 $d > 0$。
令 $\eta_1, \ldots, \eta_n$ 为 $X$ 的不可约分支的泛点（性质，Lemma
\ref{properties-lemma-Noetherian-irreducible-components}）。可以用包含
$\eta_1, \ldots, \eta_n$ 的仿射开集覆盖 $X$，见性质，Lemma
\ref{properties-lemma-point-and-maximal-points-affine}。由于 $X$ 拟紧，
可以找到有限仿射开覆盖 $X = \bigcup_{j = 1, \ldots, m} U_j$，
使得对所有 $j = 1, \ldots, m$ 都有
$\eta_1, \ldots, \eta_n \in U_j$。选择包含
$\eta_1, \ldots, \eta_n$ 的仿射开集
$U \subset U_1 \cap \ldots \cap U_m$（由刚引用的引理可以做到）。
于是态射 $U \to X$ 仿射，因为对所有 $j$，$U \to U_j$ 都仿射，
见态射，引理 \ref{morphisms-lemma-characterize-affine}。令
$Z = X \setminus U$。按构造，$\dim(Z) < \dim(X)$。由归纳假设，
可以找到 $Z$ 的仿射分层
$Z = \bigcup_{i \in \{0, \ldots, n\}} Z_i$，其中 $n \leq \dim(Z)$。
% P05-EN-0158: frozen source reverses the assignment and writes U = X_{n+1}; the new top stratum should be X_{n+1} = U.
置 $U = X_{n + 1}$，并对 $i \leq n$ 置 $X_i = Z_i$，即得结论。
\end{proof}

\begin{proposition}
\label{more-morphisms-proposition-vanishing-affine-stratification-number}
设 $X$ 是仿射分层数为 $n$ 的非空拟紧拟分离概形。则对每个拟凝聚
$\mathcal{O}_X$-模 $\mathcal{F}$，当 $p > n$ 时有
$H^p(X, \mathcal{F}) = 0$。
\end{proposition}

\begin{proof}
我们对仿射分层数 $n$ 归纳。若 $n = 0$，则 $X$ 仿射，结论就是
概形的上同调，Lemma
\ref{coherent-lemma-quasi-coherent-affine-cohomology-zero}。假设 $n > 0$。
由 定义 \ref{more-morphisms-definition-affine-stratification-number}，存在仿射概形
$U$ 和仿射开浸入 $j : U \to X$，使得补集 $Z$ 的仿射分层数为 $n - 1$。
由于 $U$ 和 $j$ 仿射，对 $p > 0$ 有
$H^p(X, j_*(\mathcal{F}|_U)) = 0$，见概形的上同调，Lemmas
\ref{coherent-lemma-relative-affine-cohomology} 和
\ref{coherent-lemma-relative-affine-vanishing}。用 $\mathcal{K}$ 和
$\mathcal{Q}$ 表示映射 $\mathcal{F} \to j_*(\mathcal{F}|_U)$ 的核和余核。
于是得到正合列
$$
0 \to \mathcal{K} \to \mathcal{F} \to j_*(\mathcal{F}|_U)
\to \mathcal{Q} \to 0
$$
它由拟凝聚 $\mathcal{O}_X$-模组成（见概形，Section
\ref{schemes-section-quasi-coherent}）。把该正合列分成短正合列并使用上同调
长正合列的标准论证表明，只需证明当 $p \geq n$ 时，
$H^p(X, \mathcal{K}) = 0$ 且 $H^p(X, \mathcal{Q}) = 0$。
由于 $\mathcal{F} \to j_*(\mathcal{F}|_U)$ 在 $U$ 上限制为同构，
$\mathcal{K}$ 和 $\mathcal{Q}$ 都支撑在 $Z$ 上。由性质，Lemma
\ref{properties-lemma-quasi-coherent-colimit-finite-type}，可以把这些模
写成其有限型拟凝聚子模的滤过余极限。利用 $X$ 上层的上同调与滤过余极限
相容这一事实，见上同调，Lemma
\ref{cohomology-lemma-quasi-separated-cohomology-colimit}，只需证明如下断言：
若 $\mathcal{G}$ 是支撑包含在 $Z$ 中的有限型拟凝聚模，则对
$p \geq n$ 有 $H^p(X, \mathcal{G}) = 0$。令 $Z' \subset X$ 为
$\mathcal{G} \oplus \mathcal{O}_Z$ 的概形论支撑；可以并且确实把
$\mathcal{G}$ 看成 $Z'$ 上的拟凝聚模，见态射，Section
\ref{morphisms-section-support}。于是 $Z'$ 和 $Z$ 有相同的底层拓扑空间，
从而有相同的仿射分层数，即 $n - 1$。因此由归纳假设，
$H^p(X, \mathcal{G}) = H^p(Z', \mathcal{G})$（该等式由概形的上同调，
引理 \ref{coherent-lemma-relative-affine-cohomology} 得出）在
$p \geq n$ 时为零。
\end{proof}

\begin{example}
\label{more-morphisms-example-affine-stratification-number}
设 $k$ 是域，$X = \mathbf{P}^n_k$ 是 $k$ 上的 $n$ 维射影空间。
引理 \ref{more-morphisms-lemma-affine-stratification-number-bound} 由构造，Lemma
\ref{constructions-lemma-standard-covering-projective-space} 可用于它。因此
$\mathbf{P}^n_k$ 的仿射分层数至多为 $n$。另一方面，
$\mathbf{P}^n_k$ 上的某些拟凝聚模在次数 $n$ 上有非零上同调，
见概形的上同调，Lemma
\ref{coherent-lemma-cohomology-projective-space-over-ring}。使用
命题 \ref{more-morphisms-proposition-vanishing-affine-stratification-number}，
可知 $\mathbf{P}^n_k$ 的仿射分层数等于 $n$。
\end{example}








\section{普遍开态射}
\label{more-morphisms-section-universally-open}

\noindent
下面给出关于普遍开态射的一些材料。

\begin{lemma}
\label{more-morphisms-lemma-test-universally-open}
设 $f : X \to S$ 是概形态射。下列条件等价：
\begin{enumerate}
\item $f$ 普遍开；
\item 对每个局部有限表示的态射 $S' \to S$，基变换
$X_{S'} \to S'$ 是开态射；
\item 对每个 $n$，态射 $\mathbf{A}^n \times X \to \mathbf{A}^n \times S$
都是开态射。
\end{enumerate}
\end{lemma}

\begin{proof}
显然，(1) 蕴含 (2)，而 (2) 蕴含 (3)。下面证明 (3) 蕴含 (1)。
假设对某个概形态射 $g : T \to S$，基变换 $X_T \to T$ 不是开态射。
则可以找到仿射开集 $V \subset S$、$U \subset X$、$W \subset T$，
满足 $f(U) \subset V$ 和 $g(W) \subset V$，并且
$U \times_V W \to W$ 不是开态射。若能证明这蕴含
$\mathbf{A}^n \times U \to \mathbf{A}^n \times V$ 不是开态射，
则 $\mathbf{A}^n \times X \to \mathbf{A}^n \times S$ 不是开态射，
证明完成。这把问题归约到下一段证明的结果。

\medskip\noindent
设 $A \to B$ 是环映射，并且对某个 $A$-代数 $A'$，
$A' \to B' = A' \otimes_A B$ 所诱导的谱映射不是开映射。
由于主开集构成 $\Spec(B')$ 的拓扑基，可知对某个 $g \in B'$，
$D(g)$ 在 $\Spec(A')$ 中的像不是开集。写成
$g = \sum_{i = 1, \ldots, n} a'_i \otimes b_i$，其中 $n$、
$a'_i \in A'$ 和 $b_i \in B$ 如此选取。考虑
$B[x_1, \ldots, x_n]$ 中的元素
$h = \sum_{i = 1, \ldots, n} x_i b_i$。反设 $D(h)$ 在态射
$$
\Spec(B[x_1, \ldots, x_n]) \longrightarrow \Spec(A[x_1, \ldots, x_n])
$$
下的像为开集，以期导出矛盾。此时 $D(h)$ 将满射到某个拟紧开集
$U \subset \Spec(A[x_1, \ldots, x_n])$。令
$A[x_1, \ldots, x_n] \to A'$ 为把 $x_i$ 映到 $a'_i$ 的 $A$-代数同态。
它还诱导把 $h$ 映到 $g$ 的 $B$-代数同态
$B[x_1, \ldots, x_n] \to B'$。由于
$$
\xymatrix{
\Spec(B[x_1, \ldots, x_n]) \ar[d] &
\Spec(B') \ar[l] \ar[d] \\
\Spec(A[x_1, \ldots, x_n]) &
\Spec(A') \ar[l]
}
$$
是笛卡尔图，$D(g)$ 在 $\Spec(A')$ 中的像等于 $U$ 在 $\Spec(A')$
中的原像，因而是开集。这就是所需的矛盾。
\end{proof}

\begin{lemma}
\label{more-morphisms-lemma-quasi-finite-Noetherian-universally-open}
设 $f : X \to Y$ 是概形态射。若
\begin{enumerate}
\item $f$ 局部拟有限；
\item $Y$ 几何单分支且局部诺特；
\item $X$ 的每个不可约分支都支配 $Y$ 的某个不可约分支，
\end{enumerate}
则 $f$ 普遍开。
\end{lemma}

\begin{proof}
对任意 $n$，由 引理 \ref{more-morphisms-lemma-number-of-branches-and-smooth} 和性质，
引理 \ref{properties-lemma-number-of-branches-1}，概形
$\mathbf{A}^n \times Y$ 几何单分支。因此对所有 $n$，态射
$\mathbf{A}^n \times X \to \mathbf{A}^n \times Y$ 满足本引理的假设。
由 引理 \ref{more-morphisms-lemma-test-universally-open}，只需证明 $f$ 是开态射。
由态射，Lemma
\ref{morphisms-lemma-locally-finite-presentation-universally-open}，
只需证明一般化可以沿 $f$ 提升。假设 $y' \leadsto y$ 是 $Y$
中点的特化，而 $x \in X$ 映到 $y$。按照
引理 \ref{more-morphisms-lemma-etale-makes-quasi-finite-finite-at-point} 选择图
$$
\xymatrix{
u \ar[d] & U \ar[d] \ar[r] & X \ar[d] \\
v & V \ar[r] & Y
}
$$
其中 $(V, v) \to (Y, y)$ 是初等 \'etale 邻域，$U \to V$ 有限，
$u$ 是 $U$ 中唯一映到 $v$ 的点，$U \subset V \times_Y X$ 是开集，
并且 $v \mapsto y$、$u \mapsto x$。令 $E$ 为 $U$ 的一个经过 $u$
的不可约分支（至少存在一个）。由于 $U \to X$ 是 \'etale 的，
$E$ 映到 $X$ 的某个不可约分支，而后者又支配 $Y$ 的某个不可约分支
（由假设）。由于 $U \to V$ 有限，因而闭，$E$ 的像 $E' \subset V$
是经过 $v$ 的不可约闭子集，并支配 $Y$ 的某个不可约分支。由于
$V \to Y$ 是 \'etale 的，$E'$ 必为 $V$ 的一个经过 $v$ 的不可约分支。
由于 $Y$ 几何单分支，$E'$ 是 $V$ 的唯一一个经过 $v$ 的不可约分支
（引理 \ref{more-morphisms-lemma-nr-branches}）。由于 $V$ 局部诺特，缩小 $V$ 后，
可以假设 $E' = V$（集合相等）。

\medskip\noindent
由于 $V \to Y$ 是 \'etale 的，可以找到特化 $v' \leadsto v$，
其像为 $y' \leadsto y$。由上可找到映到 $v'$ 的 $u' \in U$。
于是 $u' \leadsto u$，因为 $u$ 是 $U$ 中唯一映到 $v$ 的点，
且 $U \to V$ 是闭态射。最后，$u'$ 的像 $x' \in X$ 特化到 $x$，
并映到 $y'$。证明完成。
\end{proof}

\begin{lemma}
\label{more-morphisms-lemma-characterize-universally-open-finite}
设 $A \to B$ 是环映射，并设 $B$ 作为 $A$-模由
$b_1, \ldots, b_d \in B$ 生成。置
$h = \sum x_ib_i \in B[x_1, \ldots, x_d]$。则
$\Spec(B) \to \Spec(A)$ 普遍开，当且仅当 $D(h)$ 在
$\Spec(A[x_1, \ldots, x_d])$ 中的像为开集。
\end{lemma}

\begin{proof}
若 $\Spec(B) \to \Spec(A)$ 普遍开，则 $D(h)$ 的像当然为开集。
反过来，假设 $D(h)$ 的像 $U$ 为开集。设 $A \to A'$ 是环映射。
只需证明任一主开集 $D(g) \subset \Spec(A' \otimes_A B)$ 在
$\Spec(A')$ 中的像为开集。可以写成
$g = \sum_{i = 1, \ldots, d} a'_i \otimes b_i$，其中
$a'_i \in A'$。令
% P05-EN-0159: frozen source switches the generator count d to an undefined n in both polynomial rings of this proof.
$A[x_1, \ldots, x_n] \to A'$ 为把 $x_i$ 映到 $a'_i$ 的 $A$-代数同态。
它还诱导把 $h$ 映到 $g$ 的 $B$-代数同态
$B[x_1, \ldots, x_n] \to A' \otimes_A B$。由于
$$
\xymatrix{
\Spec(B[x_1, \ldots, x_n]) \ar[d] &
\Spec(B') \ar[l] \ar[d] \\
\Spec(A[x_1, \ldots, x_n]) &
\Spec(A') \ar[l]
}
$$
% P05-EN-0160: frozen source uses Spec(B') in this diagram without defining B' in the lemma; the base-changed ring is A' tensor_A B.
是笛卡尔图，$D(g)$ 在 $\Spec(A')$ 中的像等于 $U$ 在 $\Spec(A')$
中的原像，因而为开集。
\end{proof}

\begin{lemma}
\label{more-morphisms-lemma-descend-quasi-finite-universally-open}
设 $S = \lim S_i$ 是转移态射均仿射的概形有向系的极限。设
$0 \in I$，并设 $f_0 : X_0 \to Y_0$ 是 $S_0$ 上的概形态射。
假设 $S_0$、$X_0$、$Y_0$ 拟紧且拟分离。令
$f_i : X_i \to Y_i$ 为 $f_0$ 到 $S_i$ 的基变换，并令
$f : X \to Y$ 为 $f_0$ 到 $S$ 的基变换。若
\begin{enumerate}
\item $f$ 局部拟有限且普遍开；
\item $f_0$ 局部有限表示，
\end{enumerate}
则存在 $i \geq 0$，使得 $f_i$ 局部拟有限且普遍开。
\end{lemma}

\begin{proof}
由极限，引理 \ref{limits-lemma-descend-quasi-finite}，增大 $0$ 后，
可以假设 $f_0$ 局部拟有限。令 $x \in X$。通过拟有限态射的 \'etale
局部化，可以找到图
$$
\xymatrix{
X \ar[d] & U \ar[l] \ar[d] \\
Y & V \ar[l]
}
$$
其中 $V \to Y$ 是 \'etale 的，$U \subset X_V$ 是开集，$U \to V$ 有限，
且 $x$ 位于 $U \to X$ 的像中，见
引理 \ref{more-morphisms-lemma-etale-makes-quasi-finite-finite-at-point}。缩小 $V$ 后，
可以假设 $V$ 和 $U$ 仿射。由于 $X$ 拟紧，取有限多个这样的 $V$ 和 $U$
的不交并，便可使上述图中的 $U \to X$ 满射。由极限，Lemmas
\ref{limits-lemma-descend-finite-presentation}、
\ref{limits-lemma-descend-opens}、
\ref{limits-lemma-descend-surjective}、
\ref{limits-lemma-descend-finite-finite-presentation}、
\ref{limits-lemma-descend-etale} 和
\ref{limits-lemma-limit-affine}，必要时增大 $0$ 后，可以假设有图
$$
\xymatrix{
X_0 \ar[d] & U_0 \ar[l] \ar[d] \\
Y_0 & V_0 \ar[l]
}
$$
其中 $V_0$ 仿射，$V_0 \to Y_0$ 是 \'etale 的，
$U_0 \subset (X_0)_{V_0}$ 是开集，$U_0 \to V_0$ 有限，且
$U_0 \to X_0$ 满射。由于 $V_i \to Y_i$ 是 \'etale 的，因而普遍开，
只需证明当 $i$ 足够大时 $U_i \to V_i$ 普遍开。这把我们归约到
下一段讨论的情形。

\medskip\noindent
设 $A = \colim A_i$ 是环的滤过余极限，$A_0 \to B_0$ 是环映射。
置 $B = A \otimes_{A_0} B_0$ 和 $B_i = A_i \otimes_{A_0} B_0$。
假设 $A_0 \to B_0$ 有限且有限表示，并且 $A \to B$ 普遍开。
我们要证明当 $i$ 足够大时 $A_i \to B_i$ 普遍开。选择生成 $B_0$
这一 $A_0$-模的元素 $b_{0, 1}, \ldots, b_{0, d} \in B_0$。
在 $B_0[x_1, \ldots, x_d]$ 中置
$h_0 = \sum_{j = 1, \ldots, d} x_jb_{0, j}$。分别用 $h$ 和 $h_i$
表示 $h_0$ 在 $B[x_1, \ldots, x_d]$ 和
$B_i[x_1, \ldots, x_d]$ 中的像。由于 $A \to B$ 普遍开，$D(h)$
在 $\Spec(A[x_1, \ldots, x_d])$ 中的像 $U$ 是开集。由于它是仿射概形的像，
$U$ 当然拟紧。当 $i$ 足够大时，存在拟紧开集
$U_i \subset \Spec(A_i[x_1, \ldots, x_d])$，其在
$\Spec(A[x_1, \ldots, x_d])$ 中的原像为 $U$，见极限，Lemma
\ref{limits-lemma-descend-opens}。增大 $i$ 后，可以假设 $D(h_i)$ 映入
$U_i$；这可把同一引理应用于 $U_i$ 在 $D(h_i)$ 中的拉回得到。
最后，当 $i$ 更大时，概形态射 $D(h_i) \to U_i$ 将满射，
这是再次应用极限，引理 \ref{limits-lemma-descend-surjective} 的结果。
由 引理 \ref{more-morphisms-lemma-characterize-universally-open-finite}，可知
$A_i \to B_i$ 普遍开。
\end{proof}

\begin{lemma}
\label{more-morphisms-lemma-count-geometric-fibres}
设 $f : X \to Y$ 是局部拟有限态射。则
\begin{enumerate}
\item Lemmas
\ref{more-morphisms-lemma-base-change-fibres-nr-geometrically-irreducible-components} 和
\ref{more-morphisms-lemma-base-change-fibres-nr-geometrically-connected-components}
中的函数 $n_{X/Y}$ 相同；
\item 若 $X$ 拟紧，则 $n_{X/Y}$ 取得最大值 $d < \infty$。
\end{enumerate}
\end{lemma}

\begin{proof}
局部拟有限态射的（几何）纤维是离散的，故两个函数相同，见态射，Lemma
\ref{morphisms-lemma-locally-quasi-finite-fibres}。有界性由态射，Lemmas
\ref{morphisms-lemma-characterize-universally-bounded} 和
\ref{morphisms-lemma-locally-quasi-finite-qc-source-universally-bounded} 得出。
\end{proof}

\begin{lemma}
\label{more-morphisms-lemma-count-geometric-fibres-universally-open}
设 $f : X \to Y$ 是分离、局部拟有限且普遍开的概形态射。设
$n_{X/Y}$ 如 引理 \ref{more-morphisms-lemma-count-geometric-fibres} 中所述。
若对某个 $y \in Y$ 和 $d \geq 0$ 有 $n_{X/Y}(y) \geq d$，
则在 $y$ 的某个开邻域中有 $n_{X/Y} \geq d$。
\end{lemma}

\begin{proof}
问题在 $Y$ 上是局部的，故可以假设 $Y$ 仿射。令 $K$ 为剩余域
$\kappa(y)$ 的代数闭包。假设表示 $(X_y)_K$ 有 $\geq d$ 个连通分支。
于是对某个合适的拟紧开集 $X' \subset X$，概形 $(X'_y)_K$
有 $\geq d$ 个连通分支；略去细节。以 $X'$ 代替 $X$ 后，可以假设
$X$ 拟紧。于是 $f$ 拟有限。令 $x_1, \ldots, x_n$ 为 $X$ 中位于
$y$ 上的点。应用
引理 \ref{more-morphisms-lemma-etale-splits-off-quasi-finite-part-technical-variant}，
得到 \'etale 邻域 $(U, u) \to (Y, y)$ 以及分解
$$
U \times_Y X =
W \amalg
\ \coprod\nolimits_{i = 1, \ldots, n}
\ \coprod\nolimits_{j = 1, \ldots, m_i}
V_{i, j}
$$
如所引之处所述。注意，在这种情形下
$n_{X/Y}(y) = \sum_i m_i$；略去一些细节。由于 $f$ 普遍开，对所有
$i, j$，态射 $V_{i, j} \to U$ 都是开态射。因此缩小 $U$ 后，
可以假设对所有 $i, j$，$V_{i, j} \to U$ 都满射。这证明
$n_{U \times_Y X/U} \geq \sum_i m_i = n_{X/Y}(y) \geq d$.
由于 $n_{X/Y}$ 的构造与基变换相容，证明完成。
\end{proof}

\begin{lemma}
\label{more-morphisms-lemma-large-open}
设 $f : X \to Y$ 是分离、局部拟有限且普遍开的概形态射。设
$n_{X/Y}$ 如 引理 \ref{more-morphisms-lemma-count-geometric-fibres} 中所述。
若 $n_{X/Y}$ 取得最大值 $d < \infty$，则集合
$$
Y_d = \{y \in Y \mid n_{X/Y}(y) = d\}
$$
在 $Y$ 中开，并且态射 $f^{-1}(Y_d) \to Y_d$ 有限。
\end{lemma}

\begin{proof}
$Y_d$ 的开性直接由
引理 \ref{more-morphisms-lemma-count-geometric-fibres-universally-open} 得出。
为证明在 $Y_d$ 上有限，重复该引理证明中的论证。具体地，令
$y \in Y_d$。则 $X$ 中至多有 $d$ 个点位于 $y$ 上。设
$x_1, \ldots, x_n$ 为 $X$ 中位于 $y$ 上的点。应用
引理 \ref{more-morphisms-lemma-etale-splits-off-quasi-finite-part-technical-variant}，
得到 \'etale 邻域 $(U, u) \to (Y, y)$ 以及分解
$$
U \times_Y X =
W \amalg
\ \coprod\nolimits_{i = 1, \ldots, n}
\ \coprod\nolimits_{j = 1, \ldots, m_i}
V_{i, j}
$$
如所引之处所述。注意，在这种情形下
$d = n_{X/Y}(y) = \sum_i m_i$；略去一些细节。由于 $f$ 普遍开，
对所有 $i, j$，态射 $V_{i, j} \to U$ 都是开态射。因此缩小 $U$ 后，
可以假设对所有 $i, j$，$V_{i, j} \to U$ 都满射，并且可以假设
% P05-EN-0161: frozen source says U maps into W, although W is a summand of U times_Y X and no such map is defined; likely Y_d was intended.
$U$ 映入 $W$。这证明
$n_{U \times_Y X/U} \geq \sum_i m_i = d$。由于 $n_{X/Y}$ 的构造
与基变换相容，有 $n_{U \times_Y X/U} = d$。这意味着 $W$ 必为空，
从而 $U \times_Y X \to U$ 有限。由下降，Lemma
\ref{descent-lemma-descending-property-finite}，这蕴含 $X \to Y$
在开态射 $U \to Y$ 的像上有限。换言之，$f$ 在 $y$ 的某个开邻域上有限，
正合所需。
\end{proof}






\section{赋权}
\label{more-morphisms-section-weightings}

\noindent
本节材料取自 \cite[Exposee XVII, 6.2.4]{SGA4}。

\medskip\noindent
设 $\pi : U \to V$ 是具有有限纤维的局部拟有限概形态射。给定函数
$w : U \to \mathbf{Z}$，定义函数
$$
\textstyle{\int}_\pi w : V \longrightarrow \mathbf{Z},\quad
v \longmapsto
\sum\nolimits_{u \in U,\ \pi(u) = v} w(u) [\kappa(u) : \kappa(v)]_s
$$
注意，这些域扩张有限（态射，Lemma
\ref{morphisms-lemma-residue-field-quasi-finite}），
$[\kappa' : \kappa]_s$ 表示可分次数（域，Definition
\ref{fields-definition-insep-degree}），并且因为假设 $\pi$ 的纤维有限，
该和为有限和。还可以按如下方式计算 $\int_\pi w$ 在点 $v \in V$
处的值。选择代数闭域 $k$ 以及像为 $v$ 的态射
$\overline{v} : \Spec(k) \to V$。则有
$$
(\textstyle{\int}_\pi w)(v) =
\sum\nolimits_{\overline{u} \in U_{\overline{v}}} w(\overline{u})
$$
这里，$w(\overline{u})$ 当然表示 $w$ 在点 $\overline{u}$ 沿态射
$U_{\overline{v}} \to U$ 的像 $u$ 处的值。注意，可以把
$\overline{u} \in U_{\overline{v}}$ 看成满足
$\pi \circ \overline{u} = \overline{v}$ 的态射
$\overline{u} : \Spec(k) \to U$。事实上，由于 $U \to V$
局部拟有限且纤维有限，概形 $U_{\overline{v}}$ 是 $k$ 上有限维代数的谱，
其所有素理想都是剩余域为 $k$ 的极大理想。为验证该等式，注意由域，
引理 \ref{fields-lemma-separable-degree}，位于给定 $u$ 上的态射
$\overline{u}$ 的数目等于 $[\kappa(u) : \kappa(v)]_s$。

\begin{lemma}
\label{more-morphisms-lemma-weighting-check-after-etale-base-change}
给定笛卡尔方块
$$
\xymatrix{
U \ar[d]_\pi & U' \ar[l]^h \ar[d]^{\pi'} \\
V & V' \ar[l]_g
}
$$
其中 $\pi$ 局部拟有限且纤维有限，并给定函数
$w : U \to \mathbf{Z}$，则有
$(\int_\pi w) \circ g = \int_{\pi'} (w \circ h)$。
\end{lemma}

\begin{proof}
这直接由上面对 $\int_\pi w$ 的第二种描述得出。若从定义证明，
需使用如下事实：若 $E/F$ 是域的有限扩张，$F'/F$ 是另一个域扩张，
并把 $(E \otimes_F F')_{red} = \prod E'_i$ 写成 $F'$ 上有限域的乘积，
则有
$$
[E : F]_s = \sum [E'_i : F']_s
$$
为证明该等式，选择代数闭域扩张 $\Omega/F'$，并注意
\begin{align*}
[E : F]_s
& =
|\Mor_F(E, \Omega)| \\
& =
|\Mor_{F'}(E \otimes_F F', \Omega)| \\
& =
|\Mor_{F'}((E \otimes_F F')_{red}, \Omega)| \\
& =
\sum |\Mor_{F'}(E'_i, \Omega)| \\
& =
\sum [E'_i : F']_s
\end{align*}
这里使用了域，引理 \ref{fields-lemma-separable-degree}。
\end{proof}

\begin{definition}
\label{more-morphisms-definition-weighting}
设 $f : X \to Y$ 是局部拟有限态射。$f$ 的一个{\it 赋权}，
或一个 {\it pond\'eration}，是映射 $w : X \to \mathbf{Z}$，
使得对任意图
$$
\xymatrix{
X \ar[d]_f & U \ar[l]^h \ar[d]^\pi \\
Y & V \ar[l]_g
}
$$
若 $V \to Y$ 是 \'etale 的，$U \subset X_V$ 是开集，且 $U \to V$ 有限，
则函数 $\int_\pi (w \circ h)$ 局部常值。
\end{definition}

\noindent
当然，取 $w = 0$，便得到任意局部拟有限态射 $f$ 的一个赋权，
尽管它并不十分有趣。事实将表明，对各种用途而言，最有趣的是
{\it 正}赋权，即 $w : X \to \mathbf{Z}_{> 0}$。

\begin{lemma}
\label{more-morphisms-lemma-weighting-base-change}
设 $f : X \to Y$ 是局部拟有限态射，$w : X \to \mathbf{Z}$ 是一个赋权。
设 $f' : X' \to Y'$ 是 $f$ 沿态射 $Y' \to Y$ 的基变换。
则 $w$ 与投影 $X' \to X$ 的复合 $w' : X' \to \mathbf{Z}$ 是 $f'$ 的赋权。
\end{lemma}

\begin{proof}
考虑图表
$$
\xymatrix{
X' \ar[d]_{f'} & U' \ar[l]^{h'} \ar[d]^{\pi'} \\
Y' & V' \ar[l]_{g'}
}
$$
它是 定义 \ref{more-morphisms-definition-weighting} 中针对态射 $f'$ 的图。
对任意 $v' \in V'$，需证明 $\int_{\pi'} (w' \circ h')$ 在 $v'$
的某个开邻域中常值。由
引理 \ref{more-morphisms-lemma-weighting-check-after-etale-base-change}
（以及 \'etale 态射是开态射这一事实），可以用 $v'$ 的任意 \'etale
邻域代替 $V'$。以 $v'$ 的某个 \'etale 邻域代替 $V'$ 后，
可以假设 $U' = U'_1 \amalg \ldots \amalg U'_n$，其中每个 $U'_i$
都恰有一点 $u'_i$ 位于 $v'$ 上，并且 $\kappa(u'_i)/\kappa(v')$
纯不可分，见
引理 \ref{more-morphisms-lemma-etale-splits-off-quasi-finite-part-technical-variant}。
显然，只需证明 $\int_{U'_i \to V'} w'|_{U'_i}$ 在 $v'$
的某个邻域中常值。这把我们归约到下一段讨论的情形。

\medskip\noindent
我们有 $v' \in V'$，且 $U'$ 中恰有一点 $u'$ 位于 $v'$ 上，
并满足 $\kappa(u')/\kappa(v')$ 纯不可分。用 $x \in X$ 和 $y \in Y$
分别表示 $u'$ 和 $v'$ 的像。可以找到 \'etale 邻域
$(V, v) \to (Y, y)$ 以及开集 $U \subset X_V$，使得
$\pi : U \to V$ 有限，并且其中恰有一点 $u \in U$ 位于 $v$ 上；
该点沿投影 $h : U \to X$ 映到 $x \in X$，而且
$\kappa(u)/\kappa(v)$ 纯不可分。由上面使用的引理，这可以做到。
考虑态射
$$
U'' = U \times_X U' \longrightarrow V \times_Y V' = V''
$$
由于 $u$ 和 $u'$ 都映到 $x \in X$，存在映到 $(u, u')$ 的点
$u'' \in U''$。用 $v'' \in V''$ 表示 $u''$ 的像。以 $V'', v''$
代替 $V', v'$ 后，可以假设复合 $V' \to Y' \to Y$ 分解经过 \'etale
邻域的映射 $(V', v') \to (V, v)$，并且所诱导的态射
$X'_{V'} = X_{V'} \to X_V$ 把 $u'$ 映到 $u$。在基变换
$X'_{V'} = X_{V'}$ 中有两个开子概形，即 $U'$ 以及
$U \subset X_V$ 的原像 $U_{V'}$。按构造，二者都恰有一个点位于
$v'$ 上，而且对二者该点都是 $u'$。因此缩小 $V'$ 后，可以假设这两个
开子概形相同；事实上，$U' \setminus (U' \cap U_{V'})$ 和
$U_{V'} \setminus (U' \cap U_{V'})$ 在 $V'$ 中有闭像，而这些像不含
$v'$。因此 $U' = U_{V'}$，从而得到
引理 \ref{more-morphisms-lemma-weighting-check-after-etale-base-change} 中的笛卡尔图。
由假设，$\int_\pi (w \circ h)$ 局部常值，故结论成立。
\end{proof}

\begin{lemma}
\label{more-morphisms-lemma-weighting-open}
设 $f : X \to Y$ 是局部拟有限态射，$w : X \to \mathbf{Z}$ 是 $f$
的赋权。若 $X' \subset X$ 是开集，则 $w|_{X'}$ 是
$f|_{X'} : X' \to Y$ 的赋权。
\end{lemma}

\begin{proof}
直接由定义得出。
\end{proof}

\begin{lemma}
\label{more-morphisms-lemma-weighting-composition}
设 $f : X \to Y$ 和 $g : Y \to Z$ 是局部拟有限态射。设
$w_f : X \to \mathbf{Z}$ 是 $f$ 的赋权，而
$w_g : Y \to \mathbf{Z}$ 是 $g$ 的赋权。则函数
$$
X \longrightarrow \mathbf{Z},\quad
x \longmapsto w_f(x) w_g(f(x))
$$
是 $g \circ f$ 的赋权。
\end{lemma}

\begin{proof}
对 $x \in X$，置 $w_{g \circ f}(x) = w_f(x) w_g(f(x))$。考虑图
$$
\xymatrix{
X \ar[d]_{g \circ f} & U \ar[l] \ar[d]^\pi \\
Z & W \ar[l]
}
$$
其中 $W \to Z$ 是 \'etale 的，$U \subset X_W$ 是开集，且 $U \to W$ 有限。
我们要证明 $\int_\pi w_{g \circ f}|_U$ 局部常值。选择一点 $w \in W$。
由 引理 \ref{more-morphisms-lemma-weighting-check-after-etale-base-change}
（以及 \'etale 态射是开态射这一事实），只需证明用一个 \'etale
邻域代替 $(W, w)$ 后，$\int_\pi w_{g \circ f}|_U$ 常值。
以一个 \'etale 邻域代替 $(W, w)$ 后，可以假设
$U = U_1 \amalg \ldots \amalg U_n$，其中每个 $U_i$ 都恰有一点 $u_i$
位于 $w$ 上，并且 $\kappa(u_i)/\kappa(w)$ 纯不可分，见
引理 \ref{more-morphisms-lemma-etale-splits-off-quasi-finite-part-technical-variant}。
显然，只需证明 $\int_{U_i \to W} w_{g \circ f}|_{U_i}$ 在 $w$
的某个 \'etale 邻域中常值。这把我们归约到下一段讨论的情形。

\medskip\noindent
我们有 $w \in W$，且该开集中恰有一点 $u \in U$ 位于 $w$ 上，
并满足 $\kappa(u)/\kappa(w)$ 纯不可分。考虑点 $v = f(u) \in Y$。
用一个初等 \'etale 邻域代替 $(W, w)$ 后，可以假设存在 $v$
的开邻域 $V \subset Y_W$，使得 $V \to W$ 有限，见
引理 \ref{more-morphisms-lemma-etale-makes-quasi-finite-finite-at-point}。于是
$f_W^{-1}(V) \cap U$ 是 $u$ 的开邻域，其中
$f_W : X_W \to Y_W$ 是 $f$ 到 $W$ 的基变换。因此在 Zariski 意义下
缩小 $W$ 后，可以假设 $f_W(U) \subset V$。于是得到态射
$$
U \xrightarrow{a} V \xrightarrow{b} W
$$
由于 $V \to W$ 分离（因为它有限），$U \to V$ 有限。由于 $w_f$ 和
$w_g$ 分别是 $f$ 和 $g$ 的赋权，$\int_a w_f|_U$ 在 $V$ 上局部常值，
而 $\int_b w_g|_V$ 在 $W$ 上局部常值。因此再次缩小 $W$ 后，
可以假设这两个函数常值，设其值分别为 $n$ 和 $m$。于是立即可知
$\int_\pi w_{g \circ f}|_U = \int_{b \circ a} w_{g \circ f}|_U$
常值为 $nm$，正合所需。
\end{proof}

\begin{lemma}
\label{more-morphisms-lemma-weighting-universally-open}
设 $f : X \to Y$ 是局部拟有限态射，$w : X \to \mathbf{Z}$ 是一个赋权。
若对所有 $x \in X$ 都有 $w(x) > 0$，则 $f$ 普遍开。
\end{lemma}

\begin{proof}
该性质在基变换下保持，见 引理 \ref{more-morphisms-lemma-weighting-base-change}，
故只需证明 $f$ 是开态射。又因为可以用 $X$ 的任一开集代替 $X$，
只需证明 $f(X)$ 为开集。令 $y \in f(X)$。选择 $x \in X$，使得
$f(x) = y$。只需证明 $f(X)$ 包含 $y$ 的一个开邻域，并且只需在用
$y$ 的一个 \'etale 邻域代替 $Y$ 后证明这一点。通过拟有限态射的
\'etale 局部化，见 第 \ref{more-morphisms-section-etale-localization} 节，可以假设存在
$x$ 的开邻域 $U \subset X$，使得 $\pi = f|_U : U \to Y$ 有限。
于是 $\int_\pi w|_U$ 局部常值，并且在 $y$ 处取正值。因此
$\pi(U)$ 包含 $y$ 的一个开邻域，证明完成。
\end{proof}

\begin{lemma}
\label{more-morphisms-lemma-weighting-flat-quasi-finite}
设 $f : X \to Y$ 是概形态射。假设 $f$ 局部拟有限、局部有限表示且平坦。
则 $f$ 有一个正赋权 $w : X \to \mathbf{Z}_{> 0}$，它把位于
$y \in Y$ 上的 $x \in X$ 映到
$$
w(x) =
\text{length}_{\mathcal{O}_{X, x}}
(\mathcal{O}_{X, x}/\mathfrak m_y \mathcal{O}_{X, x})
[\kappa(x) : \kappa(y)]_i
$$
其中 $[\kappa' : \kappa]_i$ 是不可分次数
（域，定义 \ref{fields-definition-insep-degree}）。
\end{lemma}

\begin{proof}
考虑 定义 \ref{more-morphisms-definition-weighting} 中的图。设 $u \in U$
在 $X, Y, V$ 中的像分别为 $x, y, v$。我们断言
$$
\text{length}_{\mathcal{O}_{X, x}}
(\mathcal{O}_{X, x}/\mathfrak m_y \mathcal{O}_{X, x}) =
\text{length}_{\mathcal{O}_{U, u}}
(\mathcal{O}_{U, u}/\mathfrak m_v \mathcal{O}_{U, u})
$$
及
$$
[\kappa(x) : \kappa(y)]_i =
[\kappa(u) : \kappa(v)]_i
$$
第一个等式成立，是因为
$\mathcal{O}_{X, x} \to \mathcal{O}_{U, u}$ 是平坦局部同态，满足
$\mathfrak m_y \mathcal{O}_{U, u} = \mathfrak m_v \mathcal{O}_{U, u}$ 和
$\mathfrak m_x \mathcal{O}_{U, u} = \mathfrak m_u$（因为
$\mathcal{O}_{Y, y} \to \mathcal{O}_{V, v}$ 和
$\mathcal{O}_{X, x} \to \mathcal{O}_{U, u}$ 都非分歧），再应用代数，
引理 \ref{algebra-lemma-pullback-module} 即得该等式。第二个等式成立，
是因为 $\kappa(v)/\kappa(y)$ 是有限可分扩张，而 $\kappa(u)$ 是
$\kappa(x) \otimes_{\kappa(y)} \kappa(v)$ 的一个因子，故不可分次数不变。
由此可见，引理中的函数之形成与 \'etale 基变换相容。这把问题归约到
下一段的讨论。

\medskip\noindent
假设 $f$ 是有限、平坦且有限表示的态射。我们要证明 $\int_f w$
在 $Y$ 上局部常值。事实上，$f$ 有限局部自由（态射，Lemma
\ref{morphisms-lemma-finite-flat}），而我们将证明 $\int_f w$ 等于 $f$
的次数（它是 $Y$ 上的局部常值函数）。具体地，对 $y \in Y$ 有
\begin{align*}
(\textstyle{\int}_f w)(y)
& =
\sum\nolimits_{f(x) = y}
\text{length}_{\mathcal{O}_{X, x}}
(\mathcal{O}_{X, x}/\mathfrak m_y \mathcal{O}_{X, x})
[\kappa(x) : \kappa(y)]_i
[\kappa(x) : \kappa(y)]_s \\
& =
\sum\nolimits_{f(x) = y}
\text{length}_{\mathcal{O}_{X, x}}
(\mathcal{O}_{X, x}/\mathfrak m_y \mathcal{O}_{X, x})
[\kappa(x) : \kappa(y)] \\
& =
\text{length}_{\mathcal{O}_{Y, y}}((f_*\mathcal{O}_X)_y/
\mathfrak m_y (f_*\mathcal{O}_X)_y)
\end{align*}
最后一个等号由代数，引理 \ref{algebra-lemma-pushdown-module} 得出。
最后所得数正是 $f_*\mathcal{O}_X$ 在 $y$ 处的秩，正合所需。
\end{proof}

\begin{lemma}
\label{more-morphisms-lemma-weighting-quasi-finite-Noetherian}
设 $f : X \to Y$ 是概形态射。假设
\begin{enumerate}
\item $f$ 局部拟有限；
\item $Y$ 几何单分支且局部诺特。
\end{enumerate}
则存在赋权 $w : X \to \mathbf{Z}_{\geq 0}$，它把位于 $y \in Y$
上的 $x \in X$ 映到 $\mathcal{O}_{X, x}^{sh}$ 在
$\mathcal{O}_{Y, y}^{sh}$ 上的“泛可分次数”。
\end{lemma}

\begin{proof}
由代数，引理 \ref{algebra-lemma-quasi-finite-strict-henselization}，
$\mathcal{O}_{Y, y}^{sh} \to \mathcal{O}_{X, x}^{sh}$ 有限。由于 $Y$
几何单分支，$\mathcal{O}_{Y, y}^{sh}$ 中有唯一的极小素理想
$\mathfrak p$，见代数进阶，Lemma
\ref{more-algebra-lemma-geometrically-unibranch}。写成
$$
(\kappa(\mathfrak p) \otimes_{\mathcal{O}_{Y, y}^{sh}}
\mathcal{O}_{X, x}^{sh})_{red} =
\prod K_i
$$
其中右侧是域的有限乘积。置
$w(x) = \sum [K_i : \kappa(\mathfrak p)]_s$。

\medskip\noindent
该定义显然不受 \'etale 局部化影响。因此为证明 $w$ 是赋权，
只需证明当 $f$ 是有限态射时，$\int_f w$ 局部常值。注意，
$\int_f w$ 在 $Y$ 的泛点 $\eta$ 处的值，就是 $X \to Y$ 在 $\eta$
上的几何纤维 $X_{\overline{\eta}}$ 的点数。此外，由于 $Y$ 单分支，
$Y$ 的一点 $y$ 是唯一一个泛点 $\eta$ 的特化。因此只需证明
$(\int_f w)(y)$ 等于 $X_{\overline{\eta}}$ 的点数。转到 $y$
的仿射邻域后，可以假设 $X \to Y$ 由有限环映射 $A \to B$ 给出。
假设 $\mathcal{O}_{Y, y}^{sh}$ 是利用到代数闭域 $k$ 的映射
$\kappa(y) \to k$ 构造的。则
$$
\mathcal{O}_{Y, y}^{sh} \otimes_A B =
\prod\nolimits_{f(x) = y}
\prod\nolimits_{\varphi \in \Mor_{\kappa(y)}(\kappa(x), k)}
\mathcal{O}_{X, x}^{sh}
$$
这是由代数，引理 \ref{algebra-lemma-finite-over-henselian} 和上面使用的
引理得到的。注意，$\mathcal{O}_{Y, y}^{sh}$ 的极小素理想 $\mathfrak p$
映到 $A$ 中对应于 $\eta$ 的素理想。因此所需等式成立，因为几何纤维的
点数在域扩张下不变。
\end{proof}

\begin{lemma}
\label{more-morphisms-lemma-weighting-specialization}
设 $f : X \to Y$ 是局部拟有限概形态射，$w : X \to \mathbf{Z}$ 是 $f$
的赋权。设 $y \leadsto y'$ 是 $Y$ 中点的特化。假设对所有满足
$x \in X$ 且 $f(x) = y$ 的点都有 $w(x) = 0$。则对所有满足
$x' \in X$ 且 $f(x') = y'$ 的点都有 $w(x') = 0$。
\end{lemma}

\begin{proof}
令 $x' \in X$ 为映到 $y'$ 的点。由
引理 \ref{more-morphisms-lemma-etale-makes-quasi-finite-finite-at-point}，可以选择图
$$
\xymatrix{
U \ar[r] \ar[dr]_\pi & X_V \ar[d] \ar[r] & X \ar[d]^f \\
& V \ar[r]^\varphi & Y
}
$$
其中 $V \to Y$ 是 \'etale 的，$v' \in V$ 映到 $y'$，
% P05-EN-0162: frozen source says U is open in Y_V, but the diagram and quasi-finite localization require U open in X_V.
$U \subset Y_V$ 是开集，$U \to V$ 有限，并且恰有一点 $u' \in U$
映到 $v'$，且该点在 $X$ 中映到 $x'$。注意，由于 $V \to Y$ 是开态射，
存在点 $v \in V$ 特化到 $v'$ 并映到 $y$。由 Definition
\ref{more-morphisms-definition-weighting}，函数 $\int_\pi (w \circ h)$ 局部常值，
其中 $h : U \to X$ 是上方两个水平箭头的复合。由对 $w$ 的假设，
$\int_\pi (w \circ h)$ 在 $v$ 处为零，故在 $v'$ 处也为零。然而，
它在 $v'$ 处的值是 $w(x')[\kappa(u'):\kappa(v')]_s$，结论随即得出。
\end{proof}




\section{赋权的进一步性质}
\label{more-morphisms-section-more-weightings}

\noindent
我们再证明赋权的若干基本性质。虽然乍看之下赋权可能表现得十分不规则，
但事实上定义 \ref{more-morphisms-definition-weighting} 中所施加的条件相当强。

\begin{lemma}
\label{more-morphisms-lemma-jumps-w}
设 $f : X \to Y$ 为局部拟有限态射，
$w : X \to \mathbf{Z}$ 为 $f$ 的一个赋权。则函数 $w$ 的各水平集在 $X$ 中局部可构造。
\end{lemma}

\begin{proof}
在以下证明中，我们将不再另行说明地使用引理 \ref{more-morphisms-lemma-weighting-open} 和
\ref{more-morphisms-lemma-weighting-base-change}。我们还将使用概形及拓扑空间的可构造子集的
初等性质，见拓扑章第 \ref{topology-section-constructible} 节和
性质章第 \ref{properties-section-constructible} 节。由此读者可见，该问题关于
$X$ 和 $Y$ 都是局部的；细节从略。因此可设 $X$ 和 $Y$ 都是仿射的。
如果能够找到一个有限呈现满态射 $Y' \to Y$，使得 $w$ 的各水平集在
$X' = Y' \times_Y X$ 中的拉回均为局部可构造子集，那么结论即由
态射章定理 \ref{morphisms-theorem-chevalley} 得出。

\medskip\noindent
设 $X$ 和 $Y$ 为仿射概形。由引理
\ref{more-morphisms-lemma-quasi-finite-separated-pass-through-finite}，可选取一个浸入
$X \to T$，其中 $T \to Y$ 为有限态射。根据态射章引理
\ref{morphisms-lemma-massage-finite}，将 $Y$ 替换为一个在 $Y$ 上有限局部自由的
满态射 $Y'$，将 $X$ 替换为 $Y' \times_Y X$，并将 $T$ 替换为一个在 $Y'$ 上
有限局部自由且以 $Y' \times_Y T$ 为闭子概形的概形之后，可设 $T$ 在 $Y$ 上
有限局部自由，并且包含与 $Y$ 同构地映射的闭子概形 $T_i$，使得
$T = \bigcup_{i = 1, \ldots, n} T_i$（集合论意义下）。由于 $T_i \subset T$
是可构造闭子集（它是概形的一个有限呈现态射 $Y \to T$ 的像），故对于
$I \subset \{1, \ldots, n\}$，交集 $\bigcap_{i \in I} T_i$ 是 $T$ 的可构造闭子集，
因而映为 $Y$ 的可构造闭子集。

\medskip\noindent
对于一个各部分均非空的不交并分解
$\{1, \ldots, n\} = I_1 \amalg \ldots \amalg I_r$，考虑子集
$Y_{I_1, \ldots, I_r} \subset Y$，其元素为满足如下条件的点 $y \in Y$：
$T_y = \{x_1, \ldots, x_r\}$ 恰由 $r$ 个点组成，且
$x_j \in T_i \Leftrightarrow i \in I_j$。由以上说明，这些子集构成 $Y$ 的一个
可构造分拆。由代数章引理 \ref{algebra-lemma-constructible-is-image}，存在一个
在 $Y$ 上有限呈现的仿射概形 $Y'$，使得 $Y' \to Y$ 的像恰为
$Y_{I_1, \ldots, I_r}$。因此可设对于某个不交并分解
$\{1, \ldots, n\} = I_1 \amalg \ldots \amalg I_r$，有
$Y = Y_{I_1, \ldots, I_r}$。此时，令 $T(j) = \bigcap_{i \in I_j} T_i$，则
$T = T(1) \amalg \ldots \amalg T(r)$ 是 $T$ 分解为两两不交闭子集（因而也是
开子集）的分解。与局部闭子概形 $X$ 相交，得到类似的开闭分解
$X = X(1) \amalg \ldots \amalg X(r)$。态射 $X(j) \to Y$ 为浸入。
由于 $w$ 是赋权，$w|_{X(j)}$ 局部常值\footnote{事实上，若
$f : X \to Y$ 为浸入且 $w$ 是 $f$ 的一个赋权，则在 $w$ 非零的轨迹上，
$f$ 的限制为开映射。}，结论随即成立。
\end{proof}

\begin{lemma}
\label{more-morphisms-lemma-weighting-on-dense-set}
设 $f : X \to Y$ 为概形的局部拟有限态射，
$w : X \to \mathbf{Z}$ 为 $f$ 的一个赋权，$E \subset X$ 为稠密子集。
若对所有 $x \in E$ 均有 $w(x) = 0$，则 $w = 0$。
\end{lemma}

\begin{proof}
建议先阅读引理 \ref{more-morphisms-lemma-weighting-specialization}。
取 $x \in X$，其像为 $y \in Y$。我们将证明 $w(x) = 0$。由引理
\ref{more-morphisms-lemma-etale-makes-quasi-finite-finite-at-point}，可选取图表
$$
\xymatrix{
U \ar[r] \ar[dr]_\pi & X_V \ar[d] \ar[r] & X \ar[d]^f \\
& V \ar[r]^\varphi & Y
}
$$
其中 $V \to Y$ 为 \'etale 态射，点 $v \in V$ 映到 $y$，
开子集 $U \subset Y_V$ 满足 $U \to V$ 为有限态射，并且存在唯一一点
$u \in U$ 映到 $v$，且它在 $X$ 中进一步映到 $x$。我们可以并且确实设
$V$（从而 $U$）为仿射的。由定义 \ref{more-morphisms-definition-weighting}，函数
$\int_\pi (w \circ h)$ 局部常值，其中 $h : U \to X$ 是上方横向箭头的复合。
由引理 \ref{more-morphisms-lemma-jumps-w}，$w \circ h$ 为 $0$ 的点集 $U_0$ 是 $U$ 的
可构造子集；又因它包含 $h^{-1}(E)$（使用 $h$ 为开映射），故在 $U$ 中稠密。
由拓扑章引理 \ref{topology-lemma-dense-in-constructible}，$U_0$ 包含一个稠密
开子集；换言之，$w \circ h$ 非零的点集包含在某个无处稠密闭子集
$Z \subset U$ 中。于是由态射章引理
\ref{morphisms-lemma-image-nowhere-dense-finite}，$\pi(Z)$ 是 $V$ 的无处稠密
闭子集。因此 $\int_\pi (w \circ h)$ 在一个稠密开子集上为零，从而处处为零。
然而它在 $v$ 处的值为 $w(x)[\kappa(u):\kappa(v)]_s$，结论得证。
\end{proof}

\begin{lemma}
\label{more-morphisms-lemma-jumps-int-w}
设 $f : X \to Y$ 为有限呈现的局部拟有限态射，
$w : X \to \mathbf{Z}$ 为 $f$ 的一个赋权。则函数 $\int_f w$ 的各水平集在
$Y$ 中局部可构造。
\end{lemma}

\begin{proof}
由引理 \ref{more-morphisms-lemma-weighting-check-after-etale-base-change}，函数 $\int_f w$ 的构造
与任意基变换交换；由引理 \ref{more-morphisms-lemma-weighting-base-change}，基变换后仍得到赋权。
这意味着，若能找到一个有限呈现满态射 $Y' \to Y$，则只须在基变换到 $Y'$
之后证明结论，见态射章定理 \ref{morphisms-theorem-chevalley}。

\medskip\noindent
该问题在 $Y$ 上是局部的，故可设 $Y$ 为仿射概形。于是 $X$ 拟紧且拟分离
（因为 $f$ 为有限呈现）。设 $X = U \cup V$，其中两个开子集均拟紧。
则有
$$
\textstyle{\int}_f w =
\textstyle{\int}_{f|_U} w|_U +
\textstyle{\int}_{f|_V} w|_V -
\textstyle{\int}_{f|_{U \cap V}} w|_{U \cap V}
$$
因此，若结论对 $w|_U$、$w|_V$ 和 $w|_{U \cap V}$ 成立，则它也对 $w$ 成立。
由归纳原理（概形上同调章引理 \ref{coherent-lemma-induction-principle}），
只须在 $X$ 为仿射概形时证明本引理。

\medskip\noindent
设 $X$ 和 $Y$ 都是仿射概形。由引理
\ref{more-morphisms-lemma-quasi-finite-separated-pass-through-finite}，可选取一个开浸入
$X \to T$，其中 $T \to Y$ 为有限态射。因为仍可沿适当的 $Y' \to Y$ 作基变换，
所以可使用态射章引理 \ref{morphisms-lemma-massage-finite}，将问题约化到如下情形：
态射 $T \to Y$ 所诱导的所有剩余域扩张（因而当然也包括 $X \to Y$ 所诱导的）
均为平凡扩张。在此情形下，$\int_f w$ 就是对每条纤维中的 $w$ 值求和。
由引理 \ref{more-morphisms-lemma-jumps-w}，$w$ 的各水平集在 $X$ 中局部可构造；由性质章引理
\ref{properties-lemma-stratification-locally-finite-constructible}，函数 $w$ 仅取
有限多个值。因此结论由态射章定理 \ref{morphisms-theorem-chevalley} 以及关于整数
之和的若干初等论证得出。
\end{proof}

\begin{lemma}
\label{more-morphisms-lemma-semicontinuous-w}
设 $f : X \to Y$ 为局部拟有限态射，
$w : X \to \mathbf{Z}_{> 0}$ 为 $f$ 的一个正赋权。则 $w$ 上半连续。
\end{lemma}

\begin{proof}
取 $x \in X$，其像为 $y \in Y$。选取一个 \'etale 邻域
$(V, v) \to (Y, y)$ 以及开子集 $U \subset X_V$，使得
$\pi : U \to V$ 为有限态射，并且存在唯一一点 $u \in U$ 映到 $v$，且
$\kappa(u)/\kappa(v)$ 为纯不可分扩张。见引理
\ref{more-morphisms-lemma-etale-makes-quasi-finite-finite-multiple-points-var}。于是
$(\int_\pi w|_U)(v) = w(u)$。由定义 \ref{more-morphisms-definition-weighting}，将 $V$ 替换为
$v$ 的一个邻域后，对所有 $u' \in U$ 均有
$w|_U(u') \leq w|_U(u) = w(x)$。事实上，$w|_U(u')$ 是表达式
$(\int_\pi w|_U)(\pi(u'))$ 中的一个加项。由于 \'etale 态射 $U \to X$ 为开映射，
这就证明了本引理。
\end{proof}

\begin{lemma}
\label{more-morphisms-lemma-semicontinuous-int-w}
设 $f : X \to Y$ 为纤维有限的分离局部拟有限态射，
$w : X \to \mathbf{Z}_{> 0}$ 为 $f$ 的一个正赋权。则 $\int_f w$ 下半连续。
\end{lemma}

\begin{proof}
取 $y \in Y$，设 $x_1, \ldots, x_r \in X$ 为位于 $y$ 上方的点。应用引理
\ref{more-morphisms-lemma-etale-splits-off-quasi-finite-part-technical-variant}，得到一个 \'etale 邻域
$(U, u) \to (Y, y)$ 以及分解
$$
U \times_Y X =
W \amalg
\ \coprod\nolimits_{i = 1, \ldots, n}
\ \coprod\nolimits_{j = 1, \ldots, m_i}
V_{i, j}
$$
如所引之处。注意 $(\int_f w)(y) = \sum w(v_{i, j})$，其中
$w(v_{i, j}) = w(x_i)$。依定义，$\int_{V_{i, j} \to U} w|_{V_{i, j}}$ 局部常值，
故缩小 $U$ 后可设这些函数均为常值函数，其值为 $w(v_{i, j})$。于是
$$
\textstyle{\int}_{U \times_Y X \to U} w|_{U \times_Y X} =
\textstyle{\int}_{W \to U} w|_W +
\sum \textstyle{\int}_{V_{i, j} \to U} w|_{V_{i, j}} =
\textstyle{\int}_{W \to U} w|_W + (\int_f w)(y)
$$
该值 $\geq (\int_f w)(y)$。又因 $U \to Y$ 为开映射，且积分的构造与基变换交换
（引理 \ref{more-morphisms-lemma-weighting-check-after-etale-base-change}），结论得证。
\end{proof}

\begin{lemma}
\label{more-morphisms-lemma-max-int-w}
设 $f : X \to Y$ 为局部拟有限态射，其中 $X$ 拟紧；设
$w : X \to \mathbf{Z}$ 为 $f$ 的一个赋权。则 $\int_f w$ 取得其最大值。
\end{lemma}

\begin{proof}
由引理 \ref{more-morphisms-lemma-jumps-w} 和性质章引理
\ref{properties-lemma-stratification-locally-finite-constructible}，$w$ 在 $X$ 上仅取
有限多个值。由态射章引理
\ref{morphisms-lemma-locally-quasi-finite-qc-source-universally-bounded}，
$X \to Y$ 的几何纤维一致有界。这表明 $\int_f w$ 有界。
\end{proof}

\begin{lemma}
\label{more-morphisms-lemma-max-int-finite}
设 $f : X \to Y$ 为分离局部拟有限态射，
$w : X \to \mathbf{Z}_{> 0}$ 为 $f$ 的一个正赋权。假设 $\int_w f$ 取得最大值
$d$，并令 $Y_d \subset Y$ 为满足 $(\int_f w)(y) = d$ 的点 $y$ 所组成的开集。
则态射 $f^{-1}(Y_d) \to Y_d$ 为有限态射。
\end{lemma}

\begin{proof}
由引理 \ref{more-morphisms-lemma-semicontinuous-int-w}，$Y_d$ 为开集。取 $y \in Y_d$，设
$x_1, \ldots, x_n$ 为 $X$ 中位于 $y$ 上方的点。应用引理
\ref{more-morphisms-lemma-etale-splits-off-quasi-finite-part-technical-variant}，得到一个 \'etale 邻域
$(U, u) \to (Y, y)$ 以及分解
$$
U \times_Y X =
W \amalg
\ \coprod\nolimits_{i = 1, \ldots, n}
\ \coprod\nolimits_{j = 1, \ldots, m_i}
V_{i, j}
$$
如所引之处。注意 $d = \sum w(v_{i, j})$，其中 $w(v_{i, j}) = w(x_i)$。
依定义，$\int_{V_{i, j} \to U} w|_{V_{i, j}}$ 局部常值，故缩小 $U$ 后可设这些
函数均为常值函数，其值为 $w(v_{i, j})$。于是
$$
\textstyle{\int}_{U \times_Y X \to U} w|_{U \times_Y X} =
\textstyle{\int}_{W \to U} w|_W +
\sum \textstyle{\int}_{V_{i, j} \to U} w|_{V_{i, j}} =
\textstyle{\int}_{W \to U} w|_W + (\int_f w)(y)
$$
该值 $\geq (\int_f w)(y) = d$，故 $W$ 必为空集。因此
$U \times_Y X \to U$ 为有限态射。由下降章引理
\ref{descent-lemma-descending-property-finite}，这说明 $X \to Y$ 在开态射
$U \to Y$ 的像上为有限态射。换言之，$f$ 在 $y$ 的某个开邻域上为有限态射，
正如所需。
\end{proof}

\begin{lemma}
\label{more-morphisms-lemma-open-and-closed-in-finite}
设 $A \to B$ 为有限且有限呈现的环同态。则存在一个有限呈现环同态
$A \to A_{univ}$ 以及一个幂等元 $e_{univ} \in B \otimes_A A_{univ}$，
使得对于任意环同态 $A \to A'$ 和幂等元 $e \in B \otimes_A A'$，
都存在一个将 $e_{univ}$ 映到 $e$ 的环同态 $A_{univ} \to A'$。
\end{lemma}

\begin{proof}
选取 $b_1, \ldots, b_n \in B$，使其作为 $A$-模生成 $B$。对每个 $i$，选取首一
多项式 $P_i \in A[x]$，使得 $P_i(b_i) = 0$ 在 $B$ 中成立，见代数章引理
\ref{algebra-lemma-finite-is-integral}。于是 $B$ 是有限自由 $A$-代数
$B' = A[x_1, \ldots, x_n]/(P_1(x_1), \ldots, P_n(x_n))$ 的商。令
$J \subset B'$ 为满射 $B' \to B$ 的核。由于 $B$ 是有限生成 $A$-代数，
$J =(f_1, \ldots, f_m)$ 有限生成，见代数章引理
\ref{algebra-lemma-compose-finite-type}。选取 $B'$ 的一个 $A$-基
$b'_1, \ldots, b'_N$。考虑代数
$$
A_{univ} = A[z_1, \ldots, z_N, y_1, \ldots, y_m]/I
$$
其中 $I$ 是由下列表达式关于基元素 $b'_1, \ldots, b'_N$ 的系数在
$A[z_1, \ldots, z_n, y_1, \ldots, y_m]$ 中所生成的理想：
$$
(\sum z_j b'_j)^2 - \sum z_j b'_j + \sum y_k f_k
$$
该表达式位于 $B'[z_1, \ldots, z_N, y_1, \ldots, y_m]$ 中。依构造，元素
$\sum z_j b'_j$ 在代数 $B \otimes_A A_{univ}$ 中映为幂等元 $e_{univ}$。
此外，若 $e \in B \otimes_A A'$ 为幂等元，则可将 $e$ 提升为
$B' \otimes_A A'$ 中形如 $\sum b'_j \otimes a'_j$ 的元素，并且可以找到
$a''_k \in A'$，使得
$$
(\sum b'_j \otimes a'_j)^2 - \sum b'_j \otimes a'_j + \sum f_k \otimes a''_k
$$
在 $B' \otimes_A A'$ 中为零。因此得到一个 $A$-代数同态
$A_{univ} \to A$，它将 $z_j$ 映到 $a'_j$、将 $y_k$ 映到 $a''_k$，并将
$e_{univ}$ 映到 $e$。证明完毕。
\end{proof}

\begin{lemma}
\label{more-morphisms-lemma-open-and-closed-in-quasi-finite}
设 $X \to Y$ 为仿射概形间拟有限且有限呈现的态射。则存在一个有限呈现态射
$Y_{univ} \to Y$ 以及一个开子概形 $U_{univ} \subset Y_{univ} \times_Y X$，
使得 $U_{univ} \to Y_{univ}$ 为有限态射，并具有如下性质：给定任意仿射概形间的
态射 $Y' \to Y$ 以及开子概形 $U' \subset Y' \times_Y X$，若
$U' \to Y'$ 为有限态射，则存在态射 $Y' \to Y_{univ}$，使得 $U_{univ}$ 的逆像
为 $U'$。
\end{lemma}

\begin{proof}
回忆有限型态射为拟有限态射，当且仅当其相对维数为 $0$，见态射章引理
\ref{morphisms-lemma-locally-quasi-finite-rel-dimension-0}。在 $d = 0$ 时应用引理
\ref{more-morphisms-lemma-Noetherian-approximation-dimension-d}，可约化到 $X$ 和 $Y$ 均为
Noether 概形的情形。由代数章引理
\ref{algebra-lemma-quasi-finite-open-integral-closure}，可选取一个开浸入
$X \to X'$，使得 $X' \to Y$ 为有限态射。注意，若有如 (2) 中的 $Y' \to Y$
和 $U'$，则
$$
U' \to Y' \times_Y X \to Y' \times_Y X'
$$
是两个在 $Y'$ 上有限的概形之间的开浸入，因而也是闭浸入。由此可知，$U'$ 对应于
下列代数中的一个幂等元：
$$
\Gamma(Y', \mathcal{O}_{Y'})
\otimes_{\Gamma(Y, \mathcal{O}_Y)}
\Gamma(X', \mathcal{O}_{X'})
$$
该幂等元所对应的开闭子集包含在开集 $Y' \times_Y X$ 中。令
$Y'_{univ} \to Y$ 及幂等元
$$
e'_{univ} \in
\Gamma(Y_{univ}, \mathcal{O}_{Y_{univ}})
\otimes_{\Gamma(Y, \mathcal{O}_Y)}
\Gamma(X', \mathcal{O}_{X'})
$$
为引理 \ref{more-morphisms-lemma-open-and-closed-in-finite} 对环同态
$\Gamma(Y, \mathcal{O}_Y) \to \Gamma(X', \mathcal{O}_{X'})$ 所构造的二元组
（这里使用 $Y$ 为 Noether 概形，以知 $X'$ 在 $Y$ 上有限呈现）。令
$U'_{univ} \subset Y'_{univ} \times_Y X'$ 为相应的开闭子概形。于是
$$
U'_{univ} \setminus Y'_{univ} \times_Y X
$$
是 $U'_{univ}$ 的闭子集，故其像 $T \subset Y'_{univ}$ 为闭集。若令
$Y_{univ} = Y'_{univ} \setminus T$，并令 $U_{univ}$ 为 $U'_{univ}$ 在
$Y_{univ} \times_Y X$ 上的限制，则本引理成立。
\end{proof}

\begin{lemma}
\label{more-morphisms-lemma-descend-weighting}
设 $Y = \lim Y_i$ 为仿射概形的有向极限。取 $0 \in I$，并设
$f_0 : X_0 \to Y_0$ 为仿射概形间拟有限且有限呈现的态射。令
$f : X  \to Y$ 以及对 $i \geq 0$ 的 $f_i : X_i \to Y_i$ 为 $f_0$ 的基变换。
若 $w : X \to \mathbf{Z}$ 为 $f$ 的一个赋权，则当 $i$ 充分大时，存在 $f_i$
的一个赋权 $w_i : X_i \to \mathbf{Z}$，其到 $X$ 的拉回为 $w$。
\end{lemma}

\begin{proof}
由引理 \ref{more-morphisms-lemma-jumps-w}，$w$ 的各水平集是 $X$ 的可构造子集 $E_k$。
由性质章引理 \ref{properties-lemma-stratification-locally-finite-constructible}，
这说明函数 $w$ 仅取有限多个值。因此存在 $i$，使得对所有 $k$，$E_k$ 均下降为
$X_i$ 中的一个可构造子集 $E_{i, k}$；此外，可设
$X_i = \coprod E_{i, k}$。这是因为 $X$ 的拓扑空间是谱空间 $X_i$ 沿具有谱过渡
映射的有向系在拓扑空间范畴中的极限。见极限章第 \ref{limits-section-descent} 节和
拓扑章第 \ref{topology-section-limits-spectral} 节。定义
$w_i : X_i \to \mathbf{Z}$，使其各水平集为可构造集 $E_{i, k}$。

\medskip\noindent
按引理 \ref{more-morphisms-lemma-open-and-closed-in-quasi-finite} 选取 $Y_{i, univ} \to Y_i$ 和
$U_{i, univ} \subset Y_{i, univ} \times_{Y_i} X_i$。由该构造的泛性质，为了证明
$w_i$ 是赋权，只须证明
$$
\tau_i = \textstyle{\int}_{U_{i, univ} \to Y_{i, univ}} w_i|_{U_{i, univ}}
$$
在 $Y_{i, univ}$ 上局部常值。由引理 \ref{more-morphisms-lemma-jumps-int-w}，该函数具有可构造
水平集，但它（尚）不一定局部常值。令 $Y_{univ} = Y_{i, univ} \times_{Y_i} Y$，
并令 $U_{univ} \subset Y_{univ} \times_Y X$ 为 $U_{i, univ}$ 的逆像。由于 $w$
到 $Y_{univ} \times_Y X$ 的拉回是 $Y_{univ} \times_Y X \to Y_{univ}$ 的一个赋权
（引理 \ref{more-morphisms-lemma-weighting-base-change}），故
$$
\tau = \textstyle{\int}_{U_{univ} \to Y_{univ}} w_i|_{U_{univ}}
$$
在 $Y_{univ}$ 上局部常值。因此 $\tau$ 的各水平集均开且闭。最后，
$Y_{univ} = \lim_{i' \geq i} Y_{i', univ}$，并且 $\tau$ 的各水平集是相应定义的
$\tau_{i'}$ 的各水平集的逆极限。因此，上述引文蕴含：当 $i'$ 充分大时，
$\tau_{i'}$ 的各水平集也为开集。对这样的指标 $i'$，可知 $w_{i'}$ 是 $f_{i'}$
的一个赋权，正如所需。
\end{proof}







\section{赋权与仿射分层数}
\label{more-morphisms-section-bounds-asn}

\noindent
本节对一类由仿射概形覆盖的概形之仿射分层数给出一个界。

\begin{lemma}
\label{more-morphisms-lemma-affineness-of-large-open}
设 $f : X \to Y$ 为仿射概形间拟有限且有限呈现的态射，
$w : X \to \mathbf{Z}_{> 0}$ 为 $f$ 的一个正赋权。设 $d < \infty$ 为
$\int_f w$ 的最大值。则 $Y$ 的开子集
$$
Y_d = \{y \in Y \mid (\textstyle{\int}_f w)(y) = d \}
$$
为仿射概形。
\end{lemma}

\begin{proof}
由引理 \ref{more-morphisms-lemma-max-int-w}，$\int_f w$ 取得其最大值；由引理
\ref{more-morphisms-lemma-semicontinuous-int-w}，集合 $Y_d$ 为开集。因此本引理的陈述确有意义。

\medskip\noindent
约化到 Noether 情形；请跳过本段。回忆有限型态射为拟有限态射，当且仅当其相对维数
为 $0$，见态射章引理
\ref{morphisms-lemma-locally-quasi-finite-rel-dimension-0}。在 $d = 0$ 时应用引理
\ref{more-morphisms-lemma-Noetherian-approximation-dimension-d}，可找到仿射 Noether 概形间的
拟有限态射 $f_0 : X_0 \to Y_0$ 以及态射 $Y \to Y_0$，使得 $f$ 是 $f_0$ 的
基变换。于是可将 $Y = \lim Y_i$ 写成在 $Y_0$ 上有限型的仿射概形的有向极限，
见代数章引理 \ref{algebra-lemma-ring-colimit-fp}。由引理
\ref{more-morphisms-lemma-descend-weighting}，可找到 $i$，使我们的赋权 $w$ 下降为 $f_0$ 的基变换
$f_i : X_i \to Y_i$ 的一个赋权 $w_i$。若本引理对 $f_i, w_i$ 成立，则它也对
$f$ 成立，因为 $\int_f w$ 的构造与基变换交换，见引理
\ref{more-morphisms-lemma-weighting-check-after-etale-base-change}。

\medskip\noindent
设 $X$ 和 $Y$ 为 Noether 概形。令 $X' \to Y'$ 为 $f$ 沿态射
$g : Y' \to Y$ 的基变换。$\int_f w$ 的构造以及由此得到的开集 $Y_d$ 都与基变换
交换。若 $g$ 有限且满，则 $Y'_d \to Y_d$ 有限且满。在此情形下，证明 $Y_d$
为仿射概形等价于证明 $Y'_d$ 为仿射概形，见概形上同调章引理
\ref{coherent-lemma-image-affine-finite-morphism-affine-Noetherian}。

\medskip\noindent
由引理 \ref{more-morphisms-lemma-quasi-finite-separated-pass-through-finite}，可选取一个浸入
$X \to T$，其中 $T$ 在 $Y$ 上有限。我们对有限态射 $T \to Y$ 应用态射章引理
\ref{morphisms-lemma-massage-finite}。该引理说明，存在有限满态射 $Y' \to Y$，
使得 $Y' \times_Y T$ 是某个具有特殊形式且在 $Y'$ 上有限的概形 $T'$ 的闭子概形。
由第一段的讨论，可将 $Y$ 替换为 $Y'$，将 $T$ 替换为 $T'$，并将 $X$ 替换为
$Y' \times_Y X$。因此可设存在一个浸入 $X \to T$（未必开或闭）以及闭子概形
$T_i \subset T$，$i = 1, \ldots, n$，其中
\begin{enumerate}
\item $T \to Y$ 有限（且局部自由）；
\item $T_i \to Y$ 为同构；并且
\item $T = \bigcup_{i = 1, \ldots, n} T_i$ 在集合论意义下成立。
\end{enumerate}
令 $Y' = \coprod Y_k$ 为 $Y$ 的各不可约分支（视为 $Y$ 的整闭子概形）的不交并。
此处利用 $Y$ 为 Noether 概形，我们可以再次沿 $Y' \to Y$ 作基变换。因此还可设
$Y$ 为整 Noether 概形。

\medskip\noindent
去除重复项后，我们还可并且确实设当 $i \not = j$ 时 $T_i \not = T_j$。由于
$Y$ 以及所有 $T_i$ 都是整概形，这意味着若 $T_i$ 与 $T_j$ 相交，则它们的交是
映到 $Y$ 的真闭子集中的闭子集。

\medskip\noindent
注意 $V_i = X \cap T_i$ 是局部闭子集，而且是 $X$ 的闭子概形，故为仿射概形。
令 $\eta \in Y$ 和 $\eta_i \in T_i$ 为泛点。若 $\eta \not \in Y_d$，则
$Y_d = \emptyset$，结论成立。设 $\eta \in Y_d$。以
$I \in \{1, \ldots, n\}$ 表示满足 $\eta_i \in V_i$ 的指标 $i$ 所成的子集。
对 $i \in I$，局部闭子集 $V_i \subset T_i$ 包含不可约空间 $T_i$ 的泛点，
因而为开集。另一方面，由于 $f$ 为开映射（引理
\ref{more-morphisms-lemma-weighting-universally-open}），对任意 $x \in X$，均可找到 $i \in I$
以及特化 $\eta_i \leadsto x$。于是 $x \in T_i$，从而 $x \in V_i$。换言之，
在集合论意义下 $X = \bigcup_{i \in I} V_i$。我们断言
$Y_d = \bigcap_{i \in I} \Im(V_i \to Y)$；这将完成证明，因为 $Y$ 的各仿射开集
$\Im(V_i \to Y)$ 的交仍为仿射概形。

\medskip\noindent
对 $y \in Y$，设 $X$ 中的 $f^{-1}(\{y\}) = \{x_1, \ldots, x_r\}$。对每个
$i \in I$，至多存在一个 $j(i) \in \{1, \ldots, x_r\}$，使得
$\eta_i \leadsto x_{j(i)}$。事实上，$j(i)$ 存在且等于 $j$，当且仅当
$x_j \in V_i$。若 $i \in I$ 使得 $j = j(i)$ 存在，则 $V_i \to Y$ 在
$x_j \mapsto y$ 的一个邻域上为同构。因此，将源和靶分别替换为
$x_j \mapsto y$ 的邻域后，$\bigcup_{i \in I,\ j(i) = j} V_i \to Y$ 为有限态射。
于是赋权的定义给出 $w(x_j) = \sum_{i \in I,\ j(i) = j} w(\eta_i)$。故
$$
(\textstyle{\int}_f w)(\eta) =
\sum\nolimits_{i \in I} w(\eta_i) \geq
\sum\nolimits_{j(i)\text{ exists}} w(\eta_i) =
\sum\nolimits_j w(x_j) = (\textstyle{\int}_f w)(y)
$$
等号成立，当且仅当 $y$ 属于 $\bigcap_{i \in I} \Im(V_i \to Y)$，这正是所要证明的。
\end{proof}

\begin{proposition}
\label{more-morphisms-proposition-asn-weighting}
设 $f : X \to Y$ 为概形的拟有限满态射，
$w : X \to \mathbf{Z}_{> 0}$ 为 $f$ 的一个正赋权。假设 $X$ 为仿射概形，而
$Y$ 分离\footnote{$Y$ 的对角态射为仿射态射即已足够。}且非空。则 $Y$ 的仿射分层数
至多为 $\int_f w$ 的不同取值个数减 $1$。
\end{proposition}

\begin{proof}
注意，由于 $Y$ 分离，态射 $X \to Y$ 为仿射态射（态射章引理
\ref{morphisms-lemma-affine-permanence}）。由引理 \ref{more-morphisms-lemma-max-int-w}，函数
$\int_f w$ 取得其最大值 $d$。我们对 $d$ 作归纳。考虑 $Y$ 的开子概形
$Y_d = \{y \in Y \mid (\int_f w)(y) = d\}$；回忆
$f^{-1}(Y_d) \to Y_d$ 为有限态射，见引理 \ref{more-morphisms-lemma-max-int-finite}。
由引理 \ref{more-morphisms-lemma-affineness-of-large-open}，对每个仿射开集 $W \subset Y$，
$Y_d \cap W$ 均为仿射概形（这里使用 $W \times_Y X$ 在 $X$ 上仿射，因而为仿射
概形）。故 $Y_d \to Y$ 为仿射态射。于是
$f^{-1}(Y_d) = Y_d \times_Y X$ 在 $X$ 上仿射，从而是仿射概形。又
$f^{-1}(Y_d) \to Y_d$ 为满态射，故由极限章引理 \ref{limits-lemma-affine}，
$Y_d$ 为仿射概形。若 $Y_d = Y$，则 $Y$ 的仿射分层数为 $0$，结论成立。
否则，令 $Y' = Y \setminus Y_d$，视为 $Y$ 的既约闭子概形，并令
$X' = Y' \times_Y X$。由于 $X'$ 在 $X$ 中闭，故为仿射概形；由于 $Y'$ 在 $Y$
中闭，故为分离概形。态射 $f' : X' \to Y'$ 为满态射，而 $w$ 诱导 $f'$ 的一个
赋权 $w'$（见引理 \ref{more-morphisms-lemma-weighting-base-change}），且 $\int_{f'} w'$ 是
$\int_f w$ 在 $Y'$ 上的限制。由归纳假设，$Y'$ 具有长度 $\leq$
$\int_{f'} w'$ 的不同取值个数减 $1$ 的仿射分层，证明完成。
\end{proof}









\section{完全分解态射}
\label{more-morphisms-section-completely-decomposed}

\noindent
Nishnevich 为了定义如今所谓的 Nishnevich 拓扑，研究了 \'etale 态射的完全分解族
这一概念；例如见 \cite{Nishnevich}。

\begin{definition}
\label{more-morphisms-definition-cd-morphism}
若概形态射 $f : X \to Y$ 满足：对每个点 $y \in Y$，都存在点 $x \in X$，
使得 $f(x) = y$ 且域扩张 $\kappa(x)/\kappa(y)$ 为平凡扩张，则称它为
{\it 完全分解的}\footnote{这一术语可能并非标准术语。}。若具有固定靶的概形态射族
$\{f_i : X_i \to Y\}_{i \in I}$ 满足 $\coprod f_i : \coprod Y_i \to X$
完全分解，则称该族{\it 完全分解}。
\end{definition}

\noindent
我们先给出若干基本引理。

\begin{lemma}
\label{more-morphisms-lemma-composition-cd}
两个完全分解概形态射的复合仍完全分解。若
$\{f_i : X_i \to Y\}_{i \in I}$ 完全分解，并且对每个 $i$，族
$\{X_{ij} \to X_i\}_{j \in J_i}$ 完全分解，则族
$\{X_{ij} \to Y\}_{i \in I, j \in J_i}$ 完全分解。
\end{lemma}

\begin{proof}
从略。
\end{proof}

\begin{lemma}
\label{more-morphisms-lemma-base-change-cd}
完全分解概形态射的基变换仍完全分解。若
$\{f_i : X_i \to Y\}_{i \in I}$ 完全分解，且 $Y' \to Y$ 为概形态射，则
$\{X_i \times_Y Y' \to Y'\}_{i \in I}$ 完全分解。
\end{lemma}

\begin{proof}
设 $f : X \to Y$ 和 $g : Y' \to Y$ 为概形态射。取点 $y' \in Y'$，其在 $Y$
中的像为 $y = g(y')$。若点 $x \in X$ 满足 $f(x) = y$ 且
$\kappa(x) = \kappa(y)$，则存在唯一一点 $x' \in X' = X \times_Y Y'$，
它在 $Y'$ 中映到 $y'$、在 $X$ 中映到 $x$，并且
$\kappa(x') = \kappa(y')$，见概形章引理
\ref{schemes-lemma-points-fibre-product}。本引理由此事实容易得出；细节从略。
\end{proof}

\begin{lemma}
\label{more-morphisms-lemma-decompose}
\begin{reference}
\cite[Lemma 2.1.2]{EHIK}
\end{reference}
设 $f : X \to Y$ 为概形态射。假设 $f$ 完全分解、$f$ 局部有限呈现，并且 $Y$ 拟紧且
拟分离。则存在 $n \geq 0$ 以及态射 $Z_i \to Y$，$i = 1, \ldots, n$，
具有如下性质：
\begin{enumerate}
\item $\coprod Z_i \to Y$ 为满态射；
\item 对所有 $i$，$Z_i \to Y$ 均为浸入；
\item 对所有 $i$，$Z_i \to Y$ 均为有限呈现；并且
\item 对所有 $i$，基变换 $X \times_Y Z_i \to Z_i$ 均具有截面。
\end{enumerate}
\end{lemma}

\begin{proof}
取 $y \in Y$。由假设，存在 $Y$ 上的态射
$\sigma : \Spec(\kappa(y)) \to X$。可将 $\Spec(\kappa(y))$ 写成 $Y$ 上仿射概形
$Z$ 的有向极限，其中 $Z \to Y$ 为有限呈现浸入。事实上，选取仿射开集
$y \in \Spec(A) \subset Y$，并设 $y$ 对应于 $A$ 的素理想 $\mathfrak p$。则
$\kappa(\mathfrak p)$ 是环 $(A/I)_f$ 的滤过余极限，其中
$I \subset \mathfrak p$ 为有限生成理想，且 $f \in A$、$f \not \in \mathfrak p$。
态射 $Z = \Spec((A/I)_f) \to Y$ 是有限呈现浸入；由于 $Y$ 拟分离，
$Z \to Y$ 拟紧，见概形章引理 \ref{schemes-lemma-quasi-compact-permanence}。
由极限章命题 \ref{limits-proposition-characterize-locally-finite-presentation}，
对某个这样的 $Z$，存在 $Y$ 上的态射 $\sigma' : Z \to X$，它在
$\kappa(\mathfrak p)$ 的谱上与 $\sigma$ 一致。由于 $\sigma'$ 是 $Y$ 上的态射，
我们得到投影 $X \times_Y Z \to Z$ 的一个截面。

\medskip\noindent
由此可知，$Y$ 是若干有限呈现浸入 $Z \to Y$ 的像之并，其中
$X \times_Y Z \to Z$ 具有截面。由于 $Z \to Y$ 的像可构造（态射章引理
\ref{morphisms-lemma-chevalley}），且 $Y$ 在可构造拓扑下为紧空间（性质章引理
\ref{properties-lemma-quasi-compact-quasi-separated-spectral} 及拓扑章引理
\ref{topology-lemma-constructible-hausdorff-quasi-compact}），可知有限多个这样的像
即已足够。
\end{proof}

\begin{lemma}
\label{more-morphisms-lemma-descend-cd}
设 $S = \lim_{\lambda \in \Lambda} S_\lambda$ 为具有仿射过渡态射的概形有向系的
极限。取 $0 \in \Lambda$，并设 $f_0 : X_0 \to Y_0$ 为 $S_0$ 上的概形态射。
对 $\lambda \geq 0$，令 $f_\lambda : X_\lambda \to Y_\lambda$ 为 $f_0$ 到
$S_\lambda$ 的基变换，并令 $f : X \to Y$ 为 $f_0$ 到 $S$ 的基变换。若
\begin{enumerate}
\item $f$ 完全分解；
\item $Y_0$ 拟紧且拟分离；并且
\item $f_0$ 局部有限呈现，
\end{enumerate}
则存在 $\lambda \geq 0$，使得 $f_\lambda$ 完全分解。
\end{lemma}

\begin{proof}
由于 $Y_0$ 拟紧且拟分离，而概形 $Y$ 在 $Y_0$ 上仿射，故其拟紧且拟分离。
按引理 \ref{more-morphisms-lemma-decompose} 选取 $n \geq 0$ 以及 $Z_i \to Y$，
$i = 1, \ldots, n$。以 $\sigma_i : Z_i \to X$ 表示由 $Z_i$ 的选取而存在的
$Y$ 上态射。增大 $0 \in \Lambda$ 后，可设存在有限呈现态射
$Z_{i, 0} \to Y_0$，其到 $S$ 的基变换为 $Z_i \to Y$，见极限章引理
\ref{limits-lemma-descend-finite-presentation}。由极限章引理
\ref{limits-lemma-descend-immersion}，必要时再次增大 $0$ 后，可设
$Z_{i, 0} \to Y_0$ 为浸入。由于 $\coprod Z_i \to Y$ 为满态射，必要时增大 $0$
后，可设 $\coprod Z_{i, 0} \to Y_0$ 为满态射，见极限章引理
\ref{limits-lemma-descend-surjective}。注意
$Z_i = \lim_{\lambda \geq 0} Z_{i, \lambda}$，其中
$Z_{i, \lambda} = Y_\lambda \times_{Y_0} Z_{i, 0}$。将复合
$$
Z_i \xrightarrow{\sigma_i} X \to X_0
$$
视为 $Y_0$ 上的态射。由于 $f_0$ 局部有限呈现，由极限章命题
\ref{limits-proposition-characterize-locally-finite-presentation}，可找到
$\lambda \geq 0$，使得存在 $Y_0$ 上的态射
$\sigma'_{i, \lambda} : Z_{i, \lambda} \to X_0$，其与
$Z_i \to Z_{i, \lambda}$ 的预复合就是所显示的箭头。当然，此时
$\sigma'_{i, \lambda}$ 确定一个 $Y_\lambda$ 上的态射
$\sigma_{i, \lambda} : Z_{i, \lambda} \to
X_\lambda = X_0 \times_{Y_0} Y_\lambda$。由于
$\coprod Z_{i, \lambda} \to Y_\lambda$ 为满态射，可知
$X_\lambda \to Y_\lambda$ 完全分解。
\end{proof}




\section{丰富可逆模族}
\label{more-morphisms-section-families-ample-invertible-modules}

\noindent
我们继续态射章第 \ref{morphisms-section-families-ample-invertible-modules} 节的讨论。

\begin{lemma}
\label{more-morphisms-lemma-ample-family-ample-relative}
设 $f : X \to Y$ 为概形态射。假设
\begin{enumerate}
\item $Y$ 具有一个丰富可逆模族；
\item $X$ 上存在一个 $f$-丰富可逆模。
\end{enumerate}
则 $X$ 具有一个丰富可逆模族。
\end{lemma}

\begin{proof}
令 $\mathcal{L}$ 为 $X$ 上的一个 $f$-丰富可逆模。特别地，这蕴含 $f$ 拟紧，见
态射章定义 \ref{morphisms-definition-relatively-ample}。由态射章定义
\ref{morphisms-definition-family-ample-invertible-modules}，$Y$ 拟紧，故 $X$ 拟紧
（因而 $X$ 本身满足态射章定义
\ref{morphisms-definition-family-ample-invertible-modules} 的第一个条件）。
取 $x \in X$，其像为 $y \in Y$。由假设 (2)，可找到一个可逆
$\mathcal{O}_Y$-模 $\mathcal{N}$ 以及截面 $t \in \Gamma(Y, \mathcal{N})$，
使得 $t$ 不消失的轨迹 $Y_t$ 为仿射概形。于是 $\mathcal{L}$ 在
$f^{-1}(Y_t) = X_{f^*t}$ 上丰富，故可找到截面
$s \in \Gamma(X_{f^*t}, \mathcal{L})$，使得 $(X_{f^*t})_s$ 为仿射概形且包含
$x$。由性质章引理 \ref{properties-lemma-invert-s-sections}，对某个
$n \geq 0$，乘积 $(f^*t)^n s$ 延拓为截面
$s' \in \Gamma(X, f^*\mathcal{N}^{\otimes n} \otimes \mathcal{L})$。
最后，$f^*\mathcal{N}^{\otimes n + 1} \otimes \mathcal{L}$ 的截面
$s'' = f^* ts'$ 在 $X \setminus X_{f^*t}$ 的每一点均消失，故
$X_{s''} = (X_{f^*t})_s$ 为仿射概形，正如所需。
\end{proof}

\begin{lemma}
\label{more-morphisms-lemma-resolution-property-goes-up-affine}
设 $f : X \to Y$ 为概形的仿射或拟仿射态射。若 $Y$ 具有一个丰富可逆模族，
则 $X$ 也具有一个丰富可逆模族。
\end{lemma}

\begin{proof}
由态射章引理 \ref{morphisms-lemma-quasi-affine-O-ample}，这是引理
\ref{more-morphisms-lemma-ample-family-ample-relative} 的特殊情形。
\end{proof}






\section{吹起与丰富可逆模族}
\label{more-morphisms-section-ample-family-by-blowing-up}

\noindent
我们证明 \cite{Gross-thesis} 中的一个结果。

\begin{lemma}
\label{more-morphisms-lemma-get-ample-family}
设 $X$ 为概形。给定 $X$ 上的有效 Cartier 除子
$D_1, \ldots, D_m$ 以及可逆模 $\mathcal{L}_1, \ldots, \mathcal{L}_m$，满足
$\bigcap D_i = \emptyset$ 且 $\mathcal{L}_i|_{X \setminus D_i}$ 丰富。
则 $X$ 具有一个丰富可逆模族。
\end{lemma}

\begin{proof}
取 $x \in X$。选取指标 $i \in \{1, \ldots, m\}$，使得 $x \not \in D_i$。
令 $U_i = X \setminus D_i$。由 $\mathcal{L}_i|_{U_i}$，可找到 $n \geq 1$ 以及
截面 $s \in \Gamma(U_i, \mathcal{L}_i^{\otimes n})$，使得 $s$ 不消失的轨迹
$(U_i)_s$ 为仿射概形（性质章定义 \ref{properties-definition-ample}）。
由于 $U_i$ 是典范截面 $1 \in \mathcal{O}_X(D_i)$ 不消失的轨迹，故由性质章引理
\ref{properties-lemma-invert-s-sections}，存在 $N \geq 0$，使得 $s$ 延拓为截面
$$
s' \in \Gamma(X, \mathcal{L}_i^{\otimes n} \otimes_{\mathcal{O}_X}
\mathcal{O}_X(N D_i))
$$
将 $N$ 替换为 $N + 1$ 后，可见 $s'$ 在 $D_i$ 的每一点都消失，因而
$X_{s'} = (U_i)_s$ 为仿射概形。这证明 $X$ 具有一个丰富可逆模族，见态射章定义
\ref{morphisms-definition-family-ample-invertible-modules}。
\end{proof}

\begin{lemma}
\label{more-morphisms-lemma-blow-up-ample-family}
\begin{reference}
\cite[Proposition 1.3.1]{Gross-thesis}
\end{reference}
设 $X$ 为仅有有限多个不可约分支的拟紧拟分离概形。则存在一个拟紧稠密开子集
$U \subset X$ 以及一个 $U$-容许吹起 $X' \to X$，使得概形 $X'$ 具有一个丰富
可逆模族。
\end{lemma}

\begin{proof}
令 $\eta_1, \ldots, \eta_n \in X$ 为 $X$ 的各不可约分支的泛点。由性质章引理
\ref{properties-lemma-point-and-maximal-points-affine} 以及 $X$ 拟紧这一事实，
可找到有限仿射开覆盖 $X = U_1 \cup \ldots \cup U_m$，使得每个 $U_i$ 都包含
$\eta_1, \ldots, \eta_n$。特别地，拟紧开子集
$U = U_1 \cap \ldots \cap U_m$ 在 $X$ 中稠密。令
$\mathcal{I}_i \subset \mathcal{O}_X$ 为有限型拟凝聚理想层，使得
$U_i = X \setminus Z_i$，其中 $Z_i = V(\mathcal{I}_i)$，见性质章引理
\ref{properties-lemma-quasi-coherent-finite-type-ideals}。令
$$
f : X' \longrightarrow X
$$
为 $X$ 沿理想层 $\mathcal{I} = \mathcal{I}_1 \cdots \mathcal{I}_m$ 的吹起。
注意，$f$ 是 $U$-容许吹起，因为 $V(\mathcal{I})$ 在集合论意义下是与 $U$
不交的各 $Z_i$ 之并。此外，$f$ 为射影态射，且 $\mathcal{O}_{X'}(1)$ 是
$f$-相对丰富的，见除子章引理 \ref{divisors-lemma-blowing-up-projective}。
由除子章引理 \ref{divisors-lemma-blowing-up-two-ideals}，对每个 $i$，态射
$f'$ 分解为 $X' \to X'_i \to X$，其中 $X'_i \to X$ 是沿 $\mathcal{I}_i$
的吹起，而 $X' \to X'_i$ 是另一个吹起（即沿略去 $\mathcal{I}_i$ 后诸理想
$\mathcal{I}_j$ 之积的拉回所作的吹起）。由此可知
$D_i = f^{-1}(Z_i) \subset X'$ 是有效 Cartier 除子，见除子章引理
\ref{divisors-lemma-blow-up-pullback-effective-Cartier} 和
\ref{divisors-lemma-blowing-up-gives-effective-Cartier-divisor}。我们有
$X' \setminus D_i = f^{-1}(U_i)$。由于 $\mathcal{O}_{X'}(1)$ 为 $f$-丰富，
$\mathcal{O}_{X'}(1)$ 在 $X' \setminus D_i$ 上的限制丰富。由引理 \ref{more-morphisms-lemma-get-ample-family}，
$X'$ 具有一个丰富可逆模族。
\end{proof}

\begin{proposition}
\label{more-morphisms-proposition-envelope-with-resolution-property}
设 $X$ 为拟紧拟分离概形。则存在一个有限呈现、固有且完全分解（定义
\ref{more-morphisms-definition-cd-morphism}）的态射 $f : Y \to X$，使得概形 $Y$ 具有一个
丰富可逆模族。
\end{proposition}

\begin{proof}
由极限章命题 \ref{limits-proposition-approximate}，存在仿射态射 $X \to X_0$，
其中 $X_0$ 为 $\mathbf{Z}$ 上有限型概形。下面构造一个具有所有所需性质的态射
$Y_0 \to X_0$。令 $Y = X \times_{X_0} Y_0$，并令 $f$ 等于投影
$f : Y \to X$，即可得出结论。事实上，$f$ 为有限呈现（态射章引理
\ref{morphisms-lemma-base-change-finite-presentation}），$f$ 为固有态射
（态射章引理 \ref{morphisms-lemma-base-change-proper}），且 $f$ 完全分解
（引理 \ref{more-morphisms-lemma-base-change-cd}）。最后，由于 $Y \to Y_0$ 为仿射态射
（它是 $X \to X_0$ 的基变换），由引理
\ref{more-morphisms-lemma-resolution-property-goes-up-affine}，$Y$ 具有一个丰富可逆模族。
这将问题约化到下一段讨论的情形。

\medskip\noindent
设 $X$ 在 $\mathbf{Z}$ 上有限型。特别地，$\dim(X) < \infty$。我们对
$\dim(X)$ 作归纳。若 $\dim(X) = 0$，则 $X$ 为仿射概形并具有消解性质。
一般地，由引理 \ref{more-morphisms-lemma-blow-up-ample-family}，存在稠密开子集 $U \subset X$
以及一个 $U$-容许吹起 $X' \to X$，使得 $X'$ 具有一个丰富可逆模族。由于
$f : X' \to X$ 在 $U$ 上为同构，$U$ 的每一点都可提升为 $X'$ 中具有相同剩余域
的一点。令 $Z = X \setminus U$，赋予既约诱导概形结构。由于 $U$ 在 $X$ 中稠密
（见上文），有 $\dim(Z) < \dim(X)$。由归纳假设，可找到一个固有、完全分解的态射
$W \to Z$，使得 $W$ 具有一个丰富可逆模族。于是
$Y = X' \amalg W \to X$ 即为所需态射。
\end{proof}




\section{闭浸入的广延判据}
\label{more-morphisms-section-extensive-criterion-closed-immersions}

\noindent
本节给出概形态射成为闭浸入的一个判据。

\begin{lemma}
\label{more-morphisms-lemma-affine-extensive-criterion}
仿射概形态射 $f : X \to Y$ 为闭浸入，当且仅当对每个单射环同态 $A \to B$
以及交换方块
$$
\xymatrix{
\Spec(B) \ar[d] \ar[r] & X \ar[d]^f \\
\Spec(A) \ar[r] \ar@{..>}[ur] & Y
}
$$
都存在提升 $\Spec(A) \to X$，使两个三角形交换。
\end{lemma}

\begin{proof}
设态射 $f$ 由环同态 $\phi : R \to S$ 给出。则 $f$ 为闭浸入，当且仅当
$\phi$ 为满射。

\medskip\noindent
首先假设 $\phi$ 为满射。设 $\psi : A \to B$ 为单射环同态，并且给定交换图表
$$
\xymatrix{
R \ar[r]^\alpha \ar[d]^\phi & A \ar[d]^\psi \\
S \ar[r]^\beta \ar@{..>}[ur] & B
}
$$
定义提升 $S \to A$，令 $s \mapsto \alpha(r)$，其中 $r \in R$ 满足
$\phi(r) = s$。由于 $\psi$ 为单射且方块交换，这一定义适定。取环的谱给出交换环
范畴与仿射概形范畴之间的反等价，故 $f$ 具有所需的提升性质。

\medskip\noindent
其次，假设 $\phi$ 对所有单射环同态 $\psi: A \to B$ 均具有提升。注意
$\phi(R)$ 是 $S$ 的子环，故得到交换方块
$$
\xymatrix{
R \ar[r] \ar[d]^\phi  & \phi(R) \ar[d] \\
S \ar@{=}[r] \ar@{..>}[ur] & S
}
$$
其中存在提升 $S \to \phi(R)$。因此包含映射 $\phi(R) \to S$ 必为同构，
这说明 $\phi$ 为满射，证明完成。
\end{proof}

\begin{lemma}
\label{more-morphisms-lemma-scheme-affine-if-aff-is-split-mono}
设 $X$ 为概形。若概形章引理 \ref{schemes-lemma-morphism-into-affine} 的典范态射
$X \to \Spec(\Gamma(X, \mathcal{O}_X))$ 具有收缩，则 $X$ 为仿射概形。
\end{lemma}

\begin{proof}
记 $S = \Spec(\Gamma(X, \mathcal{O}_X))$，并以 $f : X \to S$ 表示引理中给出的
态射。设 $s : S \to X$ 为一个收缩；于是 $\text{id}_X = sf$。从而
$f s f = \text{id}_S f$。由于 $f$ 诱导同构
$\Gamma(S, \mathcal{O}_S) \to \Gamma(X, \mathcal{O}_X)$，这说明 $fs$ 与
$\text{id}_S$ 在 $\Gamma(S, \mathcal{O}_S)$ 上诱导相同的映射。由于 $S$
为仿射概形，遂有 $f s = \text{id}_S$。因此 $f$ 为同构，$X$ 为仿射概形，
正如所要证明。
\end{proof}

\begin{lemma}
\label{more-morphisms-lemma-aff-is-injective-on-sheaves}
设 $X$ 为概形，并令 $f : X \to S = \Spec(\Gamma(X, \mathcal{O}_X))$ 为概形章
引理 \ref{schemes-lemma-morphism-into-affine} 的典范态射。包含在
$f^\sharp : \mathcal{O}_S \to f_*\mathcal{O}_X$ 的核中的最大拟凝聚
$\mathcal{O}_S$-模为零。若 $X$ 拟紧，则 $f^\sharp$ 为单射。特别地，若 $X$
拟紧，则 $f$ 为支配态射。
\end{lemma}

\begin{proof}
令 $M \subset \Gamma(S, \mathcal{O}_S)$ 为与包含在 $f^\sharp$ 的核中的最大
拟凝聚 $\mathcal{O}_S$-模相对应的子模。则任意元素 $a \in M$ 均被
$f^\sharp$ 映到零。然而，$f^\sharp(a)$ 是下列模中的元素
$$
\Gamma(S, f_*\mathcal{O}_X) = \Gamma(X, \mathcal{O}_X) =
\Gamma(S, \mathcal{O}_S)
$$
且它对应于 $a$ 本身！故 $a = 0$。于是 $M = 0$，这证明了第一个断言。
注意，这等价于态射 $f : X \to S$ 在概形论意义下为满态射。

\medskip\noindent
若 $X$ 拟紧，则由态射章引理
\ref{morphisms-lemma-quasi-compact-scheme-theoretic-image}，
$\Ker(f^\sharp)$ 拟凝聚。因此 $\Ker(f^\sharp) = 0$，且 $f^\sharp$ 为单射。
在此情形下，由态射章引理
\ref{morphisms-lemma-quasi-compact-scheme-theoretic-image} 的第 (4) 部分，
$f$ 为支配态射。
\end{proof}

\begin{lemma}
\label{more-morphisms-lemma-extensive-criterion}
设 $f: X \to Y$ 为概形的拟紧态射。则 $f$ 为闭浸入，当且仅当对每个单射环同态
$A \to B$ 以及交换方块
$$
\xymatrix{
\Spec(B) \ar[d] \ar[r] & X \ar[d]^f \\
\Spec(A) \ar[r] \ar@{..>}[ur] & Y
}
$$
都存在提升 $\Spec A \to X$，使图表交换。
\end{lemma}

\begin{proof}
假设 $f$ 为闭浸入。设 $A \to B$ 为单射环同态，并考虑交换方块
$$
\xymatrix{
\Spec(B) \ar[d] \ar[r]  & X \ar[d]^f \\
\Spec(A) \ar[r] \ar@{..>}[ur] & Y
}
$$
则 $\Spec(A) \times_Y X \to \Spec(A)$ 为闭浸入，故得到理想 $I \subset A$
以及交换图表
$$
\xymatrix{
\Spec(B) \ar[d] \ar[r] & \Spec(A/I) \ar[r] \ar[d] & X \ar[d]^f \\
\Spec(A) \ar[r] \ar@{..>}[ur] & \Spec(A) \ar[r] & Y
}
$$
由引理 \ref{more-morphisms-lemma-affine-extensive-criterion}，得到一个提升。

\medskip\noindent
假设 $f$ 具有引理中所述的提升性质。证明 $f$ 为闭浸入这一问题在 $Y$ 上是局部的，
故可并且确实设 $Y$ 为仿射概形。特别地，$Y$ 拟紧，从而 $X$ 拟紧。因此存在有限
仿射开覆盖 $X = U_1 \cup \ldots \cup U_n$。态射
$$
\pi : U = \coprod U_i \longrightarrow X
$$
的源为仿射概形，并且所诱导的环同态
$\Gamma(X, \mathcal{O}_X) \to \Gamma(U, \mathcal{O}_U)$ 为单射。由假设，
下图中存在提升：
$$
\xymatrix{
U \ar[r]^\pi \ar[d] & X \ar[d]^f \\
\Spec(\Gamma(X, \mathcal{O}_X)) \ar[r]^-{f'} \ar@{..>}[ur]^h & Y
}
$$
其中 $f'$ 是与环同态
$\Gamma(Y, \mathcal{O}_Y) \to \Gamma(X, \mathcal{O}_X)$ 相对应的仿射概形态射。
由 $\pi$ 为满射可知，态射 $h$ 是典范态射
$X \to \Spec(\Gamma(X, \mathcal{O}_X))$ 的一个收缩；细节从略。因此，由引理
\ref{more-morphisms-lemma-scheme-affine-if-aff-is-split-mono}，$X$ 为仿射概形。再由引理
\ref{more-morphisms-lemma-affine-extensive-criterion}，可知 $f$ 为闭浸入。
\end{proof}
